% The following TeXfile comprises part of

%     MEASURE THEORY, by D.H.Fremlin
%          (University of Essex, Colchester, England;
%                   e-mail  fremdh@essex.ac.uk)

% Please contact me with any difficulties not covered by the remarks
% below.

% The file is set up for processing by plain TeX in the presence of
% files  amstex.tex ,  amssym.tex  and  amsppt.sty  suitably supported by the
% appropriate fonts.   If these are unavailable, a readable version
% can be obtained by removing the comment symbol  %  from the line

%        %   \amstexsupportedfalse

% below.   Some systems do not allow access to files  amstex.tex  and
% amsppt.sty , but instead employ a command "amstex <filename.tex>";
% for such systems it is sometimes right to delete or disable the
% commands

%        \input amstex
%        \documentstyle{amsppt}

% (I have little experience with this and should welcome comments.)

% Various formats of output are obtainable.   The principal variations
% are

%    1. "Standard" output uses the command
%               \smallprint
% and gives output adapted to A4 paper (21 cm x 29.7 cm).

%    2.  Using the command
%               \USsmallprint
% gives output adapted to 8.5" x 11" .

%    3.  Using the command
%               \noproofs
% produces an abridged version, in which proofs are omitted.

%    4.   Using the command
%               \resultsonly
% produces a super-abridged version, omitting comments and exercises
% as well as proofs.

% These commands are grouped together at the end of the introductory
% definitions;  they may be activated by deleting the initial comment
% symbols  % .

% Any diagrams included will be called for by inputting postscript files.
% (You can look for these by hunting for   \picture  in the code below.)
% To get these you will need the relevant  .ps  files with the expected
% names  ( \picture{mt133g}{...}  demands the file  mt133g.ps , for instance).


%(If you can read TeX fluently, you may wish to fast-forward to a line
%beginning   %file starts here   or   \frfilename
%if one of these is present.)

%TeXfiles with this heading are designed for processing with AmSTeX,
%accessed through TeXfiles  amstex.tex (with amssym.tex  and
% amsppt.sty  and possibly  amsppt.sti)  during compilation.   
%If these are unavailable, a readable
%version can (I hope) be produced by deleting the  %  before
% \amstexsupportedfalse  below.

%Some special difficulties arise with the importation of  .ps  files
%in diagrams.   Look for  \picture  below for notes on these.

%The following \count registers are used for special purposes:
%   \count1=\stylenumber  to specify style
%   \count2=\footnotenumber  to count footnotes

%AmSTeX declaration
\newif\ifamstexsupported
     \amstexsupportedtrue
%     \amstexsupportedfalse               %use when AMSTeX unavailable
\newif\ifDHFamstexloaded        %for detection of a non-standard AMSTeX
     \DHFamstexloadedfalse    
\newif\ifDHFamspptloaded        
     \DHFamspptloadedfalse    
\newif\iffourteensavailable
     \fourteensavailabletrue
    \fourteensavailablefalse
\newif\ifathome
     \athometrue     %to cope with a quirk     of my son's Ghostscript
     \athomefalse
\newif\iffirstpage
     \firstpagetrue

%Style assignments
\countdef\stylenumber=1
\countdef\footnotenumber=2
\def\fullsize{\magnification=\magstep1
     \def\pagewidth{390pt}
     \def\pageheight{9truein}
     \hsize\pagewidth
     \vsize\pageheight
     \stylenumber=1}
\def\smallprint{\magnification=\magstep0
     \overfullrule0pt
     \def\pagewidth{468pt}
     \def\pageheight{9truein}
     \hsize\pagewidth
     \vsize\pageheight
     \stylenumber=2}
\def\printersize{\magnification=\magstep1
     \baselineskip24pt
     \overfullrule0pt
     \def\pagewidth{330pt}
     \def\pageheight{8.5truein}
     \hsize\pagewidth
     \vsize\pageheight
     \stylenumber=3}
\def\refsize{\magnification=\magstep1
     \overfullrule0pt
     \def\pagewidth{390pt}
     \def\pageheight{9truein}
     \hsize\pagewidth
     \vsize\pageheight
     \stylenumber=4}
\def\rrversion{\magnification=\magstep1
     \overfullrule0pt
     \def\pagewidth{330pt}
     \def\pageheight{9truein}
     \hsize\pagewidth
     \vsize\pageheight
     \stylenumber=5}
\def\bigsize{\magnification=\magstep2
     \overfullrule0pt
     \def\pagewidth{325pt}
     \def\pageheight{9truein}
     \hsize\pagewidth
     \vsize\pageheight
     \stylenumber=6}
\def\bbigsize{\magnification=\magstep3
     \overfullrule0pt
     \def\pagewidth{271pt}
     \def\pageheight{9truein}
     \hsize\pagewidth
     \vsize\pageheight
     \stylenumber=7}
\def\USsmallprint{\magnification=\magstep0
     \overfullrule0pt
     \def\pagewidth{482pt}
     \def\pageheight{8.7truein}
     \hsize\pagewidth
     \vsize\pageheight
     \stylenumber=8}
\def\halfmagsize{\magnification=\magstephalf
     \def\pagewidth{427pt}
     \def\pageheight{9truein}
     \hsize\pagewidth
     \vsize\pageheight
     \stylenumber=9}
           % register  stylenumber  is used to control the output format

%addresses, letterheads
\def\deptaddress{\hskip 3truein Department of Mathematical Sciences\par
   \noindent e-mail: {\tt fremdh}\hskip 2.9truein\today\par
   \medskip}
\long\def\univaddress#1{\hskip 3truein Department of Mathematical Sciences\par
   \hskip 3truein University of Essex\par
   \hskip 3truein Colchester CO4 3SQ\par
   \hskip 3truein England\par
   %\hskip 3truein van Vleck Hall\par
   %\hskip 3truein 480 Lincoln Drive\par
   %\hskip 3truein Madison\par
   %\hskip 3truein WI 53706-1388, U.S.A.\par
   \hskip 3.2truein#1\par
   \noindent e-mail: {\tt david\@fremlin.org}
       %{\tt fremlin\@math.wisc.edu}
   \bigskip}
\long\def\homeaddress#1{\hskip 3truein 25 Ireton Road\par
   \hskip 3truein Colchester CO3 3AT\par
   \hskip 3truein England\par
   \hskip 3.2truein#1\par
   \noindent e-mail: {\tt david\@fremlin.org}
   \bigskip}
\long\def\letterhead#1#2#3{\univaddress{#1}
   {\parindent=0pt
    #2
   }
   \medskip

   \noindent Dear #3,
     }
\def\shortaddress{\noindent Departmentof Mathematical Sciences, University of Essex,
   Colchester CO4 3SQ, England\par
   \noindent e-mail: david\@fremlin.org}
\def\valediction#1#2{\bigskip
   \hskip 3truein #1,
   \vskip 1truein
   \hskip 3.3truein #2
   \bigskip}

%picture importation
   %It is clear that there is no general standard.   The syntax I
   %use when calling for a  .ps  file is
   %    \def\Caption{<...>}
   %    \picture{<filename>}{<dim>}
   %which is supposed to import  <filename>.ps  at height  <dim> ,
   %centered, and adding a one-line centered caption  <...> .
   %The basic macro listed in  \def\picture  is that which
   %used to work at Essex.   I now generally use the macro in
   % \def\picturewithpsfig , which requires a previous input
   % psfig.tex  or  psfig.sty ;
   %if you interpolate  \atUEssex  early in your file it will
   %give you the latter.
   %   I find that  psfig  is erratic in its centering, especially
   %under magnification, so I have written a routine
   % \centerpicturewithpsfig  (which may be called from
   % \centerpicture ) with the syntax
   %    \centerpicture{<filename>}{<dim>}{<dim>}
   %the first dimension being width and the second height.

   % A second package is called for by  \usegraphicx .   This requires
   % miniltx.tex  and  graphicx.sty ;  miniltx.tex  is currently (April 2016)
   % available from
   % http://mirror.ox.ac.uk/sites/ctan.org/macros/plain/graphics/miniltx.tex
   % The advantage of this is that it seems to allow manipulation of the
   % dvi>pdf stage using  dvipdfmx .

\def\picture#1#2{\bigskip
  \bigskip
  \immediate\write0{importing diagram #1.ps at height #2
         with dvitops}
  \vbox{\vskip #2
    \special{dvitops: import #1.ps \pagewidth #2}
    \Centerline{\Caption}
    }
   \medskip
   \def\Caption{}}
%an alternative version, without caption
\def\centerpicture#1#2#3{\bigskip
  \bigskip
  \immediate\write0{importing diagram #1.ps at width #2
         and height #3 with dvitops}
  \vbox{\vskip #3
    \special{dvitops: import #1.ps \pagewidth #3}
  }}
%to put pictures on same line
\def\sideshiftedpicture#1#2#3#4{\bigskip
   \bigskip
   \immediate\write0{importing diagram #1.ps at height #4
        and width #3 with dvitops}
   \vbox{\special{dvitops:  import #1.ps #3 #4}}
   \hskip #2 plus 0pt minus 0pt
   }
\def\startsideshiftedpicture{}
\newcount\scaleup
\newcount\scaledown
\scaleup=1
\scaledown=1
\def\scalesideshiftedpicture#1#2#3#4{
  \dimen4=#4
  \dimen3=#3
  \multiply\dimen4 by\scaleup
  \divide\dimen4 by \scaledown
  \multiply\dimen3 by\scaleup
  \divide\dimen3 by \scaledown
  \sideshiftedpicture{#1}{#2}{\dimen3}{\dimen4}}

\def\picturewithpsfig#1#2{\bigskip
   \immediate\write0{importing diagram #1.ps at height #2 with psfig}
   \vbox{\hsize\pagewidth
   \centerline{\psfig{figure=#1.ps,height=#2}}
   \medskip
   \centerline{\Caption}}
   \medskip
   \def\Caption{}}
\def\centerpicturewithpsfig#1#2#3{
  \immediate\write0{importing diagram #1.ps at width #2
         and height #3 with psfig}
  \dimen 3 = #2
  \dimen 4 = #3
  \dimen 2 = \pagewidth
  \advance\dimen2 by -\dimen3
  \divide\dimen2 by 2
  \discrversionB{\shrinkonestep}{}
  \discrversionH{\shrinkonestep}{}
  \discrversionF{\enlargeonestep}{}
  \discrversionG{\enlargeonestep\enlargeonestep}{}
  \dimen 6 = \dimen4
  \divide\dimen6 by 3
  \vskip\dimen6
  \ifdim\dimen2>0\hskip\dimen2\fi
     \vbox{\psfig{file=#1.ps,width=\dimen3,height=\dimen4}}
  \discrversionG{\vskip-\dimen6}{}
  }
\def\enlargeonestep{\dimen5 = 5truein
   \dimen6 = 7truein
   \ifdim\dimen3<\dimen5
      \ifdim\dimen4<\dimen6
         \divide\dimen3 by 5
         \multiply\dimen3 by 6
         \divide\dimen4 by 5
         \multiply\dimen4 by 6
      \fi\fi}
\def\shrinkonestep{\divide\dimen3 by 6
         \multiply\dimen3 by 5
         \divide\dimen4 by 6
         \multiply\dimen4 by 5}
\def\sideshiftedpicturewithpsfig#1#2#3#4{
   \immediate\write0{importing diagram #1.ps at height #4
      and width #3 with psfig}
   \dimen 4 = #4
   \dimen 3 = #3
   \discrversionB{\shrinkonestep}{}
   \discrversionH{\shrinkonestep}{}
   \discrversionF{\enlargeonestep}{}
   \discrversionG{\enlargeonestep\enlargeonestep}{}
%   \vbox{\psfig{figure=#1.ps,height=\dimen4,width=\dimen3}}
   \psfig{figure=#1.ps,height=\dimen4,width=\dimen3}
   \hskip #2 plus 0pt minus 0pt}

\def\picturewithgraphicx#1#2{\bigskip
   \immediate\write0{importing diagram #1.ps at height #2 with graphicx}
   \vbox{\hsize\pagewidth
   \centerline{\includegraphics[height=#2]{#1.ps}}
   \medskip
   \centerline{\Caption}}
   \medskip
   \def\Caption{}}
\def\centerpicturewithgraphicx#1#2#3{
  \immediate\write0{importing diagram #1.ps at width #2
         and height #3 with graphicx}
  \dimen 3 = #2
  \dimen 4 = #3
  \dimen 2 = \pagewidth
  \advance\dimen2 by -\dimen3
  \divide\dimen2 by 2
  \discrversionB{\shrinkonestep}{}
  \discrversionH{\shrinkonestep}{}
  \discrversionF{\enlargeonestep}{}
  \discrversionG{\enlargeonestep\enlargeonestep}{}
  \dimen 6 = \dimen4
  \divide\dimen6 by 3
  \vskip\dimen6
  \ifdim\dimen2>0\hskip\dimen2\fi
     \vbox{\includegraphics[width=\dimen3,height=\dimen4]{#1.ps}}
  \discrversionG{\vskip-\dimen6}{}
  }
\def\enlargeonestep{\dimen5 = 5truein
   \dimen6 = 7truein
   \ifdim\dimen3<\dimen5
      \ifdim\dimen4<\dimen6
         \divide\dimen3 by 5
         \multiply\dimen3 by 6
         \divide\dimen4 by 5
         \multiply\dimen4 by 6
      \fi\fi}
\def\shrinkonestep{\divide\dimen3 by 6
         \multiply\dimen3 by 5
         \divide\dimen4 by 6
         \multiply\dimen4 by 5}
\def\sideshiftedpicturewithgraphicx#1#2#3#4{
   \immediate\write0{importing diagram #1.ps at height #4
      and width #3 with graphicx}
   \dimen 4 = #4
   \dimen 3 = #3
   \discrversionB{\shrinkonestep}{}
   \discrversionH{\shrinkonestep}{}
   \discrversionF{\enlargeonestep}{}
   \discrversionG{\enlargeonestep\enlargeonestep}{}
%   \vbox{\includegraphics[height=\dimen4,width=\dimen3]{#1.ps}}
   \includegraphics[height=\dimen4,width=\dimen3]{#1.ps}
   \hskip #2 plus 0pt minus 0pt}

%alternative picture importation code;  I've been told about this
  %but have not used it successfully myself
%\def\Caption{}
%\input epsf
%\epsfverbosetrue
%\def\picture#1#2{\bigskip
%    \bigskip
%    \immediate\write0{importing diagram #1.ps at height #2 with epsf}
%    \epsfysize=#2
%    \centerline{\epsffile{#1.ps}}
%    \medskip
%    \Centerline{\Caption}
%    \medskip
%    }

%in addition it may be necessary to amend the imported  .ps  file
%by replacing the   %%BoundingBox   header by  %%BoundingBox:

\def\atUEssex{
   \input psfig.sty
   \def\picture##1##2{\picturewithpsfig{##1}{##2}}
   \def\sideshiftedpicture##1##2##3##4{
     \sideshiftedpicturewithpsfig{##1}{##2}{##3}{##4}}
   \def\centerpicture##1##2##3{
       \centerpicturewithpsfig{##1}{##2}{##3}}
   \def\startsideshiftedpicture{
       \sideshiftedpicture{empty}{1pt}{1pt}{1pt}}
   }
\def\usegraphicx{
   \input miniltx
   \def\Gin@driver{dvips.def}
   \input graphicx.sty
   \def\picture##1##2{\picturewithgraphicx{##1}{##2}}
   \def\sideshiftedpicture##1##2##3##4{
     \sideshiftedpicturewithgraphicx{##1}{##2}{##3}{##4}}
   \def\centerpicture##1##2##3{
       \centerpicturewithgraphicx{##1}{##2}{##3}}
   \def\startsideshiftedpicture{
       \sideshiftedpicture{empty}{1pt}{1pt}{1pt}}
}


%special fonts
\def\biggerscriptfonts{\scriptfont0=\eightrm
     \scriptfont1=\eighti
     \scriptfont2=\eightsy
     \scriptfont3=\eightex
     \scriptfont\bffam=\eightbf
     \scriptfont\msbfam=\eightmsb
     \scriptfont\eufmfam=\eighteufm
     }
\def\Loadcyrillic{\ifamstexsupported
      \input cyracc.def
      \font\tencyr=wncyr10
      \def\cyr{\tencyr\cyracc} \def\itcyr{\tenitcyr\cyracc}
    \else\def\cyr{\smc}\def\cdprime{u\i}\def\Cdprime{UI}\fi}
%http://wkweb3.cableinet.co.uk/linux/LDP/HOWTO/Cyrillic-HOWTO-7.html
%cyrillic.tex

\def\Loadfourteens{\iffourteensavailable
       \font\fourteenrm=cmr14
       \font\fourteenbf=cmbx14
       \font\fourteenit=cmsl14
   \else\font\fourteenrm=cmr10 scaled \magstep2
      \font\fourteenbf=cmbx10 scaled \magstep2
      \font\fourteenit=cmsl10 scaled \magstep2
   \fi}
\def\smallerfonts{\eightpoint
   \textfont\msbfam=\eightmsb
   \textfont\eufmfam=\eighteufm
   \textfont\bffam=\eightbf
   \textfont\eurmfam=\eighteurm
   \textfont\eusmfam=\eighteusm
   }
\font\twelveex=cmex10 scaled \magstep1


%Copyright declarations
\newcount\oldcopyrightdate
\newcount\newcopyrightdate
\oldcopyrightdate=0
\def\copyrightdate#1{\newcopyrightdate=#1
       \ifnum\newcopyrightdate=\oldcopyrightdate\else
       \oldcopyrightdate=#1
       \oldfootnote{}{\eightrm \copyright\ #1 D.~H.~Fremlin}
       \immediate\write0{Copyright #1}
       \fi}

%hyphenation
\hyphenation{equi-ver-id-ical}

%miscellaneous definitions
\def\Aut{\mathop{\text{Aut}}}
\def\backstep#1{\hskip -0.#1em plus 0 pt minus 0 pt}
\def\Bang{\vthsp\hbox{\bf X\hskip0.07em\llap X\hskip0.07em\llap X\rm}}
\def\BanG{\Bang{\hskip 1.111pt plus 3.333pt minus 0pt}}
%\def\BbbN{\mathchoice{\hbox{$\Bbb N$\hskip0.02em}}
%  {\hbox{$\Bbb N$\hskip0.02em}}
%  {\hbox{$\scriptstyle\Bbb N$\hskip0.04em}}
%  {\hbox{$\scriptscriptstyle\Bbb N$\hskip0.04em}}}
%\def\BbbQ{\mathchoice{\hbox{$\Bbb Q$\hskip0.02em}}
%  {\hbox{$\Bbb Q$\hskip0.02em}}
%  {\hbox{$\scriptstyle\Bbb Q$\hskip0.04em}}
%  {\hbox{$\scriptscriptstyle\Bbb Q$\hskip0.04em}}}
%\def\BbbR{\mathchoice{\hbox{$\Bbb R$\hskip0.02em}}
%  {\hbox{$\Bbb R$\hskip0.02em}}
%  {\hbox{$\scriptstyle\Bbb R$\hskip0.04em}}
%  {\hbox{$\scriptscriptstyle\Bbb R$\hskip0.04em}}}
%  %do not use with $\Bbb R^X$
%\def\BbbZ{\mathchoice{\hbox{$\Bbb Z$\hskip0.02em}}
%  {\hbox{$\Bbb Z$\hskip0.02em}}
%  {\hbox{$\scriptstyle\Bbb Z$\hskip0.04em}}
%  {\hbox{$\scriptscriptstyle\Bbb Z$\hskip0.04em}}}
\def\BbbN{\openBbb{N}}
\def\BbbQ{\openBbb{Q}}
\def\BbbR{\openBbb{R}}
  %do not use with $\Bbb R^X$  do not use in $_{\Bbb R^r}$
\def\BbbZ{\openBbb{Z}}
\def\Bbbone{\hbox{1\hskip0.13em\llap{1}}}
\def\biggerint{\textfont3=\twelveex\intop\nolimits}
  % you can get this size into a displayed formula with
  %  $$\hbox{$\biggerint_{\Cal S}\pmb{u}\,d\pmb{v}$}$$
  % but it may then seem cramped
  % or try \hbox{$\biggerint_{\Cal S}\,$}
\def\biggeriint{\textfont3=\twelveex\iint}
  % remarkable fact:  this seems to infect all \int within a given
  % $..$ ; but  $\biggerint$$\int$ seems to leave the $\int$ normal
  % really  \textfont3=\twelveex  is more reliable, since  \biggerint
  % as written gets turned off rather easily (e.g. by  \comment,
  % or by an extra  {...})
\def\bover#1#2{{{#1}\over{#2}}}
\def\Bover#1#2{\backstep4\hbox{
     \biggerscriptfonts${\shortstrut#1}\over{\shortstrut#2}$}}
  % when there are superscripts in the denominator, as in \Bover1{2^m}
  % or  \Bover1{\sqrt{2\pi}^r} , these can sometimes come out with better
  % horizontal alignment if you use  \Bover1{{\sshortstrut 2}^m}  or
  % \Bover1{{\sqrt{2\pi}}^r}
\def\bsover#1#2{\textstyle\bover{#1}{#2}}
\def\bsp{/\penalty-100}
  % \penalty seems to behave oddly in a \tt string in \fullsize mode
\def\Caption{}
\def\Center#1{

   \Centerline{#1}

   \noindent}
\def\Centerline#1{\vskip 5pt
     {\def\int{\text{}$$\textfont3=\twelveex\intop\nolimits}
      \def\dashint{\text{}$$\textfont3=\twelveex
         \mskip2mu\vrule height2.65pt width4.5pt depth-2.35pt\mskip-11mu
	 \intop\nolimits\mskip-2mu}
      \def\idashint{\text{}$$\int\mskip-7mu\dashint}	 
      \def\dashiint{\text{}$$\dashint\mskip-6mu\int}	 
      \def\dashidashint{\text{}$$\dashint\mskip-6mu\dashint}	 
      \def\underlineint{
           \text{}$$\textfont3=\twelveex\underline{\intop\nolimits}}
      \def\overlineint{
           \text{}$$\textfont3=\twelveex\overline{\intop\nolimits}}
      \def\iint{\text{}$$\textfont3=\twelveex
           \intop\nolimits\hskip-0.5em\intop\nolimits}
      \def\sum{\text{}$$\sumop}
      \def\prod{\text{}$$\prodop}
      \def\bigcap{\text{}$$\bigcapop}
      \def\bigcup{\text{}$$\bigcupop}
   \centerline{#1}}\vskip 5pt}
   \mathchardef\bigcapop="1354
   \mathchardef\bigcupop="1353
   \mathchardef\prodop="1351
   \mathchardef\sumop="1350
   %in  \Centerline{$\overline{}$, $\cmmnt{}$} must use \sumop etc
\def\cf{\mathop{\text{cf}}}
\def\circumflex#1{{\accent"5E #1}}
\def\closeplus{\hbox{+}}
\long\def\cmmnt#1{\ifresultsonly\else{#1}\fi}
\def\coint#1{\left[#1\right[}
\def\curl{\mathop{\text{curl}}}
\def\dashint{\mathchoice{-\mskip-20mu\intop\nolimits}
   {{\vrule height2.75pt width4.5pt depth-2.45pt\mskip-10mu\intop\nolimits}}
      {-\mskip-16.5mu\intop\nolimits} %not yet used
      {-\mskip-16.5mu\intop\nolimits}} %not yet used
\def\idashint{\mathchoice{\int\mskip-6mu\dashint}
      {\int\mskip-5mu\dashint}
      {\int\mskip-5mu\dashint} %not yet used
      {\int\mskip-5mu\dashint}} %not yet used
\def\dashiint{\mathchoice{\dashint\mskip-10mu\int}
      {\dashint\mskip-6mu\int}
      {\dashint\mskip-6mu\int} %not yet used
      {\dashint\mskip-6mu\int}} %not yet used
\def\dashidashint{\mathchoice{\dashint\mskip-10mu\dashint}
      {\dashint\mskip-6mu\dashint}
      {\dashint\mskip-6mu\dashint} %not yet used
      {\dashint\mskip-6mu\dashint}} %not yet used
\def\dashidashiint{\mathchoice{\dashint\mskip-8mu\dashint\mskip-10mu\int}
      {\dashint\mskip-4mu\dashint\mskip-8mu\int} 
      {\dashint\mskip-6mu\int} %not yet used
      {\dashint\mskip-6mu\int}} %not yet used
\def\dashiidashint{\mathchoice{\dashint\mskip-10mu\int\mskip-8mu\dashint}
      {\dashint\mskip-4mu\dashint\mskip-8mu\int} 
      {\dashint\mskip-6mu\int} %not yet used
      {\dashint\mskip-6mu\int}} %not yet used
\def\idashidashint{\mathchoice{\int\mskip-8mu\dashint\mskip-8mu\dashint}
      {\dashint\mskip-6mu\int} %not yet used
      {\dashint\mskip-6mu\int} %not yet used
      {\dashint\mskip-6mu\int}} %not yet used
\def\diam{\mathop{\text{diam}}}
\def\discrpage{\ifnum\stylenumber=1{\frnewpage}  %redefined in mt.tex
           \else\ifnum\stylenumber=3{\frnewpage}
           \else\ifnum\stylenumber=5{\frnewpage}
           \else\bigskip\fi\fi\fi}
\long\def\discrversionA#1#2{\ifnum\stylenumber=1#1\else#2\fi}
\long\def\discrversionB#1#2{\ifnum\stylenumber=2#1\else#2\fi}
\long\def\discrversionC#1#2{\ifnum\stylenumber=3#1\else#2\fi}
\long\def\discrversionD#1#2{\ifnum\stylenumber=4#1\else#2\fi}
\long\def\discrversionE#1#2{\ifnum\stylenumber=5#1\else#2\fi}
\long\def\discrversionF#1#2{\ifnum\stylenumber=6#1\else#2\fi}
\long\def\discrversionG#1#2{\ifnum\stylenumber=7#1\else#2\fi}
\long\def\discrversionH#1#2{\ifnum\stylenumber=8#1\else#2\fi}
           % facility for special interpolations in particular output
           % versions
\def\displaycause#1{\noalign{\noindent (#1)}}
\def\diverg{\mathop{\text{div}}}
\def\dom{\mathop{\text{dom}}}
\def\eae{=_{\text{a.e.}}}

\def\esssup{\mathop{\text{ess sup}}}
\mathchardef\Exists="3239

\def\family#1#2#3{\langle#3\rangle_{#1\in#2}}
\def\familyi#1#2{\family{i}{#1}{#2}}
\def\familyiI#1{\familyi{I}{#1}}

\def\frfilename#1{\ifnum\stylenumber=1{\noindent filename #1}
               \else\ifnum\stylenumber=6{\noindent{\fiverm filename #1}}
               \else\ifnum\stylenumber=7{\noindent{\fiverm filename #1}}
               \else\fi\fi\fi}
\def\Footnote#1#2{\footnote"#1"{#2}}
\mathchardef\Forall="3238
\def\frakc{\frak c\phantom{.}}
\def\frnewpage{\vskip \pageheight plus 0pt minus\pageheight
           \eject
           \gdef\topparagraph{}
           \gdef\bottomparagraph{}
           }
\def\frsmallcirc{{\hskip 1pt plus 0pt minus 0pt}
     {\raise 1pt\hbox{$\smallcirc$}}{\hskip 1pt plus 0pt minus 0pt}}

\def\geae{\ge_{\text{a.e.}}}
\def\grad{\mathop{\text{grad}}}
\def\grheada{{\bf ($\pmb{\alpha}$)}}
\def\grheadb{{\bf ($\pmb{\beta}$)}}
\def\grheadc{{\bf ($\pmb{\gamma}$)}}
\def\grheadd{{\bf ($\pmb{\delta}$)}}
\def\grheade{{\bf ($\pmb{\epsilon}$)}}
\def\grheadz{{\bf ($\pmb{\zeta}$)}}
\def\grv#1{{\accent"12 #1}}

\def\hfillall{\hskip\pagewidth plus0pt minus\pagewidth}

\def\Imag{\mathop{\Cal I\text{m}}}
\long\def\inset#1{{\narrower\narrower{\vskip 1pt plus 1pt minus 0pt}
     #1
     {\vskip 1pt plus 1pt minus 0pt}}}
\def\Int#1#2{\displaystyle{\int_{#1}^{#2}}}
\def\interior{\mathop{\text{int}}}

\def\leae{\le_{\text{a.e.}}}
\long\def\leaveitout#1{}
\def\LLcorner{\mathchoice{\,\hbox{\vrule height7pt width0.5pt depth0pt
       \vrule height0.5pt width5pt depth0pt}\,}
   {\,\hbox{\vrule height7pt width0.5pt depth0pt
       \vrule height0.5pt width5pt depth0pt}\,}
   {\,\hbox{\vrule height5pt width0.3pt depth0pt
       \vrule height0.3pt width4pt depth0pt}\,}
   {\,\hbox{\vrule height4pt width0.25pt depth0pt
       \vrule height0.25pt width3pt depth0pt}\,}
  }
\def\Ln{\mathop{\text{Ln}}}
\def\Loadfont#1{\font#1}
\def\loibr{\hbox{\hskip0.1em]\hskip0.05em}}
\def\Matrix#1{\pmatrix#1\endpmatrix}
\def\Mu{\text{M}}
\def\noamslogo{\def\amslogo{}}
\def\normalsubgroup{\vartriangleleft}
\def\Nu{\text{N}}
\def\ocint#1{\left]#1\right]}
\def\ofamily#1#2#3{\langle#3\rangle_{#1<#2}}
\let\oldcases=\cases
\let\oldfootnote=\footnote
\def\On{\mathop{\text{On}}}
\def\ooint#1{\left]#1\right[}
\def\openBbb#1{\Bbb#1\mskip1mu}
\def\pd#1#2{\bover{\partial#1}{\partial#2}}
\def\Pd#1#2{\Bover{\partial#1}{\partial#2}}
\def\pdd#1#2{{{\partial^2#1}\over{\partial#2^2}}}
\def\Pdd#1#2{\Bover{\partial^2#1}{\partial#2^2}}
\def\pde#1#2#3{{{\partial^2#1}\over{\partial#2\partial#3}}}
\def\Pde#1#2#3{\Bover{\partial^2#1}{\partial#2\partial#3}}
\def\Prf{\vthsp\hbox{\bf P\hskip0.07em\llap P\hskip0.07em\llap P\rm}}
\def\pushbottom#1{\hbox{\vrule height0pt depth#1 width0pt}}
\def\Qed{\vthsp\hbox{\bf Q\hskip0.07em\llap Q\hskip0.07em\llap Q\rm}}
\def\QeD{\Qed{\hskip 1.111pt plus 3.333pt minus 0pt}}
\def\Quer{\vthsp\hbox{\bf ?\hskip0.05em\llap ?\hskip0.05em\llap
   ?\thinspace\rm}}
\def\query{\discrversionA{\immediate\write0{query}
    \global\advance\footnotenumber by 1
    \oldfootnote{$^{\the\footnotenumber}$}{query}}{}}
\def\Real{\mathop{\Cal R\text{e}}}
\def\restr{\mathchoice
  {\hbox{$\restriction$}\mskip 1mu}
  {\hbox{$\restriction$}\mskip 1mu}
  {\hbox{$\scriptstyle\restriction$}\mskip 1mu}
  {\hbox{$\scriptscriptstyle\restriction$}\mskip 1mu}}
\def\restrp{\mathchoice
  {\hbox{$\restriction$}\mskip 2mu}
  {\hbox{$\restriction$}\mskip 2mu}
  {\mskip-1mu\hbox{$\scriptstyle\restriction$}\mskip 2mu}
  {\mskip-1mu\hbox{$\scriptscriptstyle\restriction$}\mskip 1mu}}
    %use in \restrp V \restrp W \restrp Y
    % \restrp\Cal P \restrp\Bbb R \restrp\frak
    % \restrp\Omega \restrp\Tau \restrp\Upsilon
    % \restrp[ \restrp\widehat{}
\def\Rho{\text{P}}
\def\roibr{\hbox{\hskip0.05em[\hskip0.1em}}
\def\RoverC{\hbox{\biggerscriptfonts${{\Bbb R}\atop{\Bbb C}}$}}
\def\sequence#1#2{\family{#1}{\Bbb N}{#2}}
\def\sequencen#1{\sequence{n}{#1}}
\def\setover#1#2{\overset{#1}\to{#2}}
\def\shortstrut{\hbox{\vrule height7.5pt depth2.5pt width0pt}}
\def\sshortstrut{\hbox{\vrule height5.8pt depth2.5pt width0pt}}
\def\smallcirc{{\scriptstyle\circ}}
\def\ssbullet{{\scriptscriptstyle\bullet}}
\def\symmdiff{\triangle}
\def\thicken#1{#1\hskip0.03em\llap #1\hskip0.03em\llap
   #1\hskip0.03em\llap #1}
\def\Tau{{\text{T}}}
\def\tbf#1{\text{\bf #1}}
\def\terror#1#2#3{#1:  `#2' for `#3'}
\def\today{\number\day.\number\month.\number\year}
\def\UEMDRR{University of Essex Mathematics Department Research Report}
\def\UEMDWP{University of Essex Mathematics Department Working Paper}
\def\versiondate#1{
    \immediate\write0{Version of #1}
    \def\vsdate{#1}
               \ifnum\stylenumber=1\hfill Version of #1
                              \medskip
               \else\ifnum\stylenumber=2\hfill Version of #1
                              \medskip
               \else\ifnum\stylenumber=6\hfill{\fiverm Version of #1}
                              \medskip
               \else\ifnum\stylenumber=7\hfill{\fiverm Version of #1}
                              \medskip
               \else\ifnum\stylenumber=8\hfill Version of #1
                              \medskip
               \else\fi\fi\fi\fi\fi}
\def\versionhead{\ifnum\stylenumber=3 \hfill PRINTER'S VERSION

                        \bigskip
                   \else\ifnum\stylenumber=4 \hfill REFEREE'S VERSION

                        \bigskip
                   \else \fi\fi}
\def\versionref{\ifnum\stylenumber=1\workingref
           \else\ifnum\stylenumber=2\workingref
           \else\ifnum\stylenumber=8\workingref
           \else\printerref\fi\fi\fi}
\def\vthsp{\hskip0.001em}

%Missing AmSTeX can be coped with by using the following lines:

\ifamstexsupported\else
   \immediate\write0{AmSTeX-unsupported version of  fremtex.tex }
   \fourteensavailablefalse
   \def\@{@}
   \def\\{\cr}
   \def\AmSTeX{AmS\TeX}
   \def\Bbb#1{\underline{\hbox{\bf #1}}}
   \def\biggerscriptfonts{}
   \def\bumpeq{=_{\text{approx}}}
   \def\Cal#1{{\cal#1}}
   \def\curlyeqprec{\preceq}
   \def\eightit{\it}
   \def\eightrm{\rm}
   \def\eightsmc{\smc}
   \def\eurm#1{\underline{\cal#1}}
   \def\Footnote#1#2{\footnote#1#2}
   \def\frac#1#2{{#1}\over{#2}}
   \def\frak#1{{\hbox{\bf #1\hskip0.1em\llap #1}}}
   \def\frnewpage{\vfill\eject}
   \def\iint{\int\!\!\!\int}
   \def\iiint{\int\!\!\!\int\!\!\!\int}
   \def\iiiint{\int\!\!\!\int\!\!\!\int\!\!\!\int}
   \def\llcorner{\lfloor}
   \def\loadeurm{}
   \def\loadeusm{}
%   \def\loadeusmx#1{\def\eusm#1{\underline{\cal#1}}}
   \def\ltimes{|\hskip-0.5em\times}
   \def\eusm#1{\underline{\cal#1}}
   \def\ltimes{|\times}
   \def\Loadfont#1{}
   \def\Matrix#1{\pmatrix{#1}}
   \def\noamslogo{}
   \def\nrightarrow{\not\rightarrow}
   \def\preccurlyeq{\preceq}
   \def\query{\discrversionA{\immediate\write0{query}
                             \footnote{$^*$}{query}}{}}
   \def\restriction{|}
   \def\rtimes{\times|}
   \def\setover#1#2{#2^{#1}}
   \def\shortstrut{\strut}
   \def\smallerfonts{\eightrm}
   \def\smallfrown{\wedge}
   \def\smc{}
   \def\spcheck{\check{\phantom{x}}}
   \def\sphat{\hat{\phantom{x}}}
   \def\square{\mathchoice\sqr64\sqr64\sqr43\sqr33}
   \def\sqr#1#2{{\vcenter{\vbox{\hrule height.#2pt
          \hbox{\vrule width.#2pt height#1pt \kern#1pt
                  \vrule width.#2pt}
          \hrule height.#2pt}}}}
   \def\pmb#1{\underline{#1}}
   \def\text#1{\hbox{#1}}
   \def\therefore{\ssbullet\raise.9ex\hbox{$\ssbullet$}\ssbullet}
   \def\twelvebf{\bf}
   \def\vartriangleleft{\triangleleft}
   \def\Vdash{|\vdash}
   \def\vDash{|=}
   \fi

%End of lines for missing AmSTeX

\ifamstexsupported\input amstex
%              \documentstyle{amsppt}
               \input amsppt.sty
               \PSAMSFonts
	       \tenpoint
               %\font\eightmsb=msym8
               \font\eightmsb=msbm10
               \font\eighteufm=eufm10 at 8 pt
               \font\eighteurm=eurm8
               \font\eighteusm=eusm8
               \font\varseveneufm=eufm10 at 7 pt
               \scriptfont\eufmfam=\varseveneufm
   \def\.#1{{\accent"5F #1}}
   \def\=#1{{\accent"16 #1}}
   \gdef\nopagenumbers{\NoPageNumbers}
   \else\fi

\newcount\footmarkcount
\ifDHFamstexloaded
    \loadmsBm
    \def\eightBbb#1{{\fam\msBfam#1}}
    \loadmsFm
    \def\eightfrak#1{{\fam\msFfam#1}}
  \else
    \def\eightBbb#1{{\scriptstyle\Bbb#1}}
    \def\eightfrak#1{{\scriptstyle\frak#1}}
  \fi
\ifathome\athome\fi

%interword glue (\tenrm standard):
%  (space) 3.333 + 1.667 - 1.111
%  (,)     3.333 + 2.083 - 0.889
%  (.)     4.444 + 5.    - 0.370

%very odd fact:  \discrversionA{2.}{}  doesn't work,
%   \discrversionA{ 2.}{} does

%The false line feed introduced by my Fortran programs, leading to
%a surplus  ae  at the top of the next page, seems to be associated
%with  Fk(\032)3081  or something similar in the PostScript file.


\immediate\write0{Extract from MEASURE THEORY,
   by D.H.Fremlin, University of Essex, Colchester.
   This material is copyright.   It is issued under the terms of the
   Design Science License as published in
   http://dsl.org/copyleft/dsl.txt.   This is a development version and the
   source files are not permanently archived.
   For further information contact david\@fremlin.org.}

%hyphenations
\hyphenation{Archi-me-dean}

%new style
\def\smallprintersize{\magnification=\magstep0
     \overfullrule0pt
     \def\pagewidth{468pt}
     \def\pageheight{9truein}
     \hsize\pagewidth
     \vsize\pageheight
     \nologo
     \def\copyrightdate##1{}
     \def\folio{\ifnum\pageno<0{}\else\number\pageno\fi}
     \def\mtfootline{}
     \stylenumber=10
     \pageno=-1000}

%With-proofs, results-only versions
\newif\ifwithproofs
\withproofstrue
\newif\ifresultsonly
\resultsonlyfalse
\def\noproofs{\withproofsfalse}
\def\resultsonly{\resultsonlytrue\withproofsfalse}
\long\def\proof#1{\ifwithproofs{\medskip

     \noindent{\bf proof} #1}\else\fi}
\long\def\prooflet#1{\ifwithproofs{#1}\else\fi}
%\long\def\Prooflet#1{\prooflet{\Prf\ #1 \Qed\ }}
\long\def\endnotes#1{\ifresultsonly\else{#1}\fi}
\long\def\exercises#1{\ifresultsonly\else{#1}\fi}
\long\def\references#1{\ifresultsonly\else{#1}\fi}
\long\def\dvro#1#2{\ifresultsonly{#1}\else{#2}\fi}

%Headline assignments

\def\topparagraph{}
\def\bottomparagraph{}
\def\newparagraph{}
\newif\ifpageturned

\newdimen\rhoffset
\newdimen\lhoffset
\lhoffset=0cm
\rhoffset=0cm

\newcount\volumeno
\volumeno=5
\newcount\luluvolumeno
\luluvolumeno=0

\let\plainpagebody=\pagebody
\def\pagebody{\ifodd\pageno\global\hoffset=\rhoffset
         \else\global\hoffset=\lhoffset\fi
     \plainpagebody
     \turnpage}
\def\turnpage{
     \iffirstpage\else{\baselineskip=24pt\line{\mtfootline}}\fi
%     \immediate\write0{turning page \folio,
%           top paragraph \topparagraph,   bottom \bottomparagraph,
%           new \newparagraph, pagetotal \the\pagetotal}
     \global\pageturnedtrue
     \global\firstpagefalse
     \global\let\topparagraph=\bottomparagraph
     \global\let\bottomparagraph=\newparagraph
     }

\newdimen\pagetotalplus
\newdimen\pagetotalminus

%wheader parameters:  #1 new paragraph name
 % #2#3#4  \vskip parameters if page break not forced (no. of points)
 % #5  spare height demanded if page break not to be forced
\long\def\wheader#1#2#3#4#5{\gdef\newparagraph{#1}
  \pagetotalplus=\the\pagetotal
  \advance\pagetotalplus by #5
  \pagetotalminus=\the\pagetotal
  \advance\pagetotalminus by -\pageshrink
  %\showthe\pagetotal \showthe\pageshrink \showthe\pagegoal
  \global\pageturnedfalse
  \ifdim\pagetotalminus<\pagegoal\ifdim\pagetotalplus>\pagegoal
%     \immediate\write0{forcing page break}
     \eject
     \else\vskip #2pt plus #3pt minus #4pt\fi
     \else\vskip #2pt plus #3pt minus #4pt\fi
  \ifpageturned\global\let\topparagraph=\newparagraph
%     \immediate\write0{paragraph change on turning page}
     \fi
%  \immediate\write0{top paragraph \topparagraph, bottom
%       \bottomparagraph, pagetotal \the\pagetotal}
  \ifdim\pagetotal>\pagegoal\else\gdef\bottomparagraph{#1}\fi
%   \immediate\write0{top paragraph \topparagraph, bottom
%       \bottomparagraph, pagetotal \the\pagetotal}
  }
\def\header#1{
     \wheader{#1}{10}{4}{4}{24pt}}
\def\leader#1#2{\immediate\write0{#1 #2}
     \wheader{#1}{14}{7}{5}{24pt}
     {\bf #1 #2}}
\def\vleader#1#2#3{\immediate\write0{#2 #3}
     \wheader{#2}{14}{7}{5}{#1}
     {\bf #2 #3}}
\def\Notesheader#1{\immediate\write0{#1 Notes and comments}
     \wheader{#1 Notes}{16}{8}{5}{48pt}
     \noindent{\bf #1 Notes and comments}}
\def\vNotesheader#1#2{\immediate\write0{#2 Notes and comments}
     \wheader{#2 Notes}{16}{8}{5}{#1}
     \noindent{\bf #2 Notes and comments}}
\def\spheader#1#2#3#4#5{\header{#1#2#3#4#5}{\bf (#5)}}
\def\spstheader#1#2#3#4#5{\header{*#1#2#3#4#5}{\bf (#5)}}
\def\vspheader#1#2#3#4#5#6{
     \wheader{#2#3#4#5#6}{10}{4}{4}{#1}{\bf (#6)}}
\def\sqheader#1#2#3#4#5{\header{#1#2#3#4#5}{$\pmb{>}${\bf (#5)}}}
\def\errheader#1#2#3{
     \wheader{#3}{6}{4}{2}{24pt}
     {\bf p #1 l #2}}
\long\def\woddheader#1#2#3#4#5{
   \ifodd\pageno\wheader{#1}{#2}{#3}{#4}{#5}
   \else\vskip #2pt plus #3pt minus #4pt\fi}

\def\indexmedskip{
  \pagetotalplus=\the\pagetotal
  \advance\pagetotalplus by 18 pt
  \pagetotalminus=\the\pagetotal
  \advance\pagetotalminus by -\pageshrink
  \ifdim\pagetotalminus<\pagegoal\ifdim\pagetotalplus>\pagegoal
      \else\medskip\fi\else\medskip\fi}

\def\indexheader#1{\immediate\write0{#1}\vindexheader{#1}{24}}
\def\vindexheader#1#2{
  \gdef\newparagraph{#1}
  \pagetotalplus=\the\pagetotal
  \advance\pagetotalplus by #2 pt
  \pagetotalminus=\the\pagetotal
  \advance\pagetotalminus by -\pageshrink
  %\showthe\pagetotal \showthe\pageshrink \showthe\pagegoal
  \global\pageturnedfalse
%  \ifdim\pagetotal>\pagegoal
%      \immediate\write0{warning: pagetotal greater than pagegoal}
%      \immediate\write0{...   last paragraph reached is #1}
%      \immediate\write0{...   check for underfull vbox here}
%      \else\fi
  \ifdim\pagetotalminus<\pagegoal\ifdim\pagetotalplus>\pagegoal{\eject}
      \fi\fi
  \ifdim\pagetotal>\pagegoal\else\gdef\bottomparagraph{#1}\fi
  \ifpageturned{\global\let\topparagraph=\newparagraph
%     \immediate\write0{paragraph change on turning page}
      }\else\fi
%  \immediate\write0{top paragraph \topparagraph, bottom
%    \bottomparagraph}
  }

\def\leftheadline{\rlap{\folio}\hfill{\eightit\chaptername}\hfill
       \llap{\eightbf\topparagraph}}
\def\rightheadline{\rlap{\eightbf\bottomparagraph}\hfill
       {\eightit\headlinesectionname}\hfill\llap{\folio}}

%Footline assignments
\def\mtfootline{\ifodd\pageno\rightfootline\else\leftfootline\fi}
\def\rightfootline{\hfill\eightsmc D.H.Fremlin}
\def\leftfootline{\ifwithproofs{\eightsmc Measure Theory\hfill}
    \else{\eightsmc Measure Theory \eightit (abridged version)\hfill}
    \fi}

%Style assignments
\def\newchapter#1{\gdef\topparagraph{}
   \gdef\bottomparagraph{Chap.\ #1 {\it intro.}}
   \gdef\newparagraph{Chap.\ #1 {\it intro.}}
   \footmarkcount=0
   \gdef\sectionname{Introduction}
   \gdef\headlinesectionname{Introduction}
   \centerline{\bf Chapter #1}
   \medskip
   \centerline{\bf \chaptername}
   \medskip
   }
\def\newsection#1{\gdef\bottomparagraph{\S#1 {\it intro.}}
   \noindent{\bf #1 \sectionname}
   \smallskip
   \let\headlinesectionname=\sectionname}
\def\newvolume#1{\gdef\topparagraph{}
   \gdef\bottomparagraph{Vol.\ #1 {\it intro.}}
   \Loadfont{\twelvebf=cmbx12}
   \def\chaptername{Volume #1}
   \def\sectionname{\volumename}
   \def\headlinesectionname{\volumename}
   \centerline{\twelvebf Volume #1:  \volumename}
   \bigskip
   }

%font variations
\newif\ifvolwp
\volwpfalse
\def\tenrmstretch#1{\fontdimen3\tenrm=#1}
\font\elevenex=cmex10 at 11pt

%contents and introductions
\newif\iflargelogo
\def\vtmpb#1{}
\def\newvolume#1{}
\def\pagereference#1#2{\hfill\hbox{\ifresultsonly#2\qquad \hskip 1truein
                                                 \else#1\qquad \fi}}
\long\def\section#1#2#3#4#5#6{\ifresultsonly
    \qquad #1 #2\pagereference{#4}{#5}
  \else\vskip 2pt plus 1pt minus 0pt
    \vbox{\qquad #1 #2\pagereference{#4}{#5}\par
    \abstr{#6}
    }\vskip -1pt plus 1pt minus 0pt
  \fi}
\long\def\chapintrosection#1#2#3{\ifresultsonly
    \qquad Introduction\pagereference{#2}{#3}
  \else\vskip 1pt plus 1pt minus 0pt
    \qquad Introduction\pagereference{#2}{#3}
    \vskip 1pt plus 1pt minus 0pt
  \fi}
\def\abstr#1{\ifresultsonly{}
   \else{\advance\leftskip by\parindent
     \advance\rightskip by\parindent
     \advance\leftskip by\parindent
     \advance\rightskip by\parindent
     \advance\leftskip by\parindent
     \advance\rightskip by\parindent
     \advance\leftskip by\parindent
     \vskip1pt plus 0pt minus0pt
     \parindent=0pt
     \smallerfonts #1%
     \hbox{\vrule height0pt depth2pt width0pt}
     \vskip1pt plus 0pt minus0pt}\fi}
     
%language variations
\newif\ifUSEnglish
\USEnglishfalse
%\USEnglishtrue

%"Later editions only"
\newcount\latereditionpage
\newcount\latereditionmark
\latereditionpage=-1
\latereditionmark=0
\newcount\newpageno

%removing the \write0 instruction altogether disorganizes it!
\def\Latereditions{\newpageno=\pageno%
  \ifdim\pagetotal>\pagegoal\advance\newpageno by 1\fi%
     \divide\newpageno by 2%
  \immediate\write0{%
%     "Later editions only";
%      pagetotal \the\pagetotal, pagegoal \the\pagegoal,
%      pageno \number\pageno,\par newpageno \number\newpageno,
%      latereditionpage \number\latereditionpage%
  }\ifnum\latereditionpage=\newpageno%
     \ifDHFamspptloaded$^{\number\latereditionmark}$%
     \else\footnote{Later editions only.}\fi%
  \else\footnote{Later editions only.}%
  \global\latereditionpage=\newpageno%
  \global\latereditionmark=\footmarkcount\fi}%
%  \immediate\write0{newpageno \number\newpageno,
%      latereditionpage \number\latereditionpage}

   
%general definitions
\def\action{\mathchoice{\lower-.3ex\hbox{$\ssbullet$}}
                       {\lower-.3ex\hbox{$\ssbullet$}}
                       {\ssbullet}
                       {\ssbullet}}  %see also \varaction
\def\add{\mathop{\text{add}}\nolimits}
\long\def\allowmorestretch#1#2{
  \ifdim\pagewidth<#1pt\else
    \fontdimen3\tenrm=2pt\fontdimen4\tenrm=1.67pt\fi
  #2
    \fontdimen3\tenrm=1.67pt\fontdimen4\tenrm=1.11pt}
\def\am{atom\discretionary{-}{}{}lessly-measurable}
\def\AmuA{\Aut_{\bar\mu}\mskip-2mu\frak A}
\def\ASint{\mathchoice{\Bbb A\mskip-5mu}
                  {\Bbb A\mskip-3mu}
                  {\Bbb A\mskip-2mu}
                  {\Bbb A\mskip-2mu}\Sint}
\def\Asisfour{$(\frak A,\bar\mu,T,\family{t}{T}{\frak A_t})$}
\def\Asissix{$(\frak A,\bar\mu,T,
   \penalty-100\family{t}{T}{\frak A_t},
   \penalty-100\Cal T,\family{\tau}{\Cal T}{\frak A_{\tau}})$}
\def\Asisfourt{$(\frak A,\bar\mu,\coint{0,\infty},
   \langle\frak A_t\rangle_{t\ge 0})$}
\def\Asissixt{$(\frak A,\bar\mu,\coint{0,\infty},
   \penalty-100\langle\frak A_t\rangle_{t\ge 0},
   \penalty-100\Cal T,\family{\tau}{\Cal T}{\frak A_{\tau}})$}
\def\Atop#1#2{{\scriptstyle #1}\atop{\scriptstyle #2}}
\def\Aut{\mathop{\text{Aut}}}

\def\Bcap{\Booleanop{\cap}}
\def\Bcup{\Booleanop{\cup}}
\def\Booleanop#1{
   \ifamstexsupported{\mathchoice
    {\hbox{\hskip.2em\scriptfont2=\eightsy$\scriptstyle#1$\hskip.2em}}
    {\hbox{\hskip.2em\scriptfont2=\eightsy$\scriptstyle#1$\hskip.2em}}
 %   {{\mskip2mu\scriptstyle#1\mskip2mu}}  %until 15.10.17
    {{\scriptscriptstyle#1}}  %from 15.10.17
    {{\scriptscriptstyle#1}}}
   \else{\mathchoice{{\scriptstyle#1}}{{\scriptstyle#1}}
    {{\scriptstyle#1}}{{\scriptscriptstyle#1}}}
    \fi}
\def\Bsetminus{\Booleanop{\setminus}}
\def\Bsubseteq{\hskip.1em\Booleanop{\subseteq}\hskip.1em}
\def\Bsubset{\Booleanop{\subset}}
\def\Bsupseteq{\Booleanop{\supseteq}}
\def\Bsupset{\Booleanop{\supset}}
\def\Bsymmdiff{\Booleanop{\symmdiff}}
\def\Bcapshort{\Booleanopshort{\cap}}
\def\Bcupshort{\Booleanopshort{\cup}}
\def\Booleanopshort#1{
   \ifamstexsupported{\mathchoice
    {\hbox{\scriptfont2=\eightsy$\scriptstyle#1$}}
    {\hskip-1pt\hbox{\scriptfont2=\eightsy$\scriptstyle#1$}}
    {{\scriptstyle#1}}
    {{\scriptscriptstyle#1}}}
   \else{\mathchoice{{\scriptstyle#1}}{{\scriptstyle#1}}
    {{\scriptstyle#1}}{{\scriptscriptstyle#1}}}
    \fi}
\def\Bsetminusshort{\Booleanopshort{\setminus}}
\def\BSint{\mathchoice{\Bbb B\mskip-5mu}
                  {\Bbb B\mskip-3mu}
                  {\Bbb B\mskip-2mu}
                  {\Bbb B\mskip-2mu}\Sint}
%\def\Bsubseteqshort{\hskip.1em\Booleanopshort{\subseteq}\hskip.1em}
\def\Bsubseteqshort{\hskip.1em\Booleanopshort{\subseteq}}
\def\Bsubsetshort{\Booleanopshort{\subset}}
\def\Bsupseteqshort{\Booleanopshort{\supseteq}}
\def\Bsymmdiffshort{\Booleanopshort{\symmdiff}}

\def\closeBooleanop#1{
   \ifamstexsupported{\mathchoice
    {\hbox{\scriptfont2=\eightsy$\scriptstyle#1$}}
    {\hbox{\scriptfont2=\eightsy$\scriptstyle#1$}}
    {{\scriptstyle#1}}
    {{\scriptscriptstyle#1}}}
   \else{\mathchoice{{\scriptstyle#1}}{{\scriptstyle#1}}
    {{\scriptstyle#1}}{{\scriptscriptstyle#1}}}
    \fi}
\def\closeBcup{\closeBooleanop{\cup}}
\def\closeBcap{\closeBooleanop{\cap}}
\def\closeBsetminus{\closeBooleanop{\setminus}}
\def\closeBsymmdiff{\closeBooleanop{\symmdiff}}

\def\biggerHeint{\textfont3=\twelveex
   \raise 0.6pt\hbox to 0pt{\hskip 2.6pt{\sevenrm H}\hss}
   \intop\nolimits}
\def\biggerPfint{\textfont3=\twelveex
   \raise 0.6pt\hbox to 0pt{\hskip 1.3pt{\sevenrm Pf}\hss}
   \intop\nolimits\hskip 1.2pt}
\def\biggerSint{\textfont3=\twelveex
   \raise 0.6pt\hbox to 0pt{\hskip 3.6pt{\sevenrm S}\hss}
   \intop\nolimits}
\def\bu{\mathop{\text{bu}}}
\def\Bvalue#1{\mathchoice
   {\hbox{$[\![#1]\!]$}}
   {\hbox{$[\![#1]\!]$}}
   {\hbox{$\scriptstyle[\![#1]\!]$}}
   {\hbox{$\scriptscriptstyle[\![#1]\!]$}}}
\def\cadlag{c\`adl\`ag}
\def\CalBa{\Cal{B}\eurm{a}}
   %note: neither  \cal  nor  \eusm  has lower case
\def\CalFr{\Cal F_{\text{Fr}}}
\def\callal{c\`all\`al}
\def\CalNwd{\Cal{N}\hskip-1.2pt\eurm{w}\hskip-.8pt\eurm{d}}

\def\CalRbg{\Cal{R}\mskip-1mu\eurm{b}\mskip-1mu\eurm{g}}
\def\CalSmz{\Cal{S}\mskip-2mu\eurm{mz}}
\def\CalUn{\Cal{U}\eurm{n}}
\def\capacity{\mathop{\text{cap}}\nolimits}
\def\Caratheodory{Cara\-th\'eo\-dory}
\def\Cdlg{C_{\text{dlg}}}
\def\cff{\cf\mskip1mu}
\def\checkpmb#1{\check{\pmb #1}}
\def\ci{\mathop{\text{ci}}}
\def\cl{\mathop{\text{cl}}}
\def\Clll{\hbox{$C$
  \hskip-4.5pt\raise4pt\hbox{\char'23
  \hskip-1.6pt\vrule height6.5pt width.5pt depth0pt
  \hskip.7pt\vrule height6.5pt width.5pt depth0pt
  \hskip-2.5pt\lower5.4pt\hbox{\char'23}}}}
\def\clMsimp{\overline{M}_{\text{simp}}}
\def\closeuparrow{\mathchar"0222}
\def\closedownarrow{\mathchar"0223}
\def\clstar{\text{cl*}}
\def\cov{\mathop{\text{cov}}}
\def\covar#1#2{\mathchoice
  {[#1{\lower4pt\hbox{${\scriptstyle *}\atop{\vrule height 18pt depth -10pt}$}}#2]}
  {[#1{*\atop{\vrule height 8pt depth 0pt}}#2]}
  {[#1{*\atop{\vrule height 6pt depth 0pt}}#2]}
  {[#1{*\atop{\vrule height 6pt depth 0pt}}#2]}}
\def\covSh{\mathop{\text{cov}}\nolimits_{\text{Sh}}}
\def\cdownarrow{c\mskip2mu^{\downarrow}}
\def\cycle#1{(\overleftarrow{#1})}
\def\ddaggerloc{$(\ddagger)_{\text{loc}}$}
\def\ddownarrow{d\mskip2mu^{\downarrow}}
%\def\DiniD{\setover{
%  \hbox{\hskip.1em\vrule height-1pt width.5em depth1.4pt\hfil}}D}
%\def\Dinid{\hskip.1em
%  \hbox{\hskip.05em\vrule height-1.3pt width.5em depth1.7pt}\hskip-.65em D}
\def\Dinid{\hbox{$D$\hskip-8.3pt\vrule height-1pt width7pt depth1.85pt}
   \hskip1pt}
\def\DiniD{\hbox{$D$\hskip-8.3pt\vrule height9.2pt width7ptdepth-8.35pt}
   \hskip1pt}

\def\discrcenter#1#2{\ifdim\pagewidth>#1{#2}\else{

  \Centerline{#2}

  \noindent}\fi}
    %note:  in \discrcenter, make sure that #2 ends in a space
    % (if you want one to appear) and that the continuation
    % follows without
    % a space after the final } of #2, so that it will be
    % truly unindented
\def\discrpage{\ifnum\stylenumber=1{\frnewpage}  %redefined in mt.tex
        \else\ifnum\stylenumber=3{\frnewpage}
        \else\ifnum\stylenumber=5{\frnewpage}
        \else\wheader{}{18}{14}{2}{72pt}\fi\fi\fi}
\def\displaycause#1{\noalign{\noindent (#1)}}
\def\Displaycause#1{\cr\displaycause{#1}&}
\def\dotproduct{\hskip.1em\tbf{.}\hskip.1em}
\def\doubleheadrightarrow{\rightarrow\mskip-12mu\rightarrow}
\long\def\doubleinset#1{\inset{\inset{\parindent=-20pt #1}}}
\def\duparrow{d\mskip2mu^{\uparrow}}
\def\dvAformerly#1{\discrversionA{\footnote{Formerly #1.}}{}}
\def\dvAnew#1{\discrversionA{\footnote{New #1.}}{}}
\def\dvArevised#1{\discrversionA{\footnote{Revised #1.}}{}}
\def\dvrocolon#1{\dvro{#1: }{#1}}
\def\eightVVdash{\mskip5mu\vrule height 6pt depth 2pt width 0.4pt
  \mskip2mu\vrule height 6pt depth 2pt width 0.4pt
  \vrule height 2.25pt depth -1.8pt width 3.2pt\mskip2mu}
\def\embedsinto{\mskip5mu\raise1.5pt\hbox{$\subset$}\mskip-15mu
   \lower3.5pt\hbox{$\rightarrow$}\mskip5mu}
\def\energy{\mathop{\text{energy}}\nolimits}
%\eqalign:  see \S478 for breaks during \eqalign
\def\equivGT{\equiv_{\text{GT}}}
\def\equivT{\equiv_{\text{T}}}
\def\Eta{\text{H}}
\def\Expn{\Bbb E}
\long\def\Extract{\oldfootnote{}{
   \noindent{\eightit
       \ifresultsonly{Extract from {\eightsmc Measure Theory,}
          results-only version,
          by {\eightsmc D.H.Fremlin,}  University of Essex,
Colchester.}
       \else\ifwithproofs{Extract from {\eightsmc Measure Theory,}
          by {\eightsmc D.H.Fremlin,}  University of Essex,
Colchester.}
       \else{Extract from {\eightsmc Measure Theory,} abridged version,
          by {\eightsmc D.H.Fremlin,}  University of Essex,
Colchester.}
       \fi\fi
   This material is copyright.   It is issued under the terms of the
   Design Science License as published in
   {\eighttt http://dsl.org/copyleft/dsl.txt}.
   This is a development version and the
   source files are not permanently archived, but current versions are
normally
   accessible through
   {\eighttt https://www1.essex.ac.uk/maths/people/fremlin/mt.htm}.
   For further information contact {\eighttt david\@fremlin.org}.}
   }}
\def\Fn{\mathop{\text{Fn}}\nolimits}
\def\FN{\mathop{\text{FN}}\nolimits}
\def\formerly#1{\footnote{Formerly #1.}}
  %use only once in a chapter, never in same volume as reference
\def\fraction#1{\mathchoice{\hbox{$<$}#1\hbox{$>$}}
     {\hbox{$<$}#1\hbox{$>$}}
     {\hbox{$\scriptstyle<$}#1\hbox{$\scriptstyle>$}}
     {\hbox{$\scriptscriptstyle<$}#1\hbox{$\scriptscriptstyle>$}}}
\def\frakmctbl{\frak m_{\text{countable}}}
\def\frmedcirc{{\raise 2.5pt\hbox{$_{\bigcirc}$}}}
\def\grad{\mathop{\text{grad}}\nolimits}
\def\hatpmb#1{\hat{\pmb{#1}}}
\def\hattilde#1{\skew2\hat{\tilde#1}}
\def\hattildepmb#1{\skew1\hat{\tildepmb#1}}
\def\hcov{\mathop{\text{hcov}}\nolimits}
\def\Heint{\mathchoice{\hbox to 0pt{\hskip 3.1pt H\hss}}
     {\raise 0.6pt\hbox to 0pt{\hskip 1.9pt{\sevenrm H}\hss}}
     {\raise 0.4pt\hbox to 0pt{\hskip 1.5pt{\fiverm H}\hss}}
     {\raise 0.4pt\hbox to 0pt{\hskip 1.5pt{\fiverm H}\hss}}
     \intop\nolimits}
\def\high#1{\lceil#1\rceil}
\def\highoverline#1{\overline{\vrule height6pt width0pt depth 0pt#1}}
\def\Hint#1{({\it Hint\/}: #1)}
\def\hL{\mathop{\text{hL}}}
\def\Hmp{\text{Hom}_{\bar\mu,\pi}}
\def\hp{\mathop{\text{hp}}\nolimits}
\def\hwtr{hereditarily weakly $\theta$-{\vthsp}refinable}
\def\iiv{ii_{\pmb{v}}}
\def\iivu{\iiv(\pmb{u})}
\def\imp{inverse-{\vthsp}measure-{\vthsp}preserving}
\def\innerprod#1#2{(#1|#2)}
\def\intstar{\text{int*}}
\def\CalIS{\Cal I(\Cal S)}
\def\jump{\mathop{\text{jump}}\nolimits}
\def\Jump{\mathop{\text{Jump}}\nolimits}
\def\Krein{Kre\v{\i}n}
\def\Lambdapv{\Lambda_{\text{pv}}}
\def\leRB{\le_{\text{RB}}}
\def\leRK{\le_{\text{RK}}}
\def\link{\mathord{\text{link}}}
\def\locmafp{\widehat{\otimes}_{\text{loc}}}
\def\Locmafp{\widehat{\bigotimes}^{\text{loc}}}
\def\low#1{\lfloor#1\rfloor}
\def\lti{left-\vthsp translation-\vthsp invariant}
\def\magnitude{\mathop{\text{mag}}}
\def\Mah{\mathop{\text{Mah}}}
\def\Mahcr{\Mah_{\text{cr}}}
\def\MahcrR{\Mah_{\text{crR}}}
\def\MahqR{\Mah_{\text{qR}}}
\def\MahR{\Mah_{\text{R}}}
\def\MahspcrR{\Mah_{\text{spcrR}}}
\def\Malm{M_{\text{alm}}}
\def\Marik{Ma\v r\'\i k} 
\def\Mbv{M_{\text{bv}}}
\def\Mcdlg{M_{\text{cdlg}}}
\def\Mdm{M_{\text{D-M}}}
\def\med{\mathop{\text{med}}\nolimits}
\def\Mfa{M_{\text{fa}}}
\def\Mhsr{\mathop{\text{Mhsr}}\nolimits}
\def\Mjf{\mathop{M_{\text{j-f}}}\nolimits}
\def\Mlbv{M_{\text{lbv}}}
\def\Mlob{M_{\text{lo-b}}}
\def\Mlns{M_{\text{ln-s}}}
\def\Mmo{M_{\text{mo}}}
\def\Mlmo{M_{\text{lmo}}}
\def\Mns{M_{\text{n-s}}}
\def\Mnsu{M_{\text{n-s}}^{\uparrow}}
%\def\Mnsu{M_{\text{n-s}\uparrow}}
\def\MMah{M_{\text{Mah}}}
\def\Mob{M_{\text{o-b}}}
\def\Mpob{M_{\text{po-b}}}
\def\Mpv{M_{\text{pv}}}
\def\MSi{M_{\text{S-i}}}
\def\Msimp{M_{\text{simp}}}
\def\Mth{Ma\-ha\-ram-{\vthsp}type-{\vthsp}homogeneous}
\def\Mvlm{M_{\text{vlm}}}
\def\negquad{\hskip-1em\relax}
\def\NN{\BbbN^{\Bbb N}}
\def\non{\mathop{\text{non}}}
\def\notBooleanop#1{
   \ifamstexsupported{\mathchoice
{\hbox{\hskip.2em$\not\scriptfont2=\eightsy\scriptstyle#1$\hskip.2em}}
{\hbox{\hskip.2em$\not\scriptfont2=\eightsy\scriptstyle#1$\hskip.2em}}
    {{\not\scriptstyle#1}}
    {{\not\scriptscriptstyle#1}}}
   \else{\mathchoice{{\not\scriptstyle#1}}{{\not\scriptstyle#1}}
    {{\not\scriptstyle#1}}{{\not\scriptscriptstyle#1}}}
    \fi}
\def\notBsubseteq{\notBooleanop{\subseteq}}
\def\notBsupseteq{\notBooleanop{\supseteq}}
\def\notVVdash{\not\mskip4mu\vrule height 7.5pt depth 2.5pt width 0.5pt
  \mskip2.5mu\vrule height 7.5pt depth 2.5pt width 0.5pt
  \vrule height 2.75pt depth -2.25pt width 4pt\mskip-2mu}
\def\nuprime{\nu^{\mskip0.7mu\prime}} %better  \nu'  in superscripts
\def\nw{\mathop{\text{nw}}}
\def\Osclln{\mathop{\text{Osclln}}\nolimits}
\def\otp{\mathop{\text{otp}}}
\def\partstar{\mathchoice{\partial\hbox{*}}
     {\partial\hbox{*}}
     {\partial\hbox{\sevenrm *}}
     {\partial\hbox{\fiverm *}}}
\def\per{\mathop{\text{per}}}
\def\Pfint{\mathchoice{\hbox to 0pt{\hskip 3.1pt Pf\hss}}
     {\raise 0.6pt\hbox to 0pt{\hskip 1.3pt{\sevenrm Pf}\hss}}
     {\raise 0.4pt\hbox to 0pt{\hskip 1.0pt{\fiverm Pf}\hss}}
     {\raise 0.4pt\hbox to 0pt{\hskip 1.0pt{\fiverm Pf}\hss}}
     \intop\nolimits
     \mathchoice{\hskip 2pt}{\hskip 1.2pt}{\hskip 1pt}{\hskip 1pt}}
\def\Pou{\mathop{\text{Pou}}\nolimits}
\def\prforc{\preccurlyeq_{\text{forcing}}}
\def\prGT{\preccurlyeq_{\text{GT}}}
\def\prT{\preccurlyeq_{\text{T}}}
\def\pssqa{power set $\sigma$-quotient algebra}
\def\qm{quasi-{\vthsp}measurable}
\def\Reverse#1{\setover
   {\lower0.4ex\hbox{$\scriptscriptstyle\leftrightarrow$}}{#1}}
\def\rhoKR{\rho_{\text{KR}}} 
\def\rhotv{\rho_{\text{tv}}} 
\def\Rint{\mathchoice{\hbox to 0pt{\hskip 3.1pt R\hss}}
     {\raise 0.6pt\hbox to 0pt{\hskip 1.9pt{\sevenrm R}\hss}}
     {\raise 0.4pt\hbox to 0pt{\hskip 1.5pt{\fiverm R}\hss}}
     {\raise 0.4pt\hbox to 0pt{\hskip 1.5pt{\fiverm R}\hss}}
     \intop\nolimits}
\def\RO{\text{RO}}
\def\rti{right-{\vthsp}translation-{\vthsp}invariant}
\def\rvm{real-{\vthsp}valued-{\vthsp}measurable}
\def\sat{\mathop{\text{sat}}}
\def\sfamilysS#1{\family{\sigma}{\Cal S}{#1_{\sigma}}}
\def\sfcc{$\sigma$-finite-cc}
\def\sgn{\mathop{\text{sgn}}}
\def\shr{\mathop{\text{shr}}\nolimits}
\def\Sii{\text{S}ii}
\def\Siiv{\Sii_{\pmb v}}
\def\Sint{\mathchoice{\hbox to 0pt{\hskip 4.1pt\ninerm S\hss}} 
     {\raise 0.6pt\hbox to 0pt{\hskip 3.0pt{\sevenrm S}\hss}}
     {\raise 0.4pt\hbox to 0pt{\hskip 1.5pt{\fiverm S}\hss}} %unchecked
     {\raise 0.4pt\hbox to 0pt{\hskip 1.5pt{\fiverm S}\hss}} %unchecked
     \intop\nolimits}
\def\Sinteightpoint{\raise 0.6pt\hbox to 0pt{\hskip 2.3pt{\fiverm S}\hss}
     \intop\nolimits}
\def\SL{\text{SL}}
\def\sm{semi-martingale} %hyphen to remind me that it's the "virtual" kind
\def\Sous{\Cal S\text{ous}}  %used in vol 6
\def\spread{\mathop{\text{spr}}}
\def\ssplrarrow{^{\scriptscriptstyle\leftrightarrow}}
\def\ssptilde{{\raise1.2ex\hbox{$\scriptscriptstyle\sim$}}}
%\def\sspvec{{\raise1.2ex\hbox{$\scriptscriptstyle\rightarrow$}}}
\def\sspvec{\mskip3mu\vec{\null}\mskip5mu}
\def\Sti{\mathop{\text{Sti}}\nolimits}
\def\strbar{^{\scriptscriptstyle\pmb{-}}}
\def\strplus{^{\scriptscriptstyle\pmb{+}}}
\def\strprime{^{\mathchoice{\mskip1.5mu\pmb{\prime}}
                  {\mskip1.5mu\pmb{\prime}}
                  {\mskip1.5mu\pmb{\prime}}
                  {\mskip1mu\pmb{\prime}}}}
%\def\strprime{\,\pmb{\prime}}
\def\supp{\mathop{\text{supp}}}
\def\tausm{\tau_{\sigma\text{-m}}}
\def\tbfoneminus#1{\rlap{$\tbf 1^{(#1)}$}\rlap{\hskip7pt$_<$}
   \phantom{\tbf 1^{(#1)}}}
%\def\tbfoneminus#1{\rlap{$\tbf 1^{(#1)}$}\rlap{\hskip7pt$_-$}
%   \phantom{\tbf 1^{(#1)}}}
\def\tensorhat{\widehat{\otimes}}
\def\Tensorhat{\widehat{\bigotimes}}
\def\tensorhatsigma{\widehat{\otimes}_{\sigma}}
\def\tildeClll{\hbox{$\tilde C$
  \hskip-4.5pt\raise4pt\hbox{\char'23
  \hskip-1.6pt\vrule height6.5pt width.5pt depth0pt
  \hskip.7pt\vrule height6.5pt width.5pt depth0pt
  \hskip-2.5pt\lower5.4pt\hbox{\char'23}}}}
\def\tildehat#1{\skew2\tilde{\hat#1}}
\def\tildehatpmb#1{\tilde{\hatpmb#1}}
\def\tildepmb#1{\tilde{\pmb{#1}}}
\def\tinyplus{{\vrule height 1.6pt depth -1.3pt width 2.5pt}
    \hskip-1.3pt{\vrule height 2.6pt depth -0.1pt width 0.3pt}}
\def\Tr{\mathop{\text{Tr}}}
\def\trs{^{\top}\hskip-1pt}
\def\triplepc#1#2#3{(#1,#2,\hbox{$<$}#3)}
\def\TSi{\frak T_{\text{S-i}}}
\def\undphi{\underline{\phi}{}}
\def\undpsi{\underline{\psi}{}}
\def\undtheta{\underline{\theta}{}}
\def\updown{\mathchoice
  {\hbox{\hskip-1.5pt\char'23
    \hskip-1.7pt\vrule height6.5pt width.7pt depth1.3pt
    \hskip.9pt\vrule height6.5pt width.7pt depth1.3pt
    \hskip-3.3pt\lower6.7pt\hbox{\char'23}\hskip-1.4pt}}
  {\hbox{\hskip-2pt\char'23
    \hskip-1.7pt\vrule height6.5pt width.7pt depth1.3pt
    \hskip.9pt\vrule height6.5pt width.7pt depth1.3pt
    \hskip-2.5pt\lower6.7pt\hbox{\char'23}\hskip-1.4pt}}
  {\hbox{\hskip-2pt\char'23  %OK for subscript, fails for superscript
    \hskip-1.7pt\vrule height6.5pt width.7pt depth0pt
    \hskip.9pt\vrule height6.5pt width.7pt depth0pt
    \hskip-2.5pt\lower5.4pt\hbox{\char'23}\hskip-1.4pt}}
  {\hbox{\hskip-2pt\char'23  %not tested
    \hskip-1.7pt\vrule height6.5pt width.7pt depth1.3pt
    \hskip.9pt\vrule height6.5pt width.7pt depth1.3pt
    \hskip-2.5pt\lower6.7pt\hbox{\char'23}\hskip-1.4pt}}}
\def\updownarrows{\uparrow\mskip-4mu\downarrow}
\def\updowncgt{$\updown$-convergent}
\def\upr{\mathop{\text{upr}}\nolimits}
\def\ureg{\mathop{\text{ureg}}}

\def\Var{\mathop{\text{Var}}\nolimits}

\def\varaction{\mathchoice{\lower-.2ex\hbox{$\smallcirc$}}
                       {\lower-.2ex\hbox{$\smallcirc$}}
                       {\ssbullet}
                       {\ssbullet}}
\def\varcheck#1{{\setover
       {\lower 0.4ex\hbox{$\scriptscriptstyle\vee$}}{#1}}}
\def\varcheckf{{\setover{\,\lower 0.4ex
              \hbox{$\scriptscriptstyle\vee$}}{f}}}
\def\varchecki{\mskip-3mu\text{\it{\accent"14{\hbox{\hskip-3.5pt\i}}}}\mskip2mu}
\def\varspcheck{^{\scriptscriptstyle\vee}}
%\def\varspcheck{{\setover{\scriptscriptstyle\vee}{\phantom{g}}}}
%\def\varspcheckf{{\setover{\lower 0.4ex
%              \hbox{$\scriptscriptstyle\vee$}}{\phantom{f}}}}
\def\varhat#1{{\setover{\lower 0.4ex
              \hbox{$\scriptscriptstyle\wedge$}}{#1}}}
\def\varhatf{{\setover{\,\lower 0.4ex
              \hbox{$\scriptscriptstyle\wedge$}}{f}}}
\def\varinnerprod#1#2{#1\dotproduct#2}
\def\varsphat{^{\scriptscriptstyle\wedge}}
% or use  \hbox{\^{\null}}
%\def\varsphat{{\setover{\scriptscriptstyle\wedge}{\phantom{g}}}}
%\def\varsphatf{{\setover{\lower 0.4ex
%              \hbox{$\scriptscriptstyle\wedge$}}{\phantom{f}}}}
\def\varBspcheck{\mathchoice
  {\hskip-2pt
  {\fourteenrm\hbox{\lower2pt\hbox{\v{\null}}}}\hskip-1pt}
  {\hskip-2pt
  {\fourteenrm\hbox{\lower2pt\hbox{\v{\null}}}}\hskip-1pt}
  {\hskip-2.5pt
  {\fourteenrm\hbox{\lower4pt\hbox{\v{\null}}}}\hskip-1pt}
  {\hskip-2.5pt
  {\fourteenrm\hbox{\lower4pt\hbox{\v{\null}}}}\hskip-1pt}}
\def\varCspcheck{\mathchoice
  {\hskip-1pt\check{\phantom x}}   %not tested
  {\hskip-1pt\check{\phantom x}}
  {\hskip-1.2pt\check{\phantom x}}
  {\hskip-1.5pt\check{\phantom x}}}  %not tested
\def\varCsphat{\mathchoice
  {\hskip-1pt\hat{\phantom x}}   %not tested
  {\hskip-1pt\hat{\phantom x}}
  {\hskip-1.2pt\hat{\phantom x}}  %not tested
  {\hskip-1.5pt\hat{\phantom x}}}  %not tested
\def\VeqL{$V\hbox{$=$}L$}
\def\VVdash{\mskip5mu\vrule height 7.5pt depth 2.5pt width 0.5pt
  \mskip2.5mu\vrule height 7.5pt depth 2.5pt width 0.5pt
  \vrule height 2.75pt depth -2.25pt width 4pt\mskip2mu}
\def\VVdP{\VVdash_{\Bbb P}}
\def\VVdPk{\VVdash_{\Bbb P_{\kappa}}}
\def\wdistr{\mathop{\text{wdistr}}}
\def\weight{\mathop{\text{wt}}}
\def\w1lim{\mathop{\text{w1lim}}}
\def\wsid{weakly $(\sigma,\infty)$-{\vthsp}distributive}
\def\1lim{\mathop{1\mskip-1.5mu\text{lim}}}
\def\2vm{two-{\vthsp}valued-{\vthsp}measurable}

%Cara\-th\'e\-o\-dory

\def\chaptername{}
\def\sectionname{}
\def\headlinesectionname{}

\loadeusm\loadeurm
\Loadfourteens
\Loadfont{\ninerm=cmr9}

%note:  $..\cmmnt{=..}..$  is better as  \cmmnt{\mskip5mu
 
  
%----------------------------------------------------------%  
  
\smallprint  
%\USsmallprint  
%\noproofs  
%\resultsonly  
%\atUEssex  
\usegraphicx
  
\Extract  
  
%----------------------------------------------------------%  

\frfilename{mt52.tex} 
\versiondate{7.10.13} 
\copyrightdate{2005} 
 
\def\chaptername{Cardinal functions of measure theory} 
 
\newchapter{52} 
 
From the point of view of this book, the most important cardinals are 
those associated with measures and measure algebras, especially, of 
course, Lebesgue measure and the usual measure $\nu_I$ of $\{0,1\}^I$. 
In this 
chapter I try to cover the principal known facts about these which are 
theorems of ZFC.   I start with a review of the theory for general measure 
spaces in \S521,  
including some material which returns to the classification 
scheme of Chapter 21, exploring relationships between (strict) 
localizability, magnitude and Maharam type.    
\S522 examines Lebesgue measure and the surprising 
connexions found by {\smc Bartoszy\'nski 84} and 
{\smc Raisonnier \& Stern 85} between the cardinals associated with the 
Lebesgue null 
ideal and the corresponding ones based on the ideal of meager subsets of 
$\Bbb R$.   \S523 looks at the 
measures $\nu_I$ for uncountable sets $I$, giving formulae for the 
additivities and cofinalities of their null ideals, and bounds for 
their covering numbers, uniformities and shrinking numbers.    
Remarkably, these cardinals 
are enough to tell us most of what we want to know concerning the 
cardinal functions of general Radon measures and semi-finite measure 
algebras (\S524).   These three sections are heavily dependent on the 
Galois-Tukey connections and Tukey functions of \S\S512-513.    
Precalibers do 
not seem to fit into this scheme, and the relatively partial information 
I have is in \S525.   The second half of the chapter deals with special 
topics which can be approached with the methods so far developed.    
In \S526 I return to the ideal of subsets of $\Bbb N$ 
with asymptotic density zero, seeking to locate it in the Tukey 
classification.   Further $\sigma$-ideals which are of interest in 
measure theory are the `skew products' of \S527.    
In \S528 I examine some interesting 
Boolean algebras, the `amoeba algebras' first introduced by 
{\smc Martin \& Solovay 70}, giving the results of {\smc Truss 88} 
on the connexions 
between different amoeba algebras and localization posets.   Finally, 
in \S529, I look at a handful of other structures, concentrating on 
results involving cardinals already described. 
 
\discrpage 
 
 
\frfilename{mt521.tex}
\versiondate{3.3.14}
\copyrightdate{2007}

\def\chaptername{Cardinal functions of measure theory}
\def\sectionname{Basic theory}



\newsection{521}

In the first half of this section (down to 521L) I collect facts
about the cardinal
functions $\add$, $\cf$, $\non$, $\cov$, $\shr$ and $\shr^+$ when applied
to the null ideal $\Cal N(\mu)$ of a measure $\mu$, and also the
$\pi$-weight of a measure.
In particular I look at their relations with the
constructions introduced earlier in this treatise:
measure algebras and function
spaces (521B), subspace measures (521F), direct sums (521G), \imp\
functions and image measures (521H), products (521J), perfect measures
(521K) and compact measures (521L).   The list is long just because I
have four volumes' worth of miscellaneous concepts to examine;  nearly all
the individual arguments are elementary.

In the second half of the section, I give
a handful of easy results which may clarify some
patterns from earlier volumes.   In 521M-521P %521M 521N 521O 521P
I look again at `strict localizability' as considered in Chapter 21,
importing the concept of `magnitude' of a measure space from \S332,
hoping to throw light on the examples of \S216.
In 521E I consider the topological densities of measure algebras.
In 521R-521S I
explore possibilities for the `countably separated' measure spaces of
\S\S343-344, examining in particular their Maharam types.   Finally,
in 521T, I review some measures which arose in \S464 while analyzing the
$L$-space $\ell^{\infty}(I)^*$.

\leader{521A}{Proposition} Let $(X,\Sigma,\mu)$ be a measure space.

(a) If $\Cal E\subseteq\Sigma$ and $\#(\Cal E)<\add\mu$ then
$\bigcup\Cal E\in\Sigma$ and

\Centerline{$\mu(\bigcup\Cal E)
=\sup\{\mu(\bigcup\Cal E_0):\Cal E_0\subseteq\Cal E$ is finite$\}$.}

(b) $\omega_1\le\add\mu\le\add\Cal N(\mu)$.

(c) If $\mu$ is the measure defined by \Caratheodory's method from
an outer measure $\theta$ on $X$, then $\add\mu=\add\Cal N(\mu)$.

(d) If $\mu$ is complete and locally determined,
$\add\mu=\add\Cal N(\mu)$.

\proof{{\bf (a)} Induce on $\#(\Cal E)$.   If $\Cal E$ is finite, the
result is trivial.   For the inductive step to
$\#(\Cal E)=\kappa\ge\omega$, enumerate $\Cal E$ as
$\ofamily{\xi}{\kappa}{E_{\xi}}$.   For each $\xi<\kappa$, set
$H_{\xi}=E_{\xi}\setminus\bigcup_{\eta<\xi}E_{\eta}$ for each
$\xi<\kappa$.   Then the inductive hypothesis tells us that
$H_{\xi}\in\Sigma$ for every $\xi$.   Set $E=\bigcup\Cal
E=\bigcup_{\xi<\kappa}H_{\xi}$;  because
$\ofamily{\xi}{\kappa}{H_{\xi}}$ is
disjoint, and $\kappa<\add\mu$, $E\in\Sigma$ and

$$\mu E
=\sum_{\xi<\kappa}\mu H_{\xi}
=\sup_{I\subseteq\kappa\text{ is finite}}
  \mu(\bigcup_{\xi\in I}H_{\xi})
\le\sup_{\Cal E_0\subseteq\Cal E\text{ is finite}}\mu(\bigcup\Cal E)
\le\mu E.$$

\medskip

{\bf (b)} By the definition of `measure' (112A), $\mu$ is
$\omega_1$-additive.   Suppose that $\Cal A\subseteq\Cal N(\mu)$ and
$\#(\Cal A)<\add\mu$.   For each $A\in\Cal A$, choose a measurable
negligible $E_A\supseteq A$.   Then (a) tells us that
$E=\bigcup_{A\in\Cal A}E_A$ has measure zero, so
$\bigcup\Cal A\subseteq E$ is negligible.   As $\Cal A$ is arbitrary,
$\add\Cal N(\mu)\ge\add\mu$.

\medskip

{\bf (c)} Now suppose that $\mu$ is defined by \Caratheodory's method
from $\theta$.
Let $\familyiI{E_i}$ be a disjoint family in $\Sigma$, where
$\#(I)<\add\Cal N(\mu)$, with union $E$.

Let $A\subseteq X$ be any set.   Then
$\theta(A\cap E)=\sum_{i\in I}\theta(A\cap E_i)$.   \Prf\ Of course

$$\eqalignno{\theta(A\cap E)
&\ge\sup_{J\subseteq I\text{ is finite}}
  \theta(A\cap\bigcup_{i\in J}E_i)
=\sup_{J\subseteq I\text{ is finite}}
  \sum_{i\in J}\theta(A\cap E_i)\cr
\displaycause{induce on $\#(J)$, using the fact that
$\theta B=\theta(B\cap E_i)+\theta(B\setminus E_i)$ for every
$B\subseteq X$ and $i\in J$}
&=\sum_{i\in I}\theta(A\cap E_i).\cr}$$

\noindent If $\sum_{i\in I}\theta(A\cap E_i)$ is infinite, we can stop.
Otherwise, recalling that $\Cal N(\mu)=\theta^{-1}[\{0\}]$,
$J=\{i:A\cap E_i\notin\Cal N(\mu)\}$ is countable,
and $\bigcup_{i\in I\setminus J}A\cap E_i$ is negligible, because
$\#(I)<\add\Cal N(\mu)$;  so

\Centerline{$\theta(A\cap E)=\theta(A\cap\bigcup_{i\in J}E_i)
\le\sum_{i\in J}\theta(A\cap E_i)=\sum_{i\in I}\theta(A\cap E_i)$}

\noindent and we have equality.\  \Qed

It follows that
$\theta(A\cap E)+\theta(A\setminus E)\le\theta A$.   \Prf\ For any finite
$J\subseteq I$,

$$\eqalign{\theta(A\setminus E)+\sum_{i\in J}\theta(A\cap E_i)
&=\theta(A\setminus E)+\theta(A\cap\bigcup_{i\in J}E_i)\cr
&\le\theta(A\setminus\bigcup_{i\in J}E_i)+\theta(A\cap\bigcup_{i\in J}E_i)
=\theta A.\cr}$$

\noindent Taking the supremum over $J$, we have the result.\ \Qed

As $A$ is arbitrary, $E\in\Sigma$;  and setting $A=E$, we see that
$\mu E=\sum_{i\in I}\mu E_i$.   As $\familyiI{E_i}$ is arbitrary,
$\add\mu\ge\add\Cal N(\mu)$ and the two additivities are equal.

\medskip

{\bf (d)} Now this follows immediately from (c), by 213C.
}%end of proof of 521A

\leader{521B}{Proposition} Let $(X,\Sigma,\mu)$ be a measure space
and $(\frak A,\bar\mu)$ its measure algebra.

(a) If $\Cal E\subseteq\Sigma$ and $\#(\Cal E)<\add\mu$, then
$(\bigcup\Cal E)^{\ssbullet}=\sup_{E\in\Cal E}E^{\ssbullet}$ and
$(X\cap\bigcap\Cal E)^{\ssbullet}=\inf_{E\in\Cal E}E^{\ssbullet}$
in $\frak A$.

(b) Suppose that $A\subseteq[-\infty,\infty]^X$ is a non-empty family of
$\Sigma$-measurable functions with $\#(A)<\add\mu$, and that
$g(x)=\sup_{f\in A}f(x)$ in $[-\infty,\infty]$ for every $f\in A$.
Then $g$ is $\Sigma$-measurable.

(c) Write $\eusm L^0$ for the family of $\mu$-virtually measurable
real-valued functions defined almost everywhere in $X$, and $L^0$ for
the corresponding space of equivalence classes, as in \S241.   Suppose
that $A\subseteq\eusm L^0$ is such that $0<\#(A)<\add\mu$ and
$\{f^{\ssbullet}:f\in A\}$ is bounded above in $L^0$.   Set
$g(x)=\sup_{f\in A}f(x)$ whenever this is defined in $\Bbb R$;  then
$g\in\eusm L^0$ and $g^{\ssbullet}=\sup_{f\in A}f^{\ssbullet}$ in $L^0$.

\allowmorestretch{500}{
(d)(i) If, in (b), $A$ consists of non-negative integrable functions and
is upwards-directed, then $\int g\,d\mu=\sup_{f\in A}\int fd\mu$.
}

\quad(ii) If, in (b), $f_1\wedge f_2=0$ a.e.\ for all distinct $f_1$,
$f_2\in A$, then $\int g\,d\mu=\sum_{f\in A}\int fd\mu$.

\proof{{\bf (a)} As in 521Aa, $\bigcup\Cal E\in\Sigma$, and of course
$(\bigcup\Cal E)^{\ssbullet}$ is an upper bound for
$\{E^{\ssbullet}:E\in\Cal E\}$.   If $F\in\Sigma$ and $F^{\ssbullet}$ is an
upper bound for $\{E^{\ssbullet}:E\in\Cal E\}$, then, applying 521Aa to
$\{E\setminus F:E\in\Cal E\}$, we see that $\bigcup\Cal E\setminus F$ is
negligible, so $(\bigcup\Cal E)^{\ssbullet}\Bsubseteq F^{\ssbullet}$.
Thus $(\bigcup\Cal E)^{\ssbullet}$ is the least upper bound of
$\{E^{\ssbullet}:E\in\Cal E\}$.

Applying this to $\{X\setminus E:E\in\Cal E\}$ we see that
$(X\cap\bigcap\Cal E)^{\ssbullet}=\inf_{E\in\Cal E}E^{\ssbullet}$.

\medskip

{\bf (b)} For any $\alpha\in\Bbb R$,

\Centerline{$\{x:g(x)>\alpha\}
=\bigcup_{f\in A}\{x:f(x)>\alpha\}\in\Sigma$}

\noindent by 521Aa.

\medskip

{\bf (c)} Take any $h\in\eusm L^0$ such that
$f^{\ssbullet}\le h^{\ssbullet}$ for every $f\in A$.   For each
$f\in A$, let $E_f$ be a conegligible measurable subset of
$\{x:x\in\dom f\cap\dom h$, $f(x)\le h(x)\}$ such that $f\restr E_f$ is
measurable.   Set $E=\bigcap_{f\in A}E_f$;  then $E$ is measurable and
$g$ is defined everywhere in $E$ and $g\restr E$ is measurable (as in
(b)).    Also $E$ is conegligible, so
$g\in\eusm L^0$, and of course $f^{\ssbullet}\le g^{\ssbullet}$ for
every $f\in A$, while $g^{\ssbullet}\le h^{\ssbullet}$.   But this
argument works for every $h$ such that $h^{\ssbullet}$ is an upper bound
for $\{f^{\ssbullet}:f\in A\}$, so $g^{\ssbullet}$ must be actually the
supremum of $\{f^{\ssbullet}:f\in A\}$.

\medskip

{\bf (d)(i)} If $\sup_{f\in A}\int fd\mu$ is infinite, this is trivial.
Otherwise, $\{f^{\ssbullet}:f\in A\}$ is bounded above in $L^1$ and
therefore in $L^0$.   By (c), $g^{\ssbullet}$ is its supremum in $L^0$,
therefore in $L^1$;  so

\Centerline{$\int g=\int g^{\ssbullet}
=\sup_{f\in A}\int f^{\ssbullet}=\sup_{f\in A}\int f$,}

\noindent as in 365Df.

\medskip

\quad{\bf (ii)} Apply (i) to $A^*=\{\sup I:I\in[A]^{<\omega}\}$.
}%end of proof of 521B

\leader{521C}{}\cmmnt{ Just because null ideals are $\sigma$-ideals of
sets, we can read off some of the elementary properties of their cardinal
functions from 511J.   But the presence of a measure gives us a new way
to use shrinking numbers, which will be useful later.

\medskip

\noindent}{\bf Proposition} Let $(X,\Sigma,\mu)$ be a measure space, and
$A\subseteq X$.

(a) If $\gamma<\mu^*A$ there is a $B\subseteq A$
such that $\#(B)<\shr^+\Cal N(\mu)$ and $\mu^*B>\gamma$.

(b) There is a $B\subseteq A$ such that
$\#(B)\le\max(\omega,\shr\Cal N(\mu))$ and $\mu^*B=\mu^*A$.

%exercise on  \non^+Cal N(\mu) ?
% on \omega_1-sat (X,\Sigma,\Cal I) ?

\proof{{\bf (a)} Set $\kappa=\shr^+\Cal N(\mu)$.
Let $\Cal E$ be the
family of those measurable subsets of $X$ such that there is a
$B\in[A\cap E]^{<\kappa}$ with $\mu^*B=\mu E$.   Then $\Cal E$ is
closed under finite unions (132Ed).   \Quer\ If $\mu^*B\le\gamma$
for every $B\in[A]^{<\kappa}$, then $\mu E\le\gamma$ for every
$E\in\Cal E$.   By 215Ab,
there is a non-decreasing sequence $\sequencen{E_n}$ in $\Cal E$ such
that $E\setminus\bigcup_{n\in\Bbb N}E_n$ is negligible for every
$E\in\Cal E$.   Now $\mu(\bigcup_{n\in\Bbb N}E_n)\le\gamma<\mu^*A$
and $A'=A\setminus\bigcup_{n\in\Bbb N}E_n$ is not negligible.   Let
$B\in[A']^{<\kappa}$ be a non-negligible set.    Then $\mu^*B\le\gamma$ is
finite, so $B$ has a measurable envelope $F$ (132Ee),
which belongs to $\Cal E$;
but $F\setminus\bigcup_{n\in\Bbb N}E_n\supseteq B$
is not negligible.\ \BanG\
So we have a $B\in[A]^{<\kappa}$ with $\mu^*B>\gamma$, as required.

\medskip

{\bf (b)} If $\mu^*A=0$ take $B=\emptyset$.   Otherwise, let
$\sequencen{\gamma_n}$ be a sequence in $\coint{0,\mu^*A}$ with supremum
$\mu^*A$.   For each $n\in\Bbb N$, (a) tells us that there is a set
$B_n\subseteq A$ such that $\#(B_n)\le\shr(\mu)$ and
$\mu^*B_n>\gamma_n$;  set $B=\bigcup_{n\in\Bbb N}B_n$.
}%end of proof of 521C

\leader{521D}{Proposition} Let $(X,\Sigma,\mu)$ be a measure space and
$(\frak A,\bar\mu)$ its measure algebra.

(a) $\pi(\frak A)\le\pi(\mu)\le\max(\pi(\frak A),\cf\Cal N(\mu))$\cmmnt{
(definitions:  511Dc, 511Gb)}.

(b) If $\mu X>0$, then $\non\Cal N(\mu)\le\pi(\mu)$.

(c) If $(X,\Sigma,\mu)$
has locally determined negligible sets\cmmnt{ (definition:
213I)}, then $\shr\Cal N(\mu)\le\pi(\mu)$.

(d) Suppose that
there is a topology $\frak T$ on $X$ such that
$(X,\frak T,\Sigma,\mu)$ is a
quasi-Radon measure space.   Then, writing $\frak A^+$ for
$\frak A\setminus\{0\}$, the
partially ordered sets $(\Sigma\setminus\Cal N(\mu),\supseteq)$
and $(\frak A^+,\Bsupseteqshort)$ are Tukey equivalent and
$\pi(\mu)=\pi(\frak A)$.

\proof{ Let $\Cal H\subseteq\Sigma\setminus\Cal N(\mu)$ be a
coinitial set with cardinal $\pi(\mu)$.

\medskip

{\bf (a)(i)} If $a\in\frak A$ is non-zero, there is an
$E\in\Sigma$ such that $E^{\ssbullet}=a$, and now $E$ is not negligible, so
there is an $H\in\Cal H$ such that $H\subseteq E$ and
$0\ne H^{\ssbullet}\Bsubseteq a$.   Thus $\{H^{\ssbullet}:H\in\Cal H\}$ is
coinitial with $\frak A^+$ and witnesses that
$\pi(\frak A)\le\#(\Cal H)=\pi(\mu)$.

\medskip

\quad{\bf (ii)} Let $B\subseteq\frak A^+$ be a coinitial set
with cardinal $\pi(\frak A)$, and $\Cal E$ a cofinal subset of $\Cal N(\mu)$ of
size $\cf\Cal N(\mu)$.   For $b\in B$, let $F_b\in\Sigma$ be such that
$F_b^{\ssbullet}=b$, and consider
$\Cal G=\{F_b\setminus E:b\in B$, $E\in\Cal E\}$.   Then
$\Cal G\subseteq\Sigma\setminus\Cal N(\mu)$ is coinitial with
$\Sigma\setminus\Cal N(\mu)$.   \Prf\ If $\mu F>0$, there is a $b\in B$
such that $b\Bsubseteq F^{\ssbullet}$.   In this case, $F_b\setminus F$ is
negligible, so there is an $E\in\Cal E$ such that
$F_b\setminus F\subseteq E$ and $F\supseteq F_b\setminus E\in\Cal G$.\ \Qed

It follows that $\pi(\mu)\le\#(\Cal G)\le\#(B\times\Cal E)$
is at most the cardinal product
$\pi(\frak A)\cdot\cf\Cal N(\mu)\le\max(\omega,\pi(\frak A),\cf\Cal N(\mu))$.
But if $\cf\Cal N(\mu)$ is finite it is $1$, so in fact
$\pi(\mu)\le\pi(\frak A)\cdot\cf\Cal N(\mu)=\max(\pi(\frak A),\cf\Cal N(\mu))$.

\medskip

{\bf (b)} For each $H\in\Cal H$ choose $x_H\in\Cal H$.   Then
$A=\{x_H:H\in\Cal H\}$ must meet every non-negligible measurable set,
so (as $\mu X>0$) cannot itself be negligible.   Thus

\Centerline{$\non\Cal N(\mu)\le\#(A)\le\#(\Cal H)=\pi(\mu)$.}

\medskip

{\bf (c)} Suppose that $B\subseteq X$ is non-negligible.
Because $(X,\Sigma,\mu)$ has locally determined
negligible sets there is an $E\in\Sigma$ such that $\mu E>0$ and
$B\cap E$ is not negligible, and now $B\cap E$ has a measurable envelope
$F$ say (132Ee again).   Set
$\Cal H'=\{H:H\in\Cal H$, $B\cap H\ne\emptyset\}$ and for
$H\in\Cal H'$ choose $x_H\in B\cap H$;  set $A=\{x_H:H\in\Cal H'\}$, so
that $A\subseteq B$ and $\#(A)\le\pi(\mu)$.
\Quer\ If $A$ is negligible, then $F\setminus A$ includes a non-negligible
measurable set so includes a member $H$ of $\Cal H$.   As $\mu H>0$ and
$F$ is a measurable envelope of $B$, $H$ meets $B$ and belongs to
$\Cal H'$, and $x_H\in A\cap H$.\ \BanG\  Thus $A$ is not negligible.
As $B$ is arbitrary, $\shr\Cal N(\mu)\le\pi(\mu)$.

\medskip

{\bf (d)} For $E\in\Sigma\setminus\Cal N(\mu)$ let $F_E$ be a closed
non-negligible subset of $E$ and set
$\phi(E)=F_E^{\ssbullet}\in\frak A^+$;  for $a\in\frak A^+$, let
$\psi(a)$ be a self-supporting measurable set such that
$\psi(a)^{\ssbullet}=a$ (414F).   Then if $\phi(E)\Bsupseteq a$,
$\psi(a)\setminus F_E$ is negligible so
$E\supseteq F_E\supseteq\psi(a)$.   Thus $(\phi,\psi)$ is a Galois-Tukey
connection and
$(\Sigma\setminus\Cal N(\mu),\supseteq,\Sigma\setminus\Cal N(\mu))
\prGT(\frak A^+,\Bsupseteqshort,\frak A^+)$.

Moreover, if $\psi(a)\supseteq E$, then $a\Bsupseteq\phi(E)$, so
$(\psi,\phi)$ also is a Galois-Tukey connection and
$(\frak A^+,\Bsupseteqshort,\frak A^+)\prGT
(\Sigma\setminus\Cal N(\mu),\discretionary{}{}{}\supseteq\nobreak,
\discretionary{}{}{}\Sigma\setminus\Cal N(\mu))$.

Thus $(\Sigma\setminus\Cal N(\mu),\supseteq,\Sigma\setminus\Cal N(\mu))
\equivGT(\frak A^+,\Bsupseteqshort,\frak A^+)$, that is,
$(\Sigma\setminus\Cal N(\mu),\supseteq)
\equivT(\frak A^+,\Bsupseteqshort)$.   By 513E(e-i), inverted,

\Centerline{$\pi(\mu)=\ci(\Sigma\setminus\Cal N(\mu))
=\ci(\frak A^+)=\pi(\frak A)$.}
}%end of proof of 521D

\leader{521E}{}\cmmnt{ It will be useful later in the chapter to be able to
calculate the topological density of measure-algebra topologies.

\medskip

\noindent}{\bf Proposition} Let $(\frak A,\bar\mu)$ be a semi-finite
measure algebra.

(a) Give $\frak A$ its measure-algebra topology\cmmnt{ (323A)}.

\quad(i) If $\frak B$ is a subalgebra of $\frak A$, it is topologically
dense iff it $\tau$-generates $\frak A$\cmmnt{, that is, $\frak A$ is the
order-closed subalgebra of itself generated by $\frak B$}.

\quad(ii) If
$\frak A$ is finite, then its topological density is $\#(\frak A)$;  if
$\frak A$ is infinite, its topological density is equal to its Maharam
type $\tau(\frak A)$.

(b) Let $\frak A^f$ be the set of elements of $\frak A$ with
finite measure, with its strong measure-algebra
topology\cmmnt{ (323Ad)}.   Then the
topological density of $\frak A^f$ is $\#(\frak A^f)=\#(\frak A)$ if
$\frak A$
is finite, and $\max(c(\frak A),\tau(\frak A))$ if $\frak A$ is infinite.

\proof{{\bf (a)(i)}\grheada\ Suppose that $\frak B$ is topologically dense.
Let $\frak C$ be the order-closed subalgebra
of $\frak A$ generated by $\frak B$.
If $a\in\frak A^f$ and $c\in\frak A$,
there is a $b\in\frak C$ such that $b\Bcap a=c\Bcap a$.   \Prf\ For each
$n\in\Bbb N$, there is an
$a_n\in\frak B$ such that $\bar\mu(a\Bcap(a_n\Bsymmdiff c))\le 2^{-n}$.
Set $b=\inf_{n\in\Bbb N}\sup_{m\ge n}a_m\in\frak C$;  then
$b\cap a=c\Bcap a$ (apply 323F to $\sequencen{a\Bcap a_n}$).\ \Qed

It follows that $\frak A^f\subseteq\frak C$.   \Prf\ If $c\in\frak A^f$,
then whenever $c\Bsubseteq a\in\frak A^f$ there is a $b_a\in\frak C$
such that $b_a\Bcap a=c$.   Now (because $\bar\mu$ is semi-finite)
$c=\inf\{b_a:c\Bsubseteq a\in\frak A^f\}\in\frak C$.\ \Qed

Finally, again because $\bar\mu$ is semi-finite,

\Centerline{$c=\sup\{a:a\in\frak A^f,\,a\Bsubseteq c\}\in\frak C$}

\noindent for every $c\in\frak A$, and $\frak A=\frak C$.  Thus
$\frak B\,\,\tau$-generates $\frak A$.

\medskip

\qquad\grheadb\ Suppose that $\frak B\,\,\tau$-generates $\frak A$.
Then the topological closure of $\frak B$ is order-closed (323D(c-i)) and
a subalgebra (323B), so must be $\frak A$, and $\frak B$ is topologically
dense.

\medskip

\quad{\bf (ii)}\grheada\
If $\frak A$ is finite, this is trivial, just because
the measure-algebra topology is Hausdorff (323Ga).   So let us
henceforth suppose that $\frak A$ is infinite, so that both
$\tau(\frak A)$ and the topological density $d_{\frak T}(\frak A)$ of
$\frak A$ are infinite.

\medskip

\qquad\grheadb\
Let $A\subseteq\frak A$ be a set with cardinal $\tau(\frak A)$
which $\tau$-generates $\frak A$, and let $\frak B$ be the subalgebra of
$\frak A$ generated by $A$.   Then
$\#(\frak B)=\#(A)=\tau(\frak A)$ (331Gc), and
$\frak B$ is topologically dense in $\frak A$, by (i);  so
$d_{\frak T}(\frak A)\le\tau(\frak A)$.

\medskip

\qquad\grheadc\ Let $A\subseteq\frak A$ be a topologically dense set
with cardinal $d_{\frak T}(\frak A)$, and $\frak B$ the subalgebra of
$\frak A$ generated by $A$.   Then $\frak B$ is topologically dense, so it
$\tau$-generates $\frak A$, and

\Centerline{$\tau(\frak A)\le\#(\frak B)=\#(A)=d_{\frak T}(\frak A)$;}

\noindent with ($\beta$), this means that we have equality, as claimed.

\medskip

{\bf (b)(i)} The case of finite $\frak A$ is again trivial;  suppose that
$\frak A$ is infinite.   Let $\familyiI{a_i}$ be a partition of unity in
$\frak A$ consisting of non-zero elements of finite measure.

\medskip

\quad{\bf (ii)} The topological density $d_{\text{top}}(\frak A^f)$
is at most $\max(c(\frak A),\tau(\frak A))$.
\Prf\ For each $i$, the topological density of $\frak A_{a_i}$, with its
measure-algebra topology, is at most
$\max(\omega,\tau(\frak A_{a_i}))\le\tau(\frak A)$ ((a) above and 514Ed);
let $B_i\subseteq\frak A_{a_i}$ be a dense subset of this size or less.
Set $B=\bigcup_{i\in I}B_i$, $D=\{\sup B':B'\in[B]^{<\omega}\}$.
Then the metric closure $\overline{D}$ of $D$ in $\frak A^f$ is closed
under $\Bcup$ and includes $\frak A_{a_i}$ for every $i$.   If now
$a\in\frak A^f$, $a=\sup_{i\in I}a\Bcap a_i\in\overline{D}$.   So

\Centerline{$d_{\text{top}}(\frak A^f)\le\#(D)
\le\max(\omega,\#(I),\tau(\frak A))
\le\max(c(\frak A),\tau(\frak A))$.  \Qed}

\medskip

\quad{\bf (iii)} $c(\frak A)\le d_{\text{top}}(\frak A^f)$.   \Prf\ Let
$\family{j}{J}{b_j}$
be any disjoint family in $\frak A^+$.
For each $j$, let $b_j'\Bsubseteq b_j$ be a non-zero element of non-zero
finite measure.   Set
$G_j=\{a:a\in\frak A^f$, $\bar\mu(a\Bsymmdiff b'_j)<\bar\mu b'_j\}$ for
$j\in J$.   Then $\family{j}{J}{G_j}$ is a disjoint family of non-empty
open sets in $\frak A^f$, so $\#(J)\le d_{\text{top}}(\frak A^f)$ (5A4Ba).   As
$\family{j}{J}{b_j}$ is arbitrary,
$c(\frak A)\le d_{\text{top}}(\frak A^f)$.\ \Qed

\medskip

\quad{\bf (iv)} $\tau(\frak A)\le d_{\text{top}}(\frak A^f)$.   \Prf\
Let $A\subseteq\frak A^f$ be a dense set with cardinal
$d_{\text{top}}(\frak A^f)$.  Let $\frak B$ be the order-closed
subalgebra of $\frak A$ generated by $B=A\cup\{a_i:i\in I\}$.   For any
$i\in I$, set
$A_i=\{a\Bcap a_i:a\in A\}$.   Now $A_i$ is topologically dense in
$\frak A_{a_i}$ (use 3A3Eb), so the order-closed subalgebra of
$\frak A_{a_i}$ it
generates is the whole of $\frak A_{a_i}$ (323H);  by 314H,
$\frak A_{a_i}=\{b\Bcap a_i:b\in\frak B\}$.   As
$a_i\in A\subseteq\frak B$, $\frak B$ includes $\frak A_{a_i}$.   As
$\sup_{i\in I}a_i=1$, $\frak B=\frak A$.   Thus

\Centerline{$\tau(\frak A)\le\#(B)
\le\max(\omega,\#(I),d_{\text{top}}(\frak A^f))
=d_{\text{top}}(\frak A^f)$}

\noindent (using (iii) for the last equality).\ \Qed

\medskip

\quad{\bf (v)} Putting these together, we have the result.
}%end of proof of 521E

\leader{521F}{Proposition} Let $(X,\Sigma,\mu)$ be a measure space,
$A$ a subset of $X$ and $\mu_A$ the subspace measure on $A$.

(a) $\Cal N(\mu_A)\prT\Cal N(\mu)$, so
$\add\Cal N(\mu_A)\ge\add\Cal N(\mu)$ and
$\cf\Cal N(\mu_A)\le\cf\Cal N(\mu)$.

(b) $(A,\in,\penalty-100\Cal N(\mu_A))\prGT(X,\in,\Cal N(\mu))$, so
$\non\Cal N(\mu_A)\ge\non\Cal N(\mu)$ and
$\cov\Cal N(\mu_A)\le\cov\Cal N(\mu)$.

(c) $\add\mu_A\ge\add\mu$.

(d) $\shr\Cal N(\mu_A)\le\shr\Cal N(\mu)$
and $\shr^+\Cal N(\mu_A)\le\shr^+\Cal N(\mu)$.

(e) If either $A\in\Sigma$ or $(X,\Sigma,\mu)$
has locally determined negligible sets, $\pi(\mu_A)\le\pi(\mu)$.

(f) If $\mu_A$ is semi-finite, then
$\tau(\mu_A)\le\tau(\mu)$.

\proof{{\bf (a)}
Because $\Cal N(\mu_A)=\Cal PA\cap\Cal N(\mu)$ (214Cb),
the embedding $\Cal N(\mu_A)\embedsinto\Cal N(\mu)$ is a
Tukey function, and $\Cal N(\mu_A)\prT\Cal N(\mu)$.
By 513Ee, $\add\Cal N(\mu_A)\ge\add\Cal N(\mu)$ and
$\cf\Cal N(\mu_A)\le\cf\Cal N(\mu)$.

\medskip

{\bf (b)} Next, setting $\phi(x)=x$ for $x\in A$ and $\psi(F)=F\cap A$ for
$F\in\Cal N(\mu)$, $(\phi,\psi)$ witnesses that
$(A,\in,\Cal N(\mu_A))\prGT(X,\in,\Cal N(\mu))$.
By 512D and 512Ed,

\Centerline{$\non\Cal N(\mu_A)=\add(A,\in,\Cal N(\mu_A))
\ge\add(X,\in,\Cal N(\mu))=\non\Cal N(\mu)$,}

\Centerline{$\cov\Cal N(\mu_A)=\cov(A,\in,\Cal N(\mu_A))
\le\cov(X,\in,\Cal N(\mu))=\cov\Cal N(\mu)$.}

\medskip

{\bf (c)} If $\ofamily{\xi}{\kappa}{F_{\xi}}$ is a disjoint family in
$\Sigma_A=\dom\mu_A$, where $\kappa<\add\mu$, then for each $\xi<\kappa$ we
have an $E_{\xi}\in\Sigma$ such that $F_{\xi}=A\cap E_{\xi}$ and
$\mu_AF_{\xi}=\mu E_{\xi}$ (214Ca).   Set
$E'_{\xi}=E_{\xi}\setminus\bigcup_{\eta<\xi}E_{\eta}$ for $\xi<\kappa$;
then $E'_{\xi}\in\Sigma$ for each $\xi$, and
$\ofamily{\xi}{\kappa}{E'_{\xi}}$ is disjoint.   Set
$E=\bigcup_{\xi<\kappa}E'_{\xi}=\bigcup_{\xi<\kappa}E_{\xi}$ and
$F=A\cap E=\bigcup_{\xi<\kappa}F_{\xi}$.   Then

\Centerline{$\sum_{\xi<\kappa}\mu_AF_{\xi}
\le\mu_AF
\le\mu E
=\sum_{\xi<\kappa}\mu E'_{\xi}
\le\sum_{\xi<\kappa}\mu E_{\xi}
=\sum_{\xi<\kappa}\mu_AF_{\xi}$,}

\noindent and we have equality.   As
$\ofamily{\xi}{\kappa}{F_{\xi}}$ is arbitrary, $\add\mu_A\ge\add\mu$.

\medskip

{\bf (d)}
If $B\in\Cal PA\setminus\Cal N(\mu_A)$, there is a $C\subseteq B$ such that
$C\notin\Cal N(\mu)$ and $\#(C)\le\shr\Cal N(\mu)$ (resp.\
$\#(C)<\shr^+\Cal N(\mu))$;  now
$C\notin\Cal N(\mu_A)$;  as $B$ is arbitrary,
$\shr\Cal N(\mu_A)\le\shr\Cal N(\mu)$ (resp.\
$\shr^+\Cal N(\mu_A)\le\shr^+\Cal N(\mu)$).

\medskip

{\bf (e)} Let $\Cal H\subseteq\Sigma\setminus\Cal N(\mu)$ be a coinitial
set with cardinal $\pi(\mu)$.   Set
$\Cal G=\{H\cap A:H\in\Cal H\}\setminus\Cal N(\mu)$.   Then
$\mu_AG$ is defined and non-zero for every $G\in\Cal G$.
Now $\Cal G$ is coinitial with $\dom\mu_A\setminus\Cal N(\mu_A)$.
\Prf\ If
$\mu_AB>0$, there is an $E\in\Sigma$ such that $B=E\cap A$.   If
$A\in\Sigma$, then $B\in\Sigma$ and there is an $H\in\Cal H$ such that
$H\subseteq B$, while of course $H\in\Cal G$.   If $(X,\Sigma,\mu)$ has
locally determined negligible sets, then, as in the proof of 521Dc,
there is a non-negligible set $F\in\Sigma$ which is a measurable envelope
of a subset of $B$.   Now there is an $H\in\Cal H$ included in
$F\cap E$, in which case $H\cap A$ is included in $B$
and belongs to $\Cal G$.\ \QeD\   So

\Centerline{$\pi(\mu_A)\le\#(\Cal G)\le\#(\Cal H)=\pi(\mu)$.}

\medskip

{\bf (f)} Writing $\frak A$, $\frak A_A$ for the measure algebras of
$\mu$ and $\mu_A$, we have a Boolean homomorphism
$\pi:\frak A\to\frak A_A$ defined by saying that
$\pi E^{\ssbullet}=(F\cap A)^{\ssbullet}$ for every $E\in\Sigma$.
(The point is just that $F\cap A\in\Cal N(\mu_A)$ whenever
$F\in\Cal N(\mu)$.)   Now $\pi$ is order-continuous.   \Prf\ Suppose that
$C\subseteq\frak A$ is non-empty and downwards-directed and
$\inf C=0$ in $\frak A$.   \Quer\ If
$b\in\frak A_A$ is a non-zero lower bound of $\pi[C]$, then, because
$\nu_A$ is semi-finite, there is a $G\in\dom\mu_A$ such that
$0<\mu_AG<\infty$ and $G^{\ssbullet}\Bsubseteq b$.   Let $E\in\Sigma$ be
such that $G=E\cap A$ and $\mu E=\mu_AG$ (214Ca).   Then $E^{\ssbullet}$
cannot be a lower bound of $C$;  let $a\in C$ be such that
$E^{\ssbullet}\Bsetminus a\ne 0$.   In this case, there is an $F\in\Sigma$
such that $F\subseteq E$ and
$F^{\ssbullet}=E^{\ssbullet}\Bsetminus a$, so that
$\pi F^{\ssbullet}$ is disjoint from $\pi a\Bsupseteq b$, and
$F\cap G=(F\cap A)\cap G$ must be negligible.   We know that $\mu F>0$,
so $\mu(E\setminus F)<\mu E$;  but also $G\setminus(E\setminus F)$ is
negligible, so

\Centerline{$\mu_AG=\mu^*G\le\mu(E\setminus F)<\mu E=\mu_AG$.  \Bang}

\noindent It follows that $\inf\pi[C]=0$ in $\frak A_A$;  as
$C$ is arbitrary, $\pi$ is order-continuous.\ \Qed

Now let $B\subseteq\frak A$ be such that $B\,\,\tau$-generates $\frak A$
and $\#(B)=\tau(\frak A)$.   Writing $\frak B$ for the
order-closed subalgebra of $\frak A_A$ generated by $\pi[B]$,
we see that $\pi^{-1}[\frak B]$ is an order-closed subalgebra of
$\frak A$ including $B$, so must be the whole of $\frak A$, and
$\frak A_A=\pi[\frak A]=\frak B$.   Accordingly

\Centerline{$\tau(\mu_A)=\tau(\frak A_A)\le\#(\pi[B])
\le\#(B)=\tau(\frak A)=\tau(\mu)$,}

\noindent as claimed.
}%end of proof of 521F

\leader{521G}{Proposition} Let $\familyiI{(X_i,\Sigma_i,\mu_i)}$ be a
non-empty
family of measure spaces with direct sum $(X,\Sigma,\mu)$.   Then

\Centerline{$\add\Cal N(\mu)=\min_{i\in I}\add\Cal N(\mu_i)$,
\quad$\add\mu=\min_{i\in I}\add\mu_i$,}

\Centerline{$\cov\Cal N(\mu)=\sup_{i\in I}\cov\Cal N(\mu_i)$,
\quad$\non\Cal N(\mu)=\min_{i\in I}\non\Cal N(\mu_i)$,}

\Centerline{$\shr\Cal N(\mu)=\sup_{i\in I}\shr\Cal N(\mu_i)$,
\quad$\shr^+\Cal N(\mu)=\sup_{i\in I}\shr^+\Cal N(\mu_i)$,}

\Centerline{$\tau(\mu)\le\max(\omega,\sup_{i\in I}\tau(\mu_i),
\min\{\lambda:\#(I)\le 2^{\lambda}\})$\dvAnew{2015}}

\noindent and $\pi(\mu)$ is the cardinal sum
$\sum_{i\in I}\pi(\mu_i)$.   If $I$ is finite, then

\Centerline{$\cf\Cal N(\mu)=\max_{i\in I}\cf\Cal N(\mu_i)$.}

\proof{ Concerning each of
$\add\Cal N(\mu)$, $\add\mu$, $\cov\Cal N(\mu)$,
$\non\Cal N(\mu)$, $\shr\Cal N(\mu)$ and $\shr^+\Cal N(\mu)$,
521F provides an inequality in one direction.   The reverse
inequalities are equally straightforward, especially if we note that
$\Cal N(\mu)\cong\prod_{i\in I}\Cal N(\mu_i)$, so that
512Hc is relevant.
For $\tau(\mu)$, 514Ef provides the formula at once if we note that the
measure algebra of $\mu$ is isomorphic to
the simple product of the measure algebras of
the $\mu_i$ (cf.\ 322M). %\query

As for $\pi(\mu)$, if for each $i\in I$ we choose a coinitial set
$\Cal H_i$ of $\Sigma_i\setminus\Cal N(\mu_i)$ with cardinal $\pi(\mu_i)$, then

\Centerline{$\Cal H=\{H\times\{i\}:i\in I$, $H\in\Cal H_i\}$}

\noindent is coinitial with $\Sigma\setminus\Cal N(\mu)$ and witnesses that
$\pi(\mu)\le\sum_{i\in I}\pi(\mu_i)$.   (As in 214L, I am thinking of
$X$ as $\bigcup_{i\in I}X_i\times\{i\}$.)   Conversely, if $\Cal H$ is
coinitial with $\Sigma\setminus\Cal N(\mu)$ and for each $i\in I$ we set
$\Cal H_i=\{H:H\times\{i\}\in\Cal H\}$, we shall have $\Cal H_i$ coinitial
with $\Sigma_i\setminus\Cal N(\mu_i)$, so that

\Centerline{$\sum_{i\in I}\pi(\mu_i)\le\sum_{i\in I}\#(\Cal H_i)
\le\#(\Cal H)=\pi(\mu)$}

\noindent and $\pi(\mu)=\sum_{i\in I}\pi(\mu_i)$.
}%end of proof of 521G

\leader{521H}{Proposition} Let $(X,\Sigma,\mu)$ and $(Y,\Tau,\nu)$ be
measure spaces, and $f:X\to Y$ an \imp\ function.

(a)(i)
$(X,\in,\Cal N(\mu))\discretionary{}{}{}\prGT(Y,\in,\Cal N(\nu))$, so
$\non\Cal N(\mu)\ge\non\Cal N(\nu)$ and
$\cov\Cal N(\mu)\le\cov\Cal N(\nu)$.

\quad(ii) If there is a topology on $Y$ such that $\nu$ is a
topological measure
inner regular with respect to the closed sets, then $\pi(\nu)\le\pi(\mu)$.

\quad(iii) If $\nu$ is $\sigma$-finite, then $\tau(\nu)\le\tau(\mu)$.

(b) If $\nu$ is the image measure $\mu f^{-1}$, then
$\add\nu\ge\add\mu$.   If, moreover, $\mu$ is complete,
$\Cal N(\nu)\prT\Cal N(\mu)$, so $\add\Cal N(\mu)\le\add\Cal N(\nu)$ and
$\cf\Cal N(\mu)\ge\cf\Cal N(\nu)$;  also
$\shr\Cal N(\mu)\ge\shr\Cal N(\nu)$
and $\shr^+\Cal N(\mu)\ge\shr^+\Cal N(\nu)$.

\proof{{\bf (a)(i)} Set $\psi(F)=f^{-1}[F]$ for
$F\in\Cal N(\nu)$.   Then $(f,\psi)$ is a Galois-Tukey connection from
$(X,\in,\Cal N(\mu))$ to $(Y,\in,\Cal N(\nu))$, so

\Centerline{$\cov\Cal N(\mu)=\cov(X,\in,\Cal N(\mu))
\le\cov(Y,\in,\Cal N(\nu))=\cov\Cal N(\nu)$,}

\Centerline{$\non\Cal N(\mu)=\add(X,\in,\Cal N(\mu))
\ge\add(Y,\in,\Cal N(\nu))=\non\Cal N(\nu)$}

\noindent (512D, 512Ed again).

\medskip

\quad{\bf (ii)} Let $\Cal H$ be a coinitial subset of
$\Sigma\setminus\Cal N(\mu)$ with cardinal $\pi(\mu)$.   Set
$\Cal G=\{\overline{f[H]}:H\in\Cal H\}$.   Because $\nu$ is a topological
measure, $\Cal G\subseteq\Tau$;  and if $H\in\Cal H$, then

\Centerline{$\nu\overline{f[H]}=\mu f^{-1}[\overline{f[H]}]
\ge\mu H>0$,}

\noindent so $\Cal G\subseteq\Tau\setminus\Cal N(\nu)$.   If
$F\in\Tau\setminus\Cal N(\nu)$, there is a closed set $F'\subseteq F$ such
that $0<\nu F'=\mu f^{-1}[F']$;  there is an $H\in\Cal H$ such that
$H\subseteq f^{-1}[F']$;  now $G=\overline{f[H]}$ belongs to $\Cal G$ and
is included in $F'\subseteq F$.   So $\Cal G$ is coinitial with
$\Tau\setminus\Cal N(\nu)$ and

\Centerline{$\pi(\nu)\le\#(\Cal G)\le\#(\Cal H)=\pi(\mu)$.}

\medskip

\quad{\bf (iii)} Let $(\frak A,\bar\mu)$ and $(\frak B,\bar\nu)$ be the
measure algebras of $\mu$, $\nu$ respectively.   Then we have a
sequentially order-continuous
measure-preserving Boolean homomorphism $\pi:\frak B\to\frak A$ defined by
setting $\pi F^{\ssbullet}=f^{-1}[F]^{\ssbullet}$ for every $F\in\dom\nu$
(324M).   If $\frak A$ is finite then $\frak B$ must be finite with
$\#(\frak B)\le\#(\frak A)$, and consequently
$\tau(\frak B)\le\tau(\frak A)$ (331Xc, or otherwise).
So let us suppose that $\frak A$ is infinite.

Writing $\frak A^f$, $\frak B^f$ for the respective ideals of
elements of finite measure, $\pi\restr\frak B^f$ is a function from
$\frak B^f$ to $\frak A^f$ which is an isometry for the measure metrics on
$\frak B^f$ and $\frak A^f$.   So the topological density
$d_{\text{top}}(\frak B^f)$ is
equal to $d_{\text{top}}(\pi[\frak B^f])$ and less than or equal to
$d_{\text{top}}(\frak A^f)$ (5A4B(h-ii)).

\woddheader{521H}{0}{0}{0}{72pt}
Observe next that $(\frak B,\bar\nu)$ is $\sigma$-finite because $\nu$ is,
and that $(\frak A,\bar\mu)$ therefore also is (324Kd).   So we get

$$\eqalignno{\tau(\nu)
&=\tau(\frak B)
\le\max(\omega,c(\frak B),\tau(\frak B))
=\max(\omega,d_{\text{top}}(\frak B^f))\cr
\displaycause{521Eb}
&\le\max(\omega,d_{\text{top}}(\frak A^f))
=\max(\omega,c(\frak A),\tau(\frak A))
=\tau(\frak A)
=\tau(\mu),\cr}$$

\noindent as required.

\medskip

{\bf (b)} If $\ofamily{\xi}{\kappa}{F_{\xi}}$ is a disjoint family in
$\Tau$, where $\kappa<\add\mu$, then
$\ofamily{\xi}{\kappa}{f^{-1}[F_{\xi}]}$ is a disjoint family in $\Sigma$,
so

$$\nu(\bigcup_{\xi<\kappa}F_{\xi})
=\mu f^{-1}[\bigcup_{\xi<\kappa}F_{\xi}]
=\mu(\bigcup_{\xi<\kappa}f^{-1}[F_{\xi}]
=\sum_{\xi<\kappa}\mu f^{-1}[F_{\xi}]
=\sum_{\xi<\kappa}\nu F_{\xi}.$$

\noindent
As $\ofamily{\xi}{\kappa}{F_{\xi}}$ is arbitrary, $\add\nu\ge\add\mu$.

Now suppose that $\mu$ is complete.   In this case, $F\in\Tau$
whenever $F\subseteq Y$ and $f^{-1}[F]\in\Cal N(\mu)$, so that
$\Cal N(\nu)$ is precisely
$\{F:F\subseteq Y,\,f^{-1}[F]\in\Cal N(\mu)\}$.   It is now easy to
check that $F\mapsto f^{-1}[F]:\Cal N(\nu)\to\Cal N(\mu)$ is a Tukey
function.   So $\add\Cal N(\nu)\ge\add\Cal N(\mu)$
and $\cf\Cal N(\nu)\le\cf\Cal N(\mu)$, by 513Ee again.

Take any non-negligible $A\subseteq Y$.   Then
$f^{-1}[A]\notin\Cal N(\mu)$, so there is a set $B\subseteq f^{-1}[A]$
such that $\#(B)\le\shr\Cal N(\mu)$ and $B\notin\Cal N(\mu)$.   In this
case, $f[B]\subseteq A$, $f[B]\notin\Cal N(\nu)$ and
$\#(f[B])\le\shr\Cal N(\mu)$.   As $A$ is arbitrary,
$\shr\Cal N(\nu)\le\shr\Cal N(\mu)$.
The same argument, with $<$ instead of $\le$ at appropriate points, shows
that $\shr^+\Cal N(\nu)\le\shr^+\Cal N(\mu)$.
}%end of proof of 521H

\leader{521I}{Corollary} Let $(X,\Sigma,\mu)$ be an atomless strictly
localizable measure space.   Then $\non\Cal N(\mu)\ge\non\Cal N$ and
$\cov\Cal N(\mu)\le\cov\Cal N$, where $\Cal N$ is the null ideal of
Lebesgue measure on $\Bbb R$.

\proof{{\bf (a)} If $\mu X=0$ this is trivial.

\medskip

{\bf (b)} If $0<\mu X<\infty$, let $\nu$ be the completion of the
normalized measure
$\Bover1{\mu X}\mu$.   Then $\nu$ is complete and atomless,
so by 343Cb there is an \imp\ function from $(X,\nu)$ to
$([0,1],\mu_1)$, where $\mu_1$ is Lebesgue measure on $[0,1]$.
Also $\Cal N(\nu)=\Cal N(\mu)$.   By 521Ha,
$\non\Cal N(\nu)\ge\non\Cal N(\mu_1)$ and
$\cov\Cal N(\nu)\le\cov\Cal N(\mu_1)$.
Now $([0,1],\Cal N(\mu_1))$ is isomorphic to
$(\Bbb R,\Cal N)$.   \Prf\ Take a bijection $h:\Bbb R\to[0,1]$
such that $h(x)=\bover12(1+\tanh x)$ for $x\in\Bbb R\setminus\Bbb Q$;  then
$h$ is a suitable isomorphism.\ \QeD\
So $\non\Cal N(\mu)=\non\Cal N(\nu)\ge\non\Cal N$ and
$\cov\Cal N(\mu)\le\cov\Cal N$.

\medskip

{\bf (c)} If $X$ has infinite measure, let $\familyiI{X_i}$ be a
decomposition of $X$ into sets of finite measure.   For each $i\in I$
let $\mu_i$ be the subspace measure on $X_i$.
Then every $\mu_i$ is atomless, so, putting (b) and 521G together,

\Centerline{$\non\Cal N(\mu)
=\min_{i\in I}\non\Cal N(\mu_i)\ge\non\Cal N$,}

\Centerline{$\cov\Cal N(\mu)
=\sup_{i\in I}\cov\Cal N(\mu_i)\le\cov\Cal N$.}
}%end of proof of 521I

\leader{521J}{}\cmmnt{ For product spaces the situation is more
complicated, because the product measure introduces `new' negligible sets
which are not directly definable in terms of the null ideals of the
factors.   In the next three sections, however, we shall find out
quite a lot about the cardinal functions of Radon measures,
and this information, when it comes, can be used to give results about
general products of probability measures.

\medskip

\noindent}{\bf Proposition} Let $\familyiI{(X_i,\Sigma_i,\mu_i)}$ be a
non-empty family of probability spaces with product $(X,\Sigma,\mu)$.

(a)

\Centerline{$\non\Cal N(\mu)\ge\sup_{i\in I}\non\Cal N(\mu_i)$,
\quad$\cov\Cal N(\mu)\le\min_{i\in I}\cov\Cal N(\mu_i)$,}

\Centerline{$\add\mu=\add\Cal N(\mu)\le\min_{i\in I}\add\Cal N(\mu_i)$,
\quad$\cf\Cal N(\mu)\ge\sup_{i\in I}\cf\Cal N(\mu_i)$,}

\Centerline{$\shr\Cal N(\mu)\ge\sup_{i\in I}\shr\Cal N(\mu_i)$,
\quad$\shr^+\Cal N(\mu)\ge\sup_{i\in I}\shr^+\Cal N(\mu_i)$,}

\Centerline{$\pi(\mu)\ge\sup_{i\in I}\pi(\mu_i)$.}

(b) Set $\kappa=\#(\{i:i\in I$, $\Sigma_i\ne\{\emptyset,X_i\}\}$.
Then $[\kappa]^{\le\omega}\prT\Cal N(\mu)$;  consequently
$\add\mu=\add\Cal N(\mu)$ is $\omega_1$ if $\kappa$ is uncountable,
while $\cf\Cal N(\mu)$ is at least $\cff[\kappa]^{\le\omega}$.

(c) Now suppose that $I$ is countable and that
we have for each $i\in I$ a probability space
$(Y_i,\Tau_i,\nu_i)$ and an \imp\ function
$f_i:X_i\to Y_i$ which represents an isomorphism of the measure algebras of
$\mu_i$ and $\nu_i$.   Let $(Y,\Tau,\nu)$ be the product of
$\familyiI{(Y_i,\Tau_i,\nu_i)}$.   Then

\Centerline{$\Cal N(\mu)\prT\Cal N(\nu)\times\prod_{i\in I}\Cal N(\mu_i)$.}

\noindent Consequently

\Centerline{$\add\Cal N(\mu)
\ge\min(\add\Cal N(\nu),\min_{i\in I}\add\Cal N(\mu_i))$,}

\noindent and if $I$ is finite

\Centerline{$\cf\Cal N(\mu)
\le\max(\cf\Cal N(\nu),\max_{i\in I}\cf\Cal N(\mu_i))$.}

(d) If $I$ is finite, then

\Centerline{$\non\Cal N(\mu)=\max_{i\in I}\non\Cal N(\mu_i)$,
\quad$\cov\Cal N(\mu)=\min_{i\in I}\cov\Cal N(\mu_i)$.}

\proof{{\bf (a)} Note that $\add\mu=\add\Cal N(\mu)$ by 521Ad.   Now
with one exception the inequalities are immediate if we
apply 521H to the canonical maps from $X$ to $X_i$.   The odd one out is
the last, because we do not have a simple general result concerning the
$\pi$-weight of an image measure.   But in the present case we can argue as
follows.   Let $\Cal H\subseteq\Sigma\setminus\Cal N(\mu)$ be a coinitial
set with cardinal $\pi(\mu)$, and take $i\in I$.   Then we can identify
$(X,\Sigma,\mu)$ with $(X',\Sigma',\mu')\times(X_i,\Sigma_i,\mu_i)$
where $(X',\Sigma',\mu')$ is the product of the family
$\langle(X_j,\Sigma_j,\mu_j)\rangle_{j\in I,j\ne i}$ (254N).
If $H\in\Cal H$, $\mu(X\setminus H)
=\int\mu_i^*(X_i\setminus H[\{x'\}])\mu'(dx')$ (252D) is less than
$1$, so there is an $x'_H\in X'$ such that
$\mu_i^*(X_i\setminus H[\{x'_H\}])<1$, $(\mu_i)_*H[\{x'_H\}]>0$
and there is a
$G_H\in\Sigma_i\setminus\Cal N(\mu_i)$ such that
$G_H\subseteq H[\{x'_H\}]$.   Set $\Cal G=\{G_H:H\in\Cal H\}$.   If
$\mu_iF>0$, then $\mu(X'\times F)>0$ and there is
an $H\in\Cal H$ included in $X'\times F$;  in which case
$G_H\subseteq H[\{x'_H\}]\subseteq F$.   So $\Cal G$ is coinitial with
$\Sigma_i\setminus\Cal N(\mu_i)$ and

\Centerline{$\pi(\mu_i)\le\#(\Cal G)\le\#(\Cal H)\le\pi(\mu)$.}

\noindent As $i$ is arbitrary, $\sup_{i\in I}\pi(\mu_i)\le\pi(\mu)$, as
claimed.

\medskip

{\bf (b)(i)} If $\kappa\le\omega$ then the constant function with value
$\emptyset$ is a Tukey function from $[\kappa]^{\le\omega}$ to
$\Cal N(\mu)$.   Otherwise, set
$J=\{i:i\in I$, $\Sigma_i\ne\{\emptyset,X_i\}\}$ and for $i\in J$ choose
a non-empty $C_i\in\Sigma_i$ such that $\mu_iC_i\le\bover12$.   Index $J$
as $\langle i_{\xi n}\rangle_{\xi<\kappa,n\in\Bbb N}$ and for $\xi<\kappa$
set $E_{\xi}=\{x:x(i_{\xi n})\in C_{i_{\xi n}}$ for every $n\in\Bbb N\}$,
so that $\mu E_{\xi}=\prod_{n\in\Bbb N}\mu_{i_{\xi n}}C_{i_{\xi n}}=0$.
Define $\phi:[\kappa]^{\le\omega}\to\Cal N(\mu)$ by setting
$\phi K=\bigcup_{\xi\in K}E_{\xi}$ for countable $K\subseteq\kappa$.
Then $\phi$ is a Tukey function.   \Prf\ If $E\in\Cal N(\mu)$, there is a
negligible $E'\supseteq E$ which is determined by coordinates in a
countable set $I'$ (254Oc).   Set
$L=\{\xi:\xi<\kappa$, $i_{\xi n}\in I'$ for some $n\in\Bbb N\}$;  then $L$
is countable.   If $\xi<\kappa$ and $E_{\xi}\subseteq E'$, $E_{\xi}$ is
determined by coordinates in $\{i_{\xi n}:n\in\Bbb N\}$;  as neither
$E_{\xi}$ nor $X\setminus E'$ is empty, this must meet $I'$, and
$\xi\in L$.   So $\{K:K\in[\kappa]^{\le\omega}$, $\phi K\subseteq E\}$ is
bounded above by $L\in[\kappa]^{\le\omega}$.   As $E$ is arbitrary, $\phi$
is a Tukey function.\ \QeD\   Accordingly
$[\kappa]^{\le\omega}\prT\Cal N(\mu)$.

\medskip

\quad{\bf (ii)} It follows that
$\add\Cal N(\mu)\le\add[\kappa]^{\le\omega}\le\omega_1$ if $\kappa$ is
uncountable, and that $\cf\Cal N(\mu)\ge\cff[\kappa]^{\le\omega}$.

\medskip

{\bf (c)(i)} Recall that any \imp\ function $f$ between measure spaces
induces a measure-preserving Boolean homomorphism
$F^{\ssbullet}\mapsto(f^{-1}[F])^{\ssbullet}$ between the measure algebras
(324M).
For $i\in I$ and $C\in\Sigma_i$ choose $\psi_i(C)\in\Tau_i$ such that
$(f_i^{-1}[\psi_i(C)])^{\ssbullet}=C^{\ssbullet}$ in the measure algebra of
$\mu_i$, that is, $C\symmdiff f_i^{-1}[\psi_i(C)]\in\Cal N(\mu_i)$.
Next, for $E\in\Cal N(\mu)$ and $n\in\Bbb N$, choose a family
$\langle C_{Enmi}\rangle_{m\in\Bbb N,i\in I}$ such that
$C_{Enmi}\in\Sigma_i$ for every $m\in\Bbb N$ and $i\in I$,
$E\subseteq\bigcup_{m\in\Bbb N}\prod_{i\in I}C_{Enmi}$, and
$\sum_{m\in\Bbb N}\prod_{i\in I}\mu_iC_{Enmi}\le 2^{-n}$;  set

\Centerline{$\phi(E)
=\bigl(\bigcap_{n\in\Bbb N}\bigcup_{m\in\Bbb N}
   \prod_{i\in I}\psi_i(C_{Enmi}),
\familyiI{\bigcup_{m,n\in\Bbb N}
(C_{Enmi}\setminus f_i^{-1}[\psi_i(C_{Enmi})])}\bigr)$.}

\noindent Because

$$\eqalign{
\nu(\bigcap_{n\in\Bbb N}\bigcup_{m\in\Bbb N}\prod_{i\in I}\psi_i(C_{Enmi}))
&\le\inf_{n\in\Bbb N}\sum_{m\in\Bbb N}\prod_{i\in I}
   \nu_i\psi_i(C_{Enmi})\cr
&=\inf_{n\in\Bbb N}\sum_{m\in\Bbb N}\prod_{i\in I}
   \mu_iC_{Enmi}
=0,\cr}$$

\noindent $\phi$ is a function from $\Cal N(\nu)$ to
$\Cal N(\nu)\times\prod_{i\in I}\Cal N(\mu_i)$.

Now $\phi$ is a Tukey function.   \Prf\ Suppose that $W\in\Cal N(\nu)$ and
that $E_i\in\Cal N(\mu_i)$ for every $i\in I$.   Define $f:X\to Y$ by
setting $f(x)=\familyiI{f_i(x(i))}$ for $x\in X$;
then $f$ is \imp\ (254H).   So
$V=f^{-1}[W]\cup\bigcup_{i\in I}\{x:x(i)\in E_i\}$ is negligible.   (This
is where we need to know that $I$ is countable.)   Suppose that
$E\in\Cal N(\mu)$ is such that
$\phi(E)\le(W,\familyiI{E_i})$;  take $x\in E$ such that
$x(i)\notin E_i$ for every $i\in I$, and $n\in\Bbb N$.   Then there is an
$m\in\Bbb N$ such that $x\in\prod_{i\in I}C_{Enmi}$.   For each $i\in I$,

\Centerline{$C_{Enmi}\setminus f_i^{-1}[\psi_i(C_{Enmi})]\subseteq E_i$,
\quad$x(i)\in C_{Enmi}\setminus E_i$,}

\noindent so $f_i(x(i))\in\psi_i(C_{Enmi})$;  thus
$f(x)\in\prod_{i\in I}\psi_i(C_{Enmi})$.   As $n$ is arbitrary,

\Centerline{$f(x)
\in\bigcap_{n\in\Bbb N}\bigcup_{m\in\Bbb N}\prod_{i\in I}\psi_i(C_{Enmi})
\subseteq W$}

\noindent and $x\in V$.   As $x$ is arbitrary, $E\subseteq V$.   As
$(W,\familyiI{E_i})$ is arbitrary, $\phi$ is a Tukey function.\ \Qed

So $\Cal N(\mu)\prT\Cal N(\nu)\times\prod_{i\in I}\Cal N(\mu_i)$.

\medskip

\quad{\bf (ii)} Accordingly

\Centerline{$\add\Cal N(\mu)
\ge\add(\Cal N(\nu)\times\prod_{i\in I}\Cal N(\mu_i))
=\min(\add\Cal N(\nu),\min_{i\in I}\add\Cal N(\mu_i))$}

\noindent and

\Centerline{$\cf\Cal N(\mu)
\le\cf(\Cal N(\nu)\times\prod_{i\in I}\Cal N(\mu_i))
=\max(\cf\Cal N(\nu),\max_{i\in I}\cf\Cal N(\mu_i))$}

\noindent if $I$ is finite.

\medskip

{\bf (d)(i)} For each $i\in I$ let $A_i\subseteq X_i$ be a non-negligible
set with cardinal $\non\Cal N(\mu_i)$.   Then $A=\prod_{i\in I}A_i$ is not
negligible (251Wm), while
$\#(A)\le\max(\omega,\max_{i\in I}\non\Cal N(\mu_i))$.   If all the
$\non\Cal N(\mu_i)$ are finite, then they are all equal to $1$, and $A$ is
a singleton.   So we must in any case have
$\#(A)=\max_{i\in I}\non\Cal N(\mu_i)$, and
$\non\Cal N(\mu)\le\max_{i\in I}\non\Cal N(\mu_i)$.   By (a), we have
equality.

\medskip

\quad{\bf (ii)} Suppose that $I=\{0,1\}$, and that $\Cal E$ is a cover of
$X=X_0\times X_1$ by negligible sets.   For each $E\in\Cal E$, set
$C_E=\{x:x\in X_0$, $E[\{x\}]\notin\Cal N(\mu_1)\}$;  then $C_E$ is
negligible.   If $\#(\Cal E)<\cov\Cal N(\mu_0)$, then there is an
$x\in X_0\setminus\bigcup_{E\in\Cal E}C_E$;  in which case
$\{E[\{x\}]:E\in\Cal E\}$ witnesses that $\cov\Cal N(\mu_1)\le\#(\Cal E)$.
So $\#(\Cal E)$ must be at least
$\min(\cov\Cal N(\mu_0),\cov\Cal N(\mu_1))$.   As $\Cal E$ is arbitrary,
$\cov\Cal N(\mu)\ge\min(\cov\Cal N(\mu_0),\cov\Cal N(\mu_1))$.

Now an induction on $\#(I)$ (using the associative law 254N) shows that
$\cov\Cal N(\mu)\ge\min_{i\in I}\cov\Cal N(\mu_i)$ whenever $I$ is finite.
Using (a) again, we have equality here also.
}%end of proof of 521J

\cmmnt{\medskip

\noindent{\bf Remark} The simplest applications of (c) here will be when
the $\mu_i$ are \Mth, so that we can take the $\nu_i$ to be the usual
measures on powers $\{0,1\}^{\kappa_i}$ of $\{0,1\}$, and $\nu$ will be
isomorphic to the
usual measure on $\{0,1\}^{\kappa}$ where $\kappa$ is the cardinal sum
$\sum_{i\in I}\kappa_i$.   The cardinal functions of these measures are
dealt with in \S523.   For non-homogeneous $\mu_i$ we shall still be able
to arrange for the $\nu_i$ to be completion regular Radon measures on
dyadic spaces, so that the product measure $\nu$ is again a Radon measure
$Y$ (532F), and (once we have identified its measure algebra -- see 334E,
334Ya) approachable by the methods of \S524.
}%end of comment

\leader{521K}{}\cmmnt{ I turn now to `perfect' and `compact' measure
spaces.    (See \S451 for the basic theory of these.)

\medskip

\noindent}{\bf Proposition} Let $(X,\Sigma,\mu)$ be a perfect
semi-finite measure space which is not purely atomic.   Then

\Centerline{$\add\Cal N(\mu)\le\add\Cal N$,
\quad$\cf\Cal N(\mu)\ge\cf\Cal N$,}

\Centerline{$\shr\Cal N(\mu)\ge\shr\Cal N$,
\quad$\shr^+\Cal N(\mu)\ge\shr^+\Cal N$,
\quad$\pi(\mu)\ge\pi(\mu_L)$}

\noindent where $\Cal N$ is the null ideal of
Lebesgue measure on $\Bbb R$ and $\mu_L$ is Lebesgue measure on $\Bbb R$.

\proof{{\bf (a)} Suppose first that $\mu$ is a complete atomless
probability measure.   Then there is a function $f:X\to[0,1]$ which is
\imp\ for $\mu$ and Lebesgue measure $\mu_1$ on $[0,1]$ (343Cb again);
and in fact
$\mu_1$ is the image measure $\mu f^{-1}$.   \Prf\ By
451O, $\mu f^{-1}$ is
a Radon measure;  since it extends $\mu_1$ it must actually be equal to
$\mu_1$, by 415H.\ \QeD\
So $\add\Cal N(\mu)\le\add\Cal N(\mu_1)$,
$\cf\Cal N(\mu)\ge\cf\Cal N(\mu_1)$,
$\shr\Cal N(\mu)\ge\shr\Cal N(\mu_1)$,
and $\shr^+\Cal N(\mu)\ge\shr^+\Cal N(\mu_1)$ and $\pi(\mu)\ge\pi(\mu_1)$,
by 521H.
As in the proof of 521I, $([0,1],\Cal N(\mu_1))$ is isomorphic to
$(\Bbb R,\Cal N)$.   Of course $\mu_1$ is not isomorphic to $\mu_L$.
But $\mu_L$ is isomorphic to a direct sum of countably many copies of
$\mu_1$, so by 521G we know that $\pi(\mu_L)$ is the cardinal product
$\omega\cdot\pi(\mu_1)$;  as $\pi(\mu_1)$ is surely infinite, this is
$\pi(\mu_1)$ again.   So we have the result in the special case.

\medskip

{\bf (b)} Now suppose that $(X,\Sigma,\mu)$ is any semi-finite perfect
measure space which is not purely atomic.   Then the completion
$\hat\mu$ of $\mu$ is still a semi-finite perfect measure which is not
purely atomic (212Gd, 451G(c-i)), and $\Cal N(\mu)=\Cal N(\hat\mu)$
(212Eb).
Because $\hat\mu$ is semi-finite and not purely atomic, there is a set
$E\in\Sigma$ of non-zero finite measure such that the subspace measure
$\hat\mu_E$ is atomless.   Set $\nu=\Bover1{\mu E}\hat\mu_E$, so that
$\nu$ is an atomless complete perfect probability measure on $E$, while
$\Cal N(\nu)=\Cal N(\mu_E)$.   Putting (a) together with 521F, we get

\Centerline{$\add\Cal N(\mu)=\add\Cal N(\hat\mu)
\le\add\Cal N(\hat\mu_E)=\add\Cal N(\nu)\le\add\Cal N$}

\noindent and similarly for $\cf$, $\shr$ and $\shr^+$.
}%end of proof of 521K

\leader{521L}{Proposition} (a) Let $(X,\Sigma,\mu)$ be a strictly
localizable measure space and $(Y,\Tau,\nu)$ a
locally compact semi-finite
measure space, and suppose that they have isomorphic measure algebras.
Then $(X,\in,\Cal N(\mu))\prGT(Y,\in,\Cal N(\nu))$;  consequently
$\cov\Cal N(\mu)\le\cov\Cal N(\nu)$ and
$\non\Cal N(\nu)\le\non\Cal N(\mu)$.

(b) Let $(X,\Sigma,\mu)$ be a \Mth\ compact
probability space with Maharam type $\kappa$.   Then
$\cov\Cal N(\mu)\penalty-100=\cov\Cal N_{\kappa}$
and $\non\Cal N(\mu)=\non\Cal N_{\kappa}$, where $\Cal N_{\kappa}$ is the
null ideal of the usual measure\cmmnt{ $\nu_{\kappa}$} on
$\{0,1\}^{\kappa}$.

(c) Let $(X,\Sigma,\mu)$ be a compact strictly localizable measure space
with measure algebra $\frak A$.   Then

\Centerline{$d(\frak A)=\min\{\#(A):A\subseteq X$ has full outer
measure$\}$.}

\proof{{\bf (a)} This follows immediately from 521Ha, because by 343B
there is an \imp\ function from $X$ to $Y$.

\medskip

{\bf (b)} The point is that $\nu_{\kappa}$ is a compact measure
(342Jd, 451Ja), so that we can apply (a) in both directions to see that
$\cov\Cal N(\mu)=\cov\Cal N_{\kappa}$ and
$\non\Cal N(\mu)=\non\Cal N_{\kappa}$.

\medskip

{\bf (c)} The case $\mu X=0$ is trivial;  suppose that $\mu X>0$.   Let
$\Cal K$ be a compact class such that $\mu$ is inner regular with
respect to $\Cal K$.

\medskip

\quad{\bf (i)} Suppose that $\ofamily{\xi}{d(\frak A)}{C_{\xi}}$ is a
family of centered sets in $\frak A$ covering $\frak A^+$.
For each $\xi<d(\frak A)$, set
$\Cal K_{\xi}=\{K:K\in\Cal K\cap\Sigma$, $K^{\ssbullet}\in C_{\xi}\}$;
then $\Cal K_{\xi}$ has the finite intersection property so there is a
point $x_{\xi}\in X\cap\bigcap\Cal K_{\xi}$.   Set
$A=\{x_{\xi}:\xi<d(\frak A)\}$.   If $K\in\Cal K\cap\Sigma$ and
$K\cap A=\emptyset$, then $K\notin\bigcup_{\xi<d(\frak A)}\Cal K_{\xi}$
so $K^{\ssbullet}=0$;  it follows that every measurable subset of
$X\setminus A$ is negligible and $A$ has full outer measure, while
$\#(A)\le d(\frak A)$.

\medskip

\quad{\bf (ii)} Let $\hat\mu$ be the completion of $\mu$, $\hat\Sigma$
its domain and $\theta:\frak A\to\hat\Sigma$ a lifting (341K, 212Gb,
322Da).   Take any $A\subseteq X$ of full outer measure for $\mu$;  then
it also has full outer measure for $\hat\mu$ (212Eb).   For $x\in A$,
set $C_x=\{a:a\in\frak A$, $x\in\theta a\}$;  then $\family{x}{A}{C_x}$
is a family of centered sets in $\frak A$ with union
$\frak A^+$, so $d(\frak A)\le\#(A)$.
}%end of proof of 521L

\leader{521M}{Proposition}\dvArevised{2010}
Let $(X,\Sigma,\mu)$ be a complete locally
determined measure space of magnitude at most $\add\mu$.   Then it is
strictly localizable.

\proof{ Write $\kappa$ for $\add\mu$.
Let $(\frak A,\bar\mu)$ be the measure algebra of $\mu$.   Then
there is a partition of unity $D\subseteq\frak A$ consisting of
elements
of finite measure;  as $\#(D)\le c(\frak A)\le\kappa$, there is a
family $\ofamily{\xi}{\kappa}{a_{\xi}}$ running over $D\cup\{0\}$.
For each $\xi<\kappa$, choose $E_{\xi}\in\Sigma$ such that
$E_{\xi}^{\ssbullet}=a_{\xi}$, and set
$F_{\xi}=E_{\xi}\setminus\bigcup_{\eta<\xi}E_{\eta}$.
Because $E_{\xi}\setminus F_{\xi}=\bigcup_{\eta<\xi}E_{\xi}\cap E_{\eta}$
is the union of fewer than $\add\mu$ negligible sets, it is negligible, and
$F_{\xi}\in\Sigma$, with $F_{\xi}^{\ssbullet}=a_{\xi}$.
Now
$\ofamily{\xi}{\kappa}{F_{\xi}}$ is a disjoint family of sets of
finite measure.   If $E\in\Sigma$ and $\mu E>0$, there is some
$\xi<\kappa$ such that $E^{\ssbullet}\Bcap a_{\xi}\ne 0$, and now
$\mu(E\cap F_{\xi})>0$.
Thus $\ofamily{\xi}{\kappa}{F_{\xi}}$ satisfies the condition of
213Oa, and $\mu$ is strictly localizable.
}%end of proof of 521M

\leader{521N}{Proposition} Let $(X,\Sigma,\mu)$ be a complete locally
determined localizable measure space of magnitude at most $\frak c$.
Then it is strictly localizable.

\proof{ Again let $(\frak A,\bar\mu)$ be the measure algebra of $\mu$,
and take a partition of unity $D\subseteq\frak A$ consisting of
elements
of finite measure;  as $\#(D)\le c(\frak A)\le\frak c$, there is an
injective function $h:D\to\Cal P\Bbb N$.   This time, because $\frak A$
is Dedekind complete, we can set $b_n=\sup\{d:d\in D$, $n\in h(d)\}$ for
each $n\in\Bbb N$.   If $d\in D$, then
$d=\inf_{n\in h(d)}b_n\Bsetminus\sup_{n\in\Bbb N\setminus h(d)}b_n$.
So if we choose $E_n\in\Sigma$ such that $E_n^{\ssbullet}=b_n$ for
each $n$, and set $F_d
=\bigcap_{n\in h(d)}E_n\setminus\bigcup_{n\in\Bbb N\setminus h(d)}E_n$
for $d\in D$, $\family{d}{D}{F_d}$ will be a disjoint family in
$\Sigma$
and $F_d^{\ssbullet}=d$ for every $d$.   Now $\mu F_d=\bar\mu d$ is
always finite;  and if $E\in\Sigma$ is non-negligible, there is a
$d\in D$ such that $0\ne\bar\mu(E^{\ssbullet}\Bcap d)=\mu(E\cap F_d)$.
Thus $\family{d}{D}{F_d}$ satisfies the condition of 213Oa, and $\mu$
is strictly localizable.
}%end of proof of 521N

\leader{521O}{Proposition} (a) If $(X,\Sigma,\mu)$ is a semi-finite
measure space, its magnitude is at most $\max(\omega,2^{\#(X)})$.

(b) If $(X,\Sigma,\mu)$ is a strictly localizable measure space, its
magnitude is at most $\max(\omega,\penalty-100\#(X))$.

(c) There is an infinite semi-finite measure space $(X,\Sigma,\mu)$
with magnitude $2^{\#(X)}$.

(d) If $\familyiI{A_i}$ is a disjoint family of subsets of $X$ and
$\#(I)>\max(\omega,\magnitude(\mu))$ then there is an $i\in I$ such that
$X\setminus A_i$ has full outer measure.

\proof{{\bf (a)-(b)} These are elementary.   If $(X,\Sigma,\mu)$ is
semi-finite, with measure algebra $\frak A$, then

\Centerline{$c(\frak A)\le\#(\frak A)\le\#(\Sigma)
\le\#(\Cal PX)=2^{\#(X)}$.}

\noindent If $\mu$ is strictly localizable, with decomposition
$\familyiI{X_i}$, then $\familyiI{X_i^{\ssbullet}}$ is a partition of
unity in $\frak A$ consisting of elements of finite measure, so

\Centerline{$c(\frak A)
\le\max(\omega,\#(\{i:i\in I$, $\mu X_i>0\}))\le\max(\omega,\#(X))$}

\noindent by 332E.

\medskip

{\bf (c)} Let $\ofamily{\xi}{\omega_1}{X_{\xi}}$ be a disjoint family
of sets such that
$\#(X_{\xi})=\#(\Cal P(\bigcup_{\eta<\xi}X_{\eta}))$ for every
$\xi<\omega_1$;  for each $\xi$, let
$h_{\xi}:\Cal P(\bigcup_{\eta<\xi}X_{\eta})\to X_{\xi}$ be an
injection.
Set $X=\bigcup_{\xi<\omega_1}X_{\xi}$.   For $A\subseteq X$ define
$f_A:\omega_1\to X$ by setting
$f(\xi)=h_{\xi}(A\cap\bigcup_{\eta<\xi}X_{\eta})$ for each $\xi$;  let
$J_A$ be $f_A[\omega_1]$ and $\mu_A$ the countable-cocountable measure
on $J_A$.    Observe that $\#(J_A)=\omega_1$ for every $A\subseteq X$,
and that if $A$, $B\subseteq X$ are distinct then $J_A\cap J_B$ is
countable.   So if we set $\mu E=\sum_{A\subseteq X}\mu_A(E\cap J_A)$
whenever $E\subseteq X$ is such that $E\cap J_A$ is countable or
cocountable in $J_A$ for every $A$, then $\mu$ will be a complete locally
determined measure on $X$.   Since $\mu J_A=1$ and $\mu(J_A\cap J_B)=0$
whenever $A$, $B\subseteq X$ are distinct, $\mu$ has magnitude
$2^{\#(X)}$.

\medskip

{\bf (d)} \Quer\ Otherwise, there is for each $i\in I$ a measurable set $F_i$ of
non-zero measure such that $F_i\subseteq A_i$.   Again writing $\frak A$ for the
measure algebra of $\mu$,
$\familyiI{F_i^{\ssbullet}}$ is a disjoint family in $\frak A\setminus\{0\}$
so $\#(I)\le c(\frak A)\le\max(\omega,\magnitude(\mu))$.\ \Bang
}%end of proof of 521O

\leader{521P}{Proposition} (a) If $2^{\lambda}<2^{\kappa}$ whenever
$\frak c\le\lambda<\kappa$ and $\cf\lambda>\omega$, then
the magnitude\cmmnt{ $\magnitude\mu$} of $\mu$ is at most
$\max(\omega,\#(X))$ for every localizable measure space $(X,\Sigma,\mu)$.

(b) Suppose that $2^{\frak c}=2^{\frak c^+}$.   Then there is a
localizable measure space $(Y,\Tau,\nu)$ with $\#(Y)=\frak c$ and
$\magnitude\nu=\frak c^+$.

\medskip

\noindent{\bf Remark}\cmmnt{ \TeX, for once, is obscure;}
$2^{\frak c^+}$ here is $\#(\Cal P(\frak c^+))$\cmmnt{, not
$(2^{\frak c})^+$}.

\proof{{\bf (a)} If $\magnitude\mu\le\omega$ we can stop.   Otherwise,
set $\kappa=\magnitude\mu$.   Let $(\frak A,\bar\mu)$
be the measure algebra of $\mu$, so that $\kappa=c(\frak A)$ and there
is a disjoint family $\ofamily{\xi}{\kappa}{a_{\xi}}$ in
$\frak A^+$ (332F).   If $\tilde\mu$ is the c.l.d.\
version of $\mu$, we can identify
$\frak A$ with the measure algebra of $\tilde\mu$ (322Db).

\medskip

\quad{\bf case 1} If $\kappa\le\frak c$, $\tilde\mu$ is strictly
localizable (521N), so has a lifting $\theta$ (341K again);  but now
$\ofamily{\xi}{\kappa}{\theta a_{\xi}}$ is a disjoint family of
non-empty subsets of $X$, so $\#(X)\ge\kappa$.

\medskip

\quad{\bf case 2} If $\kappa>\frak c$, of course $X$ is uncountable
(521Oa).   For $\xi<\kappa$, choose $E_{\xi}\in\Sigma$ such that
$E_{\xi}^{\ssbullet}=a_{\xi}$.
\Quer\ If $\#(X)<\kappa$, there is a set $Y\subseteq X$ such that
$\#(Y)$ has uncountable cofinality and $I_Y=\{\xi:\xi<\kappa$,
$\mu^*(E_{\xi}\cap Y)>0\}$ has cardinal greater than
$\max(\frak c,\#(X))$.   \Prf\ If $\cf(\#(X))$ is uncountable, take
$Y=X$.   Otherwise, let $\sequencen{Y_n}$ be an increasing sequence of
subsets of $X$, with union $X$, such that $\#(Y_n)$ is an uncountable
successor cardinal less than $\#(X)$ for every $n$.   If $\xi<\kappa$,
there is some $n$ such that $E_{\xi}\cap Y_n$ is non-negligible, that
is, $\xi\in I_{Y_n}$.   So the non-decreasing sequence
$\sequencen{I_{Y_n}}$ has union $\kappa$, and
there is some $n\in\Bbb N$ such that
$\#(I_{Y_n})>\max(\frak c,\#(X))$.   Now we can take $Y=Y_n$.\ \Qed

For every $J\subseteq I_Y$, set $b_J=\sup_{\xi\in J}a_\xi$ and let
$F_J\in\Sigma$ be such that $F_J^{\ssbullet}=b_J$.   If $J$,
$K\subseteq I_Y$ are distinct, there is a $\xi\in J\symmdiff K$, in
which case
$a_\xi\Bsubseteq b_J\Bsymmdiff b_K$, $E_\xi\setminus(F_J\symmdiff F_K)$
is negligible and $Y\cap(F_J\symmdiff F_K)$ is non-empty.   Thus
$J\mapsto Y\cap F_J:\Cal PI_Y\to\Cal PY$ is injective, and

\Centerline{$2^{\#(I_Y)}\le 2^{\#(Y)}\le 2^{\#(X)}\le 2^{\#(I_Y)}$.}

\noindent Setting $\lambda=\max(\frak c,\#(Y))$, $\kappa'=\#(I_Y)$ we
now have $\frak c\le\lambda<\kappa'$, $\cf\lambda>\omega$ and
$2^{\lambda}=2^{\kappa'}$, which is supposed to be impossible.\ \Bang

So in this case also we have $\#(X)\ge\kappa$.

\medskip

{\bf (b)(i)} Set $I=\Cal P\frak c^+$ and
$X=\{0,1\}^I\cong\{0,1\}^{2^{\frak c}}$.   Putting 5A4Be and 5A4C(a-ii)
together, we see that there is a set $Y\subseteq X$, with cardinal
at most $\frak c^{\omega}=\frak c$, which meets every non-empty
G$_{\delta}$ subset of $X$.   In particular, if $\Cal K\subseteq I$ is
countable and $x\in X$ there is a $y\in Y$ such that
$y\restr\Cal K=x\restr\Cal K$.

\medskip

\quad{\bf (ii)} Let $\mu$ be the complete locally determined
localizable
measure on $X$ described in 216E, with $C=\frak c^+$.   Then $Y$ has
full outer measure in $X$.   \Prf\ (I follow the notation and argument
of 216E.)   If
$\mu E>0$, then, by the argument of part (g) of the proof of 216E,
there are a $\gamma<\frak c^+$ and a $K\in[I]^{\le\omega}$ such that
$F_{\gamma K}\subseteq E$, where
$F_{\gamma K}=\{x:x\restr K=x_{\gamma}\restr K\}$.   But $Y$ was
chosen to meet every such set.   As $E$ is arbitrary, $Y$ has full outer
measure.\ \Qed

\medskip

\quad{\bf (iii)} $\magnitude\mu=\frak c^+$.   \Prf\ In the language of
216E, we have a family $\ofamily{\gamma}{\frakc^+}{G_{\{\gamma\}}}$ of
$\mu$-atoms of measure $1$, each pair with negligible intersection,
and every non-negligible measurable set meets some $G_{\{\gamma\}}$ in a
non-negligible set.\ \Qed

\medskip

\quad{\bf (iv)} Now let $\nu$ be the subspace measure on $Y$.   By
214Ie, $\nu$ is complete, locally determined and localizable.   By
322I, we can identify the measure algebras of $\mu$ and $\nu$, so
$\magnitude\nu=\magnitude\mu=\frak c^+$, while $\#(Y)=\frak c$.
}%end of proof of 521P

\leader{521Q}{Free \dvrocolon{products}}\cmmnt{ We have some simple
calculations associated with the measure algebra free products of \S325.

\medskip

\noindent}{\bf Proposition} (a) Let $(\frak A,\bar\mu)$ and
$(\frak B,\bar\nu)$ be semi-finite measure algebras and
$(\frak C,\bar\lambda)$ their localizable measure algebra free product.
Then

\Centerline{$c(\frak C)\le\max(\omega,c(\frak A),c(\frak B))$,}

\Centerline{$\tau(\frak C)\le\max(\omega,\tau(\frak A),\tau(\frak B))$.}

\wheader{521Q}{0}{0}{0}{36pt}
(b) Let $\familyiI{(\frak A_i,\bar\mu_i)}$ be a family of probability
algebras, and $(\frak C,\bar\lambda)$ their probability algebra free
product.   Then

\Centerline{$\tau(\frak C)
\le\max(\omega,\#(I),\sup_{i\in I}\tau(\frak A_i))$.}

\proof{{\bf (a)(i)} Let $A\subseteq\frak A$, $B\subseteq\frak B$ be
partitions of unity consisting of elements of
finite measure (322Ea).
Then $C=\{a\otimes b:a\in A$, $b\in B\}$ is a disjoint family in
$\frak C$, and

$$\eqalignno{\sup C
&=\sup\{(a\otimes 1)\Bcap(1\otimes b):a\in A,\,b\in B\}
=(\sup_{a\in A}a\otimes 1)\Bcap(\sup_{b\in B}1\otimes b)\cr
\displaycause{313Bc}
&=(\sup A\otimes 1)\Bcap(1\otimes\sup B)\cr
\displaycause{325Da}
&=(1\otimes 1)\Bcap(1\otimes 1)
=1;\cr}$$

\noindent that is, $C$ is a partition of unity.   As every member of $C$
has finite measure,

\Centerline{$c(\frak C)
\le\max(\omega,\#(C))
=\max(\omega,\#(A),\#(B))
=\max(\omega,c(\frak A),c(\frak B))$}

\noindent by 332E.

\medskip

\quad{\bf (ii)} As for Maharam types, I am just repeating the result stated
and proved in 334B.

\medskip

{\bf (b)} This is 334D.
}%end of proof of 521Q

\leader{521R}{Proposition} If $(X,\Sigma,\mu)$ is any measure space,
its Maharam type is at most $2^{\#(X)}$.

\proof{ If $\frak A$ is the measure algebra of $\mu$,

\Centerline{$\tau(\frak A)\le\#(\frak A)\le\#(\Sigma)\le\#(\Cal PX)
=2^{\#(X)}$.}

}%end of proof of 521R

\leader{521S}{Proposition} (a) A countably separated measure space
has Maharam type at most $2^{\frakc}$.

(b) There is a countably separated quasi-Radon probability space with
Maharam type $2^{\frakc}$.

(c) A countably separated semi-finite measure space has magnitude at
most $2^{\frakc}$.

(d) There is a countably separated semi-finite measure space with
magnitude $2^{\frakc}$.

\proof{ Set $\kappa=2^{\frakc}$.

\medskip

{\bf (a)} If $(X,\Sigma,\mu)$ is countably separated, there is an
injective function from $X$ to $\Bbb R$ (343E), so $\#(X)\le\frak c$;
now use 521R.

\medskip

{\bf (b)} As in (b-i) of the proof of 521P, there is a set
$Y\subseteq X=\{0,1\}^{\kappa}$, with cardinal
$\frak c$, which meets every non-empty G$_{\delta}$ subset of $X$, and
therefore has full outer measure for the usual measure
$\nu_{\kappa}$ of $X$.

In $[0,1]$ let $\family{y}{Y}{C_y}$
be a disjoint family of sets of full outer measure for Lebesgue
measure $\mu_1$ on $[0,1]$ (419I), and set
$C=\{(y,t):y\in Y,\,t\in C_y\}\subseteq Z=X\times[0,1]$.   Now $C$ has
full outer measure for the product measure $\lambda$ on $Z$.   \Prf\
Suppose that $W\subseteq Z$ and $\lambda W>0$.   Then
$\int\mu_1 W[\{x\}]\nu_{\kappa}(dx)>0$ (252D), so
$\{x:\mu_1W[\{x\}]>0\}$ has non-zero measure and meets $Y$.   Taking
$y\in Y$ such that $\mu_1W[\{y\}]>0$, $\{y\}\times(C_y\cap W[\{y\}])$
is a non-empty subset of $C\cap W$.\ \QeD\

The measure algebra $\frak A$ of the subspace measure $\lambda_C$ on $C$
can therefore be identified with the measure algebra of $\lambda$ (322Jb),
and has Maharam type $\kappa$.   Because $\family{y}{Y}{C_y}$ is disjoint,
each horizontal section of $C$ is a singleton and $C$ is separated by the
measurable sets $C\cap(X\times[0,q])$ for $q\in\Bbb Q\cap[0,1]$.
Thus $\lambda_C$ is countably separated.

If we give $Z$ the product topology, then $\lambda$ is a Radon measure
(417T, or otherwise), so $\lambda_C$ is quasi-Radon for the subspace
topology (415B).

\medskip

{\bf (c)} As in (a), $\#(X)\le\frak c$, so we can use 521Oa.

\medskip

{\bf (d)(i)} The first step is to build a measure space of magnitude
$2^{\frakc}$ and cardinal $\frak c$, as follows.
Let $h:\frak c\to([\frak c]^{\le\omega})^2$ be a surjection;
take its two components to be $h_1$ and $h_2$.   For $D\subseteq\frak c$
set $F_D=\{\xi:\xi<\frak c$, $h_2(\xi)=D\cap h_1(\xi)\}$.   For
$I\in[\frak c]^{\le\omega}$ set
$A_I=\{\xi:\xi<\frak c$, $I\not\subseteq h_1(\xi)\}$, and set
$\Cal A=\bigcup\{\Cal PA_I:I\in[\frak c]^{\le\omega}\}$;  note that
$\Cal A$ is a $\sigma$-ideal of subsets of $\frak c$.

If $D\subseteq\frak c$, $F_D\notin\Cal A$.   \Prf\ If
$I\in[\frak c]^{\le\omega}$, there is a $\xi<\frak c$ such that
$h(\xi)=(I,I\cap D)$,   Now $\xi\in F_D\setminus A_I$;  as $I$ is
arbitrary, $F_D\notin\Cal A$.\ \QeD\  So we can define a measure $\nu_D$
on $\frak c$ by saying that

$$\eqalign{\nu_D(E)
&=1\text{ if }E\subseteq\frak c\text{ and }F_D\setminus E\in\Cal A,\cr
&=0\text{ if }E\subseteq\frak c\text{ and }F_D\cap E\in\Cal A,\cr
&\text{undefined otherwise},\cr}$$

\noindent and $\nu_DF_D=1$.

If $D$, $D'\subseteq\frak c$ are distinct, $F_D\cap F_{D'}\in\Cal A$.
\Prf\ Take $\eta\in D\symmdiff D'$.   If $\xi\in F_D\cap F_{D'}$, then
$D\cap h_1(\xi)=h_2(\xi)=D'\cap h_1(\xi)$, so $\eta\notin h_1(\xi)$ and
$\xi\in A_{\{\eta\}}$.   Thus
$F_D\cap F_{D'}\subseteq A_{\{\eta\}}\in\Cal A$.\ \Qed

So if we set $\nu=\sum_{D\subseteq\frak c}\nu_D$, as defined in
234G, $\nu$ is a measure on
$\frak c$ such that $\nu F_D=1$ and $\nu(F_D\cap F_{D'})=0$ for all
distinct $D$, $D'\subseteq\frak c$.   Also $\nu$ is semi-finite, because if
$\nu E>0$ there is a $D\subseteq\frak c$ such that $\nu_DE>0$, in which
case $\nu(E\cap F_D)=\nu_D(E\cap F_D)=1$.   So $\nu$ is a semi-finite
measure on $\frak c$ with magnitude $\#(\Cal P\frak c)=2^{\frakc}$.
Because every $\nu_D$ is complete, so is $\nu$ (234Ha).

\medskip

\quad{\bf (ii)} As in (b),
let $\ofamily{\xi}{\frak c}{C_{\xi}}$ be a disjoint family of
subsets of $[0,1]$ all with full outer measure for Lebesgue measure
$\mu_1$.   Set $Z=\frak c\times[0,1]$ with its c.l.d.\ product measure
$\lambda=\nu\times\mu_1$, and
$C=\{(\xi,t):\xi<\frak c$, $t\in C_{\xi}\}\subseteq Z$.   Then
$C$ has full outer measure, by the argument of (b) above.   So, as in
(b), the measure algebra $\frak A$ of the subspace measure $\lambda_C$
on $C$ can be identified with the measure algebra of
$\lambda$.   The map $E\mapsto C\cap(E\times[0,1])$ induces a
measure-preserving homomorphism from the measure algebra of $\nu$ to
$\frak A$, so $\magnitude\lambda_C=c(\frak A)$ is at least $2^{\frakc}$;
by (c), it is exactly $2^{\frakc}$.
Also as in (b), $\lambda_C$ is countably separated.
}%end of proof of 521S

\leader{521T}{}\cmmnt{ In \S464 I looked at the $L$-space $M$ of
bounded additive functionals on $\Cal PI$ for infinite sets $I$, of
which $I=\Bbb N$ is of course by far the most important, and found a
band decomposition of $M$ as
$M_{\tau}\oplus(M_{\text{m}}\cap M_{\tau}^{\perp})
\oplus M_{\text{pnm}}$,
where $M_{\tau}$ consists of  the `completely additive' functionals
(and may be identified with $\ell^1(I)$), $M_{\text{m}}$ consists of the
`measurable' functionals (that is, those integrated by the usual
measure
on $\Cal PI$), and $M_{\text{pnm}}=M_{\text{m}}^{\perp}$ consists of
the `purely non-measurable' functionals.
Any non-negative functional $\theta\in M$ can be identified with a Radon
measure $\mu_{\theta}$ on the Stone-\v{C}ech compactification $\beta I$
(464P).   The purely atomic measures on $I$ correspond to members of
$M_{\tau}$.   Among the others, the general rule is that `simple'
measures must correspond to functionals in $M_{\text{pnm}}$;  see 464Pa
and 464Xa.   The next proposition, strengthening 464Qb, shows that
this rule is followed by Maharam types.

\medskip

\noindent}{\bf Proposition} Let $I$ be a set, and suppose that a
non-zero $\theta\in(M_{\text{m}}\cap M_{\tau}^{\perp})^+$, as defined
in \S464, corresponds to the Radon measure $\mu_{\theta}$ on $\beta I$.
Let $\nu$ be the usual measure on $\Cal PI$.   Then
the Maharam type of $\mu_{\theta}$ is at least $\cov\Cal N(\nu)$.

\proof{ Of course $I$ has to be infinite, since not every
additive functional on $\Cal PI$ is completely additive;  so
$\cov\Cal N(\nu)$ is not $\infty$.   By 464Qc, we know that

\Centerline{$\{(a,b):a$, $b\subseteq I$,
$\theta a=\Bover12\theta I$,
$\theta(a\cap b)=\Bover14\theta I\}$}

\noindent is conegligible for the product measure $\nu\times\nu$ on
$(\Cal PI)^2$.   Set

\Centerline{$A_0=\{a:a\subseteq I$, $\theta a=\Bover12$,
$\{b:\theta(a\cap b)=\Bover14\theta I\}$ is
$\nu$-conegligible$\}$;}

\noindent then $A_0$ is $\nu$-conegligible.   Now take a set
$A\subseteq A_0$ which is maximal subject to the requirement that
$\theta(a\cap b)=\bover14\theta I$ for all distinct $a$, $b\in A$.   If
$a$ is any subset of $I$, then either $a\in A$, or $a\notin A_0$, or
there is a $b\in A\setminus\{a\}$
such that $\theta(a\cap b)\ne\bover14\theta I$;  so
$\Cal PI$ is the union of

\Centerline{$(\Cal PI\setminus A_0)
  \cup\bigcup_{b\in A}\{a:\theta(a\cap b)\ne\Bover14\theta I\}$}

\noindent and $\cov\Cal N(\nu)\le 1+\#(A)$.
As $\cov\Cal N(\nu)$ is surely
infinite, it is in fact less than or equal to $\#(A)$.

Now consider the open-and-closed sets $\widehat{a}\subseteq\beta I$ for
$a\in A$.   If $a$, $b\in A$ are distinct,

\Centerline{$\mu_{\theta}(\widehat{a}\symmdiff\widehat{b})
=\mu_{\theta}(\widehat{a\symmdiff b})
=\theta(a\symmdiff b)=\bover12\theta I>0$.}

\noindent So in the measure algebra $\frak A$ of $\mu_{\theta}$,
$\{\widehat{a}^{\ssbullet}:a\in A\}$ is a discrete set with cardinal
at least $\cov\Cal N(\nu)$, and the topological density of
$\frak A$ is at least $\cov\Cal N(\nu)$ (5A4B(h-ii) again).
By 521E, $\tau(\mu_{\theta})=\tau(\frak A)\ge\cov\Cal N(\nu)$.
}%end of proof of 521T

\exercises{\leader{521X}{Basic exercises (a)}
%\spheader 521Xa
Let $\Cal B(\Bbb R)$ be the Borel $\sigma$-algebra of $\Bbb R$,
and $\mu$ the restriction of Lebesgue measure to $\Cal B(\Bbb R)$.
Show that $\add\mu=\omega_1$.   \Hint{if $\frak c=\omega_1$,
$[\Bbb R]^{\omega_1}\not\subseteq\Cal B(\Bbb R)$;  if $\frak c>\omega_1$,
$[\Bbb R]^{\omega_1}\cap\Cal B(\Bbb R)=\emptyset$;
or use 423M and 423R\formerly{4{}23Q}.}
%521A

\spheader 521Xb Let $(X,\Sigma,\mu)$ be a semi-finite
locally compact measure
space.   Show that $\add\mu$ is the least cardinal of any set
$\Cal E\subseteq\Sigma$ such that $\bigcup\Cal E\notin\Sigma$, or
$\infty$ if there is no such $\Cal E$.   \Hint{451Q.}
%521A+

\spheader 521Xc Let $(X,\Sigma,\mu)$ be a complete locally
determined measure space, and $\kappa$ a cardinal such that
$\kappa<\cov\Cal N(\mu_E)$ for every non-negligible measurable set
$E\subseteq X$, writing $\mu_E$ for the subspace measure.
Suppose that $A\subseteq X$ is such that both $A$ and
$X\setminus A$ are expressible as the union of at most $\kappa$ members of
$\Sigma$.   Show that $A\in\Sigma$.
%521B+

\sqheader 521Xd\dvAnew{2013}(i) Find a probability space $(X,\Sigma,\mu)$,
with measure algebra $\frak A$, such that $\pi(\frak A)<\pi(\mu)$.
%\pi(\frak A)=1$
(ii) Find a probability space $(X,\Sigma,\mu)$, with null ideal
$\Cal N(\mu)$, such that $\cf\Cal N(\mu)<\pi(\mu)$.
%\cf\Cal N=1
(iii) Find a probability space $(X,\Sigma,\mu)$
such that $\pi(\mu)<\cf\Cal N(\mu)$.   \Hint{513X(q-iii).}
%direct sum with $\cf\Cal N(\mu_n)=\pi(\mu_n)=\omega_n$.
%521D

\spheader 521Xe Let $(X,\Sigma,\mu)$ be a measure space and $\nu$ an
indefinite-integral measure over $\mu$.   Show that
$\add\Cal N(\nu)\ge\add\Cal N(\mu)$, $\cf\Cal N(\nu)\le\cf\Cal N(\mu)$,
$\non\Cal N(\nu)\ge\non\Cal N(\mu)$, $\cov\Cal N(\nu)\le\cov\Cal N(\mu)$,
$\shr\Cal N(\nu)\le\shr\Cal N(\mu)$,
$\shr^+\Cal N(\nu)\le\shr^+\Cal N(\mu)$, $\pi(\nu)\le\pi(\mu)$,
$\tau(\nu)\le\tau(\mu)$.
%322K 521F

\spheader 521Xf\dvAnew{2013} Let $(X,\Sigma,\mu)$ be a
semi-finite measure space which is not purely atomic.   Show that
$\pi(\mu)\ge\pi(\mu_L)$, where $\mu_L$ is Lebesgue measure on $\Bbb R$.
%521H

\spheader 521Xg Let $(X,\Sigma,\mu)$ be an atomless measure space with
locally determined negligible sets (definition: 213I).   Show that
$\non\Cal N(\mu)\ge\non\Cal N$, where $\Cal N$ is the null ideal of
Lebesgue measure.
%521I

\spheader 521Xh Let $(X,\Sigma,\mu)$ and $(Y,\Tau,\nu)$ be complete
locally determined measure spaces, neither of measure $0$,
and $\mu\times\nu$ the c.l.d.\ product measure on $X\times Y$.   Show that

\Centerline{$\non\Cal N(\mu\times\nu)
=\max(\non\Cal N(\mu),\non\Cal N(\nu))$,}

\Centerline{$\cov\Cal N(\mu\times\nu)
\le\min(\cov\Cal N(\mu),\cov\Cal N(\nu))$}

\noindent with equality if either $\mu$ or $\nu$ is strictly localizable,

\Centerline{$\add(\mu\times\nu)=\add\Cal N(\mu\times\nu)
\le\min(\add\Cal N(\mu),\add\Cal N(\nu))$,}

\Centerline{$\cf\Cal N(\mu\times\nu)
\ge\max(\cf\Cal N(\mu),\cf\Cal N(\nu))$,}

\Centerline{$\shr\Cal N(\mu\times\nu)
\ge\max(\shr\Cal N(\mu),\shr\Cal N(\nu))$,}

\Centerline{$\shr^+\Cal N(\mu\times\nu)
\ge\max(\shr^+\Cal N(\mu),\shr^+\Cal N(\nu))$,}

\Centerline{$\pi(\mu\times\nu)\ge\max(\pi(\mu),\pi(\nu))$.\dvAnew{2013}}
%521J

\spheader 521Xi Let $(X,\Sigma,\mu)$ be a probability space, and
$\mu^{\Bbb N}$ the product measure on $X^{\Bbb N}$.
(i) Show that $X$ has a set of full
outer measure with cardinal at most $\non\Cal N(\mu^{\Bbb N})$.
(ii) Show that if $\Cal A\subseteq\Sigma\setminus\Cal N(\mu)$ and
$\#(\Cal A)<\cov\Cal N(\mu^{\Bbb N})$, then there is a countable set
which meets every member of $\Cal A$.
%521J

\spheader 521Xj Show that the direct sum of $\frak c$ or fewer
countably separated measure spaces is countably separated.
%521S

\spheader 521Xk Show that $2^{\frak c}<2^{\frak c^+}$ iff every
countably separated complete locally determined localizable measure
space is strictly localizable.  \Hint{521P, 521S, 252Yp.}
%521S 521P mt52bits

\spheader 521Xl Show that if $(X,\Sigma,\mu)$ is a purely atomic
countably separated semi-finite measure space then its magnitude is at
most $\max(\omega,\#(X))$ and its Maharam type is countable.
%521S

\spheader 521Xm Suppose that $2^{\kappa}\le\frak c$ for every
$\kappa<\frak c$.   Show that there is a countably separated
semi-finite measure space with magnitude $2^{\frakc}$.
%521Sd mt52bits

\spheader 521Xn\dvAnew{2016} For a measure space $(X,\Sigma,\mu)$ with
null ideal $\Cal N(\mu)$, write $\hcov(\mu)$ for
$\inf_{E\in\Sigma\setminus\Cal N(\mu)}\penalty-100\cov(E,\Cal N(\mu))$.
(Count $\inf\emptyset$ as $\infty$, as usual.)
Show that if $(X,\Sigma,\mu)$ and $(Y,\Tau,\nu)$ are semi-finite measure
spaces, neither having zero measure,
with c.l.d.\ product $(X\times Y,\Lambda,\lambda)$,
then $\hcov(\lambda)=\min(\hcov(\mu),\hcov(\nu))$.
%521J 521Xh out of order query

\leader{521Y}{Further exercises (a)}
%\spheader 521Ya
Find a probability space $(X,\Sigma,\mu)$, a set $Y$ and
a function $f:X\to Y$ such that, setting $\nu=\mu f^{-1}$,
$\add\Cal N(\mu)>\add\Cal N(\nu)$, $\cf\Cal N(\mu)<\cf\Cal N(\nu)$,
$\shr\Cal N(\mu)<\shr\Cal N(\nu)$ and $\pi(\mu)<\pi(\nu)$.
%521H  mt52bits

\spheader 521Yb Find a strictly localizable measure space
$(X,\Sigma,\mu)$, a set $Y$, and a
function $f:X\to Y$ such that, setting $\nu=\mu f^{-1}$, $\nu$ is
semi-finite and $\tau(\mu)<\tau(\nu)$.
%521H mt52bits

\spheader 521Yc Let $(X,\Sigma,\mu)$ and $(Y,\Tau,\nu)$ be localizable
measure spaces, and suppose that
\penalty-100$\max(\magnitude(\nu),\discretionary{}{}{}\tau(\nu))\le\frak c$.
Show that the c.l.d.\
product measure on $X\times Y$ is localizable.
%521N mt52bits
%what if just  \tau(\nu)\le\frak c ?

\spheader 521Yd Show that there is a probability space
$(X,\Sigma,\mu)$ with Maharam type greater than $\#(X)$.   \Hint{523Ib.}
%521R

\spheader 521Ye
Let $\kappa$ be an infinite cardinal.   Let us say that a measure
space $(X,\Sigma,\mu)$ is {\bf $\kappa$-separated} if there is a family
$\Cal E\subseteq\Sigma$, with cardinal at most $\kappa$, separating the
points of $X$.
(i) Show that there is a disjoint family $\Cal A$ of
subsets of $\{0,1\}^{\kappa}$, all of full outer measure for the usual
measure of $\{0,1\}^{\kappa}$, such that $\#(\Cal A)=2^{\kappa}$.
(ii) Show that every $\kappa$-separated measure space has Maharam type
at most $2^{2^{\kappa}}$,  %521Sa
and that there is a $\kappa$-separated quasi-Radon
probability space with Maharam type $2^{2^{\kappa}}$.   %521Sb
(iii) Show
that every semi-finite $\kappa$-separated measure space has magnitude at
most $2^{2^{\kappa}}$, %521Sc
and that there is a semi-finite $\kappa$-separated
measure space with magnitude greater than $2^{\kappa}$.   %521Sd
(iv) Suppose that $\frak c\le\kappa\le\lambda$ and
$2^{\kappa}=2^{\lambda}$.   Show that the usual measure on
$\{0,1\}^{\lambda}$ is $\kappa$-separated.  %mt52bits
\leaveitout{?? (v) Show that for every
uncountable cardinal $\kappa$ there is a compact $\kappa$-separated
Maharam-type-homogeneous Radon
probability space with Maharam type $\kappa$ which is not isomorphic to
$\{0,1\}^{\kappa}$ with its usual measure.  See \S344 notes.}
%521S
}%end of exercises

\leaveitout{
%Problem:  which of the following are possible?
%   \add\Cal N(\mu\times\nu) < \min(\add\Cal N(\mu),\add\Cal N(\nu))
%   \cf\Cal N(\mu\times\nu) > \max(\cf\Cal N(\mu),\cf\Cal N(\nu))
%   \shr\Cal N(\mu\times\nu) > \max(\shr\Cal N(\mu),\shr\Cal N(\nu))

\leader{}{Problem} Must there be a family
$\ofamily{\xi}{\omega_1}{\Cal F_{\xi}}$ of distinct ultrafilters on
$\Bbb N$ such that $\bigcap_{\xi<\omega_1}\Cal F_{\xi}$ is measurable?
}%end of leaveitout

\endnotes{\Notesheader{521}
The cardinal functions of an ideal can be thought of as measures of the
`complexity' of that ideal.   In a measure space, it is natural to suppose
that a subspace measure (at least, on a measurable subspace) will be
`simpler' than the original measure;  in 521F we see that the additivity
and uniformity tend to
rise and the covering number, cofinality, shrinking
number and $\pi$-weight tend to fall.
Similarly, an image measure ought to be simpler than
its parent;  but here, while additivity rises and cofinality and shrinking
number fall, uniformity falls and covering number rises (521H).
Also there is a trap if the original measure is not complete (521Ya), and
$\pi$-weight is more complicated (521H(a-ii)).   There is a similar
problem concerning topological $\pi$-weight, which led to the concept
of network weight (5A4Ai, 5A4Bc);  and just as network weight matches
topological weight for compact Hausdorff spaces (5A4C(a-i)), an appropriate
hypothesis on our measures can make their $\pi$-weights more coherent
(521H(a-ii)).

Direct sums should not be more complex than their most complex component;
521G confirms this prejudice except in respect of cofinality.   Since we
are looking, in effect, at the cofinality of a product of partially ordered
sets, we can expect at least as many difficulties as are to be found in pcf
theory (\S5A2).   We should like to be able to bound the complexity of a
product in terms of the complexities of the factors;  here there seem to be
some interesting questions, and 521J and 521Xh are, I hope, only a start.

Consider the statement

\Centerline{($\dagger$) `$\magnitude\mu\le\#(X)$ for every localizable
measure space $(X,\Sigma,\mu)$'.}

\noindent From 521P we see that the generalized continuum hypothesis
implies ($\dagger$), and also that there are simple models of set
theory in which ($\dagger$) is false ({\smc Kunen 80}, VIII.4.7;
{\smc Jech 03}, 15.18).   I do
not know whether there is a natural combinatorial statement
equiveridical with ($\dagger$).   If we amend ($\dagger$) to

\Centerline{`$\magnitude\mu\le\#(X)$ for every countably separated
localizable measure space $(X,\Sigma,\mu)$'}

\noindent we find ourselves with a statement equiveridical with
`$2^{\frak c}=2^{\frak c^+}$' (cf.\ 521Xk).

I give space to `countably separated' measures because these can be
identified with the topological measures on subsets of $\Bbb R$, and I
do not think it is immediately apparent just how complicated these can
be.   In fact, as shown by the proofs of parts (b) and (d) in 521S, most
of the phenomena which can arise in any measure space with cardinal less
than or equal to $\frak c$ can appear in countably separated measure
spaces.   In 521Sb I add `quasi-Radon' to show that the very strong
restrictions on countably separated Radon probability measures (522Wa)
depend on their perfectness, not on their $\tau$-additivity.

The constructions in 521Oc and 521Sd both depend on almost-disjoint
families of sets.   Those described here are elementary.   In many
models of set theory, we have much more striking results, of which
521Xm is a simple example.

Some new considerations intrude rather abruptly
in 521T, but the argument here is both elementary and important, quite
apart from its use in helping us to understand the classification
scheme in \S464.
}%end of notes

\discrpage

\frfilename{mt522.tex}
\versiondate{11.7.23}
\copyrightdate{2003}

\def\chaptername{Cardinal functions of measure theory}
\def\sectionname{Cicho\'n's diagram}

\def\finint{\text{\tt finint}}
\def\disj{\text{\tt disj}}

\newsection{522}

In this section I describe some extraordinary relationships between the
cardinals associated with the ideals of meager and negligible sets in
the real line.   I concentrate on the strikingly symmetric pattern of
Cicho\'n's diagram (522B);  the first half
of the section is taken up with proofs
of the facts encapsulated in this diagram.   I include a handful of
results characterizing some of the most important cardinals here (522C,
522S), notes on Martin cardinals associated with the diagram
(522T) and the Freese-Nation number of $\Cal P\Bbb N$ (522U), and a
brief discussion of cofinalities (522V).

\leader{522A}{Notation} In this section, I will use the symbols $\Cal M$
and $\Cal N$ for the ideals of meager and negligible subsets of
$\Bbb R$ respectively.   Associated with these we have the eight
cardinals $\add\Cal M$,
$\cov\Cal M$, $\non\Cal M$, $\cf\Cal M$, $\add\Cal N$, $\cov\Cal N$,
$\non\Cal N$ and $\cf\Cal N$.   In addition we have two
cardinals associated with the partially ordered set $\NN$:  the {\bf
bounding number} $\frak b=\add_{\omega}\NN$\cmmnt{ (see 513H for the
definition of $\add_{\omega}$, and 522C for alternative descriptions of
$\frak b$)} and the {\bf dominating number}
$\frak d=\cf\NN$\cmmnt{;  and finally I should include $\frak c$
itself as an eleventh cardinal in the list to be examined here}.
I use the notions of Galois-Tukey connection and Tukey function, and the
associated relations $\prGT$, $\equivGT$ and $\prT$, as described in
\S\S512-513.

\leader{522B}{Cicho\'n's diagram} The diagram itself is the following:

\def\tmphrule{\hskip0.4em\raise
2.5pt\hbox{\leaders\hrule\hskip1.7em\hfil}\hskip0.4em}
\def\tmpstrut{\vrule height10.5pt depth5.5pt width0pt}

$$\vbox{\offinterlineskip
\halign{\hfil#\hfil&\hfil#\hfil&\hfil#\hfil&\hfil#\hfil
  &\hfil#\hfil&\hfil#\hfil&\hfil#\hfil
  &\hfil#\hfil&\hfil#\hfil&\hfil#\hfil&\hfil#\hfil\cr
&\tmpstrut&$\cov\Cal N$&\tmphrule&$\non\Cal M$&\tmphrule
  &$\cf\Cal M$&\tmphrule&$\cf\Cal N$&\tmphrule&$\frak c$\cr
&\tmpstrut&\vrule&&\vrule&&\vrule&&\vrule\cr
&\tmpstrut&\vrule&&$\frak b$&\tmphrule&$\frak d$&&\vrule\cr
&\tmpstrut&\vrule&&\vrule&&\vrule&&\vrule\cr
$\tmpstrut\omega_1$&\tmphrule&$\add\Cal N$&\tmphrule
  &$\add\Cal M$&\tmphrule&$\cov\Cal M$&\tmphrule&$\non\Cal N$\cr
}}$$

\noindent The cardinals here increase from bottom left to
top right;  that is,

\Centerline{$\omega_1\le\add\Cal N\le\add\Cal M\le\frak b\le\frak d
\le\cf\Cal M\le\cf\Cal N\le\frak c$,}

\noindent etc.   In addition, we have two equalities:

\Centerline{$\add\Cal M=\min(\frak b,\cov\Cal M)$,
\quad$\cf\Cal M=\max(\frak d,\non\Cal M)$.}

\cmmnt{\noindent In the rest of this section I will prove all the
inequalities
declared here, seeking to demonstrate reasons for the remarkable
symmetry of the diagram.   I will make heavy use of the ideas of \S512.
Of course many of the elementary results can be proved directly without
difficulty;  but for the most interesting part of the argument
(522K-522Q below) Tukey functions seem to be the right way to proceed.

I start with the easiest results.   It will be helpful to have
descriptions of $\frak b$ and $\frak d$ in terms of other partially
ordered sets.
}%end of comment

\leader{522C}{Lemma} (i) On $\NN$ define a relation $\le^*$ by saying that
$f\le^*g$ if
if $\{n:f(n)>g(n)\}$ is finite.   Then $\le^*$ is a pre-order on $\NN$;
$\add(\NN,\le^*)=\frak b$ and $\cf(\NN,\le^*)=\frak d$.

(ii) On $\NN$ define a relation $\preceq$ by saying that $f\preceq g$ if
either $f\le g$ or $\{n:g(n)\le f(n)\}$ is finite.
Then $\preceq$ is a partial order on $\NN$, $\add(\NN,\preceq)=\frak b$
and $\cf(\NN,\preceq)=\frak d$.

(iii) $(\NN,\le^*)\equivT(\NN,\preceq)$.

\proof{{\bf (a)} The checks that $\le^*$ is a pre-order and that $\preceq$
is a partial order are elementary.   Write $\iota$ for the identity map
from $\NN$ to itself.

\medskip

{\bf (b)} For $f\in\NN$ and $A\subseteq\NN$ say that $f\le\strprime A$ if
there is a $g\in A$ such that $f\le g$ (see 512F).   Now
$(\NN,\le\strprime,[\NN]^{\le\omega})\prGT(\NN,\le^*,\NN)$.   \Prf\ For
$g\in\NN$, set
$\psi(g)=\{h:g\le^*h\le^*g\}\in[\NN]^{\omega}$.   If $f\le^*g$,
then $f\le f\vee g\in\psi(g)$, so $f\le\strprime\psi(g)$.   Thus
$(\iota,\psi)$ is a Galois-Tukey connection and
$(\NN,\le\strprime,[\NN]^{\le\omega})\prGT(\NN,\le^*,\NN)$.\ \Qed

\medskip

{\bf (c)} $(\NN,\le^*,\NN)\prGT(\NN,\preceq,\NN)$.   \Prf\ If $f$,
$g\in\NN$ and $f\preceq g$, then $f\le^*g$;  so $(\iota,\iota)$
is a Galois-Tukey connection from
$(\NN,\le^*,\NN)$ to $(\NN,\preceq\NN)$, as in 512Cd.\ \Qed

\medskip

{\bf (d)} If $A\subseteq\NN$ is countable, there is a $\psi(A)\in\NN$
such that $g\preceq\psi(A)$ for every $g\in A$.   \Prf\ If $A$ is empty,
this is trivial.   Otherwise, let $\sequencen{g_n}$ be a sequence
running over $A$, and set $\psi(A)(i)=1+\max_{n\le i}g_n(i)$ for every
$i\in\Bbb N$.\ \Qed

It follows that
$(\NN,\preceq,\NN)\prGT(\NN,\le\strprime,[\NN]^{\le\omega})$.
\Prf\ If $A\in[\NN]^{\le\omega}$ and
$f\le\strprime A$, then there is some $g\in A$ such that $f\le g$, so
that $f\preceq\psi(A)$.   Thus $(\iota,\psi)$ is a Galois-Tukey
connection
from $(\NN,\preceq,\NN)$ to $(\NN,\le\strprime,[\NN]^{\le\omega})$.\
\Qed

\medskip

{\bf (e)} Putting (b)-(d) together, we see that

\Centerline{$(\NN,\le\strprime,[\NN]^{\le\omega})
\equivGT(\NN,\le^*,\NN)\equivGT(\NN,\preceq,\NN)$.}

\noindent In particular, $(\NN,\le^*)\equivT(\NN,\preceq)$ (513Ea).
By 512D,

\Centerline{$\add(\NN,\le\strprime,[\NN]^{\le\omega})
=\add(\NN,\le^*,\NN)=\add(\NN,\preceq,\NN)$,}

\Centerline{$\cov(\NN,\le\strprime,[\NN]^{\le\omega})
=\cov(\NN,\le^*,\NN)=\cov(\NN,\preceq,\NN)$.}

\noindent But by 513Ia we have

\Centerline{$\add(\NN,\le\strprime,[\NN]^{\le\omega})
=\add_{\omega}(\NN)=\frak b$,}

\noindent so $\frak b=\add(\NN,\le^*)=\add(\NN,\preceq)$.   In
the other direction,

$$\eqalignno{\frak d
&=\cf\NN=\cov(\NN,\le,\NN)
\le\max(\omega,\cov(\NN,\le\strprime,[\NN]^{\le\omega}))\cr
\displaycause{512Gf}
&\le\max(\omega,\cov(\NN,\le,\NN))\cr
\displaycause{512Gc}
&=\max(\omega,\frak d)=\frak d.\cr}$$

\noindent So

\Centerline{$\frak d=\cov(\NN,\le\strprime,[\NN]^{\le\omega})
=\cf(\NN,\le^*)=\cf(\NN,\preceq)$,}

\noindent and the proof is complete.
}%end of proof of 522C

\leader{522D}{Proposition} $\frak b\le\frak d$.

\proof{ Use 511He and 522C.
}%end of proof of 522D

\leader{522E}{Proposition} $\add\Cal N\le\cov\Cal N$,
$\add\Cal M\le\cov\Cal M$, $\non\Cal M\le\cf\Cal M$ and
$\non\Cal N\le\cf\Cal N$.

\proof{ We need only observe that both $\Cal M$ and $\Cal N$ are proper
ideals of $\Cal P\Bbb R$ with union $\Bbb R$, and use 511Jc.}

\leader{522F}{Proposition} $\omega_1\le\add\Cal N$ and
$\cf\Cal N\le\frak c$.

\proof{ Of course $\omega_1\le\add\Cal N$ because $\Cal N$ is a
$\sigma$-ideal of sets.    As for $\cf\Cal N$, we know that the family
of negligible Borel sets is cofinal with $\Cal N$ (134Fb) and has at
most $\frak c$ members (Fa), so $\cf\Cal N\le\frak c$.
}%end of proof of 522F

\leader{522G}{Proposition}\cmmnt{ ({\smc Rothberger 1938a})}
$\cov\Cal N\le\non\Cal M$ and $\cov\Cal M\le\non\Cal N$.

\proof{ The point is just that there is a comeager negligible set
$E\subseteq\Bbb R$.   \Prf\ Enumerate $\Bbb Q$ as $\sequencen{q_n}$, and
set

\Centerline{$E
=\bigcap_{n\in\Bbb N}\bigcup_{m\ge n}\ooint{q_n-2^{-n},q_n+2^{-n}}$.
\Qed}

\noindent Because $x\mapsto a+x$ and $x\mapsto a-x$ are
measure-preserving homeomorphisms, $a+E$ is negligible and $a-E$ is
comeager for every $a\in\Bbb R$.   Let $A\subseteq\Bbb R$ be a
non-meager set with cardinal $\non\Cal M$.   Then
$A\cap(a-E)\ne\emptyset$ for every $a\in\Bbb R$, that is,
$\{x+E:x\in A\}$ covers $\Bbb R$;  so $\cov\Cal N\le\#(A)=\non\Cal M$.

For the other inequality, note that $F=\Bbb R\setminus E$ is
conegligible and meager;  so the same argument shows that
$\cov\Cal M\le\non N$.
}%end of proof of 522G

\leader{522H}{Proposition} $\add\Cal M\le\frak b$ and
$\frak d\le\cf\Cal M$.

\proof{{\bf (a)}
Start by choosing a countable base $\Cal U$ for the topology of
$\Bbb R$, not containing $\emptyset$, and enumerate it as
$\sequence{k}{U_k}$.   For each $k\in\Bbb N$ let $\sequence{l}{V_{kl}}$
be a disjoint sequence of non-empty open subsets of $U_k$;  finally,
enumerate $V_{kl}\cap\Bbb Q$ as $\sequence{i}{x_{kli}}$ for each $k$,
$l\in\Bbb N$.

\medskip

{\bf (b)} Fix $f\in\NN$ for the moment.   Set
$E_k(f)=\{x_{kli}:l\in\Bbb N,\,i\le f(l)\}\subseteq U_k$ for each
$k\in\Bbb N$.
This is nowhere dense because if $G$ is a non-empty open set, either
$G\cap\bigcup_{l\in\Bbb N}V_{kl}=\emptyset$ and
$G\cap E_k(f)=\emptyset$, or there is an $l$ such that $G\cap V_{kl}$ is
non-empty, in which case $G\cap V_{kl}\cap E_k(f)$ is finite and
$G\setminus\overline{E_k(f)}\supseteq G\cap V_{kl}\setminus E_k(f)$ is
non-empty.

Now choose $\sequencen{k_n}$, $\sequencen{l_n}$ inductively as follows.
Given $\langle k_i\rangle_{i<n}$, $\bigcup_{i<n}E_{k_i}(f)$ is
nowhere dense, so there is an $l_n\in\Bbb N$ such that
$\overline{U}_{l_n}\subseteq U_n\setminus\bigcup_{i<n}E_{k_i}(f)$.
Now if $U_n\subseteq\overline{\bigcup_{i\le n}U_{l_i}}$, set $k_n=0$;
otherwise, take $k_n$ such that
$U_{k_n}\subseteq U_n\setminus\overline{\bigcup_{i\le n}U_{l_i}}$, and continue.
At the end of the induction, set
$\phi(f)=\overline{\bigcup_{n\in\Bbb N}E_{k_n}}$.

The construction ensures that $\overline{U}_{l_n}\cap E_{k_m}=\emptyset$
for all $m$ and $n$, so that $U_{l_n}$ is always a non-empty open subset
of $U_n\setminus\phi(f)$;
accordingly $\phi(f)$ is nowhere dense.   If  $G\subseteq\Bbb R$ is a
non-empty open set meeting $\phi(f)$, there
is a $k\in\Bbb N$ such that $E_k(f)\subseteq G\cap\phi(f)$.   \Prf\ Let
$n\in\Bbb N$ be such that $U_n\subseteq G$ and
$U_n\cap\phi(f)\ne\emptyset$.   Then there is an $i\in\Bbb N$ such that
$U_n\cap E_{k_i}(f)\ne\emptyset$;
as $E_{k_i}(f)\cap\overline{U}_{l_j}$ is empty for every $j$,
$U_n\not\subseteq\overline{\bigcup_{j\le n}U_{l_j}}$ and

\Centerline{$E_{k_n}(f)\subseteq U_{k_n}\subseteq U_n\subseteq G$,}

\noindent so we can take $k=k_n$.\ \Qed

\medskip

{\bf (c)} In the other direction, given $M\in\Cal M$, choose a sequence
$\sequencen{F_n(M)}$ of closed nowhere dense closed sets covering $M$.
For $k$, $l$, $n\in\Bbb N$ set $g_{Mnk}(l)=\min\{j:x_{klj}\notin F_n(M)\}$.
Since $\add(\NN,\preceq)=\add_{\omega}(\NN,\le)\ge\omega_1$
(522C(ii) with 513Ib, or use the construction in part (d) of the proof of
522C), $\{g_{Mnk}:n$, $k\in\Bbb N\}$ is bounded above in $(\NN,\preceq)$;
take $\psi(M)\in\NN$ such that $g_{Mnk}\le\psi(M)$ for all $n$ and $k$.

\medskip

{\bf (d)} Now $(\phi,\psi)$ is a Galois-Tukey connection from
$(\NN,\preceq,\NN)$ to
$(\Cal M,\subseteq,\Cal M)$.   \Prf\ Suppose that $f\in\NN$ and
$M\in\Cal M$ are such that $\phi(f)\subseteq M$.   Because $\phi(f)$ is
closed and not empty and included in $\bigcup_{n\in\Bbb N}F_n(M)$,
Baire's theorem (3A3G or 4A2Ma) tells us that there are $n\in\Bbb N$ and
an open set $G$ such that $\emptyset\ne G\cap\phi(f)\subseteq F_n(M)$.
By the last remark in (b), there is a $k\in\Bbb N$ such that
$E_k(f)\subseteq G\cap\phi(f)$.
But this means that, for any $l\in\Bbb N$, $x_{kli}\in G\cap\phi(f)$
for every
$i\le f(l)$, while if $j=g_{Mnk}(l)$ then $x_{klj}\notin G\cap\phi(f)$.   So
$f(l)\le g_{Mnk}(l)$ for every $l$, and $f\preceq g_{Mnk}\preceq\psi(M)$.\ \Qed

\medskip

{\bf (e)} So in fact $\phi$ is a Tukey function from $(\NN,\preceq)$ to
$(\Cal M,\subseteq)$ (513Ea).   It follows at once that

\Centerline{$\add\Cal M\le\add(\NN,\preceq)$,
\quad$\cf(\NN,\preceq)\le\cf\Cal M$}

\noindent (513Ee), that is, $\add\Cal M\le\frak b$ and $\frak d\le\cf\Cal M$,
by 522C(ii).
}%end of proof of 522H

\leader{522I}{Proposition} $\frak b\le\non\Cal M$ and
$\cov\Cal M\le\frak d$.

\proof{ Again Let $\preceq$ be the partial order on $\NN$ described in
522C(ii).
Then $(\Bbb R\setminus\Bbb Q,\in,\Cal M)\prGT(\NN,\preceq,\NN)$.   \Prf\
Let $\phi:\Bbb R\setminus\Bbb Q\to\NN$ be a homeomorphism (4A2Ub\Latereditions).
For $f\in\NN$, set $K_f=\{g:g\le f\}$;  then $K_f$ is compact, so
$\phi^{-1}[K_f]$ is compact.   Because $\phi^{-1}[K_f]$ is disjoint from
$\Bbb Q$, it is nowhere dense.   Set

\Centerline{$\psi(f)
=\bigcup\{\phi^{-1}[K_g]:g\in\NN$, $\{n:g(n)\ne f(n)\}$ is finite$\}\}$.}

\noindent Because there are only countably many functions eventually
equal to $f$, $\psi(f)\in\Cal M$.

Suppose that $x\in\Bbb R\setminus\Bbb Q$ and $f\in\NN$ are such that
$\phi(x)\preceq f$.   Set $g=\phi(x)\vee f$;  then $g(n)=f(n)$ for all
but
finitely many $n$, and $\phi(x)\le g$, so
$x\in\phi^{-1}[K_g]\subseteq\psi(f)$.   This shows that $(\phi,\psi)$ is
a Galois-Tukey connection from $(\Bbb R\setminus\Bbb Q,\in,\Cal M)$ to
$(\NN,\preceq,\NN)$, so that
$(\Bbb R\setminus\Bbb Q,\in,\Cal M)\prGT(\NN,\preceq,\NN)$.\ \Qed

It follows (using 522C(ii) and 512D) that

$$\eqalign{\frak b
&=\add(\NN,\preceq,\NN)
\le\add(\Bbb R\setminus\Bbb Q,\in,\Cal M)\cr
&=\min\{\#(A):A\subseteq\Bbb R\setminus\Bbb Q,\,A\notin\Cal M\}
\le\non\Cal M,\cr}$$

$$\eqalign{\frak d&
=\cov(\NN,\preceq,\NN)
\ge\cov(\Bbb R\setminus\Bbb Q,\in,\Cal M)
=\min\{\#(\Cal A):\Cal A\subseteq\Cal M,\,
   \Bbb R\setminus\Bbb Q\subseteq\bigcup\Cal A\}\cr&
=\min\{\#(\Cal A\cup\{\Bbb Q\}):\Cal A\subseteq\Cal M,\,
   \Bbb R=\bigcup(\Cal A\cup\{\Bbb Q\})\}
\ge\cov\Cal M.\cr}$$
}%end of proof of 522I

\leader{522J}{Theorem}\cmmnt{ (see {\smc Truss 77} and
{\smc Miller 81})} $\add\Cal M=\min(\frak b,\cov\Cal M)$ and
$\cf\Cal M=\max(\frak d,\non\Cal M)$.

\proof{ My aim this time is to prove that

\Centerline{$(\Cal M,\subseteq,\Cal M)
\prGT(\Bbb R,\in,\Cal M)^{\perp}
\ltimes(\NN,\le\strprime,[\NN]^{\le\omega})$,}

\noindent defining $\le\strprime$ as in the proof of 522H and
$\ltimes$ as in 512I.

\medskip

{\bf (a)} Let $\sequencen{q_n}$ be a sequence running over $\Bbb Q$ with
cofinal repetitions.   For
$f\in\NN$, set

\Centerline{$E_f
=\Bbb R\setminus\bigcap_{n\in\Bbb N}\bigcup_{m\ge n}
  \ooint{q_m-2^{-f(m)},q_m+2^{-f(m)}}$,}

\noindent so that $E_f$ is a meager set disjoint from $\Bbb Q$.
Observe that if $\sequencen{H_n}$ is any sequence of closed sets
disjoint from $\Bbb Q$, then there is an $f\in\NN$ such that
$\bigcup_{n\in\Bbb N}H_n\subseteq E_f$.   \Prf\ For each $n\in\Bbb N$,
let $f(n)$ be such that
$\ooint{q_n-2^{-f(n)},q_n+2^{-f(n)}}$ does not meet
$\bigcup_{m\le n}H_m$.\ \Qed

For $M\in\Cal M$, choose a sequence $\sequencen{F_n(M)}$ of nowhere
dense closed sets covering $M$.   For $x\in\Bbb R$, if
$\Bbb Q\cap(\bigcup_{n\in\Bbb N}F_n(M)-x)$ is not empty, set
$p_M(x)(n)=0$ for every $n\in\Bbb N$;  otherwise, take $p_M(x)=f$ for
some $f\in\NN$ such that $E_f\supseteq\bigcup_{n\in\Bbb N}F_n(M)-x$.
Now set $\phi(M)=(\bigcup_{n\in\Bbb N}F_n(M)+\Bbb Q,p_M)$.   This
defines $\phi:\Cal M\to\Cal M\times(\NN)^{\Bbb R}$.

\medskip

{\bf (b)} In the other direction, define
$\psi:\Bbb R\times[\NN]^{\le\omega}\to\Cal M$ by setting
$\psi(x,B)=\bigcup_{f\in B}(x+E_f)$ for $x\in\Bbb R$ and
$B\in[\NN]^{\le\omega}$.   Now $(\phi,\psi)$ is a Galois-Tukey
connection from $(\Cal M,\subseteq,\Cal M)$ to
$(\Bbb R,\in,\Cal M)^{\perp}
 \ltimes(\NN,\le\strprime,[\NN]^{\le\omega})$.
\Prf\ $(\Bbb R,\in,\Cal M)^{\perp}=(\Cal M,\not\ni,\Bbb R)$, so

\Centerline{$(\Bbb R,\in,\Cal M)^{\perp}
  \ltimes(\NN,\le\strprime,[\NN]^{\le\omega})
=(\Cal M\times(\NN)^{\Bbb R},T,\Bbb R\times[\NN]^{\le\omega})$,}

\noindent where $((M,p),(x,B))\in T$ iff $x\notin M$ and $p(x)\le g$ for
some $g\in B$.   Now suppose that $M\in\Cal M$ and
$(x,B)\in\Bbb R\times[\NN]^{\le\omega}$ are such that
$(\phi(M),(x,B))\in T$.   Then
$x\notin\bigcup_{n\in\Bbb N}F_n(M)+\Bbb Q$, so
$\Bbb Q\cap(\bigcup_{n\in\Bbb N}F_n(M)-x)=\emptyset$, while
$p_M(x)\le g$ for some $g\in B$.   But this means that

\Centerline{$E_g\supseteq E_{p_M(x)}\supseteq\bigcup_{n\in\Bbb N}F_n(M)-x
\supseteq M-x$,
\quad$M\subseteq E_g+x\subseteq\psi(x,B)$.}

\noindent As $M$ and $(x,B)$ are arbitrary, $(\phi,\psi)$
is a Galois-Tukey connection, as claimed.\ \Qed

\medskip

{\bf (c)} It follows that

$$\eqalignno{\cf\Cal M
&=\cov(\Cal M,\subseteq,\Cal M)\cr
\displaycause{512Ea, as before}
&\le\cov((\Bbb R,\in,\Cal M)^{\perp}
  \ltimes(\NN,\le\strprime,[\NN]^{\le\omega}))\cr
&=\max(\cov(\Bbb R,\in,\Cal M)^{\perp},
  \cov(\NN,\le\strprime,[\NN]^{\le\omega}))\cr
\displaycause{512Jb}
&=\max(\add(\Bbb R,\in,\Cal M),\frak d)
=\max(\non\Cal M,\frak d)\cr}$$

\noindent by 512Ed and the calculation in part (e) of the proof of 522H.
On the other hand

$$\eqalignno{\min(\cov\Cal M,\frak b)
&=\min(\cov\Cal M,\add(\NN,\le\strprime,[\NN]^{\le\omega}))\cr
&=\min(\cov(\Bbb R,\in,\Cal M),
  \add(\NN,\le\strprime,[\NN]^{\le\omega}))\cr
&=\min(\add(\Bbb R,\in,\Cal M)^{\perp},
  \add(\NN,\le\strprime,[\NN]^{\le\omega}))\cr
&=\add((\Bbb R,\in,\Cal M)^{\perp}
  \ltimes(\NN,\le\strprime,[\NN]^{\le\omega}))\cr
\displaycause{512Jc}
&\le\add(\Cal M,\subseteq,\Cal M)\cr
\displaycause{(b) above and 512Db}
&=\add\Cal M\cr}$$

\noindent(512Ea, as ever).   Since we already know from 522E and 522H that
$\add\Cal M\le\min(\frak b,\cov\Cal M)$ and that
$\max(\frak d,\non\Cal M)\penalty-100\le\cf\Cal M$, we have the result.
}%end of proof of 522J

\leader{522K}{Localization}\cmmnt{ The last step in proving the facts
announced in 522B depends on the following construction.}   Let $I$ be
any set.   Write $\Cal S_I$ for the family of sets
$S\subseteq\Bbb N\times I$ such that each vertical section
$S[\{n\}]$ has at most $2^n$ members.    For $f\in I^{\Bbb N}$ and
$S\subseteq\Bbb N\times I$ say that $f\subseteq^*S$ if
$\{n:n\in\Bbb N,\,(n,f(n))\notin S\}$ is finite\cmmnt{;  that is,
$f\setminus S$ is finite, if we identify $f$ with its graph}.
I will say that
the supported relation $(I^{\Bbb N},\subseteq^*,\Cal S_I)$ is the {\bf
$I$-localization relation}.   \cmmnt{By far the most important
case (and the only one needed in this section) is when $I$ is
countably infinite;  when $I=\Bbb N$ I will generally write $\Cal S$ rather
than $\Cal S_{\Bbb N}$.}

\cmmnt{Members of $\Cal S_I$, or similar sets, are sometimes called
{\bf slaloms}.   The particular formula `$\#(S[\{n\}])\le 2^n$' is
convenient
for the results of this section, but it is worth knowing that all
functions diverging to $\infty$ give rise to equivalent partially
ordered sets.}

\leader{*522L}{Lemma} Let $I$ be an infinite set.
For any $\alpha\in\NN$ write

\Centerline{$\Cal S_I^{(\alpha)}=\{S:S\subseteq \Bbb N\times I,\,
  \#(S[\{n\}])\le\alpha(n)$ for every $n\in\Bbb N\}$.}

\noindent Then $(I^{\Bbb N},\subseteq^*,\Cal S_I^{(\alpha)})
\equivGT(I^{\Bbb N},\subseteq^*,\Cal S_I^{(\beta)})$ whenever $\alpha$,
$\beta\in\NN$ and
$\lim_{n\to\infty}\alpha(n)=
\penalty-100\lim_{n\to\infty}\beta(n)
\penalty-100=\infty$.

\proof{ Let $g\in\NN$ be a strictly increasing sequence such that
$\beta(n)\le\alpha(i)$ whenever $n\in\Bbb N$ and $i\ge g(n)$, and let
$h_n:I\to I^{g(n+1)\setminus g(n)}$ be a bijection for each
$n$.   Define $\phi:I^{\Bbb N}\to I^{\Bbb N}$ by setting
$\phi(f)(n)=h_n^{-1}(f\restr g(n+1)\setminus g(n))$ for $f\in I^{\Bbb N}$
and $n\in\Bbb N$.   Define
$\psi:\Cal S_I^{(\beta)}\to\Cal P(\Bbb N\times I)$ by setting
$\psi(S)=\bigcup_{(n,i)\in S}h_n(i)$, identifying each
$h_n(i)\in I^{g(n+1)\setminus g(n)}
\subseteq(g(n+1)\setminus g(n))\times I$ with a subset of
$\Bbb N\times I$.   Now for $g(n)\le j<g(n+1)$,
$\psi(S)[\{j\}]=\{h_n(i)(j):i\in S[\{n\}]\}$ has at most
$\beta(n)\le\alpha(j)$ members, while $\psi(S)[\{j\}]=\emptyset$ for
$j<g(0)$, so $\psi(S)\in\Cal S_I^{(\alpha)}$ for every
$S\in\Cal S_I^{(\beta)}$.

If $f\in I^{\Bbb N}$ and $S\in\Cal S_I^{(\beta)}$ and
$\phi(f)\subseteq^*S$, then
there is an $n_0\in\Bbb N$ such that $\phi(f)(n)\in S[\{n\}]$ for every
$n\ge n_0$.   So

\Centerline{$f\restr g(n+1)\setminus g(n)
=h_n(\phi(f)(n))\subseteq\psi(S)$}

\noindent for every $n\ge n_0$, $(m,f(m))\in\psi(S)$ for every
$m\ge g(n_0)$ and $f\subseteq^*\psi(S)$.   This means that $(\phi,\psi)$
is a Galois-Tukey connection from $(\NN,\subseteq^*,\Cal S_I^{(\alpha)})$
to $(\NN,\subseteq^*,\Cal S_I^{(\beta)})$.   Similarly,
$(\NN,\subseteq^*,\Cal S_I^{(\beta)})
\prGT(\NN,\subseteq^*,\Cal S_I^{(\alpha)})$
and the two supported relations are equivalent.
}%end of proof

\leader{522M}{Proposition} Let $(\NN,\subseteq^*,\Cal S)$ be the
$\Bbb N$-localization relation.   Then
$(\NN,\subseteq^*,\Cal S)\penalty-100\equivGT(\Cal N,\subseteq,\Cal N)$.

\proof{{\bf (a)} Let $\langle G_{ij}\rangle_{i,j\in\Bbb N}$ be a
stochastically independent family of open subsets of $[0,1]$ such that the
Lebesgue measure $\mu G_{ij}$ of $G_{ij}$ is $2^{-i}$ for all $i$,
$j\in\Bbb N$.   For $f\in\NN$, set
$\phi(f)=\bigcap_{n\in\Bbb N}\bigcup_{m\ge n}G_{m,f(m)}$.
Then $\phi(f)$ is negligible.

For each $E\in\Cal N$, choose a non-empty compact self-supporting set
$K_E\subseteq[0,1]\setminus E$ (416Dc).   Let $\sequencen{W_{En}}$
enumerate a
base for the relative topology on $K_E$ not containing $\emptyset$;
because $K_E$ is self-supporting, no $W_{En}$ is negligible.   Set

\Centerline{$I_{Eni}=\{j:j\in\Bbb N$, $W_{En}\cap G_{ij}=\emptyset\}$}

\noindent for $n$, $i\in\Bbb N$.   Then

\Centerline{$\sum_{i=0}^{\infty}2^{-i}\#(I_{Eni})
=\sum\{\mu G_{ij}:i$, $j\in\Bbb N$, $G_{ij}\cap W_{En}=\emptyset\}$}

\noindent is finite, by the Borel-Cantelli lemma (273K).   For each $n$,
let $k(E,n)\in\Bbb N$ be such that
$2^{-i}\#(I_{Eni})\le 2^{-n-1}$ for $i\ge k(E,n)$, and set

\Centerline{$\psi(E)
=\bigcup_{n\in\Bbb N}\{(i,j):i,\,j\in\Bbb N,\,i\ge k(E,n),
  \,j\in I_{Eni}\}$.}

\noindent Then

$$\#(\{j:(i,j)\in\psi(E)\})
\le\sum_{n\in\Bbb N,k(E,n)\le i}\#(I_{Eni})
\le\sum_{n\in\Bbb N,k(E,n)\le i}2^{-n-1}2^i
\le 2^i$$

\noindent for every $i\in\Bbb N$, so $\psi(E)\in\Cal S$.

Now $(\phi,\psi)$ is a Galois-Tukey connection from
$(\NN,\subseteq^*,\Cal S)$ to $(\Cal N,\subseteq,\Cal N)$.   \Prf\
Suppose that $f\in\NN$ and $E\in\Cal N$ are such that
$\phi(f)\subseteq E$.   Then
$K_E\cap\bigcap_{n\in\Bbb N}\bigcup_{m\ge n}G_{m,f(m)}=\emptyset$.
By Baire's theorem, there is some $m\in\Bbb N$ such that
$\bigcup_{i\ge m}G_{i,f(i)}\cap K_E$ is not dense in $K_E$, that is,
there is an $n\in\Bbb N$ such that
$W_{En}\cap\bigcup_{i\ge m}G_{i,f(i)}=\emptyset$ so $f(i)\in I_{Eni}$
for every $i\ge m$.   But this means that $(i,f(i))\in\psi(E)$ for every
$i\ge\max(m,k(E,n))$, so that $f\subseteq^*\psi(E)$.   As $f$ and $E$
are arbitrary, we have the result.\ \Qed

Thus $(\NN,\subseteq^*,\Cal S)\prGT(\Cal N,\subseteq,\Cal N)$.

\medskip

{\bf (b)} Let $\Cal H$ be the family of finite unions of bounded open
intervals in $\Bbb R$ with rational endpoints.   Then $\Cal H$ is
countable.   For each $n\in\Bbb N$, let $\sequence{i}{H_{ni}}$ be an
enumeration of $\{H:H\in\Cal H,\,\mu H\le 4^{-n}\}$.   Now for each
$E\in\Cal N$ there is an $f\in\NN$ such that
$E\subseteq\bigcap_{n\in\Bbb N}\bigcup_{m\ge n}H_{m,f(m)}$.   \Prf\ For
each $n\in\Bbb N$, let $\sequence{i}{J_{ni}}$ be a sequence of open
intervals with rational endpoints
such that $E\subseteq\bigcup_{i\in\Bbb N}J_{ni}$ and
$\sum_{i=0}^{\infty}\mu J_{ni}\le 2^{-n-1}$.   Re-enumerating
$\langle J_{ni}\rangle_{n\in\Bbb N,i\in\Bbb N}$ as $\sequence{i}{J_i}$,
we have a sequence of open intervals with rational endpoints such that
$\sum_{i=0}^{\infty}\mu J_i\le 1$ and $E\subseteq\bigcup_{i\ge n}J_i$
for every $n$.   Let $\sequencen{k(n)}$ be a strictly increasing
sequence such that $k(0)=0$ and
$\sum_{i=k(n)}^{\infty}\mu J_i\le 4^{-n}$ for every $n\in\Bbb N$.   Then
$V_n=\bigcup_{k(n)\le i<k(n+1)}J_i$ belongs to $\Cal H$ and has measure
at most $4^{-n}$ for each $n$, so we can define $f\in\NN$ by saying that
$H_{n,f(n)}=V_n$ for each $n$, and we shall have an appropriate
function.\ \Qed

We can therefore find a function $\phi:\Cal N\to\NN$ such that
$E\subseteq\bigcap_{n\in\Bbb N}\bigcup_{m\ge n}H_{m,\phi(E)(m)}$ for
every $E\in\Cal N$.   In the reverse direction, define

\Centerline{$\psi(S)
=\bigcap_{n\in\Bbb N}\bigcup\{H_{mi}:m\ge n,\,(m,i)\in S\}$}

\noindent for $S\in\Cal S$;  because

\Centerline{$\sum_{(m,i)\in S}\mu H_{mi}
\le\sum_{m=0}^{\infty}2^m4^{-m}<\infty$,}

\noindent $\psi(S)\in\Cal N$.

Now $(\phi,\psi)$ is a Galois-Tukey connection from
$(\Cal N,\subseteq,\Cal N)$ to $(\NN,\subseteq^*,\Cal S)$.   \Prf\ If
$E\in\Cal N$ and $S\in\Cal S$ are such that $\phi(E)\subseteq^*S$, then

\Centerline{$E
\subseteq\bigcap_{n\in\Bbb N}\bigcup_{m\ge n}H_{m,\phi(E)(m)}
\subseteq\bigcap_{n\in\Bbb N}\bigcup_{m\ge n,(m,i)\in S}H_{mi}
=\psi(S)$.  \Qed}

\noindent So $(\Cal N,\subseteq,\Cal N)\prGT(\NN,\subseteq^*,\Cal S)$
and the proof is complete.
}%end of proof of 522M

\leader{522N}{Lemma} Let $X$ be a topological
space with a countable $\pi$-base.   Then there is for each $n\in\Bbb N$
a countable family $\Cal U_n$ of open subsets of $X$ such
that every dense open subset of $X$
includes some member of $\Cal U_n$ and $\bigcap\Cal V\ne\emptyset$ for every
$\Cal V\in[\Cal U_n]^{\le n}$.

\proof{ Induce on $n$.   Start by taking $\Cal U$ to be a countable
$\pi$-base for the topology of $X$ which is closed under finite unions.
Set $\Cal U_0=\{\emptyset\}$.   For the inductive step to $n+1$, let
$\sequence{i}{H_i}$ be a sequence running over $\Cal U_n$, and set

\Centerline{$\Cal J_i
=\{J:J\subseteq i,\,\bigcap_{j\in J}H_j\ne\emptyset\}$}

\noindent for $i\in\Bbb N$,

\Centerline{$\Cal U_{n+1}
=\{U\cup H_i:i\in\Bbb N,\,U\in\Cal U$,
$U\cap\bigcap_{j\in J}H_j\ne\emptyset$ whenever $J\in\Cal J_i\}$.}

\noindent Then $\Cal U_{n+1}$ is a countable family of open sets.   If
$G\subseteq X$ is a dense open set, let $i\in\Bbb N$ be such that
$H_i\subseteq G$.   Then
$\Cal J_i$ is finite, so we can find a $U\in\Cal U$ such that
$U\subseteq G$ and $U\cap\bigcap_{j\in J}U_j\ne\emptyset$ for every
$J\in\Cal J_i$;  then $U\cup H_i$ belongs to $\Cal U_{n+1}$ and is
included in $G$.
If $\Cal V\subseteq\Cal U_{n+1}$ and $\#(\Cal V)\le n+1$, then if
$\Cal V$ is empty we certainly have $\bigcap\Cal V\ne\emptyset$.
Otherwise, express $\Cal V$ as $\{U_k\cup H_{i(k)}:k\le n\}$ where
$U_k\cap\bigcap_{j\in J}H_j\ne\emptyset$ whenever $J\in\Cal J_{i(k)}$;
do this in such a way that
$i(k)\le i(n)$ for every $k<n$.   By the inductive hypothesis,
$\bigcap_{k<n}H_{i(k)}\ne\emptyset$;  if $i(k)=i(n)$ for some $k<n$,
then of course $\bigcap_{k\le n}H_{i(k)}\ne\emptyset$;  otherwise,
$U_n\cap\bigcap_{k<n}H_{i(k)}\ne\emptyset$.   In either case,
$\bigcap\Cal V$ is non-empty.   So $\Cal U_{n+1}$ has the required
properties and the induction continues.
}%end of proof of 522N
% for `endowments' see mt52bits

\leader{522O}{Proposition} Let $(\NN,\subseteq^*,\Cal S)$ be the
$\Bbb N$-localization relation.   Then
$(\Cal M,\subseteq,\Cal M)\penalty-100\prGT(\NN,\subseteq^*,\Cal S)$.

\proof{ Let $\sequencen{U_n}$ enumerate
a $\pi$-base for the topology of
$\Bbb R$ not containing $\emptyset$.   By 522N, there is for each
$n\in\Bbb N$ a countable family $\Cal V_n$
of open subsets of $U_n$ such that $\bigcap\Cal V\ne\emptyset$ for every
$\Cal V\in[\Cal V_n]^{\le 2^n}$ and every dense open subset of $U_n$
includes some member of $\Cal V_n$.   Enumerate $\Cal V_n$
as $\sequence{m}{V_{nm}}$.

For each $M\in\Cal M$, let $\sequencen{F_n(M)}$ be a non-decreasing
sequence of nowhere dense sets covering $M$, and let $\phi(M)\in\NN$ be
such that $F_n(M)\cap V_{n,\phi(M)(n)}=\emptyset$ for every $n$.
In the other direction, for $S\in\Cal S$ set

\Centerline{$\psi(S)
=\Bbb R\setminus\bigcap_{n\in\Bbb N}
  \bigcup_{m\ge n}(U_m\cap\bigcap_{i\in S[\{m\}]}V_{mi})$;}

\noindent then because $\bigcap_{i\in S[\{m\}]}V_{mi}$ is non-empty for
every $n$, $\bigcup_{m\ge n}(U_m\cap\bigcap_{i\in S[\{m\}]}V_{mi})$
is a dense open set for every $n$, and $\psi(S)$ is meager.

Now $(\phi,\psi)$ is a Galois-Tukey connection from
$(\Cal M,\subseteq,\Cal M)$ to $(\NN,\subseteq^*,\Cal S)$.   \Prf\
Suppose that $M\in\Cal M$ and $S\in\Cal S$ are such that
$\phi(M)\subseteq^*S$.   Let $n\in\Bbb N$ be such that
$\phi(M)(k)\in S[\{k\}]$ for every $k\ge n$.   Then

\Centerline{$F_m(M)\cap\bigcap_{i\in S[\{k\}]}V_{ki}
\subseteq F_k(M)\cap V_{k,\phi(M)(k)}=\emptyset$}

\noindent whenever $k\ge\max(m,n)$, so

\Centerline{$F_m(M)
\subseteq\Bbb R\setminus\bigcup_{k\ge\max(m,n)}\bigcap_{i\in
S[\{k\}]}V_{ki}
\subseteq\psi(S)$}

\noindent for every $m$, and $M\subseteq\psi(S)$.\ \Qed

So we have the result.
}%end of proof of 522O

\leader{522P}{Corollary} $\Cal M\prT\Cal N$.

\proof{ Putting 522M and 522O and 512Cb together, we see that
$(\Cal M,\subseteq,\Cal M)
\prGT(\Cal N,\penalty-100\subseteq,\penalty-100\Cal N)$, that is,
$\Cal M\prT\Cal N$.
}%end of proof of 522P

\leader{522Q}{Theorem}\cmmnt{ ({\smc Bartoszy\'nski 84},
{\smc Raisonnier \& Stern 85})} $\add\Cal N\le\add\Cal M$ and
$\cf\Cal M\le\cf\Cal N$.

\proof{ 522P, 513Ee.
}%end of proof of 522Q

\leader{522R}{The exactness of Cicho\'n's diagram} The list of
inequalities displayed in Cicho\'n's diagram is complete in the
following sense:  it is known that all assignments of the values
$\omega_1$, $\omega_2$ to the eleven cardinals of the diagram which are
allowed by the diagram together with the two equalities
$\add\Cal M=\min(\frak b,\cov\Cal M)$,
$\cf\Cal M=\max(\frak d,\non\Cal M)$ are relatively consistent with the
axioms of ZFC.   \cmmnt{So, for instance, it is possible to have

\Centerline{$\omega_1=\add\Cal N=\cov\Cal N=\add\Cal M=\frak b
=\non\Cal M$,}

\Centerline{$\cov\Cal M=\frak d=\cf\Cal M=\non\Cal N=\cf\Cal N
=\frak c=\omega_2$.}

\noindent In \S\S552 and 554 below I will describe forcing constructions
exhibiting a few of these combinations;
for the rest, I refer you to {\smc Bartoszy\'nski \& Judah 95},
\S\S5.2, 7.5 and 7.6.   I remark also that not all
the forcing methods used are effective beyond
$\omega_2$, so that if we allow $\frak c=\omega_3$ then some puzzles
remain.}%end of comment

\leader{522S}{The cardinals $\non\Cal M$,
\dvrocolon{$\cov\Cal M$}}\cmmnt{ All the
cardinals in Cicho\'n's diagram appear in many different ways in
set-theoretic real analysis.   But $\add\Cal N$, the additivity of
Lebesgue measure, the bounding number $\frak b$, and $\cov\Cal M$, the
Nov\'ak number of $\Bbb R$, seem
to be particularly important.   The additivity of measure will play a
large role in the next section.   Here I will give two striking
characterizations of $\cov\Cal M$ and a dual characterization of
$\non\Cal M$.

\medskip

\noindent}{\bf Theorem} (a) $n(\Bbb R)=\cov\Cal M=\frakmctbl$.

(b) \cmmnt{({\smc Bartoszy\'nski 87})} $\cov\Cal M$ is the least
cardinal of any set $A\subseteq\NN$ such
that for every $g\in\NN$ there is an $f\in A$ such that $f(n)\ne g(n)$
for every $n\in\Bbb N$.

(c)\dvAnew{2015} ({\smc Bartoszy\'nski 87}) $\non\Cal M$ is the
least cardinal of any set $A\subseteq\NN$ such that for every $g\in\NN$
there is an $f\in A$ such that $\{n:f(n)=g(n)\}$ is infinite.

\proof{{\bf (a)} Because $\Bbb R$ is a Baire space, the Nov\'ak number
$n(\Bbb R)$ is equal to $\cov\Cal M$ (512Eb).   By 517P(d-ii) or
517P(d-iii), $n(\Bbb R)=\frakmctbl$.

\medskip

{\bf (b)} Let $\kappa$ be the smallest cardinal of any $A\subseteq\NN$
such that for every $g\in\NN$ there is an $f\in A$ such that
$f\cap g=\emptyset$, identifying the functions $f$ and $g$ with their
graphs in $\Bbb N\times\Bbb N$.

\medskip

\quad{\bf (i)} $\cov\Cal M\le\kappa$.   \Prf\
Suppose that $A\subseteq\NN$ and that $\#(A)<\cov\Cal M$.
Set $P=\bigcup_{n\in\Bbb N}\BbbN^n$, ordered by extension of functions.
Then $P$ is a non-empty countable partially ordered set.   For each
$f\in A$ set
$Q_f=\{p:p\in P,\,p\cap f\ne\emptyset\}$;  then $Q_f$ is cofinal with
$P$.   Set $\Cal Q=\{Q_f:f\in A\}$.   Then

\Centerline{$\#(\Cal Q)\le\#(A)<\cov\Cal M=\frakmctbl
\le\frak m^{\uparrow}(P)$,}

\noindent so there is an upwards-linked $R\subseteq P$ meeting every
member of $\Cal Q$.   Now $g_0=\bigcup R\subseteq\Bbb N\times\Bbb N$ is
a function;   taking $g\in\NN$ to be any extension of $g_0$ to the whole
of $\Bbb N$, $g\cap f\ne\emptyset$ for every $f\in A$.

As $A$ is arbitrary, this shows that $\kappa\ge\cov\Cal M$.\ \Qed

In particular, $\kappa\ge\omega_1$, as can also be seen by elementary
arguments.

\medskip

\quad{\bf (ii)} Let $\sequencen{K_n}$ be any sequence of non-empty
countable sets, and write $F$ for the set of all functions $f$ such that
$\dom f$ is an infinite subset of $\Bbb N$ and $f(n)\in K_n$ for every
$n\in\dom f$.   Then if $A\in[F]^{<\kappa}$ there is a
$g\in\prod_{n\in\Bbb N}K_n$ such that $f\cap g\ne\emptyset$ for every
$f\in A$.   \Prf\ For each $n\in\Bbb N$, let $F_n$ be
$\bigcup\{\prod_{i\in I}K_i:I\in[\Bbb N]^{n+1}\}$.
For $f\in F$ and $n\in\Bbb N$ take any $(n+1)$-element subset of $f$ and call
it $p_f(n)$, so that $p_f(n)\in F_n$.   Now each $F_n$ is countably
infinite, and

\Centerline{$A'=\{p_f:f\in A\}\subseteq\prod_{n\in\Bbb N}F_n\cong\NN$}

\noindent has cardinal less than $\kappa$, so there is a
$\phi\in\prod_{n\in\Bbb N}F_n$ such that $\phi\cap p_f\ne\emptyset$ for
every $f\in A$.

Now choose $\sequence{k}{i_k}$ inductively so that
$i_k\in\dom\phi(k)\setminus\{i_j:j<k\}$ for each $k\in\Bbb N$, and take
$g\in\prod_{n\in\Bbb N}K_n$ such that $g(i_k)=\phi(k)(i_k)$ for every
$k$.   Then for any $f\in A$ there is a $k\in\Bbb N$ such that
$\phi(k)=p_f(k)\subseteq f$, so that $g(i_k)=f(i_k)$ and
$f\cap g\ne\emptyset$, as required.\ \Qed

\medskip

\quad{\bf (iii)} If $A\subseteq\NN$ and $f_0\in\NN$ and $\#(A)<\kappa$,
then there is a function $g\in\NN$ such that $g(0)>0$,             		
$g(n+1)\ge f_0(g(n))$ for every $n$ and
$\{n:f(g(n))\le g(n+1)\}$ is infinite for every $f\in A$.   \Prf\ For
$f\in A$ set

\Centerline{$f^*(0)=0$,
\quad$\tilde f(n)=\max_{i\le n}f(i)$,
\quad$f^*(n+1)=n+\tilde f(\tilde f(f^*(n)))$}

\noindent for each $n$, so that $f\le\tilde f$, $\tilde f$ and $f^*$ are
non-decreasing, and $f^*$ is unbounded.   Consider
$B=\{f^*\restr\Bbb N\setminus n:f\in A,\,n\in\Bbb N\}$;  then
$\#(B)\le\max(\#(A),\omega)<\kappa$,
so by (ii) (or otherwise) there is
an $h\in\NN$ meeting every member of $B$.   Now $h\cap f^*$ is infinite
for every $f\in A$.   Set

\Centerline{$g(0)=1+h(0)$,
\quad$g(n+1)=1+\max_{i\le n+1}h(i)+\max_{i\le n}f_0(g(i))$}

\noindent for $n\in\Bbb N$, so that
$h(n)<g(n)$ and $f_0(g(n))\le g(n+1)$ for every
$n$, and $g$ is non-decreasing.

\Quer\ Suppose, if possible, that $f\in A$ is such that
$\{n:f(g(n))\le g(n+1)\}$ is finite.
Let $n_0\in\Bbb N$ be such that $f(g(n))\ge g(n+1)$ for every $n\ge n_0$.
If $i\ge n_0$ then

\Centerline{$\tilde f(g(i))\ge f(g(i))\ge g(i+1)$,}

\noindent so if $i\ge n_0$ and $j\in\Bbb N$ are such that $f^*(j)\ge g(i)$,
then

\Centerline{$f^*(j+1)
\ge\tilde f(\tilde f(f^*(j)))
\ge\tilde f(\tilde f(g(i)))
\ge\tilde f(g(i+1))
\ge g(i+2)$}

\noindent because $\tilde f$ is non-decreasing.   But $f^*$ is unbounded;
taking $k$ such that $f^*(k)\ge g(n_0)$, we have $f^*(k+i)\ge g(n_0+2i)$
for every $i\in\Bbb N$;  because both $f^*$ and $g$ are
non-decreasing, this means that $f^*(n)\ge g(n)$ whenever
$n\ge\max(k,2k-n_0)$.   But there must be such an $n$ with
$f^*(n)=h(n)<g(n)$, so this is impossible.\ \Bang

Thus $g$ has the required property.\ \Qed

\medskip

\quad{\bf (iv)} Now suppose that $P$ is a countable partially ordered
set, $\Cal Q$ is a family of cofinal subsets of $P$ with
$\#(\Cal Q)<\kappa$, and $p_0\in P$.   Let $\langle p_i\rangle_{i\ge 1}$
be such that $\sequence{i}{p_i}$ runs over $P$ with cofinal repetitions.
Let $f\in\NN$ be a  strictly increasing function such that whenever
$n\in\Bbb N$ and $i<n$ then there is a $j\in f(n)\setminus n$ such that
$p_i\le p_j$.   For each $Q\in\Cal Q$ let $f_Q\in\NN$ be a strictly
increasing function such that whenever $n\in\Bbb N$ and $i<n$ there is a
$j\in f_Q(n)\setminus n$ such that $p_i\le p_j\in Q$.
By (iii), we can find $g\in\NN$ such that $g(0)>0$,
$g(n+1)\ge f(g(n))$ for every $n$ and $I_Q=\{n:g(n+1)\ge f_Q(g(n))\}$ is
infinite for every $Q\in\Cal Q$.

For each $n\in\Bbb N$, set $J_n=g(n+1)\setminus g(n)$, and let $\Phi_n$
be the set of functions $h:g(n)\to J_n$ such that
$p_i\le p_{h(i)}$ for every $i<g(n)$;  because
$g(n+1)\ge f(g(n))$ this is non-empty.   For $Q\in\Cal Q$ and $n\in I_Q$
let $\phi_Q(n)\in\Phi_n$ be such that $p_{\phi_Q(n)(i)}\in Q$ for every
$i<g(n)$;  such a function exists because $g(n+1)\ge f_Q(g(n))$.   Now
all the $\Phi_n$ are countable (indeed finite), so (ii) tells us that
there is a $\phi\in\prod_{n\in\Bbb N}\Phi_n$ such that $\phi\cap\phi_Q$
is non-empty for every $Q\in\Cal Q$.

Define $\sequencen{i_n}$ by setting $i_0=0$ and $i_{n+1}=\phi(n)(i_n)$
for $n\in\Bbb N$;  because $\dom\phi(n)=g(n)>0$ and $\phi(n)(i)<g(n+1)$
whenever $i<g(n)$, $i_n$ is well-defined for each $n$.   Because
$\phi(n)\in\Phi_n$ for each $n$, $p_{i_n}\le p_{i_{n+1}}$ for each $n$.
If $Q\in\Cal Q$ there is some $n$ such that $\phi(n)=\phi_Q(n)$, so that

\Centerline{$p_{i_{n+1}}=p_{\phi(n)(i_n)}=p_{\phi_Q(n)(i_n)}\in Q$.}

\noindent But this means that $R=\{p_{i_k}:k\in\Bbb N\}$ is an
upwards-linked (indeed, totally ordered) subset of $P$ meeting every
member of $\Cal Q$ and containing $p_0$.   As $p_0$ and $\Cal Q$ are
arbitrary, $\frak m^{\uparrow}(P)\ge\kappa$.   As $P$ is arbitrary,
$\frakmctbl\ge\kappa$ and $\kappa=\frakmctbl=\cov\Cal M$, as claimed.

\medskip

{\bf (c)} Write $\lambda$ for the least cardinal of any set
$A\subseteq\NN$ such that for every $g\in\NN$
there is an $f\in A$ such that $f\cap g$ is infinite.
Of course $\lambda$ is infinite.

\medskip

\quad{\bf (i)} Let $A_0\subseteq\Bbb R$ be a non-meager set with cardinal
$\non\Cal M$.   Then $A_0\setminus\Bbb Q$ is not a meager subset of
$\Bbb R$ and therefore is not a meager subset of $\Bbb R\setminus\Bbb Q$,
since any subset of $\Bbb R\setminus\Bbb Q$ which is nowhere dense in
$\Bbb R\setminus\Bbb Q$ is also nowhere dense in $\Bbb R$.   So
$\Bbb R\setminus\Bbb Q$ has a non-meager subset with cardinal $\non\Cal M$.
But $\Bbb R\setminus\Bbb Q$ is homeomorphic to $\NN$ (4A2Ub again),
so there is a non-meager set $A\subseteq\NN$ with cardinal $\non\Cal M$.

Now take any $g\in\NN$.   Then
$\bigcup_{i\ge n}\{f:f\in\NN$, $f(i)=g(i)\}$ is a dense open subset of
$\NN$ for every $n\in\Bbb N$, so

\Centerline{$\{f:f\in\NN$, $f\cap g$ is infinite$\}
=\bigcap_{n\in\Bbb N}\bigcup_{i\ge n}\{f:f(i)=g(i)\}$}

\noindent is a dense G$_{\delta}$ set and meets $A$;  that is, there is an
$f\in A$ such that $f\cap g$ is infinite.   Accordingly
$\lambda\le\non\Cal M$.

\medskip

\quad{\bf (ii)} Let $A\subseteq\NN$ be a set with cardinal $\lambda$ such that
for every $g\in\NN$
there is an $f\in A$ such that $f\cap g$ is infinite.   Write $S_2$ for
$\bigcup_{n\in\Bbb N}\{0,1\}^n$.

\medskip

\qquad\grheada\ There is a set $A_1\subseteq\NN$, with cardinal at most
$\lambda$, such
that whenever $g\in\NN$ and $D\subseteq\Bbb N$ is infinite, there is
an $f\in A_1$ such that $f\cap g\restr D$ is infinite.
\Prf\ For $n\in\Bbb N$ set $F_n=\bigcup_{I\in[\Bbb N]^{n+1}}\Bbb N^I$.
Because each $F_n$ is countably infinite, there is a family
$H\subseteq\prod_{n\in\Bbb N}F_n$, with cardinal $\lambda$, such that for every
$\psi\in\prod_{n\in\Bbb N}F_n$ there is a $\phi\in H$ such that
$\phi\cap\psi$ is infinite.   For $\phi\in H$ define
$k_{\phi}\in\NN$ by setting
$k_{\phi}(n)=\min(\dom \phi(n)\setminus\{k_{\phi}(i):i<n\})$ for each $n$, so
that $k_{\phi}\in\NN$ is injective, and choose $f_{\phi}\in\NN$ such that
$f_{\phi}(k_{\phi}(n))=\phi(n)(k_{\phi}(n))$ for every $n$.
Now suppose that $g\in\NN$ and $D\subseteq\NN$ is infinite.
Define $\psi\in\prod_{n\in\Bbb N}F_n$ by taking $\psi(n)(i)=g(i)$
whenever $i\in D$ and $\#(D\cap i)\le n$.   We know that there is a
$\phi\in H$ such that $K=\{n:\phi(n)=\psi(n)\}$ is infinite, and that for
every $n\in K$

\Centerline{$k_{\phi}(n)\in\dom \phi(n)=\dom\psi(n)\subseteq D$,}

\Centerline{$f_{\phi}(k_{\phi}(n))=\phi(n)(k_{\phi}(n))=\psi(n)(k_{\phi}(n))
=g(k_{\phi}(n))$.}

\noindent As $k_{\phi}[K]$ is an infinite subset of $D$,
$f_{\phi}\cap g\restr D$ is infinite.
So we can take $A_1=\{f_{\phi}:\phi\in H\}$.\ \Qed

\medskip

\qquad\grheadb\ Because $S_2$ is countably infinite, we can copy $A$
onto a set $A_2\subseteq S_2^{\Bbb N}$, with cardinal
$\lambda$, such that for every
$\psi\in S_2^{\Bbb N}$ there is a
$\phi\in A_2$ such that $\phi\cap\psi$ is infinite.   For each
$\phi\in A_2$ let $h_{\phi}\in\NN$ be a strictly increasing function such
that $h_{\phi}(k)+\#(\phi(n))\le h_{\phi}(k+2)$ whenever
$k\in\Bbb N$ and $h_{\phi}(k)\le n<h_{\phi}(k+1)$.

\medskip

\qquad\grheadc\
For $\phi\in A_2$ and $f\in A_1$ define $x_{\phi f}$,
$y_{\phi f}\in\{0,1\}^{\Bbb N}$ by saying (in the notation of 5A1C) that

\Centerline{$x_{\phi f}
=\sigma_0^{\smallfrown}\sigma_1^{\smallfrown}\ldots$,
\quad$y_{\phi f}
=\tau_0^{\smallfrown}\tau_1^{\smallfrown}\ldots$,}

\noindent where

$$\eqalign{\sigma_k
&=\phi( f(k))\text{ if }h_{\phi}(2k)\le f(k)<h_{\phi}(2k+1),\cr
&=\fraction{0}\text{ otherwise},\cr
\tau_k
&=\phi( f(k))\text{ if }h_{\phi}(2k+1)\le f(k)<h_{\phi}(2k+2),\cr
&=\fraction{0}\text{ otherwise}.\cr}$$

\noindent Note that

\Centerline{$\#(\sigma_k)\le h_{\phi}(2k+2)-h_{\phi}(2k)$,
\quad$\#(\tau_k)\le h_{\phi}(2k+3)-h_{\phi}(2k+1)$}

\noindent for every $k$.   Write

\Centerline{$A_3=\{x_{\phi f}:\phi\in A_2$, $f\in A_1\}
\cup\{y_{\phi f}:\phi\in A_2$, $f\in A_1\}$}

\noindent so that $A_3\subseteq\{0,1\}^{\Bbb N}$ has cardinal at most
$\lambda$.
\medskip

\qquad\grheadd\ $A_3$ is a non-meager subset of $\{0,1\}^{\Bbb N}$.
\Prf\ Let $H$ be a dense G$_{\delta}$ subset of $\{0,1\}^{\Bbb N}$.
Then we can express $H$ as $\bigcap_{n\in\Bbb N}G_n$ where
$\sequencen{G_n}$ is a non-increasing sequence of dense open subsets of
$\{0,1\}^{\Bbb N}$.   For each $n$,

\Centerline{$G'_n=\bigcap_{m\le n}\bigcap_{\tau\in\{0,1\}^m}
\{x:x\in\{0,1\}^{\Bbb N}$, $\tau^{\smallfrown}x\in G_n\}$}

\noindent is a dense open subset of $\{0,1\}^{\Bbb N}$ so there is an
$\upsilon_n\in S_2$ such that
$\{x:\upsilon_n\subseteq x\in\{0,1\}^{\Bbb N}\}\subseteq G'_n$.
Let $\phi\in A_2$ be such that $C=\{n:n\in\Bbb N$, $\phi(n)=\upsilon_n\}$
is infinite.   Set
$D_0=\{k:C\cap h_{\phi}(2k+1)\setminus h_{\phi}(2k)\ne\emptyset\}$,
$D_1=\{k:C\cap h_{\phi}(2k+2)\setminus h_{\phi}(2k+1)\ne\emptyset\}$;
then at least one of $D_0$, $D_1$ is infinite.

Suppose that $D_0$ is infinite.   For $k\in D_0$ set
$n_k=\min(C\cap h_{\phi}(2k+1)\setminus h_{\phi}(2k))$.  Then there is
an $f\in A_1$ such that

\Centerline{$E_0=\{k:k\in D_0$, $f(k)=n_k\}$}

\noindent is infinite.   In this case,
for $k\in E_0$, $h_{\phi}(2k)\le f(k)<h_{\phi}(2k+1)$ so
$\#(\phi( f(k)))\le h_{\phi}(2k+2)-h_{\phi}(2k)$ and
$\sigma_k=\phi( f(k))=\phi(n_k)=\upsilon_{n_k}$.   Accordingly
$\sigma_k^{\smallfrown}\sigma_{k+1}^{\smallfrown}\ldots\in G'_{n_k}$.

At the same time, writing $\tau$ for
$\sigma_0^{\smallfrown}\ldots^{\smallfrown}\sigma_{k-1}$,

\Centerline{$\#(\tau)\le\sum_{i=0}^{k-1}h_{\phi}(2i+2)-h_{\phi}(2i)
\le h_{\phi}(2k)\le n_k$}

\noindent and

\Centerline{$x_{\phi f}
=\tau^{\smallfrown}\upsilon_{n_k}^{\mskip5mu\smallfrown}
\sigma_{k+1}^{\mskip18mu\smallfrown}
\ldots$}

\noindent belongs to $G_{n_k}$.   Since $\sequence{k}{n_k}$ is strictly
increasing, $x_{\phi f}$ belongs to $G_n$ for infinitely many $n$ and
therefore belongs to $H$.

Similarly, if $D_1$ is infinite, we can set
$m_k=\min(C\cap h_{\phi}(2k+2)\setminus h_{\phi}(2k+1))$ for $k\in D_1$,
take $f\in A_1$ such that
$E_1=\{k:k\in D_1$, $f(k)=m_k\}$ is infinite, and for $k\in E_1$ see
that $\#(\tau_0^{\smallfrown}\ldots^{\smallfrown}\tau_{k-1})\le m_k$
and $\tau_k=\upsilon_{m_k}$ so that $y_{\phi f}\in G_{m_k}$.
Thus in either case we have a member of $A_3$
belonging to $H$;  as $H$ is arbitrary, $A_3$ is non-meager.\ \Qed

\medskip

\qquad\grheade\ Finally, let $\theta:\{0,1\}^{\Bbb N}\to[0,1]$ be the
standard surjection defined by setting
$\theta(x)=\sum_{i=0}^{\infty}2^{-i-1}x(i)$ for $x\in\{0,1\}^{\Bbb N}$.
Then $\theta$ is continuous and irreducible,
so the inverse image of a dense open subset of
$[0,1]$ is a dense open subset of $\{0,1\}^{\Bbb N}$, and $\theta[A_3]$ is
non-meager in $[0,1]$.   As the interior of $[0,1]$ in $\Bbb R$ is dense in
$[0,1]$, $\theta[A_3]$ is non-meager in $\Bbb R$ and

\Centerline{$\non\Cal M\le\#(\theta[A_3])\le\#(A_3)\le\lambda$.}

\noindent Together with (i) above, this shows that $\non\Cal M=\lambda$, as
claimed.
}%end of proof of 522S

\leader{522T}{Martin numbers} Following the identification of
$\cov\Cal M$ with $\frakmctbl$, we can amalgamate the diagrams in 522B
and 517Ob, as follows:

$$\vbox{\offinterlineskip
\halign{\hfil#\hfil&\hfil#\hfil&\hfil#\hfil&\hfil#\hfil
  &\hfil#\hfil&\hfil#\hfil&\hfil#\hfil&\hfil#\hfil&\hfil#\hfil
  &\hfil#\hfil&\hfil#\hfil&\hfil#\hfil&\hfil#\hfil\cr
&\tmpstrut&&&$\cov\Cal N$&\tmphrule&$\non\Cal M$
  &\tmphrule&$\cf\Cal M$&\tmphrule&$\cf\Cal N$&\tmphrule&$\frak c$\cr
&\tmpstrut&&&\vrule&&\vrule&&\vrule&&\vrule\cr
&\tmpstrut&&&\vrule&&$\frak b$&\tmphrule&$\frak d$&&\vrule\cr
&\tmpstrut&&&\vrule&&\vrule&&\vrule&&\vrule\cr
&\tmpstrut&&&$\add\Cal N$&\tmphrule&$\add\Cal M$
  &\tmphrule&$\cov\Cal M$&\tmphrule&$\non\Cal N$\cr
&\tmpstrut&&&\vrule&&\vrule\cr
&\tmpstrut&&&$\frak m_{\sigma\text{-linked}}$
  &\tmphrule&$\frak p$\cr
&\tmpstrut&&&\vrule&&\vrule\cr
$\tmpstrut\omega_1$&\tmphrule&$\frak m$
  &\tmphrule&$\frak m_{\text{K}}$
  &\tmphrule&$\frak m_{\text{pc}\omega_1}$\cr
}}$$

\proof{ The two new inequalities to be proved are
$\frak m_{\sigma\text{-linked}}\le\add\Cal N$ and
$\frak p\le\add\Cal M$.

\medskip

{\bf (a)} Let $\Cal S^{\infty}$ be the `($\Bbb N,\infty)$-localization poset'

\Centerline{$\{p:p\subseteq\Bbb N\times\Bbb N$, $\#(p[\{n\}])\le 2^n$
for every $n$, $\sup_{n\in\Bbb N}\#(p[\{n\}])<\infty\}$,}

\noindent ordered by $\subseteq$.   For $p\in\Cal S^{\infty}$ set
$\|p\|=\max_{n\in\Bbb N}\#(p[\{n\}])$.   Then $\Cal S^{\infty}$ is
$\sigma$-linked upwards.   \Prf\ If $p$, $q\in\Cal S^{\infty}$,
$\|p\|\le n$, $\|q\|\le n$ and $p[\{i\}]=q[\{i\}]$ for every $i\le n$,
then $p\cup q\in\Cal S^{\infty}$.   So for any $n\in\Bbb N$ and
$\langle J_i\rangle_{i\le n}\in\prod_{i\le n}[\Bbb N]^{\le 2^i}$ we have an
upwards-linked set

\Centerline{$\{p:p\in\Cal S^{\infty}$, $\|p\|\le n$,
$p[\{i\}]=J_i$ for every $i\le n\}$;}

\noindent as there are only countably many such families
$\langle J_i\rangle_{i\le n}$,
$\Cal S^{\infty}$ is $\sigma$-linked upwards.\ \Qed

Accordingly $\frak m_{\sigma\text{-linked}}
\le\frak m^{\uparrow}(\Cal S^{\infty})$.   Next,
$\frak m^{\uparrow}(\Cal S^{\infty})\le\add(\NN,\subseteq^*,\Cal S)$, where
$(\NN,\subseteq^*,\Cal S)$ is the $\Bbb N$-localization relation.
\Prf\ Suppose that
$A\subseteq\NN$ and $\#(A)<\frak m^{\uparrow}(\Cal S^{\infty})$.   For
each $f\in A$, set $Q_f=\{p:p\in\Cal S^{\infty}$, $f\subseteq^*p\}$.
If $p\in\Cal S^{\infty}$ and $\|p\|=n$, then
$p\subseteq p\cup\{(i,f(i)):i\ge n\}\in Q_f$;  so $Q_f$ is cofinal with
$\Cal S^{\infty}$.   As
$\#(\{Q_f:f\in A\})<\frak m^{\uparrow}(\Cal S^{\infty})$, there is an
upwards-directed $R\subseteq\Cal S^{\infty}$ meeting $Q_f$ for every
$f\in A$.   Set $S=\bigcup R$.   For each $n\in\Bbb N$,
$\{p[\{n\}]:p\in R\}$ is an upwards-directed family of subsets of
$\Bbb N$, all of{} size at most $2^n$, with union $S[\{n\}]$.   So
$\#(S[\{n\}])\le 2^n$;  as $n$ is arbitrary, $S\in\Cal S$.   If
$f\in A$, there is a
$p\in R\cap Q_f$, and now $f\subseteq^*p\subseteq S$.
As $A$ is arbitrary, we have the result.\ \Qed

Now

$$\eqalignno{\frak m_{\sigma\text{-linked}}
&\le\frak m^{\uparrow}(\Cal S^{\infty})
\le\add(\NN,\subseteq^*,\Cal S)
=\add(\Cal N,\subseteq,\Cal N)\cr
\displaycause{522M, 512Db}
&=\add\Cal N,\cr}$$

\noindent as required.

\medskip

{\bf (b)(i)} Let $\Cal U$ be a countable base for the topology of
$\Bbb R$, not containing $\emptyset$.   Consider the set $P$ of pairs
$(\sigma,F)$ where $\sigma\in\bigcup_{n\in\Bbb N}\Cal U^n$ and
$F\subseteq\Bbb R$ is nowhere dense, together with the relation $\le$
where $(\sigma,F)\le(\sigma',F')$ if $\sigma'$ extends $\sigma$,
$F'\supseteq F$ and $F\cap\sigma'(i)=\emptyset$ whenever
$i\in\dom\sigma'\setminus\dom\sigma$.   Then $\le$ is a partial order on
$P$.   \Prf\ If $(\sigma,F)\le(\sigma',F')\le(\sigma'',F'')$ then we
surely have $\sigma\subseteq\sigma'\subseteq\sigma''$ and
$F\subseteq F'\subseteq F''$.   If
$i\in\dom\sigma''\setminus\dom\sigma$, then either
$i\in\dom\sigma'\setminus\dom\sigma$ and $\sigma''(i)=\sigma'(i)$ must
be disjoint from $F$, or $i\in\dom\sigma''\setminus\dom\sigma'$ and
$\sigma''(i)$ must be disjoint from $F'\supseteq F$.   Thus in either
case $F\cap\sigma''(i)=\emptyset$;  as $i$ is arbitrary,
$(\sigma,F)\le(\sigma'',F'')$.   Thus $\le$ is transitive.   Evidently
it is also reflexive and anti-symmetric, so it is a partial order.\ \Qed

\medskip

\quad{\bf (ii)} $(P,\le)$ is $\sigma$-centered upwards.   \Prf\ If
$(\sigma,F_0),\ldots,(\sigma,F_k)$ are members of $P$ with a common
first member, then they have a common upper bound
$(\sigma,\bigcup_{i\le k}F_i)$ in $P$.   So for any $n\in\Bbb N$ and
$\sigma\in\Cal U^n$ the set
$\{(\sigma,F):F\subseteq\Bbb R$ is nowhere dense$\}$ is upwards-centered
in $P$;  as $\bigcup_{n\in\Bbb N}\Cal U^n$ is countable, $P$ is
$\sigma$-centered upwards.\ \Qed

\medskip

\quad{\bf (iii)} For each $V\in\Cal U$ and $n\in\Bbb N$ set

\Centerline{$Q_{nV}=\{(\sigma,F):(\sigma,F)\in P$,
$V\cap\bigcup_{n\le i<\dom\sigma}\sigma(i)\ne\emptyset\}$.}

\noindent Then $Q_{nV}$ is cofinal with $P$.   \Prf\ If
$(\sigma,F)\in P$, set $m=\max(n,\dom\sigma)+1$, and take $U\in\Cal U$
such that $U\subseteq V\setminus F$.   Setting

$$\eqalignno{\sigma'(i)
&=\sigma(i)\text{ for }i<\dom\sigma,\cr
&=U\text{ for }\dom\sigma\le i<m,\cr}$$

\noindent we find that $(\sigma,F)\le(\sigma',F)\in Q_{nV}$.\ \Qed

For each nowhere dense set $H\subseteq\Bbb R$,

\Centerline{$Q'_H=\{(\sigma,F):(\sigma,F)\in P$, $H\subseteq F\}$}

\noindent is cofinal with $P$.   \Prf\ For any $(\sigma,F)\in P$, we
have $(\sigma,F)\le(\sigma,F\cup H)\in Q'_H$.\ \Qed

\medskip

\quad{\bf (iv)} Now suppose that $\Cal A\subseteq\Cal M$ and
$\#(\Cal A)<\frak p$.   Then each member of $\Cal A$ is covered by a
sequence of nowhere dense sets, so there is a family $\Cal H$ of nowhere
dense sets with the same union as $\Cal A$ and with
$\#(\Cal H)\le\max(\omega,\#(\Cal A))$.   In this case

\Centerline{$\Cal Q
=\{Q_{nV}:n\in\Bbb N$, $V\in\Cal U\}\cup\{Q'_H:H\in\Cal H\}$}

\noindent is a family of cofinal subsets of $P$ and

\Centerline{$\#(\Cal Q)\le\max(\omega,\#(\Cal A))<\frak p
\le\frak m^{\uparrow}(P)$.}

\noindent There is therefore an upwards-directed $R\subseteq P$ meeting
every member of $\Cal Q$.   If $(\sigma,F)$ and $(\sigma',F')$ belong to
$R$, they must be upwards-compatible in $P$, and in particular $\sigma$
and $\sigma'$ have a common extension;  we therefore have a function
$\phi=\bigcup_{(\sigma,F)\in R}\sigma$ from a subset of $\Bbb N$ to
$\Cal U$.   If $n\in\Bbb N$ and $V\in\Cal U$, then there is a
$(\sigma,F)\in R\cap Q_{nV}$, so that there is some $i\ge n$ such that
$\phi(i)=\sigma(i)$ meets $V$.   As $V$ is arbitrary, the open set
$W_n=\bigcup_{i\in\dom\phi,i\ge n}\phi(i)$ is dense;  as $n$ is arbitrary,
$M=\Bbb R\setminus\bigcap_{n\in\Bbb N}W_n$ is meager.   Now
$H\subseteq M$ for every $H\in\Cal H$.   \Prf\ There is a
$(\sigma,F)\in R\cap Q'_H$.   Set $n=\dom\sigma$.   If
$i\in\dom\phi\setminus n$, there is a $(\sigma',F')\in R$ such that
$i\in \dom\sigma'$;  because $R$ is upwards-directed, we may suppose
that $(\sigma,F)\le(\sigma',F')$.   But in this case
$\phi(i)=\sigma'(i)$ must be disjoint from $F$ and therefore from $H$.
As $i$ is arbitrary, $H\cap W_n=\emptyset$ and $H\subseteq M$.\ \Qed

As $H$ is arbitrary, $\bigcup\Cal A=\bigcup\Cal H\subseteq M$.   As
$\Cal A$ is arbitrary, $\add\Cal M\ge\frak p$, as claimed.
}%end of proof of 522T

\cmmnt{\medskip

\noindent{\bf Remark} In fact $\frak m^{\uparrow}(\Cal S^{\infty})$ is
exactly equal to $\add\Cal N$;  see 528N.
}%end of comment

\leader{*522U}{\dvrocolon{FN($\Cal P\Bbb N$)}}\cmmnt{ For any cardinal
which is known to lie between $\omega_1$ and $\frak c$, it is natural,
and often profitable, to try to locate it on Cicho\'n's diagram.   For
the Freese-Nation number of $\Cal P\Bbb N$, which appeared more than once
in \S518, we have the following results.

\medskip

\noindent}{\bf Proposition}\cmmnt{ ({\smc Fuchino Koppelberg \&
Shelah 96}, {\smc Fuchino Geschke \& Soukup 01})}
(a) $\FN(\Cal P\Bbb N)\ge\frak b$.

(b) $\FN(\Cal P\Bbb N)\ge\cov\Cal N$.

(c) If $\FN(\Cal P\Bbb N)=\omega_1$ then $\shr\Cal M=\omega_1$, so

\Centerline{$\frak m=\frak m_{\text{K}}=\frak m_{\text{pc}\omega_1}
=\frak m_{\sigma\text{-linked}}=\frak p=\add\Cal N=\add\Cal M
=\frak b=\cov\Cal N=\non\Cal M=\omega_1$.}

(d) If $\FN(\Cal P\Bbb N)=\omega_1$ and $\kappa\ge\frakmctbl$ is such
that $\cff[\kappa]^{\le\omega}\le\kappa\le\frak c$, then
$\kappa=\frak c$.
So if $\FN(\Cal P\Bbb N)=\omega_1$ and $\frakmctbl<\omega_{\omega}$,
then

\Centerline{$\frakmctbl=\non\Cal N=\frak d=\cf\Cal M=\cf\Cal N
=\frak c$.}

(e) There is a set $A\subseteq\Bbb R$ with cardinal $\frakmctbl$ such that
every meager set meets $A$ in a set with cardinal less than
$\FN^*(\Cal P\Bbb N)$.

\proof{{\bf (a)} Let $\le^*$ and $\preceq$
be the pre-order and partial order on $\NN$ described in
522C, so that $\frak b=\add(\NN,\preceq)$.   Write $\kappa$ for
$\FN(\Cal P\Bbb N)$;  by
518D, $\kappa=\FN(\NN,\le)$ and we have a Freese-Nation function
$\phi:\NN\to[\NN]^{<\kappa}$ for
$\le$.   For $f\in\NN$, set $\psi(f)=\bigcup\{\phi(g):g\le^*f\le^*g\}$;
then $\#(\psi(f))\le\kappa$.

\Quer\ Suppose, if possible, that $\kappa<\frak b$.   Choose
a family $\langle f_{\xi}\rangle_{\xi\le\kappa}$ in $\NN$
inductively, as follows.   Given
$\ofamily{\eta}{\xi}{f_{\eta}}$ where $\xi\le\kappa$,
$\bigcup_{\eta<\xi}\psi(f_{\eta})$ has cardinal at most $\kappa<\frak b$,
so has a $\preceq$-upper bound $f'_{\xi}$;  now set
$f_{\xi}(i)=f'_{\xi}(i)+1$ for every $i$, and continue.

Next choose $\ofamily{\xi}{\kappa}{h_{\xi}}$ in $\phi(f_{\kappa})$
as follows.   For each $\xi<\kappa$,
$f_{\xi}\in\phi(f_{\xi})$ (511Hh) so $f_{\xi}\in\psi(f_{\xi})$,
$f_{\xi}\preceq f'_{\kappa}$ and $f_{\xi}\preceq f_{\kappa}$.   So if we set
$g_{\xi}=f_{\xi}\wedge f_{\kappa}$ then
$g_{\xi}\le^*f_{\xi}\le^*g_{\xi}$ while
$g_{\xi}\le f_{\kappa}$.   There is therefore an
$h_{\xi}\in\phi(g_{\xi})\cap\phi(f_{\kappa})$ such that
$g_{\xi}\le h_{\xi}\le f_{\kappa}$.   Now if $\eta<\xi<\kappa$,
$h_{\eta}\in\phi(g_{\eta})\subseteq\psi(f_{\eta})$ so
$h_{\eta}\preceq f'_{\xi}$.   Accordingly

$$\eqalign{\{i:h_{\xi}(i)\le h_{\eta}(i)\}&
\subseteq\{i:h_{\xi}(i)\le f'_{\xi}(i)\}\cup\{i:f'_{\xi}(i)<h_{\eta}(i)\}\cr&
\subseteq\{i:h_{\xi}(i)<f_{\xi}(i)\}\cup\{i:f'_{\xi}(i)<h_{\eta}(i)\}\cr&
\subseteq\{i:g_{\xi}(i)<f_{\xi}(i)\}\cup\{i:f'_{\xi}(i)<h_{\eta}(i)\}\cr}$$

\noindent is finite and $h_{\xi}\ne h_{\eta}$.   But this means that
$\{h_{\xi}:\xi<\kappa\}$ has cardinal $\kappa$ and
$\#(\phi(f_{\kappa}))=\kappa$, contrary to the choice of $\phi$.\ \Bang

Thus $\frak b\le\kappa=\FN(\Cal P\Bbb N)$, as claimed.
In particular, $\FN(\Cal P\Bbb N)$ is uncountable.

\medskip

{\bf (b)(i)} We need to know the following fact:  if $\Cal E$ is a
family of non-negligible Lebesgue measurable subsets of $\Bbb R$, and
$\#(\Cal E)<\cov\Cal N$, there is a countable set meeting every member of
$\Cal E$.   \Prf\ For each $E\in\Cal E$, $\Bbb R\setminus(\Bbb Q+E)$ is
negligible (439Eb), so there is an
$x\in\Bbb R\cap\bigcap_{E\in\Cal E}\Bbb Q+E$;  now $\Bbb Q+x$ is countable
and meets every member of $\Cal E$.\ \Qed

\medskip

\quad{\bf (ii)} Set $\kappa=\FN(\Cal P\Bbb N)$.   If $\Cal C$ is the
family of closed sets in $\Bbb R$, then
$(\Cal C,\subseteq)\cong(\frak T,\supseteq)$, so
$\FN(\Cal C)=\FN(\frak T)=\kappa$ (518D).
Let $f:\Cal C\to[\Cal C]^{<\kappa}$ be a
Freese-Nation function.

\medskip

\quad{\bf (iii)} \Quer\ If $\kappa<\cov\Cal N$, write $K$ for the set of
infinite successor cardinals $\lambda\le\kappa$, and for
$\lambda\in K$ set
$D_{\lambda}=\{x:x\in\Bbb R$, $\#(f(\{x\}))<\lambda\}$.   As
$\Bbb R=\bigcup_{\lambda\in K}D_{\lambda}$, there must be
some $\lambda\in K$ such that $D_{\lambda}$ cannot be covered
by $\kappa$ negligible sets.   Choose
$\langle M_{\xi}\rangle_{\xi\le\lambda}$ and
$\langle H_{\xi n}\rangle_{\xi<\lambda,n\in\Bbb N}$ inductively, as
follows.   $M_{\xi}=\emptyset$.   Given that $M_{\xi}\subseteq\Cal C$
and $\#(M_{\xi})\le\kappa$, (i)
tells us that there is a countable set $A_{\xi}\subseteq\Bbb R$
meeting every non-negligible member of $M_{\xi}$;  let
$\sequencen{H_{\xi n}}$ be a sequence of closed subsets of
$\Bbb R\setminus A_{\xi}$ such that $\bigcup_{n\in\Bbb N}H_{\xi n}$ is
conegligible.   Now set

\Centerline{$M_{\xi+1}
=M_{\xi}\cup\{H_{\xi n}:n\in\Bbb N\}\cup\bigcup_{F\in M_{\xi}}f(F)
\in[\Cal C]^{\le\kappa}$.}

\noindent At non-zero limit ordinals $\xi\le\lambda$, set
$M_{\xi}=\bigcup_{\eta<\xi}M_{\eta}$.

By the choice of $\lambda$, there is an $x\in D_{\lambda}$ which does not
belong to any negligible set belonging to $M_{\lambda}$,
nor to any of the sets
$\Bbb R\setminus\bigcup_{n\in\Bbb N}H_{\xi n}$ for $\xi<\lambda$.   Now
$\#(M_{\lambda}\cap f(\{x\}))<\lambda$;  because $\lambda$ is regular,
there is a $\xi<\lambda$ such that
$M_{\lambda}\cap f(\{x\})\subseteq M_{\xi}$.
Let $n\in\Bbb N$ be such that $x\in H_{\xi n}$.   Then there must be an
$F\in f(\{x\})\cap f(H_{\xi n})$ such that $x\in F\subseteq H_{\xi n}$.
In this case, $H_{\xi n}\in M_{\xi+1}$ and
$F\in M_{\xi+2}\subseteq M_{\lambda}$, so in fact
$F\in M_{\xi}$.   Because $x\in F$, $F$ cannot be negligible, so
$A_{\xi}\cap F\ne\emptyset$;  but $H_{\xi n}$ was chosen to be disjoint
from $A_{\xi}$.\ \Bang

\medskip

\quad{\bf (iv)} Thus $\kappa\ge\cov\Cal N$, as claimed.
\medskip

{\bf (c)} Let $A\subseteq\Bbb R$ be a non-meager set.

\medskip

\quad{\bf (i)} By 518D, $\FN(\frak T)=\omega_1$, where $\frak T$ is the
topology of $\Bbb R$.   Let $f:\frak T\to[\frak T]^{\le\omega}$ be a
Freese-Nation function.   There is a set $M$ such that

\inset{($\alpha$) whenever $G\in M\cap\frak T$ then $f(G)\subseteq M$;

($\beta$) whenever $t\in M\cap\Bbb R$ then $\Bbb R\setminus\{t\}\in M$;

($\gamma$) whenever $\Cal G\subseteq M$ is a countable family of dense
open subsets of $\Bbb R$, $M\cap A\cap\bigcap\Cal G$ is non-empty;

($\delta$) $\#(M)\le\omega_1$.}

\noindent\Prf\ Build a non-decreasing family
$\ofamily{\xi}{\omega_1}{M_{\xi}}$ of countable sets as follows.
$M_0=\emptyset$.   Given that $M_{\xi}$ is countable, let $M_{\xi+1}$
be a countable set including $M_{\xi}$ such that

\inset{($\alpha$) whenever $G\in M_{\xi}\cap\frak T$ then
$f(G)\subseteq M_{\xi+1}$;

($\beta$) whenever $t\in M_{\xi}\cap\Bbb R$ then
$\Bbb R\setminus\{t\}\in M_{\xi+1}$;

($\gamma$) $M_{\xi+1}\cap A\cap\bigcap\{G:G\in M_{\xi}$ is a dense open
subset of $\Bbb R\}$ is not empty.}

\noindent For countable limit ordinals $\xi>0$, set
$M_{\xi}=\bigcup_{\eta<\xi}M_{\eta}$.   At the end of the construction,
set $M=\bigcup_{\xi<\omega_1}M_{\xi}$.\ \Qed

\medskip

\quad{\bf (ii)} If $H\subseteq\Bbb R$ is an open set, there is a
countable family $\Cal G\subseteq M\cap\frak T$ such that
$M\cap\Bbb R\cap\bigcap\Cal G\subseteq H\subseteq\bigcap\Cal G$.   \Prf\
Set $\Cal G=\{G:G\in f(H)\cap M$, $H\subseteq G\}$;  then certainly
$H\subseteq\bigcap\Cal G$ and $\Cal G$ is countable.   If
$t\in M\cap\Bbb R\setminus H$, then $H\subseteq\Bbb R\setminus\{t\}$ so
there is a $G\in f(H)\cap f(\Bbb R\setminus\{t\})$ such that
$H\subseteq G\subseteq\Bbb R\setminus\{t\}$;  since
$\Bbb R\setminus\{t\}\in M$,
$G\in M$;  and $t\notin G$.   As $t$ is arbitrary,
$M\cap\Bbb R\cap\bigcap\Cal G\subseteq H$.\ \Qed

\medskip

\quad{\bf (iii)} Now consider $B=A\cap M$.   Then $\#(B)\le\omega_1$.
\Quer\ If $B$ is meager, let $\sequencen{H_n}$ be a sequence of dense
open sets such that $B\cap\bigcap_{n\in\Bbb N}H_n=\emptyset$.   For each
$n\in\Bbb N$, let $\Cal G_n$ be a countable family of dense open sets
belonging to $M$ such that
$M\cap\Bbb R\cap\bigcap\Cal G_n\subseteq H_n$ (using (ii)).   Set
$\Cal G=\bigcup_{n\in\Bbb N}\Cal G_n$;  then $\Cal G\subseteq M$ is a
countable family of dense open sets, so there is a
$t\in M\cap A\cap\bigcap\Cal G$, by condition ($\gamma$) in the
specification of $M$.   But now
$t\in M\cap A\cap\bigcap\Cal G_n\subseteq H_n$ for each $n$, so
$t\in B\cap\bigcap_{n\in\Bbb N}H_n$, which is impossible.\ \Bang

Thus $A$ has a non-meager subset with cardinal at most $\omega_1$;  as $A$ is
arbitrary, $\shr\Cal M=\omega_1$.

\medskip

{\bf (d)(i)} Again let $\frak T$ be the topology of $\Bbb R$ and
$f:\frak T\to[\frak T]^{\le\omega}$ a
Freese-Nation function.   This time, we can find a set $M$ such that

\inset{($\dagger$) for every $g\in\NN$ there is an $h\in M\cap\NN$ such
that $g(n)\ne h(n)$ for every $n\in\Bbb N$;

($\alpha$) whenever $G\in M\cap\frak T$ then
$f(G)\subseteq M$;

($\beta$) $M\cap[M]^{\le\omega}$ is cofinal with $[M]^{\le\omega}$;

($\gamma$) whenever $D\in M$ is countable, then there is a double
sequence $\langle G_{ij}\rangle_{i,j\in\Bbb N}$ belonging to
$M$ such that every $G_{ij}$ belongs to $\frak T$,
$\sequence{j}{G_{ij}}$ is disjoint for each $i\in\Bbb N$ and whenever
$G\in D$ is an open subset of $\Bbb R$ with infinite complement, there
is an $i\in\Bbb N$ such that $G_{ij}\setminus G$ is non-empty for every
$j\in\Bbb N$;

($\delta$) whenever $\langle G_{ij}\rangle_{i,j\in\Bbb N}\in M$ is a
double sequence of open subsets of $\Bbb R$, and $h\in M\cap\NN$, then
$\bigcup_{i\in\Bbb N}G_{i,h(i)}\in M$;

($\epsilon$) $\#(M)=\kappa$.}

\noindent\Prf\ Build a non-decreasing family
$\ofamily{\xi}{\omega_1}{M_{\xi}}$ of sets with cardinal $\kappa$ as
follows.   Start with a set $M_0\subseteq\NN$ such that $\#(M_0)=\kappa$
and for every $g\in\NN$ there is an $h\in M_0$ such that $g(n)\ne h(n)$
for every $n\in\Bbb N$ (using 522Sb).
Given that $\#(M_{\xi})=\kappa$, then $\cf[M_{\xi}]^{\le\omega}=\kappa$.
Let $M_{\xi+1}\supseteq M_{\xi}$ be such that

\inset{($\alpha$) whenever $G\in M_{\xi}\cap\frak T$ then
$f(G)\subseteq M_{\xi+1}$;

($\beta$) $M_{\xi+1}\cap[M_{\xi}]^{\le\omega}$ is cofinal with
$[M_{\xi}]^{\le\omega}$;

($\gamma$) whenever $D\in M_{\xi}$ is countable, then there is a double
sequence $\langle G_{ij}\rangle_{i,j\in\Bbb N}\in M_{\xi+1}$
such that every $G_{ij}$ is an open set in $\Bbb R$,
$\sequence{j}{G_{ij}}$ is disjoint for each $i\in\Bbb N$ and whenever
$G\in D$ is an open subset of $\Bbb R$ with infinite complement, there
is an $i\in\Bbb N$ such that $G_{ij}\setminus G$ is non-empty for every
$j\in\Bbb N$;

($\delta$) whenever $\langle G_{ij}\rangle_{i,j\in\Bbb N}\in M_{\xi}$ is
a double sequence of open subsets of $\Bbb R$, and
$h\in M_{\xi}\cap\NN$, then
$\bigcup_{i\in\Bbb N}G_{i,h(i)}\in M_{\xi+1}$;

($\epsilon$) $\#(M_{\xi+1})=\kappa$.}

\noindent For limit ordinals $\xi>0$, set
$M_{\xi}=\bigcup_{\eta<\xi}M_{\eta}$.   At the end of the construction,
set $M=\bigcup_{\xi<\omega_1}M_{\xi}$.   Then

\Centerline{$M\cap[M]^{\le\omega}
=\bigcup_{\xi<\omega_1}M_{\xi+1}\cap[M_{\xi}]^{\le\omega}$}

\noindent is cofinal with
$\bigcup_{\xi<\omega_1}[M_{\xi}]^{\le\omega}=[M]^{\le\omega}$, and it is
easy to see that the other conditions are satisfied.\ \Qed

\medskip

\quad{\bf (ii)} \Quer\ Now suppose, if possible, that there is a
$t\in\Bbb R$ such that $\Bbb R\setminus I\notin M$ for any finite set
$I$ containing $t$.   Set

\Centerline{$\Cal G=\{G:G\in M\cap f(\Bbb R\setminus\{t\})$, $t\notin G\}$.}

\noindent Then $\Cal G$ is a countable subset of
$M$ and $\Bbb R\setminus G$ is infinite for every $G\in\Cal G$.   Let
$D\in M$ be a countable set including $\Cal G$.   Then we have a double
sequence $\langle G_{ij}\rangle_{i,j\in\Bbb N}\in M$ such that
$\sequence{j}{G_{ij}}$ is a disjoint sequence of open sets for each $i\in\Bbb N$ and whenever
$G\in D$ is an open subset of $\Bbb R$ with infinite complement, there
is an $i\in\Bbb N$ such that $G_{ij}\setminus G$ is non-empty for every
$j\in\Bbb N$.   In particular, this last clause is true for every
$G\in\Cal G$.   For each $i\in\Bbb N$ choose $g(i)\in\Bbb N$ such that
$t\notin G_{ij}$ for any $j\ne g(i)$;   let $h\in M\cap\NN$ be such that
$h(i)\ne g(i)$ for every $i$, and set
$H=\bigcup_{i\in\Bbb N}G_{i,h(i)}\in M$;  note that $t\notin H$.   Now
there is a
$G\in f(H)\cap f(\Bbb R\setminus\{t\})$ such that
$H\subseteq G$ and $t\notin G$.   As $f(H)\subseteq M$,
$G\in M$, so $G\in\Cal G$.   But this means that $G_{i,h(i)}\subseteq G$ for every
$i\in\Bbb N$;  and we chose $\langle G_{ij}\rangle_{i,j\in\Bbb N}$ so
that this could not be so.\ \Bang

Thus $\Cal I=\{I:I\in[\Bbb R]^{<\omega}$, $\Bbb R\setminus I\in M\}$
covers $\Bbb R$.   As $\#(\Cal I)\le\#(M)\le\kappa$,
$\#(\Bbb R)\le\kappa$ and $\kappa=\frak c$, as claimed.

\medskip

\quad{\bf (iii)} Finally, if $\frakmctbl<\omega_{\omega}$, then we can
take $\kappa=\frakmctbl$, by 5A1F(e-iv), and get $\frakmctbl=\ldots=\frak c$.

\medskip

{\bf (e)} Because $\FN(\frak T)=\FN(\Cal P\Bbb N)$, 518E tells us that
there is a set $A\subseteq\Bbb R$, with cardinal $n(\Bbb R)=\frakmctbl$,
such that
$\#(A\cap F)<\FN^*(\frak T)=\FN^*(\Cal P\Bbb N)$ for every nowhere dense
set $F\subseteq\Bbb R$.   As $\FN^*(\Cal P\Bbb N)$ certainly has
uncountable cofinality, $A$ meets every meager set in a set with cardinal less
than $\FN^*(\Cal P\Bbb N)$.
}%end of proof of 522U
%elementary submodel arguments

%is FN(PN)\ge\non\Cal M?

\leader{522V}{\dvrocolon{Cofinalities}}\cmmnt{ For any cardinal
associated
with a mathematical structure, we can ask whether there are any
limitations on what that cardinal can be.   The commonest form of such
limitations, when they appear, is a restriction on the possible
cofinalities of the cardinal.   I run through the known results
concerning the cardinals of Cicho\'n's diagram.   Most are elementary,
but part (f) requires a substantial argument.

\medskip

\noindent}{\bf Proposition} (a) $\cf\frak c\ge\frak p$.

(b) $\add\Cal N$, $\add\Cal M$ and $\frak b$ are regular.

(c) $\cf(\cf\Cal N)\ge\add\Cal N$,
$\cf(\cf\Cal M)\ge\add\Cal M$ and $\cf\frak d\ge\frak b$.

(d) $\cf(\non\Cal N)\ge\add\Cal N$, $\cf(\non\Cal M)\ge\add\Cal M$.

(e) If $\cf\Cal M=\frakmctbl$ then $\cf(\cf\Cal M)\ge\non\Cal M$;  if
$\cf\Cal N=\cov\Cal N$, then $\cf(\cf\Cal N)\ge\non\Cal N$.

(f)\cmmnt{ ({\smc Bartoszy\'nski \& Judah 89})}
$\cf(\frakmctbl)\ge\add\Cal N$.

\proof{{\bf (a)} If $\omega\le\kappa<\frak p$ then $2^{\kappa}=\frak c$,
by 517Rb, so $\cf\frak c>\kappa$ by 5A1Fd.

\medskip

{\bf (b)} Use 513C(a-i);  to see that $\frak b$ is regular, use
its characterization as the additivity of a partially ordered set in
522C(ii).

\medskip

{\bf (c)} Use 513C(a-ii);  this time, we need to know that $\frak d$ is
the cofinality of a partially ordered set for which $\frak b$ is the
additivity.

\medskip

{\bf (d)-(e)} 513Cb with 522Sa,

\medskip

{\bf (f)(i)} Write $\Cal M_1$ for the ideal of meager subsets of
$\NN$, where $\NN$ is given its usual topology.
Let $(\NN,\subseteq^*,\Cal S)$ be the $\Bbb N$-localization relation (522K),
and set
$\Cal S^{(0)}=\{S:S\in\Cal S$, $\lim_{n\to\infty}2^{-n}\#(S[\{n\}])=0\}$.
I will write \finint\ and \disj\ for the relations
$\{(A,B):A\cap B$ is finite$\}$, $\{(A,B):A\cap B=\emptyset\}$.
Following the same mild abuse of notation as in 512Aa and elsewhere, I
will write
$(\Cal S^{(0)},\finint,\NN)$ and $(\NN,\disj,\NN)$ for the supported
relations $(\Cal S^{(0)},R_1,\NN)$ and $(\NN,R_2,\NN)$, where

\Centerline{$R_1
=\{(S,f):S\in\Cal S^{(0)}$, $f\in\NN$, $\{n:(n,f(n))\in S\}$ is finite$\}$,}

\Centerline{$R_2
=\{(f,g):f,\,g\in\NN$, $f(n)\ne g(n)$ for every $n\}$.}

\medskip

\quad{\bf (ii)($\alpha$)}
$(\NN,\in,\Cal M_1)\prGT(\Cal S^{(0)},\finint,\NN)$.
\Prf\ For $f\in\NN$, set $\phi(f)=f$ (identifying $f$ with its graph, as
usual);  for $g\in\NN$, set
$\psi(g)=\{h:h\in\NN,\,h\cap g$ is finite$\}$.   Then
$\phi(f)\in\Cal S^{(0)}$ for every $f\in\NN$, and $\psi(g)\in\Cal M_1$ for
every $g\in\NN$, because all the sets $\{h:h\cap g\subseteq n\}$ are
nowhere dense.   If $f$, $g\in\NN$ and $(\phi(f),g)\in\finint$, then
$f\cap g$ is finite so $f\in\psi(g)$;  thus $(\phi,\psi)$ is a
Galois-Tukey connection from
$(\NN,\in,\Cal M_1)$ to $(\Cal S^{(0)},\finint,\NN)$.\ \Qed

\medskip

\qquad\grheadb\ $(\Cal S^{(0)},\finint,\NN)\prGT(\NN,\disj,\NN)$.   \Prf\
Let $\sequencen{I_n}$ be a partition of $\Bbb N$ such that $\#(I_n)=2^n$
for each $n$.   For $n\in\Bbb N$, let $\theta_n:\BbbN^{I_n}\to\Bbb N$ be
a bijection.   For $S\in\Cal S^{(0)}$, choose $\phi(S)\in\NN$ such that
whenever $(n,i)\in S$ then $\phi(S)\cap\theta_n^{-1}(i)\ne\emptyset$,
where once again both the function $\phi(S)$ and the function
$\theta_n^{-1}(i)$ are identified with their graphs;  this is possible
because on each set $I_n$ there are at most $2^n$ functions with domain
$I_n$ that $\phi(S)$ has to meet.   For $g\in\NN$, define $\psi(g)\in\NN$ by
saying that $\psi(g)(n)=\theta_n(g\restr I_n)$ for every $n$.

Now suppose that $S\in\Cal S^{(0)}$ and $g\in\NN$ are such that
$S\cap\psi(g)$ is infinite.   Then there is certainly an $n$ such that
$(n,\psi(g)(n))\in S$.   In this case,

\Centerline{$\emptyset\ne\phi(S)\cap\theta_n^{-1}(\psi(g)(n))
=\phi(S)\cap g\restr I_n$,}

\noindent so $\phi(S)\cap g$ is non-empty.   Turning this round, if
$(\phi(S),g)\in\disj$ then $(S,\psi(g))\in\finint$;  that is,
$(\phi,\psi)$ is a Galois-Tukey connection from
$(\Cal S^{(0)},\finint,\NN)$ to $(\NN,\disj,\NN)$.\ \Qed

\medskip

\qquad\grheadc\ $\cov(\Cal S^{(0)},\finint,\NN)=\frakmctbl$.   \Prf\

$$\eqalignno{\frakmctbl
&=n(\NN)=\cov\Cal M_1\cr
\displaycause{517Pd}
&=\cov(\NN,\in,\Cal M_1)\le\cov(\Cal S^{(0)},\finint,\NN)\cr
\displaycause{512Da and ($\alpha$) above}
&\le\cov(\NN,\disj,\NN)\cr
\displaycause{($\beta$) above}
&=\frakmctbl\cr}$$

\noindent (522Sb, with 522Sa).\ \Qed

\medskip

\quad{\bf (iii)} Suppose that $\kappa<\add\Cal N$ and that
$\ofamily{\xi}{\kappa}{S_{\xi}}$ is any family in $\Cal S^{(0)}$.   Then
there is an $S^*\in\Cal S^{(0)}$ such that
$S_{\xi}\setminus S^*$ is finite for every $\xi<\kappa$.   \Prf\ For
$\xi<\kappa$, $n\in\Bbb N$ let $f_{\xi}(n)\in\Bbb N$ be such that
$\#(S_{\xi}[\{i\}])\le 2^{i-2n}$ for every $i\ge f_{\xi}(n)$.   Because
$\kappa<\add\Cal N\le\frak b$, there is an $f\in\NN$ such that
$\{n:f_{\xi}(n)>f(n)\}$ is finite for every $\xi<\kappa$ (522C(ii));
of course we may suppose that $f(0)=0$ and that $f$ is strictly
increasing
and that $f(n)\ge 2n$ for every $n$.   Set $J_n=f(n+1)\setminus f(n)$
for each $n$.   For each $\xi<\kappa$, let $m_{\xi}$ be such that
$f_{\xi}(n)\le f(n)$ for every $n\ge m_{\xi}$;   set
$S'_{\xi}=\{(i,j):(i,j)\in S_{\xi}$,
$i\ge f(m_{\xi})\}$.   Then $S_{\xi}\setminus S'_{\xi}$ is finite and
$\#(S'_{\xi}[\{i\}])\le 2^{i-2n}$ whenever $i\in J_n$.

For each $n\in\Bbb N$, let $\Cal K_n$ be the family of those sets
$K\subseteq J_n\times\Bbb N$ such that $\#(K[\{i\}])\le 2^{i-2n}$ for
every $i\in J_n$, and $\theta_n:\Cal K_n\to\Bbb N$ a bijection;  set
$h_{\xi}(n)=\theta_n(S'_{\xi}\cap(J_n\times\Bbb N))$ for each
$\xi<\kappa$.

By 522M, $\add(\NN,\subseteq^*,\Cal S)=\add\Cal N$ is greater than
$\kappa$, so there is an $S\in\Cal S$ such that $h_{\xi}\subseteq^*S$
for every $\xi<\kappa$.   Set
$S^*=\bigcup_{(n,j)\in S}\theta_n^{-1}(j)$.   For any $n\in\Bbb N$ and
$i\in J_n$, $\#(\theta_n^{-1}(j)[\{i\}])\le 2^{i-2n}$ for every
$j\in\Bbb N$, so that
$S^*[\{i\}]=\bigcup_{(n,j)\in S}\theta_n^{-1}(j)[\{i\}]$ has cardinal at
most $2^{i-n}$.   This means that $S^*\in\Cal S^{(0)}$.

Take any $\xi<\kappa$.   As $h_{\xi}\subseteq^*S$, there is some
$m\in\Bbb N$ such that $(n,h_{\xi}(n))\in S$, that is,
$(n,\theta_n(S'_{\xi}\cap(J_n\times\Bbb N)))\in S$, for every $n\ge m$.
But this means that $S'_{\xi}\cap(J_n\times\Bbb N)\subseteq S^*$ for
every $n\ge m$, so $S'_{\xi}\setminus S^*$ is finite;  it follows at
once that $S_{\xi}\setminus S^*$ is finite.   Thus we have a suitable
$S^*$.\ \Qed

\medskip

\quad{\bf (iv)} \Quer\ Now suppose, if possible, that
$\cf(\frakmctbl)=\kappa<\add\Cal N$.   By (ii-$\gamma$), there is a
set $A\subseteq\NN$ of
size $\frakmctbl$ such that for every $S\in\Cal S^{(0)}$ there is an
$f\in A$ such that $S\cap f$ is finite.   Express $A$ as
$\bigcup_{\xi<\kappa}A_{\xi}$ where $\#(A_{\xi})<\frakmctbl$ for every
$\xi<\kappa$.   By (ii-$\gamma$) again, we can find for each $\xi<\kappa$ an
$S_{\xi}\in\Cal S^{(0)}$ such that $S_{\xi}\cap f$ is infinite for every
$f\in A_{\xi}$.   By (iii), there is an $S^*\in\Cal S^{(0)}$ such that
$S_{\xi}\setminus S^*$ is finite for every $\xi<\kappa$.   But this
means that $S^*\cap f$ must be infinite for every $f\in A_{\xi}$ and
every $\xi<\kappa$;  which contradicts the choice of $A$.\ \Bang

So we are forced to conclude that $\cf(\frakmctbl)\ge\add\Cal N$, as
stated.
}%end of proof of 522V

\leader{522W}{Other spaces}\cmmnt{ All the theorems above refer to the
specific
$\sigma$-ideals $\Cal M$ and $\Cal N$ of subsets of $\Bbb R$ or the
specific partially ordered set $\NN$.   Of course the structures
involved appear in many other guises.   In particular, we have the
following results.

\medskip

} {\bf (a)(i)} Let $(X,\Sigma,\mu)$ be an atomless countably
separated\cmmnt{ (definition:  343D)} $\sigma$-finite
perfect\cmmnt{ (definition:  342K)} measure space of non-zero measure,
and $\Cal N(\mu)$ the null ideal of $\mu$.   Then
$(X,\Cal N(\mu))$ is isomorphic to $(\Bbb R,\Cal N)$;  in particular,
$\add\Cal N(\mu)=\add\Cal N$, $\cov\Cal N(\mu)=\cov\Cal N$,
$\non\Cal N(\mu)=\non\Cal N$ and $\cf\Cal N(\mu)=\cf\Cal N$.
\prooflet{\Prf\
The first thing to note is that because $\mu$ is $\sigma$-finite there
is a probability measure $\nu$ on $X$ with the same measurable sets and
the same negligible sets as $\mu$ (215B(vii));  and of course $\nu$ is
still atomless, countably separated and perfect.   Next, the completion
$\hat\nu$ of $\nu$ is again atomless, countably separated and perfect
(212Gd, 343H(vi), 451G(c-i))
and has the same negligible sets as $\nu$ (212Eb).   In the same way,
starting from Lebesgue measure instead of $\mu$, we have a complete
atomless countably separated perfect probability measure $\lambda$ on
$\Bbb R$ with the same negligible sets as Lebesgue measure.   But now
$(X,\hat\nu)$ and $(\Bbb R,\lambda)$ are isomorphic (344I), so that
$(X,\Cal N(\mu))$ and $(\Bbb R,\Cal N)$ are
isomorphic.\ \Qed}%end of prooflet

\medskip

\quad{\bf (ii)}\cmmnt{ The most important examples of spaces
satisfying the conditions of (i) are Lebesgue measure on the unit
interval and the usual measure on $\{0,1\}^{\Bbb N}$.   But the ideas go
much farther.}   On a Hausdorff space with a countable
network\cmmnt{ (e.g., any separable metrizable space, or any analytic
Hausdorff space), any topological measure is countably separated (433B).
So} any non-zero atomless Radon measure\cmmnt{ on such a space} will
have a null ideal isomorphic to
$\Cal N$.   \cmmnt{(The measure will be $\sigma$-finite because it is
a locally finite measure on a Lindel\"of space, and perfect by 416Wa.)}

\cmmnt{\medskip

\quad{\bf (iii)} As we shall see in \S523, there are
many more measure spaces $(X,\mu)$ for which $\Cal N(\mu)$ is close
enough to $\Cal N$ to have the same additivity and cofinality, and even
uniformity and covering number match in a number of interesting cases.}

\spheader 522Wb{\bf(i)}\cmmnt{ Similarly, the structure}
$(\Bbb R,\Cal M)$ is
duplicated in any non-empty Polish space $X$ without isolated points, in
the sense that $(X,\Cal B(X),\Cal M(X))\cong(\Bbb R,\Cal B,\Cal M)$,
where
$\Cal B$ and $\Cal B(X)$ are the Borel $\sigma$-algebras of $\Bbb R$ and
$X$ respectively, and $\Cal M(X)$ is the ideal of meager subsets of $X$.
\prooflet{\Prf\ Note first that $\NN$, with its usual topology, has an
uncountable
nowhere dense closed set;  e.g., $\{f:f(2n)=0$ for every $n\}$.   Now we
know that $X$ has a dense G$_{\delta}$ set $X_1$ homeomorphic to $\NN$
(5A4Le), and $X_1$ must also have an uncountable nowhere dense closed set
$F_1$;  since $X_1\setminus F_1$ is again a non-empty Polish space
without isolated points (4A2Qd), it too has a dense G$_{\delta}$ set
$X_2$ homeomorphic to $\NN$, and $X_2$ is a dense G$_{\delta}$ set in
$X$ with
uncountable complement.   Similarly, $\Bbb R$ has a dense G$_{\delta}$
subset $H$ which is homeomorphic to $\NN$ and has uncountable
complement.

Let $\Cal M(X_2)$, $\Cal M(H)$ be the ideals of meager subsets of $X_2$
and $H$ when they are given their subspace topologies.   Because $X_2$
is dense, a closed subset of $X$ is nowhere dense in $X$ iff its
intersection with $X_2$ is nowhere dense in $X_2$;  accordingly
$\Cal M(X_2)$ is precisely $\{M\cap X_2:M\in\Cal M(X)\}$.
Similarly, $\Cal M(H)=\{M\cap H:M\in\Cal M\}$.

Consider the complements $X\setminus X_2$, $\Bbb R\setminus H$.   These
are uncountable Borel subsets of Polish spaces.   They are therefore
Borel isomorphic (424G, 424Cb);  let
$\phi:X\setminus X_2\to\Bbb R\setminus H$ be a Borel isomorphism.
Next, $X_2$ and $H$ are homeomorphic to $\NN$, therefore to each other;
let $\psi:X_2\to H$ be a homeomorphism.   Finally, set
$\theta=\psi\cup\phi$, so that $\theta:X\to\Bbb R$ is a Borel
isomorphism.   For $M\subseteq X$,

$$\eqalignno{M\in\Cal M(X)
&\iff M\cap X_2\in\Cal M(X)\cr
\displaycause{because $X\setminus X_2$ is meager}
&\iff M\cap X_2\in\Cal M(X_2)
\iff\psi[M\cap X_2]\in\Cal M(H)\cr
&\iff\theta[M]\cap H\in\Cal M(H)
\iff\theta[M]\in\Cal M.\cr}$$

\noindent So $\theta$ is an isomorphism between the structures
$(X,\Cal B(X),\Cal M(X))$ and $(\Bbb R,\Cal B,\Cal M)$.\ \Qed
}%end of prooflet

\medskip

\quad{\bf (ii)} Again, the most important special cases here are
$X=[0,1]$, $X=\{0,1\}^{\Bbb N}$ and $X=\NN$.

\exercises{\leader{522X}{Basic exercises $\pmb{>}$(a)}
%\sqheader 522Xa
Let $\Cal K$ be the $\sigma$-ideal of subsets of $\NN$
generated by the compact sets.   Show that $(\Cal K,\subseteq)$ is Tukey
equivalent to the pre-ordered sets of 522C, so that
$\add\Cal K=\frak b$ and $\cf\Cal K=\frak d$.
%522C

\spheader 522Xb\dvAformerly{5{}22Xh} (O.Kalenda)
Set $(P,\sqsubseteq)=(\Bbb N,\le)\times(\NN,\preceq)$ where
$\preceq$ is the partial ordering of 522C(ii).
Show that $(\Bbb N,\le)\prT(P,\sqsubseteq)\prT(\NN,\le)$,
$(P,\sqsubseteq)\not\prT(\Bbb N,\le)$ and $(\NN,\le)\not\prT(P,\sqsubseteq)$.
\leaveitout{
for last: \Quer if  $\phi:\NN\to P$  is Tukey say
$\phi(f)=(\phi_1(f),\phi_2(f))$ , take  $n$  such that  $A=\{f:\phi_f(f)=n\}$
is non-meager therefore unbounded above in  $\NN$ , but
$\phi_2:A\to(\NN,\preceq)$
is Tukey and  $\phi_2[B]$  is bounded above in  $(\NN,\preceq)$  for every
countable  $B\subseteq A$ so every countable subset of $A$ is bounded above
in $\NN$ \Bang
}
%522C

\spheader 522Xc Let $(X,\Sigma,\mu)$ be an atomless semi-finite measure
space with $\mu X>0$.   Show that $\#(X)\ge\non\Cal N$.
\Hint{343Cc.}
%522G

\sqheader 522Xe Show that there are just 23 assignments of values to
the cardinals of Cicho\'n's diagram which are allowed by the results in
522D-522Q and have $\frak c=\omega_2$.
%522R mt52bits

\spheader 522Xg Show that if $\cov\Cal N>\omega_1$ then every
$\pmb{\Delta}^1_2$ (= PCA-\&-CPCA) set in a Polish space is universally
measurable.  \Hint{423Tb\formerly{4{}23Rb}, 521Xc.}
%522+

\spheader 522Xh Let $Z$ be the Stone space of the measure algebra $\frak A$ of
Lebesgue measure.   Show that the Nov\'ak number $n(Z)$ of $Z$ and
the Martin number $\frak m(\frak A)$ of $\frak A$ are both equal to
$\cov\Cal N$.   \Hint{341Q, 416V, 517K.}
%522+

\leader{522Y}{Further exercises (a)}
%\spheader 522Ya
Show that if $\add\Cal N=\cf\Cal N$ then
$(\Bbb R,\Cal M)$
and $(\Bbb R,\Cal N)$ are isomorphic, in the sense that there is a
permutation $f:\Bbb R\to\Bbb R$ such that $A\subseteq\Bbb R$ is meager iff
$f[A]$ is Lebesgue negligible.
%522Q

\spheader 522Yb Show that if $\cov\Cal N>\omega_1$ then
$\cov\Cal N\ge\frak m_{\text{pc}\omega_1}$.
%522T  517Ig  522Xh  525Td

\spheader 522Yc Let $P$ and $Q$ be partially ordered sets such that
$Q$ has no greatest member, $\sim$ an
equivalence relation on $P$, and $\pi:P\to Q$ a surjective
function such that, for $p_0$, $p_1\in P$, $\pi(p_0)\le\pi(p_1)$ iff
there is
a $p\sim p_0$ such that $p\le p_1$.   Suppose that $\kappa$ is a cardinal
such that no $\sim$-equivalence class has cardinal greater than $\kappa$.
Show that $\add(Q)\le\max(\FN(P),\kappa)$.
%522Ua

\spheader 522Yd Suppose that $\FN(\Cal P\Bbb N)=\omega_1$.   Show that
whenever $A\subseteq\Bbb R$ is non-meager there is a set
$B\in[A]^{\omega_1}$ such that every uncountable subset of $B$ is
non-meager\footnote{In the language of 554C, every non-meager subset of
the real line includes a Lusin set.}.
%522Uc

\spheader 522Ye Suppose that $\FN(\Cal P\Bbb N)=\frak p$ and that
$\kappa\ge\frakmctbl$ is such that $\cff[\kappa]^{<\frak p}\le\kappa$.
Show that $\kappa\ge\frak c$.
%522Ud 517Yd

\spheader 522Yf (S.Geschke)
Show that if $\FN^*(\Cal P\Bbb N)\le\frakmctbl$ then
$\non\Cal M\le\FN^*(\Cal P\Bbb N)$.
%522Uc mt52bits

\spheader 522Yg Let $\Cal S^{(0)}$ be the family described in the proof of
522Vf.   For any sets $A$, $B$ say that $A\subseteq^*B$ if $A\setminus B$
is finite, and define $\le^*$ as in 522C.   Show that
$(\Cal S^{(0)},\subseteq^*,\Cal S^{(0)})
\prGT(\NN,\le^*,\penalty-50\NN)\penalty-100\ltimes
(\NN,\subseteq^*,\Cal S^{(0)})$.
%522V mt52bits

\spheader 522Yh Suppose that we have supported relations $(A,R,B)$ and
$(A,S,A)$ such that $R\frsmallcirc S\subseteq R$, that is,
$(a,b)\in R$ whenever $(a,a')\in S$ and $(a',b)\in R$.
Show that if $\omega\le\cov(A,R,B)<\infty$ then
$\cf(\cov(A,R,B))\ge\add(A,S,A)$.
%522V

\spheader 522Yi
Let $X$ be any topological space with countable
$\pi$-weight and write $\Cal M(X)$ for the family of meager subsets of
$X$.   Show that $\Cal M(X)\prT\Cal M$.
%522W reduce to X zero-dimensional no non-empty meager open set, find
%  continuous  f:X\to\{0,1\}^{\Bbb N}  such that  E\subseteq X  is meager
%  iff  f[E]  meager
}%end of exercises

\leader{522Z}{Problem} Is it the case that
$(\Bbb R,\in,\Cal M)\equivGT(\NN,\finint,\NN)$?   
\cmmnt{(See 522S and the proof of 522V.)}

\leaveitout{
What 522S needs is equality of additivities and covering numbers.

{\bf (a)} $(\Bbb R,\in,\Cal M)\prGT(\NN,\finint\strprime,[\NN]^{\le\omega})$.
\Prf\
Take homeomorphism $\phi_0:\Bbb R\setminus\Bbb Q\to\NN$.   Set
$\phi(x)=\phi_0(x)$ for $x\in\Bbb R\setminus\Bbb Q$, $0$ if $x\in\Bbb Q$.   Set
$\psi(I)=\bigcup_{f\in I}\phi^{-1}[\{g:g\in\NN$, $g\cap f$ is finite$\}]
\in\Cal M$.
If $x\in\psi(I)$ then $(\phi(x),I)\in\finint\strprime$.\ \Qed

So $(\Bbb R,\in,\Cal M)\prGT(\NN,\finint,\NN)$ (512Ga, 512Gb),
$\add(\NN,\finint,\NN)\le\non\Cal M$ and
$\cov\Cal M\le\cov(\NN,\finint,\NN)=\cov(\NN,\disj,\NN)$.

 can it be that
$(\NN,\finint\strprime,[\NN]^{\le\omega})\prGT(\Bbb R,\in\Cal M)$?  and
what is $\add(\NN,\finint\strprime,[\NN]^{\le\omega})$?

{\bf (b)} Maybe
$(\NN,\finint,\NN)\prGT(\NN,\preceq,\NN)\ltimes(\Bbb R,\in,\Cal M)$?
 but $\add((\NN,\preceq,\NN)\ltimes(\Bbb R,\in,\Cal M))
=\min(\frak b,\non\Cal M)$ (512Jc, 522C(ii), 512Ed)
\query


}%end of leaveitout

\endnotes{
\Notesheader{522} All the significant ideas of this section may be found
in {\smc Bartoszy\'nski \& Judah 95}, with a good deal more.

For many years it appeared that `measure' and
`category' on the real line, or at least the structures
$(\Bbb R,\Cal B,\Cal N)$ and $(\Bbb R,\Cal B,\Cal M)$ where $\Cal B$ is
the Borel $\sigma$-algebra of $\Bbb R$, were in a symmetric duality.
It was perfectly well understood that the algebras
$\frak A=\Cal B/\Cal B\cap\Cal N$ and $\frak G=\Cal B/\Cal B\cap\Cal M$
-- what in this book I call the `Lebesgue measure algebra' and the
`category algebra of $\Bbb R$' -- are very different, but their
complexities seemed to be balanced, and such results as 522G
encouraged us to suppose that anything provable in ZFC relating measure
to category ought to respect the symmetry.   It therefore came as a
surprise to most of us when Bartoszy\'nski and Raisonnier \& Stern
(independently, but both drawing inspiration from ideas of {\smc Shelah
84}, themselves responding to a difficulty noted in {\smc Solovay 70})
showed that $\add\Cal N\le\add\Cal M$ in all models of set theory.   (It
was already known that $\add\Cal N$ could be strictly less than
$\add\Cal M$.)

The diagram in its present form emphasizes a new dual symmetry,
corresponding to the duality of Galois-Tukey connections (512Ab).   No
doubt this also is only part of the true picture.   It gives no hint,
for instance, of a striking difference between $\cov\Cal M$ and
$\cov\Cal N$.   While $\cov\Cal M=\frakmctbl$ must have uncountable
cofinality (522Vf), $\cov\Cal N$ can be $\omega_{\omega}$
({\smc Shelah 00}).   In 522H-522I and 522Sb-522Sc there are hints
of a different
symmetry which I have not been able to formalize convincingly (see 522Z).

I have hardly mentioned shrinking numbers here.   This is because while
$\shr\Cal M$ and $\shr\Cal N$ can be located in Cicho\'n's diagram (we
have $\non\Cal M\le\shr\Cal M\le\cf\Cal M$ and
$\non\Cal N\le\shr\Cal N\le\cf\Cal N$, by 511Jc), they are not known to
be connected organically with the rest of the diagram.   I will return to
them in a more general context in 523M.   I have also not said where the
$\pi$-weight of Lebesgue measure (see 511Gb) fits in;  this is in fact
equal to $\cf\Cal N$, as will appear in 524P.

In 522T I give two classic `Martin's axiom' arguments.   They are
typical in that the structure of the proof is to establish that there is
a suitable
partially ordered set for which a `generic' upwards-directed subset will
provide an object to witness the truth of some assertion.   `Generic',
in this context, means `meeting sufficiently many cofinal sets'.   If
there were any more definite method of finding the object sought, we
would use it;  these constructions are always even more ethereal than
those which depend on unscrupulous use of the axiom of choice.
`Really' they are names of propositions in a suitable forcing language,
since (as a
rule) we can lift Martin numbers above $\omega_1$ only by entering a
universe created by forcing.   But in this chapter, at least, I will
try to avoid such considerations, and use arguments which are
expressible in the ordinary language of ZFC, even though their
non-trivial applications depend on assumptions beyond ZFC.

Of the partially ordered sets $\Cal S^{\infty}$ and $P$ in the proof of
522T, the former comes readily to hand as soon as we cast the problem in
terms of the supported relation $(\NN,\subseteq^*,\Cal S)$;  we need
only realize that we can express members of $\Cal S$ as limits of
upwards-directed subsets of a subfamily in which there is some room to
manoeuvre, so that we have enough cofinal sets.   The latter is more
interesting.   It belongs to one of the standard types in that the
partially ordered set is made up of pairs $(\sigma,F)$ in which $\sigma$
is the `working part', from which the desired meager set

\Centerline{$M=\Bbb R\setminus\bigcap_{n\in\Bbb N}
  \bigcup_{\sigma\in R,i\ge n}\sigma(i)$}

\noindent will be constructed, and $F$ is a `side condition', designed
to ensure that the partial order of $P$ interacts correctly with the
problem.   In such cases, there is generally a not-quite-trivial step to
be made in proving that the ordering is transitive ((b-i) of the proof
of 522T).   Note that we have two classes of cofinal set to declare in
(b-iii) of the proof here;  the $Q_{nV}$ are there to ensure that $M$ is
meager, and the $Q'_H$ to ensure that it includes every member of
$\Cal H$.   And a final element which must appear in every proof of this
kind, is the check that the partial order found is of the correct
type, $\sigma$-linked in (a) and $\sigma$-centered in (b).

In 522U I suggest that it is natural to try to locate any newly defined
cardinal among those displayed in Cicho\'n's diagram.   Of course there
is no presumption that it will be possible to do this tidily, or that we
can expect any final structure to be low-dimensional;  the picture in
522T is already neater than we are entitled to expect, and the
complications in 522U (and 522Yd-522Yf) %522Yd 522Ye 522Yg 522Yf
are a warning that our luck may be running out.   However, we can
surprisingly often find relationships like the ones between
$\FN(\Cal P\Bbb N)$, $\frak b$, $\shr\Cal M$ and $\frakmctbl$ here,
which is one of my reasons for using this approach.   It is very
remarkable that under fairly weak assumptions on cardinal arithmetic
(the hypothesis
`$\frakmctbl<\omega_{\omega}$' in 522Ud is much stronger than is
necessary, since in `ordinary' models of set theory we have
$\cff[\kappa]^{\le\omega}=\kappa$ whenever $\cf\kappa>\omega$ -- see
5A6Bc and 5A6C),
%Jech 03, \S36
the axiom `$\FN(\Cal P\Bbb N)=\omega_1$' splits Cicho\'n's diagram neatly
into two halves.   For an explanation of why it was worth looking for
such a split, see {\smc Fuchino Geschke \& Soukup 01}.

For the sake of exactness and simplicity, I have maintained rigorously
the convention that $\Cal M$ and $\Cal N$ are the ideals of meager and
negligible sets in $\Bbb R$ with Lebesgue measure.   But from the point
of view of the diagram, they are `really' representatives of
classes of ideals defined on non-empty Polish spaces without isolated
points, on the one hand, and on atomless countably separated
$\sigma$-finite perfect measure spaces of non-zero measure on the other
(522W).   The most natural expression of the duality between the
supported relations $(\Bbb R,\in,\Cal M)$ and $(\Bbb R,\in,\Cal N)$
(522G) depends, of course, on the fact that both structures are
invariant under translation;  but even this is duplicated in $\BbbR^r$
and in infinite compact metrizable groups like $\{0,1\}^{\Bbb N}$.

At some stage I ought to mention a point concerning the language of this
chapter.   It is natural to think of such expressions as $\add\Cal N$ as
names for
objects which exist in some ideal universe.   Starting from such a
position, the sentence `it is possible that $\add\Cal N<\add\Cal M$' has
to be interpreted as
`there is a possible mathematical universe in which
$\add\Cal N<\add\Cal M$'.   But this can make sense only if
`$\add\Cal N$' can refer to different objects in
different universes, and has a meaning independent of any particular
incarnation.   I think that in fact we have to start again,
and say that the expression $\add\Cal N$ is not a name
for an object, but an abbreviation of a definition.   We can then speak
of the interpretations of that definition in different worlds.   In fact
we have to go much
farther back than the names for cardinals in this section.
$\Cal P\Bbb N$ and $\Bbb R$ also have to be considered primarily as
definitions.   The set $\Bbb N$
itself has a relatively privileged position;  but even here it is
perhaps safest to regard the symbol $\Bbb N$ as a name for a formula in
the language of set theory
rather than anything else.   Fortunately, one can do mathematics without
aiming at perfect consistency or logical purity, and I will make no
attempt to disinfect
my own language beyond what seems to be demanded by the ideas I am
trying to express at each moment;  but you should be aware that there
are possibilities for confusion here, and that at some point you will
need to find your own way of balancing among them.   My own practice,
when the path does not seem clear, is to re-read {\smc Kunen 80}.
}%end of notes

\discrpage

\frfilename{mt523.tex}
\versiondate{24.8.24}
\copyrightdate{2005}

\def\chaptername{Cardinal functions of measure theory}
\def\sectionname{The measure of $\{0,1\}^I$}

\newsection{523}

\def\headlinesectionname{The measure of
$\{0,1{\delimiter"5267309 }^I$}

In \S522 I tried to give an account of current knowledge concerning
the most important cardinals associated with Lebesgue measure.   The
next step is to investigate the usual measure $\nu_I$ on $\{0,1\}^I$
for an arbitrary set $I$.   Here I discuss the cardinals
associated with these measures.   Obviously they depend only on $\#(I)$,
and are trivial if $I$ is finite.   I start with the basic diagram
relating the cardinal functions
of $\nu_{\kappa}$ and $\nu_{\lambda}$ for different cardinals $\kappa$
and $\lambda$ (523B).   I take the opportunity to mention some simple facts
about the measures $\nu_I$ (523C-523D). %523C 523D
Then I look at additivities (523E), covering numbers (523F-523G),
uniformities (523H-523L), %523H 523I 523J 523K 523L
shrinking numbers (523M) and cofinalities (523N).   I end with a
description of these cardinals under the generalized continuum
hypothesis (523P).

\leader{523A}{Notation} For any set $I$, I will write $\nu_I$ for
the usual measure on $\{0,1\}^I$ and $\Cal N_I$ for its
null ideal.   \cmmnt{Recall that
$(\{0,1\}^{\omega},\Cal N_{\omega})$ is isomorphic to $(\Bbb R,\Cal N)$,
where $\Cal N$ is the Lebesgue null ideal (522Wa).}

\leader{523B}{The basic diagram} Suppose that $\kappa$ and $\lambda$ are
infinite cardinals, with $\kappa\le\lambda$.   Then we have the
following
diagram\cmmnt{ for the additivity, covering number, uniformity,
shrinking number and cofinality of the ideals $\Cal N_{\kappa}$ and
$\Cal N_{\lambda}$}:

\def\tmphrule{\hskip0.4em\raise
2.5pt\hbox{\leaders\hrule\hskip1.7em\hfil}\hskip0.4em}
\def\tmpstrut{\vrule height10.5pt depth5.5pt width0pt}

$$\vbox{\offinterlineskip
\halign{\hfil#\hfil&\hfil#\hfil&\hfil#\hfil&\hfil#\hfil
  &\hfil#\hfil&\hfil#\hfil&\hfil#\hfil
  &\hfil#\hfil&\hfil#\hfil&\hfil#\hfil
  &\hfil#\hfil&\hfil#\hfil&\hfil#\hfil&\hfil#\hfil\cr
&\tmpstrut&$\cov\Cal N_{\lambda}$&\tmphrule
  &$\cov\Cal N_{\kappa}$&\tmphrule
  &$\cf\Cal N_{\kappa}$&\tmphrule&$\cf\Cal N_{\lambda}$
  &\tmphrule&$\lambda^{\omega}$\cr
&\tmpstrut&\vrule&&\vrule&&\vrule&&\vrule\cr
&\tmpstrut&\vrule&&\vrule&&$\shr\Cal N_{\kappa}$&\tmphrule
  &$\shr\Cal N_{\lambda}$\cr
&\tmpstrut&\vrule&&\vrule&&\vrule&&\vrule\cr
$\tmpstrut\omega_1$&\tmphrule&$\add\Cal N_{\lambda}$&\tmphrule
  &$\add\Cal N_{\kappa}$&\tmphrule
  &$\non\Cal N_{\kappa}$&\tmphrule&$\non\Cal N_{\lambda}$\cr
}}$$

\noindent (As in 522B, the cardinals here increase from bottom
left to top right.)

\proof{ For the inequalities relating two cardinals associated with the
same ideal, see 511Jc;  all we need to know is that $\Cal N_{\kappa}$
and $\Cal N_{\lambda}$ are proper ideals containing singletons.   For
the inequalities relating the cardinal functions of the two different
ideals, use 521H;  $\nu_{\kappa}$ is the image of $\nu_{\lambda}$ under
the map $x\mapsto x\restr\kappa:\{0,1\}^{\lambda}\to\{0,1\}^{\kappa}$,
by 254Oa.   Of course $\omega_1\le\add\Cal N_{\lambda}$.   I leave the
final inequality $\cf\Cal N_{\lambda}\le\lambda^{\omega}$ for the
moment, since this will be part of Theorem 523N below.
}%end of proof of 523B

\leader{523C}{}\cmmnt{ In the next few paragraphs I will set out what is
known about the cardinals here.   It will be convenient to begin with
two easy lemmas.

\medskip

\noindent}{\bf Lemma} Let $I$ be any set, and $\Cal J$ a family of
subsets of $I$ such that every countable subset of $I$ is included in
some member of $\Cal J$.   Then a subset $A$ of $\{0,1\}^I$ belongs to
$\Cal N_I$ iff there is some $J\in\Cal J$ such that
$\{x\restr J:x\in A\}\in\Cal N_J$.

\proof{ For $J\subseteq I$, $x\in\{0,1\}^I$ set
$\pi_J(x)=x\restr J\in\{0,1\}^J$.
Then $\nu_J$ is the image measure $\nu_I\pi_J^{-1}$ (254Oa again),
so $A\in\Cal N_I$ whenever there
is some $J\in\Cal J$ such that $\pi_J[A]\in\Cal N_J$.   On the other
hand, if $A\in\Cal N_I$, there is a countable set $K\subseteq I$ such
that $\pi_K[A]\in\Cal N_K$ (254Od).   Now there is a $J\in\Cal J$ such
that $K\subseteq J$, so that
$\pi_J^{-1}[\pi_J[A]]\subseteq\pi_K^{-1}[\pi_K[A]]\in\Cal N_I$ and
$\pi_J[A]\in\Cal N_J$.
}%end of proof of 523C

\leader{523D}{}\cmmnt{ Because the measures $\nu_I$ are homogeneous in a
strong sense, we have the following facts which are occasionally useful.

\medskip

\noindent}{\bf Proposition} Let $\kappa$ be an infinite cardinal, and
$\Tau$ the domain of $\nu_{\kappa}$.
For $A\subseteq\{0,1\}^{\kappa}$ write $\Tau_A$ for the subspace
$\sigma$-algebra on $A$.

(a) If $E\subseteq\{0,1\}^{\kappa}$ is measurable and not negligible,
then $(E,\Tau_E,\Cal N_{\kappa}\cap\Cal PE)$ is
isomorphic to $(\{0,1\}^{\kappa},\Tau,\Cal N_{\kappa})$.

(b) If $\Cal E\subseteq\Cal N_{\kappa}$ and
$\#(\Cal E)<\cov\Cal N_{\kappa}$, then $(\nu_{\kappa})_*(\bigcup\Cal E)=0$.

(c) If $A\subseteq\{0,1\}^{\kappa}$ is non-negligible,
then there is a set $B\subseteq\{0,1\}^{\kappa}$, of full outer measure,
such that $(A,\Tau_A,\Cal N_{\kappa}\cap\Cal PA)$ is
isomorphic to $(B,\Tau_B,\Cal N_{\kappa}\cap\Cal PB)$.

(d) There is a set
$A\subseteq\{0,1\}^{\kappa}$ with cardinal $\non\Cal N_{\kappa}$ which has
full outer measure.
%old 523Ge

\proof{{\bf (a)} In fact the subspace measure on $E$ is isomorphic to
a scalar multiple of $\nu_I$ (344L).

\medskip

{\bf (b)} \Quer\ Otherwise, let $F\subseteq\bigcup\Cal E$ be a
non-negligible measurable set;  then $\{F\cap E:E\in\Cal E\}$ witnesses
that $\cov(F,\Cal N_{\kappa}\cap\Cal PF)<\cov\Cal N_{\kappa}$, which
contradicts (a).\ \Bang

\medskip

{\bf (c)} Let $E$ be a measurable envelope of $A$.   By (a), there is a
bijection $f:E\to\{0,1\}^{\kappa}$ which is an isomorphism of the
structures $(E,\Tau_E,\Cal N_{\kappa}\cap\Cal PE)$ and
$(\{0,1\}^{\kappa},\Tau,\Cal N_{\kappa})$.   Set $B=f[A]$.
Then $f\restr A$ is an isomorphism of the structures
$(A,\Tau_A,\Cal N_{\kappa}\cap\Cal PA)$ and
$(B,\Tau_B,\Cal N_{\kappa}\cap\Cal PB)$.   Moreover, since $A$ meets every
member of $\Tau_E\setminus\Cal N_{\kappa}$, $B$ meets every member of
$\Tau\setminus\Cal N_{\kappa}$, that is, $B$ has full outer measure.

\medskip

{\bf (d)} Let $A_0\subseteq\{0,1\}^{\kappa}$ be a non-negligible set of
cardinal $\non\Cal N_{\kappa}$.   By (c), there is a set $A$ of full outer
measure which is isomorphic to $A_0$ in the sense described there;  in
particular, $\#(A)=\non\Cal N_{\kappa}$.
}%end of proof of 523D

\leader{523E}{Additivities}\cmmnt{ Because the function
$\kappa\mapsto\add\Cal N_{\kappa}$ is non-increasing, it must stabilize,
that is, there is some first $\kappa_a$ such that
$\add\Cal N_{\kappa}=\add\Cal N_{\kappa_a}$ for every
$\kappa\ge\kappa_a$.   But
in fact it stabilizes almost immediately.}   If $\kappa$ is any
uncountable cardinal, then
$\add\Cal N_{\kappa}=\add\nu_{\kappa}=\omega_1$\prooflet{, by 521Jb}.
\cmmnt{ Thus among the additivities $\add\Cal N_{\kappa}$, only
$\add\Cal N_{\omega}=\add\Cal N$, the additivity of Lebesgue
measure, can have any surprises for us.}

\leader{523F}{Covering numbers} Still on the left-hand side of
the diagram, we\cmmnt{ again} have a non-increasing function
$\kappa\mapsto\cov\Cal N_{\kappa}$, and a critical value $\kappa_c$
after which it is constant.   \cmmnt{We can locate this value to some
extent through the following simple fact.  If
$\theta\dvro{}{\mskip5mu=\cov\Cal N_{\kappa_c}}
=\min\{\cov\Cal N_{\kappa}:\kappa$ is a cardinal$\}$, then
$\cov\Cal N_{\theta}=\theta$.   \prooflet{\Prf\ Let $\kappa$ be such
that $\cov\Cal N_{\kappa}=\theta$.   For $I\subseteq\kappa$, set
$\pi_I(x)=x\restr I$ for $x\in\{0,1\}^{\kappa}$.   Let
$\Cal E\subseteq\Cal N_{\kappa}$ be a cover of $\{0,1\}^{\kappa}$ of
cardinality $\theta$.   For each $E\in\Cal E$, let $J_E\subseteq\kappa$
be a countable set such that $\pi_{J_E}[E]\in\Cal N_{J_E}$.   Set
$I=\bigcup_{E\in\Cal E}J_E$,
so that $\#(I)\le\theta$ and $\pi_I[E]\in\Cal N_I$ for every
$E\in\Cal E$.   Then
$\{\pi_I[E]:E\in\Cal E\}$ is a cover of $\{0,1\}^I$ by at most
$\cov\Cal N_{\kappa}$ sets, and
$\cov\Cal N_I\le\cov\Cal N_{\kappa}$.   Since $(\{0,1\}^I,\Cal N_I)$ is
isomorphic to $(\{0,1\}^{\#(I)},\Cal N_{\#(I)})$, we also have

\Centerline{$\cov\Cal N_{\theta}\le\cov\Cal N_{\#(I)}
\le\cov\Cal N_{\kappa}\le\cov\Cal N_{\theta}$,}

\noindent and $\cov\Cal N_{\theta}=\cov\Cal N_{\kappa}=\theta$.\ \Qed}

What this means is that}

\Centerline{$\omega\le\kappa_c\cmmnt{\mskip5mu\le\theta}\le\cov\Cal N_{\kappa_c}
\le\cov\Cal N_{\omega}=\cov\Cal N\le\frak c$.}

\noindent Another way of putting the same idea is to say that

\cmmnt{\Centerline{if $\theta<\lambda$ then 
$\cov\Cal N_{\lambda}\le\cov\Cal N_{\theta}=\theta<\lambda$}

\noindent so that} 

\Centerline{if $\lambda$ is a cardinal such that
$\cov\Cal N_{\lambda}\ge\lambda$ then $\cov\Cal N_{\kappa}\ge\lambda$ for
every $\kappa$.}

%See {\smc Brendle 06},
%Theorem 3:  Con($\frak c=\omega_2$ or more, $\clubsuit$
%  + $\cov\Cal N_{\omega} = \frak c$ with a Sierpinski set?
% now $\stick$ already gives $\cov\Cal N_{\omega_1}=\omega_1$

%Theorem 5.8??
% consistency of  \cov\Cal N=\cov\Cal N_{\omega_1}=\frak c  large
%    and          \cov\Cal N_{\omega_2}=\omega_2

\leader{523G}{}\cmmnt{ When the additivity of Lebesgue measure is
large we have a further constraint on covering numbers.

\medskip

\noindent}{\bf Proposition}\cmmnt{ ({\smc Kraszewski 01})} If
$\kappa$ is a cardinal and $\cov\Cal N_{\kappa}<\add\Cal N$,
then $\cov\Cal N_{\kappa}\le\cff[\kappa]^{\le\omega}$.

\proof{ As $\{0,1\}^{\kappa}$ is covered by negligible sets, $\kappa$ is
infinite.   Let $\Cal E$ be a subset of $\Cal N_{\kappa}$ with cardinal
$\cov\Cal N_{\kappa}$ and union $\{0,1\}^{\kappa}$, and $\Cal J$ a
cofinal subset of $[\kappa]^{\omega}$ with cardinal
$\cff[\kappa]^{\le\omega}$.   For $J\in\Cal J$ and $x\in\{0,1\}^{\kappa}$
set $\pi_J(x)=x\restr J$, so that $\pi_J:\{0,1\}^{\kappa}\to\{0,1\}^J$
is \imp.   For $J\in\Cal J$ set
$\Cal E_J=\{E:E\in\Cal E$, $\pi_J[E]\in\Cal N_J\}$,
$H_J=\bigcup\Cal E_J$.   Since

\Centerline{$\#(\Cal E_J)\le\#(\Cal E)
=\cov\Cal N_{\kappa}<\add\Cal N=\add\Cal N_{\omega}=\add\Cal N_J$,}

\noindent $F_J=\bigcup\{\pi_J[E]:E\in\Cal E_J\}\in\Cal N_J$ and
$H_J\subseteq\pi_J^{-1}[F_J]\in\Cal N_{\kappa}$.   Since
$\bigcup_{J\in\Cal J}\Cal E_J=\Cal E$ (523C) covers $\{0,1\}^{\kappa}$,
$\{H_J:J\in\Cal J\}$ covers $\{0,1\}^{\kappa}$ and
$\cov\Cal N_{\kappa}\le\#(\Cal J)=\cff[\kappa]^{\le\omega}$.
}%end of proof of 523G

%Question:  can we have  \add\Cal N_{\omega}>\omega_1  but  Stick ?

\leader{523H}{\dvrocolon{Uniformities}}\cmmnt{ On the other side of the
diagram we have non-decreasing functions.   To get upper bounds for
$\non\Cal N_{\kappa}$ we have the following method.

\medskip

\noindent}{\bf Lemma}\cmmnt{ ({\smc Kraszewski 01})} Suppose that
$I$ is a set and $F$ a family of functions
with domain $I$ such that for every countable $J\subseteq I$ there
is an $f\in F$ such that $f\restr J$ is injective.   Then

\Centerline{$\non\Cal N_I
\le\max(\#(F),\sup_{f\in F}\non\Cal N_{f[I]})$.}

\proof{ If $I$ is finite the result is trivial.
Otherwise, for each $f\in F$ take a non-negligible subset $A_f$ of
$\{0,1\}^{f[I]}$ with cardinal $\non\Cal N_{f[I]}$.   Set
$A=\{yf:f\in F,\,y\in A_f\}\subseteq\{0,1\}^I$.   \Quer\ If
$A\in\Cal N_I$, there is a countable set $J\subseteq I$
such that $\{x\restr J:x\in A\}\in\Cal N_J$.   Let $f\in F$ be such that
$f\restr J$ is injective.   Then we have a
function $\phi:\{0,1\}^{f[I]}\to\{0,1\}^J$ defined by saying that
$\phi(z)=zf\restr J$ for every $z\in\{0,1\}^{f[I]}$, and (because
$f\restr J$ is injective) $\phi$ is
\imp\ for $\nu_{f[I]}$ and $\nu_J$, so $\phi[A_f]$ cannot be
$\nu_J$-negligible.   But if $y\in A_f$ then
$\phi(y)(\xi)=y(f(\xi))$ for every $\xi\in J$, so
$\phi[A_f]\subseteq\{x\restr J:x\in A\}$, which is supposed to be
negligible.\ \Bang

Thus $A$ is not negligible, and

\Centerline{$\non\Cal N_I\le\#(A)
\le\max(\omega,\#(F),\sup_{f\in F}\#(A_f))
=\max(\#(F),\sup_{f\in F}\non\Cal N_{f[I]})$}

\noindent because we are supposing that $I$ is infinite, so there is
some $f\in F$ such that $f[I]$ is infinite.
}%end of proof of 523H

\leader{523I}{Theorem} (a) For any cardinal $\kappa$,

\quad(i) $\non\Cal N_{\kappa}
  \le\max(\non\Cal N,\cff[\kappa]^{\le\omega})$,

\quad(ii)\dvAnew{2024} 
$\non\Cal N_{\kappa^+}\le\max(\kappa^+,\non\Cal N_{\kappa})$,

\quad(iii) $\non\Cal N_{2^{\kappa}}
\le\max(\frak c,\cff[\kappa]^{\le\omega})$,

%\def\tinyplus{{\vrule height 1.75pt depth -1.45pt width 2.5pt}
%\hskip-1.3pt{\vrule height 2.5pt depth 0pt width 0.3pt}}

\quad(iv)
$\non\Cal N_{2^{\kappa^{\tinyplus}}}
\le\max(\kappa^+,\non\Cal N_{2^{\kappa}})$.

(b) If $\cf\kappa>\omega$, then
$\non\Cal N_{\kappa^+}
\le\max(\cf\kappa,\sup_{\lambda<\kappa}\non\Cal N_{\lambda})$.

\proof{{\bf (a)} If $\kappa$ is finite, all these are trivial;  so suppose
otherwise.

\medskip

\quad{\bf (i)} Let $\Cal J\subseteq[\kappa]^{\le\omega}$ be a cofinal set
with cardinal $\cff[\kappa]^{\le\omega}$, and for $J\in\Cal J$ let
$f_J$ be the identity function on $J$.   Applying 523H with
$F=\{f_J:J\in\Cal J\}$ we get

$$\eqalign{\non\Cal N_{\kappa}
&\le\max(\#(\Cal J),\sup_{J\in\Cal J}\non\Cal N_J)\cr
&\le\max(\non\Cal N_{\omega},\cff[\kappa]^{\le\omega})
=\max(\non\Cal N,\cff[\kappa]^{\le\omega}).\cr}$$

\medskip

\quad{\bf (ii)} For each $\xi<\kappa^+$ choose
a function $f_{\xi}:\kappa^+\to\kappa$ which is injective on $\xi$,
and set $F=\{f_{\xi}:\xi<\kappa^+\}$.   By 523H,

\Centerline{$\non\Cal N_{\kappa^+}
\le\max(\#(F),\sup_{\xi<\kappa^+}\non\Cal N_{f_{\xi}[\kappa^+]})
\le\max(\kappa^+,\non\Cal N_{\kappa})$.}

\medskip

\quad{\bf (iii)} Take $\Cal J$ as in (i).
This time, for $J\in\Cal J$, define
$f_J:\Cal P\kappa\to\Cal PJ$ by setting $f_J(A)=A\cap J$ for every
$A\subseteq\kappa$.   Applying 523H with $F=\{f_J:J\in\Cal J\}$ we get

$$\eqalign{\non\Cal N_{2^{\kappa}}
&=\non\Cal N_{\Cal P\kappa}
\le\max(\#(\Cal J),\sup_{J\in\Cal J}\non\Cal N_{\Cal PJ})
\le\max(\cff[\kappa]^{\le\omega},\non\Cal N_{\frak c})\cr
&\le\max(\cff[\kappa]^{\le\omega},\non\Cal N,\cff[\frak c]^{\le\omega})
=\max(\cff[\kappa]^{\le\omega},\frak c)\cr}$$

\noindent (5A1F(e-iii)).

\medskip

\quad{\bf (iv)} Set $f_{\xi}(A)=A\cap\xi$ for $\xi<\kappa^+$ and
$A\subseteq\kappa^+$.   If
$\Cal J\subseteq\Cal P\kappa^+$ is countable, there is a $\xi<\kappa^+$
such that $A\cap\xi\ne A'\cap\xi$ for all distinct $A$, $A'\in\Cal J$,
that is, $f_{\xi}\restr\Cal J$ is injective.   So 523H tells us that

$$\eqalign{\non\Cal N_{2^{\kappa^{\tinyplus}}}
&=\non\Cal N_{\Cal P(\kappa^+)}
\le\max(\kappa^+,\sup_{\xi<\kappa^+}\non\Cal N_{f_{\xi}[\kappa^+]})\cr
&\le\max(\kappa^+,\sup_{\xi<\kappa^+}\non\Cal N_{\Cal P\xi})
=\max(\kappa^+,\non\Cal N_{2^{\kappa}}).\cr}$$

\medskip

{\bf (b)(i)} If $\kappa=\theta^+$ where $\theta$ is an infinite cardinal, 
$\non\Cal N_{\kappa^+}\le\max(\kappa,\non\Cal N_{\theta})$.   \Prf\
Choose an injective function $h_{\zeta}:\zeta\to\kappa$ for each
$\zeta<\kappa^+$.   For $\xi<\kappa$ define
$f_{\xi}:\kappa^+\to\kappa$ by saying that

\Centerline{$f_{\xi}(\zeta)
=\min(\kappa\setminus\{f_{\xi}(\eta):
  \eta<\zeta$, $h_{\zeta}(\eta)\le\xi\})$}

\noindent for $\zeta<\kappa^+$.   If
$J\subseteq\kappa^+$ is countable, then
$\xi=\sup_{\eta,\zeta\in J,\eta<\zeta}h_{\zeta}(\eta)$ is less than
$\kappa$, and
$f_{\xi}(\eta)\ne f_{\xi}(\zeta)$ for all distinct $\eta$, $\zeta\in J$.
Applying 523H with $F=\{f_{\xi}:\xi<\kappa\}$ and using (a-ii) above, we get

\Centerline{$\non\Cal N_{\kappa^+}
\le\max(\kappa,\non\Cal N_{\kappa})
\le\max(\kappa,\non\Cal N_{\theta})
=\max(\cf\kappa,\sup_{\lambda<\kappa}\non\Cal N_{\lambda})$.  \Qed}

\medskip

\quad{\bf (ii)} Now suppose that $\kappa$ is an uncountable limit cardinal
with uncountable cofinality.   Again choose an injective function
$h_{\zeta}:\zeta\to\kappa$ for each $\zeta<\kappa^+$.
This time, let $K\subseteq\kappa$ be a cofinal set with cardinal
$\cf\kappa$ consisting of cardinals, and for $\lambda\in K$ define
$f_{\lambda}:\kappa^+\to\lambda^+$ by the formula

\Centerline{$f_{\lambda}(\zeta)
=\min(\lambda^+\setminus\{f_{\lambda}(\eta):
  \eta<\zeta$, $h_{\zeta}(\eta)\le\lambda\})$}

\noindent for $\zeta<\kappa^+$.   If
$J\subseteq\kappa^+$ is countable, then there is a $\lambda\in K$ such that
$\lambda\ge\sup_{\eta,\zeta\in J,\eta<\zeta}h_{\zeta}(\eta)$, and
$f_{\lambda}(\eta)\ne f_{\lambda}(\zeta)$ for all distinct
$\eta$, $\zeta\in J$.
Applying 523H with $F=\{f_{\lambda}:\lambda\in K\}$, we get

\Centerline{$\non\Cal N_{\kappa^+}
\le\max(\#(F),\sup_{f\in F}\non\Cal N_{f[\kappa^+]})
\le\max(\cf\kappa,\sup_{\lambda<\kappa}\non\Cal N_{\lambda})$.}
}%end of proof of 523I

\leader{523J}{Corollary}\cmmnt{ ({\smc Kraszewski 01})}
(a) $\non\Cal N_{\omega_2}=\non\Cal N_{\omega_1}=\non\Cal N$.

(b) For any $n\in\Bbb N$,
$\non\Cal N_{\omega_{n+1}}\le\max(\omega_n,\non\Cal N)$.

(c) $\non\Cal N_{2^{\omega_1}}=\non\Cal N_{\frakc}$.

(d) If $n\in\Bbb N$ then
$\non\Cal N_{2^{\omega_n}}\le\max(\omega_n,\non\Cal N_{\frakc})$.

\proof{{\bf (a)} We have

$$\eqalignno{\non\Cal N&=\non\Cal N_{\omega}
\le\non\Cal N_{\omega_1}\le\non\Cal N_{\omega_2}\cr
\displaycause{523B}
&\le\max(\cf\omega_1,\non\Cal N_{\omega})\cr
\displaycause{523Ib}
&=\non\Cal N.\cr}$$

\medskip

{\bf (b)} Induce on $n$, using 523Ib for the inductive step.

\medskip

{\bf (c)} By 523I(a-iii),
$\non\Cal N_{2^{\omega_1}}\le\max(\omega_1,\non\Cal N_{\frakc})$;
since

\Centerline{$\omega_1\le\non\Cal N_{\frakc}
\le\non\Cal N_{2^{\omega_1}}$,}

\noindent we have the result.

\medskip

{\bf (d)} Induce on $n$, using 523I(a-iii) or 523I(a-iv) for the inductive step.
}%end of proof of 523J

\leader{523K}{Corollary}\cmmnt{ ({\smc Burke n05})}
For any sets $I$, $K$ let
$\Upsilon_{\omega}(I,K)$ be the least cardinal of any family $F$ of
functions from $I$ to $K$ such that for every countable $J\subseteq I$
there is an $f\in F$ which is injective on $J$.   (If
$\#(K)<\min(\omega,\#(I))$ take $\Upsilon_{\omega}(I,K)=\infty$.)   Then

(a) $\non\Cal N_I\le\max(\Upsilon_{\omega}(I,K),\non\Cal N_K)$ for all sets
$I$ and $K$;

(b) if $\kappa\ge\frak c$ is a cardinal, then
$\non\Cal N_{\kappa}
=\max(\Upsilon_{\omega}(\kappa,\frak c),\non\Cal N_{\frak c})$.

\proof{{\bf (a)} This is just a slightly weaker version of 523H.

\medskip

{\bf (b)} The point is that
$\Upsilon_{\omega}(\kappa,\frak c)\le\non\Cal N_{\kappa}$.   \Prf\
Let $A\subseteq\{0,1\}^{\kappa\times\omega}$ be a non-negligible set of
cardinal $\non\Cal N_{\kappa\times\omega}$.
For $x\in\{0,1\}^{\kappa\times\omega}$ define
$f_x:\kappa\to\{0,1\}^{\omega}$ by
setting $f_x(\xi)=\sequencen{x(\xi,n)}$ for each $\xi<\kappa$.
If $\xi$, $\eta<\kappa$ are distinct, then $\{x:f_x(\xi)=f_x(\eta)\}$ is
negligible, so if $J\subseteq\kappa$ is countable then
$\{x:f_x\restr J$ is injective$\}$ is conegligible and meets $A$.
Accordingly $\{f_x:x\in A\}$ witnesses that

\Centerline{$\Upsilon_{\omega}(\kappa,\frak c)
=\Upsilon_{\omega}(\kappa,\{0,1\}^{\omega})
\le\#(A)=\non\Cal N_{\kappa\times\omega}
=\non\Cal N_{\kappa}$.  \Qed}

Since we already know that

\Centerline{$\non\Cal N_{\frak c}\le\non\Cal N_{\kappa}
\le\max(\Upsilon_{\omega}(\kappa,\frak c),\non\Cal N_{\frak c})$,}

\noindent we have the result.
}%end of proof of 523K

\leader{523L}{}\cmmnt{ On the other side we can find lower bounds which
give a notion of the rate of growth of the numbers $\non\Cal N_{\kappa}$
as $\kappa$ increases.

\medskip

\noindent}{\bf Proposition} (a) If $\lambda$ and $\kappa$ are
infinite cardinals with $\kappa>2^{\lambda}$, then
$\non\Cal N_{\kappa}>\lambda$.
%old 523Gc

(b) If $\kappa$ is a strong limit cardinal of countable
cofinality then $\non\Cal N_{\kappa}>\kappa$.
%old 523Gd

\proof{{\bf (a)} Let
$A\subseteq\{0,1\}^{\kappa}$ be any set with cardinal at most $\lambda$.   For
$\xi<\kappa$ set $B_{\xi}=\{x:x\in A,\,x(\xi)=1\}$.   Because
$\kappa>2^{\#(A)}$, there is some $B\subseteq A$ such that
$I=\{\xi:B_{\xi}=B\}$ is infinite.    But what this means is that if
$\xi\in I$ then $x(\xi)=1$ for every $\xi\in B$ and $x(\xi)=0$ for every
$x\in A\setminus B$, and
$A\subseteq\{x:x$ is constant on $I\}$ is negligible.   As $A$ is
arbitrary, $\non\Cal N_{\kappa}>\lambda$.

\medskip

{\bf (b)} By (a),
$\non\Cal N_{\kappa}>\lambda$ for every $\lambda<\kappa$, so
$\non\Cal N_{\kappa}\ge\kappa$;  but also
$\cf(\non\Cal N_{\kappa})\ge\add\Cal N_{\kappa}$ (513C(b-ii)), so
$\non\Cal N_{\kappa}$ has uncountable cofinality and must be greater
than $\kappa$.
}%end of proof of 523L

\leader{523M}{Shrinking \dvrocolon{numbers}}\cmmnt{ As with
$\non\Cal N_{\ssbullet}$, the functions
$\kappa\mapsto\shr\Cal N_{\kappa}$ and $\kappa\mapsto\shr\Cal N_{\kappa}$
are non-decreasing, by 521Hb.
Some of the ideas used in 523I can be adapted to this context, but the
pattern as a whole is rather different.

\medskip

\noindent}{\bf Proposition} (a)(i) For any non-zero
cardinals $\kappa$ and $\lambda$,

\Centerline{$\shr\Cal N_{\kappa}
\le\max(\covSh(\kappa,\lambda,\omega_1,2),
  \sup_{\theta<\lambda}\shr\Cal N_{\theta})$.}

\quad(ii) For any infinite cardinal $\kappa$,
$\shr\Cal N_{\kappa}
\le\max(\shr\Cal N,\cff[\kappa]^{\le\omega})$.

\quad(iii) If $\cf\kappa>\omega$, then
$\shr\Cal N_{\kappa}
\le\max(\kappa,\sup_{\theta<\kappa}\shr\Cal N_{\theta})$.

(b) For any infinite cardinal $\kappa$,

\quad(i) $\shr\Cal N_{\kappa}\ge\kappa$;

\quad(ii) $\cf(\shr\Cal N_{\kappa})>\omega$;

\quad(iii) $\cf(\shr^+\Cal N_{\kappa})>\kappa$.

\cmmnt{\medskip

\noindent{\bf Remark} For the definition of $\covSh$, see 5A2Da.}

\proof{{\bf (a)(i)}
If $\covSh(\kappa,\lambda,\omega_1,2)=\infty$ or $\kappa$ is finite
this is trivial.   Otherwise, $\lambda\ge\omega_1$.
Take a non-negligible $A\subseteq\{0,1\}^{\kappa}$.
Let $\Cal J\subseteq[\kappa]^{<\lambda}$ be a
set with cardinal $\covSh(\kappa,\lambda,\omega_1,2)$ such that for every
$I\in[\kappa]^{<\omega_1}$ there is a $\Cal D\in[\Cal J]^{<2}$ such that
$I\subseteq\bigcup\Cal D$, that is, there is a $J\in\Cal J$ such that
$I\subseteq J$.   For each $J\in\Cal J$, $A_J=\{x\restr J:x\in A\}$ is
non-negligible;  let $B_J\subseteq A_J$ be a non-negligible
set with cardinal at most $\shr\Cal N_J$.   Let $B\subseteq A$ be a set with cardinal
at most $\max(\omega,\#(\Cal J),\sup_{J\in\Cal J}\shr\Cal N_J)$ such that
$B_J\subseteq\{x\restr J:x\in B\}$ for every $J\in\Cal J$.   If
$I\subseteq\kappa$ is countable, there is a $J\in\Cal J$ such that
$I\subseteq J$, so $\{x\restr I:x\in B\}\supseteq\{y\restr I:y\in B_J\}$ is
non-negligible;  it follows that $B$ is non-negligible, while
$\#(B)\le\max(\covSh(\kappa,\lambda,\omega_1,2),
  \sup_{\theta<\lambda}\shr\Cal N_{\theta})$.

\medskip

\quad{\bf (ii)} Taking $\lambda=\omega_1$ in (i),

\Centerline{$\shr\Cal N_{\kappa}
\le\max(\covSh(\kappa,\omega_1,\omega_1,2),\shr\Cal N_{\omega})
=\max(\cff[\kappa]^{\le\omega},\shr\Cal N)$.}

\medskip

\quad{\bf(iii)} Take $\lambda=\kappa$ in (i);  as
$[\kappa]^{\le\omega}=\bigcup_{\xi<\kappa}[\xi]^{\le\omega}$,

\Centerline{$\shr\Cal N_{\kappa}
\le\max(\covSh(\kappa,\kappa,\omega_1,2),
\sup_{\theta<\kappa}\shr\Cal N_{\theta})
=\max(\kappa,\sup_{\theta<\kappa}\shr\Cal N_{\theta})$.}

\medskip

{\bf (b)(i)} Induce on $\kappa$.   If $\kappa=\omega$
the result is trivial.   For the inductive step to $\kappa^+$, consider the
set

\Centerline{$A
=\{x:x\in\{0,1\}^{\kappa^+},\,\Exists\xi<\kappa^+,\,x(\eta)=0$ for every
$\eta\ge\xi\}$.}

\noindent Then the only set which includes $A$ and is determined by coordinates
in a countable set is $\{0,1\}^{\kappa^+}$, so $A$ has full
outer measure.   On the other hand, if $B\subseteq A$ and
$\#(B)\le\kappa$, then there is some $\zeta<\kappa^+$ such that
$x(\xi)=0$ for every $x\in B$ and every $\xi\ge\zeta$, so $B$ is
negligible.   Thus $A$ witnesses that
$\shr\Cal N_{\kappa^+}\ge\kappa^+$.   Because
$\kappa\mapsto\shr\Cal N_{\kappa}$ is non-decreasing (523B),
the inductive step
to limit cardinals $\kappa$ is trivial.

\medskip

\quad{\bf (ii)} \Quer\ Now suppose, if possible, that
$\cf(\shr\Cal N_{\kappa})=\omega$.
Then there is a sequence $\sequencen{\lambda_n}$ of cardinals
less than $\shr\Cal N_{\kappa}$ with supremum $\shr\Cal N_{\kappa}$.
For each $n\in\Bbb N$ set $I_n=\kappa\times\{n\}$, and let
$A_n\subseteq\{0,1\}^{I_n}$ be a
non-negligible set such that every non-negligible
subset of $A_n$ has more than $\lambda_n$ members.
By 523Dc, there is a set $B_n\subseteq\{0,1\}^{I_n}$ of full outer
measure such that every non-negligible
subset of $B_n$ has more than $\lambda_n$ members.   Set

\Centerline{$B
=\{x:x\in\{0,1\}^{\kappa\times\Bbb N}$, $x\restr I_n\in B_n$
for every $n\in\Bbb N\}$.}

\noindent Then the natural isomorphism between
$\{0,1\}^{\kappa\times\Bbb N}$ and $\prod_{n\in\Bbb N}\{0,1\}^{I_n}$
identifies $B$ with $\prod_{n\in\Bbb N}B_n$, so $B$ has full outer
measure in $\{0,1\}^{\kappa\times\Bbb N}$ (254Lb).   There must therefore
be a set $C\subseteq B$, of non-zero measure, such that
$\#(C)\le\shr\Cal N_{\kappa}$.   Express $C$ as
$\bigcup_{n\in\Bbb N}C_n$ where $\#(C_n)\le\lambda_n$ for every
$n$.   Then there is an $n\in\Bbb N$ such that $C_n$ is not negligible,
in which case $D_n=\{x\restr I_n:x\in C_n\}$ is non-negligible.   But
$D_n\subseteq B_n$ and $\#(D_n)\le\lambda_n$, so this is impossible.\ \Bang

\medskip

\quad{\bf (iii)} The argument of (i) shows that if $\kappa$ is a successor
cardinal, then $\shr^+\Cal N_{\kappa}>\kappa$.   So we need consider only
the case in which $\kappa$ is a limit cardinal.
\Quer\ If $\cf(\shr^+\Cal N_{\kappa})\le\kappa$, then
there is a family $\ofamily{\xi}{\kappa}{\lambda_{\xi}}$ of cardinals
less than $\shr^+\Cal N_{\kappa}$ with supremum $\shr^+\Cal N_{\kappa}$.
I use the same method as in (ii).
For each $\xi<\kappa$ set $I_{\xi}=\kappa\times\{\xi\}$, and let
$B_{\xi}\subseteq\{0,1\}^{I_{\xi}}$ be a
set of full outer measure such that every non-negligible
subset of $B_{\xi}$ has at least $\lambda_{\xi}$ members.  Set

\Centerline{$B
=\{x:x\in\{0,1\}^{\kappa\times\kappa}$, $x\restr I_{\xi}\in B_{\xi}$
for every $\xi<\kappa$.}

\noindent Then $B$ has full outer
measure in $\{0,1\}^{\kappa\times\kappa}$.   There must therefore
be a set $C\subseteq B$, of non-zero measure, such that
$\#(C)<\shr^+\Cal N_{\kappa}$.   Let $\xi<\kappa$ be such that
$\#(C)<\lambda_{\xi}$.   Then
$D=\{x\restr I_{\xi}:x\in C\}$ is non-negligible.   But
$D\subseteq B_{\xi}$ and $\#(D_{\xi})<\lambda_{\xi}$, so this is
impossible.\ \Bang

}%end of proof of 523M

\leader{523N}{\dvrocolon{Cofinalities}}\cmmnt{ For the cardinals
$\cf\Cal N_{\kappa}$ the pattern from 523I(a-i) and 523Mb continues, and
indeed we have an exact formula.

\medskip

\noindent}{\bf Theorem} For any infinite cardinal $\kappa$,

\Centerline{$\kappa\le\cf\Cal N_{\kappa}
=\max(\cf\Cal N,\cff[\kappa]^{\le\omega})
\le\kappa^{\omega}$.}

\proof{{\bf (a)} $\cf\Cal N_{\kappa}
\le\max(\cf\Cal N,\cff[\kappa]^{\le\omega})$.
\Prf\ Let $\Cal J$ be a
cofinal family in $[\kappa]^{\omega}$ with cardinal
$\cff[\kappa]^{\le\omega}$.   For each $J\in\Cal J$,
write $\pi_J(x)=x\restr J$ for $x\in\{0,1\}^{\kappa}$.   Let $\Cal E_J$
be a cofinal subset of
$\Cal N_J$ with cardinal $\cf\Cal N_J=\cf\Cal N_{\omega}=\cf\Cal N$.
Consider
$\Cal E=\{\pi_J^{-1}[E]:J\in\Cal J,\,E\in\Cal E_J\}$.   By 523C,
$\Cal E$ is cofinal with $\Cal N_{\kappa}$, so that

\Centerline{$\cf\Cal N_{\kappa}\le\#(\Cal E)
\le\max(\cf\Cal N,\cff[\kappa]^{\le\omega})$.
\Qed}

\medskip

{\bf (b)} We know that
$\cff[\kappa]^{\le\omega}\le\cf\Cal N_{\kappa}$ (521Jb) and that
$\cf\Cal N=\cf\Cal N_{\omega}\le\cf\Cal N_{\kappa}$ (523B).   So
$\cf\Cal N_{\kappa}=\max(\cf\Cal N,\cff[\kappa]^{\le\omega})$.

\medskip

{\bf (c)} For the inequalities, note that $\omega\le\cf\Cal N$ and
if $\kappa$ is uncountable then (in the language of 512Ba)

\Centerline{$\cff[\kappa]^{\le\omega}
\ge\cov(\kappa,\in,[\kappa]^{\le\omega})=\kappa$.}

\noindent On the other side,
$\cf\Cal N\le\frak c\le\kappa^{\omega}$ and
$\cff[\kappa]^{\le\omega}\le\#([\kappa]^{\le\omega})\le\kappa^{\omega}$.
}%end of proof of 523N

\cmmnt{
\leader{523O}{Cofinalities of the cardinals} In 523Mb I have
shown that $\shr\Cal N_{\kappa}$ has uncountable cofinality for infinite
$\kappa$, and rather more about $\shr^+\Cal N_{\kappa}$.
From 513Cb we have a little information concerning the cofinalities of
$\add\Cal N_{\kappa}$, $\cov\Cal N_{\kappa}$, $\non\Cal N_{\kappa}$ and
$\cf\Cal N_{\kappa}$;  but except when $\kappa=\omega$ we learn only
that $\cf\Cal N_{\kappa}$ and $\non\Cal N_{\kappa}$ have uncountable
cofinality, and that if $\cov\Cal N_{\kappa}=\cf\Cal N_{\kappa}$ then
their common cofinality is at least $\non\Cal N_{\kappa}$.   This last
remark can apply only to `small' $\kappa$, since
$\cf\Cal N_{\kappa}\ge\kappa$ (if $\kappa$ is infinite) and
$\cov\Cal N_{\kappa}\le\cov\Cal N$.}

\leader{523P}{The generalized continuum \dvrocolon{hypothesis}}\cmmnt{
In this chapter I am trying to present arguments in forms which show
their full strength and are not tied to particular axioms beyond those
of ZFC.   However it is perhaps worth mentioning that in one of the
standard universes the pattern is particularly simple.

\medskip

\noindent}{\bf Proposition} Suppose that the generalized continuum
hypothesis is true.   Then, for any infinite cardinal $\kappa$,

$$\eqalign{\add\Cal N_{\kappa}=\add\nu_{\kappa}
=\cov\Cal N_{\kappa}=\omega_1;\cr}$$

$$\eqalign{\non\Cal N_{\kappa}&=\lambda
  \text{ if }\kappa=\lambda^+\text{ where }\cf\lambda>\omega,\cr
&=\kappa^+\text{ if }\cf\kappa=\omega,\cr
&=\kappa\text{ otherwise};\cr}$$

$$\eqalign{\shr\Cal N_{\kappa}=\cf\Cal N_{\kappa}
&=\kappa^+\text{ if }\cf\kappa=\omega,\cr
&=\kappa\text{ otherwise};\cr}$$

$$\eqalign{\shr^+\Cal N_{\kappa}=(\shr\Cal N_{\kappa})^+
&=\kappa^{++}\text{ if }\cf\kappa=\omega,\cr
&=\kappa^+\text{ otherwise}.\cr}$$

\proof{ Since

\Centerline{$\omega_1\le\add\Cal N_{\kappa}=\add\nu_{\kappa}
\le\cov\Cal N_{\kappa}\le\cov\Cal N\le\frak c=\omega_1$,}

\noindent the additivity and covering number are always
$\omega_1$.

If $\kappa=\lambda^+$ where
$\cf\lambda>\omega$, then $\kappa>2^{\theta}$ for every
$\theta<\lambda$, so we have

\Centerline{$\lambda\le\non\Cal N_{\kappa}
\le\max(\frak c,\cff[\lambda]^{\le\omega})=\lambda$}

\noindent (523La, 523I(a-iii), 5A6Ab).   If $\kappa=\lambda^+$ where
$\cf\lambda=\omega$, then

\Centerline{$\lambda\le\non\Cal N_{\kappa}
\le\lambda^{\omega}\le 2^{\lambda}=\kappa$;}

\noindent but as $\non\Cal N_{\kappa}$ has uncountable cofinality
(513C(b-ii) again), $\non\Cal N_{\kappa}$ must be $\kappa$.
If $\kappa$ is
a limit cardinal, then $\kappa>2^{\theta}$ for every $\theta<\kappa$, so

\Centerline{$\kappa\le\non\Cal N_{\kappa}
\le\max(\omega_1,\cff[\kappa]^{\le\omega})$}

\noindent by 523I(a-i);  
if $\cf\kappa>\omega$ this is already enough to show that
$\non\Cal N_{\kappa}=\kappa$;  if $\cf\kappa=\omega$ then
$\non\Cal N_{\kappa}$ cannot be $\kappa$ so must be
$\kappa^+=\kappa^{\omega}$.

As for $\shr\Cal N_{\kappa}$, if $\cf\kappa=\omega$, then

\Centerline{$\kappa^+\le\shr\Cal N_{\kappa}
\le\cf\Cal N_{\kappa}=\max(\omega_1,\cff[\kappa]^{\le\omega})
\le 2^{\kappa}=\kappa^+$}

\noindent by 523M(b-ii) and 523N.   If $\cf\kappa>\omega$ then

\Centerline{$\kappa\le\shr\Cal N_{\kappa}\le\cf\Cal N_{\kappa}\le\kappa$}

\noindent by 523M(b-i), 523N and 5A6Ab.   This deals with
$\shr\Cal N_{\kappa}$ and $\cf\Cal N_{\kappa}$.   For the augmented
shrinking numbers, we know that if $\cf\kappa=\omega$ then
$\shr\Cal N_{\kappa}=\kappa^+$ is a successor cardinal so
$\shr^+\Cal N_{\kappa}=(\shr\Cal N_{\kappa})^+=\kappa^{++}$, while if
$\cf\kappa>\omega$ then 

$$\eqalignno{\shr\Cal N_{\kappa}&=\kappa<\shr^+\Cal N_{\kappa}
\Displaycause{523M(b-iii)}
\le(\shr\Cal N_{\kappa})^+\cr}$$

\noindent and $\shr^+\Cal N_{\kappa}=(\shr\Cal N_{\kappa})^+=\kappa^+$.
}%end of proof of 523P

\exercises{\leader{523X}{Basic exercises (a)}
%\spheader 523Xa
Show that

\Centerline{$(\Cal N_{\kappa},\not\ni,\{0,1\}^{\kappa})
\prGT([\kappa]^{\le\omega},\subseteq,[\kappa]^{\le\omega})
  \ltimes(\Cal N,\not\ni,\Bbb R)$}

\noindent for every infinite cardinal $\kappa$.   (See 512I for the
definition of $\ltimes$.)   Use this to prove 523I(a-i).
%523I mt52bits 

\spheader 523Xb Let $\kappa$ be an infinite cardinal, and $\Cal J$ a
family of subsets of $\kappa$ such that every countable subset of
$\kappa$ is included in some member of $\Cal J$.   Show that
$\non\Cal N_{\kappa}
\le\max(\#(\Cal J),\sup_{J\in\Cal J}\non\Cal N_J)$,
$\non\Cal N_{\Cal P\kappa}
\le\max(\#(\Cal J),\sup_{J\in\Cal J}\non\Cal N_{\Cal PJ})$,
$\shr\Cal N_{\kappa}
\le\max(\#(\Cal J),\penalty-100\sup_{J\in\Cal J}\shr\Cal N_J)$ and
$\cf\Cal N_{\kappa}
\le\max(\#(\Cal J),\penalty-100\sup_{J\in\Cal J}\cf\Cal N_J)$.
%523I(a-i) 523I(a-iii) 523M(a-ii) 523N

\spheader 523Xc Show that

\Centerline{$(\Cal N_{\kappa},\subseteq,\Cal N_{\kappa})
\prGT([\kappa]^{\le\omega},\subseteq,[\kappa]^{\le\omega})
\ltimes(\Cal N,\subseteq,\Cal N)$}

\noindent for every infinite cardinal $\kappa$.   Use this to prove
523N.
%523N  523Xa

\spheader 523Xd Let $(X,\Sigma,\mu)$ be any probability space, and for each
set $I$ write $\Cal N(\mu^I)$ for the null ideal of the product measure
on $X^I$.   Show that all the results of
523E-523I %523E 523F 523G 523H 523I
and 523L-523N %523L 523M 523N
are valid with
$\Cal N(\mu^I)$ in place of $\Cal N_I$ and $\Cal N(\mu^{\omega})$
in place of $\Cal N$, except that

----- in 523E the additivities may stabilize at $\infty$ rather than $\omega_1$;

----- in 523F we can no longer be sure that
$\cov\Cal N(\mu^{\omega})\le\frak c$;

----- in 523I(a-iii) we need to write
`$\non\Cal N(\mu^{2^{\kappa}})
\le\max(\frak c,\non\Cal N(\mu^{\omega}),\cff[\kappa]^{\le\omega})$';

----- in 523L and 523Mb
we have to assume that the measure algebra of $\mu$ is not
$\{0,1\}$, so that the product measure $\mu^{\Bbb N}$ is atomless;

%\query or possibly enough if  X  not a singleton?

----- in 523N we can no longer be sure that
$\cf\Cal N(\mu^{\omega})\le\kappa^{\omega}$.
%523N


\leader{523Y}{Further exercises (a)}
%\spheader 523Ya
Set $\frak A=\Cal P\Bbb R/\Cal N$.   %(i) 
Show that
$\frak c\le c(\frak A)\le\pi(\frak A)\le 2^{\shr\Cal N}$.
%(ii) Show that if
%$\shr^+\Cal N\ge\frak c$ and $2^{\lambda}\le\frak c$ for every
%$\lambda<\frak c$, then $c(\frak A)=2^{\frak c}$.  query
%523M

\spheader 523Yb Let $\kappa$ be an infinite cardinal.   Show that there is
a family $\Cal J\subseteq[\kappa]^{\le\omega}$ such that
$\#(\Cal J)\le\shr\Cal N_{\kappa}$ and every infinite subset of $\kappa$
meets some member of $\Cal J$ in an infinite set.
%523M mt52bits

\spheader 523Yc Suppose that $\kappa\ge\omega$ and that
$[\kappa]^{\le\omega}$ has bursting number at most $\add\Cal N$.   Show
that $\Cal N_{\kappa}\equivT[\kappa]^{\le\omega}\times\Cal N$.
%{\smc Fremlin 91, 2J.}
%523N   521Jb 

\spheader 523Yd Show that

\Centerline{$(\omega_1,\le,\omega_1)\ltimes(\omega_1,\le,\omega_1)
\not\preccurlyeq_{\text{GT}}
(\omega_1,\le,\omega_1)\times(\omega_1,\le,\omega_1)$.}
%523Xc 523Yc 523N mt52bits

\spheader 523Ye For infinite cardinals $\kappa$, write $\Cal M_{\kappa}$
for the ideal of meager subsets of $\{0,1\}^{\kappa}$.   Show that under
the same conventions as in 522B and 523B we have the diagrams

$$\vbox{\offinterlineskip
\halign{\hfil#\hfil&\hfil#\hfil&\hfil#\hfil&\hfil#\hfil
  &\hfil#\hfil&\hfil#\hfil&\hfil#\hfil
  &\hfil#\hfil&\hfil#\hfil&\hfil#\hfil
  &\hfil#\hfil&\hfil#\hfil&\hfil#\hfil&\hfil#\hfil\cr
&\tmpstrut&$\cov\Cal M_{\lambda}$&\tmphrule
  &$\cov\Cal M_{\kappa}$&\tmphrule
  &$\cf\Cal M_{\kappa}$&\tmphrule&$\cf\Cal M_{\lambda}$
  &\tmphrule&$\lambda^{\omega}$\cr
&\tmpstrut&\vrule&&\vrule&&\vrule&&\vrule\cr
&\tmpstrut&\vrule&&\vrule&&$\shr\Cal M_{\kappa}$&\tmphrule
  &$\shr\Cal M_{\lambda}$\cr
&\tmpstrut&\vrule&&\vrule&&\vrule&&\vrule\cr
$\tmpstrut\omega_1$&\tmphrule&$\add\Cal M_{\lambda}$&\tmphrule
  &$\add\Cal M_{\kappa}$&\tmphrule
  &$\non\Cal M_{\kappa}$&\tmphrule&$\non\Cal M_{\lambda}$\cr
}}$$

\noindent and

$$\vbox{\offinterlineskip
\halign{\hfil#\hfil&\hfil#\hfil&\hfil#\hfil&\hfil#\hfil
  &\hfil#\hfil&\hfil#\hfil&\hfil#\hfil
  &\hfil#\hfil&\hfil#\hfil&\hfil#\hfil&\hfil#\hfil\cr
&\tmpstrut&$\cov\Cal N_{\kappa}$&\tmphrule&$\non\Cal M_{\kappa}$
  &\tmphrule
  &$\cf\Cal M_{\kappa}$&\tmphrule&$\cf\Cal N_{\kappa}$
  &\tmphrule&$\kappa^{\omega}$\cr
&\tmpstrut&\vrule&&\vrule&&\vrule&&\vrule\cr
$\tmpstrut\omega_1$&\tmphrule&$\add\Cal N_{\kappa}$&\tmphrule
  &$\add\Cal M_{\kappa}$&\tmphrule
  &$\cov\Cal M_{\kappa}$&\tmphrule&$\non\Cal N_{\kappa}$\cr
}}$$

\noindent whenever $\omega\le\kappa\le\lambda$.   Show moreover that all
the results of
523E-523P %523E 523F 523G 523H 523I 523J 523L 523M 523N 523O 523P
have parallel forms referring to $\Cal M_{\kappa}$.
%523P

\spheader 523Yf In the language of 523Ye, show that
$\frak m_{\text{pc}\omega_1}\le\cov\Cal M_{\kappa}$ for every infinite
$\kappa$.
%\Hint{$\omega_1$ is a precaliber of $\RO(\{0,1\}^{\kappa})$
%523Ye 523P

\spheader 523Yg Show that Ostaszewski's $\clubsuit$ (4A1M) implies that
$\cov\Cal N_{\omega_1}=\cov\Cal M_{\omega_1}=\omega_1$.
%523Ye 523P 4A1N  set  E_{\xi}=\{x:x(\theta_{\xi}(n)) eventually constant\}
}%end of exercises

\leader{523Z}{Problem} Is there a proof in ZFC that
$\shr\Cal N_{\kappa}\ge\cff[\kappa]^{\le\omega}$ for every cardinal
$\kappa$?
%523M

\endnotes{
\Notesheader{523} The basic diagram 523B is natural and easy to
establish.   Of course it leaves a great deal of room, especially on the
right-hand side, where we have the increasing functions
$\non\Cal N_{\ssbullet}$, $\shr\Cal N_{\ssbullet}$ and
$\cf\Cal N_{\ssbullet}$, and rather weak constraints

\Centerline{$\lambda<\non\Cal N_{\kappa}\le\shr\Cal N_{\kappa}
\le\cf\Cal N_{\kappa}\le\kappa^{\omega}$ whenever $2^{\lambda}<\kappa$}

\noindent to control them.   However the generalized continuum
hypothesis is sufficient to determine exact values for all the cardinals
considered here (523P).

The combinatorics of $\cff[\kappa]^{\le\omega}$ and
almost-disjoint families of functions are extremely complex, and depend
in surprising ways on special axioms;  I think it possible that
the results of 523I-523J can be usefully extended.
%Kraszewski 3.18
However 523N at least reduces the measure-theoretic problem of
determining $\cf\Cal N_{\kappa}$ to a standard, if difficult, question
in infinitary combinatorics.   I do not know if there are corresponding
results concerning $\non\Cal N_{\kappa}$ and $\shr\Cal N_{\kappa}$ (see
523Kb and 523Z).

All the ideas in this section up to and including 523P can be applied to
ideals of meager sets (523Ye) and indeed to other classes of ideals
satisfying the fundamental lemma 523C; see {\smc Kraszewski 01}.
}%end of notes

\discrpage

\frfilename{mt524.tex}
\versiondate{29.9.10/10.9.13}
\copyrightdate{2003}

\def\chaptername{Cardinal functions of measure theory}
\def\sectionname{Radon measures}

\newsection{524}

It is a remarkable fact that for a Radon measure the principal cardinal
functions are determined by its measure algebra (524J), so can in most
cases be calculated in terms of the cardinals of
the last section (524P-524Q).   The proof of this seems to require a
substantial excursion involving not only measure algebras but also the
Banach lattices $\ell^1(\kappa)$ and/or the $\kappa$-localization
relation (524D, 524E).   The same machinery gives us formulae for the
cardinal functions of measurable algebras (524M).   The results of \S518
can be translated directly to give partial information on the
Freese-Nation numbers of measurable algebras (524O).   For covering
number and uniformity, we can see from 521L that strictly localizable
compact measures follow Radon measures.
I know of no such general results for any other class of measure, but
there are some bounds for cardinal functions of countably compact and
quasi-Radon measures, which I give in 524R-524T.  %524R 524S 524T

\leader{524A}{Notation} If $(X,\Sigma,\mu)$ is a measure space,
$\Cal N(\mu)$ will be the null ideal of $\mu$.
For any cardinal $\kappa$, $\nu_{\kappa}$ will be the usual measure on
$\{0,1\}^{\kappa}$, $\Tau_{\kappa}$ its domain and
$(\frak B_{\kappa},\bar\nu_{\kappa})$ its measure algebra.   As in
\S\S522-523, I will
write $\Cal N_{\kappa}$ for $\Cal N(\nu_{\kappa})$ and $\Cal N$ for the
null ideal of Lebesgue measure on $\Bbb R$, so that
$(\Bbb R,\Cal N)$ and $(\{0,1\}^{\omega},\Cal N_{\omega})$ are
isomorphic\cmmnt{ (522Wa)}.
If $\frak A$ is any Boolean algebra, I write $\frak A^+$ for
$\frak A\setminus\{0\}$ and $\frak A^-$ for $\frak A\setminus\{1\}$.
If $(A,R,B)$ is a supported relation, $R\strprime$ is the
relation $\{(a,I):a\in R^{-1}[I]\}$\cmmnt{ (see 512F)}.
For any cardinal $\kappa$, $(\kappa^{\Bbb N},\subseteq^*,\Cal S_{\kappa})$
will be the $\kappa$-localization relation\cmmnt{ (522K)}.

\leader{524B}{Proposition} Let $(X,\frak T,\Sigma,\mu)$ be a
$\sigma$-finite
Radon measure space with Maharam type $\kappa$.   Then
$\Cal N(\mu)\prT\frak B_{\kappa}^-$.

\proof{{\bf (a)} Suppose, to begin with, that $\mu X=1$ and
$\kappa\ge\omega$.   Let $\frak A$ be the measure algebra of the Radon
product measure $\tilde\lambda$ on $Y=X^{\Bbb N}$.   Then
$\frak A\cong\frak B_{\kappa}$.   \Prf\ By 417E(b-i), $\frak A$ is
isomorphic
to the measure algebra of the usual product measure $\lambda$ on $Y$,
which by 334E is isomorphic to $\frak B_{\kappa}$.\ \Qed

For $E\in\Cal N(\mu)$, let
$\sequence{i}{F_{Ei}}$ be a sequence of closed subsets of
$X$ such that $E\cap F_{Ei}=\emptyset$ and $\mu F_{Ei}\ge 1-2^{-i-1}$
for every $n\in\Bbb N$.   Then

\Centerline{$\tilde\lambda(\prod_{i\in\Bbb N}F_{Ei})
=\lambda(\prod_{i\in\Bbb N}F_{Ei})
\ge\prod_{i\in\Bbb N}(1-2^{-i-1})>0$;}

\noindent set

\Centerline{$\phi(E)=(Y\setminus\prod_{i\in\Bbb N}F_{Ei})^{\ssbullet}
\in\frak A^-$.}

For $b\in\frak A^-$ let $K_b\subseteq Y$ be a non-empty compact
self-supporting set such that $K_b^{\ssbullet}\Bcap b=0$.   Set
$\pi_i(y)=y(i)$ for $i\in\Bbb N$ and $y\in Y$.   Then each
$\pi_i[K_b]\subseteq X$ is compact and
$K_b\subseteq\prod_{i\in\Bbb N}\pi_i^{-1}[\pi_i[K_b]]$, so
$\prod_{i\in\Bbb N}\mu\pi_i[K_b]>0$ and
$\sup_{i\in\Bbb N}\mu\pi_i[K_b]=1$;  set

\Centerline{$\psi(b)
=X\setminus\bigcup_{i\in\Bbb N}\pi_i[K_b]
\in\Cal N(\mu)$.}

If $E\in\Cal N(\mu)$ and $b\in\frak A^-$ and $\phi(E)\Bsubseteq b$
and $j\in\Bbb N$, then

\Centerline{$K_b\setminus\pi_j^{-1}[F_{Ej}]
\subseteq K_b\setminus\prod_{i\in\Bbb N}F_{Ei}$}

\noindent is negligible.   As $K_b$ is self-supporting,
$K_b\setminus\pi_j^{-1}[F_{Ej}]$ is empty and
$\pi_j[K_b]\subseteq F_{Ej}$.   But this means that
$\pi_j[K_b]\cap E=\emptyset$ for every $j\in\Bbb N$, so that
$E\subseteq\psi(b)$.

This shows that $\phi$ is a Tukey function, so that
$\Cal N(\mu)\prT\frak A^-\cong\frak B_{\kappa}^-$.

\medskip

{\bf (b)} If $\kappa$ is finite, $\Cal N(\mu)$ has a greatest member and
the constant function with value $0$ is a Tukey function from
$\Cal N(\mu)$ to $\frak B_{\kappa}^-$ and the result is trivial.   If
$\kappa$ is infinite and $\mu X\ne 1$, then, because $\mu$ is
$\sigma$-finite and not trivial, there is a function
$f:X\to\ooint{0,\infty}$ such that $\int fd\mu=1$ (215B(ix)).   Let
$\nu$ be the corresponding indefinite-integral measure;  then $\nu$ is a
Radon probability measure (416Sa) with the same measurable sets and the
same negligible sets as $\mu$ (234L),
so $\Sigma/\Cal N(\nu)=\Sigma/\Cal N(\mu)$ has Maharam type $\kappa$.
In this case, (a) tells us that
$\Cal N(\mu)=\Cal N(\nu)\prT\frak B_{\kappa}^-$.
}%end of proof of 524B

\leader{524C}{Lemma} Let $P$ be a partially ordered set such that
$p\vee q=\sup\{p,q\}$ is defined for all $p$, $q\in P$.   Suppose that
$\rho$ is a metric on $P$ such that $P$ is complete\cmmnt{ (as a
metric space)} and $\vee:P\times P\to P$ is
uniformly continuous with respect to $\rho$.   Let $Q\subseteq P$ be an
open set, and $\kappa\ge d(Q)$ a cardinal.   Then
$(Q,\le\strprime,[Q]^{<\omega})
\prGT(\ell^1(\kappa),\le,\ell^1(\kappa))$.   If $Q$ is upwards-directed,
then $Q\prT\ell^1(\kappa)$.

\proof{{\bf (a)} If $Q$ is finite, then we can set $\phi(q)=0$ for every
$q\in Q$, $\psi(x)=Q$ for every $x\in\ell^1(\kappa)$ and $(\phi,\psi)$
will be a Galois-Tukey connection from
$(Q,\le\strprime,[Q]^{<\omega})$ to
$(\ell^1(\kappa),\le,\ell^1(\kappa))$.
So let us suppose that $Q$ and $\kappa$ are infinite.

\medskip

{\bf (b)} Let $\ofamily{\xi}{\kappa}{q_{\xi}}$ run over a dense subset
of $Q$.   For each $q\in Q$ let $m(q)\in\Bbb N$ be such that
$\{p:p\in P$, $\rho(p,q)\le 2^{-m(q)}\}\subseteq Q$.   For each
$n\in\Bbb N$, let $\delta_n>0$ be such that
$\rho(\sup I,\sup J)\le 2^{-n}$ whenever
$\emptyset\ne I\subseteq J\subseteq P$ and $\#(J)\le 2^n$ and
$\max_{q\in J}\min_{p\in I}\rho(p,q)\le 2\delta_n$;   such exists
because $\ofamily{i}{k}{p_i}\mapsto\sup_{i<k}p_i:P^k\to P$ is uniformly
continuous whenever $k>0$, and in particular when $k=2^n$.   Reducing
the $\delta_n$ if necessary, we may suppose that
$\delta_{n+1}\le\delta_n\le 2^{-n}$ for every $n$.

\medskip

{\bf (c)} Define $\phi:Q\to\ell^1(\kappa)$ as follows.   Given $p\in Q$,
choose a sequence $\sequencen{\xi(p,n)}$ in $\kappa$ such that
$\rho(p,q_{\xi(p,n)})\le\delta_{n+1}$ for every $n$.   Take
$\phi(p)\in\ell^1(\kappa)$ such that

\Centerline{$\phi(p)(m(p))\ge 1$,
\quad$\phi(p)(\xi(p,n))\ge 2^{-n}$ for every $n\in\Bbb N$}

\noindent (regarding $m(p)$ as a finite ordinal).

\medskip

{\bf (d)} Define $\psi:\ell^1(\kappa)\to[Q]^{<\omega}$ as follows.
Given $x\in\ell^1(\kappa)$, set
$K_n(x)=\{q_{\xi}:\xi<\kappa,\,x(\xi)\ge 2^{-n}\}$ for $n\in\Bbb N$.
Then

\Centerline{$\sum_{n=0}^{\infty}2^{-n}\#(K_n(x))
\le\sum_{\xi<\kappa}\sum\{2^{-n}:x(\xi)\ge 2^{-n}\}\le 2\|x\|_1<\infty$,}

\noindent so there
is a $k(x)\in\Bbb N$ such that $x(n)<1$ for $n\in\omega\setminus k(x)$
and also $\#(K_n(x))\le 2^n$ for $n\ge k(x)$.   Set
$\tilde K(x)=K_{k(x)}(x)$.   For $s\in\tilde K(x)$ set

\Centerline{$I(x,s,k(x))=\{s\}$,
\quad$I(x,s,n+1)
=\{q:q\in K_{n+1}(x),\,\rho(q,I(x,s,n))\le 2\delta_{n+1}\}$}

\noindent for $n\ge k(x)$, writing $\rho(q,I)$ for
$\inf_{q'\in I}\rho(q,q')$.   Because $\sequencen{K_n(x)}$ is
non-decreasing, so is $\langle I(x,s,n)\rangle_{n\ge k(x)}$.   Set
$r_{xsn}=\sup I(x,s,n)$ in $P$ for $n\ge k(x)$;  then
$\rho(r_{x,s,n+1},r_{xsn})\le 2^{-n-1}$ for every $n\ge k(x)$, by the
choice of $\delta_{n+1}$, so $r_{xs}=\lim_{n\to\infty}r_{xsn}$ is
defined in $P$.   Set $\psi(x)=Q\cap\{r_{xs}:s\in\tilde K(x)\}$.

\medskip

{\bf (e)} Now $(\phi,\psi)$ is a Galois-Tukey connection from
$(Q,\le\strprime,[Q]^{<\omega})$ to
$(\ell^1(\kappa),\le,\ell^1(\kappa))$.
\Prf\ Suppose that $p\in Q$ and $x\in\ell^1(\kappa)$ are such that
$\phi(p)\le x$.   Then $q_{\xi(p,n)}\in K_n(x)$ for every $n$, so
$s=q_{\xi(p,k(x))}\in\tilde K(x)$.   Also $q_{\xi(p,n)}\in I(x,s,n)$ for
every $n\ge k(x)$, because

\Centerline{$\rho(q_{\xi(p,n+1)},q_{\xi(p,n)})
\le\delta_{n+2}+\delta_{n+1}\le 2\delta_{n+1}$}

\noindent for every $n$.   So $q_{\xi(p,n)}\le r_{xsn}$ for every
$n\ge k(x)$.   It follows that

\Centerline{$p\vee r_{xs}=\lim_{n\to\infty}q_{\xi(p,n)}\vee r_{xsn}
=\lim_{n\to\infty}r_{xsn}=r_{xs}$}

\noindent and $p\le r_{xs}$.

By the choice of $k(x)$, we also have $\phi(p)(n)<1$ for $n\ge k(x)$, so
that $m(p)<k(x)$.   We therefore have

$$\eqalignno{\rho(r_{xs},p)
&\le\rho(q_{\xi(p,k(x))},p)
  +\sum_{n=k(x)}^{\infty}\rho(r_{x,s,n+1},r_{xsn})\cr
\displaycause{because $s=q_{\xi(p,k(x))}$ is the unique member of
$I(x,s,k(x))$, so is equal to $r_{x,s,k(x)}$}
&\le\delta_{k(x)+1}+\sum_{n=k(x)}^{\infty}2^{-n-1}
\le 2^{-k(x)-1}+2^{-k(x)}
\le 2^{-m(p)}\cr}$$

\noindent and $r_{xs}\in Q$.   So $p\le r_{xs}\in\psi(x)$ and
$p\le\strprime\psi(x)$.   As $p$ and $x$ are arbitrary,
$(\phi,\psi)$ is a Galois-Tukey connection.\ \Qed

\medskip

{\bf (f)} So $(Q,\le\strprime,[Q]^{<\omega})
\prGT(\ell^1(\kappa),\le,\ell^1(\kappa))$, as claimed.

Finally, if $Q$ is upwards-directed, then $\add Q\ge\omega$, so
$(Q,\le,Q)\equivGT(Q,\le\strprime,[Q]^{<\omega})$ (513Id) and
$(Q,\penalty-50\le\nobreak,\penalty-50Q)\penalty-100
   \prGT(\ell^1(\kappa),\le,\ell^1(\kappa))$, that is,
$Q\prT\ell^1(\kappa)$.
}%end of proof of 524C

\leader{524D}{Proposition} If $\kappa$ is any cardinal,

\Centerline{$(\frak B_{\kappa}^-,\Bsubseteqshort\strprime,
  [\frak B_{\kappa}^-]^{<\omega})
\prGT(\ell^1(\kappa),\le,\ell^1(\kappa))$.}

\proof{ If $\kappa$ is finite then $\frak B_{\kappa}^-$ is finite and
the result is trivial.   Otherwise, if we give $\frak B_{\kappa}$ its
measure metric $\rho$ (323Ad), then it is a complete metric space in
which $\Bcup$ is uniformly continuous (323Gc, 323B) and
$\frak B_{\kappa}^-=\frak B_{\kappa}\setminus\{1\}$ is an open set.
Now the topological density of $\frak B_{\kappa}^-$ and
$\frak B_{\kappa}$ is $\kappa$, by 521E;  so
524C gives the result.
}%end of proof of 524D

\leader{524E}{Proposition} Let $\kappa$ be an infinite cardinal.   Then

\Centerline{$(\ell^1(\kappa),\le\strprime,
  [\ell^1(\kappa)]^{\le\omega})
\prGT(\kappa^{\Bbb N},\subseteq^*,\Cal S_{\kappa})$.}

\proof{{\bf (a)} For each $i\in\Bbb N$, let
$\ofamily{\xi}{\kappa}{z_{i\xi}}$ run over a norm-dense subset of
$\{x:x\in\ell^1(\kappa)^+,\,\|x\|_1\le 4^{-i}\}$.   Now there is a
function $\phi:\ell^1(\kappa)\to\kappa^{\Bbb N}$ such that

\Centerline{for every $x\in\ell^1(\kappa)$, $n\in\Bbb N$ there is a
$k\in\Bbb N$ such that $x\le k\sum_{i=n}^{\infty}z_{i,\phi(x)(i)}$.}

\noindent\Prf\ Given $x\in\ell^1(\kappa)$, choose $\sequencen{x_n}$,
$\sequencen{\xi_n}$, $\sequencen{k_n}$ inductively, as follows.   Take
$k_0\ge 1$ such that $\|x^+\|_1\le k_0$;  set $x_0=k_0^{-1}x^+$ and take
$\xi_0<\kappa$ such that $\|x_0-z_{0\xi_0}\|_1<\bover14$.   Given that
$x_n\in\ell^1(\kappa)^+$, $\xi_n<\kappa$ are such that
$\|x_n-z_{n\xi_n}\|_1<4^{-n-1}$, let $k_{n+1}\ge 1$ be such that
$\|x_{n+1}\|_1\le 4^{-n-1}$ where
$x_{n+1}=(x_n-z_{n\xi_n})^++k_{n+1}^{-1}x^+$, and take
$\xi_{n+1}<\kappa$ such that
$\|x_{n+1}-z_{n+1,\xi_{n+1}}\|_1<4^{-n-2}$;  continue.
At the end of the process, set $\phi(x)=\sequencen{\xi_n}$.

Now, for any $n\in\Bbb N$, we have $x\le x^+\le k_nx_n$.   But we also
have, for any $m\ge n$, $x_{m+1}\ge x_m-z_{m\xi_m}$, so that
$x_n\le x_m+\sum_{i=n}^{m-1}z_{i\xi_i}$ for every $m\ge n$.
Since $\|x_m\|_1\le 4^{-m}$ for every $m$,
$\lim_{m\to\infty}x_m=0$ and

\Centerline{$x\le k_nx_n\le k_n\sum_{i=n}^{\infty}z_{i,\phi(x)(i)}$.}

\noindent As $x$ and $n$ are arbitrary, $\phi$ is a suitable function.\
\Qed

\medskip

{\bf (b)} Define $\psi_0:\Cal S_{\kappa}\to\ell^1(\kappa)$ by setting
$\psi_0(S)=\sum_{(i,\xi)\in S}z_{i\xi}$;  because

\Centerline{$\sum_{(i,\xi)\in S}\|z_{i\xi}\|_1
\le\sum_{i=0}^{\infty}4^{-i}\#(S[\{i\}])\le\sum_{i=0}^{\infty}2^{-i}$}

\noindent is finite, $\psi_0(S)$ is well defined in $\ell^1(\kappa)$ for
every $S\in\Cal S_{\kappa}$ (4A4Ie).   Now define
$\psi:\Cal S_{\kappa}\to[\ell^1(\kappa)]^{\le\omega}$ by setting
$\psi(S)=\{k\psi_0(S):k\in\Bbb N\}$ for $S\in\Cal S_{\kappa}$.

\medskip

{\bf (c)} If $x\in\ell^1$ and $S\in\Cal S_{\kappa}$ are such that
$\phi(x)\subseteq^*S$, then $x\le\strprime\psi(S)$.   \Prf\ Let
$n\in\Bbb N$ be such that $(i,\phi(x)(i))\in S$ for every $i\ge n$.
Then there is a $k\in\Bbb N$ such that

\Centerline{$x\le k\sum_{i=n}^{\infty}z_{i,\phi(x)(i)}
\le k\sum_{(i,\xi)\in S}z_{i\xi}=k\psi_0(S)\in\psi(S)$.  \Qed}

\medskip

{\bf (d)} Thus $(\phi,\psi)$ is a Galois-Tukey connection and

\Centerline{$(\ell^1(\kappa),\le\strprime,
[\ell^1(\kappa)]^{\le\omega})
\prGT(\kappa^{\Bbb N},\subseteq^*,\Cal S_{\kappa})$.}
}%end of proof of 524E

\leader{524F}{Lemma} Let $(X,\Sigma,\mu)$ be a
countably compact measure space with
Maharam type $\kappa$.

(a) If $\mu$ is a \Mth\ probability measure, there is a family
$\ofamily{\xi}{\kappa}{E_{\xi}}$ in $\Cal N(\mu)$ such
that $\bigcup_{\xi\in A}E_{\xi}$ has full outer
measure for every uncountable $A\subseteq\kappa$.

(b) If $\mu$ is $\sigma$-finite, there is a family $\ofamily{\xi}{\kappa}{E_{\xi}}$ in
$\Cal N(\mu)$ such that $\bigcup_{\xi\in A}E_{\xi}$ is non-negligible
for every uncountable $A\subseteq\kappa$.

\proof{ Let $\frak A$ be the measure algebra of $(X,\Sigma,\mu)$.

\medskip

{\bf (a)} If $\kappa$ is countable, we can take $E_{\xi}=\emptyset$ for
every $\xi$.   Otherwise, $\frak A$ is
$\tau$-generated by a stochastically independent family
$\ofamily{\xi}{\kappa}{a_{\xi}}$ of elements of measure $\bover12$, and
for every $G\in\Sigma$ there is a smallest countable set
$I_G\subseteq\kappa$ such that $G^{\ssbullet}$ is in the closed
subalgebra of $\frak A$ generated by $\{a_{\xi}:\xi\in I_G\}$ (254Rd or
325Mb).
For each $\xi<\kappa$ choose $F_{\xi}\in\Sigma$ such that
$F_{\xi}^{\ssbullet}=a_{\xi}$.   Let $\Cal K$ be a countably compact
class such that $\mu$ is inner regular with respect to $\Cal K$.

Let $\ofamily{\xi}{\kappa}{J_{\xi}}$ be a disjoint family of subsets of
$\kappa$ all with cardinal $\omega_1$.
For each $\xi<\kappa$ choose $\sequencen{K_{\xi n}}$,
$\sequencen{\alpha_{\xi n}}$ inductively, as follows.
$\alpha_{\xi 0}=\min J_{\xi}$.   Given $\alpha_{\xi n}$ and
$\ofamily{i}{n}{K_{\xi i}}$, let $K_{\xi n}\in\Cal K$ be such that
$K_{\xi n}\cap F_{\alpha_{\xi n}}=\emptyset$ and
$\mu(K_{\xi n})\ge\bover12(1-3^{-n-2})$;  now let $\alpha_{\xi,n+1}$ be
a member of $J_{\xi}$ not belonging to
$I_{K_{\xi i}}\cup\{\alpha_{\xi i}\}$ for any $i\le n$.   Continue.
Set

\Centerline{$E_{\xi}=\bigcup_{n\in\Bbb N}\bigcap_{m\ge n}K_{\xi m}
\subseteq\bigcup_{n\in\Bbb N}\bigcap_{m\ge n}(X\setminus F_{\alpha_{\xi m}})$,}

\noindent so that $E_{\xi}$ is negligible, because all the $\alpha_{\xi m}$
are different, so that $\sequence{m}{F_{\alpha_{\xi m}}}$ is stochastically
independent.

Now suppose that
$A\subseteq\kappa$ is uncountable, and that $F\subseteq X$ is measurable and not negligible.   Let
$K\in\Cal K$ be such that $K\subseteq F$ and $\mu K>0$;  let $\xi\in A$ be such that
$I_K\cap J_{\xi}=\emptyset$;  let $n\in\Bbb N$ be such that $\mu K\ge 3^{-n-1}$.
Set
$G_m=K\cap\bigcap_{n\le i<m}K_{\xi i}$ for $m\ge n$.   Then
$I_{G_m}\subseteq I_K\cup\bigcup_{i<m}I_{K_{\xi i}}$ does not contain
$\alpha_{\xi m}$, for any $m$.   This means that

\Centerline{$\mu G_{m+1}
=\mu(G_m\cap K_{\xi m})
\ge\mu(G_m\setminus F_{\alpha_{\xi m}})-\Bover{3^{-m-2}}{2}
=\Bover12(\mu G_m-3^{-m-2})$}

\noindent for every $m\ge n$, and an easy induction shows that
$\mu G_m\ge 3^{-m-1}$ for every $m$.   But this tells us that every
$G_m$ is non-empty;  because $\Cal K$ is a countably compact class,
$K\cap E_{\xi}\supseteq\bigcap_{m\ge n}G_m$ is non-empty, and $F$ meets
$E_{\xi}$.

As $F$ is arbitrary, $\bigcup_{\xi\in A}E_{\xi}$ has full outer measure.

\medskip

{\bf (b)} For the general case, because $\mu$ is $\sigma$-finite, there
is a countable partition of unity $\familyiI{a_i}$ in
$\frak A$ such that all the principal ideals
$\frak A_{a_i}$ are totally finite and \Mth\ (use 332A),
and we can find a partition $\familyiI{X_i}$ of $X$ into measurable sets
such that $X_i^{\ssbullet}=a_i$ for each $i$.   Moreover, the subspace
measure $\mu_{X_i}$ on $X_i$ is countably compact (451Db).
Writing $\kappa_i$ for the Maharam type of $\frak A_{a_i}$, there is a
family $\ofamily{\xi}{\kappa_i}{E_{i\xi}}$ of negligible subsets of
$X_i$ such that $\{\xi:\xi<\kappa_i,\,E_{i\xi}\subseteq E\}$ is
countable for every negligible set $E$.
(Apply (a) to a scalar multiple of
$\mu_{X_i}$.)   Now we know from 332S that
$\kappa=\sup_{i\in I}\kappa_i=\#(\{(i,\xi):i\in I,\,\xi<\kappa_i\})$.
On the other hand, for any negligible set $E\subseteq X$,
$\{(i,\xi):i\in I,\,\xi<\kappa_i,\,E_{i\xi}\subseteq E\}$ is countable.
So if we re-enumerate $\langle E_{i\xi}\rangle_{i\in I,\xi<\kappa_i}$ as
$\ofamily{\xi}{\kappa}{E_{\xi}}$ we shall have an appropriate family.
}%end of proof of 524F

\leader{524G}{Proposition} Let $(X,\frak T,\Sigma,\mu)$ be a
\Mth\ Radon probability space with Maharam type $\kappa\ge\omega$.   Then
$(\kappa^{\Bbb N},\subseteq^*,\Cal S_{\kappa})
\prGT(\Cal N(\mu),\subseteq\penalty+100,\Cal N(\mu))$.

\proof{ (Compare 522M.)

\medskip

{\bf (a)} By 524F, there is a family $\ofamily{\xi}{\kappa}{E_{\xi}}$ in
$\Cal N(\mu)$ such that $\{\xi:E_{\xi}\subseteq E\}$ is countable for
every $E\in\Cal N(\mu)$.
Next, because the measure algebra of $\mu$ is isomorphic to the measure
algebra of the usual measure on $[0,1]^{\Bbb N\times\kappa}$, there is a
stochastically independent family
$\langle G_{i\xi}\rangle_{i\in\Bbb N,\xi<\kappa}$ in $\Sigma$ such that
$\mu G_{i\xi}=2^{-i}$ for every $i\in\Bbb N$ and $\xi<\kappa$.
For $f\in\kappa^{\Bbb N}$ set

\Centerline{$\phi(f)=\bigcup_{n\in\Bbb N}E_{f(n)}
\cup\bigcap_{n\in\Bbb N}\bigcup_{m\ge n}G_{m,f(m)}\in\Cal N(\mu)$.}

\medskip

{\bf (b)} Take $E\in\Cal N(\mu)$ and set
$I_E=\{\xi:E_{\xi}\subseteq E\}$, so that $I_E$ is countable.
Define $\pi_E:X\to\{0,1\}^{\Bbb N\times I_E}$ by setting
$\pi_E(x)(i,\xi)=1$ if $x\in G_{i\xi}$, $0$ otherwise.
Then there is a non-empty compact self-supporting set $K_E$ such that
$\pi_E\restr K_E$ is continuous.   \Prf\
Then $\pi_E$ is measurable, therefore almost continuous (418J), and
there is a non-negligible measurable set $H\subseteq X\setminus E$ such
that $\pi_E\restr H$ is continuous.   Because $\mu$ is inner regular
with respect to the compact self-supporting sets, there is a
non-negligible compact self-supporting $K_E\subseteq H$, and this has
the required property.\ \Qed

$\pi_E[K_E]$ is compact.   Let $\sequencen{W_n(E)}$ run
over the family of open-and-closed subsets $W$ of
$\{0,1\}^{\Bbb N\times I_E}$ meeting $\pi_E[K_E]$.   Then
$\pi_E^{-1}[W_n(E)]$ is a non-empty relatively open subset of $K_E$ for
every $n$;  because $K_E$ is self-supporting, $\pi_E^{-1}[W_n(E)]$ is
never negligible.   Set

\Centerline{$J(E,n,i)
=\{\xi:\xi\in I_E,\,\pi_E^{-1}[W_n(E)]\cap G_{i\xi}=\emptyset\}$}

\noindent for $n$, $i\in\Bbb N$.   Because
$\langle G_{i\xi}\rangle_{i\in\Bbb N,\xi\in I_E}$ is stochastically
independent,

\Centerline{$\sum_{i=0}^{\infty}2^{-i}\#(J(E,n,i))
=\sum\{\mu G_{i\xi}:
  i\in\Bbb N,\,\xi\in I_E,\,G_{i\xi}\cap\pi_E^{-1}[W_n(E)]=\emptyset\}$}

\noindent is finite, by the Borel-Cantelli lemma (273K).   For each $n$,
let $k(E,n)\in\Bbb N$ be such that $2^{-i}\#(J(E,n,i))\le 2^{-n-1}$ for
$i\ge k(E,n)$, and set

\Centerline{$\psi(E)
=\bigcup_{n\in\Bbb N}\{(i,\xi):i\ge k(E,n),\,\xi\in J(E,n,i)\}
\subseteq\Bbb N\times\kappa$.}

\noindent Then

$$\eqalign{\#(\{\xi:(i,\xi)\in\psi(E)\}
&\le\sum_{n\in\Bbb N,k(E,n)\le i}\#(J(E,n,i))\cr
&\le\sum_{n\in\Bbb N,k(E,n)\le i}2^{-n-1}2^i
\le 2^i\cr}$$

\noindent for every $i\in\Bbb N$, and $\psi(E)\in\Cal S_{\kappa}$.

\medskip

{\bf (c)} Now $(\phi,\psi)$ is a Galois-Tukey connection from
$(\kappa^{\Bbb N},\subseteq^*,\Cal S_{\kappa})$ to
$(\Cal N(\mu),\subseteq,\Cal N(\mu))$.   \Prf\ Suppose that
$f\in\kappa^{\Bbb N}$ and $E\in\Cal N(\mu)$ are such that
$\phi(f)\subseteq E$.   Because $E_{f(n)}\subseteq\phi(f)$,
$f(n)\in I_E$ for every $n\in\Bbb N$.   Next, $K_E$ does not meet
$\phi(f)$, so $K_E\cap\bigcap_{n\in\Bbb N}\bigcup_{m\ge n}G_{m,f(m)}$ is
empty, that is,

\Centerline{$\pi_E[K_E]\cap
\bigcap_{n\in\Bbb N}\bigcup_{m\ge n}
\{w:w\in\{0,1\}^{\Bbb N\times I_E},\,w(m,f(m))=1\}=\emptyset$.}

\noindent By Baire's theorem, there is some $m\in\Bbb N$ such that

\Centerline{$\pi_E[K_E]\cap\bigcup_{i\ge m}
\{w:w\in\{0,1\}^{\Bbb N\times I_E},\,w(i,f(i))=1\}$}

\noindent is not dense in $\pi_E[K_E]$, and there is an $n\in\Bbb N$
such that

\Centerline{$W_n(E)\cap\bigcup_{i\ge m}
\{w:w\in\{0,1\}^{\Bbb N\times I_E},\,w(i,f(i))=1\}=\emptyset$.}

\noindent In this case, $f(i)\in J(E,n,i)$ for every $i\ge m$.   But
this means that $(i,f(i))\in\psi(E)$ for every $i\ge\max(m,k(E,n))$, so
that $f\subseteq^*\psi(E)$.   As $f$ and $E$ are arbitrary,
$(\phi,\psi)$ is a Galois-Tukey connection.\ \Qed

\medskip

{\bf (d)} Thus $\phi$ and $\psi$ witness that
$(\kappa^{\Bbb N},\subseteq^*,\Cal S_{\kappa})
\prGT(\Cal N(\mu),\subseteq,\Cal N(\mu))$, as claimed.
}%end of proof of 524G

\leader{524H}{Corollary} Let $\kappa$ be an infinite cardinal, and $\mu$
a \Mth\ Radon probability measure with Maharam type $\kappa$.
Then $(\frak B_{\kappa}^+,\Bsupseteqshort\strprime,
  [\frak B_{\kappa}^+]^{\le\omega})$,
$(\ell^1(\kappa),
  \le\strprime,\penalty-100[\ell^1(\kappa)]^{\le\omega})$,
$(\kappa^{\Bbb N},\subseteq^*,\Cal S_{\kappa})$
and $(\Cal N(\mu),\subseteq\nobreak,\Cal N(\mu))$ are
Galois-Tukey equivalent.

\proof{ By 512Gb, 524D and 524B,

\Centerline{$(\frak B_{\kappa}^-,\Bsubseteqshort\strprime,
  [\frak B_{\kappa}^-]^{<\omega_1})
  \prGT(\frak B_{\kappa}^-,\Bsubseteqshort\strprime,
  [\frak B_{\kappa}^-]^{<\omega})
  \prGT(\ell^1(\kappa),\le,\ell^1(\kappa))$,}

\Centerline{$(\Cal N(\mu),\subseteq,\Cal N(\mu))
\prGT(\frak B_{\kappa}^-,\Bsubseteqshort,\frak B_{\kappa}^-)$.}

\noindent So

$$\eqalignno{(\frak B_{\kappa}^+,\Bsupseteqshort\strprime,
  [\frak B_{\kappa}^+]^{\le\omega})
&\cong(\frak B_{\kappa}^-,\Bsubseteqshort\strprime,
  [\frak B_{\kappa}^-]^{\le\omega})
=(\frak B_{\kappa}^-,\Bsubseteqshort\strprime,
  [\frak B_{\kappa}^-]^{<\omega_1})\cr
&\prGT(\ell^1(\kappa),\le\strprime,[\ell^1(\kappa)]^{<\omega_1})\cr
\displaycause{512Gd}
&=(\ell^1(\kappa),\le\strprime,[\ell^1(\kappa)]^{\le\omega})
\prGT(\kappa^{\Bbb N},\subseteq^*,\Cal S_{\kappa})\cr
\displaycause{524E}
&\prGT(\Cal N(\mu),\subseteq,\Cal N(\mu))\cr
\displaycause{524G}
&\equivGT(\Cal N(\mu),\subseteq\strprime,[\Cal N(\mu)]^{\le\omega})\cr
\displaycause{513Id again}
&\prGT(\frak B_{\kappa}^-,\Bsubseteqshort\strprime,
  [\frak B_{\kappa}^-]^{\le\omega})\cr}$$

\noindent by 512Gb again.
}%end of proof of 524H

\leader{524I}{Corollary} Let $\mu$ be a \Mth\ Radon
probability measure with infinite Maharam type $\kappa$.   Then

\Centerline{$\add\Cal N(\mu)
=\add\Cal N_{\kappa}=\add_{\omega}\ell^1(\kappa)$,}

\Centerline{$\cf\Cal N(\mu)=\cf\Cal N_{\kappa}=\cf\ell^1(\kappa)$.}

\proof{ By 524H and 512Db,

$$\eqalign{\add(\ell^1(\kappa),\le\strprime,
  [\ell^1(\kappa)]^{\le\omega})
&=\add(\Cal N(\mu),\subseteq,\Cal N(\mu))\cr
&=\add(\Cal N(\nu_{\kappa}),\subseteq,\Cal N(\nu_{\kappa}))
=\add(\Cal N_{\kappa},\subseteq,\Cal N_{\kappa}).\cr}$$

\noindent But
$\add(\ell^1(\kappa),\le\strprime,[\ell^1(\kappa)]^{\le\omega})
=\add_{\omega}\ell^1(\kappa)$ (513Ia), while
$\add(\Cal N(\mu),\subseteq,\Cal N(\mu))=\add\Cal N(\mu)$ and
$\add(\Cal N_{\kappa},\subseteq,\Cal N_{\kappa})=\add\Cal N_{\kappa}$
(512Ea).   So

\Centerline{$\add_{\omega}\ell^1(\kappa)
=\add\Cal N(\mu)=\add\Cal N_{\kappa}$.}

On the other side, 512Da tells us that

\Centerline{$\cov(\Cal N_{\kappa},\subseteq,\Cal N_{\kappa})
=\cov(\Cal N(\mu),\subseteq,\Cal N(\mu))
=\cov(\ell^1(\kappa),\le\strprime,
  [\ell^1(\kappa)]^{\le\omega})$.}

\noindent But

\Centerline{$\cov(\Cal N_{\kappa},\subseteq,\Cal N_{\kappa})
=\cf\Cal N_{\kappa}$,
\quad$\cov(\Cal N(\mu),\subseteq,\Cal N(\mu))
=\cf\Cal N(\mu)$}

\noindent (512Ea).   Next, $\cf\ell^1(\kappa)>\omega$.   \Prf\ If
$\sequencen{x_n}$ is any sequence in $\ell^1(\kappa)$, then (because
$\kappa$ is infinite)
$F_n=\{x:x\le x_n\}$ is nowhere dense (for the norm topology) for any
$n\in\Bbb N$, so $\sequencen{F_n}$ cannot cover $\ell^1(\kappa)$ (4A2Ma)
and $\{x_n:n\in\Bbb N\}$ cannot be cofinal.\ \QeD\   So 512Gf tells us
that

\Centerline{$\cov(\ell^1(\kappa),\le\strprime,
  [\ell^1(\kappa)]^{\le\omega})
=\cov(\ell^1(\kappa),\le,\ell^1(\kappa))=\cf\ell^1(\kappa)$.}

\noindent Putting these together,

\Centerline{$\cf\Cal N(\mu)=\cf\Cal N_{\kappa}=\cf\ell^1(\kappa)$}

\noindent as required.
}%end of proof of 524I

\leader{524J}{Theorem} Let $(X,\frak T,\Sigma,\mu)$ and
$(Y,\frak S,\Tau,\nu)$ be
Radon measure spaces with non-zero measure and isomorphic measure
algebras.

(a) $\Cal N(\mu)$ and $\Cal N(\nu)$ are Tukey equivalent,
so $\add\mu=\add\Cal N(\mu)=\add\Cal N(\nu)=\add\nu$ and
$\cf\Cal N(\mu)=\cf\Cal N(\nu)$.

(b) $(X,\in,\Cal N(\mu))$ and $(Y,\in,\Cal N(\nu))$ are Galois-Tukey
equivalent, so
$\cov\Cal N(\mu)=\cov\Cal N(\nu)$ and $\non\Cal N(\mu)=\non\Cal N(\nu)$.

\proof{{\bf (a)} Let $\frak A$, $\frak B$ be the measure algebras of
$\mu$ and $\nu$.   Let $\familyiI{a_i}$ be a partition of unity in
$\frak A^+$ such that all the principal ideals
$\frak A_{a_i}$ are homogeneous and totally finite, and
$\familyiI{b_i}$ a matching family in $\frak B$, so that
$\frak A_{a_i}\cong\frak B_{b_i}$ for every $i$.   Because
$(X,\Sigma,\mu)$ and $(Y,\Tau,\nu)$ are strictly localizable (416B),
there are decompositions $\familyiI{X_i}$ and $\familyiI{Y_i}$ of $X$,
$Y$ respectively such that $X_i^{\ssbullet}=a_i$ and
$Y_i^{\ssbullet}=b_i$ for every $i$ (322M).   Write $\mu_{X_i}$,
$\nu_{Y_i}$ for the corresponding subspace measures;  of course these
are Radon measures (416Rb).   Then
$\Cal N(\mu_{X_i})$ and $\Cal N(\nu_{Y_i})$ are Tukey equivalent for
every $i$.   \Prf\ If the common Maharam type of $\frak A_{a_i}$ and
$\frak B_{b_i}$ is infinite, this is a consequence of 524H.   If
$\frak A_{a_i}=\{0,a_i\}$, then $\mu_{X_i}$ is purely atomic and there
is a single point $x$ of $X_i$ such that $\mu\{x\}=\mu X_i$ (414G).   In
this case $\Cal N(\mu_{X_i})$ has a greatest member $X_i\setminus\{x\}$,
and similarly $\Cal N(\nu_{\kappa_i})$ has a greatest member, so they have
Tukey equivalent cofinal subsets and are Tukey equivalent (513E(d-ii)).\
\Qed

Now $E\mapsto\familyiI{E\cap X_i}$ is a partially-ordered-set
isomorphism between $\Cal N(\mu)$ and $\prod_{i\in I}\Cal N(\mu_{X_i})$.
Similarly, $\Cal N(\nu)$ is isomorphic to
$\prod_{i\in I}\Cal N(\nu_{Y_i})$.   It now follows from 513Eg
that $\Cal N(\mu)$ and $\Cal N(\nu)$ are Tukey equivalent.   Accordingly
$\add\Cal N(\mu)=\add\Cal N(\nu)$ and $\cf\Cal N(\mu)=\cf\Cal N(\nu)$.
By 521Ad, $\add\mu=\add\Cal N(\mu)$ and $\add\nu=\add\Cal N(\nu)$.

\medskip

{\bf (b)} Immediate from 521La, applied in both directions.
}%end of proof of 524J

\leader{524K}{Corollary} Let $(X,\frak T,\Sigma,\mu)$ and
$(Y,\frak S,\Tau,\nu)$ be Radon measure spaces with measure algebras
$\frak A$, $\frak B$ respectively.   If $\frak A$ can be regularly
embedded in $\frak B$, then $\Cal N(\mu)\prT\Cal N(\nu)$.

\proof{ As usual, write $\bar\mu$ and $\bar\nu$ for the functionals on
$\frak A$, $\frak B$ respectively defined from $\mu$ and $\nu$, and let
$\pi:\frak A\to\frak B$ be a regular embedding, that is, an
order-continuous injective Boolean homomorphism.

\medskip

{\bf (a)} Consider first the case in which $\mu$ is totally finite
and $\pi$ is measure-preserving for $\bar\mu$ and $\bar\nu$.  Let
$(\tilde{X},\tilde{\frak T},\tilde{\Sigma},\tilde{\mu})$ and
$(\tilde{Y},\tilde{\frak S},\tilde{\Tau},\tilde{\nu})$ be the Stone spaces
of $(\frak A,\bar\mu)$ and $(\frak B,\bar\nu)$ respectively.   Then $\pi$
corresponds to a continuous function $f:\tilde Y\to\tilde X$ (312Q).
By 418I, the image measure $\tilde\nu f^{-1}$ is a Radon measure on
$\tilde X$.
If $a\in\frak A$ and $\widehat{a}$ is the corresponding open-and-closed set
in $\tilde X$, then

\Centerline{$\tilde\nu f^{-1}[\widehat{a}]=\tilde\nu(\widehat{\pi a})
=\bar\nu(\pi a)=\bar\mu a=\tilde\mu\widehat{a}$.}

\noindent By 415H(v), $\tilde\nu f^{-1}=\tilde\mu$.   By 521Hb,
$\Cal N(\tilde\mu)\prT\Cal N(\tilde\nu)$.   But now 524Ja tells us that

\Centerline{$\Cal N(\mu)\equivT\Cal N(\tilde\mu)\prT\Cal N(\tilde\nu)
\equivT\Cal N(\nu)$.}

\medskip

{\bf (b)} Next, consider the case in which $\mu$ and $\nu$ are
totally finite but $\pi$ is not necessarily measure-preserving.
As it is (sequentially) order-continuous, we have a measure $\mu'$ on $X$
defined by saying that $\mu'E=\bar\nu(\pi E^{\ssbullet})$ for $E\in\Sigma$,
and $\Cal N(\mu')=\Cal N(\mu)$.   Because $\mu'$ is absolutely continuous
with respect to $\mu$, it is an indefinite-integral
measure over $\mu$ (234O) and is a Radon measure on $X$ (416Sa again).   Taking
$\bar\mu'$ to be the corresponding functional on $\frak A$,
$(\frak A,\bar\mu')$ is the measure algebra of $\mu'$ and $\pi$ is
measure-preserving for $\bar\mu'$ and $\bar\nu$.   So (a) tells us that

\Centerline{$\Cal N(\mu)=\Cal N(\mu')\prT\Cal N(\nu)$.}

\medskip

{\bf (c)} Thirdly, suppose that $\mu$ is totally finite, but $\nu$ might
not be.   Set $\frak B^f=\{b:b\in\frak B$, $\bar\nu b<\infty\}$.   For
$b\in\frak B^f$, set $c_b=\sup\{a:a\in\frak A$, $b\Bcap\pi a=0\}$;  then
$b\Bcap\pi c_b=0$, because $\pi$ is order-continuous.
If $a\in\frak A\setminus\{0\}$, there is a
$b\in\frak B^f$ such that $b\Bcap\pi a\ne 0$, so that $a\notBsubseteq c_b$.
Accordingly $\sup_{b\in\frak B^f}1\Bsetminus c_b=1$ in $\frak A$;
as $\frak A$ is ccc, there is a sequence $\sequencen{b_n}$ in $\frak B^f$ such that
$\sup_{n\in\Bbb N}1\Bsetminus c_{b_n}=1$, that is,
$\bar\nu(a\Bcap\sup_{n\in\Bbb N}b_n)>0$ for every non-zero $a\in\frak A$.

For each $n\in\Bbb N$, choose $F_n\in\Tau$ such that $F_n^{\ssbullet}=b_n$
in $\frak B$, and set $Y'=\bigcup_{n\in\Bbb N}F_n$.   The subspace measure
$\nu_{Y'}$ is $\sigma$-finite, so there is a totally finite measure
$\nuprime$ on $Y'$, an indefinite-integral measure over $\nu_{Y'}$,
with the same null ideal as $\nu_{Y'}$ (use 215B(ix)).
The measures $\nu_{Y'}$ and $\nuprime$ are
both Radon measures (416Rb, 416Sa).   Setting $b=\sup_{n\in\Bbb N}b_n$ in
$\frak B$, the principal ideal $\frak B_b$ can be identified with the
measure algebra of $\nu_{Y'}$ (322I) and $\nuprime$.   Moreover, the map
$a\mapsto b\Bcap\pi a:\frak A\to\frak B_b$ is an injective order-continuous
Boolean homomorphism.   By (b) and 521Fa,

\Centerline{$\Cal N(\mu)\prT\Cal N(\nuprime)
=\Cal N(\nu_{Y'})\prT\Cal N(\nu)$.}

\medskip

{\bf (d)} For the general case, let $\familyiI{a_i}$ be a partition of
unity in $\frak A$ such that $\bar\mu a_i$ is finite for every $i$, and
set $b_i=\pi a_i$ for each $i$, so that $\familyiI{b_i}$ is a partition of
unity in $\frak B$.   As in the proof of 524J, we have corresponding
partitions $\familyiI{X_i}$, $\familyiI{Y_i}$ of $X$, $Y$ into measurable
sets;  as before, 322M tells us that $\Cal N(\mu)$ and $\Cal N(\nu)$ can be
identified with $\prod_{i\in I}\Cal N(\mu_{X_i})$ and
$\prod_{i\in I}\Cal N(\nu_{Y_i})$ respectively.
Now, for each $i$, we can identify the principal ideals
$\frak A_{a_i}$, $\frak B_{b_i}$ with the measure algebras of the subspace
measures $\mu_{X_i}$ and $\nu_{Y_i}$, and $\pi\restrp\frak A_{a_i}$ is an
order-continuous embedding of $\frak A_{a_i}$ in $\frak B_{b_i}$.   So (c)
tells us that $\Cal N(\mu_{X_i})\prT\Cal N(\nu_{Y_i})$.   Accordingly

\Centerline{$\Cal N(\mu)\cong\prod_{i\in I}\Cal N(\mu_{X_i})
\prT\prod_{i\in I}\Cal N(\nu_{Y_i})\cong\Cal N(\nu)$}

\noindent (513Eg again), and the proof is complete.
}%end of proof of 524K

\leader{524L}{}\cmmnt{ So far we have been looking at cardinals
defined from null ideals.   Of course there is an equally important series
based on measurable algebras, which turns out to be similarly strongly
associated with the cardinal functions of the ideals $\Cal N_{\kappa}$.
I have already developed a good deal of the machinery in the arguments
of this section.   But for `linking numbers' we need a new idea, which
is most clearly expressed in the context of homogeneous algebras.

\medskip

\noindent}{\bf Proposition}\cmmnt{ ({\smc Dow \& Stepr\=ans 94})} Let
$\kappa$ be an infinite cardinal.   Then for any $n\ge 2$ the
$n$-linking number $\link_n(\frak B_{\kappa})$ is the least $\lambda$
such that $\kappa\le 2^{\lambda}$.

\proof{ Let $\lambda$ be the least cardinal such that
$\kappa\le 2^{\lambda}$.

\medskip

{\bf (a)} By 514Cb, $\frak B_{\kappa}$ is isomorphic, as partially
ordered set, to a subset of
$\Cal P(\link(\frak B_{\kappa}))$, so we must have

\Centerline{$2^{\link(\frak B_{\kappa})}
\ge\#(\frak B_{\kappa})\ge\kappa$}

\noindent and $\link(\frak B_{\kappa})\ge\lambda$.   It follows at once
that $\link_n(\frak B_{\kappa})\ge\lambda$ for every $n\ge 2$
(511Ia).

\medskip

{\bf (b)} Now let $n\ge 2$, and take an injective function
$\phi:\kappa\to\{0,1\}^{\lambda}$.
Let $\Cal C$ be the family of measurable cylinders in
$\{0,1\}^{\kappa}$, that is, sets of the form
$\{x:x\in\{0,1\}^{\kappa},\,x\restr I=z\}$, where $I\subseteq\kappa$ is
finite and $z\in\{0,1\}^I$.    For each
$E\in\Tau_{\kappa}\setminus\Cal N_{\kappa}$ we can find
disjoint finite sets $I'_E$, $I''_E$, $J_E\subseteq\kappa$ and
$G_E\in\Tau_{\kappa}$ such that

\inset{setting $C_E=\{x:x\in\{0,1\}^{\kappa}$, $x(\xi)=0$ for $\xi\in I'_E$
and $x(\xi)=1$ for $\xi\in I''_E\}$, and $k_E=\#(I'_E)+\#(I''_E)$,
$\nu_{\kappa}(C_E\setminus E)\le\Bover1{4n}\nu_{\kappa}C_E
=\Bover1{4n}\cdot 2^{-k_E}$;

$G_E$ is determined by coordinates in $J_E$ and
$\nu_{\kappa}(C_E\cap(E\symmdiff G_E))\le\Bover1{4n}\cdot 2^{-nk_E}$;

$\nu_{\kappa}G_E\ge 1-\Bover1{2n}$.}

\noindent\Prf\
By 254Fe, there is a set $W$, expressible as the union of finitely
many measurable cylinders, such that
$\nu_{\kappa}(E\symmdiff W)\le\Bover1{5n}\nu_{\kappa}E$.   Now
$\nu_{\kappa}W\ge\bover9{10}\nu_{\kappa}E$ so
$\nu_{\kappa}(W\setminus E)\le\Bover1{4n}\nu_{\kappa}W$.   $W$ is
determined by coordinates in a finite set, so is expressible as a
disjoint union of non-empty measurable cylinders, and for at least one
of these we must have
$\nu_{\kappa}(C\setminus E)\le\Bover1{4n}\nu_{\kappa}C$;  take such a
one for $C_E$.
Express $C_E$ as $\{x:x\restr I_E=z_E\}$, where $I_E\subseteq\kappa$ is
finite and $z_E\in\{0,1\}^{I_E}$, and set
$I'_E=\{\xi:\xi\in I_E,\,z_E(\xi)=0\}$ and
$I''_E=\{\xi:\xi\in I_E,\,z_E(\xi)=1\}$;  then $\nu_{\kappa}C_E=2^{-k_E}$
and $\nu_{\kappa}(C_E\setminus E)\le\Bover14n\cdot 2^{-k_E}$.

Next, take a set $W'\subseteq\{0,1\}^{\kappa}$, determined by
coordinates in a finite subset $J$ of $\kappa$, such that
$\nu_{\kappa}(E\symmdiff W')\le\Bover1{4n}\cdot 2^{-nk_E}$.   Set

\Centerline{$G_E
=\{x:x\in\{0,1\}^{\kappa},\,\Exists y\in W'\cap C_E,\,
  x\restr\kappa\setminus I_E=y\restr\kappa\setminus I_E\}$,}

\noindent so that $G_E$ is determined by coordinates in
$J_E=J\setminus I_E$ and $G_E\cap C_E=W'\cap C_E$;  accordingly

\Centerline{$\nu_{\kappa}(C_E\cap(E\symmdiff G_E))
=\nu_{\kappa}(C_E\cap(E\symmdiff W'))
\le\nu_{\kappa}(E\symmdiff W')\le\Bover1{4n}\cdot 2^{-nk_E}$.}

Note that $G_E$ and $C_E$ are stochastically independent, so that

$$\eqalign{\nu_{\kappa}C_E(1-\nu_{\kappa}G_E)
&=\nu_{\kappa}(C_E\setminus G_E)
\le\nu_{\kappa}(C_E\setminus E)+\nu_{\kappa}(C_E\cap(E\setminus G_E))\cr
&\le\Bover1{4n}\nu_{\kappa}C_E+\Bover1{4n}(\nu_{\kappa}C_E)^n
\le\Bover1{2n}\nu_{\kappa}C_E\cr}$$

\noindent and $\nu_{\kappa}G_E\ge 1-\Bover1{2n}$.\ \Qed

\medskip

{\bf (c)} Let $Q$ be the set of all quadruples $(k,U,V,W)$ where
$k\in\Bbb N$ and $U$, $V$, $W$ are disjoint open-and-closed subsets of
$\{0,1\}^{\lambda}$ in its usual topology.   For $q=(k,U,V,W)\in Q$, set

\Centerline{$\Cal E_q
=\{E:E\in\Tau_{\kappa}\setminus\Cal N_{\kappa},\,
k_E=k,\,\phi[I'_E]\subseteq U,\,
\phi[I''_E]\subseteq V,\,\phi[J_E]\subseteq W\}$.}

\noindent For any $E\in\Tau_{\kappa}\setminus\Cal N_{\kappa}$, $I'_E$,
$I''_E$ and $J_E$, as chosen in (b) above, are disjoint finite sets, so
$\phi[I'_E]$,
$\phi[I''_E]$ and $\phi[J_E]$ also are, and there is a $q\in Q$ such
that $E\in\Cal E_q$.   Now if $q=(k,U,V,W)\in Q$ and $E_i\in\Cal E_q$
for $i<n$, then $\nu_{\kappa}(\bigcap_{i<n}E_i)>0$.   \Prf\ Set
$I'=\bigcup_{i<n}I'_{E_i}$, $I''=\bigcup_{i<n}I''_{E_i}$ and
$J=\bigcup_{i<n}J_{E_i}$.   Then $\phi[I']\subseteq U$,
$\phi[I'']\subseteq V$ and $\phi[J]\subseteq W$, so that $I'$, $I''$ and
$J$ must be disjoint.   Set

\Centerline{$C=\bigcap_{i<n}C_{E_i}
=\{x:x\in\{0,1\}^{\kappa},\,x(\xi)=0$ for $\xi\in I'$, $x(\xi)=1$ for
$\xi\in I''\}$;}

\noindent  then $\nu_{\kappa}C=2^{-\#(I'\cup I'')}\ge 2^{-nk}$.
Next, setting $G=\bigcap_{i<n}G_{E_i}$,

\Centerline{$\nu_{\kappa}G\ge 1-\sum_{i=0}^{n-1}(1-\nu_{\kappa}G_{E_i})
\ge\Bover12$,}

\noindent and $G$ is stochastically independent of $C$, so that
$\nu_{\kappa}(C\cap G)\ge 2^{-nk-1}$.   Finally,

\Centerline{$\nu_{\kappa}(C\cap G\setminus E_i)
\le\nu_{\kappa}(C_{E_i}\cap G_{E_i}\setminus E_i)
\le\Bover1{4n}\cdot 2^{-nk}$}

\noindent for each $i$, so

\Centerline{$\nu_{\kappa}(C\cap G\setminus\bigcap_{i<n}E_i)
\le 2^{-nk-2}<\nu_{\kappa}(C\cap G)$}

\noindent and $\nu_{\kappa}(\bigcap_{i<n}E_i)>0$.\ \Qed

\medskip

{\bf (d)} This means that if we set $A_q=\{E^{\ssbullet}:E\in\Cal E_q\}$
for each $q\in Q$, then every $A_q$ is an $n$-linked set in
$\frak B_{\kappa}$ and $\bigcup_{q\in Q}A_q=\frak B_{\kappa}^+$.
Because $\{0,1\}^{\lambda}$ is a compact topological space with a
subbase with cardinal
$\lambda\ge\omega$, it has $\lambda$ open-and-closed sets and
$\#(Q)=\lambda$.   So $\family{q}{Q}{A_q}$ witnesses that
$\link_n(\frak B_{\kappa})\le\lambda$, and the proof is complete.
}%end of proof of 524L

\leader{524M}{\bf Theorem} Let $(\frak A,\bar\mu)$ be a semi-finite
measure algebra.   Let $K$ be the set of infinite
cardinals $\kappa$ such that $\frak A$ has a homogeneous principal ideal
with Maharam type $\kappa$.

\vskip-\baselineskip

$$\leqalignno{\#(\frak A)
&=2^{c(\frak A)}\text{ if }\frak A\text{ is finite},
  &\indent\text{(a)}\cr
&=\tau(\frak A)^{\omega}\text{ if $\frak A$ is ccc and infinite}.\cr
\wdistr(\frak A)
&=\infty\text{ if }\frak A\text{ is purely atomic},
  &\indent\text{(b)}\cr
&=\add\Cal N\text{ if }K=\{\omega\},\cr
&=\omega_1\text{ otherwise}.\cr
\pi(\frak A)
&=c(\frak A)\text{ if }\frak A\text{ is purely atomic},
  &\indent\text{(c)}\cr
&=\max(c(\frak A),\cf\Cal N,\sup_{\kappa\in K}\cff[\kappa]^{\le\omega})
   \text{ otherwise}.\cr
\frak m(\frak A)
&=\infty\text{ if }\frak A\text{ is purely atomic},
  &\indent\text{(d)}\cr
&=\min_{\kappa\in K}\cov\Cal N_{\kappa}\text{ otherwise}.\cr
d(\frak A)
&=c(\frak A)\text{ if }\frak A\text{ is purely atomic},
  &\indent\text{(e)}\cr
&=\max(c(\frak A),\sup_{\kappa\in K}\non\Cal N_{\kappa})
  \text{ otherwise}.\cr
&&\indent\text{(f) For }2\le n<\omega,\cr
\link_n(\frak A)
&=c(\frak A)\text{ if }\frak A\text{ is purely atomic},\cr
&=\max(c(\frak A),\min\{\lambda:\tau(\frak A)\le 2^{\lambda}\})
   \text{ otherwise}.\cr}$$

\proof{ The case $\frak A=\{0\}$ is trivial, so I shall assume
henceforth that $\frak A\ne\{0\}$.   Let $\familyiI{a_i}$ be a partition
of unity in $\frak A^+$ such that all the principal ideals
$\frak A_{a_i}$ are homogeneous and totally finite.   For each $i\in I$,
set $\kappa_i=\tau(\frak A_{a_i})$, so that
$\frak A_{a_i}\cong\frak B_{\kappa_i}$, and let $(Z_i,\lambda_i)$ be the
Stone space of $(\frak A_{a_i},\bar\mu\restrp\frak A_{a_i})$.
Let $(\widehat{\frak A},\tilde\mu)$ be the localization of
$(\frak A,\bar\mu)$ (322Q).
$\frak A$ can be identified with
an order-dense Boolean subalgebra of $\widehat{\frak A}$, so that
$\familyiI{a_i}$ is still a partition of unity in $\widehat{\frak A}$.
Because $\frak A^f=\widehat{\frak A}^f$ (322P),
$\frak A_{a_i}$ is still
a principal ideal of $\widehat{\frak A}$, and $\widehat{\frak A}$ can be
identified with the simple product $\prod_{i\in I}\frak A_{a_i}$ (315F).

\medskip

{\bf (a)} This is elementary if $\frak A$ is finite (see 511Ic).   If
$\frak A$ is infinite, then 515Ma tells us that
$\#(\frak A)=\tau(\frak A)^{\omega}$.

\medskip

{\bf (b)}

$$\eqalignno{\wdistr(\frak A)
&=\wdistr(\widehat{\frak A})\cr
\displaycause{514Ee}
&=\min_{i\in I}\wdistr(\frak A_{a_i})\cr
\displaycause{514Ef}
&=\min_{i\in I}\wdistr(\frak B_{\kappa_i})
=\min_{i\in I}\add(\Cal N(\lambda_i))\cr
\displaycause{514Be, because $\Cal N(\lambda_i)$ is the ideal of nowhere
dense subsets of $Z$, by 322R}
&=\min_{i\in I}\add(\Cal N_{\kappa_i})\cr
\displaycause{524Ja}
&=\infty\text{ if }K=\emptyset,\cr
&=\add\Cal N\text{ if }K=\{\omega\},\cr
&=\omega_1\text{ otherwise}\cr}$$

\noindent (523E).

\medskip

{\bf (c)(i)} Consider first an algebra $\frak B_{\kappa}$, where
$\kappa\ge\omega$.   Then $\ci\frak B_{\kappa}^+>\omega$.   \Prf\ If
$\sequencen{b_n}$ is any sequence in $\frak B_{\kappa}^+$, then (because
$\frak B_{\kappa}$ is atomless) we can choose $c_n\Bsubseteq b_n$ such
that $0<\bar\nu_{\kappa}c_n\le 2^{-n-2}$ for each $n\in\Bbb N$.   Set
$c=\sup_{n\in\Bbb N}c_n$, $b=1\Bsetminus c$;  then $b\ne 0$ and
$b_n\notBsubseteq b$ for every $n$, so $\{b_n:n\in\Bbb N\}$ is not
coinitial with $\frak B_{\kappa}^+$.\ \Qed

It follows that

$$\eqalignno{\ci\frak B_{\kappa}^+
&=\cov(\frak B_{\kappa}^+,\Bsupseteqshort,\frak B_{\kappa}^+)
=\cov(\frak B_{\kappa}^+,\Bsupseteqshort\strprime,
  [\frak B_{\kappa}^+]^{\le\omega})\cr
\displaycause{512Gf}
&=\cov(\Cal N_{\kappa},\subseteq,\Cal N_{\kappa})\cr
\displaycause{524H}
&=\cf\Cal N_{\kappa}.\cr}$$

\medskip

\quad{\bf (ii)} If $\frak A$ is purely atomic, then
$\frak A_{a_i}=\{0,a_i\}$ for every
$i$, and $\pi(\frak A)=\#(I)=c(\frak A)$.   Otherwise,

$$\eqalignno{\max(c(\frak A),\sup_{i\in I}\pi(\frak A_{a_i}))
&\le\pi(\frak A)\cr
\displaycause{514Da, 514Ed}
&\le\max(\omega,c(\frak A),\sup_{i\in I}\pi(\frak A_{a_i}))\cr
\displaycause{514Ef}
&=\max(c(\frak A),\sup_{\kappa\in K}\pi(\frak B_{\kappa}))
=\max(c(\frak A),\sup_{\kappa\in K}\cf\Cal N_{\kappa})\cr
\displaycause{by (i)}
&=\max(c(\frak A),\cf\Cal N,
  \sup_{\kappa\in K}\cff[\kappa]^{\le\omega}))\cr}$$

\noindent by 523N.

\medskip

{\bf (d)} If $\frak A$ is purely atomic, then $\frak m(\frak A)=\infty$
(511If).   Otherwise,

$$\eqalignno{\frak m(\frak A)
&=\frak m(\widehat{\frak A})\cr
\displaycause{517Id}
&=\min_{i\in I}\frak m(\frak A_{a_i})
=\min_{i\in I}n(Z_i)\cr
\displaycause{517N}
&=\min_{i\in I}\cov\Cal N(\lambda_i)\cr
\displaycause{again because $\Cal N(\lambda_i)$ is the ideal of nowhere
dense subsets of $Z_i$}
&=\min_{i\in I}\cov\Cal N_{\kappa_i}\cr
\displaycause{524Jb}
&=\min_{\kappa\in K}\cov\Cal N_{\kappa},\cr}$$

\noindent as claimed.

\medskip

{\bf (e)(i)} I note first that $d(\frak A_{a_i})=\non\Cal N_{\kappa_i}$
for each $i$.   \Prf\ Let
$A\in\Cal PZ_i\setminus\Cal N(\lambda_i)$ be a set with cardinal
$\non\Cal N(\lambda_i)$.   Then $H=\interior\overline{A}$ is not empty.
Let $a\in\frak A_{a_i}^+$ be such that the corresponding open-and-closed
set $\widehat{a}$ is included in $H$.   Then $\widehat{a}$ can be
identified with the Stone space of $\frak A_a$ (312T);   because
$\frak A_{a_i}$ is homogeneous, and $A\cap\widehat{a}$ is dense in
$\widehat{a}$,

$$\eqalignno{d(\frak A_{a_i})
&=d(\frak A_a)=d(\widehat{a})\cr
\displaycause{514Bd}
&\le\#(A\cap\widehat{a})
\le\non\Cal N(\lambda_i)
=\non\Cal N_{\kappa_i}\cr
\displaycause{524Jb}
&\le d(Z_i)\cr
\displaycause{because $\Cal N(\lambda_i)$ is the ideal of nowhere dense
subsets of $Z_i$, so surely contains no dense set}
&=d(\frak A_{a_i})\cr}$$

\noindent by 514Bd again.\ \Qed

\wheader{524M}{6}{2}{2}{48pt}

\quad{\bf (ii)} If $\frak A$ is purely atomic, $d(\frak A)=c(\frak A)$.
Otherwise,

$$\eqalignno{\max(c(\frak A),\sup_{i\in I}d(\frak A_{a_i}))
&\le d(\frak A)\cr
\displaycause{514Da, 514Ed}
&=d(\widehat{\frak A})\cr
\displaycause{514Ee}
&\le\max(\omega,c(\frak A),\sup_{i\in I}d(\frak A_{a_i}))\cr
\displaycause{514Ef}
&=\max(c(\frak A),\sup_{i\in I}\non\Cal N_{\kappa_i})
=\max(c(\frak A),\sup_{\kappa\in K}\non\Cal N_{\kappa}).\cr}$$

\medskip

{\bf (f)} If $\frak A$ is purely atomic, this is elementary, since any
linked subset of $\frak A^+$ can contain at most one atom.   Otherwise,
set

\Centerline{$\theta
=\max(c(\frak A),\min\{\lambda:\tau(\frak A)\le 2^{\lambda}\})$,
\quad$\theta'=\link_n(\frak A)$.}

\noindent For any $i\in I$, $\kappa_i\le\tau(\frak A)$
(514Ed), so $\kappa_i\le 2^{\theta}$ and
$\link_n(\frak A_{a_i})=\link_n(\frak B_{\kappa_i})\le\theta$ (524L;  of
course the case $\kappa_i=0$ is trivial here).   Accordingly

$$\eqalignno{\theta'
&=\text{link}_n(\widehat{\frak A})\cr
\displaycause{514Ee}
&\le\max(\omega,c(\frak A),\sup_{i\in I}\text{link}_n(\frak A_{a_i}))\cr
\displaycause{514Ef}
&\le\theta.\cr}$$

On the other hand, $c(\frak A)\le\theta'$ (514Da).   For each $i\in I$,
$\link_n(\frak A_i)\le\theta'$ (514Ed), so $\kappa_i\le 2^{\theta'}$
(524L, in the other direction).   Let $A_i$ be a $\tau$-generating
subset of $\frak A_{a_i}$ with cardinal $\kappa_i$.   Now the order-closed
subalgebra of $\frak A$ generated by
$A=\{a_i:i\in I\}\cup\bigcup_{i\in I}A_i$ is $\frak A$, so

\Centerline{$\tau(\frak A)\le\#(A)
=\max(c(\frak A),\sup_{i\in I}\kappa_i)
\le\max(\theta',2^{\theta'})=2^{\theta'}$.}

\noindent But this means that $\theta\le\theta'$ and the two are equal.
}%end of proof of 524M

\cmmnt{\medskip

\noindent{\bf Remark} For the corresponding calculation of $\tau(\frak A)$,
when $(\frak A,\bar\mu)$ is localizable, see 332S.
}

\leader{524N}{Corollary (a)} If $(X,\Sigma,\mu)$ is a semi-finite locally
compact measure space, with $\mu X>0$, then
$\cov\Cal N(\mu)\ge\frak m_{\sigma\text{-linked}}$.

(b) If $\frak A$ is any measurable algebra, then
$\frak m(\frak A)\ge\frak m_{\sigma\text{-linked}}$.

\proof{{\bf (a)} Because $\mu$ is semi-finite and $\mu X>0$, there is an
$E\in\Sigma$ such that $0<\mu E<\infty$.   The subspace measure $\mu_E$ on
$E$ is compact, so $\nu=\Bover1{\mu E}\mu_E$ is a
compact probability measure.   Set $\kappa=\max(\omega,\tau(\nu))$.
Because $\nu$ is a compact measure, there is
a function $f:\{0,1\}^{\kappa}\to E$ which is \imp\ for $\nu_{\kappa}$ and
$\nu$ (343Cd).   Now

$$\eqalignno{\frak m_{\sigma\text{-linked}}
&\le\frak m(\frak B_{\frakc})\cr
\displaycause{because $\frak B_{\frakc}$ is $\sigma$-linked, by 524Mf}
&=\cov\Cal N_{\frakc}\cr
\displaycause{524Md}
&\le\cov\Cal N_{\kappa}\cr
\displaycause{523F}
&\le\cov\Cal N(\nu)\cr
\displaycause{521Ha}
&=\cov\Cal N(\mu_E)
\le\cov\Cal N(\mu)\cr}$$

\noindent (521Fb).

\medskip

{\bf (b)} This is now immediate from 524Md.
}%end of proof of 524N

\leader{524O}{Freese-Nation \dvrocolon{numbers}}\cmmnt{ I spell out
those facts about Freese-Nation numbers of measure algebras which can be
read off from the results in \S518.

\medskip

\noindent}{\bf Proposition} (a) Let $(\frak A,\bar\mu)$ be an infinite
measure algebra.   Then $\FN(\frak A)\ge\FN(\Cal P\Bbb N)$.

(b) Let $\frak A$ be a measurable algebra.

\quad(i) $\FN(\frak A)\le\frak c^+$.

\quad(ii) If $\tau(\frak A)\le\frak c$ then
$\FN(\frak A)\le\FN(\Cal P\Bbb N)$.

\quad(iii) If

\inset{($\alpha$) $\cf([\lambda]^{\le\omega})\le\lambda^+$ for every
cardinal $\lambda\le\tau(\frak A)$,

($\beta$) $\square_{\lambda}$ is true for every uncountable
cardinal $\lambda\le\tau(\frak A)$ of countable cofinality,}

\noindent then $\FN(\frak A)\le\FN^*(\Cal P\Bbb N)$.

(c) Suppose that the continuum hypothesis and
CTP$(\omega_{\omega+1},\omega_{\omega})$ are both true.   If $\frak A$ is
a measurable algebra, then

$$\eqalign{\FN(\frak A)&=\frak c=\omega_1
   \text{ if }\omega\le\tau(\frak A)<\omega_{\omega},\cr
&=\frak c^+=\omega_2\text{ otherwise }.\cr}$$

\proof{{\bf (a)} This is a special case of 518Ca.

\medskip

{\bf (b)(i)} Consider first the case $\frak A=\frak B_{\kappa}$ for some
cardinal $\kappa$.
For $I\subseteq\kappa$, let $\frak C_I$ be the closed subalgebra of
$\frak B_{\kappa}$ consisting of those
$a\in\frak B_{\kappa}$ expressible in the form $E^{\ssbullet}$ for some
measurable $E\subseteq\{0,1\}^{\kappa}$ determined by coordinates in
$I$.   For
$a\in\frak B_{\kappa}$, there is a smallest subset $I_a$ of $\lambda$
such that $a\in\frak C_I$ (325M again);  $I_a$ is always countable.

For each $a\in\frak B_{\kappa}$, set

\Centerline{$f(a)=\{b:I_b\subseteq I_a\}$.}

\noindent Then $\#(f(a))\le\frak c$.   If $a\Bsubseteq b$, then there is
a $c\in\frak B_{\kappa}$ such that
$a\Bsubseteq c\Bsubseteq b$ and $I_c\subseteq I_a\cap I_b$
(325M(b-ii)).   So $f$ is a Freese-Nation function.   This shows that
$\FN(\frak B_{\kappa})\le\frak c^+$.

In general, $\frak A$ is either $\{0\}$ or isomorphic to
a closed subalgebra of $\frak B_{\kappa}$ where
$\kappa=\max(\omega,\tau(\frak A))$, so
$\FN(\frak A)\le\FN(\frak B_{\kappa})\le\frak c^+$ by 518Cc.

\medskip

\quad{\bf (ii)}  $\frak A$ is
$\sigma$-linked (524Mf), so 518D(iii) tells us that
$\FN(\frak A)\le\FN(\Cal P\Bbb N)$.

\medskip

\quad{\bf (iii)} If $\frak B\subseteq\frak A$ is a countably generated
order-closed subalgebra, then $\FN(\frak B)\le\FN(\Cal P\Bbb N)$, by
(ii);  so 518I tells us that $\FN(\frak A)\le\FN^*(\Cal P\Bbb N)$.

\medskip

{\bf (c)} If $\tau(\frak A)<\omega_{\omega}$ then
$\cff[\lambda]^{\le\omega}=\lambda$ for
$\omega_1\le\lambda\le\tau(\frak A)$ (5A1F(e-iv)),
so we can use (a) and (b-iii);  otherwise use (b-i) and 518K.
}%end of proof of 524O

\leader{524P}{The Maharam \dvrocolon{classification}}\cmmnt{ If the
cardinal functions of a Radon measure space are determined by its
measure algebra, there ought to be some way of calculating them directly
from the classification of measure algebras in \S332.   In many cases
this is straightforward.

\wheader{524P}{6}{2}{2}{72pt}
\noindent}{\bf Theorem} Let $(X,\frak T,\Sigma,\mu)$ be a
Radon measure space, and $\frak A$ its measure algebra.
Let $K$ be the set of infinite cardinals $\kappa$ such that the
Maharam-type-$\kappa$ component of $\frak A$ is non-zero.

\vskip-\baselineskip

\newif\iftempa\tempafalse
\newif\iftempb\tempbfalse
\newif\iftempc\tempcfalse
\newif\iftempd\tempdfalse
\newif\iftempe\tempefalse
\pagetotalplus=\the\pagetotal
\advance\pagetotalplus by 120pt
\ifdim\pagetotalplus>\pagegoal\tempatrue\else
\advance\pagetotalplus by 36pt
\ifdim\pagetotalplus>\pagegoal\tempbtrue\else
\advance\pagetotalplus by 36pt
\ifdim\pagetotalplus>\pagegoal\tempctrue\else
\advance\pagetotalplus by 36pt
\ifdim\pagetotalplus>\pagegoal\tempdtrue\else
\advance\pagetotalplus by 36pt
\ifdim\pagetotalplus>\pagegoal\tempetrue
\fi\fi\fi\fi\fi
%\showthe\pagetotal   \showthe\pagetotalplus  \showthe\pagegoal

$$\leqalignno{\add\mu=\add\Cal N(\mu)&=\infty\text{ if }K=\emptyset,
  &\indent\text{(a)}\cr
&=\add\Cal N\text{ if }K=\{\omega\},\cr
&=\omega_1\text{ otherwise}.\cr
\iftempa\noalign{\break}\fi
\pi(\mu)=\pi(\frak A)
&=c(\frak A)\text{ if }K=\emptyset,
  &\indent\text{(b)}\cr
&=\max(c(\frak A),\cf\Cal N,\sup_{\kappa\in K}\cff[\kappa]^{\le\omega})
\text{ otherwise}.\cr
\iftempb\noalign{\break}\fi
\cov\Cal N(\mu)
&=1\text{ if }\frak A=\{0\},
  &\indent\text{(c)}\cr
&=\infty\text{ if }\frak A\text{ has an atom},\cr
&=\cov\Cal N_{\min K}\text{ otherwise}.\cr
\iftempc\noalign{\break}\fi
\non\Cal N(\mu)
&=\infty\text{ if }\frak A=\{0\},
  &\indent\text{(d)}\cr
&=1\text{ if }\frak A\text{ has an atom},\cr
&=\non\Cal N_{\min K}\text{ otherwise}.\cr
\iftempd\noalign{\break}\fi
\shr\Cal N(\mu)
&=0\text{ if }\frak A=\{0\},
  &\indent\text{(e)}\cr
&=1\text{ if }\frak A\text{ has an atom},\cr
&\ge\shr\Cal N\text{ otherwise}.\cr
\iftempe\noalign{\break}\fi
&&\indent\text{(f) If }\mu\text{ is }\sigma\text{-finite},\cr
\cf\Cal N(\mu)
&=1\text{ if }K=\emptyset,\cr
&=\max(\cf\Cal N,\cff[\tau(\frak A)]^{\le\omega})\text{ otherwise}.\cr}$$

\proof{ If $\mu X=0$ all these results are trivial, so let us suppose
henceforth that $\mu X>0$.   As in part (a) of the proof of 524J, there
is a decomposition $\familyiI{X_i}$ of $X$ such that the subspace
measures $\mu_{X_i}$ are all \Mth\ and non-zero.   Note that
$\max(\omega,\#(I))=\max(\omega,c(\frak A))$ (332E).   For each
$i\in I$, let $\kappa_i$ be the Maharam type of $\mu_{X_i}$.

\medskip

{\bf (a)} By 521Ad, $\add\mu=\add\Cal N(\mu)$.
The map $E\mapsto\familyiI{E\cap X_i}$ identifies
$\Cal N(\mu)$, as partially ordered set, with the product of the family
$\familyiI{\Cal N(\mu_{X_i})}$.   So
$\add\Cal N(\mu)=\min_{i\in I}\add\Cal N(\mu_{X_i})$ (511Hg).
Now if $i\in I$ and
$\kappa_i=0$, $X_i$ is an atom of $(X,\Sigma,\mu)$,
so there is an $x_i\in X_i$ such that $\mu(X_i\setminus\{x_i\})=0$
(414G again).
In this case, $X_i\setminus\{x_i\}$ is the largest member of
$\Cal N(\mu_{X_i})$ and $\add\Cal N(\mu_{X_i})=\infty$.   If $\kappa_i$
is infinite, then $\add\Cal N(\mu_{X_i})=\add\Cal N_{\kappa_i}$, by 524I
applied to a scalar multiple of $\mu_{X_i}$.   So
$\add\Cal N(\mu)=\min_{\kappa\in K}\add\Cal N_{\kappa}$,
interpreting this as $\infty$ if $K=\emptyset$.   But we know from 523E
that $\add\Cal N_{\kappa}=\omega_1$ if $\kappa>\omega$, while of course
$\add\Cal N_{\omega}=\add\Cal N$.   It follows at once that

$$\eqalignno{\add\Cal N(\mu)=\min_{i\in I}\add\Cal N(\mu_{X_i})
&=\infty\text{ if }K=\emptyset,\cr
&=\add\Cal N\text{ if }K=\{\omega\},\cr
&=\omega_1\text{ otherwise}.\cr}$$

\medskip

{\bf (b)} By 521Dd, $\pi(\mu)=\pi(\frak A)$;
and 524Mc gives us the formula for $\pi(\frak A)$.

\medskip

{\bf (c)} If $\Cal E$ is a cover of $X$ by negligible sets, and
$i\in I$, then
$\{E\cap X_i:E\in\Cal E\}$ is a cover of $X_i$ by negligible sets;  thus
$\cov\Cal N(\mu)\ge\sup_{i\in I}\cov\Cal N(\mu_{X_i})$.   By 524Jb,
$\cov\Cal N(\mu)\ge\sup_{i\in I}\cov\Cal N_{\kappa_i}$.   If any of the
$\kappa_i$ is zero, that is, if $\frak A$ has an atom, this is $\infty$,
and we can stop.

Otherwise, for each $i\in I$,

\Centerline{$\cov\Cal N(\mu_{X_i})=\cov\Cal N_{\kappa_i}
\le\cov\Cal N_{\min K}=\lambda$}

\noindent say, by 523B.   So we have a family
$\ofamily{\xi}{\lambda}{E_{i\xi}}$ of negligible subsets of $X_i$
covering $X_i$;  setting $E_{\xi}=\bigcup_{i\in I}E_{i\xi}$ for each
$\xi$, we have a family $\ofamily{\xi}{\lambda}{E_{\xi}}$ in
$\Cal N(\mu)$ covering $X$, so $\cov\Cal N(\mu)\le\cov\Cal N_{\min K}$.
But we already know that

\Centerline{$\cov\Cal N(\mu)\ge\sup_{i\in I}\cov\Cal N_{\kappa_i}
\ge\cov\Cal N_{\min K}$,}

\noindent so $\cov\Cal N(\mu)=\cov\Cal N_{\min K}$.

\medskip

{\bf (d)} A set $A\subseteq X$ is non-negligible iff $A\cap X_i$ is
non-negligible for some $i\in I$.   It follows at once that
$\non\Cal N(\mu)=\min_{i\in I}\non\Cal N(\mu_{X_i})$.
If any of the $X_i$ is an atom, it contains a point of non-zero measure,
so that $\non\Cal N(\mu)=1$.   If $\kappa_i\ge\omega$ for every $i$,
then we have

\Centerline{$\non\Cal N(\mu)=\min_{i\in I}\non\Cal N_{\kappa_i}
=\non\Cal N_{\min K}$}

\noindent by 524Jb and 523B again.

\medskip

{\bf (e)} If $\frak A$ is purely atomic, then $\mu$ is point-supported, so
$\shr\Cal N(\mu)=1$.   Otherwise,
let $E$ be a measurable set of non-zero finite measure such that the
subspace measure $\mu_E$ is atomless;  let $\nu$ be the normalized subspace
measure $\Bover1{\mu E}\mu_E$;  then $\nu$, like $\mu_E$, is a Radon
measure.   By 343Cb, there is a
function $f:E\to\{0,1\}^{\omega}$ which
is \imp\ for $\nu$ and $\nu_{\omega}$;   because $\{0,1\}^{\omega}$ is
separable and metrizable, $\nu f^{-1}$ is a Radon measure
(451O, or
418I-418J) and must be equal to $\nu_{\omega}$ (416Eb).   By 521Fd and
521Hb,

\Centerline{$\shr\Cal N(\mu)\ge\shr\Cal N(\mu_E)=\shr\Cal N(\nu)
\ge\shr\Cal N(\nu_{\omega})=\shr\Cal N$.}

\medskip

{\bf (f)(i)} If $K=\emptyset$ then (a) tells us that $\Cal N(\mu)$ has a
greatest member, so that $\cf\Cal N(\mu)=1$.

\medskip

\quad{\bf (ii)} Now suppose that $K$ is not empty.   Then 524Fb tells us
that there is a family
$\ofamily{\xi}{\tau(\frak A)}{E_{\xi}}$ in $\Cal N(\mu)$ such that
$\{\xi:E_{\xi}\subseteq E\}$ is countable for every $E\in\Cal N(\mu)$.
In this case, $J\mapsto\bigcup_{\xi\in J}E_{\xi}:
   [\tau(\frak A)]^{\le\omega}\to\Cal N(\mu)$ is a Tukey function,
so $\cf\Cal N(\mu)\ge\cff[\tau(\frak A)]^{\le\omega}$.
At the same time, there is an $i\in I$ such that $\kappa_i\ge\omega$.
The identity map from $\Cal N(\mu_{X_i})$
to $\Cal N(\mu)$ is a Tukey function;  but this means that

$$\eqalignno{\cf\Cal N(\mu)
&\ge\cf\Cal N(\mu_{X_i})=\cf\Cal N_{\kappa_i}\cr
\displaycause{524I again}
&\ge\cf\Cal N_{\omega}
=\cf\Cal N\cr}$$

\noindent (523B).   Thus
$\cf\Cal N(\mu)\ge\max(\cf\Cal N,\cff[\tau(\frak A)]^{\le\omega})$.

\medskip

\quad{\bf (iii)}
In the other direction, we know from 524H (again, applied to a scalar
multiple of $\mu_{X_i}$) that
$(\Cal N(\mu_{X_i}),\subseteq,\Cal N(\mu_{X_i})
\equivGT(\kappa_i^{\Bbb N},\subseteq^*,\Cal S_{\kappa_i})$ whenever
$\kappa_i$ is infinite.   Now $\tau(\frak A)\ge\kappa_i$, so the maps

\Centerline{identity$:\kappa_i^{\Bbb N}\to\tau(\frak A)^{\Bbb N}$,
\quad$S\mapsto S\cap(\Bbb N\times\kappa_i):
  \Cal S_{\tau(\frak A)}\to\Cal S_{\kappa_i}$}

\noindent form a Galois-Tukey connection from
$(\kappa_i^{\Bbb N},\subseteq^*,\Cal S_{\kappa_i})$ to
$(\tau(\frak A)^{\Bbb N},\subseteq^*,\Cal S_{\tau(\frak A)})$.
Accordingly we have

$$\eqalign{(\Cal N(\mu_{X_i}),\subseteq,\Cal N(\mu_{X_i}))
&\equivGT(\kappa_i^{\Bbb N},\subseteq^*,\Cal S_{\kappa_i})\cr
&\prGT(\tau(\frak A)^{\Bbb N},\subseteq^*,\Cal S_{\tau(\frak A)})
\equivGT(\Cal N_{\tau(\frak A)},\subseteq,\Cal N_{\tau(\frak A)}),\cr}$$

\noindent and $\Cal N(\mu_{X_i})\prT\Cal N_{\tau(\frak A)}$.

The arguments quoted assume that $\kappa_i$ is infinite;  but of course
it is still true that
$\Cal N(\mu_{X_i})\prT\Cal N_{\tau(\frak A)}$ when $\kappa_i=0$, since
then any constant function from $\Cal N(\mu_{X_i})$ to
$\Cal N_{\tau(\frak A)}$ is a Tukey function.   It follows that

\Centerline{$\Cal N(\mu)\cong\prod_{i\in I}\Cal N(\mu_{X_i})
\prT\Cal N_{\tau(\frak A)}^I$}

\noindent (513Eg once more).

\medskip

\quad{\bf (iv)}
At this point observe that as we are assuming that $K\ne\emptyset$,
$\tau(\frak A)$ is infinite;  and as $\mu$ is supposed to be
$\sigma$-finite, $I$ is countable.   So we can find a disjoint family
$\familyiI{F_i}$ of measurable subsets of $\{0,1\}^{\tau(\frak A)}$ such
that all the subspace measures $(\nu_{\tau(\frak A)})_{F_i}$ are
isomorphic to scalar multiples of $\nu_{\tau(\frak A)}$.   (Take
$F_i=\{x:x(n_i)=1$, $x(m)=0$ for $m<n_i\}$ where
$i\mapsto n_i:I\to\Bbb N$ is injective.)   In this case, the map

\Centerline{$\familyiI{E_i}\mapsto\bigcup_{i\in I}E_i:
\prod_{i\in I}\Cal N((\nu_{\tau(\frak A)})_{F_i})
\to\Cal N(\nu_{\tau(\frak A)})$}

\noindent  is a Tukey function, while $\Cal N_{\tau(\frak A)}^I$ is
isomorphic to
$\prod_{i\in I}\Cal N((\nu_{\tau(\frak A)})_{F_i})$.   Putting these
together,

\Centerline{$\Cal N(\mu)\prT\Cal N_{\tau(\frak A)}^I
\cong\prod_{i\in I}\Cal N((\nu_{\tau(\frak A}))_{F_i})
\prT\Cal N_{\tau(\frak A)}$.}

It follows that

\Centerline{$\cf\Cal N(\mu)\le\cf\Cal N_{\tau(\frak A)}
=\max(\Cal N,\cff[\tau(\frak A)]^{\le\omega})$.}

\noindent So we have inequalities in both directions and
$\cf\Cal N(\mu)=\max(\Cal N,\cff[\tau(\frak A)]^{\le\omega})$, as
claimed.
}%end of proof of 524P

\leader{*524Q}{}\cmmnt{ I do not know how to calculate
$\cf\Cal N(\mu)$ for non-$\sigma$-finite
Radon measures $\mu$ without special
assumptions.   In the presence of GCH, however, we have the following
result.

\medskip

\noindent}{\bf Proposition} Suppose that the generalized continuum
hypothesis is true.    Let $(X,\frak T,\Sigma,\mu)$ be a Radon
measure space and $(\frak A,\bar\mu)$ its measure algebra.   For each
cardinal $\kappa$, write $e_{\kappa}$ for the Maharam-type-$\kappa$
component of $\frak A$, and $\frak C_{\kappa}$ for the principal ideal
of $\frak A$ generated by $\sup_{\kappa'>\kappa}e_{\kappa'}$;  set
$\lambda=\sup\{\kappa:e_{\kappa}\ne 0\}$.   Then
$\cf\Cal N(\mu)=\max(c(\frak C_0)^+,\lambda^+)$ unless
$\lambda>c(\frak C_0)$ and there is some $\gamma<\lambda$ such that
$\cf\lambda>c(\frak C_{\gamma})$, in which case
$\cf\Cal N(\mu)=\lambda$.

\proof{{\bf (a)} Write

$$\eqalign{\theta
&=\lambda\text{ if }\lambda>c(\frak C_0)\text{ and }
\cf\lambda>\min_{\gamma<\lambda}c(\frak C_{\gamma}),\cr
&=\max(\lambda^+,c(\frak C_0)^+)\text{ otherwise}.\cr}$$

\noindent If $\mu$ is purely atomic, it is point-supported, so
$\lambda=0$ and $\frak C_0=\{0\}$ and $\theta=1=\cf\Cal N(\mu)$.   So
let us suppose henceforth that $\mu$ is not purely atomic, that is,
$\frak C_0\ne\{0\}$ and $\lambda\ge\omega$.   As in the proofs of
524J and 524P,
there is a decomposition $\familyiI{X_i}$ of $X$ such that the subspace
measures $\mu_{X_i}$ are all \Mth\ and non-zero.   Let
$\kappa_i$ be the Maharam type of $\mu_{X_i}$ for each $i$, so that
$\lambda=\sup_{i\in I}\kappa_i$.   Now
$\Cal N(\mu)\cong\prod_{i\in I}\Cal N(\mu_{X_i})$ (see the proof of
524Ja).   For $i\in I$, $\cf\Cal N(\mu_{X_i})=1$ if $\kappa_i=0$, and
otherwise is
$\max(\cf\Cal N,\cff[\kappa_i]^{\le\omega})
=\max(\omega_1,\cff[\kappa_i]^{\le\omega})$ (524Ja, 523N).   By 5A6Ab,

$$\eqalign{\cf\Cal N(\mu_{X_i})
&=1\text{ if }\kappa_i=0,\cr
&=\kappa_i\text{ if }\cf\kappa_i>\omega,\cr
&=\kappa_i^+\text{ if }\cf\kappa_i=\omega.\cr}$$

\medskip

{\bf (b)} For each cardinal $\kappa$, set
$J_{\kappa}=\{i:i\in I$, $\cf\Cal N(\mu_{X_i})>\kappa\}$, and set

$$\eqalign{\lambda_1=\sup_{i\in I}\cf\Cal N(\mu_{X_i})
&=\lambda^+\text{ if there is an }i\in I
\text{ such that }\kappa_i=\lambda\text{ and }\cf\kappa_i=\omega,\cr
&=\lambda\text{ otherwise}.\cr}$$

\noindent Then 513J tells us that
if $\lambda_1>\#(J_1)$ and there is some $\gamma<\lambda_1$ such that
$\cf\lambda_1>\#(J_{\gamma})$, then $\cf\Cal N(\mu)=\lambda_1$, and that
otherwise $\cf\Cal N(\mu)=\max(\#(J_1)^+,\lambda_1^+)$.   As we are
supposing that $\mu$ is not purely atomic,
$c(\frak C_0)\ge\omega$ and $c(\frak C_0)=\max(\omega,\#(J_1))$;  also
$\lambda^+\ge\lambda_1\ge\lambda\ge\omega$.

\medskip

\quad{\bf case 1} Suppose $\lambda_1\le\#(J_1)$.   Then $J_1$ is
infinite, so $c(\frak C_0)=\#(J_1)\ge\lambda$, and

\Centerline{$\cf\Cal N(\mu)=\#(J_1)^+=c(\frak C_0)^+=\theta$}

\noindent as required.

\medskip

\quad{\bf case 2} Suppose $\lambda_1>\max(\lambda,\#(J_1))$.   Then
there must be some $i\in I$ such that
$\cf\Cal N(\mu_{X_i})>\lambda\ge\omega$, in which case
$\kappa_i=\lambda$ has countable cofinality and $\lambda_1=\lambda^+$.
In this case, $\cf\lambda_1=\lambda_1>\#(J_1)$, so
$\cf\Cal N(\mu)=\lambda_1$.   If $\gamma<\lambda$, then
$\frak C_{\gamma}$ is non-trivial, and
$\cf\lambda=\omega\le c(\frak C_{\gamma})$;  so

\Centerline{$\theta=\max(\lambda^+,\#(J_1)^+)=\lambda_1
=\cf\Cal N(\mu)$.}

\medskip

\quad{\bf case 3} Suppose $\lambda_1=\lambda>\#(J_1)$ has countable
cofinality.   In this case we must have $\kappa_i<\lambda_1$ for every
$i$, so $\#(J_{\gamma})\ge\omega=\cf\lambda_1$ for every
$\gamma<\lambda_1$, and $\cf\Cal N(\mu)=\lambda_1^+$.   At the same time,
$\cf\lambda=\omega\le c(\frak C_{\gamma})$ for every $\gamma<\lambda$,
so

\Centerline{$\theta
=\max(\lambda^+,\#(J_1)^+)=\max(\lambda_1^+,\#(J_1)^+)=\cf\Cal N(\mu)$.}

\medskip

\quad{\bf case 4} Suppose $\lambda_1=\lambda>\#(J_1)$ has uncountable
cofinality.   In this case we have
$\lambda>\max(\omega,\#(J_1))=c(\frak C_0)$, so

$$\eqalign{\cf\Cal N(\mu)=\lambda_1
&\iff\#(J_{\gamma})<\cf\lambda_1\text{ for some }\gamma<\lambda_1\cr
&\iff\max(\omega,\#(J_{\gamma}))<\cf\lambda_1
   \text{ for some }\gamma<\lambda_1\cr
&\iff c(\frak C_{\gamma})<\cf\lambda
   \text{ for some }\gamma<\lambda\cr
&\iff\theta=\lambda
\iff\theta=\lambda_1,\cr}$$

\noindent and otherwise

\Centerline{$\cf\Cal N(\mu)
=\lambda_1^+=\max(\lambda^+,c(\frak C_0)^+)=\theta$.}

\noindent Thus $\cf\Cal N(\mu)=\theta$ in all cases.
}%end of proof of 524Q

\leader{524R}{}\cmmnt{ The results above show that most of the most
important cardinal functions of measurable algebras and Radon
measures are readily calculable from the cardinal functions of the
ideals $\Cal N_{\kappa}$ studied in \S523.   There are no
such simple formulae for other classes of space such as compact
or quasi-Radon measures (524Xj, 524Xk).    However I can give a handful of partial
results, as follows.

\medskip

\noindent}{\bf Proposition} Let $(X,\Sigma,\mu)$ be a countably compact
$\sigma$-finite measure space with Maharam type $\kappa$.   Then
$[\kappa]^{\le\omega}\prT\Cal N(\mu)$.   Consequently
$\cff[\kappa]^{\le\omega}\le\cf\Cal N(\mu)$, and if $\kappa$ is
uncountable then $\add\Cal N(\mu)=\omega_1$ and
$\cf\Cal N(\mu)\ge\cf\Cal N_{\kappa}$.

\proof{ If $\ofamily{\xi}{\kappa}{E_{\xi}}$ is a family as in 524Fb, then
$I\mapsto\bigcup_{\xi\in I}E_{\xi}:[\kappa]^{\le\omega}\to\Cal N(\mu)$
is a Tukey function, if both $[\kappa]^{\le\omega}$ and $\Cal N(\mu)$
are given their natural partial orderings of inclusion.   By 513Ee,
$\cff[\kappa]^{\le\omega}\le\cf\Cal N(\mu)$ and
$\add[\kappa]^{\le\omega}\ge\add\Cal N(\mu)$.   But if $\kappa$ is
uncountable, $\add[\kappa]^{\le\omega}=\omega_1$ so $\add\Cal N(\mu)$ is
also $\omega_1$.   At the same time, $\cf\Cal N(\mu)\ge\cf\Cal N$ (521K),
so

\Centerline{$\cf\Cal N(\mu)
\ge\max(\cf\Cal N,\cff[\kappa]^{\le\omega})=\cf\Cal N_{\kappa}$.}
}%end of proof of 524R

\leader{524S}{}\cmmnt{ In a different direction, there is something we
can say about quasi-Radon measures.

\medskip

\noindent}{\bf Proposition} Let $(X,\frak T,\Sigma,\mu)$ be a Radon
measure space, with $\mu X>0$, and $(Y,\frak S,\Tau,\nu)$ a quasi-Radon
measure space such that the measure algebras of $\mu$ and $\nu$ are
isomorphic.   Then

(a) $\Cal N(\nu)\prT\Cal N(\mu)$, so
$\add\nu=\add\Cal N(\nu)\ge\add\Cal N(\mu)=\add\mu$ and
$\cf\Cal N(\nu)\le\cf\Cal N(\mu)$;

(b) $(Y,\in,\Cal N(\nu))\prGT(X,\in,\Cal N(\mu))$, so
$\cov\Cal N(\nu)\le\cov\Cal N(\mu)$ and
$\non\Cal N(\nu)\ge\non\Cal N(\mu)$.

\proof{{\bf (a)} Let $(Z,\frak U,\Lambda,\lambda)$ be the Stone space of
the measure algebra $\frak B$ of $(Y,\Tau,\nu)$, and
$R\subseteq Z\times Y$ the relation described in 415Q/416V, so that
$R^{-1}[F]\in\Cal N(\lambda)$ for every $F\in\Cal N(\nu)$.   Let
$W\subseteq Z$ be the union of the open sets of finite measure.   Then
the subspace measure $\lambda_W$ is a Radon measure
and its measure algebra
is isomorphic to the measure algebras of $\nu$ and $\mu$
(411Pf).  %from 2005

Now $F\mapsto W\cap R^{-1}[F]:\Cal N(\nu)\to\Cal N(\lambda_W)$ is a
Tukey function.   \Prf\Quer\ Otherwise, there is a family
$\Cal A\subseteq\Cal N(\nu)$ such that $\bigcup\Cal A\notin\Cal N(\nu)$
but $\{W\cap R^{-1}[A]:A\in\Cal A\}$ is bounded above in
$\Cal N(\lambda_W)$.   Because $W$ is conegligible,
$B=\bigcup_{A\in\Cal A}R^{-1}[A]$ is negligible in $Z$.
Let $E\in\Tau$ be a measurable envelope of $\bigcup\Cal A$ (213J/213L).
Then the open-and-closed set $E^*\subseteq Z$ corresponding to
$E^{\ssbullet}\in\frak B$ is not negligible;  as $\lambda$ is inner
regular with respect to the open-and-closed sets (411Pb), there must be
a non-empty open-and-closed set $V\subseteq E^*$ which is disjoint from
$\bigcup_{A\in\Cal A}R^{-1}[A]$.   Express $V$ as $F^*$ where
$F\in\Tau$.   Then $R[V]=R[F^*]$ is disjoint from $\bigcup\Cal A$.
But $R[F^*]$ is measurable and $F\setminus R[F^*]$ is negligible
(415Qb), while $F\setminus E$ must also be negligible, so
$E\cap R[F^*]$ is a non-negligible measurable subset of
$E\setminus\bigcup\Cal A$, which is impossible.\ \Bang\Qed

This shows that $\Cal N(\nu)\prT\Cal N(\lambda_W)$.   But $\lambda_W$
and $\mu$ are Radon measures with isomorphic non-zero measure algebras,
so $\Cal N(\lambda_W)\equivT\Cal N(\mu)$ (524J) and
$\Cal N(\nu)\prT\Cal N(\mu)$.   Accordingly
$\add\Cal N(\nu)\ge\add\Cal N(\mu)$ and
$\cf\Cal N(\nu)\le\cf\Cal N(\mu)$

\medskip

{\bf (b)}   This is a special case of 521La.
}%end of proof of 524S

\leader{524T}{Corollary} Let $(Y,\frak S,\Tau,\nu)$ be a
quasi-Radon measure space, and $\frak B$ its measure algebra.
Let $K$ be the set of infinite cardinals $\kappa$ such that the
Maharam-type-$\kappa$ component of $\frak B$ is non-zero.

\vskip-\baselineskip

\newif\iftempa\tempafalse
\newif\iftempb\tempbfalse
\newif\iftempc\tempcfalse
\newif\iftempd\tempdfalse
\pagetotalplus=\the\pagetotal
\advance\pagetotalplus by 120pt
\ifdim\pagetotalplus>\pagegoal\tempatrue\else
\advance\pagetotalplus by 36pt
\ifdim\pagetotalplus>\pagegoal\tempbtrue\else
\advance\pagetotalplus by 36pt
\ifdim\pagetotalplus>\pagegoal\tempctrue\else
\advance\pagetotalplus by 36pt
\ifdim\pagetotalplus>\pagegoal\tempdtrue\else
\fi\fi\fi\fi
%\showthe\pagetotal   \showthe\pagetotalplus  \showthe\pagegoal


$$\leqalignno{\add\nu=\add\Cal N(\nu)&=\infty\text{ if }K=\emptyset,
  &\indent\text{(a)}\cr
&\ge\add\Cal N\text{ if }K=\{\omega\}.\cr
\iftempa\noalign{\break}\fi
\pi(\nu)=\pi(\frak B)
&=c(\frak B)\text{ if }K=\emptyset,&\indent\text{(b)}\cr
&=\max(c(\frak B),\cf\Cal N,\sup_{\kappa\in K}\cff[\kappa]^{\le\omega})
\text{ otherwise}.\cr
\iftempb\noalign{\break}\fi
\cov\Cal N(\nu)
&=1\text{ if }\frak B=\{0\},&\indent\text{(c)}\cr
&=\infty\text{ if }\frak B\text{ has an atom},\cr
&\le\cov\Cal N_{\min K}\text{ otherwise}.\cr
\iftempc\noalign{\break}\fi
\non\Cal N(\nu)
&=\infty\text{ if }\frak B=\{0\},&\indent\text{(d)}\cr
&=1\text{ if }\frak B\text{ has an atom},\cr
&\ge\non\Cal N_{\min K}\text{ otherwise}.\cr
\iftempd\noalign{\break}\fi
&&\indent\text{(e) If }\nu\text{ is }\sigma\text{-finite},\cr
\cf\Cal N(\nu)
&=1\text{ if }K=\emptyset,\cr
&\le\max(\cf\Cal N,\cff[\tau(\frak B)]^{\le\omega})
  \text{ otherwise}.\cr}$$

\proof{ Parts (a), (c), (d) and (e) are mostly a matter of putting 524P
and 524S together.   If there are atoms for $\mu$, they may no longer
include singletons of non-zero measure;  but they do include minimal
non-negligible closed sets, so there are non-negligible singletons and
$\cov\Cal N(\mu)$, $\non\Cal N(\mu)$ are $\infty$ and $1$ respectively.
As for (b), the proof of 524Pb still works.
}%end of proof of 524T

\leader{524U}{}\dvAnew{2015}\cmmnt{ There is an natural calculation which
I shall want to call on later.

\medskip

\noindent}{\bf Lemma} Let $(\frak A,\bar\mu)$ be a probability algebra.
Then there
is a Radon probability measure on $\{0,1\}^{\tau(\frak A)}$ with measure
algebra isomorphic to $(\frak A,\bar\mu)$.
%548K

\proof{ Write $\kappa$ for $\tau(\frak A)$.

\medskip

{\bf (a)} If $\frak A$ is finite, it is isomorphic to $\Cal PI$
where $I$ is the set of atoms of $\frak A$.   Now
$\#(I)\le 2^{\kappa}$ so we have an injection
$f:I\to\{0,1\}^{\kappa}$.   Let $\mu$ be the point-supported
probability measure on
$\{0,1\}^{\kappa}$ such that $\mu\{f(a)\}=\bar\mu a$ for every
$a\in I$;  this works.

\medskip

{\bf (b)} Otherwise, $\kappa$ is infinite.
By Maharam's theorem, we have a partition $\familyiI{a_i}$ of unity in
$\frak A$ such that, for each $i\in I$, either $a_i$ is an atom or
the principal ideal $\frak A_{a_i}$ is homogeneous with Maharam type
$\kappa_i\ge\omega$, and in the latter case
$(\frak A_{a_i},\bar\mu\restr\frak A_{a_i})$ is
isomorphic to the measure algebra of $\epsilon_i\nu_{\kappa_i}$, where I
write $\epsilon_i=\bar\mu a_i$ for each $i\in I$.
If $a_i$ is an atom, let
$\mu_i$ be a point-supported measure concentrated at a single point of
$\{0,1\}^{\kappa}$ and with mass $\epsilon_i$.
Otherwise, $\kappa_i=\tau(\frak A_{a_i})\le\kappa$;  let
$f_i:\{0,1\}^{\kappa_i}\to\{0,1\}^{\kappa}$ be a continuous injection
and set $\mu_i=\epsilon_i\nu_{\kappa_i}f_i^{-1}$ where
$\nu_{\kappa_i}f_i^{-1}$ is the image measure.
Then $\mu_i$ is a Radon probability measure on $\{0,1\}^{\kappa}$ with
measure algebra isomorphic to
$(\frak A_{a_i},\bar\mu\restr\frak A_{a_i})$.

Now take an injection $g:I\to\{0,1\}^{\Bbb N}$.
Define a measure $\mu$ on $\{0,1\}^{\Bbb N}\times\{0,1\}^{\kappa}$ by
setting

\Centerline{$\mu W=\sum_{i\in I}\mu_iW[\{g(i)\}]$}

\noindent
for those sets $W\subseteq\{0,1\}^{\Bbb N}\times\{0,1\}^{\kappa}$ for which
the sum is defined.   It is easy to check that $\mu$ is a complete
topological probability measure which is inner regular with respect to the
compact sets, that is, is a Radon probability measure.   Also, setting
$E_i=\{g(i)\}\times\{0,1\}^{\kappa}$, the subspace measure $\mu_{E_i}$ is
isomorphic to $\mu_i$ so has measure algebra isomorphic to
$(\frak A_{a_i},\bar\mu\restr\frak A_{a_i})$;
as $\familyiI{E_i}$ is a partition of a
$\mu$-conegligible set, the measure algebra of $\mu$ is isomorphic to the
simple product of the measure algebras of $\mu_{E_i}$, that is, to the
simple product
$\prod_{i\in I}(\frak A_{a_i},\bar\mu\restr\frak A_{a_i})
\cong(\frak A,\bar\mu)$.

As $\kappa$ is infinite,
$\{0,1\}^{\Bbb N}\times\{0,1\}^{\kappa}$
is homeomorphic to $\{0,1\}^{\kappa}$ and we can copy $\mu$ onto a Radon
probability measure on $\{0,1\}^{\kappa}$ with measure algebra isomorphic
to $(\frak A,\bar\mu)$.
}%end of proof of 524U

\exercises{\leader{524X}{Basic exercises (a)}
%\spheader 524Xa
Suppose that $(\frak A,\bar\mu)$ is a probability
algebra and that $\kappa=\link_n(\frak A)$, where $2\le n<\omega$.
Show that there are families $\ofamily{\xi}{\kappa}{A_{\xi}}$ in
$\frak A\setminus\{0\}$ and $\ofamily{\xi}{\kappa}{\epsilon_{\xi}}$ in
$\ocint{0,1}$ such that $\bar\mu(\inf I)\ge\epsilon_{\xi}$ whenever
$I\in[A_{\xi}]^n$ and
$\bigcup_{\xi<\kappa}A_{\xi}=\frak A\setminus\{0\}$.   \Hint{proof of
524L.}
%524M

\spheader 524Xb Let $(X,\Sigma,\mu)$ be a semi-finite measure
space with measure algebra $\frak A$, and $\Cal A$ a family of
non-negligible (not necessarily measurable) subsets of $X$ such that
every non-negligible member of $\Sigma$ includes a member of $\Cal A$.
Show that $\#(\Cal A)\ge\pi(\frak A)$.
%524P

\spheader 524Xc Show that if $\kappa$ is uncountable, there is no
function $f:[0,1]^{\kappa}\to\{0,1\}^{\kappa}$ which is almost
continuous and \imp\ for the usual measures on these spaces.   \Hint{if
$K\subseteq[0,1]^{\kappa}$ is a zero set, any continuous function from
$K$ to $\{0,1\}^{\kappa}$ is determined by coordinates in a countable
set.}
%524P

\spheader 524Xd Let $I^{\|}$ be the split interval and $\mu$ its usual
measure (343J).   Show that there are $f:\{0,1\}^{\omega}\to I^{\|}$ and
$g:I^{\|}\to\{0,1\}^{\omega}$ such that $\mu=\nu_{\omega}f^{-1}$ and
$\nu_{\omega}=\mu g^{-1}$.   \Hint{let $A\subseteq[0,1]$ be a
non-measurable set;  define
$f_0:[0,1]^2\to I^{\|}$ by setting $f_0(x,y)=y^+$ if $x\in A$, $y^-$
otherwise.}
%524P

\spheader 524Xe Let $(Z,\mu)$ be the Stone space of
$(\frak B_{\omega},\bar\nu_{\omega})$.   Show that there is no
$f:\{0,1\}^{\omega}\to Z$ such that $\mu=\nu_{\omega}f^{-1}$.
\Hint{use 515J and 322Ra to show that every non-negligible measurable
subset of $Z$ has cardinal $2^{\frak c}$.}
%524P

\spheader 524Xf Let $X$ be a Hausdorff space with a compact
topological probability measure $\mu$ with Maharam type $\kappa$, and
suppose that $w(X)<\cov\Cal N_{\kappa}$.   (i) Show
that there is an equidistributed sequence for $\mu$.  \Hint{491Eb.}
(ii) Show that if $\mu$ is strictly positive then $X$ is separable.
%524P

\spheader 524Xg Let $(X,\frak T,\Sigma,\mu)$ be a Radon probability
space with a strong lifting, and $(Z,\nu)$ the Stone space of its
measure algebra.   Show that $\shr\Cal N(\mu)\le\shr\Cal N(\nu)$
and $\shr^+\Cal N(\mu)\le\shr^+\Cal N(\nu)$.
\Hint{453Mb.}
%524P

\spheader 524Xh Let $(X,\frak T,\Sigma,\mu)$ be a Radon measure space,
and $K$ the set of infinite cardinals $\kappa$ such that the
Maharam-type-$\kappa$ component of its measure algebra $\frak A$ is
non-zero.   Show that

\Centerline{$\min\{\#(A):A\subseteq X$ has full outer measure$\}
=\sup(\{c(\frak A)\}\cup\{\non\Cal N_{\kappa}:\kappa\in K\})$.}

\spheader 524Xi Show that for any $\sigma$-ideal $\Cal I$ of sets there
is a compact probability measure $\mu$ such that $\Cal I=\Cal N(\mu)$.
\Hint{set $X=\bigcup\Cal I\cup\{x_0\}$.}
%524R

\spheader 524Xj Show that for any non-zero measurable algebra $\frak B$
and any cardinal $\kappa$, there is a complete compact probability
measure $\mu$ such that the measure algebra of $\mu$ is isomorphic to
$\frak B$, $\add\Cal N(\mu)=\omega_1$ and $\cf\Cal N(\mu)\ge\kappa$.
\Hint{524Xi.}
%524Xi 524R

\spheader 524Xk Suppose that
$\non\Cal N_{\frakc}=\cov\Cal N_{\frakc}=\cf\frak c=\frak c$.   Show that
there is a quasi-Radon probability measure $\mu$ with Maharam type
$\frak c$ such that $\add\Cal N(\mu)=\frak c$.
%524T mt52bits   \frak m=\frak c  does it; 524Md, 517Rb + 5A1Fd, 528L

\leader{524Y}{Further exercises (a)}
%\spheader 524Ya
Show that if $m\ge 2$ and
$\familyiI{\frak A_i}$ is a family of $\sigma$-$m$-linked
Boolean algebras, with $\#(I)\le\frak c$, then the free product of
$\familyiI{\frak A_i}$ is $\sigma$-$m$-linked.
%524L  %out of order query

\spheader 524Yb
Let $\frak A$ be a measurable algebra with Maharam type $\lambda$.
Show that there is a family
$\Cal V\subseteq[\lambda]^{\le\frakc}$, cofinal with
$[\lambda]^{\le\frakc}$, such that
$\#(\{A\cap V:V\in\Cal V\})<\FN^*(\frak A)$
for every countable set $A\subseteq\lambda$.
%524O 518J mt52bits

\spheader 524Yc For a Boolean algebra $\frak A$ and a
cardinal $\theta$, write $\psi_{\theta}(\frak A)$ for the smallest
size of any subalgebra $\frak C$ of $\frak A$ such that
$d(\frak C)\ge\theta$.   (If $\theta>d(\frak A)$ set
$\psi_{\theta}(\frak A)=\infty$.)
(i) Show that if $Z$ is the Stone space of $\frak A$, $\Cal I$ is the ideal
of nowhere dense sets in $Z$, and $\theta\ge 2$ then
$\psi_{\theta}(\frak A)\le\cov([Z]^{<\theta},\subseteq,\Cal I)$.
(ii) Show that if $(X,\Sigma,\mu)$ is a \Mth\ compact probability
space with Maharam type $\kappa$, and $\theta$ is uncountable, then

\Centerline{$\psi_{\theta}(\frak B_{\kappa})
=\cov([X]^{<\theta},\subseteq,\Cal N(\mu))
=\add(\Sigma\setminus\Cal N(\mu),\text{\tt{meet}},[X]^{<\theta})$,}

\noindent where
$\text{\tt{meet}}$ is the relation $\{(A,B):A\cap B\ne\emptyset\}$.
\Hint{start with $\mu=\nu_{\kappa}$.}
(iii) Show that if $(X,\Sigma,\mu)$ is a semi-finite locally compact
measure space with measure algebra $\frak A$ then
$\psi_{\omega_1}(\frak A)\le\cf([\cov\Cal N(\mu)]^{\le\omega})$.
(iv) Show that if $(X,\Sigma,\mu)$ is any probability space, with measure
algebra $\frak A$, and $\lambda$ is the product probability measure
on $X^{\Bbb N}$, then
$\cov\Cal N(\lambda)\le\psi_{\omega_1}(\frak A)$.
(v) Show that $\psi_{\add\Cal M}(\frak B_{\omega})\le\non\Cal M$, where
$\Cal M$ is the ideal of meager subsets of $\Bbb R$.
%mt52bits
}%end of exercises

\leader{524Z}{Problems (a)}
%\spheader 524Za
Let $(Z,\mu)$ be the Stone space of
$(\frak B_{\omega},\bar\nu_{\omega})$.   Is $\shr\Cal N(\mu)$
necessarily equal to $\shr\Cal N$?

\spheader 524Zb Can there be a quasi-Radon probability measure $\mu$ with
Maharam type greater than $\frak c$ such that $\add\Cal N(\mu)>\omega_1$?

\endnotes{
\Notesheader{524} The ideas of this section are derived primarily from
{\smc Bartoszy\'nski 84}, {\smc Fremlin 84b} and {\smc Fremlin 91}.
Of course it is not necessary to pass through both $\ell^1(\kappa)$ and
the $\kappa$-localization relation
$(\kappa^{\Bbb N},\subseteq^*,\Cal S_{\kappa})$.   I bring
$\ell^1(\kappa)$ into the argument (following {\smc Bartoszy\'nski 84})
because it will be useful when we come to look at other structures in
later in the chapter, and $\Cal S_{\kappa}$ because it echoes the ideas
of \S522.   But note that 524G seems to need a new idea (the
family $\ofamily{\xi}{\kappa}{E_{\xi}}$ from 524F) not required in 522M.

The difficulties of the work above arise from the fact that while there
are many \imp\ functions between Radon measure spaces, immediately
linking covering numbers and uniformities, there are far fewer
continuous \imp\ functions;  for instance, there is no almost continuous
\imp\ function from the unit interval to the split interval, let alone
to the Stone space of its measure algebra.   And the straightforward
Tukey functions between the ideals $\Cal N_{\kappa}$ of \S523 depend on
measures being images of each other, which is something we can rely
on only when our functions are almost continuous.   (But see 524Xd.)   I
do not know of any
direct construction of a Tukey function from the null
ideal of the Stone space of the Lebesgue measure algebra to $\Cal N$,
for instance.   This is why there is nearly nothing about
shrinking numbers in this section (see 524Za).

There is a significant gap in the calculations in 524P;  for the
cofinality of the null ideal I need to assume that the measure is
$\sigma$-finite.   I have no useful general recipe for $\cf\Cal N(\mu)$,
valid in ZFC, when $\mu$ is a
non-$\sigma$-finite Radon measure.   The point is
that although we can identify $\Cal N(\mu)$ with the product of a family
$\Cal N(\mu_{X_i})$ of partially ordered sets to which the arguments of
this section apply (524Q), this is not in itself
enough to determine its cofinality in the absence of special axioms.
% ccc forcing raises $2^{\omega_1}$ but not $\cf(\omega_1^{\omega_1})$
}%end of notes

\discrpage
\frfilename{mt525.tex}
\versiondate{11.9.13}
\copyrightdate{2004}

\def\chaptername{Cardinal functions of measure theory}
\def\sectionname{Precalibers}

\newsection{525}

I continue the discussion of precalibers in \S516 with results applying
to measure algebras.   I start with connexions between measure spaces
and precalibers of their measure algebras (525B-525C). %525B 525C
The next step is to look at measure-precalibers.   Elementary facts are
in 525D-525G. %525D 525E 525F 525G
When we come to ask which cardinals are precalibers of which measure
algebras, there seem to be real difficulties;  partial answers, largely
based on infinitary combinatorics, are in
525I-525O.  %525I 525J 525K 525M 525N 525O
525P is a note on a particular pair of cardinals.   Finally, 525T deals
with precaliber triples $(\kappa,\kappa,k)$ where
$k$ is finite;  I approach it through a general result on correlations
in uniformly bounded families of random variables (525S).

\leader{525A}{Notation} If $(X,\Sigma,\mu)$ is a measure space,
$\Cal N(\mu)$ will be the null ideal of $\mu$.   For any set
$I$, $\nu_I$ will be the usual measure on $\{0,1\}^I$, $\Tau_I$ its
domain, $\Cal N_I=\Cal N(\nu_I)$ its null ideal and
$(\frak B_I,\bar\nu_I)$ its measure algebra.
In this context, set $e_i=\{x:x\in\{0,1\}^I,\,x(i)=1\}^{\ssbullet}$
in $\frak B_I$ for $i\in I$.
Then $\familyiI{e_i}$ is a stochastically independent
family of elements of measure $\bover12$ in $\frak B_I$, and
$\{e_i:i\in I\}\,\,\tau$-generates $\frak B_I$;  I will say that
$\familyiI{e_i}$ is the {\bf standard generating family} in $\frak B_I$.

\leader{525B}{Proposition} Let $(X,\frak T,\Sigma,\mu)$ be a quasi-Radon
measure space, and $\frak A$ its measure algebra.   Then the downwards
precaliber triples of the partially
ordered set $(\Sigma\setminus\Cal N(\mu),\subseteq)$ are just the
precaliber triples of the Boolean algebra $\frak A$.

\proof{ Put 521Dd and 516C together.}

\leader{525C}{Theorem} Let $(X,\frak T,\Sigma,\mu)$ be a Radon measure
space and $(\frak A,\bar\mu)$ its measure algebra.

(a) A pair $(\kappa,\lambda)$ of cardinals is a precaliber pair of
$\frak A$ iff whenever $\ofamily{\xi}{\kappa}{E_{\xi}}$ is a family in
$\Sigma\setminus\Cal N(\mu)$ there is an $x\in X$ such that
$\#(\{\xi:x\in E_{\xi}\})\ge\lambda$.

(b) A pair $(\kappa,\lambda)$ of cardinals is a measure-precaliber pair
of $(\frak A,\bar\mu)$ iff whenever $\ofamily{\xi}{\kappa}{E_{\xi}}$ is
a family in $\Sigma\setminus\Cal N(\mu)$ such that
$\inf_{\xi<\kappa}\mu E_{\xi}>0$ then there is an $x\in X$ such that
$\#(\{\xi:x\in E_{\xi}\})\ge\lambda$.

(c) Suppose that $\kappa\ge\sat(\frak A)$ is an infinite regular cardinal.
Then the following are equiveridical:

\quad (i) $\kappa$ is a precaliber of $\frak A$;

\quad (ii) $\mu_*(\bigcup_{\xi<\kappa}E_{\xi})=0$ whenever
$\ofamily{\xi}{\kappa}{E_{\xi}}$ is a non-decreasing family in
$\Cal N(\mu)$;

\quad (iii) whenever $\ofamily{\xi}{\kappa}{A_{\xi}}$ is a
non-decreasing family of sets such that $\bigcup_{\xi<\kappa}A_{\xi}=X$,
then there is some $\xi<\kappa$ such that $A_{\xi}$ has full outer
measure in $X$.

\proof{{\bf (a)(i)} Suppose that
$(\kappa,\lambda)$ is a precaliber pair of
$\frak A$ and $\ofamily{\xi}{\kappa}{E_{\xi}}$ is a family in
$\Sigma\setminus\Cal N(\mu)$.   For each $\xi<\kappa$, let
$K_{\xi}\subseteq E_{\xi}$ be a non-negligible compact set.   Then there
is a $\Gamma\in[\kappa]^{\lambda}$ such that
$\{K_{\xi}^{\ssbullet}:\xi\in\Gamma\}$ is centered in $\frak A$.   But
in this case $\{X\}\cup\{K_{\xi}:\xi\in\Gamma\}$ has the finite
intersection property, and must have non-empty intersection.   If $x$ is
any point of this intersection, $\{\xi:x\in E_{\xi}\}$ includes $\Gamma$
and has size at least $\lambda$.

\medskip

\quad{\bf (ii)} Suppose that whenever $\ofamily{\xi}{\kappa}{E_{\xi}}$
is a family in
$\Sigma\setminus\Cal N(\mu)$ there is an $x\in X$ such that
$\#(\{\xi:x\in E_{\xi}\})\ge\lambda$.
Because $\mu$ is complete and strictly localizable
(416B), it has a lifting $\psi:\frak A\to\Sigma$ (341K).   Let
$\ofamily{\xi}{\kappa}{a_{\xi}}$ be a family in $\frak A\setminus\{0\}$;
then there is an $x\in X$ such that $\Gamma=\{\xi:x\in\psi a_{\xi}\}$
has cardinal at least $\lambda$.   But now
$\{\psi a_{\xi}:\xi\in\Gamma\}$ is centered in $\Sigma$ so
$\{a_{\xi}:\xi\in\Gamma\}$ is centered in $\frak A$.   As
$\ofamily{\xi}{\kappa}{a_{\xi}}$ is arbitrary, $(\kappa,\lambda)$ is a
precaliber pair of $\frak A$.

\medskip

{\bf (b)} We can use exactly the same argument, provided that in part
(i) we make sure that $\mu K_{\xi}\ge\bover12\mu E_{\xi}$, so that
$\inf_{\xi<\kappa}\bar\mu K_{\xi}^{\ssbullet}>0$.

\medskip

{\bf (c)}{\bf (i)$\Rightarrow$(iii)} Suppose that $\kappa$ is a
precaliber of $\frak A$ and $\ofamily{\xi}{\kappa}{A_{\xi}}$ is a
non-decreasing family of sets with union $X$.   \Quer\ If no $A_{\xi}$
has full outer measure, then we can choose, for each $\xi<\kappa$, a
non-negligible compact set $K_{\xi}\subseteq X\setminus A_{\xi}$.
Because $\kappa$ is a precaliber of $\frak A$, there is a set
$\Gamma\in[\kappa]^{\kappa}$ such that
$\{K_{\xi}^{\ssbullet}:\xi\in\Gamma\}$ is centered.   Now
$\{K_{\xi}:\xi\in\Gamma\}$ has the finite intersection property and
there is some $x\in\bigcap_{\xi\in\Gamma}K_{\xi}$, in which case
$x\notin\bigcup_{\xi\in\Gamma}A_{\xi}$.   But since $\Gamma$ must be
cofinal with $\kappa$, $\bigcup_{\xi\in\Gamma}A_{\xi}=X$.\ \BanG\  As
$\ofamily{\xi}{\kappa}{A_{\xi}}$ is arbitrary, (iii) is true.

\medskip

\quad{\bf (iii)$\Rightarrow$(ii)} Suppose that (iii) is true, and that
$\ofamily{\xi}{\kappa}{E_{\xi}}$ is a non-decreasing family in
$\Cal N(\mu)$.   \Quer\ If $\bigcup_{\xi<\kappa}E_{\xi}$ has non-zero
inner measure, let $E\subseteq\bigcup_{\xi<\kappa}E_{\xi}$ be a
non-negligible measurable set.   Set $A_{\xi}=E_{\xi}\cup(X\setminus E)$
for each $\xi$;  then $\ofamily{\xi}{\kappa}{A_{\xi}}$ is a
non-decreasing family with union $X$, so there is some $\xi$ such that
$A_{\xi}$ has full outer measure.   But $E\setminus E_{\xi}$ is a
non-negligible measurable set disjoint from $A_{\xi}$.\ \BanG\
As $\ofamily{\xi}{\kappa}{E_{\xi}}$ is arbitrary, (ii) is true.

\medskip

\quad{\bf (ii)$\Rightarrow$(i)} Let $Z$ be the Stone space of
$\frak A$ and $\nu$ its usual
measure (411P).   Because $\mu$ has a lifting, there is an \imp\ function
$f:X\to Z$ (341P).

Let $\ofamily{\xi}{\kappa}{F_{\xi}}$ be a non-decreasing family of
nowhere dense subsets of $Z$.    Then they are all $\nu$-negligible
(411Pa),
so $\ofamily{\xi}{\kappa}{f^{-1}[F_{\xi}]}$ is a non-decreasing family
in $\Cal N(\mu)$ and $\mu_*(\bigcup_{\xi<\kappa}f^{-1}[F_{\xi}])=0$.
But this means that if $G=\interior(\bigcup_{\xi<\kappa}F_{\xi})$,
$\nu G=\mu f^{-1}[G]=0$ and $G$ is empty.   By 516Rb, $\kappa$ is a
precaliber of $\frak A$.
}%end of proof of 525C

\leader{525D}{Proposition} Let $(\frak A,\bar\mu)$ be a measure algebra.

(a) Any precaliber triple of $\frak A$ is a measure-precaliber triple of
$(\frak A,\bar\mu)$.

(b) If $\triplepc{\kappa}{\lambda}{\theta}$ is a measure-precaliber
triple of $(\frak A,\bar\mu)$ and $\kappa$ has uncountable cofinality,
then $\triplepc{\kappa}{\lambda}{\theta}$ is a precaliber triple of
$\frak A$.

(c) If $\kappa$ is a measure-precaliber of $(\frak A,\bar\mu)$, so is
$\cf\kappa$.

\proof{{\bf (a)} is immediate from the definitions in 511E.

\medskip

{\bf (b)} If $\ofamily{\xi}{\kappa}{a_{\xi}}$ is any family in
$\frak A^+$, then there is a $\delta>0$ such that
$\Gamma=\{\xi:\bar\mu a_{\xi}\ge\delta\}$ has cardinal $\kappa$;  and
now there is a $\Gamma'\in[\Gamma]^{\lambda}$ such that
$\{a_{\xi}:\xi\in I\}$ has a non-zero lower bound for every
$I\in[\Gamma']^{<\theta}$.

\medskip

{\bf (c)} The point is that $\kappa$ is a measure-precaliber of
$(\frak A,\bar\mu)$ iff it is a precaliber of the supported relation
$(A_{\delta},\Bsupseteqshort,\frak A^+)$ for every $\delta>0$, where
$A_{\delta}=\{a:a\in\frak A,\,\bar\mu a\ge\delta\}$;  so this is just a
special case of 516Bd.
}%end of proof of 525D

\leader{525E}{Proposition} (a) Let $(\frak A,\bar\mu)$ be a probability
algebra and $\kappa$ an infinite cardinal.   Then $\kappa$ is a
precaliber of $\frak A$ iff either $\frak A$ is finite or $\kappa$ is a
measure-precaliber of $(\frak A,\bar\mu)$ and $\cf\kappa>\omega$.

(b) An infinite cardinal $\kappa$ is a precaliber of every measurable
algebra iff it is a measure-precaliber of every probability algebra and
has uncountable cofinality.

\proof{{\bf (a)} If $\kappa$ is a precaliber of $\frak A$, of course
$\kappa$ is a measure-precaliber of $(\frak A,\bar\mu)$.   Also
$\cf\kappa$ is a precaliber of $\frak A$ (516Bd again), so
$\cf\kappa\ge\sat(\frak A)$ (516Ja);  and if $\frak A$ is infinite,
$\cf\kappa>\omega$.

If $\frak A$ is finite, then any infinite cardinal is a precaliber of
$\frak A$ (516Lc).   If $\kappa$ is a measure-precaliber of
$(\frak A,\bar\mu)$ and $\cf\kappa>\omega$, then $\kappa$ is a
precaliber of $\frak A$ by 525Db.

\medskip

{\bf (b)} Recall that an algebra $\frak A$ is `measurable' iff either
$\frak A=\{0\}$ or there is a functional $\bar\mu$ such that
$(\frak A,\bar\mu)$ is a probability algebra (391B).   So the result
follows directly from (a).
}%end of proof

\leader{525F}{Proposition} Let $(\frak A,\bar\mu)$ be a probability
algebra.

(a) $\omega$ is a measure-precaliber of $(\frak A,\bar\mu)$.

(b) If $\omega\le\kappa<\frak m(\frak A)$, then $\kappa$ is a
measure-precaliber of $(\frak A,\bar\mu)$.

\proof{{\bf (a)} Let $\sequencen{a_n}$ be a sequence in $\frak A$ such
that $\inf_{n\in\Bbb N}\bar\mu a_n=\delta>0$.   Set
$a=\inf_{n\in\Bbb N}\sup_{m\ge n}a_m$;  then
$\bar\mu a=\inf_{n\in\Bbb N}\bar\mu(\sup_{m\ge n}a_m)
\ge\delta>0$, so
$a\ne 0$.   If $0\ne b\Bsubseteq a$ and $n\in\Bbb N$, there is an
$m\ge n$ such that $b\Bcap a_m\ne 0$.   We can therefore choose
inductively a strictly increasing sequence $\sequence{i}{n_i}$ such that
$a\Bcap\inf_{j\le i}a_{n_j}\ne 0$ for every $i$, so that
$\sequence{i}{a_{n_i}}$ is centered.   As $\sequencen{a_n}$ is
arbitrary, $\omega$ is a measure-precaliber of $(\frak A,\bar\mu)$.

\medskip

{\bf (b)} If $\kappa=\omega$, this is (a).   Otherwise, let
$\ofamily{\xi}{\kappa}{a_{\xi}}$ be a family in $\frak A$ with
$\inf_{\xi<\kappa}\bar\mu a_{\xi}=\delta>0$.   Set

\Centerline{$c=\inf_{J\subseteq\kappa,\#(J)<\kappa}
\sup_{\xi\in\kappa\setminus J}a_{\xi}$;}

\noindent then

\Centerline{$\bar\mu c=\inf_{J\subseteq\kappa,\#(J)<\kappa}
\bar\mu(\sup_{\xi\in\kappa\setminus J}a_{\xi})\ge\delta$.}

\noindent Choose $\ofamily{\xi}{\kappa}{I_{\xi}}$ inductively so that,
for each $\xi<\kappa$, $I_{\xi}$ is a countable subset of
$\kappa\setminus\bigcup_{\eta<\xi}I_{\eta}$ and
$c\Bsubseteq\sup_{\eta\in I_{\xi}}a_{\eta}$.

For $\xi<\kappa$, set

\Centerline{$Q_{\xi}=\{b:0\ne b\Bsubseteq c$, $\Exists\eta\in I_{\xi}$,
$b\Bsubseteq a_{\eta}\}$.}

\noindent Then $Q_{\xi}$ is coinitial with $\frak A_c^+$.   Because
$\kappa<\frak m(\frak A)\le\frak m(\frak A_c)$, there is a centered
$R\subseteq\frak A_c^+$ meeting every $Q_{\xi}$.   Now

\Centerline{$\Gamma
=\{\eta:\eta<\kappa$, $\Exists b\in R$, $b\Bsubseteq a_{\eta}\}$}

\noindent meets every $I_{\xi}$ so has cardinal $\kappa$, and
$\{a_{\eta}:\eta\in\Gamma\}$ is centered.   As
$\ofamily{\xi}{\kappa}{a_{\xi}}$ is arbitrary, $\kappa$ is a
measure-precaliber of $(\frak A,\bar\mu)$.
}%end of proof of 525F

\leader{525G}{}\cmmnt{ As is surely to be expected, questions about
precalibers of measurable algebras can
generally be reduced to questions about
precalibers of the algebras $\frak B_{\kappa}$.   Some of these can be
quickly answered in terms of the cardinals examined earlier in this
chapter.

\medskip

\noindent}{\bf Proposition} (a) Let $(\frak A,\bar\mu)$ be a totally
finite measure algebra.   Let $K$ be the set of infinite cardinals
$\kappa'$ such that the Maharam-type-$\kappa'$ component of $\frak A$ is
non-zero\cmmnt{ (cf.\ 524M)}.   If $\kappa$, $\lambda$ and $\theta$
are cardinals, of which $\kappa$ is infinite, then
$\triplepc{\kappa}{\lambda}{\theta}$ is a measure-precaliber triple of
$(\frak A,\bar\mu)$ iff it is a measure-precaliber triple of
$(\frak B_{\kappa'},\bar\nu_{\kappa'})$ for every $\kappa'\in K$.

(b) Suppose that $\omega\le\kappa<\cov\Cal N_{\kappa'}$.   Then $\kappa$
is a measure-precaliber of $\frak B_{\kappa'}$.

(c) For any cardinal $\kappa'$, $\omega_1$ is a precaliber of
$\frak B_{\kappa'}$ iff $\cov\Cal N_{\kappa'}>\omega_1$.

(d) If $\kappa$, $\kappa'$ are cardinals such that
$\non\Cal N_{\kappa'}<\cf\kappa$, then $\kappa$ is a precaliber of
$\frak B_{\kappa'}$.

\proof{{\bf (a)(i)} Suppose that $\triplepc{\kappa}{\lambda}{\theta}$ is
a measure-precaliber triple of $(\frak A,\bar\mu)$, $\kappa'\in K$ and
$\ofamily{\xi}{\kappa}{b_{\xi}}$ is a family in $\frak B_{\kappa'}$ with
$\inf_{\xi<\kappa}\bar\nu_{\kappa'}b_{\xi}=\delta>0$.   Let $a\in\frak A$
be such that the principal ideal $\frak A_a$ is homogeneous with Maharam
type $\kappa'$, so that there is an isomorphism
$\pi:\frak B_{\kappa'}\to\frak A_a$ with
$\Bover1{\bar\mu a}\bar\mu(\pi b)=\bar\nu_{\kappa'}b$ for every
$b\in\frak B_{\kappa'}$ (331L).   Then
$\inf_{\xi<\kappa}\bar\mu(\pi b_{\xi})=\delta\bar\mu a>0$, so there is a
$\Gamma\in[\kappa]^{\lambda}$ such that
$\inf_{\xi\in I}\pi b_{\xi}$ and $\inf_{\xi\in I}b_{\xi}$ are non-zero
for every $I\in[\Gamma]^{<\theta}$.   As
$\ofamily{\xi}{\kappa}{b_{\xi}}$ is arbitrary,
$\triplepc{\kappa}{\lambda}{\theta}$ is a measure-precaliber triple of
$(\frak B_{\kappa'},\bar\nu_{\kappa'})$.

\medskip

\quad{\bf (ii)} Suppose that $\triplepc{\kappa}{\lambda}{\theta}$ is a
measure-precaliber triple of $(\frak B_{\kappa'},\bar\nu_{\kappa'})$ for
every $\kappa'\in K$ and $\ofamily{\xi}{\kappa}{a_{\xi}}$ is a family in
$\frak A$ with $\inf_{\xi<\kappa}\bar\mu a_{\xi}=\delta>0$.   Let
$D\subseteq\frak A\setminus\{0\}$ be a partition of unity in $\frak A$
such that all the principal ideals $\frak A_d$, for $d\in D$, are
homogeneous.   Let $C\subseteq D$ be a finite set such that
$\sum_{d\in D\setminus C}\bar\mu d\le\bover12\delta$.   Then for every
$\xi<\kappa$ there is a $c\in C$ such that
$\bar\mu(a_{\xi}\Bcap c)\ge\bover12\delta\bar\mu c$, so (because
$\kappa$ is infinite) there are $c\in C$ and
$\Gamma_0\in[\kappa]^{\kappa}$ such that
$\bar\mu(a_{\xi}\Bcap c)\ge\bover12\delta\bar\mu c$ for every
$\xi\in\Gamma_0$.   If $c$ is an atom then
$\inf_{\xi\in I}a_{\xi}\Bsupseteq c$ is
non-zero for every $I\subseteq\Gamma_0$.   Otherwise, the Maharam type
$\kappa'$ of $\frak A_c$ belongs to $K$.
Let $\pi:\frak B_{\kappa'}\to\frak A_c$ be an isomorphism with
$\bar\mu(\pi b)=\bar\mu c\cdot\bar\nu_{\kappa'}b$ for every
$b\in\frak B_{\kappa'}$.   Set
$b_{\xi}=\pi^{-1}(a_{\xi}\Bcap c)$;  then
$\bar\nu_{\kappa'}b_{\xi}\ge\bover12\delta$ for every $\xi\in\Gamma_0$.
There is therefore a $\Gamma\in[\Gamma_0]^{\lambda}$ such that
$\inf_{\xi\in I}b_{\xi}$ and $\inf_{\xi\in I}a_{\xi}$
are non-zero for every $I\in[\Gamma]^{<\theta}$.
As $\ofamily{\xi}{\kappa}{a_{\xi}}$ is arbitrary,
$\triplepc{\kappa}{\lambda}{\theta}$ is a measure-precaliber triple of
$(\frak A,\bar\mu)$.

\medskip

{\bf (b)} We have $\cov\Cal N_{\kappa'}=\frak m(\frak B_{\kappa'})$
(524Md), so we can use 525Fb.

\medskip

{\bf (c)} If $\cov\Cal N_{\kappa'}>\omega_1$ then (b) tells us that
$\omega_1$ is a precaliber of $\frak B_{\kappa'}$.   If
$\cov\Cal N_{\kappa'}=\omega_1$, let $\ofamily{\xi}{\omega_1}{E_{\xi}}$
be a cover of $\{0,1\}^{\kappa'}$ by negligible sets;  then
$\ofamily{\xi}{\omega_1}{\bigcup_{\eta<\xi}E_{\eta}}$ is a
non-decreasing
family in $\Cal N_{\kappa'}$ with union of non-zero inner measure, so
525Cc tells us that $\omega_1$ is not a precaliber of
$\frak B_{\kappa'}$.

\medskip

{\bf (d)} If $\kappa'$ is finite this is elementary.   Otherwise,
$d(\frak B_{\kappa'})=\non\Cal N_{\kappa'}$ (524Me).   By 516Lc,
$\kappa$ is a precaliber of $\frak B_{\kappa'}$.
}%end of proof of 525G

\leader{525H}{The structure of $\frak B_I$}\cmmnt{ Several of the
arguments below will depend on the following ideas.}   Let $I$ be any
set and $\familyiI{e_i}$ the standard generating family in $\frak B_I$.
If $a\in\frak B_I$,
there is a smallest countable set $J\subseteq I$ such that $a$ belongs
to the closed subalgebra $\frak C_J$ of $\frak B_I$ generated by
$\{e_i:i\in J\}$\cmmnt{ (254Rd, 325Mb)}.   \cmmnt{(Of course
$\frak C_J$ is canonically isomorphic to $\frak B_J$;  see 325Ma.)}

Now suppose that $\family{\xi}{\Gamma}{a_{\xi}}$ is a family in
$\frak B_I$, that for each $\xi\in\Gamma$ we are given a set
$I_{\xi}\subseteq I$ such that $a_{\xi}\in\frak C_{I_{\xi}}$, and that
$J\subseteq I$ is such that $I_{\xi}\cap I_{\eta}\subseteq J$ for all
distinct $\xi$, $\eta\in\Gamma$.   Then $\family{\xi}{\Gamma}{a_{\xi}}$
is relatively stochastically independent over
$\frak C_J$.   \prooflet{\Prf\
$\family{\xi}{\Gamma}{\frak C_{I_{\xi}\setminus J}}$ is stochastically
independent, because $\family{\xi}{\Gamma}{I_{\xi}\setminus J}$ is
disjoint;   moreover, $\frak C_J$ is independent of
$\frak C_{I\setminus J}
\supseteq\bigcup_{\xi\in\Gamma}\frak C_{I_{\xi}\setminus J}$, and
$\frak C_{I_{\xi}\cup J}$ is the closed subalgebra generated by
$\frak C_{I_{\xi}\setminus J}\cup\frak C_J$ for each $\xi$.
So 458Lg tells us that
$\family{\xi}{\Gamma}{\frak C_{I_{\xi}\cup J}}$ is relatively stochastically
independent over $\frak C_J$;  {\it a fortiori},
$\family{\xi}{\Gamma}{a_{\xi}}$ is relatively stochastically
independent over $\frak C_J$.\ \QeD}
It follows that if $\Delta\subseteq\Gamma$ is finite and
$\inf_{\xi\in\Delta}\upr(a_{\xi},\frak C_J)\ne 0$, then
$\inf_{\xi\in\Delta}a_{\xi}\ne 0$\cmmnt{ (458Lf\dvAformerly{4{}58Hb})};
\cmmnt{in particular,} if
$\family{\xi}{\Gamma}{\upr(a_{\xi},\frak C_J)}$ is centered, so is
$\family{\xi}{\Gamma}{a_{\xi}}$.

\leader{525I}{Theorem} (a)(i) If $\kappa>0$ and
$\triplepc{\kappa}{\lambda}{\theta}$
is a measure-precaliber triple of $(\frak B_{\kappa},\bar\nu_{\kappa})$,
then it is a measure-precaliber triple of every probability algebra.

\quad(ii) If $\kappa>0$ and $\triplepc{\kappa}{\lambda}{\theta}$
is a precaliber triple of $\frak B_{\kappa}$, then it is a precaliber
triple of every measurable algebra.

(b) Suppose that $\cf\kappa\ge\omega_2$.   If $(\kappa,\lambda)$
is a precaliber pair of
$\frak B_{\kappa'}$ for every $\kappa'<\kappa$, then it is a precaliber
pair of every measurable algebra.

(c) Suppose that $\triplepc{\kappa}{\lambda}{\theta}$ is a
measure-precaliber triple of
$(\frak B_{\omega},\bar\nu_{\omega})$ and that $\kappa'$ is such that
$\cff[\kappa']^{\le\omega}<\cf\kappa$.   Then
$\triplepc{\kappa}{\lambda}{\theta}$
is a measure-precaliber triple of
$(\frak B_{\kappa'},\bar\nu_{\kappa'})$.

\proof{{\bf (a)(i)} Let
$(\frak A,\bar\mu)$ be any probability algebra and
$\ofamily{\xi}{\kappa}{a_{\xi}}$
a family in $\frak A^+$ such that $\inf_{\xi<\kappa}\bar\mu a_{\xi}>0$.
Let $\frak B$ be the closed
subalgebra of $\frak A$ generated by $\{a_{\xi}:\xi<\kappa\}$.   Then
$(\frak B,\bar\mu\restrp\frak B)$ is a probability algebra with
Maharam type at most $\kappa$, so is isomorphic to a closed subalgebra
of $(\frak B_{\kappa},\bar\nu_{\kappa})$ (332N).   Since
$\triplepc{\kappa}{\lambda}{\theta}$ is a measure-precaliber triple of
$(\frak B_{\kappa},\bar\nu_{\kappa})$ it is a measure-precaliber triple
of $(\frak B,\bar\mu\restrp\frak B)$ (cf.\ 516Sb), and there
is a $\Gamma\in[\kappa]^{\lambda}$ such that $\{a_{\xi}:\xi\in I\}$ is
bounded below in $\frak B^+$ and therefore in $\frak A^+$ for every
$I\in[\Gamma]^{<\theta}$.   As $\ofamily{\xi}{\kappa}{a_{\xi}}$ is
arbitrary, $\triplepc{\kappa}{\lambda}{\theta}$ is a measure-precaliber
triple of $(\frak A,\bar\mu)$.

\medskip

\quad{\bf (ii)} The same argument applies, deleting the phrase
`$\inf_{\xi<\kappa}\bar\mu a_{\xi}>0$', since if $\frak A$ is a
measurable algebra other than $\{0\}$ there is a
functional $\bar\mu$ such that $(\frak A,\bar\mu)$ is a probability
algebra.

\medskip

{\bf (b)} By (a-ii), it is enough to prove that $(\kappa,\lambda)$ is a
precaliber pair
of $\frak B_{\kappa}$.   Let $\ofamily{\xi}{\kappa}{a_{\xi}}$ be a
family in $\frak B_{\kappa}^+$.   For each $I\subseteq\kappa$, let
$\frak C_I$ be the closed subalgebra of $\frak B_{\kappa}$ generated by
$\{e_i:i\in I\}$, as in 525H.
Then for each $\xi<\kappa$ we have a countable
set $I_{\xi}\subseteq\kappa$ such that $a_{\xi}\in\frak C_{I_{\xi}}$.
Because $\cf\kappa\ge\omega_2$,
there are a $\Gamma\in[\kappa]^{\kappa}$ and a
$J\in[\kappa]^{<\kappa}$ such that $I_{\xi}\cap I_{\eta}\subseteq J$ for
all distinct $\xi$, $\eta\in\Gamma$ (5A1J(a-i)).   Because $\#(J)<\kappa$,
$(\kappa,\lambda)$ is a precaliber pair of $\frak B_J\cong\frak C_J$, so
there is a $\Gamma'\in[\Gamma]^{\lambda}$ such that
$\family{\xi}{\Gamma'}{\upr(a_{\xi},\frak C_J)}$ is centered.   It
follows that $\family{\xi}{\Gamma'}{a_{\xi}}$ is
centered (525H).   As $\ofamily{\xi}{\kappa}{a_{\xi}}$ is arbitrary, we
have the result.

\medskip

{\bf (c)} Let $\ofamily{\xi}{\kappa}{a_{\xi}}$ be a family in
$\frak B_{\kappa'}$ such that $\bar\nu_{\kappa'}a_{\xi}\ge\delta>0$ for
every $\xi<\kappa$.   Fix a cofinal family $\Cal J$ in
$[\kappa']^{\le\omega}$ with cardinal less than $\cf\kappa$.
For each $\xi<\kappa$ let $J_{\xi}\in\Cal J$ be such that
$a_{\xi}\in\frak C_{J_{\xi}}$, where this time $\frak C_{J_{\xi}}$ is
interpreted as a subalgebra of $\frak B_{\kappa'}$.   Then there must be
some $J\in\Cal J$ such that $A=\{\xi:J_{\xi}=J\}$ has cardinal $\kappa$.
Now $(\frak C_J,\bar\nu_{\kappa'}\restrp\frak C_J)$ is isomorphic to a
subalgebra of $(\frak B_{\omega},\bar\nu_{\omega})$, so has
$\triplepc{\kappa}{\lambda}{\theta}$ as a measure-precaliber triple, and
there is a $\Gamma\in[A]^{\lambda}$ such that $\{a_{\xi}:\xi\in I\}$ has
a non-zero lower bound for every $I\in[\Gamma]^{<\theta}$.   As
$\ofamily{\xi}{\kappa}{a_{\xi}}$ is arbitrary,
$\triplepc{\kappa}{\lambda}{\theta}$ is a measure-precaliber triple of
$(\frak B_{\kappa'},\bar\nu_{\kappa'})$.
}%end of proof of 525I

\leader{525J}{Corollary} Suppose that $\kappa$ is an infinite cardinal
and $\kappa<\cov\Cal N_{\kappa}$.   Then $\kappa$ is a
measure-precaliber of every probability algebra.

\proof{ By 525Gb, $\kappa$ is a measure-precaliber of
$(\frak B_{\kappa},\bar\nu_{\kappa})$;  by 525Ia, it is a
measure-precaliber of every probability algebra.
}%end of proof of 525J

\leader{525K}{Proposition} Let $\kappa>\non\Cal N_{\omega}$ be a regular
cardinal such that $\cff[\lambda]^{\le\omega}<\kappa$ for every
$\lambda<\kappa$\cmmnt{ (e.g., $\kappa=\frak c^+$, $(\frak c^+)^+$,
etc.;  or $\kappa=\omega_2$ if $\non\Cal N_{\omega}=\omega_1$)}.
Then $\kappa$ is a precaliber of every measurable algebra.

\proof{ The point is that $\kappa$ is a precaliber of
$\frak B_{\lambda}$ for every $\lambda<\kappa$.   \Prf\ If $\lambda$ is
finite, this is trivial.   Otherwise,

\Centerline{$d(\frak B_{\lambda})=\non\Cal N_{\lambda}
\le\max(\non\Cal N_{\omega},\cff[\lambda]^{\le\omega})<\kappa=\cf\kappa$}

\noindent by 524Me and 523I(a-i);  it follows that $\kappa$ is a precaliber
of $\frak B_{\lambda}$ (516Lc once more).\ \Qed

By 525Ib, $\kappa$ is a precaliber of all measurable algebras.
}%end of proof of 525K

\leader{525L}{}\cmmnt{ If $\kappa>\frak c$ is not a strong limit
cardinal we can do a little better than 525K.

\medskip

\noindent}{\bf Proposition}\cmmnt{ ({\smc D\v{z}amonja \& Plebanek 04})}
Suppose that $\lambda$ and $\kappa$ are infinite cardinals such
that $\lambda^{\omega}<\cf\kappa\le\kappa\le 2^{\lambda}${, where
$\lambda^{\omega}$ is the cardinal power}.   Then
$\kappa$ is a precaliber of every measurable algebra.

\proof{ By 525Eb and
525I(a-ii), it is enough to show that $\kappa$ is a precaliber of
$\frak B_{\kappa}$.   Let $\ofamily{\xi}{\kappa}{a_{\xi}}$ be a family
in $\frak B_{\kappa}\setminus\{0\}$.   Let
$\theta:\frak B_{\kappa}\to\Tau_{\kappa}$ be a lifting, and for each
$\xi<\kappa$ let $K_{\xi}$ be a non-empty closed subset of
$\theta a_{\xi}$ which is determined by coordinates in a countable set
$I_{\xi}$.   We may suppose that each $I_{\xi}$ is infinite;  let
$h_{\xi}:\Bbb N\to I_{\xi}$ be a bijection, and set
$g_{\xi}(x)=xh_{\xi}$ for $x\in\{0,1\}^{\kappa}$, so that
$g_{\xi}:\{0,1\}^{\kappa}\to\{0,1\}^{\Bbb N}$ is continuous and
$K_{\xi}=g_{\xi}^{-1}[g_{\xi}[K_{\xi}]]$.   As $\frak c<\cf\kappa$,
there is an $L\subseteq\{0,1\}^{\Bbb N}$ such that
$\Gamma_0=\{\xi:\xi<\kappa$, $g_{\xi}[K_{\xi}]=L\}$ has cardinal
$\kappa$.

Because $\kappa\le 2^{\lambda}$,
there is an $f:\kappa\times\lambda^{\omega}\to\Bbb N$ such that
whenever $\sequencen{\xi_n}$ is a sequence of distinct elements of
$\kappa$ there is an $\eta<\lambda^{\omega}$ such that $f(\xi_n,\eta)=n$
for every $n$ (5A1Fg).   For each $\eta<\lambda^{\omega}$, set $A_{\eta}
=\{\xi:\xi<\kappa$, $f(h_{\xi}(n),\eta)=n$ for every $n\in\Bbb N\}$;
then $\bigcup_{\eta<\lambda^{\omega}}A_{\eta}=\kappa$, while
$\cf\kappa>\lambda^{\omega}$, so there is an $\eta^*<\lambda^{\omega}$
such that $\Gamma=\Gamma_0\cap A_{\eta^*}$ has $\kappa$ members.

Fix $z\in L$.   For $\xi$, $\eta\in\Gamma$ and $i$, $j\in\Bbb N$,

\Centerline{$h_{\xi}(i)=h_{\eta}(j)
\Longrightarrow i=f(h_{\xi}(i),\eta^*)=f(h_{\eta}(j),\eta^*)=j$.}

\noindent So we can find an $x\in\{0,1\}^{\kappa}$ such that
$x(h_{\xi}(i))=z(i)$ whenever $\xi\in\Gamma$ and $i\in\Bbb N$;  that is,
$g_{\xi}(x)=z$ for every $\xi\in\Gamma$.   But this means that
$x\in g_{\xi}^{-1}[L]=K_{\xi}$ for every $\xi\in\Gamma$.
It follows that whenever $I\in[\Gamma]^{<\omega}$ then
$\bigcap_{\xi\in I}\theta a_{\xi}\ne\emptyset$ and
$\inf_{\xi\in I}a_{\xi}\ne 0$;  that is, that $\{a_{\xi}:\xi\in\Gamma\}$
is centered.   As $\ofamily{\xi}{\kappa}{a_{\xi}}$ is arbitrary,
$\kappa$ is a precaliber of $\frak B_{\kappa}$.
}%end of proof of 525L

\leader{525M}{Proposition} Let $(\frak A,\bar\mu)$ be a probability
algebra and $\kappa$ an infinite cardinal such that $\cf\kappa$ is a
measure-precaliber of $(\frak A,\bar\mu)$ and
$\lambda^{\omega}<\kappa$ for every $\lambda<\kappa$.   Then $\kappa$ is
a measure-precaliber of $(\frak A,\bar\mu)$.

\proof{ If $\kappa=\cf\kappa$, we can stop;  so henceforth I will suppose
that $\kappa$ is singular.

\medskip

{\bf (a)} Suppose first that $\frak A=\frak B_I$ for some set
$I$;  let $\familyiI{e_i}$ be the standard generating family in
$\frak B_I$.   If $\kappa$ is regular, the result is
trivial.   Otherwise, let $\ofamily{\xi}{\kappa}{a_{\xi}}$ be a family
in $\frak A^+$ such that $\inf_{\xi<\kappa}\bar\nu_Ia_{\xi}=\delta>0$.
There is a strictly increasing family
$\ofamily{\alpha}{\cf\kappa}{\kappa_{\alpha}}$ of regular uncountable
cardinals with supremum $\kappa$ such that $\kappa_0>\cf\kappa$ and if
$\alpha<\cf\kappa$ and $\lambda<\kappa_{\alpha}$ then
$\lambda^{\omega}<\kappa_{\alpha}$.   \Prf\ All we need to know is that
if $\theta<\kappa$ there is a regular uncountable cardinal $\theta'$
such that $\theta\le\theta'<\kappa$ and $\lambda^{\omega}<\theta'$
whenever $\lambda<\theta'$;  and $\theta'=(\theta^{\omega})^+$ has this
property.\ \QeD\

For each $\xi<\kappa$, let $I_{\xi}\subseteq I$ be a countable set such
that $a_{\xi}$ belongs to the closed subalgebra of $\frak A$ generated
by $\{e_i:i\in I_{\xi}\}$.   By the $\Delta$-system Lemma (5A1J(a-ii)),
there is for each $\alpha<\cf\kappa$ a set
$\Gamma_{\alpha}\subseteq\kappa_{\alpha+1}\setminus\kappa_{\alpha}$ such
that $\#(\Gamma_{\alpha})=\kappa_{\alpha+1}$ and
$\family{\xi}{\Gamma_{\alpha}}{I_{\xi}}$ is a $\Delta$-system with
root $J_{\alpha}$ say.   Set $J=\bigcup_{\alpha<\cf\kappa}J_{\alpha}$,
so that $\#(J)\le\cf\kappa$, and

\Centerline{$\Gamma'_{\alpha}
=\{\xi:\xi\in\Gamma_{\alpha},\,
  (I_{\xi}\setminus J_{\alpha})
  \cap(J\cup\bigcup_{\eta<\kappa_{\alpha}}I_{\eta})
  =\emptyset\}$;}

\noindent then $\#(\Gamma'_{\alpha})=\kappa_{\alpha+1}$ for every
$\alpha<\cf\kappa$, and $I_{\xi}\cap I_{\eta}\subseteq J$ for all
distinct $\xi$,
$\eta\in\Gamma'=\bigcup_{\alpha<\cf\kappa}\Gamma'_{\alpha}$.   Let
$\frak C_J$ be the closed subalgebra of $\frak A$ generated by
$\{e_i:i\in J\}$.   For $\xi\in\Gamma'$, set
$b_{\xi}=\upr(a_{\xi},\frak C_J)$.   By 515Ma,

\Centerline{$\#(\frak C_J)\le\max(\omega,\#(J))^{\omega}
<\kappa_{\alpha+1}=\cf\kappa_{\alpha+1}$,}

\noindent there is for each $\alpha<\cf\kappa$ a
$c_{\alpha}\in\frak C_J$ such that
$\Gamma''_{\alpha}=\{\xi:\xi\in\Gamma'_{\alpha},\,b_{\xi}=c_{\alpha}\}$
has cardinal $\kappa_{\alpha+1}$.    Note that

\Centerline{$\bar\nu_Ic_{\alpha}=\bar\nu_Ib_{\xi}\ge\bar\nu_Ia_{\xi}
\ge\delta$}

\noindent whenever $\alpha<\cf\kappa$ and $\xi\in\Gamma''_{\alpha}$.

Now recall that we are supposing that $\cf\kappa$ is a
measure-precaliber of
$\frak A$.   So there is a $\Delta\in[\cf\kappa]^{\cf\kappa}$ such that
$\{c_{\alpha}:\alpha\in\Delta\}$ is centered in $\frak A$.   Now
$\Gamma''=\bigcup_{\alpha\in\Delta}\Gamma''_{\alpha}$ has cardinal
$\kappa$, and $\family{\xi}{\Gamma''}{b_{\xi}}$ is centered.   It
follows that $\family{\xi}{\Gamma''}{a_{\xi}}$ is
centered (525H).

As $\ofamily{\xi}{\kappa}{a_{\xi}}$ is arbitrary, $\kappa$ is a
measure-precaliber of $\frak A$.

\medskip

{\bf (b)} For the general case, observe that by Maharam's theorem (332B)
$\frak A$ is isomorphic to the simple product
$\prod_{k\in K}\frak A_{d_k}$ of a countable family of homogeneous
principal ideals, where $\family{k}{K}{d_k}$ is a partition of unity in
$\frak A$.   Let $\ofamily{\xi}{\kappa}{a_{\xi}}$ be a family in
$\frak A$ such that $\inf_{\xi<\kappa}\bar\mu a_{\xi}=\delta>0$.   Let
$L\subseteq K$ be a finite set such that
$\sum_{k\in K\setminus L}\bar\mu d_k=\delta'<\delta$.   Then there is
some $k\in L$ such that

\Centerline{$\Gamma_k=\{\xi:\xi<\kappa,\,\bar\mu(a_{\xi}\Bcap d_k)
\ge\Bover{\delta-\delta'}{\#(L)}\}$}

\noindent has cardinal $\kappa$.   Since $\cf\kappa$ is a
measure-precaliber of $\frak A$, it is also a measure-precaliber of
$\frak A_{d_k}$ (cf.\ 516Sc).   Since
$(\frak A_{d_k},\bar\mu\restrp\frak A_{d_k})$ is isomorphic, up to a
scalar multiple of the measure, to $(\frak B_I,\bar\nu_I)$ for some $I$,
(a) tells us that $\kappa$ is a measure-precaliber of $\frak A_{d_k}$.
There is therefore a set $\Gamma\in[\Gamma_k]^{\kappa}$ such that
$\family{\xi}{\Gamma}{a_{\xi}\Bcap d_k}$ and
$\family{\xi}{\Gamma}{a_{\xi}}$ are centered.   As
$\ofamily{\xi}{\kappa}{a_{\xi}}$ is arbitrary, $\kappa$ is a
measure-precaliber of $\frak A$.
}%end of proof of 525M

\leader{525N}{Proposition}\cmmnt{ ({\smc Argyros \& Tsarpalias 82})}
Let $\kappa$ be either $\omega$ or a strong limit cardinal of countable
cofinality, and suppose that $2^{\kappa}=\kappa^+$.   Then $\kappa^+$ is
not a precaliber of $\frak B_{\kappa}$.

\proof{ By 523Lb, $\non\Cal N_{\kappa}>\kappa$.   So if we enumerate
$\{0,1\}^{\kappa}$ as $\ofamily{\xi}{\kappa^+}{x_{\xi}}$ and set
$E_{\xi}=\{x_{\eta}:\eta<\xi\}$ for $\xi<\kappa^+$,
$\ofamily{\xi}{\kappa^+}{E_{\xi}}$ is an increasing family in
$\Cal N_{\kappa}$ with union $\{0,1\}^{\kappa}$.
By 525Cc, $\kappa^+$ is not a precaliber of $\frak B_{\kappa}$.
}%end of proof of 525N

\leader{525O}{}\cmmnt{ As in 523P, GCH decides the most important
questions.

\medskip

\noindent}{\bf Proposition} Suppose that the generalized continuum
hypothesis is true.

(a) An infinite cardinal $\kappa$ is a measure-precaliber of every
probability algebra iff $\cf\kappa$ is not the successor of a cardinal
of countable cofinality.

(b) An infinite cardinal $\kappa$ is a precaliber of every measurable
algebra iff $\cf\kappa$ is neither $\omega$ nor the successor of a
cardinal of countable cofinality.

\proof{{\bf (a)(i)} If $\kappa$ is a measure-precaliber of every
probability algebra, so is $\cf\kappa$ (525Dc).   By 525N, $\cf\kappa$
cannot be the successor of a cardinal of countable cofinality.

\medskip

\quad{\bf (ii)} Now suppose that $\cf\kappa$ is not the successor of a
cardinal of countable cofinality.   If $\kappa=\omega$, then
certainly $\kappa$ is a measure-precaliber of every probability algebra
(525Fa).   Otherwise,
$\kappa>\lambda^{\omega}$ for every $\lambda<\kappa$ and
$\cf\kappa>\lambda^{\omega}$ for every $\lambda<\cf\kappa$ (5A6Ac).   By
525K, $\cf\kappa$ is a measure-precaliber of every probability algebra;
by 525M, so is $\kappa$.

\medskip

{\bf (b)} Put (a) and 525Eb together.
}%end of proof of 525O

\leader{*525P}{}\cmmnt{ As in 522U, the Freese-Nation number of
$\Cal P\Bbb N$ is relevant to the questions here.

\medskip

\noindent}{\bf Proposition} $(\frakmctbl,\FN^*(\Cal P\Bbb N))$ is not a
precaliber pair of $\frak B_{\omega}$.

\proof{ By 518D(iv), the Freese-Nation number of the
topology of $\{0,1\}^{\omega}$ is $\FN(\Cal P\Bbb N)$;  the regular
Freese-Nation numbers are therefore also equal.   We know that
$\frakmctbl$ is the covering number of the meager ideal of $\Bbb R$
(522Sa), and therefore also of the meager ideal of $\{0,1\}^{\omega}$
(522Wb) and of the nowhere dense ideal of $\{0,1\}^{\omega}$.
By 518E, there
is a set $A\subseteq\{0,1\}^{\omega}$, with cardinal $\frakmctbl$, such
that $\#(A\cap F)<\FN^*(\Cal P\Bbb N)$ for every nowhere dense set
$F\subseteq\{0,1\}^{\omega}$.

Fix a nowhere dense compact set $K\subseteq\{0,1\}^{\omega}$ of
non-zero measure.   For each $x\in A$, set
$a_x=(K+x)^{\ssbullet}$ in $\frak B_{\omega}$, where $+$ here is the
usual group operation corresponding to the identification
$\{0,1\}^{\omega}\cong\Bbb Z_2^{\omega}$.   Then every $a_x$ is
non-zero.   If $B\subseteq A$ and $\{a_x:x\in B\}$ is centered, then
$\{K+x:x\in B\}$ has the finite intersection property, so there is a $y$
in its intersection;  now $B\subseteq A\cap(K+y)$, and $K+y$ is nowhere
dense, so $\#(B)<\FN^*(\Cal P\Bbb N)$.   Thus
$\family{x}{A}{a_x}$ has no centered subfamily with cardinal
$\FN^*(\Cal P\Bbb N)$ and witnesses
that $(\frakmctbl,\FN^*(\Cal P\Bbb N))$ is not a precaliber pair of
$\frak B_{\omega}$.
}%end of proof of 525P

\leader{525Q}{}\cmmnt{ I turn now to some results which may be
interpreted as information on precaliber triples in which the third
cardinal is {\it finite}.

\medskip

\noindent}{\bf Lemma} Let $(\frak A,\bar\mu)$ be a semi-finite measure
algebra, $\sequencen{u_n}$ a $\|\,\|_2$-bounded sequence in
$L^2=L^2(\frak A,\bar\mu)^+$, and $\Cal F$ a non-principal
ultrafilter on $\Bbb N$.   Suppose that $p\in\coint{0,\infty}$ is such that
$\sup_{n\in\Bbb N}\|u_n^p\|_2$ is finite, and set
$v=\lim_{n\to\Cal F}u_n$, $w=\lim_{n\to\Cal F}u_n^p$, the limits being
taken for the weak topology in $L^2$.
Then $v^p\le w$.

\proof{ Of course the positive cone of $L^2$ is closed for the weak
topology so $v\ge 0$ and we can speak of $v^p$.   \Quer\ If
$v^p\not\le w$, there are
$\alpha$, $\beta\ge 0$ such that $\alpha^p>\beta$ and

\Centerline{$a=\Bvalue{v>\alpha}\Bsetminus\Bvalue{w>\beta}\ne 0$.}

\noindent Let $b\Bsubseteq a$ be such that $0<\bar\mu b<\infty$ and
consider $u=\Bover1{\bar\mu b}\chi b$.   Then, setting $q=p/(p-1)$ (of
course $p>1$),

$$\eqalignno{(\alpha\bar\mu b)^p
&\le(\int_bv)^p
=\lim_{n\to\Cal F}(\int u_n\times\chi b\times\chi b)^p
\le\lim_{n\to\Cal F}\bigl(\|u_n\times\chi b\|_p\|\chi b\|_q)^p\cr
\displaycause{by H\"older's inequality, 244Eb}
&=\lim_{n\to\Cal F}(\bar\mu b)^{p/q}\int_bu_n^p
=(\bar\mu b)^{p-1}\int_bw
\le\beta(\bar\mu b)^p
<\alpha^p(\bar\mu b)^p,\cr}$$

\noindent which is absurd.\ \BanG\  So $v^p\le w$.
}%end of proof of 525Q

\vleader{72pt}{525R}{Lemma}
Let $(\frak A,\bar\mu)$ be a probability algebra and
$\sequencen{u_n}$ a $\|\,\|_{\infty}$-bounded sequence in
$L^{\infty}(\frak A,\bar\mu)^+$
such that $\delta=\inf_{n\in\Bbb N}\int u_n>0$.   Let $k_0,\ldots,k_m$
be strictly positive integers with sum $k$.   Suppose that
$\gamma<\delta^k$.

(a) There are integers $n_0<n_1<\ldots<n_m$ such that
$\int\prod_{j=0}^mu_{n_j}^{k_j}\ge\gamma$.

(b) In fact, there is an infinite set $I\subseteq\Bbb N$ such that
$\int\prod_{j=0}^mu_{n_j}^{k_j}\ge\gamma$ whenever $n_0,\ldots,n_m$ belong
to $I$ and $n_0<n_1<\ldots<n_m$.

\proof{{\bf (a)} Let $\Cal F$ be a non-principal ultrafilter on
$\Bbb N$.   For each $j\le m$, let $v_j$ be the limit
$\lim_{n\to\Cal F}u_n^{k_j}$ for the weak topology on
$L^2(\frak A,\bar\mu)$;  let $v$ be the limit $\lim_{n\to\Cal F}u_n$.
By 525Q,

$$\eqalignno{\int\prod_{j=0}^mv_j
&\ge\int\prod_{j=0}^mv^{k_j}
=\int v^k
=(\|\chi 1\|_q\|v\|_k)^k\cr
\displaycause{where $q=\Bover{k}{k-1}$, or $\infty$ if $k=1$}
&\ge\bigl(\int v\times\chi 1\bigr)^k\cr
\displaycause{by H\"older's inequality again, if $k>1$}
&=\bigl(\int v\bigr)^k
=\lim_{n\to\Cal F}\bigl(\int u_n\bigr)^k
\ge\delta^k
>\gamma.\cr}$$

\noindent (Or use 244Xd to show more directly that
$\int v^k\ge(\int v)^k$.)   We can therefore choose $n_0,\ldots,n_m$
inductively so that

\Centerline{$\int\prod_{j=0}^su_{n_j}^{k_j}
  \times\prod_{j=s+1}^mv_j
>\gamma$}

\noindent for each $s\le m$ (interpreting the final product
$\prod_{j=m+1}^mv_j$ as $\chi 1$), since when we come to choose
$n_s$ we shall be able to use any member of

\Centerline{$\{n:n>n_j\text{ for }j<s,\,
\int u_n^{k_s}\times\prod_{j=0}^{s-1}u_{n_j}^{k_j}
  \times\prod_{j=s+1}^mv_j>\gamma\}$,}

\noindent which belongs to $\Cal F$ so is not empty.   At the end of the
induction we shall have a sequence $n_0<\ldots<n_m$  such that
$\int\prod_{j=0}^mu_{n_j}^{k_j}\ge\gamma$, as required.

\medskip

{\bf (b)} Let $\Cal J\subseteq[\Bbb N]^{m+1}$ be the family of all sets
of the form $\{n_0,\ldots,n_m\}$ where $n_0<\ldots<n_m$ and
$\int\prod_{j=0}^mu_{n_j}^{k_j}\ge\gamma$.   By (a), applied to
subsequences of $\sequencen{u_n}$, every infinite subset of $\Bbb N$
includes some member of $\Cal J$.   By Ramsey's theorem (4A1G), there is
an infinite $I\subseteq\Bbb N$ such that $[I]^{m+1}\subseteq\Cal J$,
which is what we need.
}%end of proof of 525R

\leader{525S}{Theorem}\cmmnt{ ({\smc Fremlin 88})} Let
$(\frak A,\bar\mu)$ be a probability algebra and $\kappa$ an infinite
cardinal.   Let $\ofamily{\xi}{\kappa}{u_{\xi}}$ be a
$\|\,\|_{\infty}$-bounded family in $L^{\infty}(\frak A)^+$.   Set
$\delta=\inf_{\xi<\kappa}\int u_{\xi}$.   Then for any $k\in\Bbb N$ and
$\gamma<\delta^{k+1}$ there is a $\Gamma\in[\kappa]^{\kappa}$ such that
$\int\prod_{i=0}^ku_{\xi_i}\ge\gamma$ for all
$\xi_0,\ldots,\xi_k\in\Gamma$.

\proof{{\bf (a)} It will be helpful to note straight away that it will be
enough to consider the case $(\frak A,\bar\mu)=(\frak B_I,\bar\nu_I)$ for
some set $I$.   \Prf\ There is always a $(\frak B_I,\bar\nu_I)$ in which
$(\frak A,\bar\mu)$ can be embedded.   In this case, $L^{\infty}(\frak A)$
can be identified, as $f$-algebra, with
a subspace of $L^{\infty}(\frak B_I)$,
and the embedding respects integrals.   So we can regard
$\ofamily{\xi}{\kappa}{u_{\xi}}$ as a family in
$L^{\infty}(\frak B_I)$ and perform all calculations there.\ \Qed

At the same time, the case $\delta=0$ is trivial, so let us suppose
henceforth that $\delta>0$.

\medskip

{\bf (b)} Next, having fixed on a suitable set $I$, let $\familyiI{e_i}$ be
the standard generating family in $\frak B_I$, and for $J\subseteq I$ let
$\frak C_J$ be the closed subalgebra of $\frak B_I$ generated by
$\{e_i:i\in J\}$;  following 325N, I will say that a member of $\frak C_J$
is `determined by coordinates in $J$'.   For $J\subseteq I$ let
$P_J:L^1(\frak B_I,\bar\nu_I)\to L^1(\frak C_J,\bar\nu_I\restr\frak C_J)$
be the conditional expectation operator.   Note that if $J$, $K\subseteq I$
then $P_JP_K=P_{J\cap K}$ (254Ra/458M(iii)).

It will be useful to start by looking at a particular
subset $W$ of $L^{\infty}(\frak B_I)$, being the set of linear
combinations $\sum_{i=0}^n\alpha_i\chi c_i$ where every $\alpha_i$ is
rational and every $c_i$ is determined by coordinates in a finite set.
Now $P_J[W]\subseteq W$ for every $J\subseteq I$.   \Prf\ If
$c\in\frak C_K$ where $K\subseteq I$ is finite, then

\Centerline{$P_J(\chi c)=P_JP_K(\chi c)=P_{J\cap K}(\chi c)
=\sum_{d\text{ is an atom of }\frak C_{J\cap K}}
   \Bover{\bar\nu_I(c\Bcap d)}{\bar\nu_Id}\chi d
\in W$.}

\noindent As $P_J$ is linear, this is enough.\ \QeD\  Observe also that if
$K\subseteq I$ is finite, then $P_K[W]$ is countable, being the set of
rational linear combinations of $\{\chi c:c\in\frak C_K\}$.

\medskip

{\bf (c)} Suppose, therefore, that we have a set $I$, a
$\|\,\|_{\infty}$-bounded family $\ofamily{\xi}{\kappa}{u_{\xi}}$
in $L^{\infty}(\frak B_I)^+$ with
$\inf_{\xi<\kappa}\int u_{\xi}\penalty-100=\delta>0$,  a $k\in\Bbb N$ and a
$\gamma<\delta^{k+1}$.   To begin with, let us suppose further that

\inset{($\alpha$) $u_{\xi}\le\chi 1$ for every $\xi<\kappa$,

($\beta$) $u_{\xi}\in W$, as described in (b),
for each $\xi<\kappa$;}

\noindent for each $\xi<\kappa$, let
$I_{\xi}\in[I]^{<\omega}$ be such that
$P_{I_{\xi}}u_{\xi}=u_{\xi}$.

\medskip

\quad{\bf (i)} Suppose that $\kappa=\omega$.   Because there are only
finitely many sequences $k_0,\ldots,k_m$ of strictly positive integers
with sum equal to $k+1$, we can use 525Rb a finite number of times to
find an infinite $\Gamma\subseteq\omega$ such that
$\int\prod_{j=0}^mu_{n_j}^{k_j}\ge\gamma$ whenever $\sum_{j=0}^mk_j=k+1$
and $n_0<\ldots<n_m$ belong to $\Gamma$.   But in this case we surely
have $\int\prod_{i=0}^ku_{n_i}\ge\gamma$ for all
$n_0,\ldots,n_k\in\Gamma$.

\medskip

\quad{\bf (ii)} Next, suppose that $\kappa>\omega$ is regular.   By the
$\Delta$-system Lemma (4A1Db) there is a $\Delta\in[\kappa]^{\kappa}$
such that $\family{\xi}{\Delta}{I_{\xi}}$ is a $\Delta$-system with root
$J$ say.
Since $P_J[W]$ is countable, there is a $v$ such that
$\Gamma=\{\xi:\xi\in\Delta,\,P_Ju_{\xi}=v\}$ has cardinal $\kappa$.
Of course

\Centerline{$\int v=\int u_{\xi}\ge\delta$}

\noindent for every $\xi\in\Gamma$.

By 458Lg again, $\family{\xi}{\Delta}{\frak C_{I_{\xi}}}$
is relatively independent over $\frak C_J$.
Now suppose that $\xi_0,\ldots,\xi_k$ belong to $\Gamma$.   Then

$$\eqalignno{\int\prod_{i=0}^ku_{\xi_i}
&\ge\int\prod_{i=0}^kP_Ju_{\xi_i}\cr
\displaycause{458Lh}
&=\int v^{k+1}
\ge\bigl(\int v\bigr)^{k+1}\cr
\displaycause{as in the proof of 525Ra}
&\ge\delta^{k+1}
\ge\gamma,\cr}$$

\noindent and this is what we need to know.

\medskip

\quad{\bf (iii)} Finally, suppose that $\kappa>\cf\kappa\ge\omega$.
Set $\lambda=\cf\kappa$ and let
$\ofamily{\zeta}{\lambda}{\kappa_{\zeta}}$ be a strictly increasing
family of regular cardinals greater than $\lambda$ and with supremum
$\kappa$.   For each $\zeta<\lambda$ let
$\Delta_{\zeta}\subseteq\kappa_{\zeta+1}\setminus\kappa_{\zeta}$ be a
set with cardinal $\kappa_{\zeta+1}$ such that
$\family{\xi}{\Delta_{\zeta}}{I_{\xi}}$ is a $\Delta$-system with root
$J_{\zeta}$ say.   Set $J=\bigcup_{\zeta<\lambda}J_{\zeta}$;  note that
$\#(J)\le\lambda<\kappa_{\zeta+1}$ for every $\zeta$, so that

\Centerline{$\Delta'_{\zeta}=\{\xi:\xi\in\Delta_{\zeta},\,
(I_{\xi}\setminus
J_{\zeta})\cap(J\cup\bigcup_{\eta<\kappa_{\zeta}}I_{\eta})
  =\emptyset\}$}

\noindent still has cardinal $\kappa_{\zeta+1}$.

If $\zeta<\lambda$ and $\xi\in\Delta'_{\zeta}$,
$I_{\xi}\cap J$ is included in the
finite set $J_{\zeta}$;  so
$\{P_Ju_{\xi}:\xi\in\Delta'_{\zeta}\}$ is countable, and there is a
$v_{\zeta}$ such that
$\Delta''_{\zeta}=\{\xi:\xi\in\Delta'_{\zeta},\,P_Ju_{\xi}=v_{\zeta}\}$ has
cardinal $\kappa_{\zeta+1}$.   Note that

\Centerline{$\int v_{\zeta}=\int u_{\xi}\ge\delta$}

\noindent whenever $\zeta<\lambda$ and $\xi\in\Delta''_{\zeta}$.

Because $\lambda$ is regular, we can apply (i) or (ii) above to find an
$A\in[\lambda]^{\lambda}$ such that
$\int\prod_{i=0}^kv_{\zeta_i}\ge\gamma$ whenever
$\zeta_0,\ldots,\zeta_k\in A$.   Set
$\Gamma=\bigcup_{\zeta\in A}\Delta''_{\zeta}$;  because $A$ must be
cofinal with $\lambda$, $\#(\Gamma)=\kappa$.

If $\xi$, $\eta\in\Gamma$ are distinct, then
$I_{\xi}\cap I_{\eta}\subseteq J$.   So
$\family{\xi}{\Gamma}{\frak C_{I_{\xi}}}$ is
relatively independent over $\frak C_J$.   Take any
$\xi_0,\ldots,\xi_k\in\Gamma$;  for each $i\le k$, let $\zeta_i\in A$ be
such that $\xi_i\in\Delta''_{\zeta_i}$.   By 458Lh again,

\Centerline{$\int\prod_{i=0}^nu_{\xi_i}
\ge\int\prod_{i=0}^nP_Ju_{\xi_i}
=\int\prod_{i=0}^nv_{\zeta_i}\ge\gamma$,}

\noindent so we are done (provided ($\alpha$)-($\beta$) are true).

\medskip

{\bf (d)} Now let us unwind these conditions from the bottom.

\medskip

\quad{\bf (i)} If ($\alpha$) is true, but ($\beta$) is not,
take $\epsilon\in\ooint{0,\delta}$ such that
$(\delta-\epsilon)^{k+1}>\gamma+(k+1)\epsilon$.   For each $\xi<\kappa$,
let $u'_{\xi}\in W$ be such that $u'_{\xi}\le\chi 1$ and
$\int|u_{\xi}-u'_{\xi}|\le\epsilon$.   (Such a $u'_{\xi}$ exists because
$\bigcup\{\frak C_K:K\in[I]^{<\omega}\}$ is topologically dense in
$\frak B_I$ and $u\wedge\chi 1\in W$ for every $u\in W$.)   Then
$\int u'_{\xi}\ge\delta-\epsilon$ for each $\xi$, so we can apply
(c) to $\ofamily{\xi}{\kappa}{u'_{\xi}}$ to see that there is a
$\Gamma\in[\kappa]^{\kappa}$ such that
$\int\prod_{i=0}^ku'_{\xi_i}\ge\gamma+(k+1)\epsilon$ for all
$\xi_0,\ldots,\xi_k\in\Gamma$.   Now (because $u_{\xi}$ and $u'_{\xi}$
all lie between $0$ and $\chi 1$) we have

\Centerline{$\bigl|\prod_{i=0}^ku_{\xi_i}-\prod_{i=0}^ku'_{\xi_i}\bigr|
\le\sum_{i=0}^k|u_{\xi_i}-u'_{\xi_i}|$}

\noindent (see 285O), so that

\Centerline{$\int\prod_{i=0}^ku_{\xi_i}
\ge\int\prod_{i=0}^ku'_{\xi_i}
  -\sum_{i=0}^k\int|u_{\xi_i}-u'_{\xi_i}|
\ge\gamma$}

\noindent whenever $\xi_0,\ldots,\xi_n\in\Gamma$, and the theorem is
still true.

\medskip

\quad{\bf (ii)} Finally, for the general case, set
$M=1+\sup_{\xi<\kappa}\|u_{\xi}\|_{\infty}$, and
$u'_{\xi}=\Bover1Mu_{\xi}$ for $\xi<\kappa$.   Then every $u'_{\xi}$
belongs to $[0,\chi 1]$ and $\int u'_{\xi}\ge\Bover{\delta}M$.   By (i)
there is a $\Gamma\in[\kappa]^{\kappa}$ such that
$\int\prod_{i=0}^ku'_{\xi_i}\ge\Bover{\gamma}{M^{k+1}}$ for all
$\xi_0,\ldots,\xi_k\in\Gamma$;  in which case
$\int\prod_{i=0}^ku_{\xi_i}\ge\gamma$ for all
$\xi_0,\ldots,\xi_k\in\Gamma$.

This completes the proof.
}%end of proof of 525S

\leader{525T}{Corollary}\cmmnt{ ({\smc Argyros \& Kalamidas 82})} (a)
If $\kappa$ is an infinite cardinal and
$k\in\Bbb N$, $(\kappa,\kappa,k)$ is a measure-precaliber triple of
every probability algebra.

(b) If $\kappa$ is a cardinal of uncountable cofinality and
$k\in\Bbb N$, $(\kappa,\kappa,k)$ is a precaliber triple of every
measurable algebra.   In particular, every measurable algebra satisfies
Knaster's condition.

(c) If $\kappa$ is a cardinal of uncountable cofinality,
$(\frak A,\bar\mu)$ is a probability algebra, $k\ge 1$ and
$\ofamily{\xi}{\kappa}{a_{\xi}}$ is a family in $\frak A\setminus\{0\}$,
then there are a $\delta>0$ and a $\Gamma\in[\kappa]^{\kappa}$ such that
$\bar\mu(\inf_{\xi\in I}a_{\xi})\ge\delta$ for every $I\in[\Gamma]^k$.

(d)\discrversionA{\footnote{Revised 2009.}}{}
For any measurable algebra $\frak A$,
$\frak m(\frak A)\ge\frak m_{\text{K}}$;  and if
$\frak m(\frak A)>\omega_1$, then
$\frak m(\frak A)\ge\frak m_{\text{pc}\omega_1}$.
So if $\omega\le\kappa<\frak m_K$, $\kappa$ is a measure-precaliber of
every probability algebra.

\proof{ Really this is just the special case of 525S in which
every $u_{\xi}$ belongs to $\{\chi a:a\in\frak A\}$.

\medskip

{\bf (a)} If $(\frak A,\bar\mu)$
is a probability algebra and $\ofamily{\xi}{\kappa}{a_{\xi}}$ is a
family in $\frak A$ such that
$\inf_{\xi<\kappa}\bar\mu a_{\xi}=\delta>0$, take any
$\gamma\in\ooint{0,\delta^k}$.   Setting $u_{\xi}=\chi a_{\xi}$
for each $\xi$, $\int u_{\xi}\ge\delta$ for each $\xi$,
so there is a $\Gamma\in[\kappa]^{\kappa}$ such that
$\int\prod_{i=1}^ku_{\xi_i}\ge\gamma$ for every
$\xi_1,\ldots,\xi_k\in\Gamma$;  in which case
$\inf_{\xi\in J}a_{\xi}\ne 0$ for every $J\in[\Gamma]^k$.   As
$\ofamily{\xi}{\kappa}{a_{\xi}}$ is
arbitrary, $(\kappa,\kappa,k)$ is a measure-precaliber triple of
$(\frak A,\bar\mu)$.

\medskip

{\bf (b)} This now follows at once from 525Db, since any non-zero
measurable algebra can be given a probability measure.   Taking
$\kappa=\omega_1$ and $k=2$, we have Knaster's condition.

\medskip

{\bf (c)} For the quantitative version, we have only to note that there
must be some $\alpha>0$ such that $\#(\{\xi:\bar\mu a_{\xi}\ge\alpha\}$
has cardinal $\kappa$, and take $\delta<\alpha^k$.

\medskip

{\bf (d)} By (b), $\frak A$ satisfies
Knaster's condition;  it follows at once that
$\frak m(\frak A)\ge\frak m_{\text{K}}$, while $\sat(\frak A)\le\omega_1$.
If $\frak m(\frak A)>\omega_1$, then $\omega_1$ is a precaliber of
$\frak A$ (517Ig) so $\frak m(\frak A)\ge\frak m_{\text{pc}\omega_1}$.
By 525Fb, every infinite cardinal less than $\frak m_{\text{K}}$ is a
measure-precaliber of every probability algebra.
}%end of proof of 525T

\exercises{\leader{525X}{Basic exercises (a)}
%\spheader 525Xa
Let $(X,\Sigma,\mu)$ be any measure space and $\frak A$
its measure algebra.   (i) Show that
$(\frak A^+,\Bsupseteqshort)\prT(\Sigma\setminus\Cal N(\mu),\supseteq)$.
(ii) Show that a pair $(\kappa,\lambda)$ is a downwards precaliber pair
of $\Sigma\setminus\Cal N(\mu)$ iff it is a precaliber pair of
$\frak A$.
%525B 516Xb

\sqheader 525Xb Let $\frak A$ be a measurable algebra.   Show that
$\omega_1$ is a precaliber of $\frak A$ iff {\it either} $\frak A$ is
purely atomic {\it or} $\tau(\frak A)\le\omega$ and
$\cov\Cal N_{\omega}>\omega_1$ {\it or}
$\cov\Cal N_{\omega_1}>\omega_1$.    \Hint{525G, 523F.}
%525C

\sqheader 525Xc (i) Suppose that
$\add\Cal N_{\omega}=\cov\Cal N_{\omega}=\kappa$.   Show that $\kappa$
is not a precaliber of $\frak B_{\omega}$.
(ii)\dvAnew{2009}
Suppose that $\non\Cal N_{\omega}=\frak c$.   Show that $\frak c$ is not a
precaliber of $\frak B_{\omega}$.
%525C %525N

\spheader 525Xd Let $(X,\Sigma,\mu)$ be a complete strictly localizable
measure space and $\frak A$ its measure algebra.   Show that the
supported relation $(\Sigma\setminus\Cal N(\mu),\ni,X)$ has the same
precaliber pairs as the Boolean algebra $\frak A$.
%525C 525Xa

\spheader 525Xe Suppose that $(\kappa,\lambda)$ is a precaliber pair of
every measurable algebra, that $I$ is a set, and that
$X\subseteq\BbbR^I$ is a compact set such that
$\#(\{i:x(i)\ne 0\})<\lambda$ for every $x\in X$.   Show that
$\#(\{i:x(i)\ne 0\})<\kappa$ for every $x$ belonging to the closed
convex hull of $X$ in $\BbbR^I$.   \Hint{461I.}
%525Ca  Marciszewski n84 ?

\spheader 525Xf\dvAnew{2009}
Suppose that $\lambda\le\kappa$ are infinite cardinals,
$(\frak A,\bar\mu)$ is a homogeneous probability algebra,
and that $\gamma<1$ is such that whenever
$\ofamily{\xi}{\kappa}{a_{\xi}}$ is a family in
$\frak A$ and $\bar\mu a_{\xi}\ge\gamma$ for every $\xi<\kappa$,
there is a $\Gamma\in[\kappa]^{\lambda}$ such that
$\{a_{\xi}:\xi\in\Gamma\}$ is centered.   Show that $(\kappa,\lambda)$ is
a measure-precaliber pair of $(\frak A,\bar\mu)$.   \Hint{given that
$\inf_{\xi<\kappa}\bar\mu a_{\xi}>0$, take
$(\frak C,\bar\lambda)=\Tensorhat_m(\frak A,\bar\mu)\cong(\frak A,\bar\mu)$
to be the probability
algebra free product of a large finite number of copies of
$(\frak A,\bar\mu)$, and consider $c_{\xi}=\sup_{j<m}\varepsilon_ja_{\xi}$
for $\xi<\kappa$.}
%525D

\spheader 525Xg\dvAnew{2009}
Let $\frak A$ be a Boolean algebra, and $\lambda$, $\kappa$ cardinals such
that $(\kappa,\lambda)$ is a
measure-precaliber pair of every probability algebra.   Suppose that
$A\subseteq\frak A\setminus\{0\}$ has positive intersection number.
Show that if $\ofamily{\xi}{\kappa}{a_{\xi}}$ is a
family in $A$, then there is a $\Gamma\in[\kappa]^{\lambda}$ such that
$\{a_{\xi}:\xi\in\Gamma\}$ is centered.
%525D

\spheader 525Xh Let $\kappa$ be a cardinal such that ($\alpha$)
$\lambda^{\omega}<\kappa$ for every $\lambda<\kappa$ ($\beta$)
$\lambda^{\omega}<\cf\kappa$ for every $\lambda<\cf\kappa$.   Show that
$\kappa$ is a measure-precaliber of every probability algebra.
%525M 525K  D\v{z}amonja \& Plebanek 04, 4.1

\spheader 525Xi\dvAnew{2016} Show that if $(X,\frak T,\Sigma,\mu)$ is a Radon
measure space and $\mu X>0$, then $\cov\Cal N(\mu)\ge\frak m_{\text{K}}$.
\Hint{assemble 524Md, 524Pc, 525Tb.}
%525T
%{\leader{525Y}{Further exercises (a)}

}%end of exercises

\leader{525Z}{Problem} Can we, in ZFC, find an infinite cardinal
$\kappa$ which is not a measure-precaliber of all probability algebras?
\cmmnt{From 525N we see that a negative answer will require a model of
set theory in which $2^{\kappa}>\kappa^+$ for all strong limit cardinals
$\kappa$ of countable cofinality;  for such models see
{\smc Foreman \& Woodin 91}, {\smc Cummings 92}.}

\leaveitout{(Haydon) If $\sequencen{\kappa_n}$ is a sequence of
precalibers of measurable algebras, is $\sup_n\kappa_n$ a
measure-precaliber of probability algebras?}

\endnotes{
\Notesheader{525}
There seem to be three methods of proving that a cardinal is a
precaliber of a measure algebra.   First, we have the counting arguments
of 516L;  since we know something about the centering numbers of measure
algebras (524Me), this gives us a start (see the proof of 525K).   Next,
we can try to use Martin numbers, as in 517Ig and 525F;  since we can
relate the
Martin number of a measure algebra to the cardinals of \S523 (524Md), we
get the formulation 525J.   In third place, we have arguments based on
the special
structure of measure algebras, using 525H to apply
$\Delta$-system theorems from infinitary combinatorics.   Subject to the
generalized continuum hypothesis, these ideas are enough to answer the
most natural questions (525O).   Without this simplification, they leave
conspicuous gaps.   The most important seems to be 525Z.   Even if we
know all the cardinals
$\add\Cal N_{\kappa}$, $\cov\Cal N_{\kappa}$, $\non\Cal N_{\kappa}$ and
$\cf\Cal N_{\kappa}$ of \S523, we may still not be able to determine
which cardinals are precalibers;   525Xb is an exceptional special case.

I have presented this section with a bias towards measure-precalibers
rather than precalibers.   When there is a difference, the former search
deeper.   `Cofinality $\omega_1$' has a rather special position in this
theory (525Ib), deriving from the combinatorial arguments of 5A1I.
}%end of notes

\discrpage

\frfilename{mt526.tex}
\versiondate{24.1.14}
\copyrightdate{2009}

\def\chaptername{Cardinal functions of measure theory}
\def\sectionname{Asymptotic density zero}

\newsection{526}

In \S491, I devoted some paragraphs to the ideal $\Cal Z$ of subsets
of $\Bbb N$ with asymptotic density zero, as part of an investigation
into equidistributed sequences
in topological measure spaces.   Here I return to $\Cal Z$ to examine
its place in the Tukey ordering of partially ordered sets.   We find
that it lies strictly between
$\NN$ and $\ell^1$ (526B, 526J, 526L)
but in some sense is closer to $\ell^1$
(526Ga).   On the way, I mention the ideal $\CalNwd$
of nowhere dense subsets of $\NN$
(526H-526L) %526H 526I 526J 526K 526L
and ideals of sets with negligible closures
(526I-526M). %526I 526J 526K 526L 526M

\leader{526A}{Proposition} For $I\subseteq\Bbb N$, set
$\nu I=\sup_{n\ge 1}\bover1n\#(I\cap n)$.

(a) $\nu$ is a strictly positive submeasure\cmmnt{ (definition:
392A)} on $\Cal P\Bbb N$.   We have a
metric $\rho$ on $\Cal P\Bbb N$
defined by setting $\rho(I,J)=\nu(I\symmdiff J)$ for all $I$,
$J\subseteq\Bbb N$, under which the Boolean operations $\cup$, $\cap$,
$\symmdiff$ and $\setminus$ and upper asymptotic density
$d^*:\Cal P\Bbb N\to[0,1]$ are uniformly continuous and
$\Cal P\Bbb N$ is complete.

(b) $\Cal Z$ is a separable closed subset of $\Cal P\Bbb N$.

(c) If $\Cal I\subseteq\Cal Z$ is such that $\sum_{I\in\Cal I}\nu I$ is
finite, then $\bigcup\Cal I\in\Cal Z$.

(d) With the subspace topology, $(\Cal Z,\subseteq)$
is a metrizably compactly based directed set\cmmnt{ (definition:  513K)}.

\proof{{\bf (a)} It is elementary to check that $\nu$ is a strictly
positive submeasure.   By 392H,
$\rho$ is a metric under which the Boolean operations are
uniformly continuous.   Since

\Centerline{$|d^*(I)-d^*(J)|\le d^*(I\symmdiff J)\le\nu(I\symmdiff J)$}

\noindent for all $I$, $J\subseteq\Bbb N$, $d^*$ is uniformly
continuous.
Let $\sequence{j}{I_j}$ be a sequence in $\Cal P\Bbb N$ such that
$\rho(I_j,I_{j+1})\le 2^{-j}$ for every $j\in\Bbb N$.   Set
$I=\bigcap_{m\in\Bbb N}\bigcup_{j\ge m}I_j$.   For any $m\in\Bbb N$ and
$n\ge 1$,

\Centerline{$\Bover1n\#(n\cap(I_m\symmdiff I))
\le\Bover1n\sum_{j=m}^{\infty}\#(n\cap(I_j\symmdiff I_{j+1}))
\le\sum_{j=m}^{\infty}\rho(I_j,I_{j+1})\le 2^{-m+1}$,}

\noindent so $\rho(I_m,I)\le 2^{-m+1}$.   Thus $\sequence{j}{I_j}$ is
convergent to $I$;  as $\sequence{j}{I_j}$ is arbitrary, $\Cal P\Bbb N$
is complete (cf.\ 2A4E).

\medskip

{\bf (b)} $\Cal Z$ is a closed subset of $\Cal P\Bbb N$.   \Prf\ If $I$
belongs to the closure $\overline{\Cal Z}$ of $\Cal Z$, and
$\epsilon>0$, let $J\in\Cal Z$ be such that
$\rho(I,J)\le\bover12\epsilon$, and let $m\ge 1$ be such that
$\bover1n\#(J\cap n)\le\bover12\epsilon$ for every $n\ge m$;  then
$\bover1n\#(I\cap n)\le\epsilon$ for every $n\ge m$.   As $\epsilon$ is
arbitrary, $I\in\Cal Z$;  as $I$ is arbitrary, $\Cal Z$ is closed.\ \Qed

$\Cal Z$ is separable because $[\Bbb N]^{<\omega}$ is a
countable dense set.   (If $I\in\Cal Z$ and $n\in\Bbb N$,
$\rho(a,a\cap n)\le\sup_{m>n}\bover1m\#(m\cap I)$.)

\medskip

{\bf (c)} Let $\epsilon>0$.   Then there is a finite
$\Cal I_0\subseteq\Cal I$ such that
$\sum_{I\in\Cal I\setminus\Cal I_0}\nu I\le\epsilon$.   Set
$J=\bigcup\Cal I$,
$J_0=\bigcup\Cal I_0$;  then $J_0\in\Cal Z$ so there is an
$n_0\in\Bbb N$ such
that $\#(J_0\cap n)\le n\epsilon$ for every $n\ge n_0$.
If $n\ge n_0$, then

\Centerline{$\#(J\cap n)
\le\#(J_0\cap n)+\sum_{I\in\Cal I\setminus\Cal I_0}\#(I\cap n)
\le n\epsilon+\sum_{I\in\Cal I\setminus\Cal I_0}n\nu I
\le 2n\epsilon$.}

\noindent As $\epsilon$ is arbitrary, $J\in\Cal Z$.

\medskip

{\bf (d)} $\Cal Z$ is closed under $\cup$, so is a directed set under
$\subseteq$, and $\cup:\Cal Z\times\Cal Z\to\Cal Z$ is continuous.   If
$a\in\Cal Z$, then on $\{b:b\subseteq a\}$ the topology $\frak T_{\rho}$
induced by $\rho$ agrees with the usual compact Hausdorff topology
$\frak S$ of $\Cal P\Bbb N\cong\{0,1\}^{\Bbb N}$.   \Prf\ If $n\in\Bbb N$
and $\rho(b,c)<\bover1{n+1}$, then $b\cap n=c\cap n$;  so $\frak T_{\rho}$
is finer than $\frak S$ on $\Cal P\Bbb N$.   If $\epsilon>0$, there is an
$m\in\Bbb N$ such that $\#(a\cap n)\le n\epsilon$ whenever $n\ge m$;  now
$\rho(b,c)\le\epsilon$ whenever $b$, $c\subseteq a$ and $b\cap m=c\cap m$.
So $\frak S$ is finer than $\frak T_{\rho}$ on $\{b:b\subseteq a\}$.\ \QeD\
Since $\{b:b\subseteq a\}$ is $\frak S$-compact, it is
$\frak T_{\rho}$-compact.

Now suppose that $\sequencen{a_n}$ is a sequence in $\Cal Z$ with
$\frak T_{\rho}$-limit $a$.   Then it has a subsequence
$\sequence{k}{a_{n_k}}$ such that $\rho(a,a_{n_k})\le 2^{-k}$ for every
$k$.   Set $b=\bigcup_{k\in\Bbb N}a_{n_k}$.   Then, given $\epsilon>0$,
let $r$, $m\in\Bbb N$ be such that $2^{-r}\le\epsilon$ and
$\#(n\cap(a\cup\bigcup_{k\le r}a_{n_k}))\le n\epsilon$ for every $n\ge m$;
then

$$\eqalign{\#(n\cap b)
&\le\#(n\cap(a\cup\bigcup_{k\le r}a_{n_k}))
  +\sum_{k=r+1}^{\infty}\#(n\cap a_{n_k}\setminus a)\cr
&\le n\epsilon+\sum_{k=r+1}^{\infty}2^{-k}n
\le 2n\epsilon\cr}$$

\noindent for every $n\ge m$.   So $b\in\Cal Z$ and
$\{a_{n_k}:k\in\Bbb N\}$ is bounded above in $\Cal Z$.
}%end of proof of 526A

\leader{526B}{Proposition}\cmmnt{ ({\smc Fremlin 91})}
$\BbbN^{\Bbb N}\prT\Cal Z\prT\ell^1$.

\proof{{\bf (a)} For $\alpha\in\BbbN^{\Bbb N}$, set

\Centerline{$\phi(\alpha)=\{2^ni:n\in\Bbb N$, $i\le\alpha(n)\}$.}

\noindent Then $\phi(\alpha)\in\Cal Z$, because if $k\in\Bbb N$ then

\Centerline{$\#(m\cap\phi(\alpha))
  \le\sum_{n=0}^k\alpha(n)+\lceil 2^{-k}m\rceil$}

\noindent for every $m$.   Also $\phi:\BbbN^{\Bbb N}\to\Cal Z$
is a Tukey function, because if $\phi(\alpha)\subseteq a\in\Cal Z$ then
$\alpha(n)\le\min\{i:2^ni\notin a\}$ for every $n\in\Bbb N$.   So
$\NN\prT\Cal Z$.

\medskip

{\bf (b)} Give $\Cal Z$ the metric $\rho$ of 526A.
Then $\Cal Z$ is complete and separable and
the lattice operation $\cup$ is uniformly continuous (526Aa).   By 524C,
$(\Cal Z,\subseteq\strprime,[\Cal Z]^{<\omega})
\prGT(\ell^1(\omega),\le,\ell^1(\omega))$.   Since $\Cal Z$
is upwards-directed, $(\Cal Z,\subseteq,\penalty-100\Cal Z)
\equivGT(\Cal Z,\subseteq\strprime,[\Cal Z]^{<\omega})$ (513Id) and
$(\Cal Z,\subseteq,\Cal Z)\prGT(\ell^1,\le,\ell^1)$, that is,
$\Cal Z\prT\ell^1$.
}%end of proof of 526B

\leader{526C}{}\cmmnt{ The next three lemmas are steps on the way to
Theorem 526F.   I give them in much more generality than is required by
that theorem because a couple of them will be useful later,
and I think they are interesting in themselves.   But if
you are reading this primarily for the sake of 526F, you might save time by
looking ahead to the proof there and working backwards, extracting
arguments adequate for the special case of 526E which is actually required.

\medskip

\noindent}{\bf Lemma} Let $\sequencen{(\frak A_n,\bar\mu_n)}$ be a
sequence of purely atomic probability algebras, and
$\frak A=\prod_{n\in\Bbb N}\frak A_n$ the simple product algebra.   Then
there is an order-continuous Boolean homomorphism
$\pi:\frak A\to\Cal P\Bbb N$ such that
$\limsup_{n\to\infty}\bar\mu_na(n)$ is the upper asymptotic density
$d^*(\pi a)$ for every $a\in\frak A$;
consequently, $\lim_{n\to\infty}\bar\mu_na(n)$ is the asymptotic density
$d(\pi a)$ of $\pi a$ if either is defined.

\proof{{\bf (a)} For each $n\in\Bbb N$, let $C_n$ be the set of atoms of
$\frak A_n$, and choose rational numbers $\alpha_n(c)$ such that
$\alpha_n(c)\le\bar\mu_nc$ for each $c\in C_n$, $\sum_{c\in
C_n}\alpha_n(c)>1-2^{-n}$, and $\{c:c\in C_n$, $\alpha_n(c)>0\}$ is
finite.   Express $\alpha_n(c)$ as
$r_n(c)/s_n$ for each $c\in C_n$, where $r_n(c)\in\Bbb N$ and
$s_n\in\Bbb N\setminus\{0\}$;
let $\family{c}{C_n}{I_n(c)}$ be a disjoint family of subsets of
$\Bbb N$ with $\#(I_n(c))=r_n(c)$ for each $c$, and set
$J_n=\bigcup_{c\in C_n}I_n(c)$;
let $\pi_n:\frak A_n\to\Cal PJ_n$ be the Boolean
homomorphism such that $\pi_n(c)=I_n(c)$ for each
$c\in C_n$.   Then

\Centerline{$(1-2^{-n})s_n<s_n\sum_{c\in C_n}\alpha_n(c)=\sum_{c\in
C_n}r_n(c)=\#(J_n)\le s_n$.}

\noindent Note that $\#(J_n)>0$.   Also

$$\eqalign{\#(J_n)(\bar\mu_nd-2^{-n})
&\le\#(J_n)\sum_{c\in C_n,c\subseteq d}\alpha_n(c)\cr
&\le s_n\sum_{c\in C_n,c\subseteq d}\alpha_n(c)
=\sum_{c\in C_n,c\subseteq d}r_n(c)
=\#(\pi_nd)\cr}$$

\noindent and

\Centerline{$(1-2^{-n})\#(\pi_nd)=(1-2^{-n})s_n\sum_{c\in
C_n,c\Bsubseteq d}\alpha_n(c)\le\#(J_n)\cdot\bar\mu_nd$}

\noindent for every $d\in\frak A_n$.   So, for $a\in\frak A$,

\Centerline{$\limsup_{n\to\infty}\bar\mu_na(n)
=\limsup_{n\to\infty}\Bover{\#(\pi_na(n))}{\#(J_n)}$.}

\medskip

{\bf (b)} Let $\sequencen{m_n}$ be a sequence in $\Bbb N$ such that

\Centerline{$m_n\#(J_n)\ge 2^n\max(\#(J_{n+1}),\sum_{i<n}m_i\#(J_i))$}

\noindent for every $n$.   Set $n_k=n$ if
$\sum_{i<n}m_i\le k<\sum_{i\le n}m_i$ where $n\in\Bbb N$;  then
$\lim_{k\to\infty}n_k=\infty$.   Set $l_k=\sum_{i<k}\#(J_{n_i})$, so
that $l_{k+1}-l_k=\#(J_{n_k})$ for each $k$, and let
$\phi_k:\Cal PJ_{n_k}\to\Cal P(l_{k+1}\setminus l_k)$ be a Boolean
isomorphism;  set

\Centerline{$\pi a=\bigcup_{k\in\Bbb N}\phi_k\pi_{n_k}a(n_k)$}

\noindent for $a\in\frak A$, so that $\pi:\frak A\to\Cal P\Bbb N$ is an
order-continuous Boolean isomorphism.

\medskip

{\bf (c)} Let $a\in\frak A$, and set

\Centerline{$\gamma=\limsup_{n\to\infty}\bar\mu_na(n)
=\limsup_{n\to\infty}\Bover{\#(\pi_na(n))}{\#(J_n)}$,}

\Centerline{$\gamma'=\limsup_{l\to\infty}\Bover1l\#(l\cap\pi a)$.}

\noindent Then $\gamma\le\gamma'$.   \Prf\ Setting
$l'_n=\sum_{i<n}m_i\#(J_i)$, we have $\#(l'_{n+1}\cap\pi a)\ge
m_n\#(\pi_na(n))$, while
$l'_{n+1}=l'_n+m_n\#(J_n)\le(1+2^{-n})m_n\#(J_n)$ for each $n$;  but
this means that

$$\gamma'\ge\limsup_{n\to\infty}\bover1{l'_{n+1}}\#(l'_{n+1}\cap\pi a)
\ge\limsup_{n\to\infty}\bover{m_n\#(\pi_na(n))}{(1+2^{-n})m_n\#(J_n)}
=\gamma. \text{ \Qed}$$

\noindent Also $\gamma'\le\gamma$.   \Prf\ Let $\epsilon>0$.   Let
$n^*\ge 1$ be such that $2^{-n^*}\le\epsilon$ and
$\#(\pi_na(n))\le(\gamma+\epsilon)\#(J_n)$ for every $n\ge n^*$.
Suppose that $l\ge l'_{n^*+1}$.   Then $l$ is of the form
$l'_{n+1}+j\#(J_{n+1})+i$ where $n\ge n^*$, $j<m_{n+1}$ and
$i<\#(J_{n+1})$.   Now $l'_{n+1}=l'_n+m_n\#(J_n)$, so

$$\eqalign{\#(l'_{n+1}\cap\pi a)
&\le l'_n+m_n\#(\pi_na(n))
\le l'_n+m_n\#(J_n)(\gamma+\epsilon)\cr
&\le m_n\#(J_n)(\gamma+\epsilon+2^{-n})
\le m_n\#(J_n)(\gamma+2\epsilon)\cr}$$

\noindent by the choice of $m_n$.   Accordingly

$$\eqalignno{\#(l\cap\pi a)
&\le m_n\#(J_n)(\gamma+2\epsilon)+(j+1)\#(\pi_{n+1}a(n+1))\cr
&\le m_n\#(J_n)(\gamma+2\epsilon)+j(\gamma+\epsilon)\#(J_{n+1})
  +\#(J_{n+1})\cr
&\le (\gamma+2\epsilon)l+\#(J_{n+1})
\le (\gamma+2\epsilon)l+2^{-n}m_n\#(J_n)\cr
\displaycause{by the choice of $m_n$}
&\le (\gamma+3\epsilon)l.\cr}$$

\noindent As this is true for any $l\ge l'_{n^*+1}$,
$\gamma'\le\gamma+3\epsilon$;  as $\epsilon$ is arbitrary,
$\gamma'\le\gamma$.\ \Qed

\medskip

{\bf (d)} Thus

\Centerline{$\limsup_{n\to\infty}\bar\mu_na(n)
=\limsup_{n\to\infty}\Bover1n\#(n\cap\pi a)$}

\noindent for every $a\in\frak A$.   But as $\pi$ is a Boolean
homomorphism, it follows at once that

\Centerline{$\liminf_{n\to\infty}\bar\mu_na(n)
=\liminf_{n\to\infty}\Bover1n\#(n\cap\pi a)$}

\noindent for every $a$, so that the limits are equal if either is
defined.
}%end of proof of 526C

\leader{526D}{Lemma} Let $(\frak A,\bar\mu)$ be a semi-finite measure
algebra, and $\kappa\ge\max(\omega,c(\frak A),\tau(\frak A))$ a cardinal.
Let $(\frak B_{\kappa},\bar\nu_{\kappa})$ be the measure algebra of the
usual measure on $\{0,1\}^{\kappa}$, and $\gamma>0$.   Then there is a
function $\theta:\frak A\to\frak B_{\kappa}$ such that

(i) $\theta(\sup A)=\sup\theta[A]$ for every non-empty $A\subseteq\frak A$ such
that $\sup A$ is defined in $\frak A$;

(ii) $\bar\nu_{\kappa}\theta(a)=1-e^{-\gamma\bar\mu a}$ for every
$a\in\frak A$, interpreting $e^{-\infty}$ as $0$;

(iii) if $\familyiI{a_i}$ is a disjoint family in $\frak A$, and
$\frak C_i$ is the closed subalgebra of $\frak B_{\kappa}$ generated by
$\{\theta(a):a\Bsubseteq a_i\}$ for each $i$, then $\familyiI{\frak C_i}$ is
stochastically independent.

\proof{{\bf (a)} By 495M\formerly{4{}95J}, we have exactly this result for some
probability algebra $(\frak B,\bar\lambda)$ in place of
$(\frak B_{\kappa},\bar\nu_{\kappa})$.
Set $\frak A^f=\{a:a\in\frak A$, $\bar\mu a<\infty\}$, and
give $\frak A^f$ its measure metric $\rho$ (323Ad).   Then
$\theta\restrp\frak A^f$ is uniformly continuous for $\rho$ and the measure
metric $\sigma$ of $\frak B$.   \Prf\ If $\epsilon>0$, there is a
$\delta>0$ such that $|e^{-\gamma s}-e^{-\gamma t}|\le\bover12\epsilon$
whenever $s$, $t\in\coint{0,\infty}$ and $|s-t|\le\delta$.   Now if $a$,
$a'\in\frak A$ and $\bar\mu(a\Bsymmdiff a')\le\delta$, set
$b=a\Bcap a'$;  then $\theta(b)\Bsubseteq\theta(a)$ and
$\bar\mu a-\bar\mu b\le\delta$, so

\Centerline{$\sigma(\theta(a),\theta(b))
=\bar\lambda(\theta(a)\Bsetminus\theta(b))
=\bar\lambda\theta(a)-\bar\lambda\theta(b)
=e^{-\gamma\bar\mu b}-e^{-\gamma\bar\mu a}\le\Bover12\epsilon$.}

\noindent Similarly, $\sigma(\theta(a'),\theta(b))\le\bover12\epsilon$ so
$\sigma(\theta(a),\theta(a'))\le\epsilon$.
As $\epsilon$ is arbitrary, this gives the result.\ \Qed

\medskip

{\bf (b)} By 521Eb, there is a set $B\subseteq\frak A^f$, of
cardinal at most $\kappa$, which is dense for $\rho$.   Accordingly
$\theta[B]$ is dense in $f[\frak A^f]$ for $\sigma$.
Taking $\frak D$ to be
the closed subalgebra of $\frak B$ generated by $\theta[B]$,
$\tau(\frak D)\le\kappa$ and $\theta[\frak A^f]\subseteq\frak D$.
But if $a\in\frak A\setminus\frak A^f$ then $\theta(a)=1$, so
$\theta[\frak A]\subseteq\frak D$.   Now there is a measure-preserving
Boolean homomorphism $\phi:\frak D\to\frak B_{\kappa}$ (332N), and
$\phi\theta:\frak A\to\frak B_{\kappa}$ has the properties we need.
}%end of proof of 526D

\vleader{48pt}{526E}{Lemma} Let $\sequencen{(\frak A_n,\bar\mu_n)}$ be a
sequence of finite probability algebras and $\sequencen{\gamma_n}$ a
sequence in $\ooint{0,\infty}$.   Write $P$ for the set

\Centerline{$\{p:p\in\prod_{n\in\Bbb N}\frak A_n$,
$\lim_{n\to\infty}\gamma_n\bar\mu_np(n)=0\}$,}

\noindent with the ordering inherited from the product partial order on
$\prod_{n\in\Bbb N}\frak A_n$.   Then $P\prT\Cal Z$.

\proof{{\bf (a)} By 526D, we can find for each $n$ a probability
algebra $(\frak B_n,\bar\nu_n)$ and a function
$\theta_n:\frak A_n\to\frak B_n$ such that, for all $a$, $a'\in\frak A_n$,

\Centerline{$\theta_n(a\Bcup a')=\theta_n(a)\Bcup\theta_n(a')$,}

\Centerline{$\bar\nu_n\theta_n(a)=1-\exp(-\gamma_n\bar\mu a)$.}

\noindent We may suppose that $\frak B_n$ is generated by
$\theta_n[\frak A_n]$, so is itself finite.   Set
$\frak A=\prod_{n\in\Bbb N}\frak A_n$,
$\frak B=\prod_{n\in\Bbb N}\frak B_n$,
$\theta(p)=\sequencen{\theta_n(p(n))}$ for
$p\in\frak A$;  then $\theta(\sup A)=\sup\theta[A]$ for
any non-empty subset $A$ of $\frak A$.   Set

\Centerline{$Q=\{q:q\in\prod_{n\in\Bbb N}\frak B_n$,
$\lim_{n\to\infty}\bar\nu_nq(n)=0\}$.}

\noindent Then $\theta\restr P$ is a Tukey function from $P$ to $Q$.
\Prf\ $P=f^{-1}[Q]$, because $\lim_{n\to\infty}\gamma_n\xi_n=0$ iff
$\lim_{n\to\infty}1-e^{-\gamma_n\xi_n}=0$.
So $\theta\restr P$ is a function from $P$ to $Q$.   If $q\in Q$,
$A=\{p:p\in\frak A$, $\theta(p)\Bsubseteq q\}$ has a supremum
$p_0\in\frak A$;  now $\theta(p_0)=\sup\theta[A]\Bsubseteq q$, so
$\theta(p_0)\in Q$
and $p_0\in P$ is an upper bound for $A$ in $P$.\ \Qed

\medskip

{\bf (b)} By 526C, we have an order-continuous Boolean homomorphism
$\pi:\frak B\to\Cal P\Bbb N$ such that $\pi(q)\in\Cal Z$
iff $q\in Q$.   Now $\pi\restr Q$ is a Tukey function from $Q$ to
$\Cal Z$.   \Prf\ If $d\in\Cal Z$, set
$B=\{q:q\in\frak B$, $\pi(q)\subseteq d\}$.   Because $\pi$ is an
order-continuous Boolean homomorphism, $B$ contains its supremum, and
$B$ is bounded above in $Q$.\ \Qed

\medskip

{\bf (c)} Thus $\pi\theta\restr P:P\to\Cal Z$ is a Tukey function and
$P\prT\Cal Z$.
}%end of proof of 526E

\leader{526F}{Theorem}
$(\ell^1,\le,\ell^1)\prGT(\NN,\le,\NN)\ltimes(\Cal Z,\subseteq,\Cal Z)$.

\proof{{\bf (a)} Let $Q\subseteq\NN$ be the set of strictly increasing
sequences $\alpha$ such that $\alpha(0)>0$.   For $\alpha\in Q$, set

$$\eqalign{P_{\alpha}
&=\{x:x\in\ell^1,\,\|x\|_{\infty}\le \alpha(0),\,
  \lim_{n\to\infty}2^n\sum_{i=\alpha(n)}^{\infty}x(i)^+=0\}\cr
&=\{x:x\in\ell^{\infty},\,\|x\|_{\infty}\le \alpha(0),\,
  \lim_{n\to\infty}2^n\sum_{i=\alpha(n)}^{\alpha(n+1)-1}x(i)^+=0\}\cr}$$

\noindent because

$$2^n\sum_{i=\alpha(n)}^{\infty}x(i)^+
=\sum_{m=n}^{\infty}2^{n-m}2^m
\sum_{i=\alpha(m)}^{\alpha(m+1)-1}x(i)^+
\le 2\sup_{m\ge n}2^m\sum_{i=\alpha(m)}^{\alpha(m+1)-1}x(i)^+$$

\noindent for every $n$ and $x$.

The point is that $P_{\alpha}\prT\Cal Z$.   \Prf\ For each $n\in\Bbb N$
set $k_n=2^{2n}(\alpha(n+1)-\alpha(n))$,

\Centerline{$V_n=(\alpha(n+1)\setminus \alpha(n))\times
k_n\alpha(0)\subseteq\Bbb N\times\Bbb N$,
\quad$\frak A_n=\Cal PV_n$,}

\noindent and let $\bar\mu_n$ be the uniform probability measure on
$\frak A_n$, so that $\bar\mu_nd=\#(d)/\#(V_n)$ for $d\subseteq V_n$.
For $n\in\Bbb N$ and $x\in\ell^{\infty}$ set

\Centerline{$f_n(x)=\{(i,j):\alpha(n)\le i<\alpha(n+1)$,
$j<k_n\min(\alpha(0),x(i))\}\subseteq V_n$.}

\noindent Then

\Centerline{$|\#(f_n(x))
  -k_n\sum_{\alpha(n)\le i<\alpha(n+1)}\min(\alpha(0),x(i)^+)|
\le\alpha(n+1)-\alpha(n)$,}

\noindent so if $\|x\|_{\infty}\le \alpha(0)$ then

$$\eqalign{\bigl|2^n\alpha(0)(\alpha(n+1)-\alpha(n))\bar\mu f_n(x)
  -2^n\sum_{i=\alpha(n)}^{\alpha(n+1)-1}x(i)^+\bigr|
&\le 2^n(\alpha(n+1)-\alpha(n))/k_n\cr
&=2^{-n}.\cr}$$

\noindent Accordingly

\Centerline{$P_{\alpha}=\{x:x\in\ell^{\infty}$,
$\|x\|_{\infty}\le\alpha(0)$,
$\lim_{n\to\infty}2^n(\alpha(n+1)-\alpha(n))\bar\mu_nf_n(x)
=0\}$.}

Let $\frak A=\prod_{n\in\Bbb N}\frak A_n$ be the simple product of the
Boolean algebras $\frak A_n$, and $I$ the ideal

\Centerline{$\{a:a\in\frak A$,
$\lim_{n\to\infty}
 2^n(\alpha(n+1)-\alpha(n))\bar\mu_na(n)=0\}$}

\noindent of $\frak A$.
For $x\in\ell^{\infty}$, set $f(x)=\sequencen{f_n(x)}$.
Observe that $f:\ell^{\infty}\to\frak A$ is supremum-preserving in the
sense that $f(\sup A)=\sup f[A]$ for any non-empty bounded subset $A$ of
$\ell^{\infty}$.

The last formula for $P_{\alpha}$ shows that $f(x)\in I$ for every
$x\in P_{\alpha}$.   But if $a\in I$,
$A=\{x:x\in P_{\alpha}$, $f(x)\Bsubseteq a\}$ is upwards-directed and
has a supremum $x_0\in\ell^{\infty}$, with $\|x_0\|_{\infty}\le
\alpha(0)$.   Now $f(x_0)=\sup_{x\in A}f(x)\Bsubseteq a$, so $x_0\in
P_{\alpha}$ and is an
upper bound for $A$ in $P_{\alpha}$.   Thus $f\restr P_{\alpha}$ is a
Tukey function from $P_{\alpha}$ to $I$, and $P_{\alpha}\prT I$.
By 526E, $I\prT\Cal Z$, so $P_{\alpha}\prT\Cal Z$.\ \Qed

Thus $(P_{\alpha},\le,P_{\alpha})\prGT(\Cal Z,\subseteq,\Cal Z)$;  it
follows at once that
$(P_{\alpha},\le,\ell^1)\prGT(\Cal Z,\subseteq,\Cal Z)$.

\medskip

{\bf (b)} Now, for $\alpha\in\NN$, take $\tilde P_\alpha=P_{\beta}$ where
$\beta(n)=1+n+\max_{i\le n}\alpha(i)$ for $n\in\Bbb N$.   Then
$\tilde P_{\alpha}\subseteq\tilde P_{\alpha'}$ whenever
$\alpha\le\alpha'$ in
$\NN$, $\bigcup_{\alpha\in\NN}\tilde P_\alpha=\ell^1$ and
$(\tilde P_\alpha,\le,\ell^1)\prGT(\Cal Z,\subseteq,\Cal Z)$ for every
$\alpha$;  so
$(\ell^1,\le,\ell^1)\prGT(\NN,\le,\NN)\ltimes(\Cal Z,\subseteq,\Cal Z)$,
by 512K.
}%end of proof of 526F

\leader{526G}{Corollary} Let $\Cal N$ be the ideal of Lebesgue
negligible subsets of $\Bbb R$.

(a) $\add_{\omega}\Cal Z=\add\Cal N=\add_{\omega}\ell^1$ and
$\cf\Cal Z=\cf\Cal N=\cf\ell^1$.

(b) If $\Cal A\subseteq\Cal Z$ and $\#(\Cal A)<\add\Cal N$, there is a
$J\in\Cal Z$ such that $I\setminus J$ is finite for every $I\in\Cal A$.

%note:  this works for all analytic P-ideals, see {\smc Farah 00},
%5.6.3.

\proof{{\bf (a)(i)} Putting 526B and 513Ie together, we see that

\Centerline{$\add_{\omega}\BbbN^{\Bbb N}\ge\add_{\omega}\Cal Z
\ge\add_{\omega}\ell^1$,}

\noindent that is,

\Centerline{$\frak b\ge\add_{\omega}\Cal Z\ge\add\Cal N$}

\noindent (522A, 524I).   Next, we can deduce from 526F that
$\add_{\omega}\ell^1
\ge\min(\add_{\omega}\BbbN^{\Bbb N},\add_{\omega}\Cal Z)$.   \Prf\ Let
$(\phi,\psi)$ be a Galois-Tukey connection from
$(\ell^1,\le,\ell^1)$ to

\Centerline{$(\NN,\le,\NN)\ltimes(\Cal Z,\subseteq,\Cal Z)
=(\NN\times\Cal Z^{\NN},T,\NN\times\Cal Z)$,}

\noindent where

\Centerline{$T=\{((p,f),(q,a)):p\le q$ in $\NN$, $f(q)\subseteq a\in\Cal
Z\}$.}

\noindent We can interpret $\phi$ as a pair $(\phi_1,\phi_2)$ where
$\phi_1$ is a function from $\ell^1$ to $\NN$ and $\phi_2$ is a function
from $\ell^1\times\NN$ to $\Cal Z$, and saying that $(\phi,\psi)$ is a
Galois-Tukey connection means just that

\Centerline{if $\phi_1(x)\le q$ and $\phi_2(x,q)\subseteq a$ then
$x\le\psi(q,a)$.}

Now suppose that $A\subseteq\ell^1$ and
$\#(A)<\min(\add_{\omega}\NN,\add_{\omega}\Cal Z)$.
Then there is a sequence
$\sequencen{q_n}$ in $\NN$ such that for every $x\in A$ there is an
$n\in\Bbb N$ such that $\phi_1(x)\le q_n$.   Next, for each $n\in\Bbb N$
there is a sequence
$\sequence{m}{a_{nm}}$ in $\Cal Z$ such that for every $x\in A$ there is
an $m\in\Bbb N$ such that $\phi_2(x,q_n)\subseteq a_{nm}$.   In this
case, $B=\{\psi(q_n,a_{nm}):m$, $n\in\Bbb N\}$ is a countable subset of
$\Cal Z$, and for every $x\in A$ there are $m$, $n\in\Bbb N$ such that
$\phi_1(x)\le q_n$ and $\phi_2(x,q_n)\subseteq a_{nm}$, so that
$x\le\psi(q_n,a_{nm})\in B$.   As $A$ is arbitrary,
$\add_{\omega}\ell^1\ge\min(\add_{\omega}\NN,\add_{\omega}\Cal Z)$.\
\Qed

Thus we have
$\add\Cal N\ge\min(\frak b,\add_{\omega}\Cal Z)=\add_{\omega}\Cal Z$,
and $\add_{\omega}\Cal Z=\add\Cal N$.   And we know from 524I, with
$\kappa=\omega$ there, that $\add\Cal N=\add_{\omega}\ell^1$


\medskip

\quad{\bf (ii)} On the other hand, 524I, 526F, 512Da and 512Jb tell us
that

$$\eqalign{\cf\Cal N
&=\cf\ell^1
=\cov(\ell^1,\subseteq,\ell^1)
\le\cov((\NN,\le,\NN)\ltimes(\Cal Z,\subseteq,\Cal Z))\cr
&=\max(\cov(\NN,\le,\NN),\cov(\Cal Z,\subseteq,\Cal Z))
=\max(\frak d,\cf\Cal Z).\cr}$$

\noindent But from 526B we see that $\frak d\le\cf\Cal Z\le\cf\ell^1$,
so $\cf\Cal Z=\cf\Cal N$, while 524I tells us that $\cf\Cal N=\cf\ell^1$.

\medskip

{\bf (b)} By (a), there is a countable set $\Cal D\subseteq\Cal Z$ such
that every member of $\Cal A$ is included in a member of $\Cal D$.   By
491Ae, there is a $J\in\Cal Z$ such that
$I\setminus J$ is finite for every $I\in\Cal D$;  this $J$ serves.
}%end of proof of 526G

\leader{526H}{}\cmmnt{ I turn now to ideals of nowhere dense sets.

\medskip

\noindent}{\bf Proposition}
Let $\CalNwd$ be the ideal of nowhere dense subsets
of $\NN$ and $\Cal M$ the ideal of meager subsets of $\NN$.

(a) $\CalNwd$ is isomorphic, as partially ordered set, to
$(\CalNwd)^{\Bbb N}$.

(b) $(\CalNwd,\subseteq\strprime,[\CalNwd]^{\le\omega})
\equivGT(\Cal M,\subseteq,\Cal M)$.

(c) $\CalNwd\prT\ell^1$.

(d) Let $X$ be a set and $\Cal V$ a countable family of subsets of $X$.
Set

\Centerline{$\Cal D=\{D:D\subseteq X$, for every $V\in\Cal V$ there is a
$V'\in\Cal V$ such that $V'\subseteq V\setminus D\}$.}

\noindent Then $\Cal D\prT\CalNwd$.

(e) If $X$ is any non-empty Polish space without isolated points, and
$\CalNwd(X)$ is the ideal of nowhere dense subsets of $X$, then
$\CalNwd\equivT\CalNwd(X)$.

(f) If $X$ is a compact metrizable space and
$\Cal C_{\text{nwd}}$ is the family of closed nowhere dense
subsets of $X$ with the Fell\cmmnt{ (or Vietoris)}
topology\cmmnt{ (4A2T)},
then $(\Cal C_{\text{nwd}},\subseteq)$ is a metrizably compactly based
directed set.

\cmmnt{\medskip

\noindent{\bf Remark} Recall that if $R$ is any relation then
$R\strprime$ is the relation
$\{(x,B):(x,y)\in R$ for some $y\in B\}$;  see 512F-512G.
}%end of comment

\proof{ Enumerate $S=\bigcup_{n\in\Bbb N}\BbbN^n$
as $\sequencen{\sigma_n}$.
For $\sigma\in S$ write
$I_{\sigma}=\{\alpha:\sigma\subseteq\alpha\in\NN\}$.

\medskip

{\bf (a)} Define $\phi:\CalNwd\to\CalNwd^{\Bbb N}$ by setting
$\phi(F)(n)=\{\alpha:\fraction{n}^{\smallfrown}\alpha\in F\}$ (notation:
5A1C).   Then $\phi$ is an isomorphism
between $\CalNwd$ and $\CalNwd^{\Bbb N}$.

\medskip

{\bf (b)(i)} Choose $\phi:\Cal M\to\CalNwd^{\Bbb N}$ such that
$M\subseteq\bigcup_{n\in\Bbb N}\phi(M)(n)$ for every $M\in\Cal M$.
Then $\phi$ is a Tukey function so
$\Cal M\prT\CalNwd^{\Bbb N}\cong\CalNwd$, that is,
$(\Cal M,\subseteq,\Cal M)\prGT(\CalNwd,\subseteq,\CalNwd)$.   By 513Id
and 512Gb,

\Centerline{$(\Cal M,\subseteq,\Cal M)
\equivGT(\Cal M,\subseteq\strprime,[\Cal M]^{\le\omega})
\prGT(\CalNwd,\subseteq\strprime,[\CalNwd]^{\le\omega})$.}

\woddheader{526H}{4}{2}{2}{42pt}

\quad{\bf (ii)} For $n\in\Bbb N$ and $\tau\in\BbbN^n$, define
$g_{\tau}:\NN\to\NN$ by saying that
$g_{\tau}(\alpha)=\tau^{\smallfrown}\alpha$.
Note that $g_{\tau}$ is a homeomorphism between
$\NN$ and $I_{\tau}$, so that $g_{\tau}[F]$ and $g_{\tau}^{-1}[F]$ are
nowhere dense whenever $F$ is.

Now for any $F\in\CalNwd$ we can find a $\phi(F)\in\CalNwd$ such that
$F\subseteq\phi(F)$ and for every $\sigma\in S$ either
$I_{\sigma}\cap\phi(F)=\emptyset$
or there is a $\tau\in S$, extending $\sigma$, such that
$g_{\tau}[F]\subseteq\phi(F)$.   \Prf\ Choose $\sequencen{\tau_n}$,
$\sequencen{\upsilon_n}$ inductively,
as follows.  Given that $I_{\upsilon_i}\cap(F\cup
g_{\tau_j}[F])=\emptyset$ for all $i$, $j<n$, set
$E=I_{\sigma_n}\cap(F\cup\bigcup_{j<n}g_{\tau_j}[F])$.   If
$E=\emptyset$ set $\upsilon_n=\sigma_n$ and $\tau_n=\emptyset$, so that
$g_{\tau_n}[F]=F$.
If $E$ is not empty, it is still nowhere dense, so we can find
$\upsilon_n\supseteq\sigma_n$ such that
$I_{\upsilon_n}\cap E=\emptyset$.   In this case,
$\bigcup_{i\le n}I_{\upsilon_i}$ is a closed set not including
$I_{\sigma_n}$, so we can find a $\tau_n\supseteq\sigma_n$ such that
$I_{\tau_n}\cap\bigcup_{i\le n}I_{\upsilon_i}=\emptyset$,
and $I_{\upsilon_i}\cap g_{\tau_n}[F]=\emptyset$ for $i\le n$.
Thus in both cases we shall have
$\bigcup_{i\le n}I_{\upsilon_i}\cap(F\cup\bigcup_{j\le n}g_{\tau_j}[F])
=\emptyset$,
and the induction proceeds.

Set $\phi(F)=\overline{F\cup\bigcup_{j\in\Bbb N}g_{\tau_j}[F]}$.
Because $\upsilon_i\supseteq\sigma_i$ and $\phi(F)\cap I_{\upsilon_i}$ is
empty for every $i\in\Bbb N$, $\phi(F)\in\CalNwd$.
If $\sigma\in S$ is such that $\phi(F)$ meets $I_{\sigma}$, there is
an $n\in\Bbb N$ such that $\sigma=\sigma_n$;  now we cannot have
$\upsilon_n=\sigma_n$ so we must have
$\tau_n\supseteq\sigma_n$ and $g_{\tau_n}[F]\subseteq\phi(F)$.   Thus we
have found an suitable set $\phi(F)$.\ \Qed

For each $M\in\Cal M$ let $\Cal E_M$ be a non-empty countable family of
closed nowhere dense sets covering $M$, and set
$\psi(M)=\{g_{\tau}^{-1}[E]:E\in\Cal E_M$, $\tau\in S\}$.
Then $(\phi,\psi)$ is a Galois-Tukey connection from
$(\CalNwd,\subseteq\strprime,[\CalNwd]^{\le\omega})$ to
$(\Cal M,\subseteq,\Cal M)$.   \Prf\ Suppose that $F\in\CalNwd$ and
$M\in\Cal M$ are such that $\phi(F)\subseteq M$.
If $F=\emptyset$ then certainly there is an $F'\in\psi(M)$ covering $F$.
Otherwise, $\phi(F)$ is a non-empty closed set included in the union of
the countable set $\Cal E_M$ of closed sets.   By Baire's theorem, there
must be a $\sigma\in S$ and an $E\in\Cal E_M$ such that
$\emptyset\ne\phi(F)\cap I_{\sigma}\subseteq E$.
In this case, there is a $\tau\supseteq\sigma$ such that
$g_{\tau}[F]\subseteq\phi(F)$, so that
$g_{\tau}[F]\subseteq E$ and $F\subseteq g_{\tau}^{-1}[E]\in\psi(M)$ and
$F\subseteq\strprime\psi(M)$.   As $F$ and $M$ are arbitrary,
$(\phi,\psi)$ is a Galois-Tukey connection.\ \Qed

\medskip

\quad{\bf (iii)} Thus we have

\Centerline{$(\Cal M,\subseteq,\Cal M)
\prGT(\CalNwd,\subseteq\strprime,[\CalNwd]^{\le\omega})
\prGT(\Cal M,\subseteq,\Cal M)$}

\noindent and $(\Cal M,\subseteq,\Cal M)
\equivGT(\CalNwd,\subseteq\strprime,[\CalNwd]^{\le\omega})$.

\medskip

{\bf (c)} We can use the idea of 522O.   Let $\sequencen{U_n}$ enumerate
a base for the topology of $\NN$ not containing $\emptyset$.
By 522N, there is for each
$n\in\Bbb N$ a countable family $\Cal V_n$
of open subsets of $U_n$ such that $\bigcap\Cal V\ne\emptyset$ for every
$\Cal V\in[\Cal V_n]^{\le 2^n}$ and every dense open subset of $U_n$
includes some member of $\Cal V_n$.   Enumerate $\Cal V_n$
as $\sequence{m}{U_{nm}}$.

For each $F\in\CalNwd$ let $f_F:\Bbb N\to\Bbb N$ be such that
$F\cap U_{n,f_F(n)}=\emptyset$ for every $n\in\Bbb N$, and for $n$,
$i\in\Bbb N$ set

$$\eqalign{\phi(F)(2^n(2i+1)-1)
&=2^{-n}\text{ if }f_F(n)=i,\cr
&=0\text{ otherwise}.\cr}$$

\noindent Then
$\sum_{i=0}^{\infty}\phi(F)(i)=2$ for each $F$, so
we have a function $\phi:\CalNwd\to\ell^1$.

Suppose that $x\in\ell^1$.   Set
$\Cal A=\{F:F\in\CalNwd$, $\phi(F)\le x\}$ and $E=\bigcup\Cal A$.
The set

\Centerline{$K=\{n:\#(\{i:x(2^n(2i+1)-1)\ge 2^{-n}\})\ge 2^n\}$}

\noindent is finite;   set $k=\sup(\{0\}\cup K)$.   If $n>k$,
then $\#(\{f_F(n):F\in\Cal A\})<2^n$, so
$\bigcap_{F\in\Cal A}U_{n,f_F(n)}$ is a non-empty open subset
of $U_n$ disjoint from $\bigcup_{F\in\Cal A}F=E$.   Thus
$\{n:U_n\subseteq\overline{E}\}\subseteq\{0,\ldots,k\}$ is finite,
and therefore in fact is empty, that is, $E\in\CalNwd$.

As $x$ is arbitrary, $\phi:\CalNwd\to\ell^1$ is a Tukey function, and
witnesses that $\CalNwd\prT\ell^1$.

\medskip

{\bf (d)} If $\Cal V=\emptyset$ then $\Cal D=\Cal PX$ has a greatest
element and the result is trivial (any function from $\Cal D$ to
$\CalNwd$ will be a Tukey function).   Otherwise,
choose a function $h:S\to\Cal V\cup\{X\}$ such that $h(\emptyset)=X$
and $\sequence{i}{h(\sigma^{\smallfrown}\fraction{i})}$ runs over
$\{V:h(\sigma)\supseteq V\in\Cal V\}$ for every $\sigma\in S$.
Note that $h(\tau)\subseteq h(\sigma)$ whenever $\tau\supseteq\sigma$,
and that $\{h(\tau):\sigma\subseteq\tau\in\BbbN^n\}=\{V:V\in\Cal V$,
$V\subseteq h(\sigma)\}$ whenever $m\in\Bbb N$, $\sigma\in\BbbN^m$ and
$n>m$.
For each $D\in\Cal D$ we can choose a sequence $\sequencen{\tau_{Dn}}$
in $S$ such that $\tau_{Dn}\supseteq\sigma_n$ and $D\cap h(\tau_{Dn})$
is empty and $\#(\tau_{Dn})\ge n$ for
every $n\in\Bbb N$.
Set $\phi(D)=\NN\setminus\bigcup_{n\in\Bbb N}I_{\tau_{Dn}}$, so that
$\phi(D)\in\CalNwd$.

Take any $F\in\CalNwd$, and set $D_0=\bigcup\{D:D\in\Cal D$,
$\phi(D)\subseteq F\}$.   Then $D_0\in\Cal D$.   \Prf\ Let $V\in\Cal V$.
Let $\upsilon\in\BbbN^1$ be such that
$h(\upsilon)=V$.   Take $\tau\supseteq\upsilon$ such that $F\cap
I_{\tau}=\emptyset$.   \Quer\ If $D_0\cap h(\tau)\ne\emptyset$, then
there is a $D\in\Cal D$ such that $\phi(D)\subseteq F$
and $D\cap h(\tau)\ne\emptyset$.   $I_{\tau}\cap\phi(D)$ is empty, that
is, $I_{\tau}\subseteq\bigcup_{n\in\Bbb N}I_{\tau_{Dn}}$;  because
$\#(\tau_{Dn})\ge n$ for every $n$, this can happen only because there
is some $n\in\Bbb N$ such that $\tau_{Dn}\subseteq\tau$.   But this
means that $D\cap h(\tau)\subseteq D\cap h(\tau_{Dn})=\emptyset$, which
is impossible.\ \BanG\  Thus $D_0\cap h(\tau)$ is
empty, and $h(\tau)$ is a member of $\Cal V$ included in
$V\setminus D_0$.   As $V$ is arbitrary, $D_0\in\Cal V$.\ \Qed

As $F$ is arbitrary, $\phi$ is a Tukey function and $\Cal D\prT\CalNwd$,
as claimed.

\medskip

{\bf (e)(i)} Taking $\Cal V$ to be a countable base for the topology of $X$
not containing $\emptyset$, we have

\Centerline{$\CalNwd(X)=\{F:F\subseteq X$, for every $V\in\Cal V$ there is a
$V'\in\Cal V$ such that $V'\subseteq V\setminus F\}$,}

\noindent so (d) tells us that $\CalNwd(X)\prT\CalNwd$.

\medskip

\quad{\bf (ii)} $X$ has a dense subset $Y$ which is
homeomorphic to $\NN$ (5A4Le).   Let $\CalNwd(Y)$ be the family of nowhere
dense subsets of $Y$.   For $F\in\CalNwd(Y)$ let $\phi(F)$ be its closure in
$X$.   Then $\phi$ is a Tukey function from $\CalNwd(Y)$ to $\CalNwd(X)$, so
$\CalNwd\cong\CalNwd(Y)\prT\CalNwd(X)$.

\medskip

{\bf (f)} By 4A2Tg,
the Fell topology on the family $\Cal C$ of all closed subsets of $X$
is compact and metrizable.   $E\cup F\in\Cal C_{\text{nwd}}$ for all $E$,
$F\in\Cal C_{\text{nwd}}$, and
$\cup:\Cal C_{\text{nwd}}\times\Cal C_{\text{nwd}}\to\Cal C_{\text{nwd}}$
is continuous (4A2T(b-ii)).   If $F\in\Cal C_{\text{nwd}}$, the set
$\{E:E\in\Cal C_{\text{nwd}}$, $E\subseteq F\}
=\{E:E\in\Cal C$, $E\cup F=F\}$ is closed in $\Cal C$, therefore compact.
Now suppose that $\sequence{k}{E_k}$ is a sequence in $\Cal C_{\text{nwd}}$
converging to $E\in\Cal C_{\text{nwd}}$.   If $X=\emptyset$ then of course
$\{E_k:k\in\Bbb N\}$ is bounded above in $\Cal C_{\text{nwd}}$.
Otherwise, let $\sequencen{U_n}$ run over a base for the topology of $X$
not containing $\emptyset$, and choose $\sequencen{k_n}$, $\sequencen{V_n}$
inductively, as follows.   Given $k_i\in\Bbb N$ for $i<n$, let
$V_n\subseteq U_n$ be a non-empty open set such that
$\overline{V}_n\cap(E\cup\bigcup_{i<n}E_{k_i})=\emptyset$;  given that
$E\cap\overline{V}_i=\emptyset$ for $i\le n$, choose $k_n\ge n$ such that
$E_{k_n}\cap\bigcup_{i\le n}\overline{V}_i$ is empty.   (This is possible
because $\bigcup_{i\le n}\overline{V}_i$ is compact, so the family of
sets disjoint from it is open in the Fell topology.)   Continue.   At
the end of the induction, $G=\bigcup_{n\in\Bbb N}V_n$ is a dense open set
disjoint from $\bigcup_{n\in\Bbb N}E_{k_n}$, so $X\setminus G$ is an upper
bound for $\{E_{k_n}:n\in\Bbb N\}$ in $\Cal C_{\text{nwd}}$.   Thus all the
conditions of 513K are satisfied, and $\Cal C_{\text{nwd}}$ is metrizably
compactly based.
}%end of proof of 526H

\leader{526I}{}\cmmnt{ A related type of ideal is the following.   I
express the result in more general form because it has some measure
theory in it.

\medskip

\noindent}{\bf Proposition}\cmmnt{ ({\smc Fremlin 91})} Let $X$ be a
second-countable topological space and $\mu$ a $\sigma$-finite
topological measure on $X$.
Let $\Cal E$ be the ideal of subsets of $X$ with negligible closures.
Then, writing $\CalNwd$ for the ideal of nowhere dense subsets of $\NN$,
$\Cal E\prT\CalNwd$ and $\Cal E\prT\Cal Z$.

\proof{{\bf (a)} If $\mu X=0$ then $\Cal E$ has a greatest element and
the result is trivial.
Otherwise, there is a probability measure on $X$ with the same
measurable sets and the same
negligible sets as $\mu$ (215B(vii));  so we may suppose that $\mu$
itself is a probability measure.   Let $\Cal U$ be a countable base for
the topology of $X$,
containing $X$ and closed under finite unions.

\medskip

{\bf (b)} For $k\in\Bbb N$ let $\Cal V_k$ be the countable set
$\{V:V\in\Cal U$, $\mu V>1-2^{-k}\}$.   Set

\Centerline{$\Cal E_k=\{E:E\subseteq X$, for every $V\in\Cal V_k$ there
is a $U\in\Cal V_k$ such that $U\subseteq V\setminus E\}$.}

\noindent Then $\Cal E=\bigcap_{k\in\Bbb N}\Cal E_k$.   \Prf\ Because
$X\in\Cal V_k$, $\mu\overline{E}\le 2^{-k}$ for every $E\in\Cal E_k$, so
$\bigcap_{k\in\Bbb N}\Cal E_k\subseteq\Cal E$.
On the other hand, if $E\in\Cal E$ and $k\in\Bbb N$ and $V\in\Cal V_k$,
then $\mu(V\setminus\overline{E})>1-2^{-k}$ and
$\Cal U'=\{U:U\in\Cal U$, $U\subseteq V\setminus E\}$ has union
$V\setminus\overline{E}$.   As $\Cal U'$ is countable, there is a finite
$\Cal U'_1\subseteq\Cal U'$ such that $U=\bigcup\Cal U'_1$ has measure
greater than $1-2^{-k}$, so that
$U\in\Cal V_k$ and $U\subseteq V\setminus E$.   As $V$ is arbitrary,
$E\in\Cal V_k$;  as $E$ and $k$ are arbitrary,
$\Cal E\subseteq\bigcap_{k\in\Bbb N}\Cal E_k$.\ \QeD

This means that the map $E\mapsto(E,E,E,\ldots)$ is a Tukey function
from $\Cal E$ to $\prod_{k\in\Bbb N}\Cal E_k$, so that
$\Cal E\prT\prod_{k\in\Bbb N}\Cal E_k$.   At the same time,
$\Cal E_k\prT\CalNwd$ for every $k$, by 526Hd.   So
$\Cal E\prT\CalNwd^{\Bbb N}\cong\CalNwd$ (513Eg, 526Ha).

\medskip

{\bf (c)} Let $\frak A$ be the countable subalgebra of $\Cal PX$
generated by $\Cal U$.   Then there is a Boolean homomorphism
$\pi:\frak A\to\Cal P\Bbb N$ such that $d(\pi E)$ is defined
and equal to $\mu E$ for every $E\in\frak A$.   \Prf\ This is easy to
prove directly (see 491Xu), but we can also argue as follows.
Let $\sequencen{\frak A_n}$ be a non-decreasing
sequence of finite subalgebras
of $\frak A$ with union $\frak A$.   By 526C, we have a Boolean
homomorphism $\pi':\prod_{n\in\Bbb N}\frak A_n\to\Cal P\Bbb N$ such that
$d(\pi'\sequencen{E_n})=\lim_{n\to\infty}\mu E_n$ whenever
$E_n\in\frak A_n$ for every $n$ and the limit on the right is defined.
For each $n\in\Bbb N$ let $\pi_n:\frak A\to\frak A_n$ be a Boolean
homomorphism extending the identity
homomorphism on $\frak A_n$ (314K, or otherwise);  set
$\pi E=\pi'\sequencen{\pi_nE}$ for $E\in\frak A$;  this works.\ \Qed

Let $\sequencen{V_n}$ be a sequence running over the closed sets
belonging to $\frak A$.
Let $\sequencen{k_n}$ be a strictly increasing sequence in $\Bbb
N\setminus\{0\}$ such that $\Bover1{k_n}\#(k_n\cap\pi E)\ge\mu E-2^{-n}$
whenever $E$ belongs to the subalgebra $\frak B_n$ of
$\frak A$ generated by $\{V_i:i\le n\}$.   Define
$\phi:\Cal E\to\Cal P\Bbb N$ by setting

\Centerline{$\phi(E)
=\bigcap\{k_i\cup\pi V_i:i\in\Bbb N$, $E\subseteq V_i\}$.}

\noindent Then $\phi$ is a Tukey function from $\Cal E$ to $\Cal Z$.

\medskip

\Prf\ {\bf (i)} If $E\in\Cal E$ and $\epsilon>0$ there is a $U\in\Cal U$
such that $U\subseteq X\setminus E$ and $\mu U\ge 1-\epsilon$.
Let $i\in\Bbb N$ be such that $X\setminus U=V_i$;  then
$\phi(E)\subseteq k_i\cup\pi V_i$, so

\Centerline{$d^*(\phi(E))\le d^*(\pi V_i)=\mu V_i\le\epsilon$.}

\noindent As $\epsilon$ is arbitrary, $\phi(E)\in\Cal Z$.   Thus $\phi$
is a function from $\Cal E$ to $\Cal Z$.

\medskip

\quad{\bf (ii)} Take any $\Cal A\subseteq\Cal E$, and set
$F=\overline{\bigcup\Cal A}$, $a=\bigcup_{E\in\Cal E}\phi(E)$.
If $n\in\Bbb N$ and $i\in k_n\setminus a$, then $i\notin\phi(E)$ for
every $E\in\Cal A$, so for every $E\in\Cal A$ there is a $j<n$ such
that $E\subseteq V_j$ and $i\notin\pi V_j$.
Set $F_{ni}=\bigcup\{V_j:j<n$, $i\notin\pi V_j\}$, so that $i\notin\pi
F_{ni}$, while $\bigcup\Cal A\subseteq F_{ni}$ and $F\subseteq F_{ni}$.
Set $F_n=\Bbb N\cap\bigcap_{i\in k_n\setminus a}F_{ni}$,
so that $F\subseteq F_n$
and no member of $k_n\setminus a$ belongs fo $\pi F_n$, that is,
$k_n\cap\pi F_n\subseteq a$.   Note that
$F_n\in\frak B_n$.   So we have

\Centerline{$\Bover1{k_n}\#(k_n\cap a)\ge\Bover1{k_n}\#(k_n\cap\pi
F_n)\ge\mu F_n-2^{-n}\ge\mu F-2^{-n}$.}

\noindent This means that $d^*(a)\ge\mu F$.   So if $\{\phi(E):E\in\Cal
A\}$ is bounded above in $\Cal Z$, $\Cal A$ must be bounded above in
$\Cal E$;  that is, $\phi$ is a
Tukey function.\ \Qed

Thus $\Cal E\prT\Cal Z$ also.
}%end of proof of 526I

\leader{526J}{Proposition} Let $\Cal E_{\text{Leb}}$ be the ideal of
subsets of $\Bbb R$ whose closures are Lebesgue negligible.
Then $\NN\prT\Cal E_{\text{Leb}}$ but $\Cal E_{\text{Leb}}\not\prT\NN$;
consequently $\Cal Z\not\prT\NN$, $\CalNwd\not\prT\NN$ and
$\ell^1\not\prT\NN$.

\proof{{\bf (a)} Enumerate $\Bbb Q\cap[0,1]$ as $\sequence{i}{q_i}$.
Define $\phi:\NN\to\Cal E_{\text{Leb}}$ by setting
$\phi(f)(n)=\{n+q_i:n\in\Bbb N$, $i\le f(n)\}$.   Then it is
easy to see that $\phi$ is a Tukey function, because if $F\subseteq\NN$
and $\{f(n):f\in F\}$ is unbounded, then $\bigcup_{f\in F}\phi(f)$ is
dense in $[n,n+1]$ so
does not belong to $\Cal E_{\text{Leb}}$.

\medskip

{\bf (b)} Let $\psi:\Cal E_{\text{Leb}}\to\NN$ be any function.   Let
$\mu$ be Lebesgue measure on $\Bbb R$, and choose $\sequencen{f(n)}$
inductively in $\Bbb N$ such that
$\mu^*\{t:t\in[0,1]$, $\psi(\{t\})(i)\le f(i)$ for every $i\le
n\}>\bover12$ for every $n$.   Set

\Centerline{$A_n=\{t:t\in[0,1]$, $\psi(\{t\})(i)\le f(i)$ for every
$i\le n\}$,
\quad$F=\bigcap_{n\in\Bbb N}\overline{A_n}$}

\noindent so that $\mu F\ge\bover12$.   Let $\sequencen{U_n}$ enumerate
the set of open intervals of $\Bbb R$, meeting $F$,
with rational endpoints,
and for each $n\in\Bbb N$ choose $t_n\in A_n\cap U_n$.   Then
$\psi(\{t_n\})(i)\le f(i)$ whenever $n\ge i$, so
$\{\psi(\{t_n\}):n\in\Bbb N\}$ is bounded above in $\NN$;
but $\overline{\{t_n:n\in\Bbb N\}}$ includes $F$, so
$\{\{t_n\}:n\in\Bbb N\}$ is not bounded above in $\Cal E_{\text{Leb}}$.
Thus $\psi$ cannot be a Tukey function.

\medskip

{\bf (c)} Accordingly $\Cal E_{\text{Leb}}\not\prT\NN$;  since
$\Cal E_{\text{Leb}}\prT\Cal Z\prT\ell^1$ and
$\Cal E_{\text{Leb}}\prT\CalNwd$ (526I, 526B), $\Cal Z\not\prT\NN$,
$\CalNwd\not\prT\NN$ and $\ell^1\not\prT\NN$.
}%end of proof of 526J

\leader{526K}{Proposition} Let $\CalNwd$ be the ideal of nowhere dense
subsets of $\NN$.   Then
$\Cal Z\not\prT\CalNwd$, so $\Cal Z\not\prT\Cal E_{\text{Leb}}$
and $\ell^1\not\prT\CalNwd$.

\proof{ Let $\phi:\Cal Z\to\CalNwd$ be any function.
Let $\sequencen{U_n}$ enumerate a base for the topology of $\NN$ which
contains $\emptyset$ and is closed under finite
unions.   For each $n\in\Bbb N$, set

\Centerline{$a_n=\{i:i\in\Bbb N$, $\phi(a)\cap U_n\ne\emptyset$ whenever
$i\in a\in\Cal Z\}$.}

Set

\Centerline{$a=\{\min(a_n\setminus n^2):n\in\Bbb N$, $a_n\not\subseteq
n^2\}$}

\noindent (interpreting $n^2$ in the formula above as a member of
$\Bbb N$ rather than as a subset of $\BbbN^2$).
Then $a\in\Cal Z$ and $a\cap a_n\ne\emptyset$ whenever $a_n$ is
infinite.   Set
$K=\{n:n\in\Bbb N$, $\phi(a)\cap U_n=\emptyset\}$, so that $K$ is infinite
and $\bigcup_{n\in K}U_n=\NN\setminus\overline{\phi(a)}$ is dense, while
$a_n$ is finite for every $n\in K$
(since otherwise there is an $i\in a\cap a_n$, and $\phi(a)\cap U_n$
will not be empty).   For $n\in\Bbb N$, $\bigcup_{m\in K\cap n}U_m$
belongs to $\Cal U$;  let
$r(n)\in\Bbb N$ be such that $U_{r(n)}=\bigcup_{m\in K\cap n}U_m$.
Then $r(n)\in K$ for every $n$, so $a_{r(n)}$ is always finite.
Take a strictly increasing sequence $\sequencen{k_n}$ in $K$ such that
$a_{r(n)}\subseteq k_n$ for every $n$.   For $i<k_0$, set
$b_i=\{i\}$;  for
$k_n\le i<k_{n+1}$, choose $b_i\in\Cal Z$ such that $i\in b_i$ and
$\phi(b_i)\cap U_{r(n)}$ is empty (such exists because
$i\notin a_{r(n)}$).

Examine $E=\bigcup_{i\in\Bbb N}\phi(b_i)\subseteq\NN$.   If $m\in K$
then $U_m\subseteq U_{r(n)}$ for every $n>m$ so
$U_m\cap\phi(b_i)=\emptyset$ for every $i\ge k_{m+1}$ and
$U_m\cap E=\bigcup_{i<k_{m+1}}U_m\cap\phi(b_i)$ is nowhere dense.   As
$\bigcup_{m\in K}U_m$ is a dense open set, $E$ is nowhere dense.   On
the other hand,
$\bigcup_{i\in\Bbb N}b_i=\Bbb N$.   So $\{b:b\in\Cal Z$,
$\phi(b)\subseteq E\}$ is not bounded above in $\Cal Z$, and $\phi$
cannot be a Tukey function.   As $\phi$ is arbitrary,
$\Cal Z\not\prT\CalNwd$.

Because $\Cal E_{\text{Leb}}\prT\CalNwd$ (526I) and $\Cal Z\prT\ell^1$
(526B), it follows that $\Cal Z\not\prT\Cal E_{\text{Leb}}$
and $\ell^1\not\prT\CalNwd$.
}%end of proof of 526K

\leader{526L}{Proposition}\cmmnt{ ({\smc M\'atrai p09})}
$\CalNwd\not\prT\Cal Z$, so $\CalNwd\not\prT\Cal E_{\text{Leb}}$
and $\ell^1\not\prT\Cal Z$.
%M\'atrai 29.6.09

\proof{{\bf (a)(i)} Fix a non-empty
zero-dimensional compact metrizable space $X$ without isolated
points, and write $\CalNwd(X)$ for the ideal of nowhere dense subsets of
$X$;  the bulk of the argument here will be a proof that
$\CalNwd(X)\not\prT\Cal Z$.
Let $\Cal V$ be the family of non-empty open-and-closed subsets of $X$.
For $V\in\Cal V$ write
$\CalNwd(V)=\CalNwd(X)\cap\Cal PV$ for the family of nowhere
dense subsets of $V$.
As in 526A, set
$\nu I=\sup_{n\ge 1}\Bover1n\#(I\cap n)$ for $I\subseteq\Bbb N$.
Take any function $f:\CalNwd(X)\to\Cal Z$.

\medskip

\quad{\bf (ii)} Let $Q$ be the set of pairs $\sigma=(m_{\sigma},I_{\sigma})$
where $I_{\sigma}\subseteq m_{\sigma}\in\Bbb N$;  for $\sigma$, $\tau\in Q$,
say that $\sigma\le\tau$ if either $\sigma=\tau$ or
$2m_{\sigma}\le m_{\tau}$ and $I_{\sigma}=m_{\sigma}\cap I_{\tau}$.   Then
$(Q,\le)$ is a partially ordered set.   For
$\sigma\in Q$ and $\epsilon>0$, let $D(\sigma,\epsilon)$ be the set
of those $E\subseteq X$ for which there is an
$F\in\CalNwd(X)$, including $E$, such that
$f(F)\cap m_{\sigma}\subseteq I_{\sigma}$ and
$\nu(f(F)\setminus I_{\sigma})\le\epsilon$.

\medskip

\quad{\bf (iii)} If $\sigma\in Q$, $\epsilon>0$ and $k\ge 2m_{\sigma}$,
then

\Centerline{$D(\sigma,\epsilon)
\subseteq\bigcup\{D(\tau,\epsilon):\sigma\le\tau\in Q$,
$m_{\tau}=k$, $\nu(I_{\tau}\setminus I_{\sigma})\le\epsilon\}$.}

\noindent\Prf\ If $E\in D(\sigma,\epsilon)$, let $F\in\CalNwd(X)$ be such
that $E\subseteq F$, $f(F)\cap m_{\sigma}\subseteq I_{\sigma}$ and
$\nu(f(F)\setminus I_{\sigma})\le\epsilon$.   Set
$\tau=(k,I_{\sigma}\cup(k\cap f(F))$;  then $\sigma\le\tau$ and $F$
witnesses that $E\in D(\tau,\epsilon)$, while
$\nu(I_{\tau}\setminus I_{\sigma})\le\epsilon$.\ \Qed

\medskip

\quad{\bf (iv)} If $\sigma\in Q$ and $\epsilon$, $\delta>0$, then

\Centerline{$D(\sigma,\epsilon)
  \subseteq\bigcup\{D(\tau,\delta):\sigma\le\tau\in Q$,
$\nu(I_{\tau}\setminus I_{\sigma})\le\epsilon\}$.}

\noindent\Prf\ If $E\in D(\sigma,\epsilon)$, let
$F\in\CalNwd(X)$ be
such that $E\subseteq F$, $F\cap m_{\sigma}\subseteq I_{\sigma}$ and
$\nu(f(F)\setminus I_{\sigma})\le\epsilon$.   As
$f(F)\in\Cal Z$, there is a $k\ge 2m_{\sigma}$ such that
$\nu(f(F)\setminus k)\le\delta$.   Set
$\tau=(k,I_{\sigma}\cup(k\cap f(F))$;  then
$F$ witnesses that $E\in D(\tau,\delta)$, while
$\nu(I_{\tau}\setminus I_{\sigma})\le\epsilon$.\ \Qed

\medskip

\quad{\bf (v)} Suppose that $n\ge 1$ and that
$\langle\sigma_j\rangle_{j\le n}$, $\langle\tau_j\rangle_{j\le n}$ are
finite sequences in $Q$ such
that $m_{\tau_j}\le m_{\sigma_j}$ for $j\le n$ and
$\sigma_j\le\tau_{j+1}$ for $j<n$.    Then
$\nu(\bigcup_{j<n}I_{\tau_{j+1}}\setminus I_{\sigma_j})
\le 3\max_{j<n}\nu(I_{\tau_{j+1}}\setminus I_{\sigma_j})$.
\Prf\ Note first that we certainly have
$m_{\sigma_j}\le m_{\tau_{j+1}}\le m_{\sigma_{j+1}}$ for every $j<n$.
Set $K=\bigcup_{j<n}I_{\tau_{j+1}}\setminus I_{\sigma_j}$ and
$\epsilon=\max_{j<n}\nu(I_{\tau_{j+1}}\setminus I_{\sigma_j})$.
If $m\in\Bbb N$, set $J=\{j:j<n$, $\sigma_j\ne\tau_{j+1}$,
$m_{\sigma_j}\le m\}$,
$J'=\{j:j\in J$, $m_{\tau_{j+1}}\le m\}$.   Then

$$\eqalignno{\#(m\cap K)
&\le\sum_{j\in J}\#(m\cap I_{\tau_{j+1}}\setminus I_{\sigma_j})\cr
\displaycause{because if $j<n$ and $m\le m_{\sigma_j}$, then
$m\cap I_{\tau_{j+1}}\setminus I_{\sigma_j}=\emptyset$}
&\le\epsilon m+\sum_{j\in J'}\#(I_{\tau_{j+1}}\setminus I_{\sigma_j})\cr
\displaycause{because $\#(J\setminus J')\le 1$}
&\le\epsilon(m+\sum_{j\in J'}m_{\tau_{j+1}})
\le\epsilon(m+2m)\cr
\displaycause{because
$2m_{\tau_{j+1}}\le 2m_{\sigma_{j+1}}
\le 2m_{\sigma_{j'}}\le m_{\tau_{j'+1}}\le m$
whenever $j$, $j'$ are successive members of
$J'$}
&=3\epsilon m.\cr}$$

\noindent As $m$ is arbitrary, $\nu K\le 3\epsilon$.\ \Qed

\medskip

{\bf (b)(i)} Suppose that $V\in\Cal V$ and that
$\Cal C_0,\ldots,\Cal C_n\subseteq\CalNwd(X)$ are such that
every nowhere dense subset of $V$ is included in some member of
$\bigcup_{i\le n}\Cal C_i$.   Then there is an $i\le n$ such that every
nowhere dense subset of $V$ is included in some member of $\Cal C_i$.
\Prf\Quer\ Otherwise, for each $i\le n$ we can find a nowhere dense
subset $E_i$ of $V$ not included in any member of $\Cal C_i$;  now
$E=\bigcup_{i\le n}E_i$ is a nowhere dense subset of $V$ not included in
any member of $\bigcup_{i\le n}\Cal C_i$.\  \Bang\Qed

\medskip

\quad{\bf (ii)} Suppose that $V\in\Cal V$ and that $\sequencen{\Cal C_n}$
is a sequence of subsets of $\CalNwd(X)$ such that
every nowhere dense subset of $V$ is included in some member of
$\bigcup_{n\in\Bbb N}\Cal C_n$.   Then for any $E\in\CalNwd(V)$ there are
a $U\in\Cal V$ and an $n\in\Bbb N$ such that $E\subseteq U$ and every
nowhere dense subset of $U$ is included in some member of $\Cal C_n$.
\Prf\ As $V\ne\emptyset$ we can suppose that $E\ne\emptyset$.
Let $\sequencen{U_n}$ be a non-increasing sequence in $\Cal V$ such
that $U_0=V$ and $\bigcap_{n\in\Bbb N}U_n=\overline{E}$.   \Quer\ If, for
every $n\in\Bbb N$, there is an $E_n\in\CalNwd(U_n)$ not included in any
member of $\Cal C_n$, consider $F=\bigcup_{n\in\Bbb N}E_n$;  then
$F\in\CalNwd(V)$ but $F$ is not included in any member of any $\Cal C_n$.\
\Bang\  So some $U_n$ will serve.\ \Qed

\medskip

{\bf (c)} (The key.)  Suppose that $V\in\Cal V$, $\sigma\in Q$
and $\epsilon>0$ are such that $\CalNwd(V)\subseteq D(\sigma,\epsilon)$.
Let $\sequencen{\epsilon_n}$ be a sequence in $\ooint{0,\infty}$.
Then there are an
$n\in\Bbb N$, $U_0,\ldots,U_n\in\Cal V$ and $\tau\in Q$ such
that

\Centerline{$\sigma\le\tau$,
\quad$\nu(I_{\tau}\setminus I_{\sigma})\le 8\epsilon$,}

\Centerline{$V\subseteq\bigcup_{j\le n}U_j$,
\quad$\CalNwd(U_j)\subseteq D(\tau,\epsilon_j)$ for every $j\le n$.}

\noindent\Prf\ It is enough to consider the case in which
$\sum_{n=0}^{\infty}\epsilon_n\le\epsilon$.
Let $\sequencen{x_n}$ run over a dense subset of $V$.
Choose $\sequencen{\sigma_n}$, $\sequencen{k_n}$,
$\langle U_n\rangle_{n\ge 1}$ and $\langle\tau_n\rangle_{n\ge 1}$
inductively, as follows.   Start with
$\sigma_0=\sigma$, $k_0=m_{\sigma_0}$.
Given that $\CalNwd(V)\subseteq D(\sigma_n,\epsilon)$,
we know from (a-iv) that

\Centerline{$\CalNwd(V)\subseteq D(\sigma_n,\epsilon)
\subseteq\bigcup\{D(\tau,\epsilon_{n+1}):\sigma_n\le\tau\in Q$,
   $\nu(I_{\tau}\setminus I_{\sigma_n})\le\epsilon\}$,}

\noindent so by (b-ii) we can find a $U_{n+1}\in\Cal V$ and a
$\tau_{n+1}\ge\sigma_n$ such that $x_n\in U_{n+1}$,
$\nu(I_{\tau_{n+1}}\setminus I_{\sigma_n})\le\epsilon$ and
$\CalNwd(U_{n+1})\subseteq D(\tau_{n+1},\epsilon_{n+1})$.
Next, taking $k_{n+1}=\max(m_{\tau_{n+1}},2m_{\sigma_n})$,
(a-iii) tells us that

\Centerline{$\CalNwd(V)\subseteq D(\sigma_n,\epsilon)
\subseteq\bigcup\{D(\tau,\epsilon):\sigma_n\le\tau\in Q$,
  $m_{\tau}=k_{n+1}$, $\nu(I_{\tau}\setminus I_{\sigma_n})\le\epsilon\}$,}

\noindent so from (b-i) we see that there is a $\sigma_{n+1}\in Q$ such
that $\CalNwd(V)\subseteq D(\sigma_{n+1},\epsilon)$,
$m_{\sigma_{n+1}}=k_{n+1}$, $\sigma_n\le\sigma_{n+1}$ and
$\nu(I_{\sigma_{n+1}}\setminus I_{\sigma_n})\le\epsilon$.   Continue.

At the end of the induction, set $E=V\setminus\bigcup_{n\in\Bbb N}U_{n+1}$.
Because $\{x_n:n\in\Bbb N\}$ is dense in $V$, so is
$\bigcup_{n\in\Bbb N}U_{n+1}$, and $E\in\CalNwd(V)$.   By (a-iv) and
(b-ii) again, there are a $U_0\in\Cal V$ and a $\tau_0\ge\sigma$ such
that $E\subseteq U_0$, $\CalNwd(U_0)\subseteq D(\tau_0,\epsilon_0)$ and
$\nu(I_{\tau_0}\setminus I_{\sigma})\le\epsilon$.   Now
$V\subseteq\bigcup_{n\in\Bbb N}U_n$;  since $V$ is compact, there is an
$n\in\Bbb N$ such that $V\subseteq\bigcup_{j\le n}U_j$.

I have still to define $\tau$.   Set $k=2\max(k_n,m_{\tau_0})$.   For each
$j\le n$, (a-iii) and (b-i), as before, show us that there is an
$\upsilon_j\in Q$ such that $\tau_j\le\upsilon_j$,
$m_{\upsilon_j}=k$,
$\nu(I_{\upsilon_j}\setminus I_{\tau_j})\le\epsilon_j$ and
$\CalNwd(U_j)\subseteq D(\upsilon_j,\epsilon_j)$.
Try setting $\tau=(k,\bigcup_{j\le n}I_{\upsilon_j})$.
Then surely $\CalNwd(U_j)\subseteq D(\tau,\epsilon_j)$ for each $j$.
To estimate $\nu(I_{\tau}\setminus I_{\sigma})$, set
$K=\bigcup_{j<n}I_{\tau_{j+1}}\setminus I_{\sigma_j}$,
$K'=\bigcup_{j<n}I_{\sigma_{j+1}}\setminus I_{\sigma_j}$.   By (a-v),
$\nu K$ and $\nu K'$ are both at most $3\epsilon$.   Now

$$\eqalign{I_{\tau}\setminus I_{\sigma}
&\subseteq\bigcup_{j\le n}(I_{\upsilon_j}\setminus I_{\tau_j})
   \cup(I_{\tau_0}\setminus I_{\sigma})\cr
&\mskip100mu
   \cup\bigcup_{j<n}(I_{\tau_{j+1}}\setminus I_{\sigma_j})
   \cup\bigcup_{j\le n}(I_{\sigma_j}\setminus I_{\sigma})\cr
&=\bigcup_{j\le n}(I_{\upsilon_j}\setminus I_{\tau_j})
   \cup(I_{\tau_0}\setminus I_{\sigma})
   \cup K\cup K',\cr}$$

\noindent and

$$\eqalign{\nu(I_{\tau}\setminus I_{\sigma})
&\le\sum_{j=0}^n\nu(I_{\upsilon_j}\setminus I_{\tau_j})
  +\nu(I_{\tau_0}\setminus I_{\sigma})
  +\nu K+\nu K'\cr
&\le 7\epsilon+\sum_{j=0}^n\epsilon_j
\le 8\epsilon,\cr}$$

\noindent as required.\ \Qed

\medskip

{\bf (d)} Now we can find $T\subseteq S=\bigcup_{n\in\Bbb N}\BbbN^n$,
$\family{t}{T}{\delta_t}$,
$\family{t}{T}{\sigma_t}$, $\family{t}{T}{\tau_t}$ and
$\family{t}{T}{V_t}$ such that

\inset{$T$ is a tree (that is, $t\restr k\in T$ whenever $t\in T$ and
$k\in\Bbb N$),

$\delta_t>0$ for every $t\in T$, $\sum_{t\in T}\delta_t<\infty$,

$\sigma_t\in Q$, $\tau_t\in Q$, $\sigma_t\le\tau_t$
for every $t\in T$,

$\sigma_t=\tau_{t\restr n}$ whenever $n\in\Bbb N$ and
$t\in T\cap\BbbN^{n+1}$,

$\nu(I_{\tau_t}\setminus I_{\sigma_t})\le\delta_t$ for every
$t\in T$,

$V_t\in\Cal V$, $\CalNwd(V_t)\subseteq D(\sigma_t,\delta_t)$ for every
$t\in T$,

$\bigcup\{V_t:t\in T\cap\BbbN^n\}=X$ for every $n\in\Bbb N$.}

\noindent\Prf\ Begin by choosing strictly positive $\delta_t$, for
$t\in S$, such that $\delta_{\emptyset}=1$
and $\sum_{t\in S}\delta_t$ is finite.
Now choose $T_n\subseteq\BbbN^n$ and
$\family{t}{T_n}{\sigma_t}$, $\family{t}{T_n}{V_t}$
inductively, as follows.
Start with $T_0=\{\emptyset\}$,
$\sigma_{\emptyset}=(0,\emptyset)$ and $V_{\emptyset}=X$.
Then

\Centerline{$\CalNwd(V_{\emptyset})=\CalNwd(X)=D((0,\emptyset),1)
=D(\sigma_{\emptyset},\delta_{\emptyset})$,}

\noindent so the process starts.
Given that $T_n$, $\family{t}{T_n}{\sigma_t}$ and
$\family{t}{T_n}{V_t}$
have been defined, then for each $t\in T_n$ use (c) to find $n_t\in\Bbb N$,
$\langle V_{t^{\smallfrown}\fraction{i}}\rangle_{i\le n_t}
\in\Cal V^{n_t+1}$
and $\tau_t\in Q$ such that $\sigma_t\le\tau_t$,
$\nu(I_{\tau_t}\setminus I_{\sigma_t})\le\delta_t$,
$V_t\subseteq\bigcup_{i\le n_t}V_{t^{\smallfrown}\fraction{i}}$
and
$\CalNwd(V_{t^{\smallfrown}\fraction{i}})
\subseteq D({\tau_t},\delta_{t^{\smallfrown}\fraction{i}})$ for every
$i\le n_t$.   Set
$T_{n+1}=\{t^{\smallfrown}\fraction{i}:t\in T_n$, $i\le n_t\}$ and
$\sigma_t=\tau_{t\restr n}$ for every $t\in T_{n+1}$, and continue.
At the end of the construction, set $T=\bigcup_{n\in\Bbb N}T_n$.\ \Qed

\medskip

{\bf (e)} Let $\sequencen{y_n}$ run over a dense subset of $X$.   For
$n\in\Bbb N$, take $t_n\in T\cap\BbbN^n$ such that
$y_n\in V_{t_n}$.   Since
$\CalNwd(V_{t_n})\subseteq D(\sigma_{t_n},\delta_{t_n})$, we can
choose an $F_n\in\CalNwd(X)$, containing $y_n$, such that
$\nu(f(F_n)\setminus I_{\sigma_{t_n}})\le\delta_{t_n}$.
Now $\{f(F_n):n\in\Bbb N\}$ is bounded above in $\Cal Z$.   \Prf\ Set
$K=\bigcup_{t\in T}I_{\sigma_t}$.   As $I_{\sigma_{\emptyset}}=\emptyset$,
\Centerline{$K
=\bigcup_{n\in\Bbb N}\bigcup_{t\in T\cap\BbbN^{n+1}}
  I_{\sigma_t}\setminus I_{\sigma_{t\restr n}}
=\bigcup_{t\in T}I_{\tau_t}\setminus I_{\sigma_t}$;}

\noindent as $\sum_{t\in T}\nu(I_{\tau_t}\setminus I_{\sigma_t})$ is
finite, $K\in\Cal Z$ (526Ac).   Next,

\Centerline{$\bigcup_{n\in\Bbb N}f(F_n)\setminus K
\subseteq\bigcup_{n\in\Bbb N}f(F_n)\setminus I_{\sigma_{t_n}}$;}

\noindent as

\Centerline{$\sum_{n=0}^{\infty}\nu(f(F_n)\setminus I_{\sigma_{t_n}})
\le\sum_{n=0}^{\infty}\delta_{t_n}$}

\noindent is finite, $\bigcup_{n\in\Bbb N}f(F_n)\setminus K\in\Cal Z$,
so $\bigcup_{n\in\Bbb N}f(F_n)$ also belongs to $\Cal Z$, and is an upper
bound for $\{f(F_n):n\in\Bbb N\}$.\ \Qed

\medskip

{\bf (f)}
On the other hand, $\{F_n:n\in\Bbb N\}$ is certainly not bounded above in
$\CalNwd(X)$, since $\bigcup_{n\in\Bbb N}F_n$ includes the dense set
$\{y_n:n\in\Bbb N\}$.   So $f$ cannot be a Tukey function.
Since $f$ is arbitrary, $\CalNwd(X)\not\prT\Cal Z$.

\medskip

{\bf (g)} Since $\CalNwd(X)\equivT\CalNwd$ (526He), it follows that
$\CalNwd\not\prT\Cal Z$.
Since $\NN\prT\Cal E_{\text{Leb}}\penalty-100\prT\Cal Z$ (526I, 526J)
and $\CalNwd\prT\ell^1$ (526Hc),
we see that
$\CalNwd\not\prT\Cal E_{\text{Leb}}$ and
$\ell^1\not\prT\Cal Z$.
}%end of proof of 526L

\cmmnt{\medskip

\noindent{\bf Remark} A somewhat stronger result is in
{\smc Solecki \& Todor\v{c}evi\'c 10}.}

\leader{526M}{}\cmmnt{ Having introduced ideals of sets with
negligible closures, I add a simple result which will be useful later.

\medskip

\noindent}{\bf Proposition} Let $X$ be a second-countable topological
space and $\mu$ a $\sigma$-finite topological measure on $X$.
Let $\Cal E$ be the ideal of subsets of $X$ with negligible closures,
$\Cal N(\mu)$ the null ideal of $\mu$, and $\Cal M$ the ideal of
meager subsets of $\NN$.   Then

\Centerline{$(\Cal E,\subseteq,\Cal N(\mu))\prGT(\Cal M,\not\ni,\NN)$;}

\noindent consequently $\add(\Cal E,\subseteq,\Cal N(\mu))\ge\frakmctbl$.

\proof{{\bf (a)} Suppose first that $\mu$ is a probability measure.
Let $\Cal U$ be a countable base for the topology of $X$, containing
$\emptyset$ and closed under finite unions.   For each $n\in\Bbb N$,
let $\sequence{i}{U_{ni}}$ run over
$\{U:U\in\Cal U$, $\mu U\ge 1-2^{-n}\}$.   For $f\in\NN$, set

\Centerline{$\psi(f)
=\bigcap_{n\in\Bbb N}\bigcup_{i\ge n}X\setminus U_{i,f(i)}
\in\Cal N(\mu)$.}

\noindent For $E\in\Cal E$, set

\Centerline{$\phi(E)=\{f:f\in\NN$, $E\not\subseteq\psi(f)\}$.}

\noindent Then $\phi(E)\in\Cal M$.   \Prf\ Since
$X\setminus\overline{E}$ is a conegligible open set, we can find for
each $i\in\Bbb N$ a $g(i)\in\Bbb N$ such that
$E\cap U_{i,g(i)}=\emptyset$.   Now

\Centerline{$M
=\bigcup_{n\in\Bbb N}\bigcap_{i\ge n}\{f:f\in\NN$, $f(i)\ne g(i)\}$}

\noindent belongs to $\Cal M$.   If $f\in\phi(E)$, there is an
$x\in E\setminus\psi(f)$, so that $x\in\bigcap_{i\ge n}U_{i,f(i)}$ for
some $n$;   now if $i\ge n$, we have
$x\in U_{i,f(i)}\setminus U_{i,g(i)}$, so $f(i)\ne g(i)$;  thus
$f\in M$.   Accordingly $\phi(E)\subseteq M\in\Cal M$.\ \Qed

Now $(\phi,\psi)$ is a Galois-Tukey connection from
$(\Cal E,\subseteq,\Cal N(\mu))$ to $(\Cal M,\not\ni,\NN)$, and
$(\Cal E,\subseteq,\penalty-100\Cal N(\mu))\prGT(\Cal M,\not\ni,\NN)$.

\medskip

{\bf (b)} If $\mu X=0$, then of course
$(\Cal E,\subseteq,\Cal N(\mu))\prGT(\Cal M,\not\ni,\NN)$ (take
$\phi(E)=\emptyset$ for every $E\in\Cal E$, $\psi(f)=X$ for every
$f\in\NN$).   Otherwise, there is a probability measure $\nu$ on $X$
with the same domain and the same null ideal as $\mu$, so (a) tells us
that $(\Cal E,\subseteq,\Cal N(\mu))\prGT(\Cal M,\not\ni,\NN)$.

\medskip

{\bf (c)} Accordingly

\Centerline{$\add(\Cal E,\subseteq,\Cal N(\mu))
\ge\add(\Cal M,\not\ni,\NN)=\cov\Cal M$}

\noindent (512Db).    But, writing $\Cal M(\Bbb R)$ for the ideal of
meager subsets of $\Bbb R$, $\cov\Cal M=\cov\Cal M(\Bbb R)=\frakmctbl$,
by 522Wb and 522Sa.
}%end of proof of 526M

\cmmnt{\medskip

\noindent{\bf Remark} If $X=\Bbb R$ and $\mu$ is Lebesgue measure, then
$\add(\Cal E,\subseteq,\Cal N(\mu))=\frakmctbl$ and
$\cov(\Cal E,\subseteq,\Cal N(\mu))=\non\Cal M$;   see
{\smc Bartoszy\'nski \& Shelah 92} or
{\smc Bartoszy\'nski \& Judah 95}, 2.6.14.
}%end of comment

\exercises{\leader{526X}{Basic exercises (a)}
%\spheader 526Xa
Let $\nu:\Cal P\Bbb N\to[0,1]$ be the submeasure
described in 526A.   Show that
$d^*(I)=\lim_{n\to\infty}\nu(I\setminus n)$ for every
$I\subseteq\Bbb N$.
%526A

\spheader 526Xb
For $I$, $J\subseteq\Bbb N$ say that $I\subseteq^*J$ if $I\setminus J$
is finite.
Show that $(\Cal Z,\subseteq^*,\Cal Z)
\equivGT(\Cal
Z,\discretionary{}{}{}\subseteq\strprime,\discretionary{}{}{}[\Cal
Z]^{\le\omega})$.
%526B

\spheader 526Xc\dvArevised{2023}
Let $\CalNwd$ be the ideal of nowhere dense subsets of
$\NN$ and $\Cal M$ the ideal of meager subsets of $\NN$.   Show that
$\add_{\omega}\CalNwd=\add\Cal M$,
$\non\CalNwd=\omega$, $\cov\CalNwd=\frakmctbl$ and
$\cf\CalNwd\le\cf\Cal M$.
%526H

\spheader 526Xd Let $X$ be a topological space with a countable
$\pi$-base, and $\CalNwd(X)$ the ideal of nowhere dense subsets of $X$.
Show that $\CalNwd(X)\prT\CalNwd$,
where $\CalNwd$ is the ideal of nowhere dense subsets of $\NN$.
%526H

\spheader 526Xe In 526Hf, show that
$\Cal C_{\text{nwd}}$ is a G$_{\delta}$ subset of the family $\Cal C$
of all closed subsets of $X$ with its Fell topology, so is a Polish
space in the subspace topology.
%526H

\spheader 526Xf Let $\Cal C_{\text{Leb}}$ be the family of closed
Lebesgue negligible subsets of $[0,1]$.   Show that $\Cal C_{\text{Leb}}$
with its Fell topology
is a Polish space and a metrizably compactly based directed set.
%526I

\spheader 526Xg Let $\Cal E_{\text{Leb}}$ be the ideal of subsets of
$\Bbb R$ with
negligible closures.   (i) Show that it is Tukey equivalent to the
partially ordered
set $\Cal C_{\text{Leb}}$ of 526Xf.   (ii) Show that it is isomorphic to
$\Cal E_{\text{Leb}}^{\Bbb N}$.   (iii) Show that if we write
$\Cal E_{\sigma}$ for the
$\sigma$-ideal of subsets of $\Bbb R$ generated by
$\Cal E_{\text{Leb}}$, then
$(\Cal E_{\text{Leb}},\subseteq\strprime,
[\Cal E_{\text{Leb}}]^{\le\omega})
\equivGT(\Cal E_{\sigma},\subseteq,\Cal E_{\sigma})$.
(iv) Show that $\add_{\omega}\Cal E_{\text{Leb}}=\add\Cal E_{\sigma}$
and $\cf\Cal E_{\text{Leb}}=\cf\Cal E_{\sigma}$.

\leader{526Y}{Further exercises (a)}
%\spheader 526Ya
Let $X$ be a locally compact separable
metrizable space.   Let $\Cal C_{\text{nwd}}$ be the
family of closed nowhere dense sets in $X$ with its Fell topology.   Show
that $\Cal C_{\text{nwd}}$ is a metrizably compactly based directed set.
%526H

\spheader 526Yb Let $\frak Z$ be the asymptotic density
algebra $\Cal P\Bbb N/\Cal Z$
and define $\bar d^*:\frak Z\to[0,1]$ by setting
$\bar d^*(I^{\ssbullet})=d^*(I)$ for every $I\subseteq\Bbb N$, as in
491I.   Show that if $A\subseteq\frak Z$ is non-empty, downwards-directed
and has infimum $0$, and $\#(A)<\frak p$, then
$\inf_{a\in A}\bar d^*(a)=0$.   (Compare 491Id.)
%n99527

\spheader 526Yc Show that $\wdistr(\frak Z)=\omega_1$.
%n99527, 3C

\spheader 526Yd Show that
$\frak m(\frak Z)\ge\frak m_{\sigma\text{-linked}}$.
%n99527, 3D

\spheader 526Ye Show that $\FN(\Cal P\Bbb N)\le\FN(\frak Z)
\le\max(\FN^*(\Cal P\Bbb N),(\cf\Cal N)^+)$.
%n99527, 3E

\spheader 526Yf Show that $\tau(\frak Z)\ge\frak p$.
%n99527, 3A
}%end of exercises

\endnotes{
\Notesheader{526} The `positive' results of this section are
straightforward enough, except perhaps for 526F.   As elsewhere in this
chapter, I am attempting to
describe a framework which will accommodate the many arguments which
have been found effective in discussing the cardinal functions of these
partially ordered sets.
I note that in this section I use the symbol $\Cal M$ to represent the
ideal of meager subsets of $\NN$, rather than the ideal of meager
subsets of $\Bbb R$, as
elsewhere in the chapter.   If you miss this point, however, none of the
formulae here are dangerous, because the two ideals are Tukey
equivalent, and indeed isomorphic (522Wb).

When we come to `negative' results, we have problems of a new kind.
The special character of Tukey functions is that they need not be of any
particular type.   They are not asked to be order-preserving, and
even if we have partially ordered sets with natural Polish
topologies (as in 526A, 526Xe and 526Xf, for instance),
Tukey functions between them are not
required to be Borel measurable.   This means that
in order to show that there is {\it no} Tukey function between a given
pair of partially ordered sets, we have had to consider arbitrary
functions, or seek to calculate suitable
invariants which we know to be related to the Tukey ordering, like
precaliber triples (516C), and show that they are incompatible with the
existence of a Tukey function.   For a discussion of a class of invariants
giving very sharp distinctions, see {\smc M\'atrai p09}, \S3.

Putting 526B and 526H-526L %526Hc, 526I, 526J, 526K, 526L
together, we find that we have a complete description of the Tukey ordering
on the set $\{\NN,\Cal E_{\text{Leb}},\CalNwd,\Cal Z,\ell^1\}$, given by
the diagram

\def\tmphrule{\hskip0.4em\raise
2.5pt\hbox{\leaders\hrule\hskip1.7em\hfil}\hskip0.4em}
\def\tmpstrut{\vrule height10.5pt depth5.5pt width0pt}

$$\vbox{\offinterlineskip
\halign{\hfil#\hfil&\hfil#\hfil&\hfil#\hfil&\hfil#\hfil
  &\hfil#\hfil&\hfil#\hfil&\hfil#\hfil
  &\hfil#\hfil&\hfil#\hfil&\hfil#\hfil&\hfil#\hfil\cr
&\tmpstrut&\tmpstrut&$\CalNwd$&\tmphrule&$\ell^1$\cr
&\tmpstrut&\tmpstrut&\vrule&&\vrule\cr
&\tmpstrut$\NN$&\tmphrule&$\Cal E_{\text{Leb}}$&\tmphrule
  &$\Cal Z$\cr
}}$$

\noindent if we interpret this in the same way as for Cicho\'n's diagram
(522B).   Moreover, this is exact, in that no two of the five are
Tukey equivalent, and $\Cal Z$ and $\CalNwd$ are Tukey incomparable.
Note that all five of these partially ordered sets
are either themselves metrizably compactly
based directed sets (526A, 513Xj, 513Xl) or are Tukey equivalent to
metrizably compactly based directed sets (526He-526Hf, 526Xf-526Xg).

In 526Yb-526Yf %526Yb 526Yc 526Yd 526Ye 526Yf
I list miscellaneous facts about the asymptotic density algebra.   A
remarkable description of its Dedekind completion is in 556S below.
}%end of notes

\discrpage

\frfilename{mt527.tex}
\versiondate{22.9.21}
\copyrightdate{2002}

\def\chaptername{Cardinal functions of measure theory}
\def\sectionname{Skew products of ideals}

\def\CalMk{\Cal{M}\eurm{k}}
\def\CalSmz{\Cal{S}\eurm{mz}}
\def\rti{right-{\vthsp}translation-{\vthsp}invariant}

\newsection{527}

The methods of this chapter can be applied to a large proportion of the
partially ordered sets which arise in analysis.   In this section I
look at skew products of ideals,
constructed by a method suggested by Fubini's theorem and the
Kuratowski-Ulam theorem (527E).   At the end of the section I introduce
`harmless' algebras (527M-527O).  %527M 527N 527O

\leader{527A}{Notation} If $(X,\Sigma,\mu)$ is a measure space,
$\Cal N(\mu)$ will be the null ideal of $\mu$;
$\Cal N$ will be the null ideal of Lebesgue measure on $\Bbb R$.
If $X$ is a topological space, $\Cal B(X)$ will be the Borel
$\sigma$-algebra of $X$ and $\Cal M(X)$ the $\sigma$-ideal of meager
subsets of $X$;  $\Cal M$ will be the ideal
$\Cal M(\Bbb R)$ of meager subsets of $\Bbb R$.

\vleader{48pt}{527B}{Skew products of ideals} Suppose that
$\Cal I\normalsubgroup\Cal PX$ and $\Cal J\normalsubgroup\Cal PY$ are
ideals of subsets of sets $X$, $Y$ respectively.

\spheader 527Ba I will write $\Cal I\ltimes\Cal J$  for their
{\bf skew product}
$\{W:W\subseteq X\times Y$, $\{x:W[\{x\}]\notin\Cal J\}\in\Cal I\}$.
\cmmnt{(This use of the symbol $\ltimes$ is unconnected with the usage
in \S512 except by the vaguest of analogies.)}
\cmmnt{It is easy to check that}
$\Cal I\ltimes\Cal J\normalsubgroup\Cal P(X\times Y)$.

\cmmnt{Similarly, }$\Cal I\rtimes\Cal J$ will be
$\{W:W\subseteq X\times Y$, $\{y:W^{-1}[\{y\}]\notin\Cal I\}\in\Cal J\}$.

\spheader 527Bb Suppose that $X$ and $Y$ are not empty and that $\Cal I$
and $\Cal J$ are proper ideals.   Then

\Centerline{$\add(\Cal I\ltimes\Cal J)=\min(\add\Cal I,\add\Cal J)$,
\quad$\cf(\Cal I\ltimes\Cal J)\ge\max(\cf\Cal I,\cf\Cal J)$,}

\Centerline{$\non(\Cal I\ltimes\Cal J)=\max(\non\Cal I,\non\Cal J)$,
\quad$\cov(\Cal I\ltimes\Cal J)=\min(\cov\Cal I,\cov\Cal J)$.}

\noindent\prooflet{\Prf\ {\bf (i)} If $\ofamily{\xi}{\kappa}{W_{\xi}}$
is a family in $\Cal I\ltimes\Cal J$ with $\kappa<\min(\add\Cal
I,\add\Cal J)$, set $W=\bigcup_{\xi<\kappa}W_{\xi}$.
For each $\xi<\kappa$, $H_{\xi}=\{x:W_{\xi}[\{x\}]\notin\Cal J\}$
belongs to $\Cal I$;  as $\kappa<\add\Cal I$,
$H=\bigcup_{\xi<\kappa}H_{\xi}\in\Cal I$.   For any $x\notin H$,
$W[\{x\}]=\bigcup_{\xi<\kappa}W_{\xi}[\{x\}]\in\Cal J$ because
$\kappa<\add\Cal J$;  so $W\in\Cal I\ltimes\Cal J$.   As
$\ofamily{\xi}{\kappa}{W_{\xi}}$ is arbitrary,
$\add(\Cal I\ltimes\Cal J)\ge\min(\add\Cal I,\add\Cal J)$.

In the other direction, as $X\notin\Cal I$, $\Cal J=\{F:F\subseteq Y$,
$X\times F\in\Cal I\ltimes\Cal J\}$,
so $F\mapsto X\times F$ is a Tukey function from $\Cal J$ to $\Cal
I\ltimes\Cal J$ and
$\add\Cal J\ge\add(\Cal I\ltimes\Cal J)$, $\cf\Cal J\le\cf(\Cal
I\ltimes\Cal J)$.   Similarly, $E\mapsto E\times Y$ is a Tukey function
from $\Cal I$ to
$\Cal I\ltimes\Cal J$ and $\add\Cal I\ge\add(\Cal I\ltimes\Cal J)$,
$\cf\Cal I\le\cf(\Cal I\ltimes\Cal J)$.

\medskip

\quad{\bf (ii)} Let $A\subseteq X$ and $B\subseteq Y$ be such that
$A\notin\Cal I$, $B\notin\Cal J$, $\#(A)=\non\Cal I$
and $\#(B)=\non\Cal J$.   Then $A\times B\notin\Cal I\ltimes\Cal J$, so
$\non(\Cal I\ltimes\Cal J)\le\#(A\times B)$.   But note that as $\Cal I$
and $\Cal J$ are ideals,
$A$ and $B$ are either singletons or infinite;  so $\#(A\times
B)=\max(\#(A),\#(B))$ and
$\non(\Cal I\ltimes\Cal J)\le\max(\non\Cal I,\non\Cal J)$.

In the other direction, take any
$W\in\Cal P(X\times Y)\setminus(\Cal I\ltimes\Cal J)$.   Set
$E=\{x:W[\{x\}]\notin\Cal J\}$.
Then $\#(E)\ge\non\Cal I$ and $\#(W[\{x\}])\ge\non\Cal J$ for every
$x\in E$, so $\#(W)\ge\max(\non\Cal I,\non\Cal J)$;  as $W$ is
arbitrary,
$\non(\Cal I\ltimes\Cal J)\ge\max(\non\Cal I,\non\Cal J)$.

\medskip

\quad{\bf (iii)} If $\Cal A\subseteq\Cal I$ covers $X$, then
$\{A\times Y:A\in\Cal A\}\subseteq\Cal I\ltimes\Cal J$ covers
$X\times Y$;  so
$\cov(\Cal I\ltimes\Cal J)\le\cov\Cal I$.
Similarly, $\cov(\Cal I\ltimes\Cal J)\le\cov\Cal J$.

Now suppose that $\Cal W\subseteq\Cal I\ltimes\Cal J$ and that
$\#(\Cal W)<\min(\cov\Cal I,\cov\Cal J)$.   For each $W\in\Cal W$ set
$E_W=\{x:W[\{x\}]\notin\Cal J\}$;  then
$E_W\in\Cal I$ for every $W$, so there is an $x\in
X\setminus\bigcup_{W\in\Cal W}E_W$, because $\#(\Cal W)<\cov\Cal I$.
Now $W[\{x\}]\in\Cal J$ for every $W$, so
there is a $y\in Y\setminus\bigcup_{W\in\Cal W}W[\{x\}]$, because
$\#(\Cal W)<\cov\Cal J$.   In this case
$(x,y)\in(X\times Y)\setminus\bigcup\Cal W$.   As $\Cal W$ is arbitrary,
$\cov(\Cal I\ltimes\Cal J)\ge\min(\cov\Cal I,\cov\Cal J)$ and we have
equality.\ \Qed}%end of prooflet

\spheader 527Bc \cmmnt{The idea of the operation $\ltimes$ here is
that we iterate notions of `negligible set' in a way indicated by
Fubini's theorem:  a measurable subset of $\BbbR^2$
is negligible iff almost every vertical section is negligible, that is,
iff it belongs to $\Cal N\ltimes\Cal N$.   However it is immediately
apparent that $\Cal N\ltimes\Cal N$
contains many non-measurable sets, and indeed many sets of full outer
measure (527Xa).   We are therefore led to the following idea.}
If $\Lambda$ is a family of subsets of $X\times Y$,
write $\Cal I\ltimes_{\Lambda}\Cal J\subseteq\Cal I\ltimes\Cal J$
for the ideal generated by $(\Cal I\ltimes\Cal J)\cap\Lambda$.
Note that if $\kappa\le\min(\add\Cal I,\add\Cal J)$ and
$\bigcup\Cal W\in\Lambda$
for every $\Cal W\in[\Lambda]^{<\kappa}$, then
$\add(\Cal I\ltimes_{\Lambda}\Cal J)\ge\kappa$;  in particular,
$\Cal I\ltimes_{\Lambda}\Cal J$ will be a $\sigma$-ideal whenever
$\Cal I$ and $\Cal J$ are $\sigma$-ideals and $\Lambda$ is a
$\sigma$-algebra of subsets of $X\times Y$.
\cmmnt{Typical applications will be with $\Lambda$ a Borel
$\sigma$-algebra or an algebra
of the form $\Sigma\tensorhat\Tau$.   Thus 252F tells us that

\inset{\noindent if $(X,\Sigma,\mu)$ and $(Y,\Tau,\nu)$ are measure
spaces with c.l.d.\ product $(X\times Y,\Lambda,\lambda)$ then
$\Cal N(\mu)\ltimes_{\Lambda}\Cal N(\nu)\subseteq\Cal N(\lambda)$.}

\noindent If $\mu$ and $\nu$ are $\sigma$-finite then we get

\Centerline{$\Cal N(\lambda)
=\Cal N(\mu)\ltimes_{\Sigma\tensorhat\Tau}\Cal N(\nu)$}

\noindent (252C).
If we take $\Cal B=\Cal B(\BbbR^2)$ to be the Borel $\sigma$-algebra of
$\BbbR^2$,
then all four ideals $\Cal N\ltimes_{\Cal B(\BbbR^2)}\Cal N$,
$\Cal M\ltimes_{\Cal B(\BbbR^2)}\Cal M$,
$\Cal M\ltimes_{\Cal B(\BbbR^2)}\Cal N$ and
$\Cal N\ltimes_{\Cal B(\BbbR^2)}\Cal M$ become interesting.
In the next few paragraphs I will sketch some of the ideas needed to
deal with ideals of these kinds.
}%end of comment

\leader{527C}{}\cmmnt{ We are already familiar with
$\Cal N\ltimes_{\Cal B(\BbbR^2)}\Cal N$;  I begin by repeating a result
from \S417 in this language.

\medskip

\noindent}{\bf Theorem} Let $(X,\frak T,\Sigma,\mu)$ and
$(Y,\frak S,\Tau,\nu)$ be $\sigma$-finite effectively locally finite
$\tau$-additive topological measure spaces, both measures being inner regular
with respect to the Borel sets.
Let $\tilde\lambda$ be the $\tau$-additive product measure on
$X\times Y$\cmmnt{ (417C, 417F\formerly{4{}17G})}.   Then
$\Cal N(\mu)\ltimes_{\Cal B(X\times Y)}\Cal N(\nu)
=\Cal N(\tilde\lambda)$.

\proof{{\bf (a)} To begin with, suppose that $\mu$ and $\nu$ are complete.
Then 417C(b-vi) and 417G\formerly{4{}17H} 
tell us that $\tilde\lambda$ is inner regular with
respect to the Borel sets, and that a Borel subset of
$X\times Y$ is $\tilde\lambda$-negligible iff it belongs to
$\Cal N(\mu)\ltimes\Cal N(\nu)$.   On the other hand, because $\mu$ and
$\nu$ are $\sigma$-finite, so is
$\tilde\lambda$ (251K), and every $\tilde\lambda$-negligible set is included in
a $\tilde\lambda$-negligible Borel set.   \Prf\ Suppose that
$W\in\Cal N(\tilde\lambda)$.   Let $\sequencen{W_n}$
be a cover of $X\times Y$ by sets of finite measure.   Because
$\tilde\lambda$ is inner regular with respect to the Borel sets,
we can find $V_n\in\Cal B(X\times Y)$ such
that $V_n\subseteq W_n\setminus W$ and
$\tilde\lambda V_n=\tilde\lambda W_n$ for each $n$.   Now

\Centerline{$W\subseteq(X\times Y)\setminus\bigcup_{n\in\Bbb N}V_n
\in\Cal N(\tilde\lambda)\cap\Cal B(X\times Y)$.  \Qed}

\noindent So $\Cal N(\mu)\ltimes_{\Cal B(X\times Y)}\Cal N(\nu)
=\Cal N(\tilde\lambda)$.

\medskip

{\bf (b)} For the general case, let $\hat\mu$ and $\hat\nu$ be the completions
of $\mu$ and $\nu$.   Just because they extend $\mu$ and $\nu$, they are
$\tau$-additive topological measures;  they are inner regular with respect to
the Borel sets by 412Ha;  and similarly they are inner regular with respect to
subsets of open sets of finite measure, that is, they are effectively locally
finite.   Write
$\tilde\lambda\mskip-3mu\hat{\phantom{x}}$ for the $\tau$-additive product of
$\hat\mu$ and $\hat\nu$;  this is a complete locally determined effectively
locally finite topological measure on $X\times Y$.   By 417C(b-iv), it
is inner regular with respect to
$(\Sigma\tensorhat\Tau)\vee\Cal B(X\times Y)$, and of course

\Centerline{$\tilde\lambda\mskip-3mu\hat{\phantom{x}}(E\times F)
=\hat\mu E\cdot\hat\nu F=\mu E\cdot\nu F$}

\noindent whenever $E\in\Sigma$ and $F\in\Tau$.   By the uniqueness assertion in
417Ca, $\tilde\lambda\mskip-3mu\hat{\phantom{x}}=\tilde\lambda$.   So

$$\eqalignno{\Cal N(\tilde\lambda)&
=\Cal N(\tilde\lambda\mskip-3mu\hat{\phantom{x}})
=\Cal N(\hat\mu)\ltimes_{\Cal B(X\times Y)}\Cal N(\hat\nu)
\Displaycause{by (a) above}
=\Cal N(\mu)\ltimes_{\Cal B(X\times Y)}\Cal N(\nu)\cr}$$

\noindent by 212Eb.


}%end of proof of 527C

\leader{527D}{}\cmmnt{ The case
$\Cal M\ltimes_{\Cal B(\BbbR^2)}\Cal M$ is also well known.

\medskip

\noindent}{\bf Theorem} Let $X$ and $Y$ be topological spaces, with
product $X\times Y$.   Write
$\Cal M^*=\Cal M(X)\ltimes_{\Cal B(X\times Y)}\Cal M(Y)$ and
$\Cal M_1^*=\Cal M(X)\ltimes_{\widehat{\Cal B}(X\times Y)}\Cal M(Y)$,
writing $\widehat{\Cal B}(X\times Y)$ for the Baire-property algebra of
$X\times Y$.

(a) If $\Cal M(X\times Y)\subseteq\Cal M_1^*$, then
$\Cal M^*=\Cal M_1^*=\Cal M(X\times Y)$.

(b) Let $\frak G$ be the category algebra of $Y$\cmmnt{ (514I)}.   If
$\pi(\frak G)<\add\Cal M(X)$ then $\Cal M^*=\Cal M(X\times Y)$.

\proof{{\bf (a)(i)} \Quer\ If $\Cal M_1^*\ne\Cal M(X\times Y)$,
there is a set
$W\in\Cal M_1^*\setminus\Cal M(X\times Y)$;  take
$W_1\in\widehat{\Cal B}(X\times Y)\cap(\Cal M(X)\ltimes\Cal M(Y))$
such that $W_1\supseteq W$.   By 4A3Sa\formerly{4{}A3Ra}, there is an open set
$V\subseteq X\times Y$ such that $W_1\symmdiff V$ is meager and
$V\cap V'$ is empty whenever $V'\subseteq X\times Y$
is open and $V'\cap W_1$ is meager.   As $W_1\notin\Cal M(X\times Y)$,
$V$ cannot be empty;  let $G\subseteq X$, $H\subseteq Y$ be non-empty
open sets such that $G\times H\subseteq V$.
In this case, $G\times H$ cannot be meager, so neither $G$ nor $H$ can
be meager.   (If $F\subseteq X$ is nowhere dense, then $F\times Y$ is
nowhere dense in $X\times Y$;  so
$M\times Y\in\Cal M(X\times Y)$ whenever $M\in\Cal M(X)$;  as
$G\times Y\notin\Cal M(X\times Y)$, $G\notin\Cal M(X)$.)   But now we
see that

\Centerline{$\{x:(G\times H)[\{x\}]\notin\Cal M(Y)\}=G\notin\Cal M(X)$,}

\noindent so that $G\times H\notin\Cal M_1^*$;  but
$(G\times H)\setminus W_1$ is meager, so belongs to $\Cal M_1^*$, and
$W_1$ is also supposed to belong to $\Cal M_1^*$.\ \Bang

\medskip

\quad{\bf (ii)} So $\Cal M_1^*=\Cal M(X\times Y)$.   Of course
$\Cal M^*\subseteq\Cal M_1^*$ just because
$\Cal B(X\times Y)\subseteq\widehat{\Cal B}(X\times Y)$.   In the other
direction, if $W\in\Cal M(X\times Y)$ there is a meager
F$_{\sigma}$ set $W'\supseteq W$.   Now $W'$ is a Borel set in
$\Cal M(X\times Y)=\Cal M_1^*$, so $W'\in\Cal M(X)\ltimes\Cal M(Y)$
witnesses that $W\in\Cal M^*$.   Thus $\Cal M(X\times Y)\subseteq\Cal M^*$
and the three classes are equal.

\medskip

{\bf (b)} By (a), I have only to show that $W\in\Cal M^*$ whenever
$W\subseteq X\times Y$ is meager.   Let
$D\subseteq\frak G\setminus\{0\}$ be an order-dense subset with cardinal
$\pi(\frak G)$.
Let $H$ be the smallest comeager regular open subset of $Y$, so that an
open subset of $Y$ is meager iff it is disjoint from $H$ (4A3Sa again).   For
each $d\in D$ let $V_d\subseteq Y$
be an open set such that $V_d^{\ssbullet}=d$ in $\frak G$;  since
$H^{\ssbullet}=1$, we may suppose that $V_d\subseteq H$.   Observe that
if $F\subseteq Y$ is a non-meager closed set,
then there is a $d\in D$ such that $0\ne d\Bsubseteq F^{\ssbullet}$ in
$\frak G$, in which case $V_d\setminus F$ is meager;  as
$V_d\subseteq H$, $V_d\subseteq F$.

If $W\subseteq X\times Y$ is a nowhere dense closed set, it belongs to
$\Cal M^*$.   \Prf\   Set $E=\{x:W[\{x\}]$ is not meager$\}$.   For each
$d\in D$, the set

\Centerline{$E_d=\{x:V_d\subseteq W[\{x\}]\}=\{x:(x,y)\in W$ for every
$y\in V_d\}$}

\noindent is a closed set in $X$ and $E_d\times V_d\subseteq W$;  so
$\interior E_d\times V_d$ is an open subset of $W$.   As $W$ is nowhere
dense, and $V_d\ne\emptyset$, $\interior E_d$ must be
empty, and $E_d\in\Cal M(X)$.   Next, $E=\bigcup_{d\in D}E_d$ and
$\#(D)=\pi(\frak G)<\add\Cal M(X)$, so $E\in\Cal M(X)$ and
$W\in\Cal M^*$.\ \Qed

Since $\Cal M^*$ is a $\sigma$-ideal, it follows that every meager
subset of $X\times Y$ belongs to $\Cal M^*$, as required.
}%end of proof of 527D

\leader{527E}{Corollary} If $X$ and $Y$ are separable metrizable spaces,
then $\Cal M(X\times Y)=\Cal M(X)\ltimes_{\Cal B(X\times Y)}\Cal M(Y)$.

\proof{ $\pi(\frak C)\le\pi(Y)\le w(Y)\le\omega<\add\Cal M(X)$ (514Ja,
5A4Ba, 4A2P(a-i)).
}%end of proof of 527E

\cmmnt{\medskip

\noindent{\bf Remark} The case $X=Y=\Bbb R$ is the {\bf Kuratowski-Ulam
theorem}.
}%end of comment

\leader{527F}{}\cmmnt{ If we mix measure and category, as in
$\Cal M\ltimes_{\Cal B(\BbbR^2)}\Cal N$ and
$\Cal N\ltimes_{\Cal B(\BbbR^2)}\Cal M$, we encounter some new
phenomena.
To deal with the first we need the following, which is important for
other reasons.

\medskip

\noindent}{\bf Lemma}\cmmnt{ (see {\smc Cicho\'n \& Pawlikowski 86})}
Let $X$ be a set, $\Sigma$ a $\sigma$-algebra of subsets of $X$, and
$\Cal I$ a $\sigma$-ideal of subsets of $X$ generated by
$\Sigma\cap\Cal I$;  suppose that
the quotient algebra $\Sigma/\Sigma\cap\Cal I$ is non-zero, atomless and
has countable $\pi$-weight.
Let $Y$ be a set, $\Tau$ a $\sigma$-algebra of subsets of $Y$, and
$\sequencen{\Cal H_n}$ a sequence of finite covers of $Y$ by members of
$\Tau$.   Set

\Centerline{$\Cal H^*_n=\{\bigcup_{m\ge n}H_m:H_m\in\Cal
H_m\cup\{\emptyset\}$ for every $m\ge n\}$}

\noindent for each $n\in\Bbb N$.   Then there is a sequence
$\sequencen{W_n}$ of subsets of $\NN\times X\times Y$ such that

(i) for every $n\in\Bbb N$, $W_n$ is expressible as the union of a
sequence of sets of the form $I\times E\times F$ where $I\subseteq\NN$
is open-and-closed, $E\in\Sigma$ and $F\in\Tau$;

(ii) whenever $n\in\Bbb N$, $\alpha\in\NN$ and $x\in X$ then
$\{y:(\alpha,x,y)\in W_n\}\in\Cal H_n^*$;

(iii) setting $W=\bigcap_{n\in\Bbb N}W_n$, the set
$\{(\alpha,x):\alpha\in\NN$, $x\in X$, $(\alpha,x,f(x))\notin W\}$
belongs to $[\NN]^{\le\omega}\ltimes\Cal I$ for every
$(\Sigma,\Tau)$-measurable function $f:X\to Y$.

\proof{ If $X\in\Cal I$ or $Y$ is empty, we can take every $W_n$ to be $\emptyset$;
suppose otherwise.

\medskip

{\bf (a)} Set $S=\bigcup_{n\in\Bbb N}\BbbN^n$.
There is a family $\family{\sigma}{S}{U_{\sigma}}$ such that

\inset{every $U_{\sigma}$ belongs to $\Sigma\setminus\Cal I$,

for every $\sigma\in S$,
$\sequence{i}{U_{\sigma^{\smallfrown}\fraction{i}}}$ is
a disjoint sequence of subsets of $U_{\sigma}$ and
$U_{\sigma}\setminus\bigcup_{i\in\Bbb N}
  U_{\sigma^{\smallfrown}\fraction{i}}\in\Cal I$,}

\noindent (see 5A1C for the notation here),

\inset{for every $E\in\Sigma\setminus\Cal I$ there is a $\sigma\in S$
such that $U_{\sigma}\setminus E\in\Cal I$.}

\noindent\Prf\ Let $D$ be a countable order-dense set in
$\frak A=\Sigma/\Sigma\cap\Cal I$.   Then the subalgebra $\frak B$ of
$\frak A$ generated by $D$ is countable and atomless and non-trivial.
Let $\Cal E$ be the subalgebra of $\Cal P(\NN)$ generated by the
sets
$I_{\sigma}=\{\alpha:\sigma\subseteq\alpha\in\NN\}$ for $\sigma\in S$.
This is also an atomless countable Boolean algebra, and must therefore
be isomorphic to $\frak B$ (316M).   Let $\pi:\Cal E\to\frak B$ be an
isomorphism, and set $b_{\sigma}=\pi(I_{\sigma})$ for each
$\sigma\in S$.   Set $U_{\emptyset}=X$ and
for $n\in\Bbb N$, $\sigma\in\BbbN^n$ choose a disjoint sequence
$\sequence{i}{U_{\sigma^{\smallfrown}\fraction{i}}}$ of subsets of
$U_{\sigma}$ such that
$U_{\sigma^{\smallfrown}\fraction{i}}^{\ssbullet}
=b_{\sigma^{\smallfrown}\fraction{i}}$ for
every $i$.   This construction ensures that
$U_{\sigma}\in\Sigma\setminus\Cal I$ for every $\sigma$.
If $E\in\Sigma\setminus\Cal I$, there must be a non-zero $d\in D$ such
that
$d\Bsubseteq E^{\ssbullet}$;  now $\pi^{-1}(d)\in\Cal E\setminus\{\emptyset\}$,
so there is a $\tau\in S$ such that
$I_{\tau}\Bsubseteq\pi^{-1}(d)$, $b_{\tau}\Bsubseteq E^{\ssbullet}$
and $U_{\tau}\setminus E\in\Cal I$.
Finally, if $\sigma\in S$, set
$E=U_{\sigma}\setminus\bigcup_{i\in\Bbb N}U_{\sigma^{\smallfrown}\fraction{i}}$;
then for every $\tau\in S$ either $\tau\subseteq\sigma$ and
$U_{\tau}\setminus E\supseteq U_{\sigma^{\smallfrown}\fraction{0}}
\notin\Cal I$, or
$U_{\tau}\cap U_{\sigma}=\emptyset$ and
$U_{\tau}\setminus E=U_{\tau}\notin\Cal I$,
or there is an $i\in\Bbb N$ such that
$\tau\supseteq\sigma^{\smallfrown}\fraction{i}$ and again
$U_{\tau}\setminus E=U_{\tau}\notin\Cal I$.   This means that
$E$ must belong to $\Cal I$, so that $\family{\sigma}{S}{U_{\sigma}}$
has all the required properties.\ \Qed

\medskip

{\bf (b)} Enumerate $S$ as $\sequence{k}{\tau_k}$.   Let
$\sequencen{H_n}$ be a sequence running over
$\bigcup_{n\in\Bbb N}\Cal H_n$.   For $n\in\Bbb N$, set

\Centerline{$K_n
=\{(\sigma,k):\sigma\in\BbbN^{n+2}$, $k<\#(\tau_{\sigma(n)})$,
$\sigma(n+1)=\#(\tau_k)$, $H_{\tau_{\sigma(n)}(k)}\in\Cal H_n\}$,}

\Centerline{$V_n
=\bigcup_{(\sigma,k)\in K_n}\{(\alpha,x,y):
  \tau_k\subseteq\alpha\in\NN$, $x\in U_{\sigma}$,
  $y\in H_{\tau_{\sigma(n)}(k)}\}$.}

\noindent  If $\alpha\in\NN$ and $x\in X$ the section
$\{y:(\alpha,x,y)\in V_n\}$ is either empty or
$H_{\tau_{\sigma(n)}(k)}$ where $\sigma\in\BbbN^{n+2}$,
$x\in U_{\sigma}$ and $\tau_k=\alpha\restr\sigma(n+1)$;
in either case it belongs to $\Cal H_n\cup\{\emptyset\}$.

So if we now set
$W_n=\bigcup_{m\ge n}V_m$, $W_n$ satisfies (i) and (ii) for every $n$.

\medskip

{\bf (c)} Set $W=\bigcap_{n\in\Bbb N}W_n$.   \Quer\ Suppose, if
possible, that $f:X\to Y$ is a
$(\Sigma,\Tau)$-measurable function such that
$\{(\alpha,x):(\alpha,x,f(x))\notin W\}
  \notin[\NN]^{\le\omega}\ltimes\Cal I$.   Note that

\Centerline{$\{V:V\subseteq\NN\times X\times Y$,
$\{x:(\alpha,x,f(x))\in V\}\in\Sigma$ for every $\alpha\in\NN\}$}

\noindent is a $\sigma$-algebra of subsets of $\NN\times X\times Y$
containing $I\times E\times F$ whenever $I$ is open-and-closed,
$E\in\Sigma$ and $F\in\Tau$, so
contains every $V_n$ and every $W_n$.

\wheader{527F}{0}{0}{0}{48pt}
Set

\Centerline{$A_0=\{\alpha:\alpha\in\NN$,
$\{x:(\alpha,x,f(x))\notin W\}\notin\Cal I\}$,}

\noindent so that $A_0$ is uncountable.   For each $\alpha\in A_0$,

\Centerline{$\bigcup_{n\in\Bbb N}\{x:(\alpha,x,f(x))\notin W_n\}
=\{x:(\alpha,x,f(x))\notin W\}$}

\noindent does not belong to $\Cal I$.   So there is an $n\in\Bbb N$
such that

\Centerline{$A_1
=\{\alpha:\alpha\in A_0$,
  $\{x:(\alpha,x,f(x))\notin W_n\}\notin\Cal I\}$}

\noindent is uncountable.   For each $\alpha\in A_1$, set
$G_{\alpha}=\{x:(\alpha,x,f(x))\notin W_n\}$;  then
$G_{\alpha}\in\Sigma\setminus\Cal I$, so there is a $\sigma\in S$ such
that
$U_{\sigma}\setminus G_{\alpha}\in\Cal I$.   Let $\sigma\in S$ be such
that

\Centerline{$A_2
=\{\alpha:\alpha\in A_1$, $U_{\sigma}\setminus G_{\alpha}\in\Cal I\}$}

\noindent is infinite.   Set $m=\max(n,\#(\sigma))$;  then
$U_{\sigma}\cap\{x:(\alpha,x,f(x))\in V_m\}\in\Cal I$ for every
$\alpha\in A_2$.   Set $M=\#(\Cal H_m)$.

Take $k\in\Bbb N$ such that $\#(\{\alpha\restr k:\alpha\in A_2\})\ge M$.
Let $\langle\alpha_i\rangle_{i<M}$ be a family in $A_2$ such that
$\alpha_i\restr k\ne\alpha_j\restr k$ for distinct $i$, $j<M$;  let
$\langle r_i\rangle_{i<M}$, $\langle l_i\rangle_{i<M}$ be such that
$\alpha_i\restr k=\tau_{r_i}$ for each $i$ and
$\Cal H_m=\{H_{l_i}:i<M\}$.
Let $s\in\Bbb N$ be such that $\tau_s(r_i)$ is defined and equal to
$l_i$ for $i<M$.   Let $\sigma'\in\BbbN^{m+2}$ be such that
$\sigma'\supseteq\sigma$, $\sigma'(m)=s$ and $\sigma'(m+1)=k$.   Then
$U_{\sigma'}\notin\Cal I$ and
$U_{\sigma'}\setminus G_{\alpha}\in\Cal I$ for every $\alpha\in A_2$.

Suppose that $i<M$ and $x\in U_{\sigma'}$.   Then

\Centerline{$\{y:(\alpha_i,x,y)\in V_m\}=H_{\tau_{\sigma'(m)}}(j)
=H_{\tau_s(j)}$}

\noindent where $(\sigma',j)\in K_m$, that is, $j$ is such that
$\tau_j\subseteq\alpha_i$ and $\#(\tau(j))=\sigma'(m+1)=k$.   Thus
$j=r_i$, $\tau_s(j)=l_i$ and
$\{y:(\alpha_i,x,y)\in V_m\}=H_{l_i}$.   But this means that, for any
$x\in U_{\sigma'}$,

\Centerline{$\bigcup_{i<M}\{y:(\alpha_i,x,y)\in V_m\}
=\bigcup_{i<M}H_{l_i}=Y$}

\noindent contains $f(x)$;  that is,
$U_{\sigma'}\subseteq\bigcup_{i<M}\{x:(\alpha_i,x,f(x))\in V_m\}$.
On the other hand,

\Centerline{$U_{\sigma'}\cap\{x:(\alpha_i,x,f(x))\in V_m\}
\subseteq U_{\sigma}\cap\{x:(\alpha_i,x,f(x))\in V_m\}\in\Cal I$}

\noindent for each $i<M$, while $U_{\sigma'}$ itself does not belong to
$\Cal I$.   So this is impossible.\ \Bang

Thus $\sequencen{W_n}$ satisfies (iii).
}%end of proof of 527F

\vleader{48pt}{527G}{Theorem} Let $X$ be a set, $\Sigma$ a $\sigma$-algebra of
subsets of $X$, and $\Cal I$ a $\sigma$-ideal of subsets of $X$ which is
generated by $\Sigma\cap\Cal I$;  suppose that the quotient algebra
$\Sigma/\Sigma\cap\Cal I$ is non-zero, atomless and has countable
$\pi$-weight.
Let $(Y,\Tau,\nu)$ be an atomless perfect semi-finite measure space such that
$\nu Y>0$.    Set
$\Cal K=\Cal I\ltimes_{\Sigma\tensorhat\Tau}\Cal N(\nu)$.   Then
$[\frak c]^{\le\omega}\prT\Cal K$, so $\add\Cal K=\omega_1$ and
$\cf\Cal K\ge\frak c$.

\proof{{\bf (a)} To begin with (down to the end of (d)) suppose that $\nu$ is totally
finite.   Because $\nu$ is atomless, we can for each $n\in\Bbb N$
find a finite cover $\Cal H_n$ of $Y$ by measurable sets with measures
at most $2^{-n}$.   Let $\Tau_0$ be the
$\sigma$-algebra generated by $\bigcup_{n\in\Bbb N}\Cal H_n$, so that
$\Tau_0$ is a $\sigma$-subalgebra of $\Tau$.
Construct $\sequencen{\Cal H^*_n}$,
$\sequencen{W_n}$ and $W$ from $\sequencen{\Cal H_n}$ as in
527F.   Then if
$f:X\to Y$ is $(\Sigma,\Tau_0)$-measurable,
$\{(\alpha,x):(\alpha,x,f(x))\notin W\}\in[\NN]^{\le\omega}\ltimes\Cal I$.
Note that $\nu H\le 2^{-n+1}$ for every $H\in\Cal H_n^*$,
so $\nu\{y:(\alpha,x,y)\in W_n\}\le 2^{-n+1}$ for every $\alpha\in\NN$,
$x\in X$ and $n\in\Bbb N$.   For each $\alpha\in\NN$ set
$K_{\alpha}=\{(x,y):(\alpha,x,y)\in W\}$.  Observe that
$K_{\alpha}\in\Sigma\tensorhat\Tau$ because
$W\in\Cal B(\NN)\tensorhat\Sigma\tensorhat\Tau$, and that

\Centerline{$\nu K_{\alpha}[\{x\}]
\le\inf_{n\in\Bbb N}\nu\{y:(\alpha,x,y)\in W_n\}=0$}

\noindent for every $x\in X$, so $K_{\alpha}\in\Cal K$ for every
$\alpha\in\NN$.

\medskip

{\bf (b)} Set $\hat\Sigma=\{E\symmdiff M:E\in\Sigma$, $M\in\Cal I\}$.
Then $\hat\Sigma$ is a $\sigma$-algebra of subsets of $X$ (cf.\ 212Ca)
and $\Cal I$ is a
$\sigma$-ideal in $\hat\Sigma$;  also the identity embedding of $\Sigma$
in $\hat\Sigma$ induces an isomorphism between $\Sigma/\Sigma\cap\Cal I$
and $\hat\Sigma/\Cal I$
(cf.\ 322Da).   Consequently $\hat\Sigma/\Cal I$ has countable
$\pi$-weight, therefore is ccc, and $\hat\Sigma$ is closed under
Souslin's operation (431G).

\medskip

{\bf (c)} Let $A\subseteq\NN$ be an uncountable set, and
$V\in\Sigma\tensorhat\Tau$ a set disjoint from
$\bigcup_{\alpha\in A}K_{\alpha}$.  (I aim to show that
$(X\times Y)\setminus V\notin\Cal I\ltimes\Cal N(\nu)$.)   There must be
sequences $\sequencen{C_n}$ in $\Sigma$, $\sequencen{F_n}$ in $\Tau$
such that $V$ belongs to the $\sigma$-algebra
generated by $\{C_n\times F_n:n\in\Bbb N\}$;  we can of course arrange
that $\bigcup_{n\in\Bbb N}\Cal H_n\subseteq\{F_n:n\in\Bbb N\}$.   Let
$\Tau_1$ be the
$\sigma$-subalgebra of $\Tau$ generated by $\{F_n:n\in\Bbb N\}$, so that
$\Tau_0\subseteq\Tau_1$ and $V\in\Sigma\tensorhat\Tau_1$.
Let $g:Y\to\Bbb R$ be the Marczewski functional defined by setting
$g=\sum_{n=0}^{\infty}3^{-n}\chi F_n$.
Because $\nu$ is perfect, there is a Borel set $H\subseteq g[Y]$ such
that $g^{-1}[H]$ is conegligible.   Let $h:H\to Y$ be any function such
that $g(h(t))=t$ for every
$t\in H$;  note that $h$ is $(\Cal B(H),\Tau_1)$-measurable, where $\Cal B(H)$
is the Borel $\sigma$-algebra of $H$, just because
$\overline{g[F_n]}\cap\overline{g[Y\setminus F_n]}$
is empty for every $n$.   Set $V_0=\{(x,t):x\in X$, $t\in H$,
$(x,h(t))\in V\}$;  then $V_0\in\Sigma\tensorhat\Cal B(H)$.
It follows that $V_0$ belongs to the class of sets obtainable by
Souslin's operation from sets of the form $E\times F$ where $E\in\Sigma$
and $F\subseteq H$ is relatively closed in $H$.   (Use 421F.)

Set $\tilde E=V_0^{-1}[H]$.   Because
$H$ is analytic and $\hat\Sigma$ is closed under Souslin's operation,
$\tilde E\in\hat\Sigma$ and there is a
$(\hat\Sigma,\Cal B(H))$-measurable function $f_1:\tilde E\to H$ such
that $(x,f_1(x))\in V_0$ for every $x\in\tilde E$ (423N\dvAformerly{4{}23M}).
Now $f_2=hf_1:\tilde E\to Y$ is $(\hat\Sigma,\Tau_1)$-measurable and
$(x,f_2(x))\in V$ for every $x\in \tilde E$.

For every $n\in\Bbb N$, $E_n=f_1^{-1}[F_n]$ belongs to $\hat\Sigma$, so
there is an $E_n'\in\Sigma$ such that $E_n\symmdiff E'_n\in\Cal I$.
Similarly, there is an $\tilde E'\in\Sigma$
such that $\tilde E\symmdiff\tilde E'\in\Cal I$.   Because $\Cal I$ is
generated by $\Sigma\cap\Cal I$, there is an $M_0\in\Sigma\cap\Cal I$ including
$(\tilde E\symmdiff\tilde E')\cup\bigcup_{n\in\Bbb N}(E_n\symmdiff E'_n)$.
Now $\tilde E\setminus M_0=\tilde E'\setminus M_0$
belongs to $\Sigma$.   Set $f_3=f_2\restr\tilde E\setminus M_0$;  then
$f_3^{-1}[F_n]=E'_n\setminus M_0\in\Sigma$ for every $n$, so $f_3$ is
$(\Sigma,\Tau_1)$-measurable.   Take any
$y_0\in Y$, and set $f(x)=f_3(x)$ if $x\in\tilde E\setminus M_0$, $y_0$
for other $x\in X$;  then $f$ is $(\Sigma,\Tau_1)$-measurable, therefore
$(\Sigma,\Tau_0)$-measurable.

The set $\{(\alpha,x):(\alpha,x,f(x))\notin W\}$ belongs to
$[\NN]^{\le\omega}\ltimes\Cal I$, so there must be an $\alpha\in A$ such
that $M_1=\{x:(\alpha,x,f(x))\notin W\}$ belongs
to $\Cal I$.   \Quer\ Suppose, if possible, that
$(X\times Y)\setminus V\in\Cal I\ltimes\Cal N(\nu)$.
Then there must be an $x\in X\setminus(M_0\cup M_1)$
such that $V[\{x\}]$ is conegligible.   In this case,
$V[\{x\}]\cap g^{-1}[H]$ is conegligible, so is not empty, and there
is a $y\in V[\{x\}]\cap g^{-1}[H]$.   Consider
$y'=h(g(y))$;  then $g(y')=g(y)$, so $\{n:y'\in F_n\}=\{n:y\in F_n\}$,
and $\{F:y\in F\iff y'\in F\}$ is a $\sigma$-algebra of subsets of $Y$
containing every $F_n$ and therefore
containing $V[\{x\}]$.   So $y'\in V[\{x\}]$ and $(x,g(y))\in V_0$.
This means that $x\in\tilde E$;  as $x\notin M_0$, $f(x)=f_3(x)=f_2(x)$
and $(x,f(x))\in V$.   On the other
hand, $x\notin M_1$, so $(\alpha,x,f(x))\in W$ and $(x,f(x))\in
K_{\alpha}$;  contradicting the choice of $V$ as a set disjoint from
$K_{\alpha}$.\ \Bang

This shows that $(X\times Y)\setminus V\notin\Cal I\ltimes\Cal N(\nu)$.
As $V$ is arbitrary, $\bigcup_{\alpha\in A}K_{\alpha}\notin\Cal K$.

\medskip

{\bf (d)} This is true for every uncountable $A\subseteq\NN$.   But this
means that $A\mapsto\bigcup_{\alpha\in A}K_{\alpha}$ is a Tukey function
from $[\NN]^{\le\omega}$ to
$\Cal K$, and $[\frak c]^{\le\omega}\cong[\NN]^{\le\omega}\prT\Cal K$.

\medskip

{\bf (e)} Thus the theorem is true if $\nu Y$ is finite.   For the
general case, let $Y_0\in\Tau$ be such that $0<\nu Y_0<\infty$.
Then the subspace measure $\nu_{Y_0}$ is still atomless and perfect
(214Ka, 451Dc), so $[\frak c]^{\le\omega}\prT\Cal K_0$, where
$\Cal K_0=\Cal I\ltimes_{\Sigma\tensorhat(\Tau\cap\Cal PY_0)}
\Cal N(\nu_{Y_0})$.
But $\Cal K_0=\Cal K\cap\Cal P(X\times Y_0)$,
so the identity map from $\Cal K_0$ to
$\Cal K$ is a Tukey function, and

\Centerline{$[\frak c]^{\le\omega}\prT\Cal K_0\prT\Cal K$}

\noindent in this case also.
It follows at once that $\add\Cal K\le\add[\frak c]^{\le\omega}=\omega_1$,
so that $\add\Cal K=\omega_1$ and
$\cf\Cal K\ge\cff[\frak c]^{\le\omega}=\frak c$.
}%end of proof of 527G

\leader{527H}{Corollary} $\Cal M\ltimes_{\Cal B(\BbbR^2)}\Cal N
\equivT[\frak c]^{\le\omega}$.

\proof{ By 527G, $[\frak c]^{\le\omega}\prT\Cal M\ltimes_{\Cal B(\BbbR^2)}\Cal N$.
In the other direction, all we need to observe is
that $\#(\Cal B(\BbbR^2))=\frak c$.
Let $\ofamily{\xi}{\frakc}{W_{\xi}}$ run over
$\Cal B(\BbbR^2)\cap(\Cal M\ltimes\Cal N)$, and for
$V\in\Cal M\ltimes_{\Cal B(\BbbR^2)}\Cal N$
choose $\xi_V<\frak c$ such
that $V\subseteq W_{\xi_V}$;  then
$V\mapsto\{\xi_V\}:\Cal M\ltimes_{\Cal B(\BbbR^2)}\Cal N\to[\frak c]^{\le\omega}$ is a
Tukey function, so $\Cal M\ltimes_{\Cal B(\BbbR^2)}\prT[\frak c]^{\le\omega}$.
}%end of proof of 527H

\leader{527I}{}\cmmnt{ I now turn to the ideal
$\Cal N\ltimes_{\Cal B(\BbbR^2)}\Cal M$.

\medskip

\noindent}{\bf Lemma} Let $X$ be a set, $\Sigma$ a $\sigma$-algebra of
subsets of $X$, and $Y$ a topological space with a countable $\pi$-base
$\Cal H$.   Let $\Cal W$ be the family of subsets of $X\times Y$ of the
form $\bigcup_{H\in\Cal H}E_H\times H$, where $E_H\in\Sigma$ for every
$H\in\Cal H$, and
$\Cal D_0$ the family of sets $D\subseteq X\times Y$ such that
$(X\times Y)\setminus D\in\Cal W$ and $D[\{x\}]$ is nowhere dense for
every $x\in X$;  let $\Cal L_0$ be the $\sigma$-ideal of subsets of
$X\times Y$ generated by $\Cal D_0$.   Then
$\Sigma\tensorhat\Cal B(Y)
\subseteq\{W\symmdiff L:W\in\Cal W$, $L\in\Cal L_0\}$.

\proof{ Write $\Cal V$ for $\{W\symmdiff L:W\in\Cal W$, $L\in\Cal L_0\}$.
Then $\Cal W$ and $\Cal V$ are closed under countable unions.
Next, $(X\times Y)\setminus W\in\Cal V$
for every $W\in\Cal W$.   \Prf\ Express $W$ as
$\bigcup_{H\in\Cal H}E_H\times H$ where
$E_H\in\Sigma$ for every $H\in\Cal H$.   For $H\in\Cal H$, set

\Centerline{$F_H=X\setminus\bigcup\{E_{H'}:H'\in\Cal H$, $H'\cap H\ne\emptyset\}
\in\Sigma$,}

\noindent and set $W'=\bigcup_{H\in\Cal H}F_H\times H$.
Then $W'$ and $W\cup W'$ belong to $\Cal W$.   Set
$D=(X\times Y)\setminus(W\cup W')$.   If $x\in X$ and $G\subseteq Y$ is a non-empty
open set, let $H\subseteq G$ be a non-empty member of $\Cal H$.
Then either $x\in F_H$ and $H$ is a non-empty open subset of
$G\setminus D[\{x\}]$, or there is an $H'\in\Cal H$ such
that $H\cap H'\ne\emptyset$ and $x\in E_{H'}$, in which case
$H\cap H'$ is a non-empty open subset of $G\setminus D[\{x\}]$.   As $G$ is
arbitrary, $D[\{x\}]$ is nowhere dense;  as $x$ is arbitrary, $D\in\Cal D_0$.
But now observe that $(X\times Y)\setminus W=W'\symmdiff D$ belongs to $\Cal V$.\ \Qed

It follows that the complement of any member of $\Cal V$ belongs to $\Cal V$, so
$\Cal V$ is a $\sigma$-algebra.   Now $E\times G\in\Cal V$ for every $E\in\Sigma$ and
open $G\subseteq Y$.   \Prf\ For $H\in\Cal H$, set $E_H=E$ if $H\subseteq G$,
$\emptyset$ otherwise;  set $W=\bigcup_{H\in\Cal H}E_H\times H\in\Cal W$.   Then
$W\subseteq E\times G$.   But, defining $W'$ from $W$ as just above, we see that
$W'$ is disjoint from $E\times G$.   So

\Centerline{$(E\times G)\symmdiff W\subseteq(X\times Y)\setminus(W\cup W')
\in\Cal D_0$}

\noindent and $E\times G\in\Cal V$.\ \Qed

Accordingly $\Cal V$ includes the $\sigma$-algebra generated by
$\{E\times G:E\in\Sigma$, $G\subseteq Y$ is open$\}$, which is
$\Sigma\tensorhat\Cal B(Y)$.
}%end of proof of 527I

\leader{527J}{Theorem}\cmmnt{ (see {\smc Fremlin 91})} Let $X$ be a
topological space and $\mu$ a $\sigma$-finite quasi-Radon measure on $X$
with countable Maharam
type;  let $Y$ be a topological space of countable $\pi$-weight.   Then
$\Cal N(\mu)\ltimes_{\Cal B(X\times Y)}\Cal M(Y)\prT\Cal N$.

%will "countable network" do for Y ?

\proof{ Write $\Cal L$ for $\Cal N(\mu)\ltimes_{\Cal B(X\times Y)}\Cal M(Y)$,
and fix a countable $\pi$-base $\Cal H$, not containing $\emptyset$,
for the topology of $Y$.

\medskip

{\bf (a)} We need to know that for every Borel set $V\subseteq X\times Y$
there are sets $V'$, $V''\in\Cal B(X)\tensorhat\Cal B(Y)$
such that $V'\subseteq V\subseteq V''$ and
$V''\setminus V'\in\Cal L$.   \Prf\ Let $\Cal V^*$ be the family of
all subsets of $X\times Y$ with this property.   Because
$\Cal B(X)\tensorhat\Cal B(Y)$ is a $\sigma$-algebra and $\Cal L$ is a $\sigma$-ideal
of sets, $\Cal V^*$ is a $\sigma$-algebra.   If $W\subseteq X\times Y$
is open, set

\Centerline{$U_H=\bigcup\{G:G\subseteq X$ is open, $G\times H\subseteq W\}$,
\quad$U'_H=\{x:H\cap W[\{x\}]\ne\emptyset\}$}

\noindent for $H\in\Cal H$, so all the $U_H$ and $U'_H$ are
open ($U'_H$ is just the projection of the open set
$W\cap(X\times H)$).
Set $V_1=\bigcup_{H\in\Cal H}U_H\times H$ and
$V_2=\bigcup_{H\in\Cal H}((X\setminus U'_H)\times H)$.   Then $V_1$
and $V_2$ both belong to $\Cal B(X)\tensorhat\Cal B(Y)$,
$V_1\subseteq W$ and $W\cap V_2=\emptyset$.

Let $x\in X$.   \Quer\ If the open set $V_1[\{x\}]\cup V_2[\{x\}]$ is
not dense, there is an $H\in\Cal H$ disjoint from
both $V_1[\{x\}]$ and $V_2[\{x\}]$.   In this case $x$ must belong to $U'_H$,
and there is a point $y\in H\cap W[\{x\}]$.   $(x,y)$ belongs to the
open set $(X\times H)\cap W$, so there
are open sets $G\subseteq X$, $\tilde H\subseteq Y$ such that
$(x,y)\in G\times\tilde H\subseteq(X\times H)\cap W$.   Now there is an
$H'\in\Cal H$
such that $H'\subseteq\tilde  H$, in which
case $x\in G\subseteq U_{H'}$.   But this will mean that
$H'\subseteq V_1[\{x\}]$ and
$H'$ is a non-empty subset of $H\cap V_1[\{x\}]$, which is
impossible.\ \Bang

Thus $V_1[\{x\}]\cup V_2[\{x\}]$ is dense for every $x$, and if we set
$V_3=(X\times Y)\setminus V_2$ we shall have
$V_3\setminus V_1\in\Cal L$, while both $V_1$ and $V_3$ belong to
$\Cal B(X)\tensorhat\Cal B(Y)$, and $V_1\subseteq W\subseteq V_3$.   So
$W\in\Cal V^*$.   This is true for every open set $W\subseteq X\times Y$,
so the $\sigma$-algebra $\Cal V^*$ must contain every Borel set, as required.\ \Qed

It follows that every member of $\Cal L$ is included in a member of
$\Cal L\cap(\Cal B(X)\tensorhat\Cal B(Y))$.   \Prf\ If $V\in\Cal L$
there is a Borel set $V'\supseteq V$
which belongs to $\Cal L$, and now there is a set
$V''\in\Cal B(X)\tensorhat\Cal B(Y)$ such that $V''\supseteq V'$ and
$V''\setminus V'\in\Cal L$, in which case $V''\supseteq V$ also must
belong to $\Cal L$.\ \Qed

Thus $\Cal L=\Cal N(\mu)\ltimes_{\Cal B(X)\tensorhat\Cal B(Y)}\Cal M(Y)$.

\medskip

{\bf (b)} To begin with let us suppose
that $X$ is compact and metrizable, $\mu$ is totally finite and $Y$ is a
Baire space.

\medskip

\quad{\bf (i)}
Taking $\Sigma=\Cal B(X)$, define $\Cal W$, $\Cal D_0$ and $\Cal L_0$ as in 527I.
Now let $\Cal D$ be the family of closed subsets belonging to $\Cal D_0$, and
$\Cal L_1$ the $\sigma$-ideal of subsets of $X\times Y$ generated by
$\{E\times Y:E\in\Cal N(\mu)\}\cup\Cal D$.

\medskip

\quad {\bf (ii)} $\Cal D_0\subseteq\Cal L_1$.
\Prf\ If $D\in\Cal D_0$, express $(X\times Y)\setminus D$ as
$\bigcup_{H\in\Cal H}E_H\times H$ where $E_H\in\Cal B(X)$
for every $H\in\Cal H$.   Because $\mu$ is totally finite, $\mu$ is outer
regular with respect to the open sets (412Wb).
So for each $n\in\Bbb N$ we can find a family
$\family{H}{\Cal H}{G_{nH}}$ of open sets in $X$ such that $E_H\subseteq G_{nH}$ for
every $H$ and $\sum_{H\in\Cal H}\mu(G_{nH}\setminus E_H)\le 2^{-n}$.   Set
$D_n=(X\times Y)\setminus\bigcup_{H\in\Cal H}(G_{nH}\times H)$.   Then $D_n$ is closed
and $D_n\subseteq D\in\Cal D_0$ so $D_n\in\Cal D$.   Set
$E=\bigcap_{n\in\Bbb N}\bigcup_{H\in\Cal H}(G_{nH}\setminus E_H)$;   then
$E\in\Cal N(\mu)$ and

\Centerline{$D\subseteq (E\times Y)\cup\bigcup_{n\in\Bbb N}D_n\in\Cal L_1$. \Qed}

\medskip

\quad {\bf (iii)} Of course every member of $\Cal D$ belongs to $\Cal L$, so
$\Cal L_1\subseteq\Cal L$.   But in fact $\Cal L=\Cal L_1$.   \Prf\ If $V\in\Cal L$,
there is a $V'\in(\Cal N(\mu)\ltimes\Cal M(Y))\cap(\Cal B(X)\tensorhat\Cal B(Y))$
such that $V\subseteq V'$, by (a).   By 527I, we can
express $V'$ as $W\symmdiff L$ where $W\in\Cal W$ and $L\in\Cal L_0$.   By (ii),
$\Cal L_0\subseteq\Cal L_1$, so $W\in\Cal L$.   There is therefore a negligible set
$E\subseteq X$ such that $W[\{x\}]$ is meager for every $x\in X\setminus E$.
But $W[\{x\}]$ is always open, and $Y$ is a Baire
space, so $W\subseteq E\times Y\in\Cal L_1$.   Accordingly $V'$ and $V$ belong to
$\Cal L_1$.   As $V$ is arbitrary, $\Cal L\subseteq\Cal L_1$.\ \Qed

\medskip

\quad{\bf (iv)} Let $\Cal G$ be a countable base for the topology of $X$ containing
$X$.   Let
$\Cal U_0$ be the family of those sets $U\subseteq X\times Y$ such that $U$
is expressible as a finite union of sets of the form $G\times H$ where $G\in\Cal G$
and $H\in\Cal H$, and $\Cal U$ the set of those $U\in\Cal U_0$ such that
$\pi_1[U]=X$, where $\pi_1$ is the projection from $X\times Y$ onto $X$.
Consider

\Centerline{$\Cal D'=\{D:D\subseteq X\times Y$, for every $U_0\in\Cal U$
there is a $U\in\Cal U$ such that $U\subseteq U_0\setminus D\}$.}

$\Cal D\subseteq\Cal D'$.   \Prf\
Suppose that $D\in\Cal D$ and $U_0\in\Cal U$, and consider
$\Cal U_1=\{U:U\in\Cal U_0$, $U\subseteq U_0\setminus D\}$.   For every
$x\in X$ the section $U_0[\{x\}]$ is open and not empty and the section
$D[\{x\}]$ is
nowhere dense, so there is a $y$ such that $(x,y)\in U_0\setminus D$;
now there are $G\in\Cal G$, containing $x$, and
an open $H$ containing $y$ such that $G\times H\subseteq U_0\setminus D$.
Let $H'\in\Cal H$ be such that $\emptyset\ne H'\subseteq H$.
Then $U=G\times H'\in\Cal U_1$ and $x\in\pi_1[U]$.   As $x$ is arbitrary,
$\{\pi_1[U]:U\in\Cal U_1\}$ is an open cover of $X$;  as $X$ is compact and $\Cal U_1$
is upwards-directed, there is a $U\in\Cal U_1$ such that $\pi_1[U]=X$;  in which case
$U\in\Cal U$ and $U\subseteq U_0\setminus D$.   As $U$ is arbitrary, $D\in\Cal D'$;
as $D$ is arbitrary, $\Cal D\subseteq\Cal D'$.\ \Qed

$\Cal D$ is cofinal with $\Cal D'$.   \Prf\ Let $D\in\Cal D'$.
For each $H\in\Cal H\setminus\{\emptyset\}$, $X\times H\in\Cal U$, so
there is a $U_H\in\Cal U$ such that $U_H\subseteq(X\times H)\setminus D$;
try
$D_1=(X\setminus Y)\setminus\bigcup_{H\in\Cal H\setminus\{\emptyset\}}U_H$.
$D_1$ is closed.   Since $\Cal U\subseteq\Cal U_0\subseteq\Cal W$,
$(X\times Y)\setminus D_1\in\Cal W$.   If $x\in X$, then $D_1[\{x\}]$ is a
closed set not including any member of the $\pi$-base $\Cal H$, so is
nowhere dense in $Y$;  thus $D_1\in\Cal D_0$ and (being closed) belongs to
$\Cal D$.   Of course $D\subseteq D_1$.   As
$D$ is arbitrary, $\Cal D$ is cofinal with $\Cal D'$.\ \Qed

\medskip

\quad{\bf (v)} Because $\Cal U$ is countable, 526Hd tells us that
$\Cal D'\prT\CalNwd$, where $\CalNwd$ is the ideal of nowhere dense
subsets of $\NN$;  while of course $\Cal D\equivT\Cal D'$ (513E(d-ii)).
Let $\phi:\Cal L\to\Cal N(\mu)\times\Cal D^{\Bbb N}$ be such that if
$\phi(V)=(E,\sequencen{D_n})$ then
$V\subseteq(E\times Y)\cup\bigcup_{n\in\Bbb N}D_n$;
such a function exists by (iii), and is evidently a Tukey function.

Note that the measure algebra of $\mu$, being a totally finite measure
algebra with countable Maharam type, can be regularly embedded in the
measure algebra of Lebesgue measure on either $[0,1]$ or on $\Bbb R$.
As $\mu$ is a Radon measure (416G), $\Cal N(\mu)\prT\Cal N$ (524K) and

\Centerline{$\Cal L\prT\Cal N(\mu)\times\Cal D^{\Bbb N}
\prT\Cal N\times\CalNwd^{\Bbb N}\cong\Cal N\times\CalNwd$}

\noindent (513Eg, 526Ha).   Accordingly

$$\eqalignno{(\Cal L,\subseteq,\Cal L)
&\equivGT(\Cal L,\subseteq\strprime,[\Cal L]^{\le\omega})\cr
\displaycause{513Id}
&\prGT(\Cal N\times\CalNwd,\le\strprime,
  [\Cal N\times\CalNwd]^{\le\omega})\cr
\displaycause{512Gb}
&\equivGT(\Cal N,\subseteq\strprime,[\Cal N]^{\le\omega})
  \times(\CalNwd,\subseteq\strprime,[\CalNwd]^{\le\omega})\cr
\displaycause{512Hd}
&\equivGT(\Cal N,\subseteq\Cal N)
  \times(\Cal M,\subseteq,\Cal M)\cr
\displaycause{513Id, 526Hb, 512Hb}
&\prGT(\Cal N,\subseteq,\Cal N)\times(\Cal N,\subseteq,\Cal N)\cr
\displaycause{522P}
&\equivGT(\Cal N,\subseteq,\Cal N)\cr}$$

\noindent (513Eh), and $\Cal L\prT\Cal N$.

\medskip

{\bf (c)} This proves the theorem when $X$ is compact and metrizable,
$\mu$ is totally finite and $Y$ is a Baire space.   Now suppose that $Y$
is still
a Baire space, while $(X,\mu)$ is any totally finite quasi-Radon measure
space with countable Maharam type.

\medskip

\quad{\bf (i)} There is a compact metrizable Radon measure space
$(Z,\lambda)$ such that $\lambda$ and $\mu$ have isomorphic measure
algebras.   \Prf\ Because the measure
algebra $(\frak A,\bar\mu)$ of $\mu$ is totally finite, it is isomorphic
to the simple product of a countable family
$\familyiI{(\frak A_i,\bar\mu_i)}$ of homogeneous
totally finite measure algebras (332B).   Because $\mu$ has countable
Maharam type, every $\frak A_i$ is either $\{0\}$, $\{0,1\}$ or
isomorphic to the measure algebra of Lebesgue
measure on an interval;  in any case it is isomorphic to the measure
algebra of a compact Radon measure space $(Z_i,\lambda_i)$.   Take
$(Z',\lambda')$ to be the direct sum
of the measure spaces $\familyiI{(Z_i,\lambda_i)}$;  then the measure
algebra of $(Z',\lambda')$ is isomorphic to $\frak A$.   If we give $Z'$
its disjoint union topology, it is
a locally compact $\sigma$-compact metrizable space, and its one-point
compactification $Z$ is second-countable, therefore metrizable;  taking
$\lambda$ to be the trivial extension
of $\lambda'$, $(Z,\lambda)$ is a compact metrizable Radon measure space
with measure algebra $(\frak B,\bar\lambda)\cong(\frak A,\bar\mu)$.\
\Qed

\medskip

\quad{\bf (ii)} Let $f:X\to Z$ be an \imp\ function inducing an
isomorphism $\pi:\frak B\to\frak A$ of the measure algebras (416Wb).
By 418J, $f$ is almost
continuous, so there is a Borel measurable function which is equal
almost everywhere to $f$ (418V\Latereditions);
this function will still represent $\pi$, so
we may suppose that $f$ itself is Borel
measurable.   Now if $V\in\Cal B(X)\tensorhat\Cal B(Y)$, there is a
$V'\in\Cal B(Z)\tensorhat\Cal B(Y)$ such that
$\{x:V[\{x\}]\ne V'[\{f(x)\}]\}\in\Cal N(\lambda)$.
\Prf\ Let $\tilde{\Cal V}$ be the family of subsets $V$ of $X\times Y$
for which there is a $V'\in\Cal B(Z)\tensorhat\Cal B(Y)$ such that
$\{x:V[\{x\}]\ne V'[\{f(x)\}]\}\in\Cal N(\lambda)$.
Then $\tilde{\Cal V}$ is a $\sigma$-algebra.   If $E\in\Cal B(X)$ and
$H\in\Cal B(Y)$,
then there must be an $F\in\Cal B(Z)$ such that
$F^{\ssbullet}=\pi E^{\ssbullet}$ in $\frak B$, so that
$E\symmdiff f^{-1}[F]\in\Cal N(\mu)$;  now $F\times H$ witnesses that
$E\times H$
belongs to $\tilde{\Cal V}$.   Accordingly $\tilde{\Cal V}$ must include
$\Cal B(X)\tensorhat\Cal B(Y)$.\ \Qed

\medskip

\quad{\bf (iii)} We know that $\Cal N(\mu)\prT\Cal N(\lambda)$
(524Sa), so there is a Tukey function
$\theta:\Cal N(\mu)\to\Cal N(\lambda)$.   Set
$\Cal L'=\Cal N(\lambda)\ltimes_{\Cal B(Z)\tensorhat\Cal B(Y)}\Cal M(Y)$.
Define a function $\phi:\Cal L\to\Cal L'$ as
follows.   First, for $V\in\Cal L$, choose
$\phi_0(V)\in\Cal L\cap(\Cal B(X)\tensorhat\Cal B(Y))$ including $V$
((a) above).   Next, by (ii) here, we can
choose $\phi_1(V)\in\Cal B(Z)\tensorhat\Cal B(Y)$
such that $N_V=\{x:\phi_0(V)[\{x\}]\ne\phi_1(V)[\{f(x)\}]\}$ belongs to
$\Cal N(\mu)$.   Set
$F=\{z:z\in Z$, $\phi_1(V)[\{z\}]$ is not meager$\}$;
then $F$ is a Borel set, by 4A3Ta\formerly{4{}A3Sa}, and
$f^{-1}[F]\subseteq N_V\cup\{x:\phi_0(V)[\{x\}]\notin\Cal M(Y)\}\in\Cal N(\mu)$;
so $F\in\Cal N(\lambda)$ and $\phi_1(V)\in\Cal L'$.
Finally, set
$\phi(V)=(\theta(N_V)\times Y)\cup\phi_1(V)\in\Cal L'$.

$\phi$ is a Tukey function from $\Cal L$ to $\Cal L'$.   \Prf\ Take
$W\in\Cal L'$ and consider $\Cal E=\{V:V\in\Cal L$, $\phi(V)\subseteq W\}$.
If $Y=\emptyset$ then of course $\Cal E$ is bounded above in $\Cal L$.
Otherwise, $N^*=\{z:W[\{z\}]=Y\}$ must be $\lambda$-negligible, while
$\theta(N_V)\subseteq N^*$ for every $V\in\Cal E$;
because $\theta$ is a
Tukey function, $\tilde N=\bigcup\{N_V:V\in\Cal E\}$ is $\mu$-negligible.
Take $W_1\in\Cal L'\cap(\Cal B(Z)\tensorhat\Cal B(Y))$ including $W$,
and set $\tilde W=\{(x,y):(f(x),y)\in W_1\}$;
then $\tilde W\in\Cal B(X)\tensorhat\Cal B(Y)$ because $f$ is Borel
measurable.   As

\Centerline{$\{x:\tilde W[\{x\}]\notin\Cal M(Y)\}
=f^{-1}\{z:W_1[\{z\}]\notin\Cal M(Y)\}$}

\noindent is negligible, $\tilde W\in\Cal L$.   So
$V_0=(\tilde N\times Y)\cup\tilde W$ belongs to $\Cal L$.
Now take any $V\in\Cal E$.   If
$x\in X\setminus\tilde N$, then $x\notin N_V$, so

$$\eqalign{V[\{x\}]&
\subseteq\phi_0(V)[\{x\}]
=\phi_1(V)[\{f(x)\}]
\subseteq W[\{f(x)\}]\cr&
\subseteq W_1[\{f(x)\}]
=\tilde W[\{x\}]
=V_0[\{x\}].\cr}$$

\noindent This shows that $V\subseteq V_0$;  as $V$ is arbitrary, $V_0$
is an upper bound for $\Cal E$ in $\Cal L$;  as $W$ is arbitrary, $\phi$
is a Tukey function.\ \Qed

\medskip

\quad{\bf (iv)} By (b), we know that $\Cal L'\prT\Cal N$, so (iii)
tells us that $\Cal L\prT\Cal N$, and the theorem is true in this case
also.

\medskip

{\bf (d)} We are nearly home.   If $Y$ is a Baire space and $(X,\mu)$ is
a $\sigma$-finite quasi-Radon measure space with countable Maharam type,
which is not totally finite, then
there is a measurable function $f:X\to\ooint{0,\infty}$ such that
$\int fd\mu=1$ (215B(ix)).   Let $\nu$ be the indefinite-integral
measure defined by $f$.   Then $\nu$ has the same
negligible sets as $\mu$ (234Lc), and is a quasi-Radon measure (415Ob), so

\Centerline{$\Cal L=\Cal N(\nu)\ltimes_{\Cal B(X\times Y)}\Cal M(Y)
\prT\Cal N$,}

\noindent by (c).

\medskip

{\bf (e)} Finally, suppose that $Y$ is not a Baire space.   In this
case, let $H^*$ be the smallest
comeager regular open subset of $Y$ (4A3Sa once more), and set
$\Cal L^*=\Cal N(\mu)\ltimes_{\Cal B(X\times H^*)}\Cal M(H^*)$.
Then $\Cal L\prT\Cal L^*$.
\Prf\ For every $V\in\Cal L$, let
$V'$ be such that
$V\subseteq V'\in\Cal B(X\times Y)\cap(\Cal N(\mu)\ltimes\Cal M(Y))$, and set
$\phi(V)=V'\cap(X\times H^*)$.   Then
$\phi(V)$ is a Borel subset of $X\times H^*$,
and $\phi(V)[\{x\}]=V'[\{x\}]\cap H^*$ is meager in $H^*$ whenever
$V'[\{x\}]$ is meager in $Y$, so $\phi(V)\in\Cal L^*$.   To see that
$\phi:\Cal L:\to\Cal L^*$ is
a Tukey function, take any $W\in\Cal L^*$.   There is a Borel set
$W'\in\Cal L^*$ including $W$, and now $V'=W'\cup(X\times(Y\setminus H^*))$
is a Borel subset of $X\times Y$;  since $V'[\{x\}]$ is meager in $Y$
whenever $W'[\{x\}]$ is meager in $H^*$, $V'\in\Cal L$.   Of course $V'$
is an upper bound
of $\{V:V\in\Cal L$, $\phi(V)\subseteq W\}$;  as $W$ is arbitrary,
$\phi$ is a Tukey function and $\Cal L\prT\Cal L^*$.\ \Qed

By (d), $\Cal L\prT\Cal N$ in this case also, and the proof is complete.
}%end of proof of 527J

\leader{527K}{Corollary}
$\Cal N\ltimes_{\Cal B(\BbbR^2)}\Cal M\equivT\Cal N$.

\proof{ By 527J, $\Cal N\ltimes_{\Cal B(\BbbR^2)}\Cal M\prT\Cal N$.   On
the other hand, $E\mapsto E\times\Bbb R$ is a Tukey function from
$\Cal N$ to $\Cal N\ltimes_{\Cal B(\BbbR^2)}\Cal M$, so
$\Cal N\prT\Cal N\ltimes_{\Cal B(\BbbR^2)}\Cal M$.
}%end of proof of 527K

\leader{527L}{}\cmmnt{ There are some interesting questions concerning
the saturation of skew products.
Here and in 527O I give results which will be useful later.

\medskip

\noindent}{\bf Theorem} Let $X$ be a set, $\Sigma$ a $\sigma$-ideal of
subsets of $X$, and $\Cal I\normalsubgroup\Cal PX$ a $\sigma$-ideal;
suppose that $\Sigma/\Sigma\cap\Cal I$ is
ccc.   Let $(Y,\Tau,\nu)$ be a $\sigma$-finite measure space.   Then
$(\Sigma\tensorhat\Tau)
/((\Sigma\tensorhat\Tau)\cap(\Cal I\ltimes\Cal N(\nu)))$ is ccc.

\proof{{\bf (a)} The case $\nu Y=0$ is trivial, as then
$\Cal I\ltimes\Cal N(\nu)=\Cal P(X\times Y)$.
Otherwise, there is a probability measure on $Y$ with the same
domain and null ideal as $\nu$ (215B(vii)),
so we may suppose that $\nu Y=1$.

\medskip

{\bf (b)} The family $\Cal W$ of sets $W\subseteq X\times Y$ such that
$W[\{x\}]\in\Tau$ for every $x\in X$ and $x\mapsto\nu W[\{x\}]$ is
$\Sigma$-measurable is a Dynkin class (136A), and contains $E\times F$
whenever $E\in\Sigma$ and $F\in\Tau$;  by the Monotone Class Theorem
(136B) it includes $\Sigma\tensorhat\Tau$.

\medskip

{\bf (c)} Now suppose that
$\ofamily{\xi}{\omega_1}{W_{\xi}}$ is a disjoint family in
$\Sigma\tensorhat\Tau$.
For $n\in\Bbb N$ and $\xi<\kappa$ set

\Centerline{$E_{n\xi}=\{x:\nu W_{\xi}[\{x\}]\ge 2^{-n}\}$;}

\noindent then
$\#(\{\xi:x\in E_{n\xi}\})\le 2^{-n}$ for every $x\in X$.
It follows that $A_n=\{\xi:\xi<\omega_1$, $E_{n\xi}\notin\Cal I\}$ is
countable.   \Prf\Quer\ Otherwise, write $\frak A$ for the ccc
algebra $\Sigma/\Sigma\cap\Cal I$, and
$a_{\xi}=E_{n\xi}^{\ssbullet}$ for $\xi<\omega_1$.   Then $\frak A$ is
Dedekind complete;  set $b_{\xi}=\sup_{\xi\le\eta<\omega_1}a_{\eta}$ for
$\xi<\omega_1$ and $b=\inf_{\xi<\omega_1}b_{\xi}$.   Because $\frak A$ is
ccc, there is a $\zeta<\omega_1$ such that $b=b_{\xi}$ for every
$\xi\ge\zeta$;  because $A_n$ is uncountable, $b\ne 0$.
Choose $\sequence{i}{c_i}$ and $\sequence{i}{\xi_i}$ inductively such that
$c_0=b$ and, given that $0\ne c_i\Bsubseteq b$,
$\xi_i$ is to be such
that $c_{i+1}=a_{\xi_i}\Bcap c_i\ne 0$ and $\xi_i>\xi_j$ for every $j<i$.

\noindent Now $\inf_{i\le 2^{n}}a_{\xi_i}\Bsupseteq c_{2^n+1}$
is non-zero, so there is an
$x\in\bigcap_{i\le 2^n}E_{n\xi_i}$;  but this is impossible.\ \Bang\Qed

\medskip

{\bf (d)} This is true for every $n\in\Bbb N$, so there is a $\xi<\omega_1$
such that $\xi\notin A_n$ for every $n$, that is, $E_{n\xi}\in\Cal I$ for
every $n$.   But in this case

\Centerline{$\{x:W_{\xi}[\{x\}]\notin\Cal N(\nu)\}
=\bigcup_{n\in\Bbb N}E_{n\xi}$}

\noindent belongs to $\Cal I$ and $W_{\xi}\in\Cal I\ltimes\Cal N(\nu)$.
As $\ofamily{\xi}{\omega_1}{W_{\xi}}$ is arbitrary,
$(\Sigma\tensorhat\Tau)\cap(\Cal I\ltimes\Cal N(\nu))$ is
$\omega_1$-saturated in $\Sigma\tensorhat\Tau$ and
$(\Sigma\tensorhat\Tau)/
(\Sigma\tensorhat\Tau)\cap(\Cal I\ltimes\Cal N(\nu))$ is ccc (316C).
}%end of proof of 527L

\leader{527M}{}\cmmnt{ The final theorem of this section
provides me with an opportunity to introduce a concept which will
be needed in \S547.

\medskip

\noindent}{\bf Definition} A Boolean algebra $\frak A$ is
{\bf harmless}\cmmnt{ (cf.\ {\smc Just 92})} if it is ccc
and whenever $\frak B$ is a countable subalgebra of
$\frak A$, there is a regularly embedded countable subalgebra of $\frak A$
including $\frak B$.

\leader{527N}{Lemma} (a) A Boolean algebra with a harmless
order-dense subalgebra is itself harmless.

(b) If $\frak A$ is a Dedekind complete Boolean algebra, then it is
harmless iff
every order-closed subalgebra of $\frak A$ with countable Maharam type
has countable $\pi$-weight.

(c) For any set $I$, the regular open algebra $\RO(\{0,1\}^I)$
of $\{0,1\}^I$ is harmless, so the category algebra of $\{0,1\}^I$ is
harmless.
%also any product of spaces with countable \pi-weight

(d) If $\frak A$ has countable $\pi$-weight it is harmless.

(e) If $\frak A$ is a harmless Boolean algebra, $\frak B$ is a Boolean
algebra and $\pi:\frak A\to\frak B$ is a surjective order-continuous
Boolean homomorphism, then
$\frak B$ is harmless.   In particular,
any principal ideal of a harmless Boolean algebra is harmless.

\proof{{\bf (a)} Let $\frak A$ be a Boolean algebra with a harmless order-dense
subalgebra $\frak D$.   By 316Xj or 513E(e-iii), $\frak A$ is ccc.
Let $\frak B$ be a countable subalgebra of $\frak A$.   For each
$b\in\frak B$ let
$D_b\subseteq\frak D$ be a countable set with supremum $b$ (313K, 316E).
Let $\frak D_0$
be the subalgebra of $\frak D$ generated by $\bigcup_{b\in\frak B}D_b$.
Then $\frak D_0$ is countable, so there is a countable subalgebra
$\frak D_1$ of $\frak D$, including $\frak D_0$, which is regularly
embedded in $\frak D$.   Let $\frak C$ be the subalgebra of $\frak A$
generated by $\frak B\cup\frak D_1$.   Then $\frak C$ is countable.   Now
every member of $\frak C$ is the supremum of the members of
$\frak D_1$ it includes.  \Prf\ Set

\Centerline{$C
=\{c:c\in\frak C$, $c=\sup\{d:d\in\frak D_1$, $d\Bsubseteq c\}
=\inf\{d:d\in\frak D_1$, $c\Bsubseteq d\}\}$.}

\noindent Then $C$ is
closed under union (use 313Bd) and complementation (313A), and includes
$\frak B\cup\frak D_1$, so $C=\frak C$.\ \Qed

It follows that $\frak C$ is regularly embedded in $\frak A$, because if
$C\subseteq\frak C$ has supremum $1$ in $\frak C$ then
$\bigcup_{c\in C}\{d:d\in\frak D_1$, $d\Bsubseteq c\}$ must have supremum
$1$ in $\frak C$ and therefore in $\frak D_1$ (because
$\frak D_1\subseteq\frak C$) and in $\frak D$ (because $\frak D_1$ is
regularly embedded in $\frak D$) and in $\frak A$ (because $\frak D$ is
regularly embedded in $\frak A$).   But this means that $\sup C$ must be
$1$ in $\frak A$.   As $C$ is arbitrary, $\frak C$ is
regularly embedded.   As $\frak B$ is arbitrary, $\frak A$ is harmless.

\medskip

{\bf (b)(i)} Suppose that $\frak A$ is harmless and that
$\frak B\subseteq\frak A$ is
an order-closed subalgebra of countable Maharam type.
Let $B\subseteq\frak B$ be a
countable set which $\tau$-generates $\frak B$,
and $\frak B_0$ the algebra generated
by $B$;  let $\frak C$ be a countable subalgebra of $\frak A$,
including $\frak B_0$,
which is regularly embedded in $\frak A$.   Let $\frak D$ be the set

\Centerline{$\{d:d\in\frak A$, $d=\sup\{c:c\in\frak C$, $c\Bsubseteq d\}
=\inf\{c:c\in\frak C$, $d\Bsubseteq c\}\}$.}

\noindent Then $\frak D$ is an order-closed subalgebra of $\frak A$.
\Prf\ As in (a) just above, it is a subalgebra.
If $D\subseteq\frak D$ is a non-empty set with supremum
$a$ in $\frak A$, set $C=\{c:c\in\frak C$, $c\Bsubseteq a\}$,
$C'=\{c:c\in\frak C$, $a\Bsubseteq c\}$.   Then
$a$ is an upper bound for $C$ and a lower bound for $C'$.   \Quer\ If either $a$ is
not the least upper bound of $C$, or $a$ is not the greatest lower bound of $C'$, then
$A=\{c'\Bsetminus c:c'\in C'$, $c\in C\}$ is a subset of $\frak C$ with a non-zero lower
bound in $\frak A$, so $A$ has a non-zero lower bound $c^*$ in $\frak C$.
Now if $d\in D$, $c\in\frak C$ and $c\Bsubseteq d$, then $c\in C$ so $c\Bcap c^*=0$;
as $d=\sup\{c:c\in\frak C$, $c\subseteq d\}$, $d\Bcap c^*=0$.   This is true for every
$d\in D$, so $a\Bcap c^*=0$ and $1\Bsetminus c^*\in C'$;  but $c^*$ was chosen to be
included in every member of $C'$.\ \BanG\  Thus $a\in\frak D$;  as $D$ is arbitrary,
$\frak D$ is order-closed in $\frak A$.\ \Qed

Now $B\subseteq\frak C\subseteq\frak D$.
As $\frak B$ is regularly embedded in $\frak A$ (314Ga),
$\frak B\cap\frak D$ is an order-closed subalgebra of $\frak B$
including $B$, so is the whole of $\frak B$,
and $\frak B\subseteq\frak D$.   It
follows that $\pi(\frak B)\le\pi(\frak D)$ (514Eb).   But
$\frak C$ is countable and order-dense in $\frak D$,
so $\pi(\frak D)$ and $\pi(\frak B)$ are countable.
As $\frak B$ is arbitrary, $\frak A$ satisfies the declared condition.

\medskip

\quad{\bf (ii)} Now suppose that $\frak A$ satisfies the condition.   Note first that
$\frak A$ is ccc.   \Prf\Quer\ Suppose, if possible, otherwise;  let
$\ofamily{\xi}{\omega_1}{a_{\xi}}$ be a disjoint family in
$\frak A\setminus\{0\}$.   Replacing $a_0$ by
$a_0\Bcup(1\Bsetminus\sup_{\xi<\omega_1}a_{\xi})$ if necessary, we
may suppose that $\sup_{\xi<\omega_1}a_{\xi}=1$.   The map
$I\mapsto\sup_{\xi\in I}a_{\xi}:\Cal P\omega_1\to\frak A$
is an injective order-continuous
Boolean homomorphism, so its image $\frak B$ is an order-closed subalgebra of
$\frak A$ isomorphic to $\Cal P\omega_1$.   Now
$\tau(\frak B)=\tau(\Cal P\omega_1)=\omega$ (514Ef, or otherwise),
but $\pi(\frak B)=\omega_1$;  which is supposed to be impossible.\
\Bang\Qed

If $\frak B$ is a countable subalgebra of $\frak A$,
let $\frak B_1$ be the order-closed
subalgebra of $\frak A$ which it generates.   Then
$\tau(\frak B_1)\le\omega$ so
$\pi(\frak B_1)\le\omega$, and there is a countable subalgebra $\frak C$ of
$\frak B_1$ which is order-dense in $\frak B_1$;  of course we may suppose that
$\frak B\subseteq\frak C$.   Now the identity maps from $\frak C$ to $\frak B_1$ and
from $\frak B_1$ to $\frak A$ are both order-continuous, so their composition also is,
and $\frak C$ is regularly embedded in $\frak A$.   As $\frak B$ is arbitrary,
$\frak A$ is harmless.

\medskip

{\bf (c)} All regular open algebras are Dedekind complete.
If $\frak B\subseteq\RO(\{0,1\}^I)$ is an order-closed subalgebra with
countable
Maharam type, let $\sequencen{G_n}$ be a sequence in $\frak B$ which
$\tau$-generates $\frak B$.   Every regular open subset of $\{0,1\}^I$ is
determined by coordinates in some countable set (4A2E(b-i)), so there is a
countable $J\subseteq I$ such that every $G_n$ is determined by coordinates
in $J$.   Let $\pi_J:\{0,1\}^I\to\{0,1\}^J$ be the restriction map;  then
we have an injective order-continuous Boolean homomorphism
$H\mapsto\pi_J^{-1}[H]:\RO(\{0,1\}^J)\to\RO(\{0,1\}^I)$
(4A2B(f-iii)).   Let $\frak C$ be the
image of this homomorphism, so that $\frak C$ is an order-closed subalgebra of
$\RO(\{0,1\}^I)$.   If $H_n=\pi_J[G_n]$
then $H_n$ is regular and open for each $n$ (4A2B(f-iii) again, because $\pi_J$
is surjective and $\pi_J^{-1}[H_n]=G_n$ is regular and open), so
$G_n=\pi_J^{-1}[H_n]\in\frak C$;  accordingly $\frak B\subseteq\frak C$.   Now
$\frak B$ is an order-closed subalgebra of $\frak C$ so

\Centerline{$\pi(\frak B)\le\pi(\frak C)=\pi(\{0,1\}^J)\le\omega$.}

\noindent As $\frak B$ is arbitrary, $\RO(\{0,1\}^I)$ satisfies the
condition of (b) and is harmless.

Of course it follows at once that the category algebra is harmless, because
it is isomorphic to the regular open algebra (514If-514Ig).

\medskip

{\bf (d)} Let $D$ be a countable order-dense set in $\frak A$.   If $\frak B$ is a
countable subalgebra of $\frak A$, let $\frak C$ be the subalgebra of $\frak A$
generated by $D\cup\frak B$;  then $\frak C$ is countable, includes $\frak B$ and is
order-dense, therefore regularly embedded in $\frak A$.   As $\frak B$ is arbitrary,
$\frak A$ is harmless.

\medskip

{\bf (e)} Let $\frak D\subseteq\frak B$ be a countable subalgebra.   Because
$\pi[\frak A]=\frak B$, there is a countable subalgebra $\frak C$ of $\frak A$ such
that $\pi[\frak C]=\frak D$.   Let $\frak C_1\supseteq\frak C$ be a countable
regularly embedded subalgebra of $\frak A$.   Then
$\frak D_1=\pi[\frak C_1]$ is regularly
embedded in $\frak B$.   \Prf\ Let $D\subseteq\frak D_1$ be a non-empty set such
that $1$ is not the least upper bound of $D$ in $\frak B$.   Set
$C=\frak C_1\cap\pi^{-1}[D\cup\{0\}]$;  then $1$ is not the least upper bound of
$\pi[C]$ in $\frak B$, so (because $\pi$ is order-continuous)
$1$ is not the least upper bound of $C$ in $\frak A$.
Because $\frak C_1$ is regularly embedded in $\frak A$, there is a non-zero
$c_0\in\frak C_1$ such that $c_0\Bcap c=0$ for every $c\in C$.   In particular,
$c_0\notin C$ and $\pi c_0\ne 0$.   But we also have $\pi c\Bcap\pi c_0=0$ for every
$c\in C$, that is, $d\Bcap\pi c_0=0$ for every $d\in D$, and $1$ is not the least
upper bound of $D$ in $\frak D_1$.   As $D$ is arbitrary, $\frak D_1$ is regularly
embedded.\ \QeD\   Of course $\frak D_1$ is countable.   As $\frak D$ is arbitrary,
$\frak B$ is harmless.

If $c\in\frak A$ then $a\mapsto a\Bcap c$ is an order-continuous homomorphism onto the
principal ideal $\frak A_c$ generated by $c$, so $\frak A_c$ is harmless.
}%end of proof of 527N

\leader{527O}{Theorem} Let $(X,\Sigma,\mu)$ be a $\sigma$-finite measure
space and
$Y$ a topological space such that the category algebra $\frak G$ of $Y$ is
harmless.   Write $\Cal L$ for
 $(\Sigma\tensorhat\Cal B(Y))\cap(\Cal N(\mu)\ltimes\Cal M(Y))$ and
$\frak A$ for the measure algebra of $\mu$.   Then
$\frak C=(\Sigma\tensorhat\Cal B(Y))/\Cal L$
is ccc, and is isomorphic to the Dedekind completion of the
free product $\frak A\otimes\frak G$.
If neither $\frak A$ nor $\frak G$ is trivial,
the isomorphism corresponds to embeddings
$E^{\ssbullet}\mapsto(E\times Y)^{\ssbullet}:\frak A\to\frak C$ and
$F^{\ssbullet}\mapsto(X\times F)^{\ssbullet}:\frak B\to\frak C$.

\proof{ Write $\frak S$ for the topology of $Y$.

\medskip

{\bf (a)} Let $\Cal W$ be the family of all sets of the form
$\bigcup_{n\in\Bbb N}E_n\times H_n$, where $E_n\in\Sigma$ and $H_n\subseteq Y$ is
open for every $n$.   Then for any $W\in\Cal W$ there is a $W'\in\Cal W$ such that
$W'\symmdiff((X\times Y)\setminus W)\in\Cal L$.   \Prf\ Express
$W$ as $\bigcup_{n\in\Bbb N}E_n\times H_n$ where $E_n\in\Sigma$ and $H_n\in\frak S$
for each $n$.   Let
$\frak D$ be the order-closed subalgebra of $\frak G$ generated by
$\{H_n^{\ssbullet}:n\in\Bbb N\}$.   Because $\frak G$ is harmless and Dedekind
complete, $\pi(\frak D)\le\omega$ (527Nb);  let $\sequencen{G_n}$
be a sequence in $\frak S$
such that $\{G_n^{\ssbullet}:n\in\Bbb N\}$ is a $\pi$-base for
$\frak D$;  we may suppose that any non-empty open subset of any $G_n$ is non-meager.
Let $\frak S_1$ be the second-countable topology on $Y$ generated by
$\{H_n:n\in\Bbb N\}\cup\{G_n:n\in\Bbb N\}$, and
$\Cal B_1(Y)\subseteq\Cal B(Y)$ the corresponding Borel
$\sigma$-algebra.   Then $V^{\ssbullet}\in\frak D$ for every
$V\in\frak S_1$, because
$V$ is the union of a countable family of sets all with images in
$\frak D$.   If $V\in\frak S_1$ is dense for $\frak S_1$, and
$n\in\Bbb N$ is such that $G_n$ is non-empty, $V\cap G_n\ne\emptyset$ so
$V^{\ssbullet}\Bcap G_n^{\ssbullet}\ne 0$,
by the choice of the $G_n$.   But this means that $V^{\ssbullet}=1$, that
is, $V$ is comeager for the original topology of $Y$.

Now $W$ and $(X\times Y)\setminus W$ belong to
$\Sigma\tensorhat\Cal B_1(Y)$.   By 527I, there are $W'$ and
$\sequencen{D_n}$ such that

\inset{$((X\times Y)\setminus W)\symmdiff W'
\subseteq\bigcup_{n\in\Bbb N}D_n$,

$W'$ is expressible as $\bigcup_{n\in\Bbb N}F_n\times V_n$ where
$F_n\in\Sigma$ and $V_n\in\frak S_1$ for every $n$,

every $D_n$ belongs to $\Sigma\tensorhat\Cal B_1(Y)$,

for every $x\in X$ and $n\in\Bbb N$,
$D_n[\{x\}]$ is closed and nowhere dense for $\frak S_1$.}

\noindent Evidently $W'\in\Cal W$;  but we have just seen that sets which are closed
and nowhere dense for $\frak S_1$ are meager for $\frak S$.   So every $D_n$ belongs
to $\Cal L$ and $((X\times Y)\setminus W)\symmdiff W'\in\Cal L$.\ \Qed

\medskip

{\bf (b)} It follows (as in the proof of 527I) that
$\Cal V=\{W\symmdiff D:W\in\Cal W$, $D\in\Cal L\}$ is a $\sigma$-algebra of sets, and as
$E\times H\in\Cal W$ for every $E\in\Sigma$ and $H\in\frak S$,
$\Cal V=\Sigma\tensorhat\Cal B(Y)$.

\medskip

{\bf (c)} $\frak C$ is ccc.   \Prf\Quer\ Otherwise, there is a disjoint family
$\ofamily{\xi}{\omega_1}{e_{\xi}}$
in $\frak C\setminus\{0\}$.   For each $\xi<\omega_1$, there is a
$V_{\xi}\in(\Sigma\tensorhat\Cal B(Y))\setminus\Cal L$ such that
$V_{\xi}^{\ssbullet}=e_{\xi}$, and a $W_{\xi}\in\Cal W$ such that
$V_{\xi}\symmdiff W_{\xi}\in\Cal L$.   Express $W_{\xi}$ as
$\bigcup_{n\in\Bbb N}E_{\xi n}\times H_{\xi n}$;  as $W_{\xi}\notin\Cal L$, there must
be an $n_{\xi}$ such that $E_{\xi}=E_{\xi,n_{\xi}}\notin\Cal N(\mu)$ and
$H_{\xi}=H_{\xi,n_{\xi}}$ is non-meager.    Since the measure algebra of
$\mu$ satisfies Knaster's condition (525Tb), there is an uncountable
$A\subseteq\omega_1$ such that $E_{\xi}\cap E_{\eta}\notin\Cal N(\mu)$
for all $\xi$, $\eta\in A$;  because
$\frak G$ is ccc, there are distinct $\xi$, $\eta\in A$ such that
$H_{\xi}\cap H_{\eta}$ is non-meager.   But also

\Centerline{$(E_{\xi}\cap E_{\eta})\times(H_{\xi}\cap H_{\eta})
\subseteq W_{\xi}\cap W_{\eta}\in\Cal L$}

\noindent because
$(W_{\xi}\cap W_{\eta})^{\ssbullet}=e_{\xi}\Bcap e_{\eta}=0$.
So this is impossible.\ \Bang\Qed

Thus $\frak C$ is ccc.   As it is Dedekind $\sigma$-complete (314C),
it is Dedekind complete (316Fa).

\medskip

{\bf (d)} If either $\mu X=0$ or $Y$ is meager, then
$\frak A\otimes\frak G$ and $\frak C$ are trivially isomorphic, and we can
stop.   Otherwise,
the map $E\mapsto(E\times Y)^{\ssbullet}:\Sigma\to\frak C$ is a Boolean
homomorphism with kernel $\Sigma\cap\Cal N(\mu)$, so induces a Boolean
homomorphism
$\pi_1:\frak A\to\frak C$.   Similarly, we have a Boolean homomorphism
$\pi_2:\frak G\to\frak C$ defined by setting
$\pi_2(F^{\ssbullet})=(X\times F)^{\ssbullet}$ for $F\in\Cal B(Y)$.
These now give
us a Boolean homomorphism $\phi:\frak A\otimes\frak G\to\frak C$ defined
by saying that

\Centerline{$\psi(E^{\ssbullet}\otimes F^{\ssbullet})
=\pi_1(E^{\ssbullet})\Bcap\pi_2(F^{\ssbullet})=(E\times F)^{\ssbullet}$}

\noindent for $E\in\Sigma$ and $F\in\Cal B(Y)$ (315Jb).   If
$E\in\Sigma\setminus\Cal N(\mu)$ and $F\in\Cal B(Y)\setminus\Cal M(Y)$,
then $E\times F\notin\Cal L$;  so $\phi$ is injective
(use 315Kb).   If $c\in\frak C$ is
non-zero, it is expressible as $W^{\ssbullet}$ for some
$W\in\Cal W\setminus\Cal L$;
there must now be $E\in\Sigma$, $F\in\Cal B(Y)$ such that
$E\times F\subseteq W$ and $E\times F\notin\Cal L$, so that
$\phi(E^{\ssbullet}\otimes F^{\ssbullet})$ is
non-zero and included in $w$.
Thus $\phi[\frak A\otimes\frak G]$ is isomorphic to
$\frak A\otimes\frak G$ and is an order-dense
subalgebra of the Dedekind complete Boolean algebra $\frak C$;
it follows that
$\frak C$ can be identified with the Dedekind completion of
$\frak A\otimes\frak G$.
}%end of proof of 527O

\exercises{\leader{527X}{Basic exercises $\pmb{>}$(a)}
%\sqheader 527Xa
Show that there is a set belonging to $\Cal N\ltimes\Cal N$ which has
full outer measure for Lebesgue measure in the plane.   \Hint{enumerate
the compact non-negligible subsets of the plane as
$\ofamily{\xi}{\frak c}{K_{\xi}}$ (4A3Fa);  note that the projection
$L_{\xi}$ of $K_{\xi}$ onto the first
coordinate is always non-negligible, therefore uncountable, therefore of
cardinal $\frak c$ (423L);
choose $s_{\xi}\in L_{\xi}\setminus\{s_{\eta}:\eta<\xi\}$ and
$t_{\xi}\in K_{\xi}[\{s_{\xi}\}]$ for each $\xi$;  consider
$\{(s_{\xi},t_{\xi}):\xi<\frak c\}$.}
%527B

\sqheader 527Xb Show that there is a unique construction of iterated
skew products $\Cal I_0\ltimes\Cal I_1\ltimes\ldots\ltimes\Cal I_n$ such
that

(i) whenever $X_0,\ldots,X_n$ are sets and $\Cal I_j$ is an ideal of
subsets of $X_j$ for every $j$, then $\Cal I_0\ltimes\ldots\ltimes\Cal
I_n$ is an ideal of subsets of
$X_0\times\ldots\times X_n$;

(ii) whenever $X_0,\ldots,X_n$ are sets, $\Cal I_j$ is an ideal of
subsets of $X_j$ for every $j$, and $k<n$, then the natural
identification of $X_0\times\ldots\times X_n$
with $(X_0\times\ldots\times X_k)\times(X_{k+1}\times\ldots\times X_n)$
identifies $\Cal I_0\ltimes\ldots\ltimes\Cal I_n$ with
$(\Cal I_0\ltimes\ldots\ltimes I_k)\ltimes(\Cal
I_{k+1}\ltimes\ldots\ltimes I_n)$ as defined in 527B.
%527B

\spheader 527Xc Complete the analysis in 527Bb by describing what
happens if one of $X$, $Y$ is empty or one of the ideals is not proper.
%527B

\spheader 527Xd Let $X$ be a set, $\Sigma$ a $\sigma$-algebra of subsets of
$X$, and $\Cal I$ an ideal of subsets of $X$;  let $Y$ be a topological
space, $\Cal B$ its Borel $\sigma$-algebra, $\widehat{\Cal B}$ its
Baire-property algebra, and $\Cal M$ its meager ideal.   Show that
$\Cal I\ltimes_{\Sigma\tensorhat\Cal B}\Cal M
=\Cal I\ltimes_{\Sigma\tensorhat\widehat{\Cal B}}\Cal M$.
%527B

\sqheader 527Xe Let $Z$ be the Stone space of the measure algebra of
Lebesgue measure on $[0,1]$, and $f:Z\to[0,1]$ the canonical \imp\
continuous function (416V).
Let $F\subseteq[0,1]$ be a nowhere dense set which is not negligible,
and set $W=\{(x,z):x\in[0,1]$, $z\in Z$, $x+f(z)\in F\}$.   Show that
$W$ is a nowhere dense closed
set in $[0,1]\times Z$ but does not belong to
$\Cal M([0,1])\ltimes\Cal M(Z)$.
\Hint{meager subsets of $Z$ are negligible (321K).}
%527D

\sqheader 527Xf Suppose that $I$ and $J$ are sets,
$X=\{0,1\}^I$ and $Y=\{0,1\}^J$.   Show that
$\Cal M(X)\ltimes_{\Cal B(X\times Y)}\Cal M(Y)=\Cal M(X\times Y)$.
%527D 527E

\spheader 527Xg Write $\frak X$ for the class of topological spaces
which have category algebras which are atomless and with countable
$\pi$-weight.   (i) Show that the Sorgenfrey line (415Xc)
belongs to $\frak X$.   (ii) Show that the split interval
(419L) belongs to $\frak X$.   (iii) Show that if the regular open
algebra of a topological space $X$ is atomless and has countable
$\pi$-weight, then $X\in\frak X$.   (iv) Show that any open subspace of
a space in $\frak X$ belongs to $\frak X$.   (v) Show that any dense
subspace of a space in $\frak X$ belongs to $\frak X$.
(vi) Show that any comeager subspace of a space in $\frak X$ belongs to
$\frak X$.
(vii) Show that the product of countably many spaces in $\frak X$
belongs to $\frak X$.
%527F

\spheader 527Xh Show that a measurable algebra is harmless iff it is purely
atomic.
%527M

\leader{527Y}{Further exercises (a)}
%\spheader 527Ya
Show that
$\Cal I\ltimes\Cal J\ne\Cal I\rtimes\Cal J$ for any of the four cases in
which $\{\Cal I,\Cal J\}\subseteq\{\Cal M,\Cal N\}$.
%527B

\def\bigltimes{\hbox{$\vrule height9pt width0.5pt depth1.5pt \mskip-1.5mu
          \raise3pt\hbox{$\vee$}\mskip-12mu\raise-3pt\hbox{$\wedge$}$}}

\spheader 527Yb Extend the idea of 527Xb to define an ideal
$\bigltimes_{\xi<\zeta}\Cal I_{\xi}$ of subsets of $\prod_{\xi<\zeta}X_{\zeta}$
when $\zeta$ is any ordinal and $I_{\xi}$ is an ideal of subsets of $X_{\xi}$
for every $\xi<\zeta$.
%527B

\spheader 527Yc Let $X$ be a set, $\Sigma$ a $\sigma$-algebra of subsets of $X$ and
$\Cal I$ a $\sigma$-ideal of $\Sigma$ such that $\Sigma/\Cal I$ is ccc.
Let $(Y,\Tau,\nu)$ be a probability space.   Show that
$(\Sigma\tensorhat\Tau)/(\Sigma\tensorhat\Tau)\cap(\Cal I\ltimes\Cal N(\nu))$
is ccc.
%\Hint{show that if $V\in\Sigma\tensorhat\Tau$ then $x\mapsto\nu V[\{x\}]$ is
%$\Sigma$-measurable, and hence that there is no uncountable disjoint family
%in $(\Sigma\tensorhat\Tau)\setminus(\Cal I\ltimes\Cal N(\nu))$.}
%527G

\spheader 527Yd Let $(Y,\frak T)$ be a topological space.   Show that there
is a topology $\frak S$ on $Y$, coarser than $\frak T$, such that the
weight of $(Y,\frak S)$ is equal to the $\pi$-weight of $(Y,\frak T)$ and
the two topologies have the same nowhere dense sets, the same meager ideal
and the same Baire-property algebras.
%527J

\spheader 527Ye Let $X$ be a topological space with a $\sigma$-finite
measure $\mu$ such that $\mu$ has countable Maharam type and every
measurable set can be expressed as the
symmetric difference of a Borel set and a negligible set.   Let
$Y$ be a topological space with a countable $\pi$-base.   Show that
$\Cal N(\mu)\ltimes_{\Cal B(X\times Y)}\Cal M(Y)\prT\Cal N(\mu)\times\Cal N$.
%527J

\spheader 527Yf Let $\familyiI{\frak A_i}$ be a family of harmless
Boolean algebras satisfying Knaster's condition,
and $\frak A$ their free product (315I).   Show that $\frak A$ is harmless.
%527N

\leaveitout{\spheader 527Y? Let $\mu$ be a $\sigma$-finite Borel probability measure
on a topological space $X$, and $Y$ a topological
space such that its category algebra is harmless.   Show that
$\Cal B(X\times Y)/\Cal B(X\times Y)\cap(\Cal N(\mu)\ltimes\Cal M(Y))$ can
be identified with the Dedekind completion of
$\frak A\otimes\frak G$, where $\frak A$ is the measure algebra of $\mu$
and $\frak G$ is the category algebra of $Y$.
%527O
\query false, try $X$ a \rvm\ cardinal and $\mu$ a witnessing probebility and
$Y=\{0,1\}^X$}
}%end of exercises

\endnotes{
\Notesheader{527} Skew products of ideals have been used many times for
special purposes, and we are approaching the point at which it would be
worth developing a general theory
of such products.   I am not really attempting to do this here, though
the language of 527B is supposed to point to the right questions.   My
primary aim in this section is
to show that $\Cal M\ltimes_{\Cal B(\BbbR^2)}\Cal N$ and
$\Cal N\ltimes_{\Cal B(\BbbR^2)}\Cal M$ are very different (527H, 527K).
Of course the difference appears
only when the continuum hypothesis is false (513Xf, 513Xr).

The version of the Kuratowski-Ulam theorem given in 527D is a natural
one from the point of view of this chapter, but you should be aware that
there are many more cases in which $\Cal M^*=\Cal M(X\times Y)$;  see 527Xf
and {\smc Fremlin Natkaniec \& Rec{\l}aw 00}.
The statement of 527J includes the phrase `quasi-Radon measure'.
Actually we do not really need either $\tau$-additivity or inner
regularity with respect to closed sets.
What we need is a measure $\mu$ such that $\Cal N(\mu)\prT\Cal N$ and
the Borel sets generate the measure algebra (527Ye).   The argument for
527J betrays its origin in the case
$X=Y=[0,1]$, which is of course also the natural home of
527C-527F.   %527C 527D 527E 527F
Some of the complications of the argument are due to its being written out
for spaces of countable $\pi$-weight;  an alternative approach would
start with a reduction to the case in which $Y$ is second-countable
(527Yd).

It is interesting that all four of the quotient algebras

\Centerline{$\Cal B(\BbbR^2)/\Cal B(\BbbR^2)\cap(\Cal M\ltimes\Cal M)$,
\quad$\Cal B(\BbbR^2)/\Cal B(\BbbR^2)\cap(\Cal M\ltimes\Cal N)$,}

\Centerline{$\Cal B(\BbbR^2)/\Cal B(\BbbR^2)\cap(\Cal N\ltimes\Cal M)$,
\quad$\Cal B(\BbbR^2)/\Cal B(\BbbR^2)\cap(\Cal N\ltimes\Cal N)$}

\noindent are ccc (see 527E, 527Yc, 527O, 527Bc and also 527L).
This should not be taken for
granted;  for a variety of examples of quotient algebras associated with
$\sigma$-ideals see {\smc Fremlin 03}.

}%end of notes
\discrpage

\leaveitout{Pawlikowski 86:  " \Cal R*\Cal C  can be completely embedded
in  \Cal C*\Cal R " (cor. 3.3)
}
\frfilename{mt528.tex}
\versiondate{10.2.11/11.2.11}
\copyrightdate{2007}

\def\chaptername{Cardinal functions of measure theory}
\def\sectionname{Amoeba algebras}

\def\AM{\mathop{\text{AM}}}
%\def\checked{$\surd$}

\newsection{528}

In the course of investigating the principal consequences of Martin's
axiom, {\smc Martin \& Solovay 70} introduced the partially ordered set
of open subsets of $\Bbb R$ with measure strictly less than $\gamma$,
for $\gamma>0$ (528O).   Elementary extensions of this idea lead us to a
very interesting class of partially ordered sets, which I study here in
terms of their regular open algebras, the `amoeba algebras' (528A).   Of
course the most important ones are those associated with Lebesgue
measure, and these are closely related to `localization posets'
(528I), themselves intimately connected with the localization relations of
522K.   In the second half of the section I look at the cardinal functions
of these algebras, of which the most interesting seems to be Maharam type
(528V).

As elsewhere in this chapter, I will write
$(\frak B_{\kappa},\bar\nu_{\kappa})$ for the measure algebra of the
usual measure on $\{0,1\}^{\kappa}$.   In any measure algebra
$(\frak A,\bar\mu)$
I will write $\frak A^f=\{a:a\in\frak A$, $\bar\mu a<\infty\}$.

\leader{528A}{Amoeba algebras}
Let $(\frak A,\bar\mu)$ be a measure algebra.

\spheader 528Aa If $0<\gamma\le\bar\mu 1$, the {\bf amoeba algebra}
$\AM(\frak A,\bar\mu,\gamma)$ is the regular open algebra
$\RO^{\uparrow}(P)$ where $P=\{a:a\in\frak A$, $\bar\mu a<\gamma\}$,
ordered by $\Bsubseteq$.

\spheader 528Ab
The {\bf variable-measure amoeba algebra}
$\AM^*(\frak A,\bar\mu)$\cmmnt{ ({\smc Truss 88})} is the regular open
algebra $\RO^{\uparrow}(P')$ where

\Centerline{$P'=\{(a,\alpha):a\in\frak A$,
$\alpha\in\ocint{\bar\mu a,\bar\mu 1}\}$,}

\noindent ordered by saying that

\Centerline{$(a,\alpha)\le(b,\beta)$ if $a\Bsubseteq b$ and
$\beta\le\alpha$.}

\leader{528B}{}\dvAnew{2011}\cmmnt{ It may help to have the following
simple facts set out straight away.

\wheader{528B}{4}{2}{2}{72pt}

\noindent}{\bf Lemma} Let $(\frak A,\bar\mu)$ be a measure algebra
and $0<\gamma\le\bar\mu 1$.   Set
$P=\{a:a\in\frak A$, $\bar\mu a<\gamma\}$.

(a) Two elements $a$, $b\in P$ are compatible upwards in $P$ iff
$\bar\mu(a\Bcup b)<\gamma$.

(b) Suppose that $(\frak A,\bar\mu)$ is semi-finite and atomless.

\quad(i) $P$ is separative upwards, so
$\coint{a,\infty}\in\RO^{\uparrow}(P)$ for every $a\in P$.

\quad(ii) If $A\subseteq P$ is non-empty, then the infimum
$\inf_{a\in A}\coint{a,\infty}$ is empty unless $\sup A$ is defined in
$\frak A$ and belongs to $P$, and in this case
$\inf_{a\in A}\coint{a,\infty}=\coint{\sup A,\infty}$.

\proof{{\bf (a)} $\coint{a,\infty}\cap\coint{b,\infty}
=\{c:a\Bcup b\Bsubseteq c\in P\}$ is non-empty iff $a\Bcup b\in P$.

\medskip

{\bf (b)(i)} Let $a$, $b\in P$ be
such that $a\notBsubseteq b$.   If $\bar\mu(a\Bcup b)\ge\gamma$ then $a$
and $b$ are already incompatible upwards.   Otherwise,
$\bar\mu(1\Bsetminus(a\Bcup b))\ge\gamma-\bar\mu(a\Bcup b)$.   Because
$(\frak A,\bar\mu)$ is atomless and semi-finite, there is a
$d\Bsubseteq 1\Bsetminus(a\Bcup b)$ such that
$\bar\mu d=\gamma-\bar\mu(a\Bcup b)$.   Set $c=b\Bcup d$.  Then

\Centerline{$\bar\mu c=\gamma-\bar\mu(a\Bsetminus b)<\gamma
=\bar\mu(a\Bcup c)$,}

\noindent so
$c\in\coint{b,\infty}\subseteq P$, while $a$ and $c$ are incompatible
upwards in $P$.   As $a$ and $b$ are arbitrary, $P$ is separative upwards.

By 514Me, it follows that $\coint{a,\infty}$ is a regular up-open set for
every $a\in P$.

\medskip

\quad{\bf (ii)} This is a re-phrasing of 514Mf.
}%end of proof of 528B

\leader{528C}{Proposition}\dvAformerly{5{}28B}
Suppose that $(X,\Sigma,\mu)$ is a measure
space, $(\frak A,\bar\mu)$ its measure algebra and $0<\gamma\le\mu X$.
If $\Cal E\subseteq\Sigma$ is any family such that $\mu$ is outer
regular with respect to $\Cal E$, then $\AM(\frak A,\bar\mu,\gamma)$ is
isomorphic to $\RO^{\uparrow}(\{E:E\in\Cal E$, $\mu E<\gamma\})$.

\proof{ Set $P=\{a:a\in\frak A$, $\bar\mu a<\gamma\}$,
$Q=\{E:E\in\Cal E$, $\mu E<\gamma\}$.   Because $\mu$ is outer regular
with respect to $\Cal E$, the map
$G\mapsto G^{\ssbullet}:Q\to\frak A$ maps $Q$ onto a cofinal subset $P'$
of $P$.   Moreover, two elements $E_0$ and $E_1$ of $Q$ are compatible
upwards in $Q$ iff $\mu(E_0\cup E_1)<\gamma$ iff $E_0^{\ssbullet}$
and $E_1^{\ssbullet}$ are compatible upwards in $P$.   By 514R,
$\RO^{\uparrow}(P)$ and $\RO^{\uparrow}(Q)$ are isomorphic.
}%end of proof of 528C

\leader{528D}{Proposition} (a)\cmmnt{ ({\smc Truss 88})} Let
$(\frak A,\bar\mu)$ be an atomless homogeneous probability algebra.
Then the amoeba algebras $\AM(\frak A,\bar\mu,\gamma)$ and
$\AM(\frak A,\bar\mu,\gamma')$ are isomorphic for all $\gamma$,
$\gamma'\in\ooint{0,1}$.

(b)\dvAformerly{5{}28C}
Let $(\frak A,\bar\mu)$ be a
non-totally-finite atomless quasi-homogeneous measure algebra\cmmnt{
(definition:  374G)}.   Then all the amoeba algebras
$\AM(\frak A,\bar\mu,\gamma)$, for $\gamma>0$, are isomorphic.

\proof{{\bf (a)(i)} Set $P=\{a:a\in\frak A$, $\bar\mu a<\gamma\}$, and let
$\kappa$ be the Maharam type of $\frak A$.   Then the
upwards cellularity of $P$ is at most $\kappa$.   \Prf\Quer\ Otherwise,
there is an up-antichain $A\subseteq P$ with cardinality $\kappa^+$.
Let $\epsilon>0$ be such that
$A'=\{a:a\in A$, $\bar\mu a\le\gamma-\epsilon\}$ has cardinal
$\kappa^+$.   Because the topological density of $\frak A$ is $\kappa$
(521Ea), there must be distinct $a$, $a'\in A'$ such that
$\bar\mu(a\Bsymmdiff a')<\epsilon$;  but in this case
$\bar\mu(a\Bcup a')<\gamma$, so that $a\Bcup a'$ is an upper bound for
$\{a,a'\}$ in $P$.\ \Bang\Qed

\medskip

\quad{\bf (ii)} If
$1-\sqrt{1-\gamma}\le\alpha<\gamma$ and $D$ is a countable subset of
$\ooint{\alpha,\gamma}$ such that $\sup D=\gamma$, then there is a
maximal up-antichain
$\family{(t,\xi)}{D\times\kappa}{a_{t\xi}}$ in $P$ such that
$\bar\mu a_{t\xi}=t$ for every $t\in D$, $\xi<\kappa$.   \Prf\ Start
with a stochastically independent family
$\family{(t,\xi)}{D\times\kappa}{c_{t\xi}}$ of elements of $\frak A$
with $\bar\mu c_{t\xi}=t$ for all $t\in D$, $\xi<\kappa$.   Because
$\alpha\ge 1-\sqrt{1-\gamma}$,
$A=\family{(t,\xi)}{D\times\kappa}{c_{t\xi}}$ is an up-antichain in
$P$.   Next, because $\sup D=\gamma$,
$Q=\{a:a\in P$, $\bar\mu a\in D\}$ is cofinal with $P$.   So there is a
maximal up-antichain $A'\supseteq A$ such that $A'\subseteq Q$ (513Aa).
Now (because $c^{\uparrow}(P)\le\kappa$) $\{a:a\in A'$, $\bar\mu a=t\}$
has cardinal $\kappa$ for every $t\in D$, so we can enumerate $A'$ as
$\family{(t,\xi)}{D\times\kappa}{a_{t\kappa}}$ in $P$ where
$\bar\mu a_{t\xi}=t$ for every $t\in D$ and $\xi<\kappa$.\ \Qed

\medskip

\wheader{528D}{0}{0}{0}{36pt}
\quad{\bf (iii)} There are $\alpha$, $\alpha'\in\ooint{0,1}$ such that

\Centerline{$1-\sqrt{1-\gamma}\le\alpha<\gamma$,
\quad$1-\sqrt{1-\gamma'}\le\alpha'<\gamma'$,
\quad$\Bover{\gamma-\alpha}{1-\alpha}
=\Bover{\gamma'-\alpha'}{1-\alpha'}$.}

\noindent\Prf\ We need consider only the case $\gamma\le\gamma'$.   Set

\Centerline{$\beta=\Bover1{\sqrt{1-\gamma}}-1$,
\quad$\alpha=\gamma-\beta(1-\gamma)$,
\quad$\alpha'=\gamma'-\beta(1-\gamma')$.}

\noindent Then
$\Bover{\gamma-\alpha}{1-\gamma}=\beta
=\Bover{\gamma'-\alpha'}{1-\gamma'}$.   Of course $\alpha\le\gamma$ and
$\alpha'\le\gamma'$.   On the other side,
$\alpha=1-\sqrt{1-\gamma}$, while $\beta\le\Bover1{\sqrt{1-\gamma'}}-1$
so $\alpha'\ge 1-\sqrt{1-\gamma'}$.\ \Qed

\medskip

\quad{\bf (iv)} If $a\in P$,
$\{b:a\Bsubseteq b\in P\}$ is isomorphic, as partially ordered set, to
$\{b:b\in\frak A$, $\bar\mu b<\Bover{\gamma-\bar\mu a}{1-\bar\mu a}\}$.
\Prf\ The principal ideal $\frak A_{1\Bsetminus a}$ generated by
$1\Bsetminus a$ is isomorphic, up to a scalar multiple of the measure,
to $\frak A$, and $\{b:a\Bsubseteq b\in P\}$ is isomorphic, as partially
ordered set, to $\{b:b\Bsubseteq 1\Bsetminus a$,
$\bar\mu b<\gamma-\bar\mu a\}$.\ \Qed

\medskip

\quad{\bf (v)} For each $n\in\Bbb N$, set
$\alpha_n=\gamma-2^{-n}(\gamma-\alpha)$,
$\alpha'_n=\gamma'-2^{-n}(\gamma'-\alpha')$;  then

\Centerline{$\Bover{1-\alpha'_n}{\gamma'-\alpha_n'}
=1+\Bover{1-\gamma'}{\gamma'-\alpha_n'}
=1+2^n\Bover{1-\gamma'}{\gamma'-\alpha'}
=\Bover{1-\alpha_n}{\gamma-\alpha_n}$,
\quad$\Bover{\gamma'-\alpha_n'}{1-\alpha_n'}
=\Bover{\gamma-\alpha_n}{1-\alpha_n}$}

\noindent for every $n\in\Bbb N$.   Set
$P'=\{a:a\in\frak A$, $\bar\mu a<\gamma'\}$.    By (b), we have a
maximal up-antichain $\family{(n,\xi)}{\Bbb N\times\xi}{a_{n\xi}}$ in
$P$ such that $\bar\mu a_{n\xi}=\alpha_n$ for all $n\in\Bbb N$ and
$\xi<\kappa$;  similarly, there is a maximal up-antichain
$\family{(n,\xi)}{\Bbb N\times\xi}{a'_{n\xi}}$ in $P'$ such that
$\bar\mu a'_{n\xi}=\alpha'_n$ for all $n\in\Bbb N$ and $\xi<\kappa$.
Now, for each $n\in\Bbb N$ and $\xi<\kappa$, $\coint{a_{n\xi},\infty}$,
taken in $P$, is isomorphic, as partially ordered set, to
$[a'_{n\xi},\infty[$, taken in $P'$, by (d).   So

$$\eqalignno{\RO^{\uparrow}(P)
&\cong\prod_{n\in\Bbb N,\xi<\kappa}
  \RO^{\uparrow}(\coint{a_{n\xi},\infty})\cr
\displaycause{514Nf}
&\cong\prod_{n\in\Bbb N,\xi<\kappa}
  \RO^{\uparrow}(\coint{a'_{n\xi},\infty})
\cong\RO^{\uparrow}(P').\cr}$$

\medskip

{\bf (b)} Suppose that $\beta$, $\gamma>0$.   As in Lemma 332I, we have a
partition $D$ of unity in $\frak A$ such that $\bar\mu a=\beta$ for
every $a\in D$.   Similarly, we have a partition $D'$ of unity such that
$\bar\mu a=\gamma$ for every $a\in D'$.   By 332E,
$\#(D)=\#(D')=c(\frak A)$.   Let $h:D\to D'$ be a bijection.   If $d\in D$,
the principal ideals $\frak A_d$, $\frak A_{h(d)}$ have the same
Maharam type, because $(\frak A,\bar\mu)$ is quasi-homogeneous (374H),
and are therefore isomorphic as measure algebras, up to a
scalar factor of the measure;  let $\pi_d:\frak A_d\to\frak A_{h(d)}$ be
a Boolean isomorphism such that
$\bar\mu(\pi_da)=\Bover{\gamma}{\beta}\bar\mu a$ for every
$a\Bsubseteq d$.   Now we have a function $\pi:\frak A^f\to\frak A^f$
defined by saying that $\pi a=\sup_{d\in D}\pi_d(a\Bcap d)$ whenever
$\bar\mu a<\infty$, and $\pi$ is a Boolean ring automorphism such that
$\bar\mu\pi a=\Bover{\gamma}{\beta}\bar\mu a$ for every $a\in\frak A^f$.
But now $\pi$ includes an isomorphism between the partially ordered sets
$\{a:\bar\mu a<\beta\}$ and $\{a:\bar\mu a<\gamma\}$, so their regular
open algebras $\AM(\frak A,\bar\mu,\beta)$ and
$\AM(\frak A,\bar\mu,\gamma)$ are isomorphic.
}%end of proof of 528D

\leader{528E}{Lemma}\dvAnew{2011}
Let $(\frak A,\bar\mu)$ be an atomless semi-finite
measure algebra.   Then there is a family
$\family{\alpha}{[0,\bar\mu 1]}{c_{\alpha}}$ in $\frak A$ such that
$c_{\alpha}\Bsubseteq c_{\beta}$ and $\bar\mu c_{\alpha}=\alpha$
whenever $0\le\alpha\le\beta\le\bar\mu 1$,
and $\alpha\mapsto c_{\alpha}$ is continuous for
the measure-algebra topology of $\frak A$.

\proof{ Because $(\frak A,\bar\mu)$ is semi-finite, there is a
non-decreasing sequence $\sequencen{e_n}$ in $\frak A^f$ such that
$\sup_{n\in\Bbb N}\bar\mu e_n=\bar\mu 1$, starting from
$e_0=0$;  set
$e=\sup_{n\in\Bbb N}e_n$, so that $\bar\mu e=\bar\mu 1$.   Then
$(\frak A_e,\bar\mu\restrp\frak A_e)$ is $\sigma$-finite and atomless.
Let $(\frak C,\bar\lambda)$ be the
measure algebra of Lebesgue measure on $\coint{0,\bar\mu 1}$.
For each $n\in\Bbb N$ set
$e'_n=e_{n+1}\Bsetminus e_n$ and
$d_n=\coint{\bar\mu e_n,\bar\mu e_{n+1}}^{\ssbullet}\in\frak C$.

Because $\frak A$ is atomless, 332P tells us that
there is for each $n\in\Bbb N$ a measure-preserving Boolean
homomorphism $\pi_n$
from the principal ideal $\frak C_{d_n}$ to a principal ideal
of $\frak A_{e'_n}$, which must be $\frak A_{e'_n}$ itself because
$\bar\mu e'_n=\bar\lambda d_n$;  by 324Kb, $\pi_n$ is order-continuous.
Assembling these, we have an order-continuous
measure-preserving Boolean homomorphism $\pi:\frak C\to\frak A_e$ defined
by setting $\pi d=\sup_{n\in\Bbb N}\pi_n(d\Bcap d_n)$ for every
$d\in\frak C$.   Now set $c_{\alpha}=\pi\coint{0,\alpha}^{\ssbullet}$ for
$\alpha\le\bar\mu 1$.   Because $\pi$ is continuous for the measure-algebra
topologies of $\frak C$ and $\frak A_e$ (324Fc), or otherwise,
$\alpha\mapsto c_{\alpha}$ is continuous.
}%end of proof of 528E

\leader{528F}{Proposition} Let $(\frak A,\bar\mu)$ be a semi-finite
measure algebra, and $\gamma\in\ooint{0,\infty}$.

(a) Suppose that $e\in\frak A$ and $\bar\mu e\ge\gamma$.
If $\frak A_e$ is atomless, then
$\AM(\frak A_e,\bar\mu\restrp\frak A_e,\gamma)$ can be regularly
embedded in $\AM(\frak A,\bar\mu,\gamma)$.

(b)\dvAnew{2011}
Suppose that $\frak A$ is atomless, and that $\gamma<\bar\mu 1$.
Let $\sequence{k}{e_k}$ be a non-decreasing sequence in $\frak A$
with supremum $1$, and
suppose that $\bar\mu e_k\ge\gamma$ for every $k\in\Bbb N$.   Then
we have a sequence $\sequence{k}{\pi_k}$ such that
$\pi_k:\AM(\frak A_{e_k},\bar\mu\restrp\frak A_{e_k},\gamma)\to
\AM(\frak A,\bar\mu,\gamma)$ is a regular embedding for every
$k\in\Bbb N$, and
$\bigcup_{k\in\Bbb N}
\pi_k[\AM(\frak A_{e_k},\bar\mu\restrp\frak A_{e_k},\gamma)]
\,\,\tau$-generates $\AM(\frak A,\bar\mu,\gamma)$.

(c)\dvAformerly{5{}28E} Now suppose that $(\frak A,\bar\mu)$ is atomless
and quasi-homogeneous, and that $\gamma<\bar\mu 1$.   Then
$\AM(\frak A,\bar\mu,\gamma)$ can be regularly
embedded in $\AM^*(\frak A,\bar\mu)$.

\proof{{\bf (a)} Set $P=\{a:a\in\frak A$, $\bar\mu a<\gamma\}$ and
$Q=P\cap\frak A_e$.   By 528E, we have a continuous
order-preserving function
$\alpha\mapsto c_{\alpha}:[0,\bar\mu e]\to\frak A_e$ such that
$\bar\mu c_{\alpha}=\alpha$ for each $\alpha$.   If $a\in\frak A$,
then the function $\beta\mapsto\bar\mu(c_{\beta}\Bsetminus a)$
is a continuous non-decreasing function from $[0,\bar\mu e]$ onto
$[0,\bar\mu(c_{\bar\mu e}\Bsetminus a)]$, and
we can set
$\delta(a,\alpha)=\min\{\beta:\bar\mu(c_{\beta}\Bsetminus a)=\alpha\}$
whenever $0\le\alpha\le\bar\mu(c_{\bar\mu e}\Bsetminus a)$.
In this case,

\Centerline{$\bar\mu((a\Bcap e)\Bcup c_{\delta(a,\alpha)})
=\bar\mu(a\Bcap e)+\bar\mu(c_{\delta(a,\alpha)}\Bsetminus a)
=\alpha+\bar\mu(a\Bcap e)$.}

\noindent Note that $\delta(a,\alpha)\le\delta(a',\alpha')$ whenever
$a\Bsubseteq a'$ and
$\alpha\le\alpha'\le\bar\mu(c_{\bar\mu 1}\Bsetminus a')$.

If $a\in P$, then

\Centerline{$\bar\mu(c_{\bar\mu e}\Bsetminus a)
=\bar\mu c_{\bar\mu e}-\bar\mu(a\Bcap c_{\bar\mu e})
\ge\bar\mu e-\bar\mu(a\Bcap e)
\ge\bar\mu a-\bar\mu(a\Bcap e)
=\bar\mu(a\Bsetminus e)$.}

\noindent So $\delta(a,\bar\mu(a\Bsetminus e))$ is defined, and we have a
function $f$ given by the formula

\Centerline{$f(a)=(a\Bcap e)\Bcup c_{\delta(a,\bar\mu(a\Bsetminus e))}$}

\noindent for $a\in P$.
In this case $\bar\mu f(a)=\bar\mu a$, so $f(a)\in Q$, for each $a$,
and $f$, like $\delta$, is order-preserving.   Of course $f(a)=a$ for
$a\in Q$.

If $a\in P$, $b\in Q$ and $f(a)\Bsubseteq b$, there is an $a'\in P$ such
that $a\Bsubseteq a'$ and $b=f(a')$.   \Prf\
Set $a'=a\Bcup(b\Bsetminus f(a))$.   Then

\Centerline{$\bar\mu a'
=\bar\mu a+\bar\mu(b\Bsetminus f(a))
=\bar\mu f(a)+\bar\mu(b\Bsetminus f(a))
=\bar\mu b<\gamma$,}

\noindent so $a'\in P$.   Also
$b\Bsubseteq f(a)\Bcup(a'\Bcap e)\Bsubseteq f(a')$;  as
$\bar\mu b=\bar\mu a'=\bar\mu f(a')$, $b=f(a')$.\ \QeD\   So if
$Q_0\subseteq Q$ is cofinal with $Q$, $f^{-1}[Q_0]$ will be cofinal with
$P$ (as in the proof of 514P), and we have an order-continuous Boolean
homomorphism $\pi:\RO^{\uparrow}(Q)\to\RO^{\uparrow}(P)$
defined by setting $\pi H=\interior\overline{f^{-1}[H]}$ for every
$H\in\RO^{\uparrow}(Q)$.   Finally, $f[P]=f[Q]=Q$.   So $\pi$ is injective
and is a regular embedding of
$\AM(\frak A_e,\bar\mu\restrp\frak A_e,\gamma)=\RO^{\uparrow}(Q)$ into
$\AM(\frak A,\bar\mu,\gamma)=\RO^{\uparrow}(P)$.

\medskip

{\bf (b)(i)}
For each $k\in\Bbb N$, set $Q_k=P\cap\frak A_{e_k}$ and
choose functions $f_k:P\to Q_k$
and $\pi_k:\RO^{\uparrow}(Q_k)\to\RO^{\uparrow}(P)$ as in (a)
above.   If we write
$\coint{c,\infty}=\{a:c\Bsubseteq a\in P\}$ for every $c\in P$, then
$\frak A_{e_k}\cap\coint{c,\infty}=\{b:c\Bsubseteq b\in Q_k\}$ for
$k\in\Bbb N$ and $c\in Q_k$;  in this case,
$\frak A_{e_k}\cap\coint{c,\infty}\in\RO^{\uparrow}(Q_k)$,
by 528B(b-i).

\medskip

\quad{\bf (ii)} Let $\frak G$ be the order-closed subalgebra of
$\RO^{\uparrow}(P)$ generated by
$\bigcup_{k\in\Bbb N}\pi_k[\RO^{\uparrow}(Q_k)]$.   If $a\in P$, there is a
non-empty $G\in\frak G$ included in
$\coint{a,\infty}\in\RO^{\uparrow}(P)$.   \Prf\ Because
$a\Bsubseteq\sup_{k\in\Bbb N}e_k$ and $\sequence{k}{e_k}$ is
non-decreasing, there is an infinite $I\subseteq\Bbb N$ such that
$\sum_{k\in I}\bar\mu(a\Bsetminus e_k)<\gamma-\bar\mu a$.
Set $b=\sup_{k\in I}f_k(a)$.   Then

\Centerline{$\bar\mu b
\le\bar\mu a+\sum_{k\in I}\bar\mu(a\Bsetminus f_k(a))
\le\bar\mu a+\sum_{k\in I}\bar\mu(a\Bsetminus e_k)
<\gamma$}

\noindent because $f_k(a)\Bsupseteq a\Bcap e_k$ for every $k$, by the
construction in (a).   Thus $b\in P$.   Also

\Centerline{$\bar\mu(a\Bsetminus b)
\le\inf_{k\in\Bbb N}\bar\mu(a\Bsetminus f_k(a))=0$,}

\noindent so $a\Bsubseteq b$.

Set

\Centerline{$V_k
=\frak A_{e_k}\cap\coint{f_k(a),\infty}\in\RO^{\uparrow}(Q_k)$}

\noindent for every $k$.   Then
$\pi_kV_k=\interior\overline{f_k^{-1}[V_k]}$ belongs to
$\frak G$ for each $k$, and
$G=\inf_{k\in\Bbb N}\pi_kV_k=\bigcap_{k\in\Bbb N}\pi_kV_k$
(514M(d-ii)) belongs to $\frak G$.
Because every $f_k$ is order-preserving,
$f_k(b')\Bsupseteq f_k(a)$ and $f_k(b')\in V_k$ for every
$b'\Bsupseteq b$;  thus
$b\in\interior f_k^{-1}[V_k]$ for every $k$, and $b\in G$.   This shows
that $G\ne\emptyset$.

\Quer\ Suppose, if possible, that $G\not\subseteq\coint{a,\infty}$.
Then there is a $c\in G$ such that $a\Bsetminus c\ne 0$.   If
$\bar\mu(c\Bcup a)>\gamma$, set $c'=c$.   Otherwise, let
$\delta>0$ be such that

\Centerline{$\delta+\bar\mu c<\gamma<\delta+\bar\mu(c\Bcup a)
<\bar\mu 1$.}

\noindent Because $(\frak A,\bar\mu)$ is atomless and semi-finite,
there is a $d\Bsubseteq 1\Bsetminus(c\Bcup a)$ such that
$\bar\mu d=\delta$.   Set $c'=c\Bcup d$;  then $c\Bsubseteq c'\in P$ so
$c'\in G$, while $\bar\mu(c'\Bcup a)>\gamma$, as in the previous case.

Because $I$ is infinite, $\sup_{k\in I}e_k=1$ and
there is a $k\in I$
such that $\bar\mu((c'\Bcup a)\Bcap e_k)\ge\gamma$.   In this case,
$c'\in\pi_kV_k\subseteq\overline{f_k^{-1}[V_k]}$, so
$\coint{c',\infty}$ meets $f_k^{-1}[V_k]$ and there is a
$c''\Bsupseteq c'$ such that $c''\in P$ and $f_k(c'')\in V_k$, that is,
$f_k(c'')\Bsupseteq f_k(a)$.   Now, however,

\Centerline{$f_k(c'')\Bsupseteq(c'\Bcap e_k)\Bcup(f_k(a)\Bcap e_k)
\Bsupseteq(c'\Bcup a)\Bcap e_k$}

\noindent has measure at least $\gamma$, and cannot belong to
$Q_k$.\ \Bang\Qed


\medskip

\quad{\bf (iii)} Since $\{\coint{a,\infty}:a\in P\}$ is a base for the
topology of $P$, it is a $\pi$-base for $\RO^{\uparrow}(P)$, and
$\frak G$ includes a $\pi$-base  for $\RO^{\uparrow}(P)$.   But this means
that every member of $\RO^{\uparrow}(P)$ is the supremum of the members of
$\frak G$ it includes, and belongs to $\frak G$.   Thus
$\frak G=\RO^{\uparrow}(P)$, as claimed.

\medskip

{\bf (c)(i)} This time, let
$\family{\alpha}{[0,\bar\mu 1]}{c_{\alpha}}$ be a family in $\frak A$ such
that $c_{\alpha}\Bsubseteq c_{\beta}$ and $\bar\mu c_{\alpha}=\alpha$
whenever $0\le\alpha\le\beta\le\bar\mu 1$.
Set $P=\{(a,\alpha):a\in\frak A$, $\alpha\in\ocint{\bar\mu a,\bar\mu 1}\}$.
Let $\sequencen{\gamma_n}$ be a strictly increasing sequence with
supremum $\bar\mu 1$ and $\gamma_0=0$.   For each $n\in\Bbb N$, set
$P_n=\{(a,\alpha):\gamma_n\le\bar\mu a<\alpha\le\gamma_{n+1}\}$ and
$Q_n=\{a:a\in\frak A$, $\bar\mu a<\gamma_{n+1}\}$, so that $P_n$ is an
up-open set in $P$.
Note that $\bigcup_{n\in\Bbb N}P_n$ is dense in $P$ for the up-topology,
since if $(a,\alpha)\in P$ then $(a,\min(\alpha,\gamma_{n+1}))\in P_n$
where $\gamma_n\le\bar\mu a<\gamma_{n+1}$.   Also

\Centerline{$\RO^{\uparrow}(Q_n)
=\AM(\frak A,\bar\mu,\gamma_{n+1})\cong\AM(\frak A,\bar\mu,\gamma)$.}

\noindent\Prf\ If $\bar\mu 1=\infty$, this is 528Db.   If
$(\frak A,\bar\mu)$ is totally finite, then $\frak A$ is homogeneous, so
we can apply 528Da to an appropriate multiple of the measure $\bar\mu$.\
\Qed

For $a\in\frak A^f$, the function
$\alpha\mapsto\bar\mu(c_{\alpha}\Bsetminus a):[0,\bar\mu 1]\to[0,\infty]$
is continuous and non-decreasing, and

\Centerline{$\bar\mu(c_{\bar\mu 1}\Bsetminus a)
\ge\bar\mu c_{\bar\mu 1}-\bar\mu a
=\bar\mu 1-\bar\mu a
=\bar\mu(1\Bsetminus a)$.}

\noindent So we can define $\delta(a,\alpha)$, for
$a\in\frak A^f$ and $0\le\alpha\le\bar\mu(1\Bsetminus a)$, by saying that

\Centerline{$\delta(a,\alpha)
=\min\{\beta:\bar\mu(c_{\beta}\Bsetminus a)=\alpha\}
=\min\{\beta:\bar\mu(a\Bcup c_{\beta})=\bar\mu a+\alpha\}$.}

\noindent As in (a), $\delta(a,\alpha)\le\delta(a',\alpha')$ whenever
$a\Bsubseteq a'$ and $\alpha\le\alpha'$.    For $(a,\alpha)\in P_n$, set

\Centerline{$f_n(a,\alpha)=a\Bcup c_{\delta(a,\gamma_{n+1}-\alpha)}$,}

\noindent so that
$\bar\mu f_n(a,\alpha)=\bar\mu a+\gamma_{n+1}-\alpha<\gamma_{n+1}$ and
$f_n(a,\alpha)\in Q_n$.   Of course $f_n(a,\gamma_{n+1})=a$ if
$(a,\gamma_{n+1})\in P_n$, that is, if $a\in Q_n$ and
$\bar\mu a\ge\gamma_n$.

\medskip

\quad{\bf (ii)}\grheada\ $f_n:P_n\to Q_n$ is order-preserving.   \Prf\ If
$(a,\alpha)\le(a',\alpha')$ in $P_n$, then
$\delta(a,\gamma_{n+1}-\alpha)\le\delta(a',\gamma_{n+1}-\alpha')$, so
$f_n(a,\alpha)\Bsubseteq f_n(a',\alpha')$.\ \Qed

\medskip

\qquad\grheadb\ If $p\in P_n$, $b\in Q_n$ and $f_n(p)\Bsubseteq b$,
there is a
$p'\in P_n$ such that $p\le p'$ and $b\Bsubseteq f_n(p')$.   \Prf\
Express $p$ as $(a,\alpha)$.   Consider
$a'=a\Bcup(b\Bsetminus f_n(p))$.   Then

\Centerline{$\bar\mu a'=\bar\mu a+\bar\mu b-\bar\mu f_n(p)
=\bar\mu a+\bar\mu b-\bar\mu a-\gamma_{n+1}+\alpha<\alpha$,}

\noindent so $(a',\alpha)\in P$.   Of course
$(a,\alpha)\le(a',\alpha)$, so $p'=(a',\alpha)\in P_n$.   Also
$f_n(p')\Bsupseteq f_n(p)$ and

\Centerline{$f_n(p')\Bsupseteq a'\Bsupseteq b\Bsetminus f_n(p)$,}

\noindent so $b\Bsubseteq f_n(p')$.\ \Qed

\medskip

\qquad\grheadc\ $f_n[P_n]$ is cofinal with $Q_n$.   \Prf\ If
$b\in Q_n$, take $b'\in\frak A$ such that $b\Bsubseteq b'$ and
$\gamma_n\le\bar\mu b'<\gamma_{n+1}$.   Then $(b',\gamma_{n+1})\in P_n$
and

\Centerline{$b\Bsubseteq b'=f_n(b',\gamma_{n+1})\in f_n[P_n]$.  \Qed}

\medskip

\quad{\bf (iii)} By 514P, $\RO^{\uparrow}(Q_n)$ can be regularly embedded
in
$\RO^{\uparrow}(P_n)$.   Now $\AM(\frak A,\bar\mu,\gamma)$ is isomorphic
to $\RO^{\uparrow}(Q_n)$, so there is an injective
order-continuous Boolean homomorphism
$\pi_n:\AM(\frak A,\bar\mu,\gamma)\to\RO^{\uparrow}(P_n)$.
Putting these together, we have an injective order-continuous Boolean
homomorphism $\pi:\AM(\frak A,\bar\mu,\gamma)
\to\prod_{n\in\Bbb N}\RO^{\uparrow}(P_n)$ defined by setting
$\pi u=\sequencen{\pi_n(u)}$ for $u\in\AM(\frak A,\bar\mu,\gamma)$.
On the other hand, since $\sequencen{P_n}$ is a disjoint sequence of
up-open subsets of $P$ with dense union,

\Centerline{$\prod_{n\in\Bbb N}\RO^{\uparrow}(P_n)
\cong\RO^{\uparrow}(P)=\AM^*(\frak A,\bar\mu)$}

\noindent by 315H.   So we have a regular embedding of
$\AM(\frak A,\bar\mu,\gamma)$ into $\AM^*(\frak A,\bar\mu)$, as claimed.
}%end of proof of 528F

\leader{528G}{Proposition} Let $(\frak A,\bar\mu)$ be a
measure algebra, and $\frak C$ a $\sigma$-subalgebra of $\frak A$
such that $\sup(\frak C\cap\frak A^f)=1$ in $\frak A$.   Then
$\AM^*(\frak C,\bar\mu\restrp\frak C)$ can be regularly embedded in
$\AM^*(\frak A,\bar\mu)$.

\proof{{\bf (a)} For each $a\in\frak A^f$ we have a `conditional
expectation' $u_a\in L^1(\frak C)$ defined by saying that
$\int_cu_a=\bar\mu(a\Bcap c)$ for every $c\in\frak C^f$.   (Apply
365O\formerly{3{}65P} to
the identity map from $\frak C^f$ to $\frak A^f$.)   Note that as
the supremum of $\frak C^f$ in $\frak A$ is $1$,

\Centerline{$\int u_a=\sup_{c\in\frak C^f}\int_cu_a
=\sup_{c\in\frak C^f}\bar\mu(a\Bcap c)=\bar\mu a$.}

\noindent Also, of course,
$0\le\bar\mu(a\Bcap c)\le\bar\mu c$ for every $c\in\frak C^f$, so
$0\le u_a\le\chi 1$ in $L^{\infty}(\frak C)$.   Next, let
$u_a^*$ be the decreasing rearrangement of $u_a$, that is, the element of
$L^{\infty}(\frak A_L)$ (where $\frak A_L$ is the measure algebra of
Lebesgue measure on $\coint{0,\infty}$) such that
$\Bvalue{u^*>\alpha}=\coint{0,\bar\mu\Bvalue{u>\alpha}}^{\ssbullet}$
for every $\alpha\ge 0$ (373Da).

\medskip

{\bf (b)} Set

\Centerline{$P=\{(a,\alpha):a\in\frak A$,
$\alpha\in\ocint{\bar\mu a,\bar\mu 1}\}$,
\quad$Q=\{(c,\alpha):c\in\frak C$,
$\alpha\in\ocint{\bar\mu c,\bar\mu 1}\}$.}

\noindent
Define a function $f$ on $P$ by saying that $f(a,\alpha)=(c,\beta)$ if

\Centerline{$c=\Bvalue{u_a=1}=\max\{d:d\in\frak C$, $d\Bsubseteq a\}$,}

\Centerline{$\beta
=\max\{\beta':\beta'\ge 0$,
  $\beta'+\int_{\beta'}^{\infty}u^*_a\le\alpha\}$.}

\noindent Note that $\beta>\bar\mu c$ because

\Centerline{$\bar\mu c+\int_{\bar\mu c}^{\infty}u^*_a
=\int u^*_a=\int u_a<\alpha$,}

\noindent using 373Fa for the equality in the middle, while
$\beta\le\alpha\le\bar\mu 1$;  so $(c,\beta)$ belongs to $Q$.

\medskip

{\bf (c)(i)} If $p\le p'$ in $P$,
then $f(p)\le f(p')$ in $Q$.   \Prf\ Express $p$, $p'$, $f(p)$ and $f(p')$
as $(a,\alpha)$, $(a',\alpha')$, $(c,\beta)$, $(c',\beta')$ respectively.
Then $c\Bsubseteq a\Bsubseteq a'$ so $c\Bsubseteq c'$.   Next,
$\chi a\le\chi a'$ so $u_a\le u_{a'}$ and $u^*_a\le u^*_{a'}$
(373Db);  accordingly

\Centerline{$\alpha'\le\alpha=\beta+\int_{\beta}^{\infty}u^*_a
\le\beta+\int_{\beta}^{\infty}u^*_{a'}$}

\noindent and $\beta'\le\beta$.\ \Qed

\medskip

\quad{\bf (ii)}
If $p\in P$, $q\in Q$ and $f(p)\le q$, then there is a $p'\ge p$ such that
$f(p')\ge q$.   \Prf\ Express $p$, $f(p)$ and $q$ as $(a,\alpha)$,
$(c,\beta)$ and $(d,\gamma)$ respectively.   Set $a'=a\Bcup d$.   Then

$$\eqalignno{\bar\mu a'
&=\bar\mu a+\bar\mu d-\bar\mu(a\Bcap d)
=\int u_a+\bar\mu d-\int_du_a
\ge\int u_a^*+\bar\mu d-\int_0^{\bar\mu d}u_a^*\cr
\displaycause{apply 373E with $v=\chi d$}
&=\bar\mu d+\int_{\bar\mu d}^{\infty}u_a^*
<\beta+\int_{\beta}^{\infty}u_a^*\cr
\displaycause{because $\Bvalue{u_a^*=1}=[0,\bar\mu c]^{\ssbullet}$
and $\bar\mu c\le\bar\mu d<\gamma\le\beta$}
&=\alpha.\cr}$$

\noindent So $(a',\alpha)\in P$.   Next, computing the integrals
$\int_bu_a\vee\chi d$ for $b$ belonging to $\frak C^f$
and either included in $d$ or
disjoint from it, we see that $u_{a'}=u_a\vee\chi d$ so that
$\Bvalue{u_{a'}=1}=\Bvalue{u_a=1}\Bcup d=d$.   Accordingly

\Centerline{$\bar\mu a'=\int u_{a'}=\int u^*_{a'}
=\bar\mu d+\int_{\bar\mu d}^{\infty}u^*_{a'}
<\gamma+\int_{\gamma}^{\infty}u^*_{a'}$}

\noindent (as noted above for $u_a$, we have
$\Bvalue{u_{a'}^*=1}=[0,\bar\mu d]^{\ssbullet}$), and if we set
$\alpha'=\min(\alpha,\gamma+\int_{\gamma}^{\infty}u^*_{a'})$ then
$p'=(a',\alpha')\ge p$ and $f(p')\ge q$, as required.\ \Qed

\medskip

\quad{\bf (iii)}
Since $P$ and $Q$ have a common least element $(0,\bar\mu 1)$ which is
invariant under $f$, $f$ satisfies the second condition of 514P and
$\AM^*(\frak C,\bar\mu\restrp\frak C)=\RO^{\uparrow}(Q)$ is
regularly embedded in $\AM^*(\frak A,\bar\mu)=\RO^{\uparrow}(P)$.
}%end of proof of 528G

\leader{528H}{Proposition} Let $(\frak A,\bar\mu)$ be a
semi-finite measure algebra, not $\{0\}$,
and let $\kappa\ge\max(\omega,\tau(\frak A),c(\frak A))$ be a cardinal.
Then $\AM^*(\frak A,\bar\mu)$ can be regularly embedded in
$\AM(\frak B_{\kappa},\bar\nu_{\kappa},\bover12)$.

\proof{{\bf (a)} To begin with (down to the end of (g) below), assume that
$\frak A$ is atomless.
Let $(\frak A^{\Bbb N},\bar\mu_{\infty})$ be the simple
product of a sequence of copies of $(\frak A,\bar\mu)$
(322L), so that
$\bar\mu_{\infty}\pmb{a}=\sum_{n=0}^{\infty}\mu a_n$ if
$\pmb{a}=\sequencen{a_n}\in\frak A^{\Bbb N}$.   Note that as $\frak A$
is certainly infinite, $\tau(\frak A^{\Bbb N})=\tau(\frak A)$ and
$c(\frak A^{\Bbb N})=c(\frak A)$ (514Ef).   By 526D, there is a function
$\theta:\frak A^{\Bbb N}\to\frak B_{\kappa}$ such that

\inset{$\theta(\sup A)=\sup\theta[A]$ for every non-empty
$A\subseteq\frak A^{\Bbb N}$ with a supremum in $\frak A$,

$\bar\nu_{\kappa}\theta(\pmb{a})=1-\exp(-\bar\mu_{\infty}\pmb{a})$
for every $\pmb{a}\in\frak A^{\Bbb N}$,

whenever $\familyiI{\pmb{a}^{(i)}}$ is a disjoint family in
$\frak A^{\Bbb N}$ and $\frak C_i$ is the closed subalgebra of
$\frak B_{\kappa}$
generated by $\{\theta(\pmb{a}):\pmb{a}\Bsubseteq\pmb{a}^{(i)}\}$ for
each $i$, then $\familyiI{\frak C_i}$ is stochastically independent.
}

\medskip


{\bf (b)} For $b\in\frak B_{\kappa}$, set
$g(b)=\sup\{\pmb{a}:\pmb{a}\in\frak A^{\Bbb N}$,
$\theta(\pmb{a})\Bsubseteq b\}$.

\medskip

\quad{\bf (i)} It is immediate from its definition that
$g:\frak B_{\kappa}\to\frak A^{\Bbb N}$ is order-preserving.

\medskip

\quad{\bf (ii)} Because $\theta$ is
supremum-preserving, $\theta(g(b))\Bsubseteq b$ for every $b\in\frak B_{\kappa}$.

\wheader{528H}{6}{2}{2}{36pt}

\quad{\bf (iii)} If $b\in\frak B_{\kappa}\setminus\{1\}$ then

\Centerline{$1-\bar\nu_{\kappa} b\le 1-\bar\nu_{\kappa}\theta(g(b))
=\exp(-\bar\mu_{\infty}g(b))$,}

\noindent so $\bar\mu_{\infty}g(b)\le-\ln(1-\bar\nu_{\kappa} b)$ is finite.

\medskip

\quad{\bf (iv)} $\pmb{a}\subseteq g(\theta(\pmb{a}))$ for every
$\pmb{a}\in\frak A^{\Bbb N}$;  and if
$\pmb{a}\in\frak A^{\Bbb N}$ has finite measure then
$g(\theta(\pmb{a}))=\pmb{a}$, because if $\pmb{a}'\notBsubseteq\pmb{a}$
then
$\bar\nu_{\kappa}\theta(\pmb{a}\Bcup\pmb{a}')
  >\bar\nu_{\kappa}\theta(\pmb{a})$.

\medskip

\quad{\bf (v)} If
$b\in\frak B_{\kappa}$ and $\epsilon>0$, there is a $\delta>0$ such
that $\bar\mu_{\infty}(g(b')\Bsetminus g(b))\le\epsilon$ whenever
$\bar\nu_{\kappa}(b'\Bsetminus b)\le\delta$.   \Prf\Quer\
Otherwise, $g(b)\ne 1_{\frak A^{\Bbb N}}$ so $b\ne 1_{\frak B_{\kappa}}$
and we can find a sequence $\sequencen{b_n}$ in $\frak B_{\kappa}$
such that
$\bar\nu_{\kappa}(b_n\Bsetminus b)\le 2^{-n-2}(1-\bar\nu_{\kappa}b)$ and
$\bar\mu_{\infty}(g(b_n)\Bsetminus g(b))\ge\epsilon$ for every
$n\in\Bbb N$.   For each $n$, set
$b^*_n=b\Bcup\sup_{m\ge n}b_m$;  then

\Centerline{$\bar\mu_{\infty}g(b^*_n)
=\bar\mu_{\infty}g(b)+\bar\mu_{\infty}(g(b^*_n)\Bsetminus g(b))
\ge\bar\mu_{\infty}g(b)+\bar\mu_{\infty}(g(b_n)\Bsetminus g(b))
\ge\epsilon+\bar\mu_{\infty}g(b)$.}

\noindent Note that $\bar\nu_{\kappa}b^*_0<1$ so $g(b^*_0)$ has finite
measure.

The sequences $\sequencen{b^*_n}$,
$\sequencen{g(b^*_n)}$ and $\sequencen{\theta(g(b^*_n))}$
are all non-increasing.   Set $\pmb{a}=\inf_{n\in\Bbb N}g(b^*_n)$, so that

\Centerline{$\theta(\pmb{a})
\Bsubseteq\inf_{n\in\Bbb N}\theta(g(b^*_n))
\Bsubseteq\inf_{n\in\Bbb N}b^*_n=b$}

\noindent because $\bar\nu_{\kappa}(b^*_n\Bsetminus b)\le 2^{-n-1}$ for
every $n$.   It follows that $\pmb{a}\Bsubseteq g(b)$.   At the same time,

\Centerline{$\bar\mu_{\infty}\pmb{a}
=\lim_{n\to\infty}\bar\mu_{\infty}g(b^*_n)
>\bar\mu_{\infty}g(b)$,}

\noindent which is impossible.\ \Bang\Qed

\medskip

{\bf (c)} Define $\psi:\frak A^{\Bbb N}\to\frak A$ by setting
$\psi(\pmb{a})=\sup_{n\in\Bbb N}a_n$ whenever
$\pmb{a}=\sequencen{a_n}\in\frak A^{\Bbb N}$.

\medskip

\quad{\bf (i)} $\psi$ is supremum-preserving and $\psi(0)=0$.

\medskip

\quad{\bf (ii)} If $\pmb{a}$, $\pmb{a}'\in\frak A^{\Bbb N}$ then

\Centerline{$\bar\mu(\psi(\pmb{a})\Bsymmdiff\psi(\pmb{a}'))
\le\bar\mu_{\infty}(\pmb{a}\Bsymmdiff\pmb{a}')$
\quad$\bar\mu(\psi(\pmb{a})\Bsetminus\psi(\pmb{a}'))
\le\bar\mu_{\infty}(\pmb{a}\Bsetminus\pmb{a}')$.}

\medskip

\quad{\bf (iii)} Now if $b\in\frak B_{\kappa}$, $a\in\frak A$,
$a\Bsupseteq\psi(g(b))$ and
$\bar\mu a<\alpha\in\Bbb R$, there is a $b'\Bsupseteq b$ such that
$a\Bsubseteq\psi(g(b'))$, $\bar\mu\psi(g(b'))<\alpha$ and
$\bar\nu_{\kappa}(b'\setminus b)\le 1-\exp(-\bar\mu(a\Bsetminus\psi(g(b)))$.

\Prf\ Take $\alpha'$ such that $\bar\mu a<\alpha'<\alpha$.
By (b-v), there is a $\delta>0$ such that
$\bar\mu_{\infty}(g(b')\Bsetminus g(b))\le\alpha'-\bar\mu a$ whenever
$\bar\nu_{\kappa}(b'\Bsetminus b)\le\delta$.
Set $\pmb{a}^{(n)}=\sequence{i}{a_{ni}}$ for each
$n\in\Bbb N$, where $a_{ni}=a\setminus\psi(g(b))$ if $i=n$, $0$
otherwise.   For each $n\in\Bbb N$, let $\frak C_n$ be the closed
subalgebra of $\frak B_{\kappa}$ generated by

\Centerline{$\{\theta(\pmb{a}):\pmb{a}\in\frak A^{\Bbb N}$,
$\pmb{a}\Bcap\pmb{a}^{(m)}=0$ for every $m\ge n\}$,}

\noindent and let $T_n:L^1(\frak B_{\kappa},\bar\nu_{\kappa})
\to L^1(\frak B_{\kappa},\bar\nu_{\kappa})$ be
the corresponding conditional-expectation operator (365Q\formerly{3{}65R}).
Then $\sequencen{\frak C_n}$ is non-decreasing;  also
$\bar\nu_{\kappa}\theta(\pmb{a}^{(n)}\Bcap c)
=\bar\nu_{\kappa}\theta(\pmb{a}^{(n)})\cdot\bar\nu_{\kappa} c$ for every $n\in\Bbb N$ and
$c\in\frak C_n$, by the final clause of (a).   By L\'evy's martingale
theorem (275I, 367Jb), $\sequencen{T_n(\chi b)}$ is
$\|\,\|_1$-convergent.   We can therefore find an $n\in\Bbb N$ such that
$\|T_n(\chi b)-T_{n+1}(\chi b)\|_1\le\delta\exp(-\bar\mu a)$.
Set $b'=b\Bcup\theta(\pmb{a}^{(n)})$.   Then
$g(b')\Bsupseteq g(b)\Bcup\pmb{a}^{(n)}$, so
$\psi(g(b'))\Bsupseteq a_{nn}\Bcup\psi(g(b))=a$.  Also

\Centerline{$\bar\nu_{\kappa}(b'\Bsetminus b)
\le\bar\nu_{\kappa}\theta(\pmb{a}^{(n)})
=1-\exp(-\bar\mu_{\infty}\pmb{a}^{(n)})
=1-\exp(-\bar\mu(a\Bsetminus\psi(g(b)))$.}

\Quer\ If $\bar\mu\psi(g(b'))\ge\alpha$, set
$\pmb{e}=g(b')\Bsetminus(g(b)\Bcup\sup_{m\in\Bbb N}\pmb{a}^{(m)})$.
Since

\Centerline{$\psi(g(b)\Bcup\sup_{m\in\Bbb N}\pmb{a}^{(m)})
=\psi(g(b))\Bcup a=a$,}

\noindent $\psi(\pmb{e})\Bsupseteq\psi(g(b'))\Bsetminus a$  and

\Centerline{$\bar\mu_{\infty}\pmb{e}\ge\alpha-\bar\mu a
>\alpha'-\bar\mu a$;}

\noindent as
$\pmb{e}\Bsubseteq g(\theta(\pmb{e}))$ and $\pmb{e}\Bcap g(b)=0$,
$\bar\nu_{\kappa}(\theta(\pmb{e})\Bsetminus b)>\delta$.
On the other hand,

$$\eqalignno{(1-\bar\nu_{\kappa}\theta(\pmb{a}^{(n)}))
  \bar\nu_{\kappa}(b\Bcap\theta(\pmb{e}))
&=(1-\bar\nu_{\kappa}\theta(\pmb{a}^{(n)}))
  \int_{\theta(\pmb{e})}T_n(\chi b)\cr
\displaycause{because $\pmb{e}\Bcap\pmb{a}^{(m)}=0$ for every $m$, so
$\theta(\pmb{e})\in\frak C_n$}
&=\int\chi(1\Bsetminus\theta(\pmb{a}^{(n)}))\cdot
   \int T_n(\chi b)\times\chi\theta(\pmb{e})\cr
&=\int\chi(1\Bsetminus\theta(\pmb{a}^{(n)}))
  \times T_n(\chi b)\times\chi\theta(\pmb{e})\cr
\displaycause{because
$T_n(\chi b)\times\chi\theta(\pmb{e})\in L^0(\frak C_n)$ and
$\chi(1\Bsetminus\theta(\pmb{a}^{(n)}))$ are stochastically independent}
&=\int_{\theta(\pmb{e})\Bsetminus\theta(\pmb{a}^{(n)})}T_n(\chi b),\cr
\cr
(1-\bar\nu_{\kappa}\theta(\pmb{a}^{(n)}))\bar\nu_{\kappa}(\theta(\pmb{e}))
&=\bar\nu_{\kappa}(\theta(\pmb{e})\Bsetminus\theta(\pmb{a}^{(n)})
=\bar\nu_{\kappa}(b\Bcap\theta(\pmb{e})\Bsetminus\theta(\pmb{a}^{(n)})\cr
\displaycause{because
$\theta(\pmb{e})\Bsubseteq\theta(g(b'))\Bsubseteq b'
=b\Bcup\theta(\pmb{a}^{(n)})$}
&=\int_{\theta(\pmb{e})\Bsetminus\theta(\pmb{a}^{(n)})}
  T_{n+1}(\chi b)\cr}$$

\noindent because $\theta(\pmb{e})$ and $\theta(\pmb{a}^{(n)})$ both
belong to $\frak C_{n+1}$.   So

$$\eqalign{\delta\exp(-\bar\mu a)
&=\delta\exp(-\bar\mu_{\infty}\pmb{a}^{(n)})
<\bar\nu_{\kappa}(\theta(\pmb{e})\Bsetminus b)
  \exp(-\bar\mu_{\infty}\pmb{a}^{(n)})\cr
&=(1-\bar\nu_{\kappa}\theta(\pmb{a}^{(n)}))
  \bar\nu_{\kappa}(\theta(\pmb{e})\Bsetminus b))\cr
&=\int_{\theta(\pmb{e})\Bsetminus\theta(\pmb{a}^{(n)})}
  T_{n+1}(\chi b)-T_n(\chi b)\cr
&\le\|T_n(\chi b)-T_{n+1}(\chi b)\|_1
\le\delta\exp(-\bar\mu a),\cr}$$

\noindent which is impossible.\ \Bang

So $\bar\mu\psi(g(b'))<\alpha$, as required.\ \Qed

\medskip

{\bf (d)} Fix $c\in\frak B_{\kappa}$ with measure $\bover12$;
then the principal ideal of
$\frak B_{\kappa}$ generated by $c$ is isomorphic to
$\frak B_{\kappa}$ with the measure halved.   We therefore have a Boolean
isomorphism $\pi:\frak B_{\kappa}\to(\frak B_{\kappa})_c$ such that
$\bar\nu_{\kappa}\pi b=\bover12\bar\nu_{\kappa}b$ for every
$b\in\frak B_{\kappa}$.   Set
$h(b)=\psi(g(\pi^{-1}(b\Bcap c)))$ for $b\in\frak B_{\kappa}$.   Then
$h:\frak B_{\kappa}\to\frak A$ is order-preserving and $h(b)=h(b\Bcap c)$
for
every $b\in\frak B_{\kappa}$.   Translating the results of (b) and (c), we see
that

\inset{if $b\in\frak B_{\kappa}$ and $\epsilon>0$, there is a
$\delta>0$ such that $\bar\mu(h(b')\Bsetminus h(b))\le\epsilon$ whenever
$\bar\nu_{\kappa}(b'\Bsetminus b)\le\delta$,

if $b\in\frak B_{\kappa}$, $a\in\frak A$, $a\Bsupseteq h(b)$ and
$\bar\mu a<\alpha\in\Bbb R$, there is a $b'\Bsupseteq b$ such that
$a\Bsubseteq h(b')$, $\bar\mu h(b')<\alpha$ and
$\bar\nu_{\kappa}(b'\Bsetminus b)
\le\bover12(1-\exp(-\bar\mu(a\Bsetminus h(b)))$.}

\noindent  Note also that

$$\eqalign{\bar\mu h(b)
&=\bar\mu\psi(g(\pi^{-1}(b\Bcap c)))
\le\bar\mu_{\infty}g(\pi^{-1}(b\Bcap c))\cr
&\le-\ln(1-\bar\nu_{\kappa}\pi^{-1}(b\Bcap c))
\le-\ln(1-2\bar\nu_{\kappa}(b\Bcap c))\cr}$$

\noindent if we take $\ln(0)$ to be $-\infty$.

\medskip

{\bf (e)(i)} Set $\gamma_0=\bover12(1-\exp(-\bar\mu 1))$, interpreting
$\exp(-\infty)$ as $0$, so that $0<\gamma_0\le\bover12$.
Let $P$ be the partially ordered set
$\{(a,\alpha):a\in\frak A$, $\alpha\in\ocint{\bar\mu a,\bar\mu 1}\}$
and $Q$ the partially ordered set
$\{b:b\in\frak B_{\kappa}$, $\bar\nu_{\kappa} b<\gamma_0\}$, so that
$\AM^*(\frak A,\bar\mu)=\RO^{\uparrow}(P)$ and
$\AM(\frak B_{\kappa},\bar\nu_{\kappa},\gamma_0)=\RO^{\uparrow}(Q)$.
For $b\in Q$, set
$\alpha_b=\sup\{\bar\mu h(b'):b\Bsubseteq b'\in Q\}$.
Then $\alpha_b>\bar\mu(h(b))$.   \Prf\ We have

\Centerline{$\bar\mu h(b)\le-\ln(1-2\bar\nu_{\kappa}(b\Bcap c))
<-\ln(1-2\gamma_0)=\bar\mu 1$,}

\noindent so $h(b)\ne 1$.   Because $\frak A$ is atomless,
there is an $a\in\frak A$,
disjoint from $h(b)$, such that $0<\bar\mu a<-\ln(2\bar\nu_{\kappa}b)$.
Set $\pmb{a}=\sequencen{a_n}$ where $a_0=a$, $a_n=0$ for $n\ge 1$.   Then

\Centerline{$\bar\nu_{\kappa}\theta(\pmb{a})
=1-\exp(-\bar\mu a)<1-2\bar\nu_{\kappa}b$,}

\noindent so $b'=b\Bcup\pi\theta(\pmb{a})\in Q$, while
$h(b')\Bsupseteq h(b)\Bcup a\Bsupset h(b)$.\ \Qed

\medskip

\quad{\bf (ii)} If $b\in Q$ and $\bar\mu h(b)<\alpha$, there is
a $b_1\in Q$ such that $b\Bsubseteq b_1$, $h(b_1)=h(b)$ and
$\alpha_{b'}\le\alpha$.   \Prf\ Let $\delta>0$ be such that
$\bar\mu(h(b')\Bsetminus h(b))\le\alpha-\bar\mu h(b)$ whenever
$\bar\nu_{\kappa}(b'\Bsetminus b)\le\delta$.   Because
$\gamma_0\le\bover12$, there is a $b_1\in Q$ such that
$b\Bsubseteq b_1$, $b\Bcap c=b_1\Bcap c$ and
$\bar\mu b_1\ge\gamma_0-\delta$.   Then $h(b_1)=h(b)$.   If $b'\in Q$
and $b'\Bsupseteq b_1$, then
$\bar\nu_{\kappa}(b'\Bcap c\Bsetminus b)\le\delta$, so

\Centerline{$\bar\mu(h(b')\Bsetminus h(b))
=\bar\mu(h(b'\Bcap c)\Bsetminus h(b))\le\alpha-\bar\mu h(b)$}

\noindent and $\bar\mu h(b')\le\alpha$;  thus $\alpha_{b_2}\le\alpha$.\
\Qed

\medskip

{\bf (f)} By (e-i), we can define $f:Q\to P$ by setting
$f(b)=(h(b),\alpha_b)$ for $b\in Q$.

\medskip

\quad{\bf (i)} $f$ is order-preserving because $h$ is.

\medskip

\quad{\bf (ii)} If $P_1\subseteq P$ is up-open and cofinal with $P$,
$f^{-1}[P_1]$ is cofinal with $Q$.   \Prf\ Take any $b\in Q$.   Set

\Centerline{$\alpha
=\min(\alpha_b,\bar\mu h(b)-\ln(1-2\gamma_0+2\bar\nu_{\kappa}b))
>\bar\mu h(b)$,}

\noindent so that $f(b)\le(h(b),\alpha)$ in $P$.   Then
there is an $(a,\beta)\in P_1$ such that $(h(b),\alpha)\le(a,\beta)$,
that is, $h(b)\Bsubseteq a$ and $\beta\le\alpha$.   In this case, there
is a $b_1\in\frak B_{\kappa}$ such that
$b_1\Bsupseteq b$, $h(b_1)\Bsupseteq a$, $\bar\mu h(b_1)<\beta$ and

\Centerline{$\bar\nu_{\kappa}(b_1\Bsetminus b)
\le\Bover12(1-\exp(-\bar\mu(a\Bsetminus h(b))))
<\gamma_0-\bar\nu b$}

\noindent because $\bar\mu(a\Bsetminus h(b))
<-\ln(1-2\gamma_0+2\bar\nu_{\kappa}b)$.   So
$b_1\in Q$.   By (e-ii), there is a $b_2\in Q$ such that
$b_2\Bsupseteq b_1$, $h(b_2)=h(b_1)$ and $\alpha_{b_2}\le\beta$.
Now $b\Bsubseteq b_2$, while $f(b_2)=(h(b_1),\alpha_{b_2})\ge(a,\beta)$.
As $P_1$ is up-open, $f(b_2)\in P_1$;  as $b$ is
arbitrary, $f^{-1}[P_1]$ is cofinal with $Q$.\ \Qed

\medskip

\quad{\bf (iii)} $f[Q]$ is cofinal with $P$.
\Prf\ Take $(a,\alpha)\in P$.
Set $\pmb{a}=\sequencen{a_n}\in\frak A^{\Bbb N}$ where $a_0=a$ and
$a_n=0$ for $n\ge 1$, and set $b=\pi\theta(\pmb{a})$.   Note that

\Centerline{$\bar\nu_{\kappa}b=\Bover12(1-\exp(-\bar\mu a))<\gamma_0$,}

\noindent so $b\in Q$.   Because
$\bar\mu_{\infty}\pmb{a}<\infty$, $g(\pi^{-1}b)=\pmb{a}$ and $h(b)=a$.
By (e-ii) again, we can now find a
$b_1\Bsupseteq b$ in $Q$ such that $h(b_1)=h(b)$ and
$\alpha_{b_1}\le\alpha$.   So $f(b_1)\ge(a,\alpha)$.   As $(a,\alpha)$
is arbitrary, $f[Q]$ is cofinal.\ \Qed

\medskip

{\bf (g)} By 514O, $\AM^*(\frak A,\bar\mu)=\RO^{\uparrow}(P)$ can be
regularly embedded in

\Centerline{$\RO^{\uparrow}(Q)
=\AM(\frak B_{\kappa},\bar\nu_{\kappa},\gamma_0)
\cong\AM(\frak B_{\kappa},\nu_{\kappa},\bover12)$}

\noindent by 528Da.

\medskip

{\bf (h)} All this has been done on the assumption that $\frak A$ is
atomless, as required in (e).   For the general case, consider the
localizable measure algebra free product
$(\frak C,\bar\lambda)$ of $(\frak A,\bar\mu)$ and
$(\frak B_{\omega},\bar\nu_{\omega})$ (325E).   By 521Qa, we have

\Centerline{$\max(\omega,c(\frak C),\tau(\frak C))
\le\max(\omega,c(\frak A),c(\frak B_{\omega}),\tau(\frak A),
  \tau(\frak B_{\omega}))
\le\kappa$.}

\noindent Also $\frak C$ is atomless because $\frak B_{\omega}$ is
isomorphic to a closed subalgebra of $\frak C$ (325Dd) and is atomless
(316R\formerly{3{}16Xi}).
By (a)-(g), $\AM^*(\frak C,\bar\lambda)$ can be regularly embedded in
$\AM(\frak B_{\omega},\bar\nu_{\omega},\bover12)$.   Now consider the
canonical embedding $\varepsilon_1:\frak A\to\frak C$.
This is order-continuous and measure-preserving (325Da), so identifies the
Dedekind $\sigma$-complete Boolean algebra $\frak A$ with a
$\sigma$-subalgebra of $\frak C$;  also $\frak A^f$ has supremum $1$ both
in $\frak A$ and $\frak C$.   By 528G, $\AM^*(\frak A,\bar\mu)$ can be
regularly embedded in $\AM^*(\frak C,\bar\lambda)$ and therefore in
$\AM(\frak B_{\omega},\bar\nu_{\omega},\bover12)$.
}%end of proof of 528H

\leader{528I}{Definition} For any set $I$, the
{\bf ($I,\infty)$-localization poset} is the set

\Centerline{$\Cal S_I^{\infty}
=\{p:p\subseteq\Bbb N\times I$, $\#(p[\{n\}])\le 2^n$
for every $n$, $\sup_{n\in\Bbb N}\#(p[\{n\}])$ is finite$\}$,}

\noindent ordered by $\subseteq$.   For $p\in\Cal S_I^{\infty}$ set
$\|p\|=\max_{n\in\Bbb N}\#(p[\{n\}])$.    I will write $\Cal S^{\infty}$
for $\Cal S^{\infty}_{\Bbb N}$\cmmnt{,
already introduced in the proof of 522T}.

\leader{528J}{Proposition} Let $\kappa$ be an infinite cardinal,
$\Cal S_{\kappa}^{\infty}$ the $(\kappa,\infty)$-localization poset,
and $(\frak A,\bar\mu)$ a
semi-finite measure algebra, not $\{0\}$, with
$\kappa\ge\max(\omega,c(\frak A),\tau(\frak A))$.
Then the variable-measure amoeba algebra
$\AM^*(\frak A,\bar\mu)$ can be regularly embedded in
$\RO^{\uparrow}(\Cal S_{\kappa}^{\infty})$.

\proof{{\bf (a)} To begin with (down to the end of (d) below), suppose that
$\frak A$ is atomless.   Let $P$ be the partially ordered set
$\{(a,\alpha):a\in\frak A$, $\alpha\in\ocint{\bar\mu a,\mu 1}\}$, so that
$\AM^*(\frak A,\bar\mu)=\RO^{\uparrow}(P)$.   Give
$\frak A^f$ its measure metric, so that its topological density is at
most $\kappa$ (521Eb).
Set $\gamma_0=\bover12\bar\mu 1$ and for $n\ge 1$ set
$\gamma_n=2^{-2n-1}\bar\mu 1$ if $\bar\mu 1<\infty$, $4^{-n}$ otherwise.
For each $n$, let $D_n$ be a dense subset of
$\{a:a\in\frak A^f$, $\bar\mu a\le\gamma_n\}$, containing $0$, with cardinal at
most $\kappa$, and let
$\ofamily{\xi}{\kappa}{d_{n\xi}}$ be a family running over $D_n$ with
cofinal repetitions.

\medskip

{\bf (b)} If $p\in\Cal S_{\kappa}^{\infty}$ set

\Centerline{$a_p=\sup_{(n,\xi)\in p}d_{n\xi}$,
\quad$\alpha_p
=\bar\mu a_p+\sum_{n=0}^{\infty}(2^n-\#(p[\{n\}]))\gamma_n$.}

\noindent Then

\Centerline{$\bar\mu a_p<\alpha_p
\le\sum_{n=0}^{\infty}2^n\gamma_n=\bar\mu 1$,}

\noindent so we can define $f:\Cal S_{\kappa}^{\infty}\to P$ by setting
$f(p)=(a_p,\alpha_p)$.    $f$ is order-preserving, because if
$p\subseteq p'$ in $\Cal S_{\kappa}^{\infty}$ then

$$\eqalign{\alpha_{p'}
&=\bar\mu a_{p'}+\sum_{n=0}^{\infty}(2^n-\#(p'[\{n\}]))\gamma_n\cr
&\le\bar\mu a_p+\sum_{(n,\xi)\in p'\setminus p}\bar\mu d_{n\xi}
  +\sum_{n=0}^{\infty}(2^n-\#(p'[\{n\}]))\gamma_n\cr
&\le\bar\mu a_p
  +\sum_{n=0}^{\infty}\#(p'[\{n\}]\setminus p[\{n\}])\gamma_n
  +\sum_{n=0}^{\infty}(2^n-\#(p'[\{n\}]))\gamma_n
=\alpha_p.\cr}$$

\medskip

{\bf (c)} Suppose that $p\in\Cal S_{\kappa}^{\infty}$ and
$f(p)\le(a,\alpha)\in P$.   Take $\alpha'\in\ooint{\bar\mu a,\alpha}$.
For $n\in\Bbb N$, set $k_n=2^n-\#(p[\{n\}])$.   Then $a_p\Bsubseteq a$ and

\Centerline{$\bar\mu a
<\alpha\le\alpha_p=\bar\mu a_p+\sum_{n=0}^{\infty}k_n\gamma_n$.}

\noindent So there is an $r\in\Bbb N$ such that

\Centerline{$\bar\mu a
<\bar\mu a_p+\sum_{n=0}^{\infty}\gamma_n\min(r,k_n)$;}

\noindent take $r$ so large that, in addition,
$\sum_{n=r+1}^{\infty}2^n\gamma_n\le\alpha-\alpha'$.

For each $n$, set $k'_n=\min(r,k_n)$ and
$C_n=\{\sup D:D\in[D_n]^{\le k'_n}\}$.
Then (because $\frak A$ is atomless) $C_n$ is dense in
$\{c:c\in\frak A^f$, $\bar\mu c\le k'_n\gamma_n\}$.   We can therefore
choose $\sequencen{c_n}$ inductively in such a way that

\Centerline{$c_n\in C_n$,
\quad$\bar\mu(a\Bcup\sup_{m<n}c_m)<\alpha'$,}

\Centerline{$\bar\mu(a\Bsetminus(a_p\Bcup\sup_{m<n}c_m))
<\sum_{m=n}^{\infty}k'_m\gamma_m$}

\noindent for every $n\in\Bbb N$.   \Prf\
For the inductive step to $n\ge 0$, set
$b=a\Bsetminus(a_p\Bcup\sup_{m<n}c_m)$.   Take $b'\Bsubseteq b$ such that
$\bar\mu b'=\min(k'_n\gamma_n,\bar\mu b)$, so that

\Centerline{$\bar\mu(a\Bcup\sup_{m<n}c_m\Bcup b')
=\bar\mu(a\Bcup\sup_{m<n}c_m)<\alpha'$,}

\Centerline{$\bar\mu(b\setminus b')<\sum_{m=n+1}^{\infty}k'_m\gamma_m$.}

\noindent Let $c_n\in C_n$ be such that

\Centerline{$\alpha'
>\bar\mu(a\Bcup\sup_{m<n}c_m\Bcup b')+\bar\mu(c_n\Bsetminus b')
\ge\bar\mu(a\Bcup\sup_{m\le n}c_m)$,}

\Centerline{$\sum_{m=n+1}^{\infty}k'_m\gamma_m
>\bar\mu(b\setminus b')+\bar\mu(b'\Bsetminus c_n)
=\bar\mu(a\Bsetminus(a_p\Bcup\sup_{m\le n}c_m))$}

\noindent and the induction proceeds.\ \Qed

For each $n$, we can find a set $D'_n\subseteq D_n$, with cardinal{} $k'_n$,
such that $c_n=\sup D'_n$.   Because $\ofamily{\xi}{\kappa}{d_{n\xi}}$
runs over $D_n$ with cofinal repetitions, we can find a set
$I_n\subseteq\kappa\setminus p[\{n\}]$ such that $\#(I_n)=k'_n$ and
$c_n=\sup_{\xi\in I_n}d_{n\xi}$.   Set
$q=p\cup\{(n,\xi):n\in\Bbb N$, $\xi\in I_n\}$.   Then

\Centerline{$\#(q[\{n\}])\le\#(p[\{n\}])+k'_n\le\min(2^n,\|p\|+r)$}

\noindent for every $n$, so $q\in\Cal S_{\kappa}^{\infty}$ and
$p\subseteq q$.    Now

\Centerline{$a_q=a_p\Bcup\sup_{n\in\Bbb N,\xi\in I_n}d_{n\xi}
=a_p\Bcup\sup_{n\in\Bbb N}c_n\Bsupseteq a$}

\noindent because

\Centerline{$\bar\mu(a\Bsetminus a_q)
\le\inf_{n\in\Bbb N}\bar\mu(a\Bsetminus(a_p\Bcup\sup_{m<n}c_m))
\le\inf_{n\in\Bbb N}\sum_{m=n}^{\infty}2^m\gamma_m
=0$.}

\noindent Also

\Centerline{$\bar\mu a_q
=\sup_{n\in\Bbb N}\bar\mu(a_p\Bcup\sup_{m<n}c_m)
\le\alpha'<\alpha$,}

\noindent while $\#(q[\{n\}])=\#(p[\{n\}])+k_n=2^n$ whenever $n\le r$,
so

\Centerline{$\alpha_q
=\bar\mu a_q+\sum_{n=r+1}^{\infty}(2^n-\#(q[\{n\}]))\gamma_n
\le\alpha'+\sum_{n=r+1}^{\infty}2^n\gamma_n
\le\alpha$.}

\medskip

{\bf (d)} What (c) shows is that if $p\in\Cal S_{\kappa}^{\infty}$ and
$f(p)\le(a,\alpha)$ in $P$, then there is a $q\supseteq p$ in
$\Cal S_{\kappa}^{\infty}$ such that $(a,\alpha)\le f(q)$.
Next, $\Cal S^{\infty}_{\kappa}$ has a least element $\emptyset$, and
$f(\emptyset)=(0,\bar\mu 1)$ is the least element of $P$.   So 514P tells
us that
$\RO^{\uparrow}(P)=\AM^*(\frak A,\bar\mu)$ can be regularly embedded in
$\RO^{\uparrow}(\Cal S_{\kappa}^{\infty})$.

\medskip

{\bf (e)} As for the general case, we can use the same trick as in
part (h) of the proof of 528H.   Let $(\frak C,\bar\lambda)$ be the
localizable measure algebra free product of $(\frak A,\bar\mu)$ and
$(\frak B_{\omega},\bar\nu_{\omega})$;  as before, $\frak C$ is atomless,
$\max(\omega,c(\frak C),\tau(\frak C))\le\kappa$ and $(\frak A,\bar\mu)$
is embedded in $(\frak C,\bar\lambda)$ as a $\sigma$-subalgebra with
sufficient elements of finite measure.   So
$\AM^*(\frak A,\bar\mu)$ is regularly embedded in
$\AM^*(\frak C,\bar\lambda)$ and in
$\RO^{\uparrow}(\Cal S_{\kappa}^{\infty})$.
}%end of proof of 528J

\leader{528K}{Theorem}\cmmnt{ ({\smc Truss 88})} Let
$(\frak A,\bar\mu)$ be an atomless $\sigma$-finite measure algebra in which
every non-zero principal ideal has Maharam type $\kappa$,
and $0<\gamma<\bar\mu 1$.   Then each of the algebras

\Centerline{$\AM(\frak A,\bar\mu,\gamma)$,
\quad$\AM^*(\frak A,\bar\mu)$,
\quad$\AM(\frak B_{\kappa},\bar\nu_{\kappa},\bover12)$}

\noindent can be regularly embedded in the other two, and all three can
be regularly embedded in $\RO^{\uparrow}(\Cal S_{\kappa}^{\infty})$.

\proof{ By 528H,
$\AM^*(\frak A,\bar\mu)$ can be regularly embedded in
$\AM(\frak B_{\kappa},\bar\nu_{\kappa},\bover12)$.   Take any
$e\in\frak A$ such that $\gamma<\bar\mu e<\infty$.   Then the principal
ideal $(\frak A_e,\bar\mu\restrp\frak A_e)$ is isomorphic, up to a
scalar multiple of the measure, to
$(\frak B_{\kappa},\bar\nu_{\kappa})$, so

$$\eqalignno{\AM(\frak B_{\kappa},\bar\nu_{\kappa},\textstyle{\bover12})
&\cong\AM(\frak B_{\kappa},\bar\nu_{\kappa},
  \Bover{\gamma}{\bar\mu e})\cr
\displaycause{528Da}
&\cong\AM(\frak A_e,\bar\mu\restrp\frak A_e,\gamma)\cr}$$

\noindent can be regularly embedded in $\AM(\frak A,\bar\mu,\gamma)$
(528Fa).   By 528Fc,
$\AM(\frak A,\bar\mu,\gamma)$ can be regularly embedded in
$\AM^*(\frak A,\bar\mu)$.   Finally, by 528J, $\AM^*(\frak A,\bar\mu)$
can be regularly embedded in $\RO^{\uparrow}(\Cal S_{\kappa}^{\infty})$.
Because regular embeddability is transitive (313N), these facts are
enough to prove the theorem.
}%end of proof of 528K

\leader{528L}{}\cmmnt{ It is possible without great effort to
calculate many of the cardinal functions of these algebras.

\medskip

\noindent}{\bf Lemma}
$\frak m(\AM(\frak B_{\omega},\bar\nu_{\omega},\bover12))\le\add\Cal N$,
where $\Cal N$ is the null ideal of Lebesgue measure on $\Bbb R$.

\proof{ Set $P=\{a:a\in\frak B_{\omega}$, $\bar\nu_{\omega}a<\bover12\}$.
Then $\wdistr(\frak B_{\omega})\ge\frak m^{\uparrow}(P)$.   \Prf\ Take a
family $\ofamily{\xi}{\kappa}{B_{\xi}}$ of maximal antichains in
$\frak B_{\omega}$, where
$\kappa<\frak m^{\uparrow}(P)$.   Let $C\subseteq\frak B_{\omega}$ be a
maximal disjoint set such that $\{b:b\in B_{\xi}$, $b\Bcap c\ne 0\}$ is
finite for every $\xi<\kappa$ and $c\in C$.   \Quer\ Suppose, if
possible, that $c_0=1\Bsetminus\sup C$ is not $0$.   Take $a_0\in P$
such that $\bar\nu_{\omega}(a_0\Bcup c_0)>\bover12$.
(If $\bar\nu_{\omega}c_0>\bover12$, take $a_0=0$;
otherwise, take $a_0\Bsubseteq 1\Bsetminus c_0$ such that
$\bover12-\bar\nu_{\omega}c_0<\bar\nu_{\omega}a_0<\bover12$.)
For each $\xi<\kappa$, set

\Centerline{$Q_{\xi}
=\{a:a\in P$, $\{b:b\in B_{\xi}$, $b\notBsubseteq a\}$ is finite$\}$;}

\noindent then $Q_{\xi}$ is cofinal with $P$.   There is therefore an
upwards-directed $R\subseteq P$ such that $a_0\in R$ and $R$ meets every
$Q_{\xi}$.   Set $e=\sup R$;  then $\bar\nu_{\omega}e\le\bover12$ so
$c_1=c_0\setminus e=(a_0\Bcup c_0)\setminus e$ is non-zero.

If $\xi<\kappa$, there is an $a\in R\cap Q_{\xi}$, so that

\Centerline{$\{b:b\in B_{\xi}$, $b\Bcap c_1\ne 0\}
\subseteq\{b:b\in B_{\xi}$, $b\notBsubseteq a\}$}

\noindent is finite.   But this means that we ought to have added $c_1$
to $C$.\ \Bang

Thus $C$ is a maximal antichain.   As $\ofamily{\xi}{\kappa}{B_{\xi}}$
is arbitrary, $\wdistr(\frak A)\ge\frak m^{\uparrow}(P)$.\ \Qed

Now 524Mb tells us that $\wdistr(\frak B_{\omega})=\add\Cal N$, so
$\frak m^{\uparrow}(P)\le\add\Cal N$.   Finally, by 517Db,

\Centerline{$\frak m(\AM(\frak A,\bar\mu,\gamma))=\frak m^{\uparrow}(P)
\le\add\Cal N$,}

\noindent as claimed.
}%end of proof of 528L

\leader{528M}{Lemma}
$\frak m^{\uparrow}(\Cal S^{\infty})\ge\add\Cal N$.

\proof{{\bf (a)} Recall the definition of the supported relations
$(\NN,\subseteq^*,\Cal S^{(\alpha)}$) from 522L, where
$\Cal S^{(\alpha)}=\{S:S\subseteq\Bbb N\times\Bbb N$,
$\#(S[\{n\}])\le\alpha(n)$ for every $n\in\Bbb N\}$ for $\alpha\in\NN$.
Putting 522L, %indep of \alpha
522M %OK for usual \alpha
and 512Db %\add(C,S,D)\le\add(A,R,B)
together, we have
$\add(\NN,\subseteq^*,\Cal S^{(\alpha)})=\add\Cal N$ whenever
$\lim_{n\to\infty}\alpha(n)=\infty$.

\medskip

{\bf (b)} The core of the argument is the following fact.   Suppose that
$Q\subseteq\Cal S^{\infty}$ is cofinal and up-open, $n\in\Bbb N$ and
$\sigma\in[\Bbb N\times\Bbb N]^{<\omega}$.   Let
$G\subseteq\Bbb N\times\Bbb N$ be a set with finite vertical sections.
Then there is a $k\in\Bbb N$ such that whenever
$\sigma\subseteq p\in\Cal S^{\infty}$, $p\subseteq\sigma\cup G$ and
$\|p\|\le n$,
there is a $q\in Q$ such that $p\subseteq q$ and $\|q\|\le k$.

\Prf\Quer\ Suppose, if possible, otherwise.   Then for each $j\in\Bbb N$
we can find $p_j\in\Cal S^{\infty}$ such that
$\sigma\subseteq p_j\subseteq\sigma\cup G$, $\|p_j\|\le n$ and $\|q\|>j$
whenever $p\subseteq q\in Q$.   Let $p$ be a cluster point of
$\sequence{j}{p_j}$ in $\Cal P(\sigma\cup G)$.   Then
$\#(p[\{i\}])\le\sup_{j\in\Bbb N}\#(p_j[\{i\}])\le\min(2^i,n)$ for every
$i$, so $p\in\Cal S^{\infty}$.  Because $Q$ is cofinal with $\Cal
S^{\infty}$, there is a $q\in Q$
such that $p\subseteq q$.   Set $k=n+\|q\|$.   Then
$(\sigma\cup G)\cap(k\times\Bbb N)$ is finite, so there is an $i\ge k$
such that $p_i\cap(k\times\Bbb N)=p\cap(k\times\Bbb N)\subseteq q$.
Set $q'=p_i\cup q$.   Then

$$\eqalign{\#(q'[\{j\}])
&=\#(q[\{j\}])\le\min(\|q\|,2^j)\text{ if }j<k,\cr
&\le\|p_i\|+\|q\|\le k\le 2^j\text{ otherwise}.\cr}$$

\noindent So $q'\in\Cal S^{\infty}$ and $\|q'\|\le k\le i$;  because $Q$
is up-open in $\Cal S^{\infty}$, $q'\in Q$, while $p_i\subseteq q'$.
But we
chose $p_i$ so that this could not happen.\ \Bang\Qed

\medskip

{\bf (c)} We need to know that $\Cal S^{\infty}$ is upwards-ccc.   \Prf\
For any $n\in\Bbb N$, finite $\sigma\subseteq\Bbb N\times\Bbb N$ the set
$\{p:p\in\Cal S^{\infty}$, $\|p\|\le 2^{n-1}$,
$p\cap(n\times\Bbb N)=\sigma\}$ is upwards-linked.\ \Qed

\medskip

{\bf (d)} Now let $\ofamily{\xi}{\kappa}{Q_{\xi}}$ be any family of
cofinal subsets of $\Cal S^{\infty}$, where $\kappa<\add\Cal N$, and
$p_0\in\Cal S^{\infty}$.   For each $\xi<\kappa$ let $A_{\xi}\subseteq
Q_{\xi}$ be a
maximal up-antichain;  by (c), $A_{\xi}$ is countable.   Set
$Q'_{\xi}=\bigcup\{\coint{q,\infty}:q\in A_{\xi}\}$, so that $Q'_{\xi}$
is an up-open cofinal subset of $\Cal S^{\infty}$.   Set
$A=\{p_0\}\cup\bigcup_{\xi<\kappa}A_{\xi}$.   For $q\in A$, let
$F_q\subseteq\NN$ be a finite set such that (identifying each member of
$F_q$ with its graph) $q\subseteq\bigcup F_q$;
set $F=\bigcup_{q\in A}F_q$, so that

\Centerline{$\#(F)\le\max(\omega,\kappa)<\add\Cal N\le\frak b$}

\noindent (522B).   Let $g_0\in\NN$ be a strictly increasing function
such that $\{i:f(i)>g_0(i)\}$ is finite for every $f\in F$, and also
$p_0[\{i\}]\subseteq g_0(i)$ for every $i$.   Set
$G=\{(i,j):i\in\Bbb N$, $j<g_0(i)\}$, so that
$G\subseteq\Bbb N\times\Bbb N$ has finite vertical sections.
Observe that if $q\in A$ then $q\setminus G$ is finite.

For each $\xi<\kappa$, $n\in\Bbb N$ and finite
$\sigma\subseteq\Bbb N\times\Bbb N$, let $k(\xi,\sigma,n)\in\Bbb N$ be
such that whenever $p\in\Cal S^{\infty}$ and
$\sigma\subseteq p\subseteq\sigma\cup G$
then there is a $q\in Q'_{\xi}$ such that $p\subseteq q$ and
$\|q\|\le k(\xi,\sigma,n)$;  such a $k$ exists by (b) above.   Set
$k_{\xi}(n)=\sup\{k(\xi,\sigma,n):\sigma\subseteq n\times g_0(n)\}$.
Again because $\kappa<\frak b$, there is a $g_1\in\NN$ such that
$\{n:k_{\xi}(n)>g_1(n)\}$ is finite for every $\xi<\kappa$.   Let
$\alpha\in\NN$ be a non-decreasing function
such that $\lim_{n\to\infty}\alpha(n)=\infty$ and

\Centerline{$\alpha(2g_1(n))\le n$,
\quad$\alpha(n)+\#(p_0[\{n\}])\le 2^n$,
\quad$2\alpha(n)\le n$}

\noindent for every $n$.

Because $\add(\NN,\subseteq^*,\Cal S^{(\alpha)})=\add\Cal N$, there is an
$S_0\in\Cal S^{(\alpha)}$ such that $f\subseteq^*S_0$
for every $f\in F$, so that $q\setminus S_0$ is finite for every
$q\in A$.   Replacing $S_0$ by $S_0\cap G$ if necessary, we may suppose
that $S_0\subseteq G$.

\medskip

{\bf (e)} Let $\Cal S$ be the family of subsets $S$ of $\Bbb N\times\Bbb N$
such that $\#(S[\{n\}])\le 2^n$ for every $n$, as in 522K.
Note that $p_0\cup S_0\in\Cal S$,
because $\alpha(n)+\#(p_0(n))\le 2^n$ for every $n$.
Let $C$ be the family of finite subsets $\sigma$ of
$\Bbb N\times\Bbb N$ such that $\sigma\cup S_0\in\Cal S$.
For each $\xi<\kappa$, set

\Centerline{$D_{\xi}
=\{\sigma:\sigma\in C$, $\Exists q\in A_{\xi}$,
  $q\subseteq\sigma\cup S_0\}$.}

\noindent Then $D_{\xi}$ is cofinal with $C$.   \Prf\ Let $\sigma\in C$.
Let $n_0$ be so large that $g_1(2^{n_0})\ge k_{\xi}(2^{n_0})$ and
$\sigma\subseteq n_0\times g_0(n_0)$.   Set $m=2g_1(2^{n_0})$,
$p=\sigma\cup(S_0\cap(m\times\Bbb N))\in\Cal S^{\infty}$.   Then
$\sigma\subseteq p\subseteq\sigma\cup G$ and
$\|p\|\le\max(2^{n_0},\alpha(m))=2^{n_0}$, so there is a $q\in Q'_{\xi}$
such that $p\subseteq q$ and

\Centerline{$\|q\|\le k(\xi,\sigma,2^{n_0})\le k_{\xi}(2^{n_0})
\le g_1(2^{n_0})=\Bover{m}{2}$.}

\noindent Let $q'\in A_{\xi}$ be such that $q'\subseteq q$.   Let
$m'\ge\max(m,n_0)$ be such that $q'\subseteq(m'\times\Bbb N)\cup S_0$,
and set $\tau=q\cap(m'\times\Bbb N)$, so that $\sigma\subseteq\tau$.

For $n<m$, we have

\Centerline{$S_0[\{n\}]\subseteq p[\{n\}]\subseteq q[\{n\}]
=\tau[\{n\}]$,}

\noindent so $(\tau\cup S_0)[\{n\}]=q[\{n\}]$ has at most $2^n$ members.
For $m\le n<m'$, we have

\Centerline{$\#((\tau\cup S_0)[\{n\}])\le\#(q[\{n\}])+\#(S_0[\{n\}])
\le\|q\|+\alpha(n)\le\Bover{m}2+\Bover{n}2\le 2^n$,}

\noindent while for $n\ge m'$ we have

\Centerline{$\#((\tau\cup S_0)[\{n\}])
=\#(S_0[\{n\}])\le\alpha(n)\le 2^n$.}

\noindent So $\tau\cup S_0\in\Cal S$ and $\tau\in C$.   Since

\Centerline{$q'\subseteq(q'\cap(m'\times\Bbb N))\cup S_0
\subseteq(q\cap(m'\times\Bbb N))\cup S_0=\tau\cup S_0$,}

\noindent $\tau\in D_{\xi}$.   As $\sigma$ is arbitrary, $D_{\xi}$ is
cofinal with $C$.\ \Qed

\medskip

{\bf (f)} Because $p_0\cup S_0\in\Cal S$,
$\sigma_0=p_0\setminus S_0$ belongs to $C$.
Because $\kappa<\add\Cal N\le\frakmctbl\le\frak m^{\uparrow}(C)$, there
is an upwards-directed set $E\subseteq C$ meeting every $D_{\xi}$ and
containing $\sigma_0$.   Set $S_1=S_0\cup\bigcup E$.   Then, because $E$
is upwards-directed,

\Centerline{$\#(S_1[\{n\}])
=\sup_{\sigma\in E}\#((\sigma\cup S_0)[\{n\}])\le 2^n$}

\noindent for every $n$, and $S_1\in\Cal S$.   Set
$R=\{p:p\in\Cal S^{\infty}$, $p\subseteq S_1\}$;  then
$R\subseteq\Cal S^{\infty}$ is
upwards-directed (in fact, closed under $\cup$), and $p_0\in R$ because
$\sigma_0\in E$.   Now $R$ meets $Q_{\xi}$ for every $\xi<\kappa$.
\Prf\ There is a
$\sigma\in D_{\xi}\cap E$.   But this means that there is a
$q\in A_{\xi}$ such that
$q\subseteq\sigma\cup S_0\subseteq S_1$ and $q\in R\cap Q_{\xi}$.\ \Qed

As $p_0$ and $\ofamily{\xi}{\kappa}{Q_{\xi}}$ are arbitrary,
$\frak m^{\uparrow}(\Cal S^{\infty})\ge\add\Cal N$.
}%end of proof of 528M

\leader{528N}{Theorem}\cmmnt{ ({\smc Brendle 00}, 2.3.10;
{\smc Judah \& Repick\'y 95})}
Let $(\frak A,\bar\mu)$ be an atomless $\sigma$-finite measure algebra
with countable Maharam type, and $0<\gamma<\bar\mu 1$.   Then the
algebras $\AM(\frak A,\bar\mu,\gamma)$ and $\AM^*(\frak A,\bar\mu)$
and the $(\Bbb N,\infty)$-localization poset $\Cal S^{\infty}$ (active upwards)
all have Martin numbers equal to $\add\Cal N$.

\proof{ By 517Ia and 528K, with 517Db again,

$$\eqalign{\frak m(\AM(\frak A,\bar\mu,\gamma))
&=\frak m(\text{AM}^*(\frak A,\bar\mu))
=\frak m(\AM(\frak B_{\omega},\bar\nu_{\omega},\textstyle{\bover12}))\cr
&\ge\frak m(\RO^{\uparrow}(\Cal S^{\infty}))
=\frak m^{\uparrow}(\Cal S^{\infty}).\cr}$$

\noindent As

\Centerline{$\frak m(\AM(\frak B_{\omega},\bar\nu_{\omega},\bover12))
\le\add\Cal N\le\frak m^{\uparrow}(\Cal S^{\infty})$}

\noindent (528L, 528M), all these are equal to $\add\Cal N$.
}%end of proof of 528N

\leader{528O}{Corollary} Let $\gamma>0$.   Let $\Cal G$ be the
partially ordered set

\Centerline{$\{G:G\subseteq\Bbb R$ is open, $\mu_LG<\gamma\}$,}

\noindent where $\mu_L$ is Lebesgue measure.   Then
$\frak m^{\uparrow}(\Cal G)=\add\Cal N$.

\proof{ Put 528C and 528N together.
}%end of proof of 528O

\vleader{72pt}{528P}{Proposition}
Let $(\frak A,\bar\mu)$ be an atomless
semi-finite measure algebra, and $0<\gamma<\bar\mu 1$.

(a) For any integer $m\ge 2$,

\Centerline{$c(\AM(\frak A,\bar\mu,\gamma))
=\link_m(\AM(\frak A,\bar\mu,\gamma))=\max(c(\frak A),\tau(\frak A))$.}

(b) $d(\AM(\frak A,\bar\mu,\gamma))
=\pi(\AM(\frak A,\bar\mu,\gamma))
=\max(\cff[c(\frak A)]^{\le\omega},\pi(\frak A))$.

%what about \tau  \wdistr
% I think \tau(\AM)  related to  \link\frak A  but do we have  = ?

\proof{ Set $P=\{a:a\in\frak A$, $\bar\mu a<\gamma\}$, so that
$\AM(\frak A,\bar\mu,\gamma)=\RO^{\uparrow}(P)$.

\medskip

{\bf (a)} Set $\kappa_0=\max(c(\frak A),\tau(\frak A))$,
$\kappa_1=\link_m(\RO^{\uparrow}(P))=\link_m^{\uparrow}(P)$
and $\kappa_2=c(\RO^{\uparrow}(P))=c^{\uparrow}(P)$
(514N).

\medskip

\quad{\bf (i)} The topological density of $\frak A^f$ for its
measure metric is $\kappa_0$ (521Eb), so $P$ has a metrically dense
subset $D$ with cardinal at most $\kappa_0$.   For $d\in D$, set

\Centerline{$U_d
=\{a:a\in P$, $\bar\mu(a\Bsetminus d)<\Bover1m(\gamma-\bar\mu d)\}$.}

\noindent Then $U_d$ is upwards-$m$-linked in $P$.   Also, if $a\in P$,
there is a $d\in D$ such that
$\bar\mu(a\Bsymmdiff d)<\bover1{m+1}(\gamma-\bar\mu a)$, and now
$a\in U_d$.   So $P$ is $\kappa_0$-$m$-linked upwards and
$\kappa_1\le\kappa_0$.

\medskip

\quad{\bf (ii)} By 511Hb or 511Ia, $\kappa_2\le\kappa_1$.

\medskip

\quad{\bf (iii)} We need to check that $\kappa_2$ is infinite.   \Prf\
Take $a\in\frak A$ such that $\bar\mu a=\gamma$.   For any $n\ge 1$, we
can find disjoint $a_0,\ldots,a_n\Bsubseteq a$ all of measure
$\bover1{n+1}\gamma$;  now $\langle a\Bsetminus a_i\rangle_{i\le n}$ is
an up-antichain in $P$.   So $\kappa_2=c^{\uparrow}(P)\ge n+1$;  and
this is true for every $n$.\ \Qed

Now if $(\frak A,\bar\mu)$ is totally finite, then
$c(\frak A)=\omega\le\kappa_2$.
Otherwise, there is a partition $D$ of unity in $\frak A$ such that
$\bar\mu d=\bover12\gamma$ for every $d\in D$;  now $D$ is an
up-antichain in $P$ and $\kappa_2\ge\#(D)=c(\frak A)$.   So we see that
in all cases $\kappa_2\ge c(\frak A)$.

\medskip

\quad{\bf (iv)} If $e\in\frak A^f$ and the principal ideal $\frak A_e$
is homogeneous, then $\tau(\frak A_e)\le\kappa_2$.   \Prf\Quer\
Otherwise,
set $\alpha=\bar\mu e$, $\kappa=\tau(\frak A_e)$.   Because
$\bar\mu 1>\gamma$, there is a $d\Bsubseteq 1\Bsetminus e$ such that
$\gamma<\bar\mu(e\Bcup d)<\gamma+\bar\mu e$, that is,
$0<\gamma-\bar\mu d<\alpha$.   Set
$\beta=\sqrt{1-\Bover{\gamma-\bar\mu d}{\alpha}}$.   Because
$\frak A_e$ is isomorphic, up to a scalar multiple of the measure, to
the measure algebra of the usual measure on $[0,1]^{\kappa}$, there is a
family $\ofamily{\xi}{\kappa}{c_{\xi}}$ in $\frak A_e$ such that

\Centerline{$\bar\mu c_{\xi}=\beta\alpha$,
\quad$\bar\mu(c_{\xi}\Bcap c_{\eta})=\beta^2\alpha$}

\noindent whenever $\xi$, $\eta<\kappa$ are distinct.   Set
$b_{\xi}=d\Bcup(e\Bsetminus c_{\xi})$ for $\xi<\kappa$.   Then

\Centerline{$\bar\mu(b_{\xi}\Bcup b_{\eta})
=\bar\mu d+\alpha-\beta^2\alpha=\gamma$,}

\Centerline{$\bar\mu b_{\xi}=\bar\mu d+\alpha-\beta\alpha<\gamma$}

\noindent for all distinct $\xi$, $\eta<\kappa$.   So
$\ofamily{\xi}{\kappa}{b_{\xi}}$ is an up-antichain in $P$ and
witnesses that $\kappa_2\ge\kappa$.\ \Bang\Qed

\medskip

\quad{\bf (v)} Let $E$ be a partition of unity in $\frak A$ such that
$0<\bar\mu e<\infty$ and $\frak A_e$ is homogeneous for every $e\in E$.
For $e\in E$, let $A_e\Bsubseteq\frak A_e$ be a set with cardinal at most
$\kappa_2$ which $\tau$-generates $\frak A_e$.   Then
$A=\bigcup_{e\in E}A_e\,\,\tau$-generates $\frak A$, so that

\Centerline{$\kappa_0
=\max(c(\frak A),\tau(\frak A))\le\max(c(\frak A),\#(A))
\le\max(c(\frak A),\kappa_2)=\kappa_2$,}

\noindent and the three cardinals must be equal.

\medskip

{\bf (b)} Set $\kappa_3=\max(\pi(\frak A),\cff[c(\frak A)]^{\le\omega})$,
$\kappa_4=\pi(\AM(\frak A,\bar\mu,\gamma))$ and
$\kappa_5=d(\AM(\frak A,\bar\mu,\gamma))$.

\medskip

\quad{\bf (i)} By 323Mc, $\frak A^f$ is complete in its measure metric.
By 323Ma, $\Bcup:\frak A^f\times\frak A^f\to\frak A^f$ is uniformly
continuous for the measure metric, and $P$ is an open subset of
$\frak A^f$, while the topological density of $\frak A^f$ is $\kappa_0$.
By 524C, $(P,\Bsubseteqshort\strprime,[P]^{<\omega})
\prGT(\ell^1(\kappa_0),\le,\ell^1(\kappa_0))$, where
$\Bsubseteqshort\strprime$ is defined as in 512F.   It follows that

$$\eqalignno{\kappa_4
&=\pi(\RO^{\uparrow}(P))
=\cf P\cr
\displaycause{514Nb}
&=\cov(P,\Bsubseteqshort,P)
\le\max(\omega,\cov(P,\Bsubseteqshort\strprime,[P]^{\le\omega}))\cr
\displaycause{512Gf}
&\le\max(\omega,\cov(P,\Bsubseteqshort\strprime,[P]^{<\omega}))\cr
\displaycause{512Gb}
&\le\max(\omega,\cf\ell^1(\kappa_0))\cr
\displaycause{512Da}
&=\cf\ell^1(\kappa_0)
=\cf\Cal N_{\kappa_0}\cr
\displaycause{where $\Cal N_{\kappa_0}$ is the null ideal of the
usual measure on $\{0,1\}^{\kappa_0}$, as in 524I}
&=\max(\cf\Cal N,\cff[\kappa_0]^{\le\omega})\cr
\displaycause{523N}
&=\max(\cf\Cal N,\cff[\tau(\frak A)]^{\le\omega},
  \cff[c(\frak A)]^{\le\omega})\cr
&=\max(\cf\Cal N,\cff[\tau(\frak A)]^{\le\omega},c(\frak A),
  \cff[c(\frak A)]^{\le\omega})
=\max(\pi(\frak A),\cff[c(\frak A)]^{\le\omega})\cr
\displaycause{524Mc}
&=\kappa_3.\cr}$$

\medskip

\quad{\bf (ii)} By 514Nd, $d^{\uparrow}(P)=\kappa_5$.   Let
$\ofamily{\xi}{\kappa_5}{B_{\xi}}$ be a family of
upwards-centered sets covering $P$.   For each $\xi$,
$b_{\xi}=\sup B_{\xi}$ is defined in $\frak A$ (counting $\sup\emptyset$
as $0$ if necessary), and

\Centerline{$\bar\mu b_{\xi}=\sup_{I\in[B_{\xi}]^{<\omega}}\bar\mu(\sup I)
\le\gamma$.}

\noindent
Set $D=\{b_{\xi}\Bsetminus b_{\eta}:\xi$, $\eta<\kappa_5\}$.   Then $D$
is order-dense in $\frak A$.   \Prf\ If $a\in\frak A\setminus\{0\}$,
take $a'\Bsubseteq a$ such that
$0<\bar\mu a'<\gamma$.   Then $a'\in P$, so there is some $\xi<\kappa_5$
such that $a'\in B_{\xi}$ and
$a'\Bsubseteq b_{\xi}$.   Next, let $c\Bsubseteq 1\Bsetminus b_{\xi}$ be
such that

\Centerline{$\gamma-\bar\mu b_{\xi}<\bar\mu c
<\gamma-\bar\mu b_{\xi}+\bar\mu a'$.}

\noindent Then $c\Bcup(b_{\xi}\Bsetminus a')\in P$, so there is an
$\eta<\kappa_5$ such that
$c\Bcup(b_{\xi}\Bsetminus a')\Bsubseteq b_{\eta}$.   Now
$d=b_{\xi}\Bsetminus b_{\eta}\Bsubseteq a'$;  as
$\bar\mu(b_{\xi}\Bcup c)>\gamma\ge\bar\mu b_{\eta}$,
$b_{\xi}\notBsubseteq b_{\eta}$ and $d\ne 0$.   Of course $d\in D$ and
$d\Bsubseteq a$;  as $a$ is arbitrary, $D$ is order-dense.\ \Qed

Accordingly $\pi(\frak A)\le\#(D)\le\kappa_5$.   At the same time,
$\cff[c(\frak A)]^{\le\omega}\le\kappa_5$.  \Prf\ There is a disjoint set
$E\subseteq\frak A\setminus\{0\}$ with cardinal $c(\frak A)$ (332F).   For
each $\xi<\kappa_5$, let $I_{\xi}$ be the countable set
$\{e:e\in E$, $e\Bcap b_{\xi}\ne 0\}$.
If $J\subseteq E$ is countable, let $\family{e}{J}{\epsilon_e}$ be a
strictly positive family of real numbers with sum less than $\gamma$.
For each $e\in J$ let $a_e\Bsubseteq e$ be such that
$0<\bar\mu a_e\le\epsilon_e$, and set $a=\sup_{e\in J}a_e$.   Then $a\in
P$ so there is a $\xi<\kappa_5$ such that $a\Bsubseteq b_{\xi}$ and
$J\subseteq I_{\xi}$.   As $J$ is arbitrary, $\{I_{\xi}:\xi<\kappa_5\}$
is cofinal with $[E]^{\le\omega}$, and

\Centerline{$\cff[c(\frak A)]^{\le\omega}
=\cff[E]^{\le\omega}\le\kappa_5$.\ \Qed}

Putting these together, we see that $\kappa_3\le\kappa_5$.

\medskip

\quad{\bf (iii)} By 514Da, $\kappa_5\le\kappa_4$, so the three cardinals
are equal.
}%end of proof of 528P

\leader{528Q}{Proposition} Let $\Cal S^{\infty}$ be the
$(\Bbb N,\infty)$-localization poset.

(a)
$\pi(\RO^{\uparrow}(\Cal S^{\infty}))=\cf\Cal S^{\infty}=\frak c$.

(b) For every $m\ge 2$,

\Centerline{$c(\RO^{\uparrow}(\Cal S^{\infty}))
=c^{\uparrow}(\Cal S^{\infty})
=\link_m(\RO^{\uparrow}(\Cal S^{\infty}))
=\link_m^{\uparrow}(\Cal S^{\infty})=\omega$.}

(c) $d(\RO^{\uparrow}(\Cal S^{\infty}))
=d^{\uparrow}(\Cal S^{\infty})=\cf\Cal N$.

%what about  \wdistr, \tau

\proof{{\bf (a)} If $p$, $q\in\Cal S^{\infty}$ and $p\not\subseteq q$,
take $(n,i)\in p\setminus q$;   then there is a $q'\in\Cal S^{\infty}$
such that $q'\supseteq q$, $\#(q'[\{n\}])=2^n$ and $(n,i)\notin q'$, in
which case $p$ and $q'$ are incompatible upwards in $\Cal S^{\infty}$.
So $\Cal S^{\infty}$ is separative upwards and 514Nb tells us that
$\pi(\RO^{\uparrow}(\Cal S^{\infty}))=\cf\Cal S^{\infty}$.

Next, there is an almost-disjoint family
$\ofamily{\xi}{\frakc}{h_{\xi}}$ in $\NN$ (5A1Nc).   Identifying each
$h_{\xi}$ with its graph, we can regard them as members of
$\Cal S^{\infty}$;  and any member of $\Cal S^{\infty}$ includes only
finitely many of them.   So $\cf\Cal S^{\infty}\ge\frak c$.   On the
other hand, of course,
$\cf\Cal S^{\infty}\le\#(\Cal S^{\infty})=\frak c$.   So
$\pi(\RO^{\uparrow}(\Cal S^{\infty}))=\cf\Cal S^{\infty}=\frak c$.

\medskip

{\bf (b)} If $m\ge 2$, let $Q$ be the countable set of pairs $(I,r)$
where $r\in\Bbb N$ and
$I\in[\Bbb N\times\Bbb N]^{<\omega}$, and for $(I,r)\in Q$ set

\Centerline{$A_{Ir}
=\{p:p\in\Cal S^{\infty}$, $p\cap(r\times\Bbb N)=I$,
$\|p\|\le\Bover{2^r}{m}\}$.}

\noindent Then $\bigcup_{i<m}p_i\in\Cal S^{\infty}$ for any family
$\ofamily{i}{m}{p_i}$ in $A_{Ir}$, that is, $A_{Ir}$ is
upwards-$m$-linked in $\Cal S^{\infty}$.   Also
$\bigcup_{(I,r)\in Q}A_{Ir}=\Cal S^{\infty}$, so
$\link_m^{\uparrow}(\Cal S^{\infty})\le\omega$.
Of course $c^{\uparrow}(\Cal S^{\infty})$ is infinite, and since
$c^{\uparrow}(\Cal S^{\infty})\le\link_m^{\uparrow}(\Cal S^{\infty})$
(511Hb again), both must be $\omega$.   Now 514N tells us that

\Centerline{$c(\RO^{\uparrow}(\Cal S^{\infty}))
=\link_m(\RO^{\uparrow}(\Cal S^{\infty}))=\omega$.}

\medskip

{\bf (c)} Consider the $\Bbb N$-localization relation $(\NN,\subseteq^*,\Cal S)$
of 522K.   We know from 522M and 512Da that

\Centerline{$\cov(\NN,\subseteq^*,\Cal S)
=\cov(\Cal N,\subseteq,\Cal N)=\cf\Cal N$.}

\medskip

\quad{\bf (i)} Let $\Cal A\subseteq\Cal S$ be a set with cardinal
$\cf\Cal N$ such that for every $f\in\NN$ there is an $S\in\Cal A$ such
that $f\subseteq^*S$.   Let $\Cal A^*$ be

\Centerline{$\{S:S\in\Cal S$, $S\setminus\bigcup\Cal A'$ is finite for
some finite $\Cal A'\subseteq\Cal A\}$;}

\noindent then every member of
$\Cal S^{\infty}$ is included in some member of $\Cal A^*$.   But if
$S\in\Cal A^*$ then $\{p:p\in\Cal S^{\infty}$, $p\subseteq S\}$ is
upwards-directed.   So

\Centerline{$\duparrow(\Cal S^{\infty})\le\#(\Cal A^*)\le\cf\Cal N$.}

\medskip

\quad{\bf (ii)} Now let $\Cal Q$ be a family of upwards-centered subsets
of $\Cal S^{\infty}$ covering $\Cal S^{\infty}$.   For each
$Q\in\Cal Q$, $S_Q=\bigcup Q$ belongs to $\Cal S$.   Also every
$f\in\NN$ belongs to $\Cal S^{\infty}$ so is covered by some $S_Q$.   So
$S_Q$ witnesses that
$\cf\Cal N=\cov(\NN,\subseteq^*,\Cal S)\le\#(\Cal Q)$;  as $\Cal Q$ is
arbitrary, $\cf\Cal N\le\duparrow(\Cal S^{\infty})$.

\medskip

\quad{\bf (iii)} 514Nd tells us that

\Centerline{$d(\RO^{\uparrow}(\Cal S^{\infty}))
=\duparrow(\Cal S^{\infty})$,}

\noindent so we have equality throughout.
}%end of proof of 528Q

\leader{528R}{Theorem} Let $\kappa$ be any cardinal, and
$\Cal S^{\infty}_{\kappa}$ the $(\kappa,\infty)$-localization poset.
Then $\RO^{\uparrow}(\Cal S^{\infty}_{\kappa})$ has countable Maharam
type.

\proof{{\bf (a)} If $\kappa$ is finite then
$\cf\Cal S^{\infty}_{\kappa}$ is finite
and the result is trivial.   So let us suppose from now on that $\kappa$ is
infinite.

\medskip

{\bf (b)} $\Cal S^{\infty}_{\kappa}$ is separative upwards.
\Prf\ If $p$, $q\in\Cal S^{\infty}_{\kappa}$ and $p\not\subseteq q$, take
$(n,\xi)\in p\setminus q$.   Let $J\subseteq\kappa\setminus p[\{n\}]$ be a set
of{} size $2^n-\#(q[\{n\}]$, and set $q'=q\cup(\{n\}\times J)$;
then $q\subseteq q'\in\Cal S^{\infty}$ and $p$, $q'$ are incompatible upwards in
$\Cal S^{\infty}_{\kappa}$.\ \Qed

Accordingly $\coint{p,\infty}\in\RO^{\uparrow}(\Cal S^{\infty}_{\kappa})$
for every $p\in\Cal S^{\infty}_{\kappa}$ (514Me).

\medskip

{\bf (c)} For $n\in\Bbb N$, $m<2^n$ and $\xi<\kappa$, set

\Centerline{$G_{mn\xi}
=\sup\{\coint{p,\infty}:p\in\Cal S^{\infty}_{\kappa},\,
\#(p[\{n\}])=2^n,\,(n,\xi)\in p\text{ and }\#(p[\{n\}]\cap\xi)=m\}$,}

\noindent the supremum being taken in
$\RO^{\uparrow}(\Cal S^{\infty}_{\kappa})$.

\medskip

{\bf (d)} If $n\in\Bbb N$, $m<2^n$ and $\xi<\eta<\kappa$ then
$G_{mn\xi}\cap G_{mn\eta}=\emptyset$.   \Prf\ If
$p$, $q\in\Cal S^{\infty}_{\kappa}$,
$\#(p[\{n\}])=\#(q[\{n\}])=2^n$, $(n,\xi)\in p$, $(n,\eta)\in q$ and
$\#(p[\{n\}]\cap\xi)=\#(q[\{n\}])\cap\eta)=m$ then $p[\{n\}]\ne q[\{n\}]$,
$\#(p[\{n\}]\cup q[\{n\}])>2^n$ and $\coint{p,\infty}\cap\coint{q,\infty}$ is
empty.\ \Qed

\medskip

{\bf (e)} If $\xi<\kappa$ then
$\bigcup\{G_{mn\xi}:n\in\Bbb N$, $m<2^n\}$ is dense in
$\Cal S^{\infty}_{\kappa}$.   \Prf\ If $p\in\Cal S^{\infty}_{\kappa}$, take
$n\in\Bbb N$ such that $\#(p[\{n\}])<2^n$;  then there is a
$q\in S^{\infty}_{\kappa}$ such that $p\subseteq q$, $\xi\in q[\{n\}]$ and
$\#(q[\{n\}])=2^n$.   Set $m=\#(q[\{n\}]\cap\xi)$;  then
$\coint{q,\infty}\subseteq\coint{p,\infty}\cap G_{mn\xi}$.\ \Qed

Thus $\sup\{G_{mn\xi}:n\in\Bbb N$, $m<2^n\}=1$ in
$\RO^{\uparrow}(\Cal S^{\infty})$ whenever $\xi<\kappa$.

\medskip

{\bf (f)}
Let $\frak G$ be the order-closed subalgebra of
$\RO^{\uparrow}(\Cal S^{\infty}_{\kappa})$ generated by
$\{G_{mn\xi}:n\in\Bbb N$, $m<2^n$, $\xi<\kappa\}$.
For $n\in\Bbb N$ and $\xi<\kappa$ set
$H_{n\xi}=\coint{\{(n,\xi)\},\infty}$;
then $H_{n\xi}=\sup_{m<2^n}G_{mn\xi}$.
\Prf\ Certainly $G_{mn\xi}\subseteq H_{n\xi}$ whenever $m<2^n$.
If $\{(n,\xi)\}\subseteq p\in\Cal S^{\infty}_{\kappa}$, let
$q\in\Cal S^{\infty}_{\kappa}$ be such that $p\subseteq q$ and
$\#(q[\{n\}])=2^n$;  set $m=\#(q[\{n\}]\cap\xi)$;  then
$\coint{q,\infty}\subseteq H_{n\xi}\cap G_{mn\xi}$.   Thus
$\bigcup_{m<2^n}G_{mn\xi}$ is dense in $H_{n\xi}$ and
$H_{n\xi}=\sup_{m<2^n}G_{mn\xi}\in\frak G$.\ \QeD\  Consequently
$H_{n\xi}\in\frak G$.

\medskip

{\bf (g)} If $p\in\Cal S^{\infty}_{\kappa}$ then
$\coint{p,\infty}=\inf_{(n,\xi)\in p}H_{n\xi}$ belongs to $\frak G$,
by 514Me.
So $\frak G$ includes an order-dense subset of
$\RO^{\uparrow}(\Cal S^{\infty}_{\kappa})$ and must be the whole of
$\RO^{\uparrow}(\Cal S^{\infty}_{\kappa})$;  that is,
$\RO^{\uparrow}(\Cal S^{\infty}_{\kappa})$ is $\tau$-generated by
$\{G_{mn\xi}:n\in\Bbb N$, $m<2^n$, $\xi<\kappa\}$.
With (iv) and (v), we see
that the conditions of 514F are satisified with $J=\kappa$ and
$I=\{(m,n):n\in\Bbb N$, $m<2^n\}$, so that

\Centerline{$\tau(\RO^{\uparrow}(\Cal S^{\infty}_{\kappa}))
\le\max(\omega,\#(I))=\omega$.}
}%end of proof of 528R

\leader{528S}{}\cmmnt{ The calculation of Maharam types of
amoeba algebras seems to be a good deal harder.
However it leads through an investigation of the
structure of measure algebras, which is one of the things this book is
about, so I take the space to give one of the main theorems.   It depends
on a special property of the standard generating families in algebras
$\frak B_{\kappa}$.

\medskip

\noindent}{\bf Definition}\dvAnew{2011}
Let $(\frak A,\bar\mu)$ be a measure algebra.
I will say that a {\bf well-spread basis} for $\frak A$ is a
non-decreasing sequence $\sequencen{D_n}$ of
subsets of $\frak A$ such that

\inset{(i) setting $D=\bigcup_{n\in\Bbb N}D_n$,
$\#(D)\le\max(\omega,c(\frak A),\tau(\frak A))$;

(ii) if $a\in\frak A$, $\gamma\in\Bbb R$ and $\bar\mu a<\gamma$, there is a set
$D\subseteq\bigcup_{n\in\Bbb N}D_n$ such that
$a\Bsubseteq\sup D$ and $\bar\mu(\sup D)<\gamma$;

(iii) if $n\in\Bbb N$ and $\sequence{i}{d_i}$ is a sequence in $D_n$
such that $\bar\mu(\sup_{i\in\Bbb N}d_i)<\infty$, there
is an infinite set $J\subseteq\Bbb N$ such that $d=\sup_{i\in J}d_i$
belongs to $D_n$;

(iv) whenever $n\in\Bbb N$, $a\in\frak A$ and
$\bar\mu a\le\gamma'<\gamma<\bar\mu 1$, there is a $b\in\frak A$ such that
$a\Bsubseteq b$ and $\gamma'\le\bar\mu b<\gamma$ and
$\bar\mu(b\Bcup d)\ge\gamma$ whenever $d\in D_n$ and
$d\notBsubseteq a$.}

\leader{528T}{Lemma}\dvAnew{2011}
(a) Let $\kappa$ be an infinite cardinal, and
$\ofamily{\xi}{\kappa}{e_{\xi}}$ the standard generating family in
$\frak B_{\kappa}$.
For $n\in\Bbb N$ let $C_n$ be the set of elements of
$\frak B_{\kappa}$ expressible as
$\inf_{\xi\in I}e_{\xi}\Bcap\inf_{\xi\in J}(1\Bsetminus e_{\xi})$
where $I$,
$J\subseteq\kappa$ are disjoint and $\#(I\cup J)\le n$.   Then
$\sequencen{C_n}$ is a well-spread basis for
$(\frak B_{\kappa},\bar\nu_{\kappa})$.   Moreover,

\doubleinset{(*) for each $n\ge 1$, there is a set $C'_n\subseteq C_n$,
with cardinal $\kappa$,
such that $\bar\nu_{\kappa} c=2^{-n}$ for every $c\in C'_n$, and whenever
$a\in\frak B_{\kappa}\setminus\{1\}$ and $I\subseteq C'_n$ is infinite,
there is a $c\in I$ such that $c'\notBsubseteq a\Bcup c$ whenever
$c\Bsubset c'\in C_n$.}

(b) Let $(\frak A,\bar\mu)$ be a measure algebra and $e\in\frak A$.
If $\sequencen{C_n}$ is a well-spread basis for
$(\frak A_e,\bar\mu\restrp\frak A_e)$ and
$\sequencen{D_n}$ is a well-spread basis for
$(\frak A_{1\Bsetminus e},\bar\mu\restrp\frak A_{1\Bsetminus e})$, then
$\sequencen{C_n\cup D_n}$
is a well-spread basis for $(\frak A,\bar\mu)$.

\proof{{\bf (a)(i)} $\sequencen{C_n}$ satisfies (i) of
Definition 528S just because $\tau(\frak B_{\kappa})=\#(C_n)=\kappa$ for
$n\ge 1$, while $C_0=\{1\}$.

\medskip

\quad{\bf (ii)} For $J\subseteq\kappa$, let $\frak C_J$ be the
order-closed subalgebra of $\frak B_{\kappa}$ generated by
$\{e_{\xi}:\xi\in J\}$;  recall that for every $a\in\frak B_{\kappa}$ there
is a countable set $\supp a\subseteq\kappa$ such that
$a\in\frak C_J$ iff $J\supseteq\supp a$ (254Rd/325Mb).
Of course $\#(\supp c)\le n$ whenever $n\in\Bbb N$ and $c\in C_n$.

Suppose that $a\in\frak B_{\kappa}$ and $\gamma>\bar\nu_{\kappa}a$.
Then for each $k\in\Bbb N$ we can find an $a_k\in\frak B_{\kappa}$, with
finite support, such that
$\bar\nu_{\kappa}(a\Bsymmdiff a_k)\le 2^{-k-2}(\gamma-\bar\nu_{\kappa}a)$
(254Fe/325Jc).   Set $b=\sup_{k\in\Bbb N}a_k$;  then

\Centerline{$\bar\nu_{\kappa}b
\le\bar\nu_{\kappa}a+\sum_{k=0}^{\infty}\bar\nu_{\kappa}(a_k\Bsetminus a)
<\gamma$,}

\Centerline{$\bar\nu_{\kappa}(a\Bsetminus b)
\le\inf_{k\in\Bbb N}\bar\nu_{\kappa}(a\Bsetminus a_k)=0$,}

\noindent so $a\Bsubseteq b$.
If $k\in\Bbb N$ and $\#(\supp a_k)=n_k$, then
$a_k=\sup\{c:c\in C_{n_k}$, $c\Bsubseteq a_k\}$, so
$b=\sup\{c:c\in\bigcup_{n\in\Bbb N}C_n$, $c\Bsubseteq b\}$.   Thus
528S(ii) is satisfied.

\medskip

\quad{\bf (iii)} If $n\in\Bbb N$ and $\sequence{i}{c_i}$ is a sequence
in $C_n$, there is an infinite $I\subseteq\Bbb N$ such
that $\family{i}{I}{\supp c_i}$ is a $\Delta$-system with root $K$ say
(5A1Jc).
For $i\in I$, express $c_i$ as $c'_i\Bcap c''_i$ where $c'_i\in\frak C_K$
and $c''_i\in\frak C_{(\supp c_i)\setminus K}$;  as $\frak C_K$ is finite,
there is
a $c$ such that $J=\{i:c'_i=c\}$ is infinite.   Now $c\in C_n$, and if
$m\in\Bbb N$ then

\Centerline{$\sup_{i\in J\setminus m}c_i
=c\Bcap\sup_{i\in J\setminus m}c''_i
=c$}

\noindent because $\family{i}{J\setminus m}{c''_i}$ is a stochastically
independent family of elements of $\frak B_{\kappa}$
all of measure at least $2^{-n}$, so has supremum $1$.
In particular, 528S(iii) is satisfied.

\medskip

\quad{\bf (iv)} Suppose that $n\in\Bbb N$ and $a\in\frak B_{\kappa}$.
Then there
is a $\delta>0$ such that $\bar\nu_{\kappa}(c\Bsetminus a)\ge\delta$ whenever
$c\in C_n$ and $c\notBsubseteq a$.  \Prf\Quer\
Otherwise, there is a sequence $\sequence{i}{c_i}$
in $C_n$ such that $0<\bar\nu_{\kappa}(c_i\Bsetminus a)\le 2^{-i}$ for every
$i\in\Bbb N$.   By (iii) just above, there is an infinite set
$J\subseteq\Bbb N$ such that
$c_j\Bsubseteq\sup_{i\in J\setminus m}c_i$ for every $j\in J$.
Set $j_0=\min J$, and let $m$ be such that
$2^{-m+1}<\bar\nu_{\kappa}(c_{j_0}\Bsetminus a)$;  then

\Centerline{$2^{-m+1}<\bar\nu_{\kappa}(\sup_{j\in J\setminus m}c_j\Bsetminus a)
\le\sum_{j=m}^{\infty}\bar\nu_{\kappa}(c_j\Bsetminus a)\le 2^{-m+1}$,}

\noindent which is absurd.\ \Bang\Qed

\medskip

\quad{\bf (v)} Suppose that $n\in\Bbb N$, $a\in\frak B_{\kappa}$ and
$\bar\nu_{\kappa} a\le\gamma'<\gamma<1$.   Pick
$\delta>0$, $r>n$, $k^*\in\Bbb N$, $\epsilon>0$ and $a'\in\frak B_{\kappa}$
such that

\Centerline{$\bar\nu_{\kappa}(c\Bsetminus a)\ge\delta$
whenever $c\in C_n$ and $c\notBsubseteq a$,}

\Centerline{$2^{-r}\le\gamma-\gamma'$,
\quad$(2^{-n}-2^{-r})^n\delta\ge 2^{-r+2}$,}

\Centerline{$(1-2^{-r})^{k^*}\le 1-\gamma$,}

\Centerline{$\epsilon\le\Bover12\delta$,
\quad$\epsilon\le 2^{-r}(1-2^{-r})^{k^*}$,}

\Centerline{$\supp a'$ is finite,
\quad$\bar\nu_{\kappa}(a\Bsymmdiff a')\le\epsilon$.}

Let $\sequence{i}{K_i}$ be a
disjoint sequence in $[\kappa\setminus\supp a']^r$, and set
$c_i=\inf_{\xi\in K_i}e_{\xi}$ for each $i\in\Bbb N$.
Then $\sup_{i\in\Bbb N}c_i=1$, so there is a first $k$ such that
$\bar\nu_{\kappa}(a\Bcup\sup_{i\le k}c_i)\ge\gamma$;  set $b_1=\sup_{i<k}c_i$
and $b=a\Bcup b_1$.
Surely $a\Bsubseteq b$ and $\bar\nu_{\kappa} b<\gamma$;  also

\Centerline{$(1-2^{-r})^{k^*}\le 1-\gamma\le 1-\bar\nu_{\kappa} b_1=(1-2^{-r})^k$}

\noindent so $k\le k^*$.   Moreover,

\Centerline{$\gamma-\bar\nu_{\kappa} b
\le\bar\nu_{\kappa}(b\Bcup c_k)-\bar\nu_{\kappa} b
\le\bar\nu_{\kappa}(c_k\Bsetminus b_1)=2^{-r}(1-2^{-r})^k
\le 2^{-r}\le\gamma-\gamma'$,}

\noindent so that, in particular, $\bar\nu_{\kappa} b\ge\gamma'$.

If $c\in C_n$ and $c\notBsubseteq a$ then
$\bar\nu_{\kappa}(c\Bsetminus a)\ge\delta$ so
$\bar\nu_{\kappa}(c\Bsetminus a')\ge\delta-\epsilon$.
Express $c$ as $\inf_{i\le k}c'_i$
where $\supp c'_i\subseteq K_i$ for $i<k$ and
$\supp c'_k\subseteq\kappa\setminus\bigcup_{i<k}K_i$.   Set
$J=\{i:i<k$, $c'_i\ne 1\}$;  then $\#(J)\le n$.   Now

$$\eqalignno{\bar\nu_{\kappa}(c\Bsetminus(a'\Bcup b_1))
&=\bar\nu_{\kappa}((c'_k\Bsetminus a')
     \Bcap\inf_{i\in J}(c'_i\Bsetminus c_i)
     \Bcap\inf_{i\in k\setminus J}(1\Bsetminus c_i))\cr
&=\bar\nu_{\kappa}(c'_k\Bsetminus a')
   \cdot\prod_{i\in J}\bar\nu_{\kappa}(c'_i\Bsetminus c_i)
   \cdot\prod_{i\in k\setminus J}\bar\nu_{\kappa}(1\Bsetminus c_i)\cr
\displaycause{because $\supp(c'_k\Bsetminus a')
\subseteq\supp c'_k\cup\supp a'
\subseteq\kappa\setminus\bigcup_{i<k}K_i$, so we are taking an
infimum of stochastically independent elements of $\frak B_{\kappa}$}
&\ge(\delta-\epsilon)\cdot\prod_{i\in J}(2^{-n}-2^{-r})
   \cdot\prod_{i\in k\setminus J}(1-2^{-r})\cr
\displaycause{of course every $c'_i$ belongs to $C_n$}
&\ge\Bover12(2^{-n}-2^{-r})^n(1-2^{-r})^k\delta
\ge 2^{-r+1}(1-2^{-r})^k\cr
&\ge 2^{-r}(1-2^{-r})^k+2^{-r}(1-2^{-r})^{k^*}
\ge\gamma-\bar\nu_{\kappa} b+\epsilon,\cr}$$

\noindent and

\Centerline{$\bar\nu_{\kappa}(c\Bsetminus b)\ge\gamma-\bar\nu_{\kappa} b$,}

\noindent so $\bar\nu_{\kappa}(b\Bcup c)\ge\gamma$.

As $n$, $a$, $\gamma'$ and $\gamma$ are arbitrary, $\sequencen{C_n}$
satisfies 528S(iv) and is a well-spread basis.

\medskip

\quad{\bf (vi)} As for (*), given $n\ge 1$,
take a disjoint family $\ofamily{\xi}{\kappa}{K_{\xi}}$
in $[\kappa]^n$, and set $c_{\xi}=\inf_{\eta\in K_{\xi}}e_{\eta}$ for
$\xi<\kappa$, $C'_n=\{c_{\xi}:\xi<\kappa\}$.   If $I\subseteq\kappa$ is
infinite and $\bar\nu_{\kappa} a<1$, take $\delta>0$ such that
$\delta<2^{-n}(1-\bar\nu_{\kappa} a-\delta)$, and
$a'\in\frak B_{\kappa}$ such that
$\supp a'$ is finite and $\bar\nu_{\kappa}(a\Bsymmdiff a')\le\delta$.
Then there is a
$\xi\in I$ such that $K_{\xi}\cap\supp a_{\xi}=\emptyset$.
\Quer\
If $c\in C_n$ is such that $c_{\xi}\Bsubset c\Bsubseteq a\Bcup c_{\xi}$,
there must be a $d\in\frak B_{\kappa}$, with support $K_{\xi}$, included in
$c\Bsetminus c_{\xi}$.   But now $d\Bsubseteq a$ and
$\supp d\Bcap\supp a'=\emptyset$, so

\Centerline{$2^{-n}(1-\bar\nu_{\kappa} a-\delta)
\le 2^{-n}(1-\bar\nu_{\kappa} a')
=\bar\nu_{\kappa}(d\Bsetminus a')
\le\delta+\bar\nu_{\kappa}(d\Bsetminus a)=\delta$.  \Bang}

\woddheader{528T}{4}{2}{2}{60pt}

{\bf (b)(i)} We have

$$\eqalign{\#(\bigcup_{n\in\Bbb N}C_n\cup D_n)
&\le\max(\omega,\#(\bigcup_{n\in\Bbb N}C_n),\#(\bigcup_{n\in\Bbb N}D_n))
  \cr
&\le\max(\omega,c(\frak A_e),c(\frak A_{1\Bsetminus e}),
  \tau(\frak A_e),\tau(\frak A_{1\Bsetminus e}))
=\max(\omega,c(\frak A),\tau(\frak A))\cr}$$

\noindent by 514E.   %514Ed 514Ef

\medskip

\quad{\bf (ii)} Suppose that $a\in\frak A$ and $\bar\mu a<\gamma$.
Then there are $\gamma_1$, $\gamma_2$ such that
$\bar\mu(a\Bcap e)<\gamma_1$, $\bar\mu(a\Bsetminus e)<\gamma_2$ and
$\gamma_1+\gamma_2\le\gamma$.   Let $C\subseteq\bigcup_{n\in\Bbb N}C_n$,
$D\subseteq\bigcup_{n\in\Bbb N}D_n$ be such that
$a\Bcap e\Bsubseteq\sup C$, $a\Bsetminus e\Bsubseteq\sup D$,
$\bar\mu(\sup C)<\gamma_1$ and $\bar\mu(\sup D)<\gamma_2$.   Then
$C\cup D\subseteq\bigcup_{n\in\Bbb N}C_n\cup D_n$,
$a\Bsubseteq\sup(C\cup D)$ and $\bar\mu(\sup(C\cup D))<\gamma$.

\medskip

\quad{\bf (iii)} Suppose that $n\in\Bbb N$ and
$\sequence{i}{c_i}$ is a sequence in $C_n\cup D_n$ such that
$\bar\mu(\sup_{i\in\Bbb N}c_i)<\infty$.   Then
there is an infinite $J\subseteq\Bbb N$ such that either
$c_i\in C_n$ for every $i\in J$, or $c_i\in D_n$ for every $i\in I$.
In either case, there is an infinite $I\subseteq J$ such that
$\sup_{i\in I}c_i$ belongs to $C_n\cup D_n$.

\medskip

\quad{\bf (iv)} Thus $\sequencen{C_n\cup D_n}$ satisfies (i)-(iii) of
Definition 528S.   As for 528S(iv), suppose that
$n\in\Bbb N$, $a\in\frak A$ and
$\bar\mu a\le\gamma'<\gamma<\bar\mu 1$.   We need to find a
$b\in\frak A$ such that $a\Bsubseteq b$ and

\Centerline{$\gamma'\le\bar\mu b<\gamma
\le\bar\mu(b\Bcup c)$}

\noindent whenever $c\in C_n\cup D_n$ and $c\notBsubseteq a$.

\medskip

\qquad{\bf case 1} If $e\Bsubseteq a$, then $\bar\mu e$ is finite and

\Centerline{$\bar\mu(a\Bsetminus e)\le\gamma'-\bar\mu e
<\gamma-\bar\mu e<\bar\mu(1\Bsetminus e)$.}

\noindent So there is a $b_2\in\frak A_{1\Bsetminus e}$ such that
$a\Bsetminus e\Bsubseteq b_2$ and

\Centerline{$\gamma'-\bar\mu e\le\bar\mu b_2
<\gamma-\bar\mu e\le\bar\mu(b_2\Bcup d)$}

\noindent whenever $d\in D_n$ and $d\notBsubseteq a\Bsetminus e$;  that is,

\Centerline{$\gamma'\le\bar\mu(e\Bcup b_2)
<\gamma\le\bar\mu(e\Bcup b_2\Bcup d)$}

\noindent whenever $d\in D_n$ and $d\notBsubseteq a$.   Since
$c\Bsubseteq a$ for every $c\in C_n$, we have
$\bar\mu(e\Bcup b_2\Bcup c)\ge\gamma$ whenever $c\in C_n\cup D_n$ and
$c\notBsubseteq a$, and can take $b=e\Bcup b_2$.

\medskip

\qquad{\bf case 2} Similarly, if $a\Bsupseteq 1\Bsetminus e$, we can take
$b=(1\Bsetminus e)\Bcup b_1$ for a suitable $b_1\Bsubseteq e$.

\medskip

\qquad{\bf case 3} If neither $e$ nor $1\Bsetminus e$ is included in $a$,
we have

\Centerline{$\max(\bar\mu(a\Bcap e),\gamma-\bar\mu(1\Bsetminus e))
<\min(\bar\mu e,\gamma-\bar\mu(a\Bsetminus e))$,}

\noindent so we can find $\gamma'_1$, $\gamma_1$ such that

\Centerline{$\max(\bar\mu(a\Bcap e),\gamma-\bar\mu(1\Bsetminus e))
<\gamma'_1<\gamma_1<\min(\bar\mu e,\gamma-\bar\mu(a\Bsetminus e))$}

\noindent and $\gamma_1-\gamma'_1<\gamma-\gamma'$.
Let $b_1\in\frak A_e$ be such that $a\Bcap e\Bsubseteq b_1$ and

\Centerline{$\gamma'_1\le\bar\mu b_1<\gamma_1
\le\bar\mu(b_1\Bcup c)$}

\noindent whenever $c\in C_n$ and $c\notBsubseteq a\Bcap e$.   Set
$\gamma'_2=\gamma-\gamma_1$ and $\gamma_2=\gamma-\bar\mu b_1$, so that

\Centerline{$\bar\mu(a\Bsetminus e)<\gamma'_2<\gamma_2
\le\gamma-\gamma'_1<\bar\mu(1\Bsetminus e)$.}

\noindent Let
$b_2\in\frak A_{1\Bsetminus e}$ be such that $a\Bsetminus e\Bsubseteq b_2$
and

\Centerline{$\gamma'_2\le\bar\mu b_2<\gamma_2
\le\bar\mu(b_2\Bcup d)$}

\noindent whenever $d\in D_n$ and $d\notBsubseteq a\Bsetminus e$.

Try $b=b_1\Bcup b_2$.   Then $a\Bsubseteq b$ and
$\bar\mu b=\bar\mu b_1+\bar\mu b_2$ belongs to

\Centerline{$\coint{\bar\mu b_1+\gamma'_2,\bar\mu b_1+\gamma_2}
\subseteq\coint{\gamma'_1+\gamma-\gamma_1,\gamma}
\subseteq\coint{\gamma',\gamma}$.}

\noindent If $c\in C_n$ and $c\notBsubseteq a$, then
$c\notBsubseteq a\Bcap e$, so

\Centerline{$\bar\mu(b\Bcup c)
=\bar\mu(b_1\Bcup c)+\bar\mu b_2
\ge\gamma_1+\gamma'_2
=\gamma$;}

\noindent while if $d\in D_n$ and $d\notBsubseteq a$, then

\Centerline{$\bar\mu(b\Bcup d)
=\bar\mu b_1+\bar\mu(b_2\Bcup d)
\ge\mu b_1+\gamma_2
=\gamma$.}

\noindent So in this case also we have found a suitable $b$.
}%end of proof of \zz1F

\leader{528U}{Lemma}\dvAnew{2011}
Let $(\frak A,\bar\mu)$ be an atomless semi-finite
measure algebra and $0<\gamma<\bar\mu 1$.   Let $E$, $\epsilon$,
$\preccurlyeq$ and $\Cal F$ be such that

\inset{$E$ is a partition of unity in $\frak A$ such that $\frak A_e$ is
homogeneous and $0<\epsilon\le\bar\mu e<\infty$ for every $e\in E$;

$\preccurlyeq$ is a well-ordering of $E$ such that
$\tau(\frak A_e)\le\tau(\frak A_{e'})$ whenever $e\preccurlyeq e'$ in $E$;

$\Cal F$ is a partition of $E$ such that each member of $\Cal F$ is either
a singleton or a countable set with no $\preccurlyeq$-greatest member.}

\noindent Let $P_0$ be

\Centerline{$\{a:a\in\frak A$, $\bar\mu a<\gamma$,
$\gamma\le\bar\mu(a\Bcup e)$ whenever $\{e\}\in\Cal F\}$,}

\noindent ordered by $\Bsubseteq$.   Then $\RO^{\uparrow}(P_0)$ has
countable Maharam type.

\proof{{\bf (a)(i)}
For every $e\in E$, $(\frak A_e,\bar\mu\restrp\frak A_e)$
is a non-zero atomless homogeneous totally finite measure algebra, so is
isomorphic, up to a scalar multiple of the measure, to
$(\frak B_{\kappa},\bar\nu_{\kappa})$ for some infinite cardinal $\kappa$
(331L).   So we can copy the well-spread basis for
$(\frak B_{\kappa},\bar\nu_{\kappa})$ described in 528Ta into a
well-spread basis $\sequencen{D_{en}}$ for
$(\frak A_e,\bar\mu\restrp\frak A_e)$ such that

\inset{$\#(\bigcup_{n\in\Bbb N}D_{en})=\tau(\frak A_e)$,

$\bar\mu d\ge 2^{-n}\bar\mu e$ whenever $n\in\Bbb N$ and $d\in D_{en}$,

$D_{e0}=\{e\}$,

for each $n\ge 1$ there is a set $D'_{en}\subseteq D_{en}$, with cardinal
$\tau(\frak A_e)$, such that $\bar\mu d=2^{-n}\bar\mu e$ for every
$d\in D'_{en}$, and whenever $a\in\frak A_e\setminus\{e\}$ and
$I\subseteq D'_{en}$ is infinite, there is a $d\in I$ such that
$d'\notBsubseteq a\Bcup d$ whenever $d'\in D_{en}$ and $d'\Bsupset d$,

$(\bigcup_{n\in\Bbb N}D_{en})\setminus(\bigcup_{n\ge 1}D'_{en})$ has
cardinal $\tau(\frak A_e)$.}

\noindent (The last item is not mentioned in 528T, but is clearly
achievable by thinning the sets $D'_{en}$ appropriately, besides being
automatic if we use the construction in (a-vi) of the proof of 528T.)
Note that $\langle D'_{en}\rangle_{n\ge 1}$
is a disjoint sequence of subsets of
$\frak A_e$ for each $e$, so $\langle D'_{en}\rangle_{e\in E,n\ge 1}$
is disjoint.

\medskip

\quad{\bf (ii)} For $e\in F\in\Cal F$, set

\Centerline{$D_e
=\bigcup_{n\in\Bbb N}D_{en}\setminus\bigcup_{n\ge 1}D'_{en}$,
\quad$D^*_e
=\bigcup_{e'\in F,e'\preccurlyeq e}D_e$.}

\noindent Because $F$ is countable and
$\tau(\frak A_{e'})\le\tau(\frak A_e)$ whenever $e'\preccurlyeq e$,
$\#(D^*_e)=\tau(\frak A_e)=\#(D'_{en})$ for every $n\ge 1$.
We therefore have a partition
$\family{d}{D^*_e}{I_{ed}}$ of $\bigcup_{n\ge 1}D'_{en}$ into countably
infinite sets such that $I_{ed}\cap D'_{en}$ is infinite whenever
$d\in D^*_e$ and $n\ge 1$.

Let $\theta$ be a limit ordinal such that the set $\Omega$ of limit
ordinals less than $\theta$ has cardinal
$\#(\bigcup_{e\in E}D_e)$.   (Of course we can take $\theta$ to be either
an uncountable cardinal or the ordinal product $\omega\cdot\omega$ or $0$.)
Again because every member of
$\Cal F$ is countable, we have an enumeration
$\ofamily{\xi}{\theta}{d_{\xi}}$ of $\bigcup_{e\in E,n\in\Bbb N}D_{en}$
such that whenever $\xi\in\Omega$ then there are $F\in\Cal F$ and $e\in F$
such that

\Centerline{$d_{\xi}\in D_e$,
\quad$\{d_{\xi+i}:i\ge 1\}
=\bigcup_{e'\in F,e'\succcurlyeq e}I_{e'd_{\xi}}$.}

\noindent This will mean that whenever $\xi\in\Omega$ and $F\in\Cal F$,
$e\in F$ are such that $d_{\xi}\in\frak A_e$, then
$\{i:d_{\xi+i}\in D'_{e'n}\}$ is infinite whenever $e'\in F$,
$e\preccurlyeq e'$ and $n\in\Bbb N$.

\medskip

{\bf (b)(i)} Setting $P=\{a:a\in\frak A$, $\bar\mu a<\gamma\}$,
$P_0\in\RO^{\uparrow}(P)$.   \Prf\ Evidently $P_0$ is up-open.   If
$a\in P\setminus P_0$, that is, there is some $e$ such that
$\{e\}\in\Cal F$ and $\bar\mu(a\Bcup e)<\gamma$, set $b=a\Bcup e$;  then
$a\Bsubseteq b\in P$, while $\bar\mu(b'\Bcup e)=\bar\mu b'<\gamma$ whenever
$b'\in\coint{b,\infty}$, so $\coint{b,\infty}$ does not meet $P_0$.
Accordingly $\coint{a,\infty}\not\subseteq\overline{P}_0$ and
$a\notin\interior\overline{P}_0$.   As $a$ is arbitrary,
$P_0=\interior\overline{P}_0\in\RO^{\uparrow}(P)$.\ \Qed

It follows that $\RO^{\uparrow}(P_0)$ is the principal ideal of
$\RO^{\uparrow}(P)$ generated by $P_0$ (314R(b-ii)).
Moreover, for $a\in P_0$, $\coint{a,\infty}$ is the same
whether taken in $P$ or $P_0$, and belongs to $\RO^{\uparrow}(P)$ by
528B(b-i).

\medskip

\quad{\bf (ii)} For $a\in P_0$ and $n\in\Bbb N$, set
$A_n(a)=\{d:d\in\bigcup_{e\in E}D_{en}$, $d\Bsubseteq a\}$.   Then any
sequence in $A_n(a)$ has a subsequence with an upper bound in $A_n(a)$.
\Prf\ Set $L=\{e:e\in E$, $\bar\mu(a\Bcap e)\ge 2^{-n}\epsilon\}$;  then
$L$ is finite.   If $e\in E\setminus L$ and $d\in D_{en}$, then
$d\Bsubseteq e$ and

\Centerline{$\bar\mu d\ge 2^{-n}\bar\mu e\ge 2^{-n}\epsilon
>\bar\mu(a\Bcap e)\ge\bar\mu(a\Bcap d)$,}

\noindent so
$d\notBsubseteq a$.   Thus $A_n(a)\subseteq\bigcup_{e\in L}D_{en}$.
It follows that if $\sequence{i}{c_i}$ is any sequence in $A_n(a)$, there
is an $e\in L$ such that $J=\{i:c_i\in D_{en}\}$ is infinite.   Now there
is an infinite $I\subseteq J$ such that $c=\sup_{i\in I}c_i$ belongs to
$D_{en}$.   In this case, $c\Bsubseteq a$ so $c\in A_n(a)$ is an upper
bound of $\{c_i:i\in I\}$.\ \Qed

It follows that $A_n(a)$ has only finitely many maximal elements, and any
non-decreasing sequence in $A_n(a)$ has an upper bound in $A_n(a)$.
Consequently,
every member of $A_n(a)$ is included in a maximal element of $A_n(a)$.
\Prf\Quer\ Otherwise, we should be able to
find a strictly increasing family
$\ofamily{\xi}{\omega_1}{c_{\xi}}$ in $A_n(a)$;   but now there must be a
$\xi<\omega_1$ such that $\bar\mu c_{\xi}=\bar\mu c_{\xi+1}<\gamma$ and
$c_{\xi}=c_{\xi+1}$.\ \Bang\Qed

Set $E_n(a)=\{\xi:d_{\xi}$ is a maximal element of $A_n(a)\}$, so that
$E_n(a)$ is a finite subset of $\theta$.

\medskip

\quad{\bf (iii)} For $n\in\Bbb N$, set

\Centerline{$Q_n=\{b:b\in P_0$, $A_n(b)=A_n(b')$ whenever
$b\Bsubseteq b'\in P_0\}$.}

\noindent Then whenever $a\in P_0$ and $n\in\Bbb N$ there is a
$b\in Q_n$ such that $a\Bsubseteq b$ and $A_n(a)=A_n(b)$.
\Prf\ Let $L$ be a finite subset of $E$ including
$\{e:\bar\mu(a\Bcap e)\ge 2^{-n-1}\epsilon\}$ and such that
$\bar\mu(\sup L)>\gamma$.   Then
$\sequence{m}{\bigcup_{e\in L}D_{em}}$ is a well-spread basis for
$(\frak A_{\sup L},\bar\mu\restrp\frak A_{\sup L})$.   (Induce on $\#(L)$,
using 528Tb for the inductive step.)   Since

\Centerline{$\bar\mu(a\Bcap\sup L)
<\gamma-\bar\mu(a\Bsetminus\sup L)
<\bar\mu(\sup L)$,}

\noindent there is a a $b_0\in\frak A_{\sup L}$, including $a\Bcap\sup L$,
such that

\Centerline{$\gamma-\bar\mu(a\Bsetminus\sup L)-2^{-n-1}\epsilon
\le\bar\mu b_0
<\gamma-\bar\mu(a\Bsetminus\sup L)
\le\bar\mu(b_0\Bcup d)$}

\noindent whenever $d\in\bigcup_{e\in L}D_{en}$ and $d\notBsubseteq a$.
Set $b=b_0\Bcup a$.   Then
$\bar\mu b=\bar\mu b_0+\bar\mu(a\Bsetminus\sup L)<\gamma$, so
$b\in P_0$.   If $b\Bsubseteq b'\in P_0$ and
$d\in\bigcup_{e\in E}D_{en}\setminus A_n(a)$, then either $e\in L$ and

\Centerline{$\bar\mu(b'\Bcup d)
\ge\bar\mu(b\Bcup d)+\bar\mu(a\Bsetminus\sup L)
\ge\gamma>\bar\mu b'$,}

\noindent or $e\notin L$,

\Centerline{$\bar\mu(d\Bsetminus a)\ge\bar\mu d-\bar\mu(a\Bcap e)
\ge 2^{-n}\bar\mu e-2^{-n-1}\epsilon\ge 2^{-n-1}\epsilon$}

\noindent and

\Centerline{$\bar\mu(b'\Bcup d)
\ge\bar\mu b_0+\bar\mu(a\Bsetminus\sup L)+2^{-n-1}\epsilon\ge\gamma
>\bar\mu b'$;}

\noindent in either case $d\notBsubseteq b'$.
Thus $A_n(b')=A_n(a)=A_n(b)$ whenever
$b\Bsubseteq b'\in P_0$, and $b\in Q_n$.\ \Qed

\medskip

{\bf (c)(i)} For $m$, $n$, $i\in\Bbb N$ and $\xi\in\Omega$, set

\Centerline{$Q_{nmi\xi}=\{b:b\in Q_n$, $\xi+i\in E_n(b)$,
$\#(E_n(b)\cap\xi)=m\}$,}

\Centerline{$G_{nmi\xi}=\sup\{\coint{b,\infty}:b\in Q_{nmi\xi}\}
\in\RO^{\uparrow}(P_0)$.}

\medskip

\quad{\bf (ii)} For any $m$, $n$, $i\in\Bbb N$,
$\family{\xi}{\Omega}{G_{nmi\xi}}$ is disjoint.   \Prf\ Suppose that
$\xi<\eta$ in $\Omega$.   If $a\in Q_{nmi\xi}$ and $b\in Q_{nmi\eta}$,
we see that $\xi+i<\eta$, $\xi+i\in E_n(a)$ and

\Centerline{$\#(E_n(b)\cap\eta)=m
=\#(E_n(a)\cap\xi)<\#(E_n(a)\cap\eta)$.}

\noindent So $E_n(a)\ne E_n(b)$ and $A_n(a)\ne A_n(b)$.   But both $a$ and
$b$ are supposed to belong to $Q_n$, so $\coint{a,\infty}$ must be disjoint
from $\coint{b,\infty}$.
As $b$ is arbitrary, $\coint{a,\infty}\cap G_{nmi\eta}=\emptyset$;
as $a$ is arbitrary, $G_{nmi\xi}\cap G_{nmi\eta}=\emptyset$.\ \Qed

\medskip

\quad{\bf (iii)} For any $\xi\in\Omega$ and $a\in P_0$, there are $m$, $n$,
$i\in\Bbb N$ and $b\in Q_{nmi\xi}$ such that $a\Bsubseteq b$.
\Prf\ Let $e\in E$ be such that $d_{\xi}\Bsubseteq e$;  let $F$ be the
member of $\Cal F$ containing $e$.   If $F=\{e\}$, then
$\bar\mu(a\Bcup e)\ge\gamma>\bar\mu a$;  set $e_0=e$, so that
$e_0\in F$, $e_0\succcurlyeq e$ and $a\Bcap e_0\ne e_0$.   Otherwise,
there are infinitely many members of $F$ greater than $e$ for the ordering
$\preccurlyeq$, because $F$ has no greatest member, so
$\bar\mu(\sup_{e'\in F,e'\succcurlyeq e}e')=\infty$, and there must be
an $e_0\in F$ such that $e_0\succcurlyeq e$ and $a\Bcap e_0\ne e_0$.

Let $n\in\Bbb N$ be such that
$2^{-n}\bar\mu e_0<\min(\gamma-\bar\mu a,\bar\mu(e_0\Bsetminus a))$.
Then $\{d_{\xi+i}:i\in\Bbb N\}$ meets $D'_{e_0n}$ in an infinite set, and
there is an $i\in\Bbb N$ such that
$d_{\xi+i}\in D'_{e_0n}$, $\bar\mu d_{\xi+i}=2^{-n}\bar\mu e_0$,
and $d\notBsubseteq(a\Bcap e_0)\Bcup d_{\xi+i}$ whenever $d\in D_{e_0n}$
and $d\Bsupset d_{\xi+i}$.   Set $a'=a\Bcup d_{\xi+i}$;  then
$d_{\xi+i}$ is a maximal member of $A_n(a')$.   Let $b\in Q_n$ be such that
$a'\Bsubseteq b$ and $A_n(b)=A_n(a')$.   Then $\xi+i\in E_n(b)$.
Set $m=\#(E_n(b)\cap\xi)$.   Then $b\in Q_{nmi\xi}$ and
$a\Bsubseteq b$.\ \Qed

Accordingly $b\in\coint{a,\infty}\cap G_{nmi\xi}$.
As $a$ is arbitrary, $\bigcup_{m,n,i\in\Bbb N}G_{nmi\xi}$ is dense in $P_0$
and $\sup_{m,n,i\in\Bbb N}G_{nmi\xi}=P_0$ in $\RO^{\uparrow}(P_0)$.

\medskip

{\bf (d)(i)} Let $\frak G$ be the order-closed subalgebra of
$\RO^{\uparrow}(P_0)$ generated by
$\{G_{nmi\xi}:m$, $n$, $i\in\Bbb N$, $\xi\in\Omega\}$.
By (c-ii) and (c-iii), the conditions of 514F
are satisfied, and $\frak G$ has countable Maharam type.

\medskip

\quad{\bf (ii)} If $d\in P_0\cap\bigcup_{e\in E,n\in\Bbb N}D_{en}$ then
$\coint{d,\infty}\in\frak G$.   \Prf\ Set

\Centerline{$H
=\sup\{G_{nmi\xi}:m$, $n$, $i\in\Bbb N$, $\xi\in\Omega$ and
$G_{nmi\xi}\subseteq\coint{d,\infty}\}\in\RO^{\uparrow}(P_0)$.}

\noindent Then $H\in\frak G$ and $H\Bsubseteq\coint{d,\infty}$.
Suppose that $a\in P_0$ and $a\Bsupseteq d$.
Let $n\in\Bbb N$ be such that $d\in\bigcup_{e\in E}D_{en}$.
Then there is a $b\in Q_n$
such that $a\Bsubseteq b$.   In this case, $d\in A_n(b)$ so there is a
maximal $d'\in A_n(b)$ including $d$;  let $\xi\in\Omega$,
$i\in\Bbb N$ be such that
$d'=d_{\xi+i}$, and set $m=\#(E_n(b)\cap\xi)$.   Then $b\in Q_{nmi\xi}$.
On the other hand, for any
$b'\in Q_{nmi\xi}$, $d\Bsubseteq d_{\xi+i}\Bsubseteq b'$, so
$\coint{b',\infty}\subseteq\coint{d,\infty}$;  as
$b'$ is arbitrary, $G_{nmi\xi}\subseteq\coint{d,\infty}$ and
$G_{nmi\xi}\subseteq H$.   Accordingly $b\in H\cap\coint{a,\infty}$.
As $a$ is arbitrary, $H$ is dense in $\coint{d,\infty}$ and must be the
whole of $\coint{d,\infty}$;  thus we have $\coint{d,\infty}=H\in\frak G$.\
\Qed

\medskip

\quad{\bf (iii)} If $a\in P_0$ there is a $b\in P_0$ such that
$a\Bsubseteq b$ and $\coint{b,\infty}\in\frak G$.   \Prf\ Let $E_0$ be a
countable subset of $E$ such that $a\Bsubseteq\sup E_0$ and
$\bar\mu(\sup E_0)>\gamma$.
Set $L=\{e:e\in E_0$, $a\Bsupseteq e\}$.   Then $E_0\setminus L$ is
non-empty, and

\Centerline{$\sum_{e\in E_0\setminus L}\bar\mu(a\Bcap e)
=\bar\mu a-\bar\mu(\sup L)<\gamma-\bar\mu(\sup L)$.}

\noindent We therefore have a family
$\family{e}{E_0\setminus L}{\gamma_e}$ such that
$\bar\mu(a\Bcap e)<\gamma_e\le\bar\mu e$ for every $e\in E_0\setminus L$
and $\sum_{e\in E_0\setminus L}\gamma_e<\gamma-\bar\mu(\sup L)$.
For each $e\in E_0$ there is a $B_e\subseteq\bigcup_{n\in\Bbb N}D_{en}$
such that $a\Bcap e\Bsubseteq\sup B_e$ and $\bar\mu(\sup B_e)\le\gamma_e$,
by 528S(ii).   Set

\Centerline{$B=L\cup\bigcup_{e\in E_0\setminus L}B_e
\subseteq\bigcup_{e\in E,n\in\Bbb N}D_{en}$}

\noindent and $b=\sup B$.   Then $a\Bsubseteq b$ and

\Centerline{$\bar\mu b
=\bar\mu(\sup L)+\sum_{e\in E_0\setminus L}\bar\mu(\sup B_e)
\le\bar\mu(\sup L)+\sum_{e\in E_0\setminus L}\gamma_e
<\gamma$,}

\noindent so $b\in P_0$.   On the other hand,

\Centerline{$\coint{b,\infty}=\bigcap_{d\in B}\coint{d,\infty}
=\inf_{d\in B}\coint{d,\infty}\in\frak G$,}

\noindent as required.\ \Qed

\medskip

\quad{\bf (iv)} As $a$ is arbitrary, $\frak G$ includes a $\pi$-base for
the Boolean algebra $\RO^{\uparrow}(P_0)$ and must be the whole of
$\RO^{\uparrow}(P_0)$.   Accordingly

\Centerline{$\tau(\RO^{\uparrow}(P_0))=\tau(\frak G)\le\omega$.}

\noindent This completes the proof.
}%end of proof of 528U

\leader{528V}{Theorem}\dvAnew{2011;  extends former 5{}28Rb}
Let $(\frak A,\bar\mu)$ be an atomless semi-finite
measure algebra and $0<\gamma<\bar\mu 1$.   Then
$\AM(\frak A,\bar\mu,\gamma)$ has countable Maharam type.

\proof{ Throughout the proof, $P$ will stand for
$\{a:a\in\frak A$, $\bar\mu a<\gamma\}$.

\medskip

{\bf (a)} Suppose that there are a partition $E$ of unity in
$\frak A$ and an $\epsilon>0$ such that $\frak A_e$ is homogeneous and
$\epsilon\le\bar\mu e<\infty$ for every $e\in E$.

\medskip

\quad{\bf (i)}  Let $\preccurlyeq$ be a well-ordering of $E$ such that
$\tau(\frak A_e)\le\tau(\frak A_{e'})$ whenever $e\preccurlyeq e'$ in $E$.
Let $\Cal F_0$ be a maximal disjoint
family of subsets of $E$ of order type
$\omega$ in the ordering induced by $\preccurlyeq$.
Then $M=E\setminus\bigcup\Cal F_0$ must be finite;  set
$\Cal F=\Cal F_0\cup\{\{e\}:e\in M\}$.

\medskip

\quad{\bf (ii)} For $L\subseteq M$, set

\Centerline{$P_L
=\{a:a\in P$, $a\Bsupseteq\sup L$, $\bar\mu(a\Bcup e)\ge\gamma$ for
$e\in M\setminus L\}$.}

\noindent Then $\langle P_L\rangle_{L\subseteq M}$ is a disjoint family of
open subsets of $P$.   Also $\bigcup_{L\subseteq M}P_L$ is dense in $P$.
\Prf\ If $a\in P$, let $L\subseteq M$ be a maximal set such that
$\bar\mu(a\Bcup\sup L)<\gamma$, and set $b=a\Bcup\sup L$;  then
$a\Bsubseteq b\in P_L$.\ \QeD\  So
$\RO^{\uparrow}(P)$ is isomorphic to the simple product
$\prod_{L\subseteq M}\RO^{\uparrow}(P_L)$ (315H again).

\medskip

\quad{\bf (iii)}
If $L\subseteq M$, then $\RO^{\uparrow}(P_L)$ has countable
Maharam type.   \Prf\ If $P_L=\emptyset$ this is trivial.   Otherwise there
is an $a\in P_L$ and $\bar\mu(\sup L)\le\bar\mu a<\gamma$.
Consider $\frak A'=\frak A_{1\Bsetminus\sup L}$,
$\gamma'=\gamma-\bar\mu(\sup L)$, $E'=E\setminus L$,
$\Cal F'=\Cal F\setminus\{\{e\}:e\in L\}$ and
$\preccurlyeq'\mskip5mu=\mskip5mu\preccurlyeq\mskip-3mu
  \cap\mskip2mu(E'\times E')$.
Then $(\frak A',\bar\mu\restrp\frak A')$,
$\gamma'$, $E'$, $\epsilon$, $\preccurlyeq'$ and $\Cal F'$ satisfy the
conditions of 528U.   Setting

\Centerline{$Q_0
=\{c:c\in\frak A'$,
$\bar\mu c<\gamma'\le\bar\mu(c\Bcup e)$ for every $e\in M\setminus L\}$,}

\noindent $\RO^{\uparrow}(Q_0)$ has countable Maharam type, by 528U.
But the map $c\mapsto c\Bcup\sup L$ is an order-isomorphism between
$Q_0$ and $P_L$, so $\RO^{\uparrow}(P_L)$ has countable Maharam type.\ \Qed

\medskip

\quad{\bf (iv)} Thus $\AM(\frak A,\bar\mu,\gamma)=\RO^{\uparrow}(P)$
is isomorphic to the product of finitely many Boolean algebras
with countable Maharam type, and has countable Maharam type (514Ef).

\medskip

{\bf (b)} Now suppose that
$(\frak A,\bar\mu)$ is localizable.

\medskip

\quad{\bf (i)} In this case, let $E$ be a partition
of unity in $\frak A$ such that $\frak A_e$ is homogeneous and
$0<\bar\mu e<\infty$ for every $e\in E$.   Let $\epsilon>0$ be such that
$\sum_{e\in E,\bar\mu e\ge\epsilon}\bar\mu e>\gamma$.
For each $k\in\Bbb N$, set

\Centerline{$E_k=\{e:e\in E$, $\bar\mu e\ge 2^{-k}\epsilon\}$,
\quad$e^*_k=\sup E_k$.}

\noindent By (a),
$\AM(\frak A_{e^*_k},\bar\mu\restrp\frak A_{e^*_k},\gamma)$
has countable Maharam type for every $k$.

\medskip

\quad{\bf (ii)} Now 528Fb tells us that we have a sequence
$\sequence{k}{\pi_k}$ such that
$\pi_k$ is a regular embedding of
$\AM(\frak A_{e^*_k},\bar\mu\restrp\frak A_{e^*_k},\gamma)$ into
$\AM(\frak A,\bar\mu,\gamma)$ for each $k$, and
$\bigcup_{k\in\Bbb N}
\pi_k[\AM(\frak A_{e^*_k},\bar\mu\restrp\frak A_{e^*_k},\gamma)]
\,\,\tau$-generates $\AM(\frak A,\bar\mu,\gamma)$.
So $\AM(\frak A,\bar\mu,\gamma)$ has countable Maharam type.
\Prf\ For each $k$, we have a countable $\tau$-generating set
$D_k\subseteq\AM(\frak A_{e^*_k},\bar\mu\restrp\frak A_{e^*_k},\gamma)$.
Let $\frak G$ be the
order-closed subalgebra of $\AM(\frak A,\bar\mu,\gamma)$ generated by
$D=\bigcup_{k\in\Bbb N}\pi_k[D_k]$.   For each $k\in\Bbb N$,
$\pi_k^{-1}[\frak G]$ is an order-closed subalgebra of
$\AM(\frak A_{e^*_k},\bar\mu\restrp\frak A_{e^*_k},\gamma)$ including
$D_k$, so is the whole of
$\AM(\frak A_{e^*_k},\bar\mu\restrp\frak A_{e^*_k},\gamma)$, that is,
$\pi_k[\AM(\frak A_{e^*_k},\bar\mu\restrp\frak A_{e^*_k},\gamma)]
\subseteq\frak G$.   Since
$\bigcup_{k\in\Bbb N}
\pi_k[\AM(\frak A_{e^*_k},\bar\mu\restrp\frak A_{e^*_k},\gamma)]
\,\,\tau$-generates $\AM(\frak A,\bar\mu,\gamma)$,
$\frak G=\AM(\frak A,\bar\mu,\gamma)$ and
$\tau(\AM(\frak A,\bar\mu,\gamma))\le\#(D)\le\omega$.\ \Qed

\medskip

{\bf (c)} Thus we have the result when $(\frak A,\bar\mu)$ is localizable.
For the general case of atomless semi-finite $(\frak A,\bar\mu)$,
let $(\widehat{\frak A},\tilde\mu)$ be the
localization of $(\frak A,\bar\mu)$ (322Q).
Since the embedding $\frak A\embedsinto\widehat{\frak A}$ identifies
$\frak A^f$ with $\widehat{\frak A}^f$ (322P),
$\{a:a\in\widehat{\frak A}$, $\tilde\mu a<\gamma\}$
can be identified with $P$, and the regular open algebras
$\AM(\frak A,\bar\mu,\gamma)$ and
$\AM(\widehat{\frak A},\tilde\mu,\gamma)$ are isomorphic.   Again because
$\frak A^f$ and $\widehat{\frak A}^f$ are isomorphic, $\widehat{\frak A}$
is atomless.   By (b), the common Maharam type of
$\AM(\frak A,\bar\mu,\gamma)$ and
$\AM(\widehat{\frak A},\tilde\mu,\gamma)$ is countable.
}%end of proof of 528V

\exercises{\leader{528X}{Basic exercises (a)}
%\spheader 528Xa
Suppose that $(X,\Sigma,\mu)$ is a measure
space and $(\frak A,\bar\mu)$ its measure algebra.   Let
$\Cal E\subseteq\Sigma$ be a family such that $\mu$ is outer
regular with respect to $\Cal E$, and $P$ the set
$\{(E,\alpha):E\in\Cal E$, $\mu E<\alpha\le\mu X\}$, ordered by saying that
$(E,\alpha)\le(F,\beta)$ if $E\subseteq F$ and $\beta\le\alpha$.   Show that
$\RO^{\uparrow}(P)$ is isomorphic to $\AM^*(\frak A,\bar\mu)$.
%528C out of order query

\spheader 528Xb
 Let $(\frak A,\bar\mu)$ be an atomless
quasi-homogeneous semi-finite measure
algebra.   Show that $\AM(\frak A,\bar\mu,\gamma)$ is homogeneous whenever
$0<\gamma<\bar\mu 1$.   \Hint{first check that
$\frak A\cong\frak A_{1\Bsetminus a}$ whenever $a\in\frak A$ and
$0<\bar\mu a<\bar\mu 1$.}
%528D

\spheader 528Xc(i)
Let $(\frak A,\bar\mu)$ be a totally finite measure algebra.
Show that $\AM(\frak A,\bar\mu,\bar\mu 1)$ is isomorphic to $\frak A$.
(ii) Let $(\frak A,\bar\mu)$ be an atomless measure algebra and
$e\in\frak A$ a non-zero element of finite measure.   Show that the
principal ideal $\frak A_e$ can be regularly embedded in
$\AM(\frak A,\bar\mu,\bar\mu e)$.
%528Fa

\spheader 528Xd\dvArevised{2011}
Show that if $(\frak A,\bar\mu)$ is a probability algebra,
$0<\gamma\le 1$ and $\kappa\ge\max(\omega,\tau(\frak A))$ then
$\AM(\frak A,\bar\mu,\gamma)$ can be regularly embedded in
$\AM(\frak B_{\kappa},\bar\nu_{\kappa},\gamma)$.
%528H 528Xb mt52bits

\spheader 528Xe Let $(X,\frak T,\Sigma,\mu)$ be a quasi-Radon measure space
and $(\frak A,\bar\mu)$ its amoeba algebra.   Show that if $0<\gamma<\mu X$
then the additivity of
$\mu$ is not a precaliber of $\AM(\frak A,\bar\mu,\gamma)$.
%528L out of order query

\spheader 528Xf
Let $(\frak A,\bar\mu)$ be an atomless
$\sigma$-finite measure algebra
and $0<\gamma<\bar\mu 1$.   Show that
$\frak m(\AM(\frak A,\bar\mu,\gamma))=\wdistr(\frak A)$.
%528L for \le 517If & 528N for \ge;  mt52bits
%528N

\spheader 528Xg Let $(\frak A,\bar\mu)$ be an atomless
semi-finite measure algebra.   (i) Show that

\Centerline{$c(\AM^*(\frak A,\bar\mu))
=\link_m(\AM^*(\frak A,\bar\mu))=\max(c(\frak A),\tau(\frak A))$}

\noindent for any integer $m\ge 2$.   (ii) Show that

\Centerline{$d(\AM^*(\frak A,\bar\mu))
=\pi(\AM^*(\frak A,\bar\mu))
=\max(\cff[c(\frak A)]^{\le\omega},\pi(\frak A))$.}
%528P

\spheader 528Xh\dvAnew{2011}
Show that for any cardinal $\kappa$ there is a
probability algebra $(\frak A,\bar\mu)$ such that
$\AM(\frak A,\bar\mu,\bover12)$ has Maharam type $\kappa$.
%528V 528Xc atom of measure $\bover12$

\leader{528Y}{Further exercises (a)}
%\spheader 528Ya
Let $(\frak A,\bar\mu)$ be an atomless quasi-homogeneous
semi-finite measure algebra.   Show that
$\AM^*(\frak A,\bar\mu)$ is homogeneous.
%528Xb 528D

\spheader 528Yb\dvAnew{2011} Let $(\frak A,\bar\mu)$ be an atomless
totally finite
measure algebra, and suppose that $\AM(\frak A,\bar\mu,\gamma)$ can be
regularly embedded in $\AM^*(\frak A,\bar\mu)$ for every
$\gamma\in\ooint{0,\bar\mu 1}$.   Show that $\frak A$ is homogeneous.
%528F out of order query mt52bits

\spheader 528Yc\dvAnew{2011} Show that $\frak B_{\omega_1}$ cannot
be regularly embedded in
$\AM(\frak B_{\omega},\bar\nu_{\omega},\bover12)$.
%528K mt52bits

\spheader 528Yd
 Let $(\frak A,\bar\mu)$ be an atomless
probability algebra and $\gamma\in\ooint{0,1}$.   Show that
$\AM(\frak A,\bar\mu,\gamma)$ is not \wsid.
%528L

\spheader 528Ye Let $\kappa$ be an infinite cardinal.   Show that
(i) $\pi(\RO^{\uparrow}(\Cal S_{\kappa}^{\infty}))
=\cf\Cal S_{\kappa}^{\infty}$
is the cardinal power $\kappa^{\omega}$;   (ii) for every $m\ge 2$,

\Centerline{$c(\RO^{\uparrow}(\Cal S_{\kappa}^{\infty}))
=c^{\uparrow}(\Cal S_{\kappa}^{\infty})
=\link_m(\RO^{\uparrow}(\Cal S_{\kappa}^{\infty}))
=\link^{\uparrow}_m(\Cal S_{\kappa}^{\infty})=\kappa$;}

\noindent (iii) $d(\RO^{\uparrow}(\Cal S_{\kappa}^{\infty}))
=\duparrow(\Cal S_{\kappa}^{\infty})
=\max(\cf\Cal N,\cff[\kappa]^{\le\omega})$.
%528Q

\spheader 528Yf\dvArevised{2011}
Let $(\frak A,\bar\mu)$ be a purely atomic semi-finite
measure algebra of cellularity
at most $\frak c$, and $0<\gamma<\bar\mu 1$.   Show that
$\AM(\frak A,\bar\mu,\gamma)$ has countable Maharam type.
%528R

\spheader 528Yg\dvAnew{2011} Let $(\frak A,\bar\mu)$ be an atomless
semi-finite measure algebra and $0<\gamma<\bar\mu 1$.   Set
$\kappa=\max(\omega,c(\frak A),\tau(\frak A))$ and
$P=\{a:a\in\frak A$, $\bar\mu a<\gamma\}$;  let $\Bbb P$ be the forcing
notion $(P,\Bsubseteqshort,0,\uparrow)$ (see 5A3A).   Show that
$\VVdash_{\Bbb P}\check\kappa<\omega_1$.
%528V

\spheader 528Yh\dvAnew{2011}
Show that if $(\frak A,\bar\mu)$ is a measure
algebra with at most $\frak c$ atoms, then
$\tau(\AM^*(\frak A,\bar\mu))\le\omega$.
%528V

\leaveitout{\spheader 5{}28Yc For $\alpha\in\NN$ and any set $I$
let $(\Cal S^{\infty}_I)^{(\alpha)}$
be the partially ordered set

\Centerline{$\{p:p\subseteq\Bbb N\times I$,
$\#(p[\{n\}])\le\alpha(n)$ for every $n$,
$\sup_{n\in\Bbb N}\#(p[\{n\}])<\infty\}$.}

\noindent Show that $\RO^{\uparrow}(\Cal S^{\infty}_I)^{(\alpha)}$ can be
regularly embedded in
$\RO^{\uparrow}(\Cal S^{\infty}_I)^{(\beta)}$ whenever $I$ is infinite and
$\alpha$, $\beta$ are both unbounded.

This certainly is wrong.   But what if both $\alpha$ and $\beta$ tend to
$\infty$?
}
}%end of exercises

\leader{528Z}{Problems (a)}
%\spheader 528Za
Let $(\frak A_L,\bar\mu_L)$ be the measure
algebra of Lebesgue measure on $\Bbb R$.   Is the amoeba algebra
$\AM(\frak A_L,\bar\mu_L,1)$ isomorphic to the amoeba algebra
$\AM(\frak B_{\omega},\bar\nu_{\omega},\bover12)$?
%528K

\spheader 528Zb Let $(\frak A,\bar\mu)$ be a probability algebra, $\frak B$
a closed subalgebra of $\frak A$, and $0<\gamma<1$.   Is it
necessarily true that
$\AM(\frak B,\bar\mu\restrp\frak B,\gamma)$ can be regularly embedded in
$\AM(\frak A,\bar\mu,\gamma)$?   \cmmnt{(See 528Xd and 528G.)}
%528G 528K

%Can the category algebra \frak G_{\omega} be regularly embedded in
% \AM(\frak B_{\omega},\bar\nu_{\omega},\bover12) ?

\endnotes{
\Notesheader{528} The ideas of
528A-528K %528A 528D 528E 528F 528H 528J 528K
are based on {\smc Truss 88}.   The original amoeba algebras of
{\smc Martin \& Solovay 70}, used in their proof that
$\add\Cal N\ge\frak m$ (528L), are closest to 528C.   For some more
about the amoeba algebras derived from Lebesgue measure, see
{\smc Bartoszy\'nski \& Judah 95}, \S3.4.
In this section I have been willing
to assume that the measure algebras involved are atomless;  amoeba
algebras are surely still interesting for other measure algebras, but
the new questions seem to be combinatoric rather than measure-theoretic.
It seems still to be unknown whether the algebras
$\AM(\frak A_L,\bar\mu_L,1)$ and
$\AM(\frak B_{\omega},\bar\nu_{\omega},\bover12)$ are actually
isomorphic, rather than just mutually embeddable (528K, 528Za).

If we think of the partially ordered sets of 528A and 528I as forcing
notions, we can study them in terms of the forcing universes they
lead to.   This is associated with the
prominence of `regular embeddings' in this section.
I will not attempt to use such methods here, but I mention
them because results such as 528Yg have been part of the impulse for
studying amoeba algebras, and led naturally to 528Ya, 528R and 528U-528V.
}%end of notes

\discrpage

\frfilename{mt529.tex}
\versiondate{26.5.11}
\copyrightdate{2003}

\def\chaptername{Cardinal functions of measure theory}
\def\sectionname{Further partially ordered sets of measure theory}

\def\Fn{\mathop{\text{Fn}}\nolimits}

\newsection{529}

I end the chapter with notes on some more structures which can be
approached by the methods used earlier.   The Banach lattices of Chapter
36 are of course partially ordered sets, and many of them can easily be
assigned places in the Tukey classification
(529C, 529D\cmmnt{, 529Xa}).   More surprising is the fact that the
Nov\'ak numbers of $\{0,1\}^I$, for large $I$, are supported by the
additivity of Lebesgue measure (529F);   this is associated with an
interesting property of the localization poset from the last section
(529E).   There is a similarly unexpected connexion between the covering
number of Lebesgue measure and `reaping numbers'
$\frak r(\omega_1,\lambda)$ for large $\lambda$ (529H).

\leader{529A}{Notation} \cmmnt{As in previous sections,} I will write
$\Cal N(\mu)$ for the null ideal of $\mu$ in a measure space
$(X,\Sigma,\mu)$, and $\Cal N$ for the null ideal of Lebesgue measure on
$\Bbb R$.

\leader{529B}{Proposition} Let $(\frak A,\bar\mu)$ be a semi-finite
measure algebra.

(a) For $p\in\coint{1,\infty}$, give
$L^p=L^p(\frak A,\bar\mu)$\cmmnt{ (definition: 366A)} its norm
topology.   Then its topological density is

$$\eqalign{d(L^p)
&=1\text{ if }\frak A=\{0\},\cr
&=\omega\text{ if }0<\#(\frak A)<\omega,\cr
&=\max(c(\frak A),\tau(\frak A))
  \text{ if }\frak A\text{ is infinite.}\cr}$$

(b) Give $L^0=L^0(\frak A)$ its topology of convergence in
measure\cmmnt{ (367L)}.   Then

$$\eqalign{d(L^0)
&=1\text{ if }\frak A=\{0\},\cr
&=\omega\text{ if }0<\#(\frak A)<\omega,\cr
&=\tau(\frak A)\text{ if }\frak A\text{ is infinite.}\cr}$$

\proof{{\bf (a)(i)} The case in which $\frak A$ is finite is elementary,
since in this case $L^p\cong\BbbR^n$, where $n$ is the
number of atoms of $\frak A$.   So henceforth let us suppose that
$\frak A$ is infinite.

\medskip

\quad{\bf (ii)} If $\frak A^f$ is the set of elements of $\frak A$ of
finite
measure, we have a natural injection $a\mapsto\chi a:\frak A^f\to L^p$, and
$\|\chi a-\chi b\|_p=\mu(a\Bsymmdiff b)^{1/p}$, so $\chi$ is a
homeomorphism
for the measure metric on $\frak A^f$ and the norm metric on $L^p$.
It follows that the density $d(A)$ of $A=\{\chi a:a\in\frak A^f\}$ for the
norm topology is equal to the density of $\frak A^f$ for the strong
measure-algebra topology, which is $\max(c(\frak A),\tau(\frak A))$, by
521Eb.   So

\Centerline{$\max(c(\frak A),\tau(\frak A))=d(A)\le d(L^p)$}

\noindent by 5A4B(h-ii).
In the other direction, if $A_0$ is a dense subset of $A$
with cardinal $d(A)$ and $D$ is the set of rational
linear combinations of members of $A_0$, $\overline{D}\supseteq S(\frak A^f)$
is dense in $L^p$ (366C), so

\Centerline{$d(L^p)\le\#(D)\le\max(\omega,\#(A_0))
=\max(c(\frak A),\tau(\frak A))$.}

\medskip

{\bf (b)} Again, the case of finite $\frak A$ is trivial, so we need
consider only infinite $\frak A$.   In this case, $\tau(\frak A)$ is
equal to the topological density $d_{\frak T}(\frak A)$
of $\frak A$ with its measure-algebra topology $\frak T$ (521Ea).

\medskip

\quad{\bf (i)} Let $A\subseteq\frak A$ be a topologically dense set of
cardinal $\tau(\frak A)$.   Set

\Centerline{$D=\{\sum_{i=0}^nq_i\chi a_i:q_0,\ldots,q_n\in\Bbb Q$,
$a_0,\ldots,a_n\in\frak A\}$,}

\noindent so that $D\subseteq L^0$ has cardinal $\tau(\frak A)$.
Because $a\mapsto\chi a:\frak A\to L^0$ is continuous (367Ra), the
closure $\overline{D}$ of $D$ includes $\{\chi a:a\in\frak A\}$.
Because $\overline{D}$ is a linear subspace of $L^0$, it includes
$S(\frak A)$.   Because $S(\frak A)$ is dense in $L^0$
(367Nc),
$\overline{D}=L^0$ and $d(L^0)\le\#(D)=\tau(\frak A)$.

\medskip

\quad{\bf (ii)} Let $B\subseteq L^0$ be a dense set with cardinal
$d(L^0)$.   Set

\Centerline{$A=\{\Bvalue{u>\bover12}:u\in B\}$,}

\noindent so that $A\subseteq\frak A$ and $\#(A)\le d(L^0)$.   Then $A$
is topologically dense in $\frak A$.   \Prf\ If $c\in\frak A$,
$a\in\frak A^f$ and $\epsilon>0$, there is a $u\in B$ such that
$\int|u-\chi c|\wedge\chi a\le\bover12\epsilon$.   But in this case,
setting $b=\Bvalue{u>\bover12}$, $|\chi(b\Bsymmdiff c)|\le 2|u-\chi c|$,
so

\Centerline{$\bar\mu(a\Bcap(b\symmdiff c))
\le 2\int|u-\chi c|\wedge\chi a\le\epsilon$.}

\noindent As $c$, $a$ and $\epsilon$ are arbitrary, $A$ is topologically
dense in $\frak A$.\ \Qed

Accordingly

\Centerline{$\tau(\frak A)=d_{\frak T}(\frak A)\le\#(A)\le d(L^0)$}

\noindent and $d(L^0)=\tau(\frak A)$, as claimed.
}%end of proof of 529B

\leader{529C}{Theorem}\cmmnt{ ({\smc Fremlin 91})} Let $U$ be an
$L$-space.   Then $U\equivT\ell^1(\kappa)$, where $\kappa=\dim U$ if $U$
is finite-dimensional, and otherwise is the topological density of $U$.

\proof{{\bf (a)} The finite-dimensional case is trivial, since in this
case $U$ and $\ell^1(\kappa)$ are isomorphic as Banach lattices.   So
henceforth let us suppose that $U$ is infinite-dimensional.   Now
$\vee:U\times U\to U$ is uniformly continuous.
\Prf\ We have only to observe that $u\vee v=\bover12(u+v+|u-v|)$ in any
Riesz space, that addition and subtraction are uniformly continuous in
any linear topological space, and that $u\mapsto|u|$ is uniformly
continuous just because $||u|-|v||\le|u-v|$ (see 354B).\ \QeD\
So 524C, with $Q=P=U$, tells us that $U\prT\ell^1(\kappa)$.

The rest of the proof will therefore be devoted to showing that
$\ell^1(\kappa)\prT U$.

\medskip

{\bf (b)} By Kakutani's theorem (369E), there is a localizable measure
algebra $(\frak A,\bar\mu)$ such that $U$ is isomorphic, as Banach
lattice, to $L^1(\frak A,\bar\mu)$.   Let $\familyiI{a_i}$ be a
partition of unity in $\frak A$ such that $0<\bar\mu a_i<\infty$ and the
principal ideal $\frak A_{a_i}$ is homogeneous for each $i$.
Set $\kappa_i=\tau(\frak A_{a_i})$, so that $\kappa_i$ is either $0$ or
infinite for every $i$, and $\kappa=\max(\#(I),\sup_{i\in I}\kappa_i)$
by 529Ba.

It will simplify the calculations to follow if we arrange that all the
$a_i$ have measure $1$.   To do this, set
$\bar\nu a=\sum_{i\in I}\Bover{\bar\mu(a\Bcap a_i)}{\bar\mu a_i}$ for
$a\in\frak A$;  that is, $\bar\nu a=\int_aw\,d\bar\mu$, where
$w=\sup_{i\in I}\Bover1{\bar\mu a_i}\chi a_i$ in $L^0(\frak A)$.   In this
case, $\int v\,d\bar\nu=\int v\times w\,d\bar\nu$ for every
$v\in L^1(\frak A,\bar\nu)$, while
$\int u\,d\bar\mu=\int\Bover1w\times u\,d\bar\nu$ for every
$u\in L^1(\frak A,\bar\mu)$ (365S\formerly{3{}65T}).   But this means that
$u\mapsto u\times\Bover1w$ is a Banach lattice isomorphism between
$L^1(\frak A,\bar\mu)$ and $L^1(\frak A,\bar\nu)$, and $U$ is
isomorphic, as $L$-space, to $L^1=L^1(\frak A,\bar\nu)$;  while
$\bar\nu a_i=1$ for every $i$.

\medskip

{\bf (c)} There are a set $J$, with cardinal $\kappa$, and a family
$\family{j}{J}{u_j}$ in $L^1$ such that $\#(J)=\kappa$, $\|u_j\|\le 2$
for every $j\in J$ and
$\|\sup_{j\in K}u_j\|\ge\bover12\sqrt{\#(K)}$ for every finite
$K\subseteq J$.   \Prf\ Set

\Centerline{$J=\{(i,0):i\in I$, $\kappa_i=0\}
  \cup\{(i,\xi):i\in I$, $\xi<\kappa_i\}$.}

\noindent Then $\#(J)=\kappa$.   If $i\in I$ and $\kappa_i=0$, set
$u_{(i,0)}=\chi a_i$.   If $i\in I$ and $\kappa_i>0$, then
$(\frak A_{a_i},\bar\nu\restrp\frak A_{a_i})$ is a homogeneous
probability algebra with Maharam type $\kappa_i\ge\omega$, so is
isomorphic to the measure algebra $(\frak C_i,\bar\lambda_i)$ of
$\ocint{0,1}^{\kappa_i}$ with its usual measure $\lambda_i$, the product
of Lebesgue measure on each copy of $\ocint{0,1}$ (334E).   For
$\xi<\kappa_i$, set

\Centerline{$h_{i\xi}(t)=\Bover1{\sqrt{t(\xi)}}$ for
$t\in\ocint{0,1}^{\kappa_i}$,}

\noindent and let $u_{(i,\xi)}\in L^1$ correspond to
$h_{i\xi}^{\ssbullet}\in L^1(\lambda_i)
\cong L^1(\frak C_i,\bar\lambda_i)$ (365B).   Of course

$$\eqalignno{\|u_{(i,\xi)}\|
&=\int h_{i\xi}(t)\lambda_i(dt)
=\int_0^1\Bover1{\sqrt{\alpha}}d\alpha\cr
\displaycause{because the coordinate map $t\mapsto t(\xi)$ is \imp}
&=2.\cr}$$

If $L\subseteq\kappa_i$ is finite and not empty, then
$\|\sup_{\xi\in L}u_{(i,\xi)}\|=\int g\,d\lambda_i$ where
$g=\sup_{\xi\in L}h_{i\xi}$, that is,
$g(t)=\sup_{\xi\in L}\Bover1{\sqrt{t(\xi)}}$ for
$t\in\ocint{0,1}^{\kappa_i}$.
Now, for any $\alpha\ge 1$,

$$\eqalignno{\lambda_i\{t:g(t)\le\alpha\}
&=\lambda_i\{t:\alpha^2 t(\xi)\ge 1\text{ for every }\xi\in L\}\cr
&=(1-\Bover1{\alpha^2})^{\#(L)}
\le\max(\Bover12,1-\Bover1{2\alpha^2}\#(L))\cr
\displaycause{induce on $\#(L)$}
&=1-\Bover1{2\alpha^2}\#(L)\text{ if }\alpha\ge\sqrt{\#(L)}.\cr}$$

\noindent So

$$\eqalignno{\|\sup_{\xi\in L}u_{(i,\xi)}\|
&=\int g\,d\lambda_i
=\int_0^{\infty}\lambda_i\{t:g(t)>\alpha\}d\alpha\cr
\displaycause{252O}
&\ge\int_{\sqrt{\#(L)}}^{\infty}\lambda_i\{t:g(t)>\alpha\}d\alpha\cr
&\ge\int_{\sqrt{\#(L)}}^{\infty}\Bover1{2\alpha^2}\#(L)d\alpha
=\Bover12\sqrt{\#(L)}.\cr}$$

What this means is that if $K\subseteq J$ is finite and all the first
coordinates of members of $K$ are the same, then
$\|\sup_{j\in K}u_j\|\ge\bover12\sqrt{\#(K)}$.   In general, if
$K\subseteq J$ is finite, then for each $i\in I$ set
$L_i=\{\xi:(i,\xi)\in K\}$.   Set $v_i=0$ if $L_i$ is empty,
$\sup_{\xi\in L_i}u_{(i,\xi)}$ otherwise, so that
$\|v_i\|\ge\bover12\sqrt{\#(L_i)}$;  now
$\sup_{j\in K}u_j=\sum_{i\in I}v_i$, so

\Centerline{$\|\sup_{j\in K}u_j\|=\sum_{i\in I}\|v_i\|
\ge\Bover12\sum_{i\in I}\sqrt{\#(L_i)}
\ge\Bover12\sqrt{\sumop_{i\in I}\#(L_i)}=\Bover12\sqrt{\#(K)}$,}

\noindent as required.\ \Qed

\medskip

{\bf (d)} We can now apply the idea of the proof of 524C, as follows.
The density of $\ell^1(\kappa)$ is of course $\kappa$, by 529Ba applied
to counting measure on $\kappa$, or otherwise.
Index a dense subset of $\ell^1(\kappa)$ as
$\family{j}{J}{y_j}$.   For each $x\in\ell^1$, choose a sequence
$\sequencen{k(x,n)}$ in $J$ such that

\Centerline{$\|x-\sum_{m=0}^ny_{k(x,m)}\|\le 8^{-n}$}

\noindent for every $n$.   Note that

\Centerline{$\|y_{k(x,n)}\|
\le\|x-\sum_{m=0}^ny_{k(x,m)}\|+\|x-\sum_{m=0}^{n-1}y_{k(x,m)}\|
\le 9\cdot 8^{-n}$}

\noindent for each $n$.
Choose $f(x)\in L^1$ such that $\|f(x)\|\ge\|x\|$ and
$f(x)\ge\sum_{n=0}^{\infty}2^{-n}u_{k(x,n)}$;
this is possible because $\{u_{k(x,n)}:n\in\Bbb N\}$ is bounded.

\medskip

{\bf (e)} $f$ is a Tukey function.   \Prf\ Take $v\in L^1$ and set

\Centerline{$A=\{x:f(x)\le v\}$,
\quad$K_n=\{k(x,n):x\in A\}$}

\noindent for $n\in\Bbb N$.   Fix $n$ for the moment.   If $j\in K_n$,
then there is an $x\in A$ such that $j=k(x,n)$ and

\Centerline{$u_j=u_{k(x,n)}\le 2^nf(x)\le 2^nv$,}

\noindent while $\|y_j\|=\|y_{k(x,n)}\|\le 9\cdot 8^{-n}$.   If
$K\subseteq K_n$ is finite, $\|2^nv\|\ge\bover12\sqrt{\#(K)}$, by (c);
so $\#(K_n)\le 2^{2n+2}\|v\|^2$.   This means that if we set
$z_n=\sum_{j\in K_n}|y_j|$ we shall have
$\|z_n\|\le 9\cdot 8^{-n}\#(K_n)\le 36\cdot 2^{-n}\|v\|^2$, while
$y_{k(x,n)}\le z_n$ for every $x\in A$.

Now $z=\sum_{n=0}^{\infty}z_n$ is defined in $\ell^1(\kappa)$, and if
$x\in A$ then $\sum_{m=0}^ny_{k(x,m)}\le z$ for every $n\in\Bbb N$, so
that $x\le z$.   Thus $A$ is bounded above in $\ell^1(\kappa)$.   As $v$
is arbitrary, $f$ is a Tukey function.\ \Qed

\medskip

{\bf (f)} Accordingly $\ell^1(\kappa)\prT L^1\cong U$, and the proof is
complete.
}%end of proof of 529C

\leader{529D}{Theorem}\cmmnt{ ({\smc Fremlin 91})} Let $\frak A$ be a
homogeneous measurable algebra with Maharam type $\kappa\ge\omega$.  Then
$L^0(\frak A)\equivT\ell^1(\kappa)$.

\proof{{\bf (a)} Let $\bar\mu$ be such that $(\frak A,\bar\mu)$ is a
probability algebra.   If we give $L^0=L^0(\frak A)$ its topology of
convergence in measure, its density is $\kappa$, by 529Bb.   Moreover,
this topology is defined by the metric
$(u,v)\mapsto\int|u-v|\wedge\chi 1$, under which the lattice operation
$\vee$ is uniformly continuous.   \Prf\ Just as in part (a) of the proof
of 529C, we have $u\vee v=\bover12(u+v+|u-v|)$ for all $u$ and $v$,
addition and subtraction are uniformly continuous, and $u\mapsto|u|$ is
uniformly continuous.\ \QeD\   So, just as in
529C, we can use 524C to see that $L^0\prT\ell^1(\kappa)$.

\medskip

{\bf (b)} For the reverse connection, I repeat ideas from the proof
of 529C.   $(\frak A,\bar\mu)$ is isomorphic to the measure algebra
$(\frak C,\bar\lambda)$ of $\ocint{0,1}^{\kappa}$ with its usual measure
$\lambda$.    For $\xi<\kappa$ and $t\in\ocint{0,1}^{\kappa}$ set
$h_{\xi}(t)=\Bover1{\sqrt{t(\xi)}}$, and set
$u_{\xi}=h_{\xi}^{\ssbullet}$ in $L^0(\lambda)\cong L^0(\frak C)$
(364Ic).
This time, observe that if $x\in\ell^1(\kappa)^+$ and
$\alpha\ge\sqrt{\|x\|}$ then

$$\eqalignno{\bar\lambda\Bvalue{\sup_{\xi<\kappa}\sqrt{x(\xi)}u_{\xi}
  \le\alpha}
&=\bar\lambda(\inf_{\xi<\kappa}\Bvalue{\sqrt{x(\xi)}u_{\xi}\le\alpha})
=\prod_{\xi<\kappa}\bar\lambda\Bvalue{\sqrt{x(\xi)}u_{\xi}\le\alpha}\cr
&=\prod_{\xi<\kappa}\lambda\{t:\sqrt{\Bover{x(\xi)}{t(\xi)}}\le\alpha\}
=\prod_{\xi<\kappa}(1-\Bover{x(\xi)}{\alpha^2})\cr
&\ge 1-\Bover1{\alpha^2}\sum_{\xi<\kappa}x(\xi)
\to 0\cr}$$

\noindent as $\alpha\to\infty$.   This means that
$\sup_{\xi<\kappa}\sqrt{x(\xi)}u_{\xi}$ is defined in
$L^0(\lambda)$ (364La).   So we can define
$f:\ell^1(\kappa)\to L^0(\lambda)$ by saying that
$f(x)=\sup_{\xi<\kappa}\sqrt{\max(0,x(\xi))}u_{\xi}$ for every
$x\in\ell^1(\kappa)$.

\medskip

{\bf (c)} $f$ is a Tukey function.   \Prf\ Take $v\in L^0(\lambda)^+$,
and set $A=\{x:f(x)\le v\}$.    Note that
$f(x\vee x')=f(x)\vee f(x')$ for all $x$, $x'\in\ell^1(\kappa)$, so $A$
is upwards-directed.   Take $\alpha>0$ such that
$\bar\lambda\Bvalue{v\le\alpha}=\beta>\bover12$.
If $x\in A$ and $x\ge 0$
then $f(x)\ge\sqrt{x(\xi)}\chi 1$ so $x(\xi)\le\alpha$ for every
$\xi$.   Now the calculation in (b) tells us that

$$\eqalignno{\beta
&\le\bar\lambda\Bvalue{\sup_{\xi<\kappa}\sqrt{x(\xi)}u_{\xi}
  \le\alpha}
=\prod_{\xi<\kappa}(1-\Bover1{\alpha^2}x(\xi))\cr
&\le\max(\Bover12,1-\Bover1{2\alpha^2}\sum_{\xi<\kappa}x(\xi))
=\max(\Bover12,1-\Bover1{2\alpha^2}\|x\|),\cr}$$

\noindent so $\|x\|\le 2\alpha^2(1-\beta)$.   As $A$ is
upwards-directed and norm-bounded and contains $0$, it is bounded above in
$\ell^1(\kappa)$ (354N).   As $v$ is arbitrary, $f$ is a Tukey
function.\ \Qed

\medskip

{\bf (d)} Accordingly
$\ell^1(\kappa)\prT L^0(\lambda)\cong L^0(\frak A)$ and
$\ell^1(\kappa)$ and $L^0(\frak A)$ are Tukey equivalent.
}%end of proof of 529D

\leader{529E}{Proposition} Let $\Cal S^{\infty}$ be the
$(\Bbb N,\infty)$-localization
poset\cmmnt{ (528I)}.   Then $\RO(\{0,1\}^{\frakc})$ can be
regularly embedded in $\RO^{\uparrow}(\Cal S^{\infty})$.

\woddheader{529E}{0}{0}{0}{60pt}

\proof{{\bf (a)} Let $\ofamily{\xi}{\frakc}{h_{\xi}}$ be a family of
eventually-different functions in $\NN$ (5A1Nc).   Set

$$\eqalign{W_0
&=\bigcup_{n\in\Bbb N\text{ is even}}
  \{(h,p):h\in\NN,\,p\in\Cal S^{\infty},\,\#(p[\{n\}])=2^n,\cr
&\mskip180mu (n,h(n))\notin p,\,(i,h(i))\in p
  \text{ for every }i>n\}\cr
&\mskip90mu \cup\{(h,p):h\in\NN,\,p\in\Cal S^{\infty},\,(i,h(i))\in p
  \text{ for every }i\in\Bbb N\},\cr
W_1
&=\bigcup_{n\in\Bbb N\text{ is odd}}\{(h,p):h\in\NN,
\,p\in\Cal S^{\infty},\,\#(p[\{n\}])=2^n,\cr
&\mskip180mu (n,h(n))\notin p,\,(i,h(i))\in p
  \text{ for every }i>n\}.\cr}$$

\noindent Observe that (i) $W_0\cap W_1=\emptyset$ (ii) if
$(h,p)\in W_j$, where $j=0$ or $j=1$, and
$p\subseteq q\in\Cal S^{\infty}$ then
$(h,q)\in W_j$ (iii) if $p\in\Cal S^{\infty}$ then

\Centerline{$\#(\{\xi:(h_{\xi},p)\in W_0\cup W_1\})
\le\|p\|$}

\noindent is finite.

\medskip

{\bf (b)} Set $Q=\Fn_{<\omega}(\frak c,\{0,1\})$, the set of functions from
finite subsets of $\frak c$ to $\{0,1\}$, ordered
by extension of functions, so that $(Q,\subseteq)$ is isomorphic to
$(\Cal U,\supseteq)$ where
$\Cal U$ is the usual base of the topology of $\{0,1\}^{\frakc}$, and
$\RO^{\uparrow}(Q)\cong\RO^{\downarrow}(\Cal U)$
can be identified with the regular open algebra of $\{0,1\}^{\frak c}$
(514Sd).   Define $f:\Cal S^{\infty}\to Q$ by setting $f(p)(\xi)=j$ if
$(h_{\xi},p)\in W_j$.   Then $f$ is order-preserving.

\medskip

{\bf (c)} For $p\in\Cal S^{\infty}$, set

\Centerline{$A_0(p)
=\{\xi:\xi<\frak c$, $\{n:n$ is even, $(n,h_{\xi}(n))\notin p\}$ is
finite$\}$,}

\Centerline{$A_1(p)
=\{\xi:\xi<\frak c$, $\{n:n$ is odd, $(n,h_{\xi}(n))\notin p\}$ is
finite$\}$,}

\Centerline{$A(p)=A_0(p)\cup A_1(p)$,}

\noindent so that $A(p)$ is finite and $\dom f(p)\subseteq A(p)$.   Now
$P_1=\{p:p\in\Cal S^{\infty}$, $A(p)=\dom f(p)\}$ is cofinal with $\Cal
S^{\infty}$.   \Prf\
Take $p\in\Cal S^{\infty}$.   Let $m$ be such that

\Centerline{$2^m\ge\|p\|+\#(A(p))$,}

\Centerline{$(n,h_{\xi}(n))\in p$ whenever
$\xi\in A_0(p)$ and $n>m$ is even,}

\Centerline{$(n,h_{\xi}(n))\in p$ whenever
$\xi\in A_1(p)$ and $n>m$ is odd.}

\noindent Let $p'\in\Cal S^{\infty}$ be such that

\Centerline{for $n\le m$, $p'[\{n\}]\supseteq p[\{n\}]$ and
$\#(p'[\{n\}])=2^n$,}

\Centerline{for $n>m$,
$p'[\{n\}]=p[\{n\}]\cup\{h_{\xi}(n):\xi\in A(p)\}$.}

\noindent Then $p\le p'$ and $A(p')=A(p)$.   Also
$A(p)=\dom f(p')$, because if $\xi\in A(p)$ then either
$(n,h_{\xi}(n))\in p'$ for every $n$ and $(h_{\xi},p')\in W_0$, or there
is a largest $n$ such that $(n,h_{\xi}(n))\notin p'$, in which case
$n\le m$ and $\#(p'[\{n\}])=2^n$, so $(h_{\xi},p')$ belongs to $W_0$ if
$n$ is even and $W_1$ otherwise.\ \Qed

\medskip

{\bf (d)} If $p\in P_1$ and $q\in Q$ extends $f(p)$, there is a
$p_1\in\Cal S^{\infty}$ such that $p_1\supseteq p$ and $f(p_1)=q$.
\Prf\ Let $m$ be such that
$2^m\ge\|p\|+\#(\dom q)$ and $h_{\xi}(n)\ne h_{\eta}(n)$ whenever $\xi$,
$\eta\in\dom q$ are distinct and $n\ge m$.
For each $\xi\in\dom q\setminus\dom f(p)=\dom q\setminus A(p)$,
$\{n:n$ is even, $(n,h_{\xi}(n))\notin p\}$
and $\{n:n$ is odd, $(n,h_{\xi}(n))\notin p\}$ are both infinite.
So we can find $m'\ge m$ such that all these sets meet
$m'\setminus m$.   Set

$$\eqalign{p'
&=p\cup\{(n,h_{\xi}(n)):n\in m'\setminus m
  \text{ is odd},\,q(\xi)=0\}\cr
&\mskip200mu
\cup\{(n,h_{\xi}(n)):n\in m'\setminus m\text{ is even},\,q(\xi)=1\}\cr
&\mskip200mu
\cup\{(n,h_{\xi}(n)):n\in\Bbb N\setminus m',\,\xi\in\dom q\},
\cr}$$

\noindent so that $p\subseteq p'\in\Cal S^{\infty}$.   Let
$p_1\in\Cal S^{\infty}$ be such that $p_1\supseteq p'$,
$p_1\setminus p'$ is finite, $\#(p_1[\{n\}])=2^n$ for every $n<m'$ and
$(n,h_{\xi}(n))\notin p_1\setminus p'$ whenever $n\in\Bbb N$ and
$\xi\in\dom q$.   Now $f(p_1)=q$, while $p\subseteq p_1$.\ \Qed

\medskip

{\bf (e)} Putting (c) and (d) together, we see that $f^{-1}[Q_0]$ must
be cofinal with $\Cal S^{\infty}$ for every cofinal $Q_0\subseteq Q$;
moreover, since $\emptyset\in P_1$ and $f(\emptyset)$ is the empty
function, $f[\Cal S^{\infty}]=Q$.   So $f$ satisfies the conditions of
514O, and $\RO^{\uparrow}(Q)\cong\RO(\{0,1\}^{\frakc})$ can be regularly
embedded in $\RO^{\uparrow}(\Cal S^{\infty})$.
}%end of proof of 529E

%questions:  can \frak B_{\frak c} be regularly embedded in
%\RO^{\uparrow}(\Cal S^{\infty})?
%can \RO(\Bbb R) be regularly embedded in
%\AM(\frak B_{\omega},\bar\nu_{\omega}\bover12)?

\leader{529F}{Corollary}\cmmnt{ ({\smc Brendle 00}, 2.3.10;
{\smc Brendle 06}, Theorem 1)}
$n(\{0,1\}^I)\ge\add\Cal N$ for every set $I$.

\proof{ If $I$ is finite, this is trivial.   Otherwise, write
$\lambda=n(\{0,1\}^I)$.   Then $\lambda\ge n(\{0,1\}^{\frakc})$.   \Prf\
If $J\subseteq I$ is a countably infinite set, then
$\{\{x:x\restr J=z\}:z\in\{0,1\}^J\}$ is a cover of $\{0,1\}^I$ by
continuum many nowhere dense sets, so $\lambda\le\frak c$.   Let
$\ofamily{\xi}{\lambda}{E_{\xi}}$ be a cover of $\{0,1\}^{\kappa}$ by
nowhere dense sets.   Then each $E_{\xi}$ is included in a nowhere dense
closed set $F_{\xi}$ determined by coordinates in a countable set
$K_{\xi}\subseteq I$ (4A2E(b-iii)).   Set
$K=\bigcup_{\xi<\lambda}K_{\xi}$, so that $\#(K)\le\frak c$.   Then all
the projections $F'_{\xi}=\{x\restr K:x\in F_{\xi}\}$ are nowhere dense
in $\{0,1\}^K$ (apply 4A2B(f-ii) to the continuous open surjections
$x\mapsto x\restr K_{\xi}:\{0,1\}^I\to\{0,1\}^{K_{\xi}}$ and
$y\mapsto y\restr K_{\xi}:\{0,1\}^K\to\{0,1\}^K_{\xi}$),
and they cover $\{0,1\}^K$.
Next, we have an injection $\phi:K\to\frak c$, and the sets
$F''_{\xi}=\{x:x\phi\in F'_{\xi}\}$ form a cover of $\{0,1\}^{\frakc}$
by nowhere dense sets;  so $n(\{0,1\}^{\frakc})\le\lambda$.\ \Qed

Because every non-empty open set $\{0,1\}^{\frakc}$ includes an open set
homeomorphic to $\{0,1\}^{\frakc}$,

$$\eqalignno{n(\{0,1\}^{\frakc})
&=\min\{n(H):H\subseteq\{0,1\}^{\frakc}\text{ is open and not empty}\}\cr
&=\frak m(\RO(\{0,1\}^{\frak c}))\cr
\displaycause{517J}
&\ge\frak m(\RO^{\uparrow}(\Cal S^{\infty}))\cr
\displaycause{where $\Cal S^{\infty}$ is the $(\Bbb N,\infty)$-localization 
poset, by 529E and 517Ia}
&=\add\Cal N\cr}$$

\noindent by 528N.
}%end of proof of 529F

\leader{529G}{Reaping numbers}\cmmnt{ (following
{\smc Brendle 00})} For cardinals $\theta\le\lambda$ let
$\frak r(\theta,\lambda)$ be the smallest cardinal of any set
$\Cal A\subseteq[\lambda]^{\theta}$ such that for every
$B\subseteq\lambda$ there is an $A\in\Cal A$ such that either
$A\subseteq B$ or $A\cap B=\emptyset$.

\leader{529H}{Proposition}\cmmnt{ ({\smc Brendle 00}, 2.7;
{\smc Brendle 06}, Theorem 5)}
$\frak r(\omega_1,\lambda)\ge\cov\Cal N$ for all uncountable
$\lambda$.

\proof{ Let $\ofamily{\xi}{\kappa}{A_{\xi}}$ be a family in
$[\lambda]^{\omega_1}$, where $\kappa<\cov(\Cal N)$.
I seek a $B\subseteq\lambda$ such that $A_{\xi}\cap B$ and
$A_{\xi}\setminus B$ are non-empty for every $\xi<\kappa$.

\medskip

{\bf (a)} If $\kappa\le\omega_1$, then choose
$\ofamily{\xi}{\kappa}{\alpha_{\xi}}$ and
$\ofamily{\xi}{\kappa}{\beta_{\xi}}$ inductively so that

\Centerline{$\alpha_{\xi}\in A_{\xi}\setminus\{\beta_{\eta}:\eta<\xi\}$,
\quad$\beta_{\xi}\in A_{\xi}\setminus\{\alpha_{\eta}:\eta\le\xi\}$}

\noindent for every $\xi<\kappa$, and set
$B=\{\beta_{\xi}:\xi<\kappa\}$;  this serves.   So henceforth let us
suppose that $\kappa>\omega_1$.

\medskip

{\bf (b)} For each $\xi<\kappa$ let $A'_{\xi}\subseteq A_{\xi}$ be a set
of order type $\omega_1$.   For each $n$, let $X_n$ be a set with cardinal
$n!$ with its discrete topology and the uniform probability measure
which gives measure $\Bover1{n!}$ to every singleton.   Give
$X=\prod_{n\in\Bbb N}X_n$ its product measure $\mu$ and its product
topology.   Because $X$ is a compact metrizable space and $\mu$ is a
Radon measure (416U), $\cov\Cal N(\mu)=\cov\Cal N$ (522Wa).
We can therefore choose a family $\ofamily{\xi}{\kappa}{x_{\xi}}$ in $X$
in such a way that each $x_{\zeta}$ is random over its predecessors in
the sense that

\Centerline{whenever $\xi<\kappa$ and
$\overline{\{x_{\eta}:\eta\in A'_{\xi}\cap\zeta\}}$
is negligible, it does not contain $x_{\zeta}$.}

\noindent For distinct $x$, $y\in X$, set
$\Delta(x,y)=\min\{i:x(i)\ne y(i)\}$.   For $x\in X$,
set $B(x)=\{\eta:x_{\eta}\ne x$, $\Delta(x_{\eta},x)$ is even$\}$.

\medskip

{\bf (c)} For every $\xi<\kappa$,
$\{x:x\in X$, $A_{\xi}\subseteq B(x)\}$ and
$\{x:x\in X$, $A_{\xi}\cap B(x)=\emptyset\}$ are
negligible.   \Prf\ There is a $\zeta<\kappa$ such that
$\zeta\in A'_{\xi}$, $A'_{\xi}\cap\zeta$ is countable and
$D=\{x_{\eta}:\eta\in A'_{\xi}\cap\zeta\}$ is dense in
$\{x_{\eta}:\eta\in A'_{\xi}\}$.   Since
$x_{\zeta}\in\overline{D}$, $\overline{D}$ has measure greater than $0$.
By 275I, applied to the sequence $\sequencen{\Sigma_n}$ where $\Sigma_n$ is
the finite algebra of subsets of $X$ determined by coordinates less than
$n$, $\overline{D}$ has a point $w$ which is a density point in the sense
that

\Centerline{$\lim_{n\to\infty}
  \Bover{\mu\{y:y\restr n=w\restr n,\,y\in\overline{D}\}}
{\mu\{y:y\restr n=w\restr n\}}=1$.}

\noindent Consequently, setting

\Centerline{$J_n=\{y(n):y\in D\}
=\{y(n):y\in\overline{D}\}
\supseteq\{y(n):y\in\overline{D}$, $y\restr n=w\restr n\}$,}

\noindent

\Centerline{$\Bover{\#(J_n)}{n!}
\ge\Bover{\mu\{y:y\restr n=w\restr n,\,y\in\overline{D}\}}
{\mu\{y:y\restr n=w\restr n\}}\to 1$}

\noindent as $n\to\infty$.

Next note that, for any $y\in X$ and $n\in\Bbb N$,

\Centerline{$\mu\{x:\,\Exists i>n,\,x(i)=y(i)\}
\le\sum_{i=n+1}^{\infty}\Bover1{i!}\le\Bover{n+2}{(n+1)(n+1)!}$.}

\noindent So

$$\eqalign{\mu\{x:\,\Exists y\in D,\,&x(n)=y(n),\,
  x(i)\ne y(i)\text{ for every }i>n\}\cr
&\mskip100mu\ge\Bover{\#(J_n)}{n!}\bigl(1-\Bover{n+2}{(n+1)(n+1)!}\bigr)
\to 1\cr}$$

\noindent as $n\to\infty$, and

\Centerline{$\mu\{x:\,\Exists y\in D\setminus\{x\},\,
\Delta(x,y)$ is even$\}
=\mu\{x:\,\Exists y\in D\setminus\{x\},\,\Delta(x,y)$ is odd$\}
=1$.}

\noindent But if $y\in D\setminus\{x\}$ and $\Delta(x,y)$ is even, we have an
$\eta\in A_{\xi}$ such that $\Delta(x,x_{\eta})$ is even, and
$\eta\in A_{\xi}\cap B(x)$;  similarly, if there is a
$y\in D\setminus\{x\}$ such that
$\Delta(x,y)$ is odd, there is an $\eta\in A_{\xi}\setminus B(x)$.   So
$\{x:A_{\xi}\subseteq B(x)\}$ and $\{x:A_{\xi}\cap B(x)=\emptyset\}$ are
both negligible.\ \Qed

\medskip

{\bf (d)} Since $\cov\Cal N(\mu)=\cov\Cal N>\kappa$, there is an
$x\in X$ such that both
$A_{\xi}\cap B(x)$ and $A_{\xi}\setminus B(x)$ are non-empty for every
$\xi<\kappa$.   So in this case also we have a suitable set $B$.
}%end of proof of 529H

\exercises{\leader{529X}{Basic exercises (a)}
%\spheader 529Xa
Let $(X,\Sigma,\mu)$ be a measure space, and $p\in\coint{1,\infty}$.
Show that $L^p(\mu)\equivT\ell^1(\kappa)$, where $\kappa=\dim L^p(\mu)$
if this is finite, $d(L^p(\mu))$ otherwise.
%529C

\spheader 529Xb Let $U$ be an $L$-space.   (i) Show that $\add U=\infty$
if $U=\{0\}$, $\omega$ otherwise.   (ii) Show that
$\add_{\omega}U=\infty$ if $U$ is finite-dimensional, $\add\Cal N$ if
$U$ is separable and infinite-dimensional, $\omega_1$ otherwise.
(iii) Show that $\cf U=1$ if $U=\{0\}$, $\omega$ if $0<\dim U<\omega$,
$\max(\cf\Cal N,\cff[d(U)]^{\le\omega})$ otherwise.   (iv) Show that
$\link^{\uparrow}_{<\kappa}(U)=1$ if $\kappa\le\omega$, $\cf U$
otherwise.
%529C

\spheader 529Xc Let $U$ be a separable Banach lattice.   Suppose that
$\ofamily{\xi}{\kappa}{u_{\xi}}$ is a family in $U$, where
$\kappa<\add\Cal N$.   Show that there is a family
$\ofamily{\xi}{\kappa}{\epsilon_{\xi}}$ of strictly positive real numbers
such that $\{\epsilon_{\xi}u_{\xi}:\xi<\kappa\}$ is order-bounded in $U$.
%529C 529Xb out of order query

\sqheader 529Xd Let $U$ be a separable Banach lattice, and $D\subseteq U$
a dense set.   Let
$A\subseteq U$ be a set with cardinal less than $\add\Cal N$.   Show that there
is a $w\in U$ such that for every $u\in A$ and every $\epsilon>0$ there is
a $v\in D$ such that $|u-v|\le\epsilon w$.
%529C out of order query

\spheader 529Xe\dvAformerly{5{}29Xc}
Let $\Cal I$ be the ideal of subsets $I$ of $\Bbb N$
such that $\sum_{n\in I}\Bover1{n+1}$ is finite.  (See 419A.)   Show
that $\ell^1\equivT\Cal I$, so that $\add_{\omega}\Cal I=\add\Cal N$ and
$\cf\Cal I=\cf\Cal N$.
%529+ mt52bits
%note:  analytic P-ideal enough

\spheader 529Xf\dvAnew{2011} Show that if
$\theta\le\theta'\le\lambda'\le\lambda$ are cardinals, then
$\frak r(\theta,\lambda)\le\frak r(\theta',\lambda')$.
%529G

\spheader 529Xg\dvAformerly{5{}29Xd;  revised 2011}(i) Show that
$\frak r(\omega,\omega)\ge\cov\Cal E\ge\max(\cov\Cal N,\frakmctbl)$, where
$\Cal E$ is the ideal of subsets of $\Bbb R$ with Lebesgue negligible
closures.
(ii) Show that if $\lambda$ is an infinite cardinal then
$\frak r(\omega,\lambda)\ge\max(\add\Cal N,\cov\Cal N_{\lambda})$, where
$\Cal N_{\lambda}$ is the null ideal of the usual measure on
$\{0,1\}^{\lambda}$.   \Hint{529F.}
%529H

\leader{529Y}{Further exercises (a)}
%\spheader 529Ya
Let $X$ be a Polish space and $\Cal K_{\sigma}$ the
family of K$_{\sigma}$ subsets of $X$.    Show that, defining $\le^*$ as in
522C,
$(\Cal K_{\sigma},\subseteq)\preccurlyeq_{\text{T}}(\NN,\le^*)$.
%529B

\spheader 529Yb Let $X$ be a topological space with a countable network,
and $c:\Cal PX\to[0,\infty]$ an outer regular submodular Choquet
capacity (definitions:  432J).   Show that if $\Cal A$ is an
upwards-directed family of subsets of $X$ such that
$\#(\Cal A)<\frak m_{\sigma\text{-linked}}$,
then $c(\bigcup\Cal A)=\sup_{A\in\Cal A}c(A)$.
%529B+

\spheader 529Yc Let $r\ge 3$ be an integer.   (i) Let
$c:\Cal P\BbbR^r\to[0,\infty]$ be Choquet-Newton capacity
(\S479).   Show that if $\Cal A$ is an
upwards-directed family of subsets of $\BbbR^r$ such that
$\#(\Cal A)<\add\Cal N$, then $c(\bigcup\Cal A)=\sup_{A\in\Cal A}c(A)$.
\Hint{479Xi.}
(ii) Let $\Cal I$ be the ideal of polar sets in $\BbbR^r$.
Show that $\add\Cal I=\add\Cal N$.
%look at  A\times[0,1]^{r-2} ;  non-zero capacity iff  \mu_1^*A>0
%529D

\spheader 529Yd
Show that, for any infinite set $I$, the regular
open algebra $\RO(\{0,1\}^I)$ of $\{0,1\}^I$ is homogeneous, so that
$\frak m(\RO(\{0,1\}^I))=n(\{0,1\})^I$.
%529E

\spheader 529Ye Show that
$\frak b\le\frak r(\omega,\omega)\le\pi(\Cal P\Bbb N/[\Bbb N]^{<\omega})$.
%529H

}%end of exercises

\endnotes{
\Notesheader{529}
Many of the ideas of the last two chapters were first embodied in
forcing arguments.   In 529E this becomes particularly transparent.
If we have an upwards-directed set $R\subseteq\Cal S^{\infty}$ which is
`generic' in the sense that it meets all the cofinal subsets of
$\Cal S^{\infty}$ definable in a language $\Cal L$ with terms for all
the functions $h_{\xi}$, as well as such obvious ones as
$\{p:\#(p[\{n\}])=2^n\}$ for each $n$, and we set $S=\bigcup\Cal R$,
then $S$ will belong to the set $\Cal S=\Cal S_{\Bbb N}$ of 522K,
and we shall have
$h_{\xi}\subseteq^*S$ for every $\xi$;  so that we have a corresponding
function $\tilde f(S)=\bigcup_{p\in R}f(p)\in\{0,1\}^{\frakc}$
defined by setting

\Centerline{$\tilde f(S)(\xi)
\equiv\sup\{i:(i,h_{\xi}(i))\notin S\}\mod 2$.}

\noindent Next, if $G\subseteq\{0,1\}^{\frakc}$ is a dense open set with
a definition in $\Cal L$, then $\tilde f(S)\in G$;  for,
setting $U_q=\{\phi:q\subseteq\phi\in\{0,1\}^{\frakc}\}$ when
$q\in Q=\Fn_{<\omega}(\frak c,\{0,1\})$, $\{q:U_q\subseteq G\}$ is
cofinal with $Q$, so $\{p:U_{f(p)}\subseteq G\}$ is cofinal with
$\Cal S^{\infty}$ (part (d) of the proof of 529E) and meets $R$.   Thus
$\tilde f(S)$ is `generic' in the sense that it belongs to every dense
open set with a name in $\Cal L$;  and it is a commonplace in the theory
of forcing that a function which transforms generic objects in one
forcing extension into generic objects in another extension corresponds
to a regular embedding of the corresponding regular open algebras.
}%end of notes

\discrpage

\references{
\frfilename{mt5r.tex}
\versiondate{28.4.23}
\copyrightdate{2008}

% except for references in ...X., ...Y., ...Z. give locations in form
%  342L  rather than  342Lb
%  \S511 {\it notes\/}, \S512 {\it intro.}]

\gdef\topparagraph{}
\gdef\bottomparagraph{}
\gdef\newparagraph{}
\gdef\chaptername{References}
\gdef\sectionname{References}
\gdef\headlinesectionname{References}

\centerline{\bf References for Volume 5}

%use subparagraph details (5A1Aa) only for X. Y. Z.
%  \S511 {\it notes\/}, \S512 {\it intro.}]

\medskip%A

%Antonovskii M.Ja.\ \& Chudnovsky D.V.\ [76] `Some questions of general
%topology and Tikhonov semifields II', Russian Math.\ Surveys 31 (1976)
%69-128.

Argyros S.A.\ \& Kalamidas N.D.\ [82] `The $K_{\alpha n}$ property on
spaces with strictly positive measure', Canadian J.\ Math.\ 34 (1982)
1047-1058.
\cmmnt{[525T.]}

Argyros S.A.\ \& Tsarpalias A.\ [82] `Calibers of compact spaces',
Trans.\ Amer.\ Math.\ Soc.\ 270 (1982) 149-162.
\cmmnt{[525N.]}

\medskip%B

Babiker A.G.\ [76] `On uniformly regular topological measure spaces',
Duke Math.\ J.\ 43 (1976) 775-789.
\cmmnt{[533Xd.]}

Balcar B.\ \& Fran\v{e}k F.\ [82] `Independent families in complete Boolean
algebras', Trans.\ Amer.\ Math.\ Soc.\ 274 (1982) 607-618.
\cmmnt{[515H.]}

Balcar B., Jech T.\ \& Paz\'ak T.\ [05]
`Complete ccc Boolean algebras, the order sequential topology,
and a problem of von Neumann',
Bull.\ London Math.\ Soc.\ 37 (2005) 885-898.
\cmmnt{[\S539 {\it intro.}, 539E, 539L, 539N, \S539 {\it notes\/}, 555K.]}
 % arXiv.math.LO/0312473

Balcar B., Jech T.\ \& Zapletal J.\ [97] `Semi-Cohen Boolean algebras',
Ann.\ Pure \& Applied Logic 87 (1997) 187-208.
\cmmnt{[\S515 {\it notes\/}, \S547 {\it notes\/}.]}

Balcar B.\ \& Vojt\'a\v{s} P.\ [77] `Refining systems on Boolean algebras',
pp.\ 45-58 in {\smc Lachlan Srebny \& Zarach 77}.
\cmmnt{[515E.]}

Bar-Hillel Y.\ [70] {\it Mathematical Logic and Foundations of Set Theory.}
North-Holland, 1970.

Bartoszy\'nski T.\ [84] `Additivity of measure implies additivity of
category', Trans.\ Amer.\ Math.\ Soc.\ 281 (1984) 209-213.
\cmmnt{[Chap.\ 52 {\it intro.}, 522Q, \S524 {\it notes\/}.]}

Bartoszy\'nski T.\ [87] `Combinatorial aspects of measure and category',
Fund.\ Math.\ 127 (1987) 225-239.
\cmmnt{[522S.]}
% http://matwbn.icm.edu.pl/ksiazki/fm/fm127/fm127118.pdf

% Bartoszy\'nski T.\ [88] `On covering of real line by null sets',
% Pacific J.\ Math.\ 131 (1988) 1-12.

% Bartoszy\'nski T.\ [p97] `Invariants of measure and category', in
% preparation for {\smc Foreman Kanamori \& Magidor ??}

% Bartoszy\'nski T.\ [00] `A note on duality between measure and category',
% Proc.\ Amer.\ Math.\ Soc.\ 128 (2000) 2745-2748.
   % (2^N,+,meager) not iso to (2^N,+,negligible)

Bartoszy\'nski T.\ \& Judah H.\ [89] `On the cofinality of the smallest
covering of the real line by meager sets', J.\ Symbolic Logic 54 (1989) 828-832.
\cmmnt{[522V.]}

Bartoszy\'nski T.\ \& Judah H.\ [95] {\it Set Theory:  on the structure
of the real line}.  A.K.Peters, 1995.
\cmmnt{[522R, \S522 {\it notes\/},
526M, \S528 {\it notes\/}, 534Q, \S534 {\it notes\/}, 546I,
553C, 554Ya, 5A1Q.]}

%Bartoszy\'nski T., Judah H.\ \& Shelah S.\ [93] `The Cicho\'n diagram',
%J.\ Symbolic Logic 58 (1993) 401-423.

%Bartoszy\'nski T.\ \& Scheepers M.\ [96] (eds.) {\it Set Theory (BEST, 1992-1994).}
%Amer.\ Math.\ Soc., 1996 (Contemporary Mathematics 192).
% ISBN 0821803069

Bartoszy\'nski T.\ \& Shelah S.\ [92] `Closed measure zero sets', Ann.\
Pure \& Applied Logic 58 (1992) 93-110.
\cmmnt{[526M.]}

Baumgartner J.E., Taylor A.D.\ \& Wagon S.\ [77] `On splitting
stationary subsets of large cardinals',
J.\ Symbolic Logic 42 (1977) 203-214.
\cmmnt{[\S541 {\it notes\/}, \S544 {\it notes\/}.]}

Bell M.G.\ [81] `On the combinatorial principle $P(c)$',
Fundamenta Math.\ 114 (1981) 149-157.
\cmmnt{[\S517 {\it notes\/}.]}

Benedikt M.\ [98] `Ultrafilters which extend measures',
J.\ Symbolic Logic 63 (1998) 638-662.
\cmmnt{[538M, 538Xm, \S538 {\it notes\/}.]}

Benedikt M.\ [99] `Hierarchies of measure-theoretic ultrafilters',
Annals of Pure and Applied Logic 97 (1999) 203-219.
\cmmnt{[538H.]}

Blair C.E.\ [77] `The Baire category theorem implies the principle of
dependent choice', Bull.\ Acad.\ Polon.\ Sci.\ (Math.\ Astron.\ Phys.) 25
(1977) 933-934.
\cmmnt{[\S566 {\it notes\/}.]}

%Blass A.\ [88] `Selective ultrafilters and homogeneity', Ann.\ Pure and
%Applied Logic 38 (1988) 252-285.

%Blass A.\ [p01b] `Combinatorial cardinal characteristics of the continuum',
%pp. ?? in {\smc Foreman Kanamori \& Magidor ??}.

Blass A.\ \& Laflamme C.\ [89] `Consistency results about filters and the
number of inequivalent growth types', J.\ Symbolic Logic 54 (1989) 50-56.
% u<g  filter dichotomy
\cmmnt{[5A6Ib, 5A6J.]}

Booth D.\ [70] `Ultrafilters on a countable set',
Ann.\ Math.\ Logic 2 (1970) 1-24.
\cmmnt{[517R.]}

Borel E.\ [1919] `Sur la classification des ensembles de mesure nulle', Bull.\ Soc.\ Math.\ France 47 (1919) 97-125.
\cmmnt{[\S534 {\it intro.}]}

Brendle J.\ [00] {\it Inner model theory and large cardinals.}
Research report, March 2000.
\cmmnt{[528N, 529F-529H.]} %529F 529G 529H

%Brendle J.\ [01] `Cardinal invariants of the continuum --- a survey',
%J.\ Research Institute for Mathematical Sciences, Kyoto University
%(JRIMS Kokyuroku) 1202 (2001) 7-32.

Brendle J.\ [06] `Cardinal invariants of the continuum and combinatorics
on uncountable cardinals',
Annals of Pure and Applied Logic 144 (2006) 43-72.
\cmmnt{[529F, 529H.]}

Burke M.R.\ [93] `Liftings for Lebesgue measure', pp.\ 119-150 in
{\smc Judah 93}.
\cmmnt{[\S535 {\it notes\/}.]}

%Burke M.R.\ [n95] `Consistent liftings', privately circulated, 1995.

%Burke M.R.\ [n96] `A theorem of Gitik and Shelah on disjoint refinements
%of sequences of subsets of the real line', note of 23.8.96.
%\cmmnt{[5x46J, \S547 {\it intro.}]}

Burke M.R.\ [n05] `Non-null sets of minimal cofinality', note of 2.3.05.
\cmmnt{[523K.]}

Burke M.R.\ \& Magidor M.\ [90] `Shelah's pcf theory and
its applications', Ann.\ Pure and Applied Logic 50 (1990) 207-254.
\cmmnt{[513J, \S5A2 {\it intro.}.]}

Burke M.R.\ \& Shelah S.\ [92]
`Linear liftings for non-complete probability spaces',
Israel J.\ Math.\ 79 (1992) 289-296.
\cmmnt{[\S535 {\it notes\/}.]}

\medskip%C

Carlson T.J.\ [84] `Extending Lebesgue measure to infinitely many sets',
Pacific J.\ Math.\ 115 (1984) 33-45.
\cmmnt{[552N, \S552 {\it notes.}]}

Carlson T.J.\ [93] `Strong measure zero and strongly meager sets',
Proc.\ Amer.\ Math.\ Soc.\ 118 (1993) 577-586.
\cmmnt{[534R.]}

Carlson T., Frankiewicz R.\ \& Zbierski P.\ [94]
`Borel liftings of the measure
algebra and the failure of the continuum hypothesis', Proc.\ Amer.\ Math.\ Soc.\
120 (1994) 1247-1250.
\cmmnt{[\S535 {\it notes\/}, 554I.]}

Choksi J.R.\ \& Fremlin D.H.\ [79] `Completion regular measures on
product spaces', Math.\ Ann.\ 241 (1979) 113-128.
\cmmnt{[532I.]}

%Cicho\'n J., Kamburelis T.\ \& Pawlikowski J.\ [85] `On dense subsets of
%the measure algebra', Proc.\ Amer.\ Math.\ Soc.\ 94 (1985) 142-146.

Cicho\'n J.\ \& Pawlikowski J.\ [86]
`On ideals of subsets of the plane and on Cohen
reals', J.\ Symbolic Logic 51 (1986) 560-569.
\cmmnt{[527F.]}

%Cohen P.\ [63] `The independence of the continuum hypothesis',
% Proc.\ Nat.\ Acad.\ Sci.\ U.S.A.\ 50 (1963) 1143-1148.

Comfort W.W.\ \& Negrepontis S.\ [74] {\it The Theory of Ultrafilters.}
Springer, 1974.
\cmmnt{[538F, 538Yb.]}

Comfort W.W.\ \& Negrepontis S.\ [82] {\it Chain Conditions in
Topology.}   Cambridge U.P., 1982.
\cmmnt{[516 {\it notes\/}.]}

Cummings J.\ [92] `A model in which GCH holds at successors but fails at
limits',
Trans.\ Amer.\ Math.\ Soc.\ 329 (1992) 115-142.
\cmmnt{[525Z.]}

\medskip%D

Dellacherie C., Meyer P.A.\ \& Weil M.\ (eds.) [73]
{\it S\'eminaire de Probabilit\'es VII.}   Springer, 1973
(Springer Lecture Notes in Mathematics 321).

Devlin K.J.\ [84] {\it Constructibility.} Springer, 1984.
\cmmnt{[5A6D.]}

%Dow A.\ \& Hart K.P.\ [p98] `The measure algebra does not always embed',
%preprint, 1998

Dow A.\ \& Stepr\=ans J.\ [94] `The $\sigma$-linkedness of the measure
algebra', Canad.\ Math.\ Bulletin 37 (1994) 42-45.
\cmmnt{[524L.]}

%Drake F.\ [74] {\it Set Theory.}  North-Holland, 1974.

%Drake F.R.\ \& Singh D.\ [96] {\it Intermediate Set Theory.}
%Wiley, 1996.  \cmmnt{[\S5A1.]}
% ISBN 0 471 96496 4

D\v{z}amonja M.\ \& Kunen K.\ [93] `Measures on compact HS spaces',
Fundamenta Math.\ 143 (1993) 41-54.
\cmmnt{[\S531 {\it notes\/}.]}

%D\v{z}amonja M.\ \& Kunen K.\ [95] `Properties of the class of
%measure-separable compact spaces', Fundamenta Math.\ 147 (1995) 261-277.

D\v{z}amonja M.\ \& Plebanek G.\ [04]
`Precalibre pairs of measure algebras',
Topology Appl.\ 144 (2004) 67-94.
\cmmnt{[525L.]}

\medskip%E

Engelking R.\ [89] {\it General Topology.} Heldermann, 1989 (Sigma
Series in Pure Mathematics 6).
\cmmnt{[\S5A4.]}

%Engelking R.\ \& Kar{\l}owicz M.\ [65]
%`Some theorems of set theory and their topological consequences',
% Fund.\ Math.\ 57 (1965) 275-285.

Erd\H{o}s P., Hajnal A., M\'at\'e A.\ \& Rado R.\ [84] {\it Combinatorial
Set Theory:  Partition Relations for Cardinals.}   Akad\'emiai Kiad\'o,
1984 (Disquisitiones Math.\ Hung.\ 13).
\cmmnt{[5A1F, 5A1H, 5A6F.]}

\medskip%F

Farah I.\ [00] {\it Analytic Quotients.}
Mem.\ Amer.\ Math.\ Soc.\ 148 (2000).
\cmmnt{[5A6H.]}

Farah I.\ [03] `How many algebras $\Cal P(\Bbb N)/\Cal I$ are there?',
Illinois J.\ Math.\ 46 (2003) 999-1033.
\cmmnt{[\S556 {\it notes\/}.]}

Farah I.\ [06] `Analytic Hausdorff gaps II: the density zero ideal',
Israel J.\ Math.\ 154 (2006) 235-246.
\cmmnt{[556S.]}

Farah I.\ \& Veli\v{c}kovi\'c B.\ [06]
`Von Neumann's problem and large cardinals',
Bull.\ London Math.\ Soc.\ 38 (2006) 907-912.
\cmmnt{[539Q.]}

Feferman S. \& L\'evy A. [63] `Independence results in set theory by Cohen's
method', Notices Amer.\ Math.\ Soc.\ 10 (1963) 53.
\cmmnt{[561A.]}

Fleissner W.G.\ [91]
`Normal measure axiom and Balogh's theorems', Topology and its Appl.\
39 (1991) 123-143.
\cmmnt{[555N.]}

% Fleissner W.G., Hansell R.W.\ \& Junnila H.J.K.\ [82] `PMEA implies
% Proposition P', Topology and its Appl.\ 13 (1982) 255-262.

% Foreman M., Kanamori A.\ \& Magidor M.\ [??] {\it Handbook of Set
% Theory}.   Kluwer, to appear.

%Foreman M., Magidor M.\ \& Shelah S.\ [88] `Martin's maximum, saturated ideals, and non-regular ultrafilters', Annals of Math.\ (2) 127 (1988) 1-47.

Foreman M.\ [10] `Ideals and generic elementary embeddings',  pp.\ 885-1147
in {\smc Foreman \& Kanamori 10}, vol.\ 2.
\cmmnt{[547Z.]}

Foreman M.\ \& Kanamori A.\  [10] {\it Handbook of Set Theory.}
Springer, 2010.

Foreman M.\ \& Wehrung F.\ [91] `The Hahn-Banach theorem implies the
existence of a non-Lebesgue measurable set', Fundamenta Math.\ 138
(1991) 13-19.
\cmmnt{[563A.]}

Foreman M.\ \& Woodin W.\ [91] `The generalized continuum hypothesis can
fail everywhere', Annals of Math.\ 133 (1991) 1-35.
\cmmnt{[525Z.]}

Fossy J.\ \& Morillon M.\ [98] `The Baire category property and some
notions of compactness', J.\ London Math.\ Soc. (2) 57 (1998) 1-19.
\cmmnt{[\S566 {\it notes\/}.]}

% Frankiewicz R.\ \& Plebanek G.\ [94] `Nonaccessible filters in
% measure algebras and functionals on $L^{\infty}(\lambda)^*$', Studia
% Math.\ 108 (1994) 191-200.

%Frankiewicz R.\ \& Plebanek G.\ [96] `On asymptotic density and
%uniformly distributed sequences', Studia Math.\ 119 (1996) 17-26.

Freese R. \& Nation J.B.\ [78] `Projective lattices',
Pacific J.\ Math.\ 75 (1978) 93-106.
\cmmnt{[\S518 {\it notes\/}.]}

Freiling C.\ [86] `Axioms of symmetry;  throwing darts at the real number
line', J.\ Symbolic Logic 51 (1986) 190-200.
\cmmnt{[537K.]}

Fremlin D.H.\ [74] {\it Topological Riesz Spaces and Measure Theory.}
Cambridge U.P., 1974.
\cmmnt{[564Xc.]}

Fremlin D.H.\ [77]
`Uncountable powers of {\bf R} can be almost Lindel\"of',
Manuscripta Math.\ 22 (1977) 77-85.
\cmmnt{[533J.]}

Fremlin D.H.\ [84a] {\it Consequences of Martin's Axiom.}
Cambridge U.P., 1984.
\cmmnt{[\S511 {\it notes\/}, \S517 {\it notes\/}, \S531 {\it notes\/}.]}

Fremlin D.H.\ [84b] `On the additivity and cofinality of Radon measures',
Mathematika 31 (1984) 323-335.
\cmmnt{[\S524 {\it notes\/}.]}

Fremlin D.H.\ [87] {\it Measure-additive Coverings and
Measurable Selectors.} Dissertationes Math.\ 260 (1987).
\cmmnt{[551A.]}

Fremlin D.H.\ [88] `Large correlated families of positive
random variables',
Math.\ Proc.\ Cambridge Phil.\ Soc.\ 103 (1988) 147-162.
\cmmnt{[525S.]}

Fremlin D.H.\ [91] `The partially ordered sets of measure
theory and Tukey's ordering,' Note di Matematica  11 (1991) 177-214.
[\S524 {\it notes\/}, 526B, 526I, 527J, 529C, 529D, 534L.]

Fremlin D.H.\ [93] `Real-valued-measurable cardinals', pp.\ 151-304
in {\smc Judah 93}.
\cmmnt{[\S537 {\it notes\/}, \S541 {\it notes\/}, \S544 {\it notes\/},
\S545 {\it notes\/}, \S5A2 {\it intro.}.]}

Fremlin D.H.\ [97] `On compact spaces carrying Radon measures of
uncountable Maharam type', Fundamenta Math.\  154 (1997) 295-304.
\cmmnt{[531T.]}

%Fremlin D.H.\ [00] `Weakly $\alpha$-favourable measure spaces',
%Fundamenta Math.\ 165 (2000) 67-94.

Fremlin D.H.\ [03] `Skew products of ideals',
J.\ Applied Analysis 9 (2003) 1-18.
\cmmnt{[\S527 {\it notes\/}.]}
%n01910.tex

%Fremlin D.H.\ [n97] `A theorem of Gitik and Shelah on disjoint
%refinements of sequences of subsets of the real line', note of 4.4.97.
%n96j01.tex

%Fremlin D.H.\ [n99] `Embedding the Lebesgue measure algebra into
%$\Cal P\Bbb N/\Cal Z$', note of 25.5.99.
%n99525.tex

%Fremlin D.H.\ [n99] `Notes on {\smc Farah p99}', note of 26.7.99.
%n99612.tex

%Fremlin D.H.\ [n99] `Notes on {\smc Solecki p96}', note of 17.6.99.
%n99617.tex

%Fremlin D.H.\ [n02] `Games and the Tukey classification', note of 29.9.02.
%n02929.tex

Fremlin D.H.\ \& Grekas S.\ [95] `Products of completion regular
measures', Fundamenta Math.\ 147 (1995) 27-37.
\cmmnt{[532D.]}

Fremlin D.H.\ \& Miller A.W.\ [88]
`On some properties of Hurewicz, Menger and Rothberger',
Fund.\ Math.\ 129 (1988) 17-33.
\cmmnt{[534Q.]}

Fremlin D.H., Natkaniec T.\ \& Rec{\l}aw I.\ [00]
`Universally Kuratowski-Ulam spaces', Fundamenta Math.\ 165 (2000) 239-247.
\cmmnt{[\S527 {\it notes\/}.]}

%Friedman H.\ [80]
%`A consistent Fubini-Tonelli theorem for nonmeasurable functions',
% Illinois J.\ Math.\ 24 (1980) 390-395.

Friedman S.D. \& Koepke P.\ [97] `An elementary approach to the fine
structure of L', Bull.\ Symbolic Logic 3 (1997) 453-468.
\cmmnt{[5A6D.]}

Fuchino S.,  Geschke S., Shelah S.\ \& Soukup L.\ [01]
`On the weak Freese-Nation property of complete Boolean algebras',
Ann.\ Pure Appl.\ Logic 110 (2001) 89-105.
\cmmnt{[518K.]}

Fuchino S., Geschke S.\ \& Soukup L.\ [01] `On the weak Freese-Nation
property of $\Cal P(\omega)$', Arch.\ Math.\ Logic 40 (2001) 425-435.
\cmmnt{[522U, \S522 {\it notes\/}.]}

Fuchino S., Koppelberg S.\ \& Shelah S.\ [96] `Partial orderings with
the weak Freese-Nation property',
Ann.\ Pure and Applied Logic 80 (1996) 35-54.
\cmmnt{[518A, 518G, 518Yb, \S518 {\it notes\/}, 522U.]}

Fuchino S.\ \& Soukup L.\ [97] `More set theory around the weak
Freese-Nation property', Fundamenta Math.\ 154 (1997) 159-176.
\cmmnt{[518I, 518K.]}

\medskip%G

Galvin F.\ [80] `Chain conditions and products',
Fundamenta Math.\ 108 (1980) 33-48.
\cmmnt{[\S553 {\it notes\/}.]}

Galvin F., Mycielski J.\ \& Solovay R.M.\ [79] `Strong measure zero sets',
Notices Amer.\ Math.\ Soc.\ 26 (1979) A280.
%(abstract 79T-E25);  often assigned to year 1973
\cmmnt{[534K.]}

Gandy R.O.\ \& Hyland J.M.E.\ [77](eds.) {\it Logic Colloquium '76.}
North-Holland, 1977.

Geschke S.\ [02] `On tightly $\kappa$-filtered Boolean algebras',
Algebra Universalis 47 (2002) 69-93.
\cmmnt{[518P, 518S, \S535 {\it notes\/}.]}

Gitik M.\ \& Shelah S.\ [89] `Forcings with ideals and
simple forcing notions', Israel J.\ Math.\ 68 (1989) 129-160.
\cmmnt{[543E, 547F.]}

Gitik M.\ \& Shelah S.\ [93]
`More on simple forcing notions and forcings with ideals',
Annals of Pure and Applied Logic 59 (1993) 219-238.
\cmmnt{[542E, 543E, 547F.]}

Gitik M.\ \& Shelah S.\ [01]
`More on real-valued measurable cardinals and forcing with ideals',
Israel J.\ Math. 124 (2001) 221-242.
\cmmnt{[548E.]}
%GiSh582

G\l\'owczy\'nski W.\ [91] `Measures on Boolean algebras',
Proc.\ Amer.\ Math.\ Soc.\ 111 (1991) 845-849.
\cmmnt{[555K.]}

G\l\'owczy\'nski W.\ [08] `Outer measure on Boolean algebras', Acta Univ.\
Carolinae (Math.\ et Phys.) 49 (2008) 3-8.
\cmmnt{[555K.]}

Goldstern M., Judah H.\ \& Shelah S.\ [93] `Strong measure zero sets
without Cohen reals',
J.\ Symbolic Logic 58 (1993) 1323-1341.
\cmmnt{[\S534 {\it notes\/}.]}

\medskip%H

% Hart J.\ \& Kunen K.\ [98] `Weak Measure Extension Axioms',
% Topology and its Applications 85 (1998) 219-246.
% http://www.math.wisc.edu/~kunen/me.ps

Haydon R.G.\ [77] `On Banach spaces which contain $\ell^1(\tau)$ and
types of measures on compact spaces', Israel J.\ Math.\ 28 (1977) 313-324.
\cmmnt{[531E, 531L.]}

%Haydon R.G.\ [78] `On dual $L^1$-spaces and injective bidual Banach
%spaces,' Israel J.\ Math.\ 31 (1978) 142-152.

Hodges W., Hyland M., Steinhorn C.\ \& Truss J.\ [96] (eds.)
{\it Logic: from Foundations to Applications,
European Logic Colloquium, 1993.}
Clarendon, 1996.
%[ISBN 0-19-853862-6]

Humke P.D.\ \& Laczkovich M.\ [05] `Symmetrically approximately
continuous functions, consistent density theorems, and Fubini type
inequalities', Trans.\ Amer.\ Math.\ Soc.\ 357 (2005) 31-44.
\cmmnt{[537Q, \S537 {\it notes\/}.]}

\medskip%I

Ihoda J.I.\ [88] `Strong measure zero sets and rapid filters',
J.\ Symbolic Logic 53 (1988) 393-402.
%model of Borel conjecture + Ramsey filter
\cmmnt{[534Q.]}

\medskip%J

Jech T.\ [73] {\it The Axiom of Choice.}  North-Holland, 1973.
\cmmnt{[561A, 561Xc, 561Yc, 561Yi, 564Yd.]}

Jech T.\ [78] {\it Set Theory.}  Academic, 1978 (ISBN 0123819504).
\cmmnt{[{\it Intro.}, 553H,
555O, \S555 {\it notes\/}, \S562 {\it notes\/}, \S5A1, 5A6B.]}
%5A1B 5A1D 5A1E 5A1F B
%ASIN B0095LMZ9O

Jech T.\ [03] {\it Set Theory, Millennium Edition.}  Springer, 2003
(ISBN 3540440852).
\cmmnt{[\S521 {\it notes\/}, \S555 {\it notes\/},
\S562 {\it notes\/}, \S567 {\it notes\/}, \S5A1, 5A3N, 5A3P, 5A6B.]}
%5A1B 5A1D 5A1E 5A1F B

%Jin R.\ \& Shelah S.\ [p96] `Compactness of Loeb spaces', preprint, 1996.

%Jovanovi\'c A.\ [p01]
%`Real valued measurability, some set-theoretic aspects',
%pp.\ ?? in {\smc Pap 0?}

Judah H.\ [93] (ed.) {\it Proceedings of the Bar-Ilan Conference on Set
Theory and the Reals, 1991.} Amer.\ Math.\ Soc.\ (Israel Mathematical
Conference Proceedings 6), 1993.

Judah H.\ \& Repick\'y M.\ [95] `Amoeba reals', J.\ Symbolic Logic 60 (1995) 1168-1185.
\cmmnt{[528N.]}

%Juh\'asz I.\ [71]  {\it Cardinal Functions in Topology.}  North-Holland,
%1975 (Math.\ Centre Tracts 34).

%Juh\'asz I.\ [80]  {\it Cardinal Functions in Topology -- Ten Years Later.}  Mathematisch Centrum, Amsterdam, 1980 (Math.\ Centre Tracts 123).   ISBN 9061961963.

Just W.\ [92] `A modification of Shelah's oracle-c.c., with
applications', Trans.\ Amer.\ Math.\ Soc.\ 329 (1992) 325-356.
\cmmnt{[527M.]}

%Just W.\ \& Krawczyk A.\ [84] `On certain Boolean algebras
%$\Cal P(\omega)/I$', Trans.\ Amer.\ Math.\ Soc.\ 285 (1984) 411-429.

Just W.\ \& Weese M.\ [96] {\it Discovering Modern Set Theory I.}
Amer.\ Math.\ Soc., 1996 (Graduate Studies in Mathematics 8).
\cmmnt{[{\it Intro.}.]}

Just W.\ \& Weese M.\ [97] {\it Discovering Modern Set Theory II.}
Amer.\ Math.\ Soc., 1997 (Graduate Studies in Mathematics 18).
\cmmnt{[5A1B, 5A1F, 5A1H.]}

\medskip%K

%Kada M. \& Yuasa Y.\ [96] `Cardinal invariants about shrinkability of unbounded
%sets', Topology and its Applications 74 (1996) 215-223.

%Kamburelis A.\ [89] `Iterations of Boolean algebras with measure',
%Arch.\ Math.\ Logic 29 (1989) 21-28.
% chargeable algebras

%Kamburelis A.\ [90] `A new proof of the Gitik-Shelah theorem', Israel
%J.\ Math.\ 72 (1990) 373-380.

Kanamori A.\ [03] {\it The Higher Infinite.}   Springer, 2003.
\cmmnt{[\S541 {\it notes\/}, \S555 {\it notes\/}, \S567 {\it notes\/},
5A1H, 5A6B, 5A6F.]}

Kechris A.S.\ [95] {\it Classical Descriptive Set Theory.}  Springer,
1995.
\cmmnt{[5A1D, \S567 {\it notes\/}.]}
%QA248.K4

Keisler H.J.\ \& Tarski A.\ [64] `From accessible to
inaccessible cardinals', Fundamenta Math.\ 53 (1964) 225-308;  errata
Fundamenta Math.\ 57 (1965) 119.
\cmmnt{[\S541 {\it intro.\/}, \S541 {\it notes\/}.]}

Kelley J.L.\ [50] `The Tychonoff product theorem implies the axiom of
choice', Fundamenta Math.\ 37 (1950) 75-76.
\cmmnt{[561D.]}

Koppelberg S.\ [75]
`Homomorphic images of $\sigma$-complete Boolean algebras',
Proc.\ Amer.\ Math.\ Soc.\ 51 (1975) 171-175.
\cmmnt{[515L.]}

Koppelberg S.\ [89] {\it General Theory of Boolean
Algebras\/}, vol.\ 1 of {\smc Monk 89}.
\cmmnt{[\S515 {\it notes\/}.]}

Koppelberg S.\ \& Shelah S.\ [96]
`Subalgebras of Cohen algebras need not be Cohen',
pp.\ 261-275 in {\smc Hodges Hyland Steinhorn \& Truss 96}.
\cmmnt{[\S515 {\it notes\/}, \S547 {\it notes\/}.]}

Kraszewski J.\ [01] `Properties of ideals on the generalized Cantor spaces',
J.\ Symbolic Logic 66 (2001) 1303-1320.
\cmmnt{[523G, 523H, 523J, \S523 {\it notes\/}.]}

%Krawczyk A.\ [n85] `Dominating reals add random reals', 1985.

Kumar A.\ [13] `Avoiding rational distances', Real Analysis Exchange 38
(2012/13), 493-498.
%{\tt http://www.math.wisc.edu/\~{}akumar/ARD.pdf}.
%{\tt
%https:/{\bsp}projecteuclid.org/journals/real-analysis-exchange/volume-38%
%/issue-2/Avoiding-Rational-{\vthsp}Distances/rae/1403894908.full}
\cmmnt{[548Xb.]}

%Kumar A.\ [n16] `On a theorem of Gitik and Shelah', note of January 2016
%({\tt http:/{\bsp}www.math.huji.ac.il/\~{}akumar/otgs.pdf}).

Kumar A.\ \& Shelah S.\ [17] `A transversal of full outer measure',
Advances in Math.\ 321 (2017) 475-485.
\cmmnt{[547P, 547Q, \S548 {\it intro.\/}, 548C.]}
%http://home.iitk.ac.in/~krashu/tfom.pdf;  printed version not seen

%Kunen K.\ [75] `Some points in $\beta N$',
% Math.\ Proc.\ Cambridge Phil.\ Soc.\ 80 (1975) 385-398.

Kunen K.\ [80] {\it Set Theory.}  North-Holland, 1980.
\cmmnt{[{\it Intro.}, 511A, \S521 {\it notes\/}, \S522 {\it notes\/},
\S541 {\it notes\/}, 551Q, \S551 {\it notes\/}, 556F,
\S562 {\it notes}, Appendix {\it intro.}, \S5A1, \S5A3, 5A6B.]}

Kunen K.\ [81] `A compact $L$-space under CH', Topology and its Appl.\
12 (1981) 283-287.
\cmmnt{[\S531 {\it notes\/}.]}

Kunen K.\ [84] `Random and Cohen reals', pp.\ 887-911 in
{\smc Kunen \& Vaughan 84}.
\cmmnt{[\S554 {\it notes}.]}

Kunen K.\ [n70] `$\Pi^1_1$ reflection at the continuum', note of January
1970.
\cmmnt{[543C, 544C, 544E, 544F.]}

Kunen K.\ \& Mill J.van [95] `Measures on Corson compact spaces',
Fundamenta Math.\ 147 (1995) 61-72.
\cmmnt{[531O.]}
%http://www.math.wisc.edu/~kunen/corson.ps

Kunen K.\ \& Vaughan J.E.\ [84] (eds.) {\it Handbook of Set-Theoretic
Topology.}   North-Holland, 1984.

Kuratowski K.\ [66] {\it Topology,} vol.\ I.   Academic, 1966.
\cmmnt{[562G.]}

\medskip%L

Lachlan A., Srebny M.\ \& Zarach A.\ [77] (eds.) {\it Set Theory and Hierarchy Theory.}  Springer, 1977 (Lecture Notes in Math.\ 619).

%Laczkovich M.\ \& Miller A.W.\ [95] `Measurability of functions with
%approximately continuous vertical sections and measurable horizontal
%sections', Colloquium Math.\ 69 (1995) 299-308.

Larson P.\ [09] `The filter dichotomy and medial limits',
J.\ Math.\ Logic 9 (2009) 159-165.
\cmmnt{[538S.]}

Larson P., Neeman I.\ \& Shelah S.\ [10] `Universally measurable sets in
generic extensions', Fundamenta Math.\ 208 (2010) 173-192.
\cmmnt{[553O.]}
	
%Larson P.\ \& Shelah S.\ [p09]
%`There may be just continuum many universally measurable sets', 2009.
%LrSh:947

Laver R.\ [76] `On the consistency of Borel's conjecture', Acta Math.\
137 (1976) 151-169.
\cmmnt{[534Q.]}

Laver R.\ [87] `Random reals and Souslin trees',
Proc.\ Amer.\ Math.\ Soc.\ 100 (1987) 531-534.
\cmmnt{[553M.]}

Levinski J.-P., Magidor M.\ \& Shelah S.\ [90] `Chang's conjecture for
$\aleph_{\omega}$', Israel J.\ Math.\ 69 (1990) 161-172.
\cmmnt{[5A6F.]}

Levy A.\ [71] `The sizes of the indescribable cardinals',
pp.\ 205-218 in {\smc Scott 71}.
\cmmnt{[\S541 {\it notes\/}, \S544 {\it notes\/}.]}

Lipecki Z.\ [09]
`Semivariations of an additive function on a Boolean ring',
Colloquium Math.\ 117 (2009) 267-279.
\cmmnt{[552Xe.]}

%Loeb P.A.\ \& Talvila E.\ [p00]
%`Covering theorems and Lebesgue integration', preprint, June 2000.

%Losada M. \& Todor\v{c}evi\'c S. [p98] `Chains and antichains in Boolean
%algebras', note of 5.12.98.

%Louveau A.\ \& Veli\v{c}kovi\'c B.\ [99] `Analytic ideals and cofinal
%types', Ann.\ Pure and Applied Logic 99 (1999) 171-195.
%\ell^1\not\prT\CalNwd

\medskip%M

Maharam D.\ [1947] `An algebraic characterization of measure algebras',
Ann.\ Math.\ 48 (1947) 154-167.
\cmmnt{[539P.]}

Makowsky J.A.\ \& Ravve E.V.\ (eds.) [98]
{\it Logic Colloquium '95.}  Springer, 1998 (Lecture Notes in Logic 11).
%ISBN 3-540-63994-2

%Malliaris M.\ \& Shelah S.\ [16] `Cofinality spectrum theorems in
%model theory, set theory and general topology',
%J.\ Amer.\ Math.\ Soc.\ 29 (2016) 237-297.

Martin D.A.\ [70] `Measurable cardinals and analytic games',
Fundamenta Math.\ 66 (1970) 287-291.
\cmmnt{[567N.]}

Martin D.A.\ [75] `Borel determinacy', Ann.\ of Math.\ (2) 102 (1975)
363-371.
\cmmnt{[\S567 {\it notes\/}.]}

Martin D.A.\ \& Solovay R.M.\ [70] `Internal Cohen
extensions', Ann.\ Math.\ Logic 2 (1970) 143-178.
\cmmnt{[Chap.\ 52 {\it intro.}, \S528 {\it intro.}, \S528 {\it notes\/}.]}

%Mathias A.R.D.\ [78] `0\# and the $p$-point problem', pp.\ 375-384 in
%{\smc M\"uller \& Scott 78}.

M\'atrai T.\ [p09] `More cofinal types of definable directed orders', 2009
({\tt https://citeseerx.ist.psu.edu/viewdoc/summary?doi=10.1.1.414.3938}).
\cmmnt{[526L, \S526 {\it notes\/}.]}
% books/matrai.p09.pdf

%Mercourakis S.\ [96] `Some remarks on countably determined measures and
%uniform distribution of sequences', Monatschefte f\"ur Math.\ 121 (1996)
%79-101.

Meyer P.-A.\ [73] `Limites m\'ediales d'apr\`es Mokobodzki',
pp.\ 198-204 in {\smc Dellacherie Meyer \& Weil 73}.
\cmmnt{[538Yl.]}

Meyer P.A.\ [75] (ed.) {\it S\'eminaire de Probabilit\'es IX.}
Springer, 1975 (Lecture Notes in Mathematics 465).

Miller A.W.\ [80] `There are no Q-points in Laver's model for the Borel
conjecture', Proc.\ Amer.\ Math.\ Soc.\ 78 (1980) 103-106.
\cmmnt{[\S538 {\it notes\/}.]}
%actually, no rapid ultrafilters, therefore no rapid filters

Miller A.W.\ [81] `Some properties of measure and category',
Trans.\ Amer.\ Math.\ Soc.\ 266 (1981) 93-114.
\cmmnt{[522J.]}

Miller A.W.\ [82] `The Baire category theorem and cardinals of countable
cofinality', J.\ Symbolic Logic 47 (1982) 275-288.
\cmmnt{[552G, 552Xb.]}

%Miller A.W.\ [83] `Mapping a set of reals onto the reals',
%J.\ Symbolic Logic 48 (1983) 575-584.

Mokobodzki G.\ [75] `R\'el\`evement bor\'elien compatible avec une
classe d'ensembles n\'egligeables.  Application \grv{a} la
d\'esint\'egration des mesures', pp.\ 437-442 in {\smc Meyer 75}.
\cmmnt{[535I.]}

Mokobodzki G.\ [7?] `D\'esintegration des mesures et rel\`evements
Boreliens de sous-espaces de $L^{\infty}(X,\Cal B,\mu)$',
Lab.\ d'Analyse Fonctionnelle, Univ.\ Paris VI (?).
\cmmnt{[535E, \S535 {\it notes\/}.]\footnote{I have not been able to
locate this paper;  I believe it was a seminar report.   I took notes from
it in 1977.}}
% "On two theorems of Mokobodzki", 23.6.77.

Monk J.D.\ [89] (ed.) {\it Handbook of Boolean Algebra.}
North-Holland, 1989.

%Monk J.D.\ [96] {\it Cardinal Invariants on Boolean Algebras.}
%Birkhauser, 1996.
% QA266

Moore J.T.\ [05] `Set mapping reflection', J.\ Math.\ Logic 5 (2005) 87-98.
\cmmnt{[517O.]}

Moschovakis Y.N.\ [70] `Determinacy and prewellorderings of the continuum',
pp.\ 24-62 in {\smc Bar-Hillel 70}.
\cmmnt{[567M.]}

%M\"uller G.H.\ \& Scott D.S.\ [78] (eds.) {\it Higher Set Theory.}
%Springer, 1978 (Lecture Notes in Mathematics 669).

Mycielski J.\ [64] `On the axiom of determinateness',
Fund.\ Math.\ 53 (1964) 205-224.
\cmmnt{[567D.]}

%Mycielski J.\ \& Steinhaus H.\ [62]
%`A mathematical axiom contradicting the axiom of choice',
%Bull.\ Acad.\ Polon.\ Sci.\ (Math., Astron., Phys.) 10 (1962) 1-3.

Mycielski J.\ \& \'Swierczkowski S.\ [64]
`On the Lebesgue measurability and the axiom of determinateness',
Fund.\ Math.\ 54 (1964) 67-71.
\cmmnt{[567F.]}

\medskip%N

Naimark M.A. [70] {\it Normed Rings.}  Wolters-Noordhoff, 1970.
\cmmnt{[\S561 {\it notes\/}.]}

Neumann J.von [1931] `Algebraische Repr\"asentanten der Funktionen ``bis
auf eine Menge vom Masse Null''\nobreak', Crelle's J.\ Math.\ 165 (1931)
109-115.
\cmmnt{[535G.]}

Neumann J.von \& Stone M.H.\ [1935] `The determination of representative
elements in the residual classes of a Boolean algebra', Fundamenta
Math.\ 25 (1935) 353-378.
\cmmnt{[\S535 {\it notes\/}.]}

%Normann D.\ [76] `Martin's axiom and medial functions', Math.\
%Scandinavica 38 (1976) 167-176.

%Nyikos P.J.\ [80] `A provisional solution to the normal Moore space
%problem', Proc.\ Amer.\ Math.\ Soc.\ 78 (1980) 429-435.

\medskip%O\medskip%Pa

%Pap E.\ [0?] (ed.) {\it Handbook of Measure Theory.}  North-Holland, 200?.

Pawlikowski J.\ [86] `Why Solovay real produces Cohen real',
J.\ Symbolic Logic 51 (1986) 957-968.
\cmmnt{[552G, 552H.]}

%Pb%Pc%Pd%Pe

Pe{\l}czy'nski A.\ [68] `On Banach spaces containing $L^1(\mu)$, Studia
Math.\ 30 (1968) 231-246.
\cmmnt{[531Ye.]}

Perovi\'c \v{Z}.\ \& Veli\v{c}kovi\'c B.\ [18] `Ranks of Maharam algebras',
Advances in Math.\ 330 (2018) 253-279.
\cmmnt{[539V.]}

%Pf%Pg%Ph%Pu%Pi

% Pierce R.S.\ [58] `A note on complete Boolean algebras',
% Proc.\ Amer.\ Math.\ Soc.\ 9 (1958) 892-896.
% the Dedekind complete case of 515L

%Pj%Pk%Pl

%Plebanek G.\ [93] `Remarks on measurable Boolean algebras and
%sequential cardinals', Fundamenta Math.\ 143 (1993) 11-22.

Plebanek G.\ [95] `On Radon measures on first-countable compact spaces',
Fundamenta Math.\ 148 (1995) 159-164.
\cmmnt{[531O.]}

Plebanek G.\ [97] `Non-separable Radon measures and small compact
spaces', Fundamenta Math.\ 153 (1997) 25-40.
\cmmnt{[531L, 531M, 531V.]}

Plebanek G.\ [00]
`Approximating Radon measures on first-countable compact spaces',
Colloquium Math.\ 86 (2000) 15-23.
\cmmnt{[533H, 533Yc.]}

Plebanek G.\ [02]
`On compact spaces carrying Radon measures of large Maharam type',
Acta Univ.\ Carolinae 43 (2002) 87-99.
\cmmnt{[531U.]}

Plebanek G.\ \& Sobota D.\ [15] `Countable tightness in the spaces of
regular probability measures', Fund.\ Math.\ 229 (2015) 159-169.
\cmmnt{[531U.]}

%Pm%Pn%Po

Pol R.\ [82] `Note on the spaces $P(S)$ of regular probability measures
whose topology is determined by countable subsets', Pacific J.\ Math.\
100 (1982) 185-201.
\cmmnt{[533Ya.]}

%Pp%Pq%Pr

Prikry K.\ [75] `Ideals and powers of cardinals', Bull.\
Amer.\ Math.\ Soc.\ 81 (1975) 907-909.
\cmmnt{[555N.]}

%Ps%Pt%Pu%Pv%Pw%Px%Py%Pz
\medskip%Q

Quickert S.\ [02] `CH and the Saks property', Fundamenta Math.\ 171
(2002) 93-100.
\cmmnt{[539L.]}

\medskip%R

Raisonnier J.\ \& Stern J.\ [85] `The strength of measurability hypotheses', Israel J.\ Math.\ 50 (1985) 337-349.
\cmmnt{[Chap.\ 52 {\it intro.}, 522Q.]}

%Roitman J.\ [79] `Adding a random or a Cohen real;  topological
%consequences and the effect on Martin's axiom', Fundamenta Math.\ 103
%(1979) 47-60.

%Ros{\l}anowski A.\ \& Shelah S.\ [p00] `Measured creatures', preprint, 7.10.00

Rothberger F.\ [1938a] `Eine \"Aquivalenz zwischen der kontinuumhypothese
unter der Existenz der Lusinschen und Sierpinschischen Mengen',
Fundamenta Math.\ 30 (1938) 215-217.
\cmmnt{[522G.]}

Rothberger F.\ [1938b] `Eine Versch\"arfung der Eigenschaft C',
Fundamenta Math.\ 30 (1938) 50-55.
\cmmnt{[\S534 {\it notes\/}.]}

Rothberger F.\ [1941] `Sur les familles ind\'enombrables de suites de nombres naturels et les probl\`emes concernant la propri\'et\'e {\bf C}', Proc.\ Cambridge Phil.\ Soc.\ 37 (1941) 109-126.
\cmmnt{[534R.]}

\medskip%S

Scott D.S.\ [71] (ed.) {\it Axiomatic Set Theory.}  Amer.\
Math.\ Soc., 1971
(Proceedings of Symposia in Pure Mathematics XIII, vol.\ 1).

Shelah S.\ [82] {\it Proper Forcing.}  Springer, 1982
(Lecture Notes in Mathematics 940).
\cmmnt{[\S538 {\it notes\/}.]}

Shelah S.\ [83] `Lifting problem of the measure algebra', Israel J.\
Math.\ 45 (1983) 90-96.   % Sh185
\cmmnt{[\S535 {\it notes\/}.]}

Shelah S.\ [84] `Can you take Solovay inaccessible away?', Israel J.\ Math.\ 48 (1984) 1-47.
\cmmnt{[\S522 {\it notes\/}.]}

Shelah S.\ [92] `Cardinal arithmetic for skeptics',
Bull.\ Amer.\ Math.\ Soc.\ 26 (1992) 197-210.
% Sh400a
\cmmnt{[5A2D.]}

Shelah S.\ [94] {\it Cardinal Arithmetic.}  Oxford U.P., 1994.
% Sh400
\cmmnt{[5A2D, 5A2G.]}

Shelah S.\ [96] `Further cardinal arithmetic', Israel J.\ Math.\ 95
(1996) 61-114.
\cmmnt{[542I.]}
%Sh 430

Shelah S.\ [98a] {\it Proper and Improper Forcing.}  Springer, 1998.
\cmmnt{[\S538 {\it notes\/}.]}
%http://projecteuclid.org/DPubS?service=UI&version=1.0&verb=Display&handle=euclid.pl/1235419814

Shelah S.\ [98b] `There may be no nowhere dense ultrafilter', pp.\ 305-324
in {\smc Makovsky \& Ravve 98}.
\cmmnt{[538H, \S538 {\it notes\/}.]}
% Sh594;  arXiv:math.LO/9611221.

% Shelah S.\ [99] `Special subsets of $^{\cf(\mu)}\mu$, Boolean
% algebras and Maharam measure algebras',
% Topology and its Applications 99 (1999) 135-235.
 % Sh620

Shelah S.\ [00] `Covering of the null ideal may have countable
cofinality', Fundamenta Math.\ 166 (2000) 109-136.
\cmmnt{[\S522 {\it notes\/}.]}
% Sh592

Shelah S.\ [03] `The null ideal restricted to some non-null set may be
$\aleph_1$-saturated',  Fundamenta Math.\ 179 (2003) 97-129.
\cmmnt{[ \S548 {\it notes\/}.]}
% Sh619

%Shelah S.\ \& Fremlin D.H.\ [93] `Pointwise compact and stable sets of
%measurable functions', J.\ Symbolic Logic 58 (1993) 435-455.
  %FrSh 406

Shelah S.\ \& Stepr\=ans J.\ [05]
`Comparing the uniformity invariants of null sets for different measures',
Advances in Math.\ 192 (2005) 403-426.
\cmmnt{[534Za.]}
% Con(\non\Cal N(\mu_H)<\non\Cal N)

%Shelah S.\ \& Woodin H.\ [90] `Large cardinals imply that every
%reasonably definable set of reals is Lebesgue measurable', Israel J.\
%Math.\ 70 (1990) 381-394.

Shipman J.\ [90] `Cardinal conditions for strong Fubini
theorems', Trans.\ Amer.\ Math.\ Soc.\ 321 (1990) 465-481.
\cmmnt{[537K.]}

%Solecki S.\ [06] `Amenability, free subgroups, and Haar null sets in
%non-locally-compact groups',
%Proc.\ London Math.\ Soc.\ 93 (2006) 693-722.

Solecki S.\ \& Todor\v{c}evi\'c S.\ [04]
`Cofinal types of topological directed orders',
Ann.\ Inst.\ Fourier 54 (2004) 1877-1911.
\cmmnt{[513K, 513O, 513Yi.]}

Solecki S.\ \& Todor\v{c}evi\'c S.\ [10]
`Avoiding families and Tukey functions on the nowhere-dense ideal',
J.\ Inst.\ Math.\ Jussieu 9 (2010) 1-31.
\cmmnt{[526L.]}

Solovay R.M.\ [66] `New proof of a theorem of Gaifman and Hales',
Bull.\ Amer.\ Math.\ Soc.\ 72 (1966) 282-284.
\cmmnt{[514Xi.]}

Solovay R.M.\ [70] `A model of set theory in which every set of reals
is Lebesgue measurable', Annals of Math.\ 92 (1970) 1-56.
\cmmnt{[\S522 {\it notes\/}.]}

Solovay R.M.\ [71] `Real-valued measurable cardinals', pp.\ 397-428 in
{\smc Scott 71}.
\cmmnt{[541J, 541P, 541Ya, \S541 {\it notes\/}, 555D, 555O.]}

%Solovay R.M., Reinhardt W.N.\ \& Kanamori A.\ [78]
%`Strong axioms of infinity and elementary embeddings',
%Annals of Mathematical Logic 13 (1978) 73-116.

%Solovay R.M.\ \& Tennenbaum S.\ [71]
%`Iterated Cohen extensions and Souslin's problem',
%Ann.\ Math.\ (2) 94 (1971) 201-245.

%Stepr\=ans J.\ [96] `Cardinal invariants associated with Hausdorff
%capacities', pp.\ 147-184 in {\smc Bartoszy\'nski \& Scheepers 96}.

%Szyma\'nski A.\ [88] `Some remarks on real-valued
%measurable cardinals', Proc.\ Amer.\ Math.\ Soc.\ 104 (1988)
%596-602.

\medskip%T\medskip

Talagrand M.\ [84] {\it Pettis Integral and Measure Theory.}
Mem.\ Amer.\ Math.\ Soc.\  307 (1984).
\cmmnt{[536C.]}

Talagrand M.\ [08] `Maharam's problem', Annals of Math.\ 168 (2008)
981-1009.
\cmmnt{[539A.]}

Tarski A.\ [1945] `Ideale in volst\"andigen Mengenk\"orpen II', Fundamenta Math.\ 33 (1945) 51-65.
\cmmnt{[541P.]}

Todor\v{c}evi\'c S.\ [85] `Remarks on chain conditions in products',
Compositio Math.\ 55 (1985) 295-302.
\cmmnt{[537G.]}

Todor\v{c}evi\'c S.\ [87] `Partitioning pairs of countable ordinals',
Acta Math.\ 159 (1987) 261-294.
\cmmnt{[554Yc.]}

Todor\v{c}evi\'c S.\ [00] `A dichotomy for $P$-ideals of countable sets',
Fundamenta Math.\ 166 (2000) 251-267.
\cmmnt{[5A6G.]}

T\"ornquist A.\ [11] `On the pointwise implementation of near-actions',
Trans.\ Amer.\ Math.\ Soc.\ 363 (2011) 4929-4944.
%{\tt http://arxiv.org/abs/0909.3869}
\cmmnt{[535Yc.]}

Truss J.K.\ [77] `Sets having calibre $\aleph_1$',
pp.\ 595-612 in {\smc Gandy \& Hyland 77}.
\cmmnt{[522J.]}

Truss J.K.\ [88] `Connections between different amoeba algebras',
Fund.\ Math.\ 130 (1988) 137-155.
\cmmnt{[Chap.\ 52 {\it intro.}, 528A, 528Da, 528K, \S528 {\it notes\/}.]}

Tukey J.W.\ [1940] {\it Convergence and Uniformity in Topology.}
Princeton U.P., 1940 (Ann.\ Math.\ Studies 1).
\cmmnt{[513F, \S513 {\it notes}.]}

%U
\medskip

Ulam S.\ [1930] `Zur Masstheorie in der allgemeinen
Mengenlehre', Fundamenta Math.\ 16 (1930) 140-150.
\cmmnt{[\S541 {\it intro.\/}, \S541 {\it notes\/}.]}

%V
\medskip

Veli\v{c}kovi\'c B.\ [92] `Forcing axioms and stationary sets', Advances in
Math.\ 94 (1992) 256-284.
\cmmnt{[517O.]}

Veli\v{c}kovi\'c B.\ [05] `CCC forcing and splitting reals',
Israel J.\ Math.\ 147 (2005) 209-221.
\cmmnt{[\S539 {\it intro.}, 539E, 539N.]}

Vojt\'a\v{s} P.\ [93] `Generalized Galois-Tukey connections between
explicit relations on classical objects of real analysis', pp.\ 619-643
in {\smc Judah 93}.
\cmmnt{[512A.]}

%W
\medskip

Wimmers E.\ [82] `The Shelah $P$-point independence theorem',
Israel J.\ Math.\ 43 (1982) 28-48.
\cmmnt{[\S538 {\it notes\/}.]}

\medskip%X\medskip%Y\medskip%Z

Zakrzewski P.\ [92] `Strong Fubini theorems from measure extension
axioms', Comm.\ Math.\ Univ.\ Carolinae 33 (1992) 291-297.
\cmmnt{[544J.]}

%Zindulka O.\ [00] `Residual measures in locally compact spaces', Topology Appl.\ 108 (2000) 253-265.

%Zindulka O.\ [p98] `Killing residual measures', preprint, 1998.

\discrpage

}%end of references

%---------------------------------------------------------------%

\end



